\documentclass[11pt,a4paper]{article}
\usepackage[T1]{fontenc}
\usepackage[utf8]{inputenc}
\usepackage{lmodern,amsmath,amssymb,amsthm,mathtools,bm,mathrsfs}
\usepackage[margin=27mm]{geometry}
\usepackage{booktabs,longtable,array,enumitem,microtype,calc}
\usepackage[numbers,sort&compress]{natbib}
\providecommand{\printbibliography}{\bibliographystyle{plainnat}\bibliography{Born-v6_59}}
\usepackage{xurl}
\usepackage[colorlinks=true,linkcolor=blue,citecolor=blue,urlcolor=blue]{hyperref}
\usepackage{bookmark}
\usepackage{alphalph}

\makeatletter
\@ifpackageloaded{xr-hyper}{}{%
  \@ifpackageloaded{xr}{}{\usepackage{xr-hyper}}%
}
\makeatother

\ExplSyntaxOn
\prop_new:N \g_bh_external_refs_prop
\cs_new_eq:NN \BHOriginalExternalDocument \externaldocument

\cs_new_protected:Npn \bh_external_document:nnnn #1#2#3#4
 {
  \tl_set:Nn \l_tmpa_tl {#3}
  \regex_replace_once:nnN
   { \A (?: xref/ )? (Paper[0-9]+[A-F]?) .* \Z }
   { \1 }
   \l_tmpa_tl
  \tl_set:Nx \l_tmpb_tl { #1 | \tl_use:N \l_tmpa_tl }
  \prop_if_in:NVTF \g_bh_external_refs_prop \l_tmpb_tl
   {
    \PackageInfo{BornHopf}
     {Skipping~duplicate~external~refs~for~#3}
   }
   {
    \prop_gput:NVn \g_bh_external_refs_prop \l_tmpb_tl {#3}
    \IfNoValueTF{#4}
     { \BHOriginalExternalDocument[#1][#2]{#3} }
     { \BHOriginalExternalDocument[#1][#2]{#3}[#4] }
   }
 }

\RenewDocumentCommand{\externaldocument}{O{} O{} m o}
 { \bh_external_document:nnnn {#1}{#2}{#3}{#4} }

\NewDocumentCommand{\BHExternalRefs}{m}
 {
  \tl_set:Nx \l_tmpa_tl { \c_sys_jobname_str }
  \regex_replace_once:nnN
   { \A (Paper[0-9]+[A-F]?) .* \Z }
   { \1 }
   \l_tmpa_tl
  \tl_set:Nn \l_tmpb_tl { #1 }
  \regex_replace_once:nnN
   { \A (Paper[0-9]+[A-F]?) .* \Z }
   { \1 }
   \l_tmpb_tl
  \str_if_eq:eeTF { \tl_use:N \l_tmpa_tl } { \tl_use:N \l_tmpb_tl }
   {
    \PackageInfo{BornHopf}
     {Skipping~external~refs~for~#1~because~this~is~the~current~document}
   }
   {
    \IfFileExists{xref/#1.aux}
     { \externaldocument[][nocite]{xref/#1}[#1.pdf] }
     {
      \IfFileExists{#1.aux}
       { \externaldocument[][nocite]{#1} }
       {
        \PackageInfo{BornHopf}
         {External~aux~file~#1.aux~not~found;~no~labels~imported}
       }
     }
   }
 }
\ExplSyntaxOff
\BHExternalRefs{Paper1A-Geometry-v5_180}
\BHExternalRefs{Paper1B-ExceptionalGeometry-v1_22}
\BHExternalRefs{Paper1C-M8Geometry-v1_144}

\providecommand{\tightlist}{\setlength{\itemsep}{0pt}\setlength{\parskip}{0pt}}
\setlength{\emergencystretch}{3em}
\setcounter{tocdepth}{2}
\numberwithin{equation}{section}
\newtheorem{theorem}{Theorem}[section]
\newtheorem{proposition}[theorem]{Proposition}
\newtheorem{lemma}[theorem]{Lemma}
\newtheorem{corollary}[theorem]{Corollary}
\theoremstyle{definition}
\newtheorem{definition}[theorem]{Definition}
\theoremstyle{remark}
\newtheorem{remark}[theorem]{Remark}
\newcommand{\R}{\mathbb R}
\newcommand{\C}{\mathbb C}
\newcommand{\Z}{\mathbb Z}
\newcommand{\Q}{\mathbb Q}
\newcommand{\HH}{\mathbb H}
\newcommand{\F}{\mathbb F}
\newcommand{\T}{\mathbb T}
\newcommand{\SU}{\mathrm{SU}}
\newcommand{\SO}{\mathrm{SO}}
\newcommand{\U}{\mathrm U}
\newcommand{\Spin}{\mathrm{Spin}}
\newcommand{\Cl}{\mathrm{Cl}}
\newcommand{\id}{\mathrm{id}}
\newcommand{\Ad}{\operatorname{Ad}}
\newcommand{\End}{\operatorname{End}}
\newcommand{\Hom}{\operatorname{Hom}}
\newcommand{\tr}{\operatorname{tr}}
\newcommand{\Tr}{\operatorname{Tr}}
\newcommand{\rank}{\operatorname{rank}}
\newcommand{\im}{\operatorname{im}}
\newcommand{\diag}{\operatorname{diag}}
\newcommand{\Spec}{\operatorname{Spec}}
\newcommand{\Sym}{\operatorname{Sym}}
\newcommand{\vol}{\operatorname{vol}}
\newcommand{\sgn}{\operatorname{sgn}}
\newcommand{\supp}{\operatorname{supp}}
\newcommand{\osc}{\operatorname{osc}}
\newcommand{\cA}{\mathcal A}
\newcommand{\cE}{\mathcal E}
\newcommand{\cP}{\mathcal P}
\newcommand{\cH}{\mathcal H}
\newcommand{\dd}{\mathrm d}
\newcommand{\Stirling}[2]{\genfrac{\{}{\}}{0pt}{}{#1}{#2}}
\title{\textbf{Born--Hopf Geometry I-D:}\\[0.4em]\large Technical proofs and computational certificates}
\author{Peter Kugelmann}
\date{Version 1.69 --- 20 November 2026}
\begin{document}
\providecommand{\RR}{\mathbb R}
\providecommand{\ZZ}{\mathbb Z}
\providecommand{\Sone}{S^1}
\providecommand{\TableNotes}[1]{\par\smallskip\footnotesize #1\par}
\providecommand{\CC}{\mathbb{C}}
\providecommand{\Rmom}{R_{\mathrm{mom}}}

\maketitle
\begin{abstract}
This volume supplies the technical proofs, explicit calculations and finite
certificates supporting the intrinsic $S^3$ compendium (Paper~I-A), the
exceptional-geometry constructions (Paper~I-B), and the reciprocal
$M_8$ geometry (Paper~I-C). It develops the mixed rank-two
classification and natural-boundary proof chains on their stated
bundle, gauge and relative-frame hypotheses; includes the full
restriction, obstruction, extension and fibre calculations; and records
analytic, spectral, geometric-functional and finite comparison results.
Mathematical claims retain their declared domains and normalizations.
Computational certificates do not establish physical occurrence or
replace an unproved global classification.  The source-specific assumptions and comparison maps remain explicit
throughout.
\end{abstract}
\tableofcontents
\clearpage

\noindent\textbf{Organization.}
The first two sections develop the complete rank-two gauge-image and
natural relative-boundary proof chains. The following sections provide
analytic, finite, metric, bundle, Dolbeault and geometric-functional
calculations and explicit classification ledgers. Subject-organized
appendices contain extended comparison proofs and finite certificates.
All restrictions to a selected source, finite test family, specified
framing or comparison map are part of the corresponding mathematical
statements.
\medskip


\section{Mixed rank-two gauge images and complete classification}
\label{p1c:sec:rank-two-full-proofs}
The final classification statement is given in Paper~I-C,
Theorem~\ref{p1b:thm:complete-mixed-rank-two}; the following is the
complete positive proof chain with the original labels and hypotheses.

\subsection{The evaluated mixed rank-two classification}
\label{p1b:sec:rank-two-closure}
Retain the actual mixed local representation and the four compatible
free-region bundles of Section~\ref{p1b:sec:v2170-mixed}. Put
\(z=z_6\), so that \(T=C_6\{z\}\oplus C_2\{e\}\oplus C_2\{t\}\)
and \(a=3z\). The coefficient \(A\) is the second-Chern coefficient
in \(c_2=A\mathcal D\). Parity below means parity of \(A\).
The scalar matched top bit is written
\(\zeta_{\rm nat}=w_t+h_t\in C_2\); it is distinct from the local
circle class \(\zeta=2\kappa\) in the earlier boundary chart.
The following source-bound proof steps retain symbolic coefficients
until their complete natural-map evaluations. The resulting evaluated
classification is stated in Theorem~\ref{p1b:thm:complete-mixed-rank-two}
and the natural tests in Theorem~\ref{p1b:thm:natural-rank-two-completion}.

\subsection{Finite restriction image and comparison obstructions}
\label{p1b:sec:rtproof29}

Retain the actual bundles and boundary identification of
Section~\ref{p1b:sec:v2170-mixed}, including the even free-primary theorem,
the local boundary normalization, and Theorems~\ref{p1b:rt:theorem123}
and~\ref{p1b:rt:theorem127}.  Write \(\eta=\eta_W\).
All new statements below are on those source-bound inputs.

\begin{lemma}
\label{p1b:rt:lemma128}
The actual integral transfer classes satisfy
\[
 hq=2c_4+\tau^2,\qquad \bar h\,\bar q=\chi^4.
\]
\end{lemma}
\begin{proof}
In \(B=M_{hC_2}\), set \(c=\operatorname{tr}(a_tu_x)\) and
\(c'=\operatorname{tr}(a_tu_p)\). The transfer projection formula gives
\(h_Bq_B=c+c'\). Integral degree four consists of the two primitive free
transfer directions and \(C_2\{t^2\}\). Permutation modules in product
degrees one and three and the sign module in degree two give no other
positive-filtration terms; the fixed-point section retains \(t^2\).
The underlying diagonal Thom restriction is \(c-c'\), and its fixed
restriction is \(t^2\), whereas \(c-c'\) has fixed restriction zero.
Consequently \(\delta=c-c'+t^2=2c-h_Bq_B+t^2\).
Restrict to the free complement, where \(\delta=0\), and reduce modulo
two.
\end{proof}

\begin{theorem}
\label{p1b:rt:theorem129}
For every compatible actual bundle,
\[
 u_{\chi,V}=0,\qquad
 J_{0,V}=\langle j_{h,V},j_{\eta,V},u_{c,V}\rangle,\qquad
 J_{0,V}^{\det=1}=\langle j_{\eta,V},u_{c,V}\rangle .
\]
\end{theorem}
\begin{proof}
Use the actual \(g_h\) and the determinant-one primary-\(q\) section through
degree six. The next obstruction has coefficients \(C_{12}\).
Its boundary restriction vanishes because
\(H^7(W;C_{12})\to H^7(R;C_{12})\) is zero for the oriented connected
pair. Adjusting six-cell choices makes the partial section defined on
\(W^{(6)}\cup R\). Form \(k=g_hs_qg_h^{-1}s_q^{-1}\) there.
Its primary is zero and its fourth coefficient is the Bott Samelson
value \(\bar h\bar q=\chi^4\).

Here the relative commutator extension must be checked, rather than
discarding a relative top obstruction. Pointwise, the based map
\(u\mapsto g_hug_h^{-1}u^{-1}\) is nullhomotopic. Thus its application
kills the seven-cell defect of the partial section. For coherence on
eight-cells, the only possible contribution is a one-parameter family
of these nullhomotopies applied to that defect in \(\pi_6SU(2)\).
Pointwise contributions on \(\pi_7SU(2)\) vanish; parameter dimensions
at least two first contribute in dimension nine, and the \(S^3\) part
of \(U(2)\) first contributes in dimension ten. Bundle-frame
variation begins in parameter dimension two: frames exist on the
one-skeleton and the structure group is connected.
The one-parameter family is the determinant-loop conjugation
operation \(A\), with coefficient
\(A_*x=\eta_3Ex\). Lemma~\ref{p1b:rt:lemma104} gives
\(A_*\nu'=0\), hence zero for every \(C_{12}\) defect.
The relative eight-cell obstruction is therefore zero. This extends
\(k\) with its specified boundary value.

Both boundary factors have zero normal degree and even determinant
circle winding. Normalization on this subgroup is a homomorphism to
the additive \(\mathscr A\), so the normalized commutator is zero.
Theorem~\ref{p1b:rt:theorem123} makes the boundary value independent of
all later choices at the fixed fourth class. This computes
\(u_{\chi,V}\). The remaining equalities follow from
Theorem~\ref{p1b:rt:theorem127}.
\end{proof}

\begin{theorem}
\label{p1b:rt:theorem130}
In a fixed parity and restriction sector, the two fifth-bit bundles
have equal full and determinant-one boundary gauge images, and equal
degree-two sphere relations modulo that image.
\end{theorem}
\begin{proof}
Use the closed complementary sphere \(Z_p\) of the actual fifth-bit
patch, away from \(R\) and \(P\). Its neighborhood is \(D^5\times S^3\)
and every rank-two complex bundle is trivial there. Reglue it by
\(\eta_3:S^4\to SU(2)\), independent of the complementary sphere.
Its first bundle difference is the nonzero fifth class dual to \(Z_p\);
it retains Chern data, boundary identification and the \(P\) sector,
so the four-class theorem identifies the other fifth-bit bundle.

An actual gauge restricts to \(n\iota_3\) on \(Z_p\). Torsion primary
restricts to zero there, and the actual mapping-torus parity theorem
gives \(n=2m\). After a gauge homotopy supported off \(R\), take a based
sphere map constant in the normal directions. Its compatibility
defect with the new clutching is the smash commutator
\[
 \langle\eta_3,2m\iota_3\rangle
 =2m\langle\eta_3,\iota_3\rangle=0
 \quad\text{in }\pi_7U(2)=C_2 .
\]
Choose a based nullhomotopy and extend over the reglued neighborhood.
This preserves the exact boundary gauge. For determinant-one gauges
all maps and homotopies can be taken in \(SU(2)\). Inverse surgery gives
the reverse inclusion. Lower data through degree four can be fixed
on a skeleton disjoint from \(Z_p\).
In particular the degree-two sphere gauge transfers with its boundary
value, proving the relation comparison.
\end{proof}

\begin{corollary}
\label{p1b:rt:corollary131}
At Chern zero every finite restriction column has \(e\)-coordinate zero,
and
\[
 s_{00}=s_{01}=s_0\le1,\qquad s_{10}=s_{11}=s_1\le2,\qquad
 f_{\varepsilon1}=2f_{\varepsilon0}.
\]
Thus
\[
 N_{\rm mix}(0)=96(2^{-s_0}+2^{-s_1})
 \in\{72,96,120,144,192\}.
\]
This is an intermediate image formula; the following even calculation gives \(s_0=1,s_1=2\) and total72. The explicit split
odd-determinant redundancy transfers to its fifth-bit mate, leaving
three distinct parity/sector membership tests.
\end{corollary}
\begin{proof}
Intersect the paired kernel planes in
Theorem~\ref{p1b:thm:v2170-zero-fibres}: the intersections are
\(\langle t\rangle\) and \(\langle a,t\rangle\).
Use that theorem's uniform fibre formula. Equivalently each sector
has two equal free-region quotient orders \(48/2^{s_\varepsilon}\).
The inherited split representatives have \(j_h=j_\eta=0\), and
Theorem~\ref{p1b:rt:theorem130} preserves them.
\end{proof}

For a finite ordinary extension computation, divide the zero-sphere
even-circle subgroup by the two actual global restrictions
\(4q_+-2\kappa\) and \(2\kappa\). Its representative set is
\[
 2d q_++p\kappa+v,\qquad d,p\in\mathbb F_2,\quad v\in T[2].
\]
It is \(\mathbb F_2^5\). Let \(E_V\) be its actual extendable subgroup.
Then
\[
 J_{0,V}=\operatorname{pr}_T E_V,\qquad
 J_{0,V}^{\det=1}=\operatorname{pr}_T(E_V\cap\{d=0\}).
\]
The second equality needs no general fibrewise algorithm:
\(H^1(W;\mathbb Z)\to H^1(R;\mathbb Z)\) sends \(h\) to \(2a_F\)
and is injective. An ordinary extension of a determinant-one boundary
representative has zero determinant winding. Its logarithm can be
chosen zero on the connected boundary; division by its scalar half
makes the extension determinant one while fixing that boundary.

The finite classifying-map formulation also gives a supplementary decidability comparison; Section~\ref{p1b:app:rank-two-decidability} states its finite target and rational H-space hypotheses. The exact images here are evaluated by the subsequent actual-map calculations.



\subsection{The remaining fourth gauge from a punctured factor}
\label{p1b:sec:rtproof30}

Retain \(K=S^1\times S^3\), the unordered free region \(W\), its
boundary \(R=K\times L\), \(L=\mathbb {RP}^3\), and
\(E_\partial=\mathbf1\oplus\lambda_L\). Write \(\rho\) for the
mixed involution. Identifying \(SU(2)\) with the unit quaternions,
conjugation by \(\rho\) is conjugation by \(i\). We use the actual
boundary group from the supplied sources:
\[
 \pi_1H=\mathbb Z\{q_+\}\oplus C_4\{\kappa\},\qquad
 \pi_4H=C_2\{e\}\oplus C_2\{t\}.
\]
Every circle loop is represented by local eigenline windings. In
particular, the kernel of determinant winding in \(\pi_1H\) is
\(C_4\{\kappa\}\), represented by loops in the diagonal
determinant-one circle. The determinant fibration has
\(\pi_2\operatorname {Map}(L,U(1))=0\); thus a determinant-one loop
in this kernel can be homotoped to its local representative through
determinant-one loops.

The retained unstable input is
\[
 \langle\iota_3,\eta_3\rangle
 =\langle\iota_3,\iota_3\rangle\eta_6
 =\mu\ne0\quad\hbox{in }\pi_7SU(2)=C_2.
 \tag{RT.30.1}\label{p1b:eq:rt-30-1}
\]
Here \(\langle\iota_3,\iota_3\rangle\) is a generator of
\(\pi_6SU(2)=C_{12}\). This is the actual G20 Samelson coefficient,
not the different vanishing conjugation coefficient
\(\eta_3E\nu'\). The natural top insertion of \(\mu\) is the
nonzero element \(t\) of the displayed \(\pi_4H\).

\begin{lemma}[A proper representative of the transfer class]
\label{p1b:rt:lemma30-transfer}
Fix \(k_0\in K\). The configurations with one point equal to
\(k_0\) form a properly embedded four-manifold
\[
 C\cong K\setminus\operatorname {int}D^4\subset W.
\]
Its rank-four normal bundle is trivial, and its mod-two Poincare--
Lefschetz dual is \(\bar c_4\). Every compatible rank-two bundle
\(V\) is trivial on a tubular neighborhood of \(C\).
\end{lemma}
\begin{proof}
The fixed particle distinguishes an ordering on this locus. In that
ordering, moving the fixed particle identifies the normal bundle with
the constant space \(T_{k_0}K\). At the boundary, the other particle
can move to maintain the chosen separation, so this normal
identification remains transverse to \(R\). The ordered dual is the
sum of the two factor orientation classes; descent is exactly the
transfer \(\operatorname {tr}(a_tu_x)=c_4\).

Use the product cell structure of \(K\), with cells in dimensions
\(0,1,3,4\), and remove a ball in its open four-cell. The punctured
four-cell retracts onto its attaching boundary; consequently
\(C\simeq S^1\vee S^3\). Since \(\pi_1BU(2)=\pi_3BU(2)=0\),
every complex rank-two bundle on \(C\) is trivial. The tubular
neighborhood retracts onto \(C\), proving the last assertion.
\end{proof}

Choose a unitary frame on that neighborhood. On each normal four-disc
insert a based representative of \(\eta_3:S^4\to SU(2)\), with
the identity at the disc boundary, and use the identity elsewhere.
This gives an actual global determinant-one gauge \(g_C\).
Its primary is zero and its fourth difference is \(\bar c_4\).
The supplied stage-four-kernel theorem makes its entire boundary
component equal to that of any other actual realization of these
lower data. It therefore computes the specified column \(u_{c,V}\).

\begin{lemma}[The constant-frame boundary is the top insertion]
\label{p1b:rt:lemma30-local}
For \(V_s=\mathbf1\oplus\lambda_W\), use its constant flat frame
in the ordering distinguished by the fixed particle. Then
\[
 {}[g_C|_R]=t\in\pi_4H.
 \tag{RT.30.2}\label{p1b:eq:rt-30-2}
\]
The boundary has zero circle and sphere coordinates.
\end{lemma}
\begin{proof}
Take an oriented Euclidean coordinate chart at \(k_0\), identified
with \(\mathbb H\). The diagonal tube can be chosen flat on this
chart. In its ordered boundary coordinates the two particles are
\(z\pm r u\), where \(z\in\mathbb H\), \(u\in S^3\), and
\(r>0\). Exchange sends \(u\) to \(-u\) and acts on the gauge
fibre by \(\rho\). The proper boundary of \(C\) is the graph
\[
 i_C:S^3\longrightarrow D^4\times L,\qquad
 v\longmapsto(rv,[v]).
 \tag{RT.30.3}\label{p1b:eq:rt-30-3}
\]
It has a product normal tube
\((v,n)\mapsto(rv+n,[v])\), \(|n|<\delta<r\).
The two centers in each lens fibre are disjoint.

We can choose the suspended Hopf representative explicitly. Set
\[
 f(n)=\cos\theta(|n|)+\sin\theta(|n|)\,
          \frac{n i\bar n}{|n|^2},\qquad
 \theta(0)=\pi,\quad\theta(\delta)=0,
 \tag{RT.30.4}\label{p1b:eq:rt-30-4}
\]
where \(\theta\) decreases continuously, put \(f(0)=-1\), and
put \(f=1\) outside the disc. This represents \(\eta_3\), and
\[
 f(-n)=f(n),\qquad f(v^{-1}n)=v^{-1}f(n)v\quad(v\in S^3).
 \tag{RT.30.5}\label{p1b:eq:rt-30-5}
\]
In the ordered boundary frame let \(B_0(z,u)\) be \(f(z-ru)\)
near \(ru\), \(\rho f(z+ru)\rho^{-1}\) near \(-ru\), and
the identity outside those two discs. This is the actual constant-frame
restriction. Its support is in a small four-disc of the fixed locus,
so compactifying \(z\) makes it a based \(S^4\)-family in \(H\).

For a comparison family put
\[
 \psi(w)=
 \begin{cases}
 f(w-r),& |w-r|<\delta,\\
 \rho f(-w-r)\rho^{-1},& |w+r|<\delta,\\
 1,&\text{otherwise},
 \end{cases}
 \qquad B_1(z,u)=\psi(u^{-1}z).
 \tag{RT.30.6}\label{p1b:eq:rt-30-6}
\]
The real quaternion \(r\) is used in this formula.
One has \(\psi(-w)=\rho\psi(w)\rho^{-1}\), and both \(w=0\)
and \(w=\infty\) give the identity. Thus radial restriction of
\(\psi\) is a based determinant-one loop in the lens gauge group.
By the actual circle calculation it is homotopic to a local diagonal
loop \(\ell(|w|)\). Interpret the homotopy as equivariant maps of
the radial compactification of \(\mathbb H\), fixing both poles,
and substitute \(w=u^{-1}z\) throughout. It gives a homotopy of
\(B_1\) to \(\ell(|z|)\), independent of \(u\).
As a based map from the compactified \(z\)-space \(S^4\), this
last map factors through the diagonal circle and is nullhomotopic.
Only the pole at infinity must stay fixed in this final nullhomotopy.
Therefore
\[
 {}[B_1]=0\quad\hbox{in }\pi_4H.
 \tag{RT.30.7}\label{p1b:eq:rt-30-7}
\]

The comparison does not erase the normal-frame change. On the product
tube \eqref{p1b:eq:rt-30-3}, equations \eqref{p1b:eq:rt-30-5}--\eqref{p1b:eq:rt-30-6} give exactly
\[
 B_1(v,n)=v^{-1}f(n)v,\qquad B_0(v,n)=f(n).
\]
Their pointwise quotient is the identity at \(v=1\) and on the
normal-disc boundary. It factors through \(S^3\wedge S^4=S^7\)
and represents \(\langle\iota_3,\eta_3\rangle=\mu\), up to
an immaterial sign. The trivial tube frame permits its homotopy into a
small three-disc of the parameter sphere and a normal four-disc.
After inclusion in \(R\), this is exactly one seven-disc insertion
of \(\mu\), hence the natural element \(t\).
There is no factor two: the parameter is the single embedded
\(S^3=\partial C\), even though its projection to \(L\) is a
double cover. Equivalently, collapse to the tube's top cell has degree
one. Thus \([B_1]-[B_0]=t\). Equations \eqref{p1b:eq:rt-30-7} and \(2t=0\)
give \([B_0]=t\), as required.
\end{proof}

\begin{theorem}[The actual fourth column for every compatible bundle]
\label{p1b:rt:theorem30-column}
For any compatible actual \(V\) at Chern coefficient \(A\),
in either restriction sector and either fifth-bit class,
\[
 \boxed{u_{c,V}=(1+(A\bmod2))t.}
 \tag{RT.30.8}\label{p1b:eq:rt-30-8}
\]
The representative has zero circle, sphere, \(a\), and \(e\)
coordinates. In particular the pure \(t\) boundary gauge extends
over \(V\) whenever \(A\) is even.
\end{theorem}
\begin{proof}
Use the frame of Lemma~\ref{p1b:rt:lemma30-transfer}. In the selected
boundary ordering, compare it with the canonical local mixed frame.
The comparison is a map \(b:S^3\to U(2)\); dependence in the
normal-disc direction can be homotoped away. Its determinant has a
logarithm on \(S^3\), so its \(\pi_3\) class can be represented
by a based \(SU(2)\) map. Write \([b]=N\iota_3\).

Determine \(N\) using the actual coherent gluing, rather than a
chosen free-region sector. The one-fixed-particle slice of the ordered
space is \(\{k_0\}\times K\). Its part outside the diagonal tube
is \(C\), and the missing four-ball has the local mixed frame.
The comparison \(b\) is precisely the clutching map on this slice.
The underlying coherent bundle has \(c_2=A\mathcal D\), and
\(\mathcal D|_{\{k_0\}\times K}=a_Eu_p\). Therefore
\[
 N=\pm A,\qquad N\bmod2=A\bmod2.
 \tag{RT.30.9}\label{p1b:eq:rt-30-9}
\]
Changing a boundary reference in the same prescribed Chern component
does not alter this parity. This argument uses only the actual Chern
class, so the restriction sector and fifth-bit surgery add no hidden
term to \eqref{p1b:eq:rt-30-9}.

In the boundary reference frame the supported gauge is the conjugate
of the constant-frame profile by \(b(v)\). Its comparison with that
profile is supported on the same product tube. Since the primary
class on each normal disc is \(\eta_3\), the difference is
\[
 \langle N\iota_3,\eta_3\rangle=N\mu.
\]
As in the preceding lemma, this becomes \(Nt\) on \(R\).
Add this to the constant-frame value \(t\); \eqref{p1b:eq:rt-30-9} gives
\eqref{p1b:eq:rt-30-8}. The supported global construction already extends over
all of \(W\), so no vanishing absolute top group is substituted
for a relative extension calculation.
\end{proof}

\begin{corollary}[Two remaining Chern-zero totals]
\label{p1b:rt:corollary30-counts}
On the supplied sources and the paired-image and \(u_\chi=0\)
results of Subsection~\ref{p1b:sec:rtproof29}, the split pair has \(J_0=\langle t\rangle\)
and rank \(s_0=1\). At even parity, both sector images contain
\(t\), and
\[
 s_1\in\{1,2\},\qquad
 \boxed{N_{\rm mix}(0)\in\{72,96\}.}
 \tag{RT.30.10}\label{p1b:eq:rt-30-10}
\]
The Chern-zero fibre multiplicities in \((00,01,10,11)\) order are
\((1,2,2,4)\) if \(s_1=1\), and \((1,2,1,2)\) if \(s_1=2\).
Each split-pair free-region class has 24 mixed gluings; using its
retained nonzero sphere relation, its additive quotient is
\(C_{12}\oplus C_2\). At odd parity the pure fourth-stage image
is zero, so each full finite image has rank at most two and each
determinant-one finite image has rank at most one. The necessary odd
totals reduce to \(\{48,72,96,120,144,192\}\).
\end{corollary}
\begin{proof}
The explicit split low columns are zero; \eqref{p1b:eq:rt-30-8} and Subsection~\ref{p1b:sec:rtproof29}
give its finite image. The even sector-one image is contained in
\(\langle a,t\rangle\), and now contains \(t\), proving its
rank alternatives. Substitute in
\(96(2^{-s_0}+2^{-s_1})\) and the retained uniform fibre formula.
The split quotient order and structure follow from its rank one and
its nonzero degree-two sphere correction. At odd parity
\(u_\chi=u_c=0\), leaving only the two low columns, or just
\(j_\eta\) in determinant one. Enumerating ranks \(0,1,2\)
in the paired total formula gives the stated necessary set.
\end{proof}

This fourth-gauge calculation determines the specified column.
The subsequent low-column and sphere proofs select the even exotic
rank and evaluate the general-Chern forgetful fibres. Their combined
classification is Theorem~\ref{p1b:thm:complete-mixed-rank-two}.



\subsection{The determinant-sensitive even restriction image}
\label{p1b:sec:rtproof31}

Retain the supplied actual rank-two classification on
\(M=K^2\), \(K=S^1\times S^3\), with exchange free quotient
\(W\), ordered free region \(U\), and boundary
\(R=K\times L\), \(L=\mathbb {RP}^3\).
Write \(p:U\to W\) for the ordered cover. The boundary model is
\(E_\partial=\mathbf1\oplus\lambda_L\), with gauge group \(H\).
The normalized boundary difference group is
\[
 \mathscr A=\mathbb Z\{b\}\oplus T,\qquad
 T=C_6\{z_6\}\oplus C_2\{e\}\oplus C_2\{t\},
 \qquad a=3z_6.
 \tag{RT.31.1}\label{p1b:eq:rt-31-1}
\]
Every actual zero-sphere global restriction belongs to
\(T[2]=\langle a,e,t\rangle\). The actual low gauge generators
have data \((h,0)\), \((0,\eta_W)\), and \((0,2q)\), with the
last one determinant one. The main supplements of Subsections~\ref{p1b:sec:rtproof29}--\ref{p1b:sec:rtproof30} give
\[
 u_\chi=0,\qquad u_c=(1+(A\bmod2))t,
 \tag{RT.31.2}\label{p1b:eq:rt-31-2}
\]
and equality of the full and determinant-one boundary images, and
sphere relations, between the two fifth-bit bundles in a fixed parity
and restriction sector. These are load-bearing inputs here.

At \(A=0\), denote the four free-region bundles by
\(V_{\varepsilon,\delta}\). Intersecting the supplied ordinary
forgetful kernels of the fifth-bit pair gives
\[
 J_{0,V}\subset
 \begin{cases}\langle t\rangle,&\varepsilon=0,\\
 \langle a,t\rangle,&\varepsilon=1.
 \end{cases}
 \tag{RT.31.3}\label{p1b:eq:rt-31-3}
\]
Thus every such finite restriction has zero \(e\)-coordinate.

\subsubsection{Primary differences before determinant forgetting}

The following refinement retains the determinant trivialization on the
underlying ordinary bundle. It does not identify this category with
principal bundles carrying a coherent \(SU(2)\) exchange lift.

\begin{lemma}[The retained ordinary primary chart]
\label{p1b:rt:lemma31-primary}
At \(c_1=c_2=0\), primary differences of ordinary bundles with a
chosen determinant trivialization take values in \(C_{12}^2\).
A normalized finite boundary difference \((z,e,t)\) changes them by
\[
 (\Delta z_t,\Delta z_E)=(2z,-2z)\pmod {12}.
 \tag{RT.31.4}\label{p1b:eq:rt-31-4}
\]
Changing the determinant trivialization by a function of winding
\((k_t,k_E)\in H^1(M;\mathbb Z)\) changes the primary pair by
\[
 6\varepsilon(k_t,k_E)\pmod {12}.
 \tag{RT.31.5}\label{p1b:eq:rt-31-5}
\]
\end{lemma}
\begin{proof}
Use the actual G20 torus mapping chart retained in the supplied ordinary
owner. At zero Chern data both torus-edge components are central in
the order-twelve Heisenberg group. Conjugation therefore leaves each
central coordinate fixed. The slice-change matrix \(N_c\) is zero,
so the primary differences remain in \(C_{12}^2\) before determinant
forgetting, in both restriction sectors.

The natural ordered-cover map sends the sphere torsion insertion to
coefficient \(2z\). The two torus-edge Gysin orientations give
\((2z,-2z)\). The cover of \(e\) is a pure top filling difference,
and that of \(t\) is zero, so neither changes an edge coordinate.
This is the finite-generator calculation of the supplied Chern-zero
forgetful proof, retained before reduction modulo six in the exotic
sector. No choice of a secondary origin is required.

The determinant boundary in the actual restriction component sends
one winding to the central element \(\kappa_P^{6\varepsilon}\).
Equivalently, its coefficient is
\(\eta_3\eta_4\eta_5=6\nu'\) on the exotic smash cell.
Applying the two torus determinant operations gives \eqref{p1b:eq:rt-31-5}, exactly
the primary action in the supplied ordinary owner. Any affine effect
on a face filling is outside this primary comparison.
\end{proof}

\begin{lemma}[The determinant of the gluing isomorphism]
\label{p1b:rt:lemma31-gluing}
Fix a Chern-zero reference coherent gluing of \(V\). If an actual
global gauge has zero sphere coordinate and determinant class \(kh\),
let \(j\in T[2]\) be its normalized boundary value. The resulting
ordinary isomorphism between the reference gluing and its modification
by \(j\) has determinant winding
\[
 (k,k)\in H^1(M;\mathbb Z).
 \tag{RT.31.6}\label{p1b:eq:rt-31-6}
\]
For a determinant-one gauge the isomorphism can be made determinant
preserving.
\end{lemma}
\begin{proof}
The ordered pullback has determinant class
\(p^*(kh)=k(a_t+a_E)\), by the transfer identity for \(h\).
Its boundary circle determinant winding is \(2k\). Every circle
class in \(\pi_1H=\mathbb Z\{q_+\}\oplus C_4\{\kappa\}\)
has a representative from local eigenline windings \((m,n)\), whose
determinant winding is \(m+n\). Choose this local gauge with
\(m+n=2k\) to remove the circle part. The residual finite boundary
gauge can be represented with determinant identically one. A collar
homotopy makes the determinant functions match exactly on the seam.

The global gauge on \(U\) and the chosen local gauge on \(N\)
then define an ordinary isomorphism between the two glued bundles.
Since the residual finite gluing has determinant one, their determinant
line bundles and reference trivializations may be identified. The
determinant of this isomorphism is consequently a function on \(M\).
Removing a codimension-four diagonal does not change \(H^1\):
excision gives \(H^1(M,U)=H^2(M,U)=0\), hence
\(H^1(M)\to H^1(U)\) is an isomorphism. Its restriction is the
already computed \(k(a_t+a_E)\), which proves \eqref{p1b:eq:rt-31-6}.

For a determinant-one global gauge, the local circle removal can use
\(m+n=0\). The pasted determinant function has winding zero and
therefore a global logarithm. Multiplying the isomorphism by the scalar
half of its inverse makes its determinant identically one. This changes
neither bundle nor primary difference.
\end{proof}

\begin{theorem}[The two even low columns]
\label{p1b:rt:theorem31-columns}
Modulo the known extendable class \(t\), the Chern-zero columns are
\[
 \boxed{j_{h,V}=\varepsilon a\pmod {\langle t\rangle},\qquad
 j_{\eta,V}=0\pmod {\langle t\rangle}.}
 \tag{RT.31.7}\label{p1b:eq:rt-31-7}
\]
In particular the two even exotic ordinary-extension tests are
\[
 \boxed{2q_++\gamma(a):\ \mathrm{YES},\qquad
 \kappa+\gamma(a):\ \mathrm{NO}.}
 \tag{RT.31.8}\label{p1b:eq:rt-31-8}
\]
Thus \((x_h,x_\eta)=(1,0)\).
\end{theorem}
\begin{proof}
Let \(z\in C_6\) be the sphere-torsion coordinate of any actual
zero-sphere gauge with determinant class \(kh\). Comparing
\eqref{p1b:eq:rt-31-4} with the primary determinant change of the isomorphism in
\eqref{p1b:eq:rt-31-6} gives
\[
 (2z,-2z)=(6\varepsilon k,6\varepsilon k)\pmod {12}.
 \tag{RT.31.9}\label{p1b:eq:rt-31-9}
\]
Because the gauge's finite restriction is in \(T[2]\), \(z\) is
zero or three; the equation forces \(z=3\varepsilon k\).
Apply this first to \(g_h\), where \(k=1\), and then to the
determinant-one \(g_\eta\), where \(k=0\). Their \(e\)-coordinates
vanish by \eqref{p1b:eq:rt-31-3}. This proves \eqref{p1b:eq:rt-31-7}.

The actual \(u_c\) gauge extends \(t\), so its multiplication removes
any residual \(t\)-coordinate of \(g_h\). The actual fourth
\(\alpha\) gauge has whole boundary circle \(2\kappa\) and zero
normalized finite value; it removes the even \(\kappa\)-ambiguity
in the pure-primary-zero circle of \(g_h\). Therefore the specified
representative \(2q_++\gamma(a)\) actually extends in the exotic
sector.

If \(\kappa+\gamma(a)\) had an ordinary extension, its boundary
determinant is one. Injectivity of \(H^1(W)\to H^1(R)\),
\(h\mapsto2a_F\), and the relative logarithmic scalar adjustment
from Subsection~\ref{p1b:sec:rtproof29} would make the extension determinant one. Its finite
\(a\)-coordinate would then violate \eqref{p1b:eq:rt-31-9} with \(k=0\).
This proves the NO statement by the same load-bearing comparison.
\end{proof}

\subsubsection{Exact even counts and determinant action}

Write \(J_{0,V}^{\mathrm{det}=1}\) for the zero-sphere image of
the determinant-one global gauge subgroup. The local mixed gauge
group is retained when forming its auxiliary gluing quotient.

\begin{theorem}[The even finite images and mixed cardinality]
\label{p1b:rt:theorem31-even}
For every even \(A\), both fifth-bit classes in each sector have
\[
 \boxed{J_{0,V}^{\mathrm{det}=1}=\langle t\rangle,\qquad
 J_{0,V}=\begin{cases}\langle t\rangle,&\varepsilon=0,\\
 \langle a,t\rangle,&\varepsilon=1.
 \end{cases}}
 \tag{RT.31.10}\label{p1b:eq:rt-31-10}
\]
Consequently the full mixed quotient has 24 classes per sector-zero
free-region bundle, 12 per sector-one free-region bundle, and
\[
 \boxed{N_{\rm mix}(A)=48+24=72\qquad(A\ \mathrm{even}).}
 \tag{RT.31.11}\label{p1b:eq:rt-31-11}
\]
The auxiliary quotient using determinant-one global gauges has 24
classes per bundle, 96 in total. The odd-determinant action on it is
trivial in sector zero and is a free involution in sector one.
At Chern zero the four ordinary-image rows have uniform lift
multiplicities
\[
 \boxed{(f_{00},f_{01},f_{10},f_{11})=(1,2,1,2).}
 \tag{RT.31.12}\label{p1b:eq:rt-31-12}
\]
\end{theorem}
\begin{proof}
Equations \eqref{p1b:eq:rt-31-2}, \eqref{p1b:eq:rt-31-3}, and \eqref{p1b:eq:rt-31-7} give the actual images
at \(A=0\). The determinant-one image omits \(j_h\); all its
remaining generators are in \(\langle t\rangle\), and \(t\) is
actually realized. Adding the actual \(j_h\) contributes \(a\)
exactly in sector one. The supplied classification states that the
compatible-bundle and restriction problem is the same within a parity,
and \(\mathscr D(V,A)=\mathscr A/J_V\) in each normal component.
Transporting these images therefore gives their stated ranks at every
even \(A\); it does not transport Chern-zero ordinary fibre labels.

The retained quotient formula is \(48/2^{\dim J_{0,V}}\). The ranks
are one and two, respectively, and there are two bundles in each sector.
This proves \eqref{p1b:eq:rt-31-11}. In the auxiliary determinant-one quotient,
translation by \(j_h\) is zero in sector zero and nonzero of order two
in sector one. A nonzero translation has no fixed points, proving the
determinant-action assertion. This auxiliary construction retains the
local mixed gauge group; its count is not a count of coherent principal
\(SU(2)\) bundles.

At Chern zero, the four disjoint ordinary-image sizes are
\(24,12,12,6\). The corresponding mixed row sizes are
\(24,24,12,12\). The supplied affine forgetful homomorphism has
uniform fibres within each row. Dividing these sizes gives \eqref{p1b:eq:rt-31-12}.
\end{proof}

\subsubsection{The Chern-zero sphere relation}

\begin{theorem}[The additive quotient types]
\label{p1b:rt:theorem31-sphere}
In the supplied Chern-zero normalized coordinates, the actual
determinant-one sphere-two gauge has finite correction
\[
 z_V=5\pmod6,\qquad e_V=0,
 \tag{RT.31.13}\label{p1b:eq:rt-31-13}
\]
in both sectors and both fifth-bit classes, modulo the extendable \(t\).
After choosing an origin, the full quotients are
\[
 \boxed{\mathscr D(V,0)\cong
 \begin{cases}C_{12}\oplus C_2,&\varepsilon=0,\\
 C_6\oplus C_2,&\varepsilon=1.
 \end{cases}}
 \tag{RT.31.14}\label{p1b:eq:rt-31-14}
\]
The auxiliary determinant-one quotient is \(C_{12}\oplus C_2\)
in both sectors.
\end{theorem}
\begin{proof}
Before determinant forgetting, write the primary coefficient of the
free boundary generator as \(p_V\in C_{12}\). The supplied normalized
forgetful map has coefficient one modulo \(m_\varepsilon\), where
\(m_0=12\) and \(m_1=6\). Hence \(p_V=1\) in sector zero and
\(p_V\in\{1,7\}\) in sector one. This uses the same central
primary gluing calculation as Lemma~\ref{p1b:rt:lemma31-primary}, now
also on the free generator; its origin is fixed before changing the
determinant. In either case \(2p_V=2\pmod {12}\).

The actual \(g_2\) is determinant one. The gluing argument of
Lemma~\ref{p1b:rt:lemma31-gluing} therefore preserves the ordinary
determinant trivialization. Its primary difference must be zero before
determinant forgetting, giving
\(2p_V+2z_V=0\pmod {12}\). Thus \(z_V=5\pmod6\), including
the exotic sector. The retained ordinary top bit forces \(e_V=0\)
for \(\delta=0\). Transfer of the exact sphere gauge under the
fifth-bit surgery gives the same sphere relation in the other class.
Changes of the free-generator origin by an order-two top class do not
alter this equation for \(2b\). Since \(t\) already extends,
its correction can be removed.

Modulo the determinant-one finite image, the resulting presentation is
\[
 \langle b,z,e\mid 6z=0,\ 2e=0,\ 2b+5z=0\rangle.
\]
Its \((b,z)\)-part is cyclic of order twelve, giving
\(C_{12}\oplus C_2\). The full sector-zero image is the same.
In sector one the full image additionally imposes \(3z=0\);
the \((b,z)\)-part is then cyclic of order six, giving
\(C_6\oplus C_2\). This proves \eqref{p1b:eq:rt-31-14}.
\end{proof}

These are exact even-parity counts on the actual classification inputs. The subsequent odd-parity and homogeneous-fibre proofs complete the corresponding odd branch and every nonzero-Chern ordinary fibre.



\subsection{Actual low gauges in the odd restriction problem}
\label{p1b:sec:rtproof32}

Retain the supplied notation \(M=T^2\times P\),
\(P=S_x^3\times S_p^3\), and the free quotient \(W\) with
boundary \(R=F\times L\). The Chern vector is
\((A,-A,-A,A)\). The normalized boundary group and its binary part are
\[
 \mathscr A=\mathbb Z\{b\}\oplus C_6\{z\}\oplus C_2\{e,t\},
 \qquad T[2]=\mathbb F_2\{a,e,t\},\quad a=3z.
 \tag{RT.32.1}\label{p1b:eq:rt-32-1}
\]
The actual low generators have data \((h,0)\), \((0,\eta_W)\),
and \((0,2q)\). Subsections~\ref{p1b:sec:rtproof29}--\ref{p1b:sec:rtproof30} prove
\(u_\chi=0\), \(u_c=(1+(A\bmod2))t\), and equality of actual
boundary images and sphere relations across the fifth-bit pair.
In particular the pure finite kernel is zero at odd parity: each low
column then has a unique finite boundary value.

\subsubsection{Retaining the slice framing at nonzero Chern classes}

\begin{lemma}[The finite primary coordinate for every parity]
\label{p1b:rt:lemma32-primary}
For any actual zero-sphere global gauge with determinant \(kh\),
the \(a\)-coordinate of its normalized boundary restriction is
\(\varepsilon k\pmod2\). In particular, at odd parity,
\[
 \boxed{j_h=\varepsilon a+h_e e+h_t t,\qquad
 j_\eta=\eta_e e+\eta_t t.}
 \tag{RT.32.2}\label{p1b:eq:rt-32-2}
\]
No value of the displayed \(e,t\) coefficients is inferred by this lemma.
\end{lemma}
\begin{proof}
Before ordinary unframing, a finite insertion \(z\) changes the two
central edge coordinates by \((2z,-2z)\), independently of the
Chern degrees: the insertion is central in the actual G20 group.
An ordinary isomorphism of determinant winding \((k_t,k_E)\), whose
restriction to the chosen sphere-product slice has degree pair \(v\),
changes the edge coordinates by
\[
 N_c v+6\varepsilon(k_t,k_E),\qquad
 N_c=\begin{pmatrix}-A&-A\\A&A\end{pmatrix}.
 \tag{RT.32.3}\label{p1b:eq:rt-32-3}
\]
This is the retained slice-conjugation and determinant action in the
current ordinary owner, used with the degree pair of the actual
isomorphism rather than an arbitrary ordinary slice change.

The gluing isomorphism from a zero-sphere global gauge has determinant
winding \((k,k)\), by the transfer and codimension-four excision
argument of Subsection~\ref{p1b:sec:rtproof31}. Its degree pair on \(P\) is zero: the free
primary coefficient is zero and the torsion primary \(\eta_W\)
restricts to zero on \(P\). Removing the local circle gauge takes
place on \(N\), disjoint from this detecting slice, and does not
alter these degrees. The possibly nonzero central component of the
slice gauge has trivial conjugation action on edge components.
Thus \(v=0\) in \eqref{p1b:eq:rt-32-3}, even when \(A\) is odd.

It follows that
\[
 (2z,-2z)=6\varepsilon(k,k)\pmod {12}.
\]
Since an actual zero-sphere finite restriction belongs to \(T[2]\),
\(z\in\{0,3\}\), so \(z=3\varepsilon k\). Apply this to
\(g_h\) and the determinant-one \(g_\eta\). This proves
\eqref{p1b:eq:rt-32-2}. Quotienting by all \(N_cv\) before this comparison would
discard the restriction \(v=0\) and would not prove the result.
\end{proof}

\subsubsection{Two explicit coherent gauges in sector zero}

Use the supplied coboundary mapping-torus construction in the trivial
restriction bundle. Identify its determinant-one gauges with actual
maps \(P\to SU(2)=S^3\), and put
\[
 \alpha(f)(P)=\rho f(BP)\rho^{-1},\qquad
 \rho=\operatorname{diag}(1,-1),\qquad
 g=h\alpha(h)^{-1},\quad k=\alpha(h).
 \tag{RT.32.4}\label{p1b:eq:rt-32-4}
\]
Here \(h\) is an actual gauge realizing the prescribed \(A\) in the
retained construction. Thus \(\alpha(g)=g^{-1}\), \(gk=h\),
and \(\alpha(k)=h\). The bundle is the pullback, under the circle
ratio, of the mapping torus with convention
\[
 (r+2\pi,e)\sim(r,g(P)e).
 \tag{RT.32.5}\label{p1b:eq:rt-32-5}
\]
The construction supplies a coherent mixed bundle for every \(A\) in
sector zero. Its Born lift on the unrolled frame is \(\rho\).

Write the two torus angles as \(s+r,s\), so exchange is
\[
 (s,r,P)\longmapsto(s+r,-r,BP).
 \tag{RT.32.6}\label{p1b:eq:rt-32-6}
\]
In the cut \(-\pi\le r\le\pi\), set \(m=s+r/2\),
\(c=e^{im}\). Conjugation by \(\rho\) on the unit quaternions
fixes \(i\) and sends \(j\) to \(-j\).

\begin{lemma}[A based contraction with an explicit formula]
\label{p1b:rt:lemma32-loop}
The formula
\[
 C(m,v)=
 (\cos v\,e^{mi}+\sin v\,j)
 (\cos v+\sin v\,j)^{-1},\qquad 0\le v\le\pi/2,
 \tag{RT.32.7}\label{p1b:eq:rt-32-7}
\]
is a loop in the unit quaternions, with
\(C(m,0)=e^{mi}\), \(C(m,\pi/2)=1\), and \(C(0,v)=1\).
\end{lemma}
\begin{proof}
Both quaternionic factors before inversion have norm one: their two
summands are orthogonal in \(\mathbb H\). Their product is therefore
unit. The endpoint and based identities follow by substitution, and
the formula is periodic in \(m\) with period \(2\pi\).
\end{proof}

Choose \(v(r)\) on \([0,\pi]\), zero near zero and \(\pi\),
and equal to \(\pi/2\) near \(\pi/2\). Define
\[
 U_H(r,s,P)=
 \begin{cases}
 C(m,v(r)),&0\le r\le\pi/2,\\
 k(P)C(m,v(r))k(P)^{-1},&\pi/2\le r\le\pi.
 \end{cases}
 \tag{RT.32.8}\label{p1b:eq:rt-32-8}
\]
The two expressions agree where \(C=1\). Choose \(\theta(r)\)
zero near zero and \(\pi/2\) near \(\pi\), and also define
\[
 U_I(r,s,P)=k(P)e^{\theta(r)j}k(P)^{-1}.
 \tag{RT.32.9}\label{p1b:eq:rt-32-9}
\]
Extend either function to negative \(r\) by
\[
 U_X(-r,s,P)=\rho U_X(r,s-r,BP)\rho^{-1},
 \qquad X=H,I,
 \tag{RT.32.10}\label{p1b:eq:rt-32-10}
\]
and put \(G_X=cU_X\).

\begin{theorem}[Actual local-boundary gauges in sector zero]
\label{p1b:rt:theorem32-sector0}
The functions \(G_H,G_I\) define actual coherent unitary gauges on
the entire sector-zero mapping-torus bundle. On a neighborhood of the
fixed locus they are respectively
\[
 G_H=\operatorname{diag}(e^{2im},1),\qquad G_I=e^{im}I.
 \tag{RT.32.11}\label{p1b:eq:rt-32-11}
\]
Their descended low data on \(W\) are \((h,0)\) and
\((h,\eta_W)\). Their normalized finite boundary values are zero.
The gauge \(G_HG_I^{-1}\) is determinant one, has primary
\(\eta_W\), and restricts to the exact local circle \(\kappa\).
Consequently, for both odd fifth-bit bundles in sector zero,
\[
 \boxed{j_h=j_\eta=0.}
 \tag{RT.32.12}\label{p1b:eq:rt-32-12}
\]
\end{theorem}
\begin{proof}
The formulas are periodic in \(s\). The negative extension is exactly
the coherence equation for \eqref{p1b:eq:rt-32-6}, since its mean angle is unchanged
by exchange. Near \(r=0\), the first function is \(e^{mi}\) on
both sides and the second is one, proving \eqref{p1b:eq:rt-32-11}.

It remains to check the mapping-torus seam, not merely the pointwise
norms. Let \(m_+=s+\pi/2\). Equations \eqref{p1b:eq:rt-32-4} and \eqref{p1b:eq:rt-32-10}
give
\[
 U_H(-\pi,s,P)=-h(P)e^{m_+i}h(P)^{-1},\qquad
 U_I(-\pi,s,P)=-h(P)j h(P)^{-1}.
\]
The corresponding positive values are conjugated by \(k(P)\).
Since \(gk=h\), these equations are
\(U_X(-\pi)=-gU_X(\pi)g^{-1}\). The scalar also changes sign,
\(c(-\pi)=-c(\pi)\), so
\[
 G_X(-\pi)=gG_X(\pi)g^{-1}.
\]
This is precisely the automorphism condition for convention
\eqref{p1b:eq:rt-32-5}. Thus the functions define whole-bundle gauges with an
actual coherent exchange action.

Their determinant is \(c^2=e^{i(2s+r)}\), the product of the two
original circle coordinates. It descends with winding \(h\).
On each detecting copy of \(P\), the determinant-one part is either
constant or a conjugate \(kqk^{-1}\) of a constant unit quaternion.
Such a map is nullhomotopic by contracting \(q\) to one, so both
free primary degrees vanish.

On the fixed-circle boundary, \eqref{p1b:eq:rt-32-11} gives \(2q_+\) and
\(q_++q_-=2q_+-\kappa\). Clutching the two boundary eigenlines
shows primary zero for the first, since \(2a_F\tau=0\), and
\(a_F\tau\) for the second. The retained restriction maps therefore
give the stated global primaries zero and \(\eta_W\). Both
restrictions are actual local mixed gauges on \(N\), hence normalize
to zero. Their quotient restricts to \(q_+-q_-=\kappa\).

This constructs the result for one actual sector-zero free-region
bundle at each parity. The exact-boundary gauge transfer under the
fifth-bit surgery of Subsection~\ref{p1b:sec:rtproof29} gives the same extensions for its mate.
At odd parity all pure finite fourth corrections vanish, so the two
low-column values are uniquely zero as stated in \eqref{p1b:eq:rt-32-12}.
\end{proof}

\subsubsection{The remaining odd total and determinant action}

\begin{theorem}[One decisive test for the odd cardinality]
\label{p1b:rt:theorem32-count}
For every odd \(A\), sector zero has finite rank zero and contributes
96 mixed classes. In the exotic sector put
\[
 r_\eta=\dim\langle j_\eta\rangle\in\{0,1\}.
\]
Its full finite rank is \(1+r_\eta\), and
\[
 \boxed{N_{\rm mix}(A)=96+\frac{48}{2^{r_\eta}}
 \in\{144,120\}\qquad(A\ \mathrm{odd}).}
 \tag{RT.32.13}\label{p1b:eq:rt-32-13}
\]
Ordinary extension of the exact boundary circle \(\kappa\) in the
odd exotic bundle decides this total:
\[
 \boxed{\kappa\text{ extends }\Longrightarrow N=144,\qquad
 \kappa\text{ does not extend }\Longrightarrow N=120.}
 \tag{RT.32.14}\label{p1b:eq:rt-32-14}
\]
The actual surgery in Theorem~\ref{p1b:rt:theorem34-kappa} proves \(j_\eta=t\), hence \(r_\eta=1\) and total120.

In every parity and both fifth-bit classes, the odd-determinant action
is redundant in sector zero and nonredundant in sector one. In sector
one it is a free involution on the auxiliary quotient using
determinant-one global gauges.
\end{theorem}
\begin{proof}
At odd parity \(u_\chi=u_c=0\). Equation \eqref{p1b:eq:rt-32-12} therefore
gives \(J_0=0\) in sector zero, with 48 classes per bundle. In
sector one, \eqref{p1b:eq:rt-32-2} gives an \(a\)-coordinate of one for \(j_h\)
and zero for \(j_\eta\), so \(j_h\) is independent of the
determinant-one image. The full rank is exactly \(1+r_\eta\).
Apply the retained per-bundle formula \(48/2^{\dim J_0}\) and the
two-bundle count in each sector to obtain \eqref{p1b:eq:rt-32-13}.

The odd pure finite kernel is zero, so \(\kappa\) extends if and
only if \(j_\eta=0\). A determinant-one boundary extension can
always be made determinant one using injectivity of
\(H^1(W)\to H^1(R)\) and its relative logarithm. This proves
\eqref{p1b:eq:rt-32-14} for the specified ordinary extension problem.

At odd parity, \(j_h=0\) in sector zero and is independent of
\(j_\eta\) in sector one, settling both determinant memberships.
The even memberships were proved in Subsection~\ref{p1b:sec:rtproof31}. In every nonredundant
case the action is a nonzero order-two translation of the auxiliary
gluing torsor, so has no fixed points. The paired-image theorem carries
each assertion to the other fifth-bit bundle.
\end{proof}

For the odd exotic low columns, only the four candidates
\(a+h_e e+h_t t\) for \(j_h\), and the four candidates
\(\eta_e e+\eta_t t\) for \(j_\eta\), remain. These eight inputs,
the exotic sphere correction, and the resulting quotient type are not
evaluated by the sector-zero construction.



\subsection{Sphere primaries and ordinary lift-fibre cardinalities}
\label{p1b:sec:rtproof33}

Retain the actual images from Subsections~\ref{p1b:sec:rtproof31}--\ref{p1b:sec:rtproof32} and the supplied complete
ordinary forgetful image. Put
\[
 h_{12}=\gcd(12,A),\qquad h_6=\gcd(6,A),
 \qquad \gcd(m,0)=m.
\]
The ordinary image sizes are
\[
 I_0(A)=\begin{cases}h_{12},&A\text{ odd},\\
 4h_{12},&A\equiv2\pmod4,\\3h_{12},&4\mid A,
 \end{cases}\qquad
 I_1(A)=\begin{cases}2h_6,&A\text{ odd},\\3h_6,&A\text{ even}.
 \end{cases}
 \tag{RT.33.1}\label{p1b:eq:rt-33-1}
\]
These count underlying ordinary bundles admitting at least one mixed
coherent lift. They are not mixed lift counts.

\subsubsection{A universal normalization of the sphere primary}

\begin{lemma}[The based ordinary cover of the free boundary generator]
\label{p1b:rt:lemma33-cover}
The generator \(b\), of evaluation degree two on \(S_F^3\), can
be normalized so that its based ordered-cover insertion has coefficient
one in \(\pi_6SU(2)=C_{12}\). Then a gluing difference
\((b,z,e,t)\) has raw primary difference
\[
 (b+2z,-b-2z)\pmod {12}
 \tag{RT.33.2}\label{p1b:eq:rt-33-2}
\]
before ordinary unframing. The primary part of the ordinary forgetful
map is therefore \(b+2z\) modulo \(h_{12}\) in sector zero and
modulo \(h_6\) in sector one.
\end{lemma}
\begin{proof}
Pull a sphere generator of \(\pi_3H\) to the ordinary normal
three-sphere and remove its evaluation map using a determinant-one
local gauge on \(N\). The remaining based map is a pure
\(S_F^3\wedge S_N^3\) insertion, with some coefficient \(c\)
in \(C_{12}\). This is an operation on the fixed local boundary
model and is independent of the exterior bundle and its normal Chern
component. The local removal is the identity on the detecting \(P\),
so introduces no slice-conjugation term in the raw primary difference.

The supplied Chern-zero sector-zero forgetful proof forces this
coefficient to be odd: finite insertions give only even residues, while
the free generator must reach the odd residues. Adding a finite sphere
insertion to the free generator changes \(c\) by an arbitrary even
residue, so normalize it to one. The finite sphere insertion has
coefficient \(2z\), while \(e,t\) have no edge primary. Central
insertion on the two edges has the opposite Gysin orientations,
giving \eqref{p1b:eq:rt-33-2}, for any background Chern degree.

Ordinary slice changes add \((-A\ell,A\ell)\). Determinant
forgetting reduces the exotic central modulus to six. Quotienting
these cyclic primary differences gives the stated \(h_m\) coordinates.
\end{proof}

\begin{theorem}[The primary sphere-two correction]
\label{p1b:rt:theorem33-sphere}
In the normalization of Lemma~\ref{p1b:rt:lemma33-cover}, the actual
determinant-one sphere-two gauge has
\[
 \boxed{z_V=5-2A\pmod6.}
 \tag{RT.33.3}\label{p1b:eq:rt-33-3}
\]
In particular its \(a\)-coordinate is one for every \(A\), in both
sectors and both fifth-bit bundles. At odd parity the sector-zero full
quotient, after choosing an origin, is
\[
 \boxed{C_{12}\oplus C_2^2.}
 \tag{RT.33.4}\label{p1b:eq:rt-33-4}
\]
\end{theorem}
\begin{proof}
The gauge has determinant zero and degree pair \((2,2)\) on \(P\),
by the actual primary \(2q\). Keeping this actual slice framing in
the ordinary gluing isomorphism, equation \eqref{p1b:eq:rt-32-3} gives primary
change \((-4A,4A)\). Its boundary sphere coordinate is two, so
\eqref{p1b:eq:rt-33-2} gives
\(2+2z_V=-4A\pmod {12}\). Dividing this congruence by two gives
\eqref{p1b:eq:rt-33-3}. Its right side is odd for every \(A\).

At odd parity the sector-zero finite image is zero by Subsection~\ref{p1b:sec:rtproof32}.
The two-primary sphere correction is nonzero because its
\(a\)-coordinate is one. The retained additive presentation therefore
has two-primary group \(C_4\oplus C_2^2\), with its surviving
\(C_3\), proving \eqref{p1b:eq:rt-33-4}. The unknown \(e,t\) coordinates do
not change this isomorphism type when the finite image is zero.
\end{proof}

\subsubsection{The central and top operations on ordinary images}

Write \(\mathcal B\) for the actual ordinary top modification,
represented by the supplied \(\pi_7U(2)=C_2\) insertion. A class has
two top values precisely when \(\mathcal B\) does not fix it in
the actual ordinary quotient.

\begin{lemma}[Homogeneous fibres within a fixed free-region image]
\label{p1b:rt:lemma33-homogeneous}
At fixed \(V\), mixed gluing translations induce a transitive action
on its ordinary image. Every ordinary class in that image consequently
has the same number of lifts from \(V\).

At even parity, the two free-region image rows are disjoint. In the
exotic sector exactly one row has two ordinary top values over every
primary and the other has one. This intrinsically labels the two rows
as the special and other rows, without choosing an exotic affine origin.
\end{lemma}
\begin{proof}
Removing ordinary evaluation as in Lemma~\ref{p1b:rt:lemma33-cover}
leaves central six-cell insertions and top seven-cell insertions.
Their primary and filling actions descend through the actual ordinary
identifications. Indeed the central component acts trivially on the
G20 filling module; its mixed degree/top commutators and their
homotopies have dimensions nine and ten, above that mapping two-type.
The supplied affine-unframing lemma says that an order-twelve carry
has no affine translation. The remaining interaction with a determinant
loop is \(A_*\nu'=\eta_3E\nu'=0\), proved in the current gauge
owner. Thus the central operation also descends through determinant
forgetting, at every coefficient, and commutes with \(\mathcal B\).
The top operation has identity linear action on filling differences.

These are precisely the ordinary operations induced by the normalized
mixed gluing translations. Gluing naturality makes the forgetful map
equivariant. Its image is the orbit of the reference bundle, so is
transitive. Equivariance from a transitive translation torsor identifies
any two image fibres by a translation, proving uniformity.

For even \(A\), the diagonal lower fifth label is preserved by all
ordinary identifications: the determinant translations have lower
coefficients \(A\bmod2=0\), and slice changes preserve them.
The two fifth-bit free-region bundles give the two different lower
labels, so their images are disjoint.

In the exotic ordinary stabilizer the retained affine character can
be absorbed into a translation of the lower coordinates. Exactly one
diagonal lower label has two top classes and the other has one, as
proved in the supplied ordinary and mixed-image owners. Central
translation changes the primary by one, preserves the lower label,
and commutes with \(\mathcal B\). It therefore carries top survival
unchanged through all primary residues. The same row is special at
every primary. This proves the asserted intrinsic row labelling;
it does not set the coordinate-dependent affine character to zero.
\end{proof}

\begin{theorem}[All even individual lift cardinalities]
\label{p1b:rt:theorem33-evenfibres}
For every even \(A\), the uniform fibre cardinalities and ordinary
image sizes per free-region bundle are as follows.
\[
\begin{array}{c|c|c|c}
\text{sector}&\text{row}&\text{ordinary image size}&
 \text{mixed lifts per ordinary class}\\\hline
0, A\equiv2\pmod4&\delta=0&2h_{12}&12/h_{12}\\
0, A\equiv2\pmod4&\delta=1&2h_{12}&12/h_{12}\\
0, 4\mid A&\delta=0&2h_{12}&12/h_{12}\\
0, 4\mid A&\delta=1&h_{12}&24/h_{12}\\
1&\text{special}&2h_6&6/h_6\\
1&\text{other}&h_6&12/h_6
\end{array}
 \tag{RT.33.5}\label{p1b:eq:rt-33-5}
\]
In sector one the special row is the one where the ordinary top
modification is nontrivial. At Chern zero it is \(\delta=0\) in
the supplied chart. Each even ordinary mixed-image class has at most
six coherent mixed lifts.
\end{theorem}
\begin{proof}
The ordinary primary image has \(h_{12}\) or \(h_6\) elements
by \eqref{p1b:eq:rt-33-2}. In sector zero, if \(A\equiv2\pmod4\), the
stabilizer has diagonal parity and retains both top values at each
diagonal lower label. If \(4\mid A\), it retains both at
\(\delta=0\) and identifies them at \(\delta=1\). These are
the supplied G20 shear and mixed-image counts. In sector one, use the
special and other row of Lemma~\ref{p1b:rt:lemma33-homogeneous}.

The full mixed row sizes are 24 in sector zero and 12 in sector one,
by Subsection~\ref{p1b:sec:rtproof31}. Divide by the image size in each disjoint row and use
uniformity to obtain \eqref{p1b:eq:rt-33-5}. Since the possible even gcd values
are \(h_{12}=2,6\) when \(A\equiv2\pmod4\),
\(h_{12}=4,12\) when \(4\mid A\), and \(h_6=2,6\), the
listed cardinalities are positive integers at most six. At \(A=0\)
they recover \((1,2,1,2)\).
\end{proof}

\begin{theorem}[Odd individual lift cardinalities]
\label{p1b:rt:theorem33-oddfibres}
For every odd \(A\), every sector-zero ordinary mixed-image class
has exactly
\[
 \boxed{96/h_{12}}
 \tag{RT.33.6}\label{p1b:eq:rt-33-6}
\]
coherent mixed lifts. In the exotic sector the two free-region rows
have disjoint ordinary images, and each ordinary image class has
\[
 \boxed{\frac{24}{2^{r_\eta}h_6}}
 \tag{RT.33.7}\label{p1b:eq:rt-33-7}
\]
lifts. The actual relative sixth-stage surgery in
Theorem~\ref{p1b:rt:theorem34-kappa} gives \(r_\eta=1\), so this expression is \(12/h_6\).
\end{theorem}
\begin{proof}
For odd \(A\) the ordinary top direction is removed. In sector zero
an individual ordinary determinant change fixes the primary and
identifies the two diagonal lower labels, by the supplied determinant
action. Each free-region image is therefore the whole primary image
of size \(h_{12}\). The central coefficient two already reaches
every primary because \(h_{12}\) is odd. Each free-region bundle
has 48 mixed classes by Subsection~\ref{p1b:sec:rtproof32} and uniform fibres by
Lemma~\ref{p1b:rt:lemma33-homogeneous}. Both free-region classes lift
every ordinary image class, giving \(2\cdot48/h_{12}\).

In sector one a return to an exact raw primary forces the two
determinant parities to be equal, so cannot identify the two diagonal
lower labels. This is the supplied odd exotic image argument.
The two rows are disjoint, each with \(h_6\) primary classes.
Their mixed size is \(24/2^{r_\eta}\) by Subsection~\ref{p1b:sec:rtproof32}. Uniformity gives
\eqref{p1b:eq:rt-33-7}. No value of \(r_\eta\) is inferred from this division.
\end{proof}

The homogeneous fibre argument here is retained before the explicit odd exotic image and sphere comparisons below. The special-row
labelling in the even exotic sector is intrinsic; no identification with
an arbitrarily preassigned affine origin is asserted.



\subsection{The exact torsion-circle extension in the exotic sector}
\label{p1b:sec:rtproof34}

Retain \(W\), \(R=F\times L\), \(F=S^1\times S^3\),
\(L=\mathbb {RP}^3\), and the boundary bundle
\(E_\partial=\mathbf1\oplus\lambda_L\). Put
\(G=SU(2)\), \(H=\mathcal G_{U(2)}(E_\partial|_L)\).
The actual boundary groups and normalized finite chart are
\[
 \pi_1H=\mathbb Z\{q_+\}\oplus C_4\{\kappa\},\qquad
 \pi_4H=C_2\{e,t\},\qquad
 T[2]=\mathbb F_2\{a,e,t\}.
 \tag{RT.34.1}\label{p1b:eq:rt-34-1}
\]
Here \(t\) is the natural lens-top insertion of
\(\mu=\nu'\eta_6\ne0\in\pi_7G\). The retained unstable value
is \(\langle\iota_3,\eta_3\rangle=\mu\).
At odd parity the pure finite fourth-stage kernel has zero boundary
image. Subsection~\ref{p1b:sec:rtproof32} constructs an actual sector-zero determinant-one gauge
of primary \(\eta_W\) with exact boundary circle \(\kappa\).

Write \(\chi\in H^1(L;C_2)\),
\(x=\bar a_F\in H^1(F;C_2)\), and
\(u=\bar u_F\in H^3(F;C_2)\). Then
\[
 H^*(R;C_2)=C_2[x,u,\chi]/(x^2,u^2,\chi^4),\qquad
 \bar\eta_W|_R=x\chi^2,\quad q|_R=2u_F.
 \tag{RT.34.2}\label{p1b:eq:rt-34-2}
\]
All coefficient systems in the determinant-one adjoint bundle are
untwisted. The remaining comparison must retain the boundary
trivialization used when the sixth bundle stage is changed.

\subsubsection{A sixth-stage bundle surgery on the proper band}

Let \(C\subset W\) be the proper circle-pair M\"obius band with
both sphere coordinates equal to a fixed point. The retained normal
model is
\[
 \nu_C\cong\mathbb R^5\oplus\gamma_C,\qquad
 \operatorname {Th}(\nu_C)\simeq P^7(2)=S^6\cup_2e^7.
 \tag{RT.34.3}\label{p1b:eq:rt-34-3}
\]
Its boundary is one copy of the fixed circle, at a chosen point of the
normal lens space, and its transverse intersection with the detecting
\(P=S^3\times S^3\) has multiplicity one. Every actual \(V\) is
trivial on a neighborhood of this band.

Use \(\alpha_6\in\pi_6BU(2)=C_2\) for the bundle class whose
loop adjoint is \(\eta_3\eta_4\in\pi_5G\). It extends to an
actual based map
\[
 f:P^7(2)\longrightarrow BU(2),\qquad f|_{S^6}=\alpha_6,
 \tag{RT.34.4}\label{p1b:eq:rt-34-4}
\]
because the attaching obstruction is \(2\alpha_6=0\).
The determinant is trivial: \(P^7(2)\) has no degree-two
cohomology. In a frame near \(C\), pull this relative bundle back
by the Thom collapse, keep its basepoint frame on the outer boundary,
and glue to the unchanged exterior of \(V\).

\begin{lemma}[Actual surgery and its boundary identification]
\label{p1b:rt:lemma34-surgery}
The resulting bundle \(V^\dagger\) has the same Chern and fifth-stage
data as \(V\), and the opposite sixth-stage restriction sector.
It remains compatible with \(E_\partial\).

The boundary identification can be chosen determinant preserving,
with zero degree-three comparison and with first nonzero comparison
cochain
\[
 \boxed{\beta_4=u\chi\in H^4(R;C_2).}
 \tag{RT.34.5}\label{p1b:eq:rt-34-5}
\]
This is the relative fourth cochain in a trivialization of the
modified boundary bundle; its lens-top obstruction cancels the
inserted sixth bundle class.
\end{lemma}
\begin{proof}
The relative modification has no cells below six, so retains Chern and
fifth-stage data. On the transverse slice \(P\), the bottom Thom
disc has degree one. Thus its sixth change is the nonzero
\(\alpha_6\), toggling the G15 sector.

On the boundary, the normal bundle is trivial. Its Thom source is
\(S^6\vee S^7\). The bottom map into \eqref{p1b:eq:rt-34-3} is the bottom
inclusion \(i\). The top map is zero or \(i\eta_6\), since
\(\pi_7P^7(2)=C_2\{i\eta_6\}\). Its bundle image therefore has
bottom \(\alpha_6\) and possible top \(\alpha_6\eta_6\), whose
loop adjoint is \(\eta_3\eta_4\eta_5=6\nu'\).

Use the actual circle-cofibre coaction in the Chern-class-one component,
proved in the current gauge owner. After one loop its determinant-one
connecting map is \((2,\Omega^2A)\). Its consecutive coefficient
values are
\[
 \eta_3\longmapsto(0,\eta_3\eta_4),\qquad
 \eta_3\eta_4\longmapsto(0,6\nu').
 \tag{RT.34.6}\label{p1b:eq:rt-34-6}
\]
Consequently the two boundary bundle insertions just described are
null in their actual lens mapping spaces: the first is killed in the
degree corresponding to \(\pi_2S_L\), and the second in that
corresponding to \(\pi_3S_L\). This supplies a boundary isomorphism,
including the possible top value, rather than only zero characteristic
classes on \(R\).

Keep the nullhomotopy in the first equation of \eqref{p1b:eq:rt-34-6}. It is the
\(\eta_3\) comparison on the product of the parameter sphere
\(S_F^3\) with the circle one-cell of \(L\). On the two-cell its
two copies paste by \(2\eta_3=0\); the circle-cofibre mixed term
on the three-cell is exactly \(A_*\eta_3=\eta_3\eta_4\).
Thus the first comparison cochain evaluates as one on
\(S_F^3\times L^{(1)}\), giving \(u\chi\). There is no
degree-three comparison in this construction. The possible Thom top
map changes only the subsequent comparison stages. The one boundary
circle of \(C\) gives coefficient one, with no double-cover
multiplier. This proves \eqref{p1b:eq:rt-34-5}.
\end{proof}

\subsubsection{The relative boundary commutator}

\begin{lemma}[Transport of an actual determinant-one gauge]
\label{p1b:rt:lemma34-transport}
Suppose an actual determinant-one gauge of \(V\) has primary
\(nq+p\eta_W\), \(p\in C_2\). Transfer it across the surgery
above and identify the new boundary with \(E_\partial\) by
Lemma~\ref{p1b:rt:lemma34-surgery}. Its normalized boundary component
changes by
\[
 \boxed{p\,t.}
 \tag{RT.34.7}\label{p1b:eq:rt-34-7}
\]
At odd parity this gives the unique finite value for the indicated
low data, independently of pure finite corrections.
\end{lemma}
\begin{proof}
The band neighborhood retracts to a circle, and its determinant-one
gauge fibre is \(S^3\). A gauge there is nullhomotopic because
\(\pi_1S^3=0\). By section homotopy extension, homotope the global
gauge to the identity on a smaller surgery neighborhood. This preserves
its global primary and its boundary component. The identity on that
neighborhood pastes with the same exterior gauge on the actual modified
bundle, giving a whole-bundle transferred gauge.

Compare its boundary with the original gauge while retaining the
relative trivialization of Lemma~\ref{p1b:rt:lemma34-surgery}. The
first comparison is in degree four. A determinant-one gauge can be
made the identity through degree two; its first difference is the
degree-three primary \(\xi\). Interchanging these two relative
clutchings on a product cell gives the commutator
\[
 S^4\wedge S^3\longrightarrow G,
 \qquad\langle\eta_3,\iota_3\rangle=\mu.
 \tag{RT.34.8}\label{p1b:eq:rt-34-8}
\]
Cell by cell its coefficient is the cup product of the fourth
comparison cochain with the mod-two primary cocycle. This follows
directly by collapsing the boundaries of the two discs: their
interchange is the defining smash commutator, and the signed cell-pair
incidences are those of the cellular cup product. The coefficient group
of \(\mu\) is \(C_2\), so only the mod-two incidences remain.

The boundary primary reduces to
\[
 \bar\xi_R=\overline{2nu_F+p a_F\tau}=p x\chi^2.
\]
Hence the entire comparison is
\[
 \beta_4\bar\xi_R=p xu\chi^3,
 \qquad\langle xu\chi^3,[R]_2\rangle=1.
 \tag{RT.34.9}\label{p1b:eq:rt-34-9}
\]
This is exactly \(p\) copies of the natural seven-cell insertion
\(t\). Comparisons starting above four, or gauge differences above
three, give products of degree at least eight and cannot contribute on
the seven-dimensional boundary. There is no lower insertion, so the
\(a,e\) coordinates and the free sphere coordinate do not change.

The fourth comparison in \eqref{p1b:eq:rt-34-5} belongs to the relative bundle
trivialization. Its unfilled lens-top obstruction is the inserted
\(\alpha_6\). The actual-boundary-gauge centrality theorem applies
after this relative bundle comparison has been retained; it cannot
erase \eqref{p1b:eq:rt-34-8} by dropping the accompanying bundle change.

A different complete boundary identification differs by an actual
boundary gauge. The even local circle subgroup containing \(\kappa\)
is central across normal components, its circle--sphere products vanish,
and the parameter-top factor is central. Thus such a change does not
alter the \(\kappa+p t\) component used below. At odd parity the
zero pure finite image gives uniqueness for fixed low data. This proves
the stated transport formula in the needed actual restriction problem.
\end{proof}

\begin{theorem}[The odd exotic torsion-primary column]
\label{p1b:rt:theorem34-kappa}
For both odd exotic fifth-bit bundles the actual low column is
\[
 \boxed{j_\eta=t.}
 \tag{RT.34.10}\label{p1b:eq:rt-34-10}
\]
The exact circle \(\kappa\) does not extend ordinarily over this
bundle, whereas \(\kappa+\gamma(t)\) does extend:
\[
 \boxed{\kappa:\ \mathrm{NO},\qquad
 \kappa+\gamma(t):\ \mathrm{YES}.}
 \tag{RT.34.11}\label{p1b:eq:rt-34-11}
\]
Here \(\gamma(t)\) denotes the specified pure boundary top insertion.
\end{theorem}
\begin{proof}
Apply the surgery to the actual sector-zero gauge of Subsection~\ref{p1b:sec:rtproof32} with
primary \(\eta_W\) and exact boundary \(\kappa\).
Lemma~\ref{p1b:rt:lemma34-transport}, with \(p=1\), supplies an actual
exotic extension of \(\kappa+\gamma(t)\). The four-compatible-class
theorem locates its bundle in the opposite sector with the same absolute
fifth-stage data. The exact-boundary transfer of Subsection~\ref{p1b:sec:rtproof29} carries the result
to its fifth-bit mate. At odd parity the pure finite image is zero, so
this value is uniquely \(t\).

If the exact \(\kappa\) extended by a unitary gauge, its determinant
would be one on \(R\). Injectivity of \(H^1(W)\to H^1(R)\) and
the relative logarithm make that extension determinant one without
changing its boundary. Its primary must be \(\eta_W\): the
restriction of \(nq+p\eta_W\) is \(2nu_F+p a_F\tau\), and
the exact circle has primary \(a_F\tau\). Dividing this proposed
extension by the one already constructed would give a pure finite
correction with normalized boundary \(t\). The odd pure finite image
is zero and \(t\ne0\) by the current boundary owner. This is a
contradiction, proving the NO assertion.
\end{proof}

The argument computes the exact torsion-circle extension. It assigns
neither secondary coordinate of the odd exotic determinant column.



\subsection{All mixed cardinalities and the residual secondary type}
\label{p1b:sec:rtproof35}

Retain Subsections~\ref{p1b:sec:rtproof31}--\ref{p1b:sec:rtproof34} and write
\(h_{12}=\gcd(12,A)\), \(h_6=\gcd(6,A)\), including
\(\gcd(m,0)=m\). The odd exotic low columns now have form
\[
 j_h=a+h_e e+h_t t,\qquad j_\eta=t.
 \tag{RT.35.1}\label{p1b:eq:rt-35-1}
\]
Keep \(h_e,h_t\) as the coefficients of these actual natural maps through this presentation. Their complete evaluations are \(h_e=0,h_t=1\), proved below rather than inferred from these cardinalities.

\begin{theorem}[The actual full and determinant-one finite images]
\label{p1b:rt:theorem35-images}
In each fifth-bit pair the finite images and per-bundle mixed sizes are
\[
\begin{array}{c|c|c|c|c}
A\bmod2&\varepsilon&J_V^{\det=1}&J_V&|\mathscr D(V,A)|\\\hline
0&0&\langle t\rangle&\langle t\rangle&24\\
0&1&\langle t\rangle&\langle a,t\rangle&12\\
1&0&0&0&48\\
1&1&\langle t\rangle&\langle a+h_e e,t\rangle&12
\end{array}
 \tag{RT.35.2}\label{p1b:eq:rt-35-2}
\]
In particular
\[
 \boxed{N_{\rm mix}(A)=
 \begin{cases}72,&A\text{ even},\\120,&A\text{ odd},\end{cases}}
 \qquad
 \boxed{N_{\rm aux,det=1}(A)=
 \begin{cases}96,&A\text{ even},\\144,&A\text{ odd}.
 \end{cases}}
 \tag{RT.35.3}\label{p1b:eq:rt-35-3}
\]
The auxiliary count uses determinant-one global gauges and the full
local mixed gauge group. It is not a coherent principal-\(SU(2)\) count.
\end{theorem}
\begin{proof}
The even rows are Subsection~\ref{p1b:sec:rtproof31}. At odd parity Subsection~\ref{p1b:sec:rtproof32} gives zero in sector
zero, and the pure finite kernel is zero. Theorem~\ref{p1b:rt:theorem34-kappa}
gives \(t\) in the exotic determinant-one image. The determinant
column has \(a\)-coordinate one, so the full rank is two, with
image \(\langle t,a+h_e e\rangle\). Removing \(h_t t\) in this
subgroup description is a basis change; it does not assign \(h_t\)
for the actual primary-zero determinant generator.

Use \(|\mathscr D(V,A)|=48/2^{\dim J_V}\) and two bundles in each
sector. This gives \(2(24+12)=72\) at even parity and
\(2(48+12)=120\) at odd parity. Using the determinant-one global
images gives \(2(24+24)=96\) and \(2(48+24)=144\). In both
parities the odd-determinant action is the identity in sector zero and
a free involution in sector one, accounting for the difference 24.
\end{proof}

\begin{theorem}[Evaluated individual lift cardinalities]
\label{p1b:rt:theorem35-fibres}
An ordinary bundle outside the retained mixed-image predicate has zero
mixed lifts. Inside that image, all individual fibre cardinalities are
as follows. The predicate requires \(c_1=0,c_2=A\mathcal D\) and a
normalized determinant-preserving torus representative with
\(z_t+z_E+A^2=0\pmod {12}\) and diagonal lower labels \(r=s\),
interpreted in the actual ordinary determinant quotient. In sector one
the reduced primary condition is modulo six, with a determinant change
lifting it to that representative equation.
\[
\begin{array}{c|c|c|c}
\text{parity/sector}&\text{row}&\text{ordinary image size}&
 \text{lifts per ordinary class}\\\hline
A\equiv2\pmod4,0&\delta=0,1\text{ separately}&2h_{12}&12/h_{12}\\
4\mid A,0&\delta=0&2h_{12}&12/h_{12}\\
4\mid A,0&\delta=1&h_{12}&24/h_{12}\\
A\text{ even},1&\text{special}&2h_6&6/h_6\\
A\text{ even},1&\text{other}&h_6&12/h_6\\
A\text{ odd},0&\text{both rows combined}&h_{12}&96/h_{12}\\
A\text{ odd},1&\text{each disjoint row}&h_6&12/h_6
\end{array}
 \tag{RT.35.4}\label{p1b:eq:rt-35-4}
\]
The special even exotic row is intrinsically the one in which the
ordinary top modification survives. At Chern zero it is \(\delta=0\)
in the supplied chart. Every even fibre has at most six lifts.
At odd parity, sector-zero fibres are 96 or 32 and exotic fibres are
12 or 4, according as \(3\nmid A\) or \(3\mid A\).
\end{theorem}
\begin{proof}
Subsection~\ref{p1b:sec:rtproof33} proves the ordinary image sizes, transitive gluing action and
uniformity of each fixed free-region fibre. Its even formulas and odd
sector-zero formula are already evaluated. In its odd exotic formula,
substitute the proved \(r_\eta=1\) to obtain
\(24/(2h_6)=12/h_6\). The two exotic rows have disjoint images;
the two odd sector-zero rows have the same image and both contribute.
The gcd values at odd parity are one and three. The retained ordinary
mixed-image theorem supplies the zero-fibre predicate outside the image.
No exotic affine chart origin is chosen by the intrinsic special-row
label.
\end{proof}

\subsubsection{Matched odd sphere relations and the remaining group bit}

Normalize the local free generator as in Subsection~\ref{p1b:sec:rtproof33}. The primary part
of the actual sphere-two correction is
\[
 z_A=5-2A\pmod6,
 \qquad w_A=a+w_e e+w_t t
 \quad\text{on the two-primary part}.
 \tag{RT.35.5}\label{p1b:eq:rt-35-5}
\]

\begin{theorem}[Sphere transport and conditional exotic type]
\label{p1b:rt:theorem35-type}
At odd parity the actual pure-primary \(2q\) sphere correction is
the same in the two restriction sectors, after the boundary
identifications of Subsection~\ref{p1b:sec:rtproof34} and the fifth-bit transfer of Subsection~\ref{p1b:sec:rtproof29}.
The coefficients retain their actual natural definitions; their complete evaluations below are \(w_e=0,w_t=1\).

The sector-zero additive quotient is \(C_{12}\oplus C_2^2\), and
the exotic determinant-one global quotient is \(C_{12}\oplus C_2\).
For the exotic full quotient, one remaining bit
\[
 \omega=w_e+h_e\in C_2
 \tag{RT.35.6}\label{p1b:eq:rt-35-6}
\]
selects its additive isomorphism type:
\[
 \boxed{\omega=1:\ C_{12},\qquad
 \omega=0:\ C_6\oplus C_2.}
 \tag{RT.35.7}\label{p1b:eq:rt-35-7}
\]
Both alternatives have the proved cardinality twelve per bundle.
Equivalently, in this actual additive gluing quotient,
\[
 \boxed{6b=\omega e,\qquad
 \operatorname {ord}(3b)=\begin{cases}4,&\omega=1,\\2,&\omega=0.
 \end{cases}}
 \tag{RT.35.8}\label{p1b:eq:rt-35-8}
\]
No value of \(\omega\) or of this order is assigned here.
\end{theorem}
\begin{proof}
Transfer the actual determinant-one gauge of pure primary \(2q\)
through Lemma~\ref{p1b:rt:lemma34-transport}. Its torsion-primary
coefficient is zero and its boundary primary is \(4u_F\), with zero
mod-two reduction. Thus its boundary comparison contributes zero,
retaining the sphere coordinate two and its complete finite correction.
At odd parity the pure finite kernel is zero, so fixed pure-primary
data determine this correction uniquely. The paired-image transfer
gives the same result for the fifth-bit mate.

In sector zero the finite image is zero and the correction has
\(a\)-coordinate one, giving the already proved
\(C_{12}\oplus C_2^2\). In the exotic determinant-one quotient
\(t\) is killed and the correction still has \(a\)-coordinate one,
giving \(C_{12}\oplus C_2\).

For the full exotic quotient also kill \(a+h_e e\). Modulo this
image, the sphere correction is \((w_e+h_e)e\); the top term is
already killed. The remaining finite two-primary group has one
generator. Its relation with the integer generator is \(2b+\omega e=0\).
If \(\omega=1\) this makes a \(C_4\); if \(\omega=0\) it gives
\(C_2^2\). The retained surviving \(C_3\) produces the two groups
in \eqref{p1b:eq:rt-35-7}. Gluing sets are torsors until an origin is chosen.
Tripling the actual sphere relation gives
\(6b=-3z_Az-w_e e-w_t t\). Since \(z_A\) is odd,
\(3z_Az=a\); kill \(t\) and \(a+h_e e\) to obtain
\(6b=\omega e\). The sphere parity of \(3b\) is nonzero,
and the surviving finite \(e\) is nonzero, proving \eqref{p1b:eq:rt-35-8}.
\end{proof}

The complete source-bound natural evaluations are given in Theorem~\ref{p1b:thm:natural-rank-two-completion}; the evaluated counts, quotient groups and ordinary fibres are summarized in Theorem~\ref{p1b:thm:complete-mixed-rank-two}.



\subsection{Transport between normal components}
\label{p1b:sec:rtproof36}

Retain the actual boundary gauge group
\(H=\mathcal G_{U(2)}(\mathbf1\oplus\lambda_L)\),
\(L=\mathbb {RP}^3\), and the natural normalized chart
\[
 \mathscr A=\mathbb Z\{b\}\oplus C_6\{z\}
                 \oplus C_2\{e,t\},\qquad a=3z.
 \tag{RT.36.1}\label{p1b:eq:rt-36-1}
\]
The evaluation of \(b\) has degree two. The natural cover sends
\(z\) to \(2\nu'\), \(e\) to \(\mu\), and \(t\) to zero.
Here \(\langle\iota_3,\iota_3\rangle=\nu'\), up to orientation,
and \(\mu=\nu'\eta_6\). The local even circle subgroup
\(2\mathbb Zq_+\oplus C_4\kappa\) is central across normal
components. These are the retained actual statements, not an untwisted
lens-space approximation.

\begin{lemma}[Conjugation by a normal representative]
\label{p1b:rt:lemma36-conjugation}
Choose a normal representative \(h_n\in H\) supported on the open
lens three-cell, whose ordered cover has degree \(2n\). In compatible
orientations, its conjugation has the following action:
\[
 b\longmapsto b+2nz,\qquad z\longmapsto z,\qquad
 e\longmapsto e,\qquad t\longmapsto t.
 \tag{RT.36.2}\label{p1b:eq:rt-36-2}
\]
On circle classes it sends \(q_+\) to \(q_++2n\kappa\), and
fixes \(2q_+\) and \(\kappa\).
\end{lemma}
\begin{proof}
Conjugation preserves sphere evaluation. Thus its difference on \(b\)
belongs to the natural \(C_6\) subgroup of \(\pi_3H\).
Pull the comparison to the ordered normal sphere. Its only contributing
term is the commutator of the evaluation \(2\iota_3\) with the
normal map \(2n\iota_3\). Collapsing the parameter and normal wedges
gives
\[
 \langle2n\iota_3,2\iota_3\rangle=4n\nu'
       \in\pi_6SU(2)=C_{12}.
 \tag{RT.36.3}\label{p1b:eq:rt-36-3}
\]
The part of \(b\) already based in the lens variable has positive
lens filtration. Its commutator with the supported three-cell class
has lens filtration greater than three, hence contributes zero. The
cover of the natural torsion is the injective map \(z\mapsto2z\),
so \eqref{p1b:eq:rt-36-3} gives precisely \(2nz\). There is no parameter circle
in this comparison and thus no \(\pi_4H\) term.

The same filtration argument fixes \(z\): it has lens filtration
three. In degree four the actual pointed calculation supplies
representatives of \(e,t\) with lens filtrations at least two and
three respectively; the degree-one lens factor has died in the
circle-cofibre sequence. Since evaluation is zero, these representatives
can be chosen trivial below those filtrations. The cellular diagonal
shows that their commutators with a lens-three class have filtrations
at least five and six, outside the three-dimensional lens space.
They are therefore fixed. Equivalently, on a cell pair the commutator
requires the product of the two lens cells, whose degree exceeds three.

The circle formula is the retained supported Bott commutator. The
scalar loop \(2q_+-\kappa\) is pointwise central; the exponent-two
target of circle--sphere products then fixes \(\kappa\) and
\(2q_+\). The multiplication chart of maps from
\(S^1\times S^3\) to \(H\) is natural for group automorphisms.
Conjugation acts on each parameter factor as just computed; the
central even circle subgroup adds no product correction on normalization.
This proves the complete displayed action on \(\mathscr A\).
\end{proof}

\begin{theorem}[Exact parity transport of the relations]
\label{p1b:rt:theorem36-transport}
Use the gluing convention in which a normal change \(n\) increases
the ordinary Chern label by \(2n\). Transport of the normalized
global restriction image from \(A\) to \(A+2n\) is the automorphism
\[
 T_n(b)=b-2nz,\qquad T_n(z)=z,\quad T_n(e)=e,\quad T_n(t)=t.
 \tag{RT.36.4}\label{p1b:eq:rt-36-4}
\]
In particular, for the actual sphere-two relation
\[
 r_A=2b+(5-2A)z+w_e e+w_t t,
 \qquad T_n(r_A)=r_{A+2n}.
 \tag{RT.36.5}\label{p1b:eq:rt-36-5}
\]
The two secondary coefficients in this relation are constant within a
parity, with the same boundary identifications as Subsections~\ref{p1b:sec:rtproof29} and~\ref{p1b:sec:rtproof34}.
The finite low columns are also constant within a parity. Hence at odd
parity the four coefficients \(h_e,h_t,w_e,w_t\) are four shared
bits, not separate inputs for each odd integer \(A\).
\end{theorem}
\begin{proof}
The free-region compatible bundle depends only on the parity and its
restriction and fifth-stage labels. Changing the chosen boundary
gluing by \(h_n\) acts on its restriction image by inverse
conjugation. Lemma~\ref{p1b:rt:lemma36-conjugation} therefore gives
\eqref{p1b:eq:rt-36-4}. This convention agrees with the retained raw primary
formula: the coefficient of \(z\) in a sphere-two relation changes
by \(-4n\), exactly
\(5-2(A+2n)-(5-2A)\).

The finite generators are fixed and
\(T_n(2b)=2b-4nz\), proving \eqref{p1b:eq:rt-36-5}.
The odd pure finite kernel is zero, so no additional low-data
correction can change this transport. The retained fifth-bit transfer
matches actual restrictions, and the determinant-one surgery transport
of Subsection~\ref{p1b:sec:rtproof34} matches the sphere relation across restriction sectors.
Together these yield the claimed common odd coefficients. The finite
determinant column in the exotic sector is unchanged by \eqref{p1b:eq:rt-36-4},
so its two coefficients are shared across all odd \(A\) as well.
\end{proof}

\begin{corollary}[All even quotient types]
\label{p1b:rt:corollary36-even}
For every even \(A\), after choosing an origin, the full mixed
quotients in sectors zero and one have additive types
\[
 C_{12}\oplus C_2,\qquad C_6\oplus C_2,
 \tag{RT.36.6}\label{p1b:eq:rt-36-6}
\]
respectively. The auxiliary quotients using determinant-one global
gauges both have type \(C_{12}\oplus C_2\).
\end{corollary}
\begin{proof}
Subsection~\ref{p1b:sec:rtproof31} computes these groups at \(A=0\), where the sphere
correction has zero \(e\) coefficient modulo the extendable \(t\).
The exact automorphism \eqref{p1b:eq:rt-36-4} transports the complete presentations,
not merely their ranks or cardinalities, to every even \(A\).
\end{proof}

This calculation evaluates the normal-component shear. It does not
evaluate the four remaining natural odd coefficients or select the odd
exotic quotient type. In the universal chart the restriction images
are equivalent by \eqref{p1b:eq:rt-36-4}; their written sphere coefficients need
not be equal. The detailed arithmetic and this chart distinction are
recorded in the appendix.



\subsection{Uniform generators and the groups of lift differences}
\label{p1b:sec:rtproof37}

Retain Subsections~\ref{p1b:sec:rtproof31}--\ref{p1b:sec:rtproof36}. At odd parity the actual sphere correction is
\(z_Az+w_e e+w_t t\), where \(z_A=5-2A\pmod6\) is odd.
Introduce matched generators and comparison bits
\[
 \widehat z=z+w_e e+w_t t,\qquad
 \omega=w_e+h_e,\quad \zeta_{\rm nat}=w_t+h_t.
 \tag{RT.37.1}\label{p1b:eq:rt-37-1}
\]
These definitions use the actual sphere-two restriction. They are a
change of presentation; the original natural coefficients are evaluated by the later complete map comparisons. Since \(z_A\) is odd, the actual
sphere relation is exactly \(2b+z_A\widehat z=0\), while
\[
 j_h=3\widehat z+\omega e+\zeta_{\rm nat} t,\qquad j_\eta=t.
 \tag{RT.37.2}\label{p1b:eq:rt-37-2}
\]
The matched generator has order six and
\(2\widehat z=2z\). Its natural ordered-cover value is
\((2,w_e)\), so this presentation does not replace the natural
cover chart of the current manuscript.

\begin{theorem}[A uniform cyclic generator before determinant reduction]
\label{p1b:rt:theorem37-generator}
Put
\[
 g_A=b-A\widehat z\quad(A\text{ odd}).
 \tag{RT.37.3}\label{p1b:eq:rt-37-3}
\]
In the quotient by the actual sphere relation,
\(2g_A=\widehat z\) and \(12g_A=0\). The odd sector-zero
full quotient and the odd exotic determinant-one global quotient are,
respectively,
\[
 C_{12}\{g_A\}\oplus C_2\{e,t\},\qquad
 C_{12}\{g_A\}\oplus C_2\{e\}.
 \tag{RT.37.4}\label{p1b:eq:rt-37-4}
\]
The odd exotic full quotient has the uniform presentation
\[
 \boxed{\langle g_A,e\mid12g_A=0,\ 2e=0,\ 6g_A=\omega e\rangle.}
 \tag{RT.37.5}\label{p1b:eq:rt-37-5}
\]
It has twelve representatives \(rg_A+se\),
\(0\le r<6\), \(s\in C_2\). The additive law in these
representatives is
\[
 (r,s)+(r',s')=
 \bigl((r+r')\bmod6,\ s+s'+\omega\lfloor(r+r')/6\rfloor\bigr).
 \tag{RT.37.6}\label{p1b:eq:rt-37-6}
\]
Thus \(\omega=1\) gives \(C_{12}\) and \(\omega=0\) gives
\(C_6\oplus C_2\). The subsequent coincident-circle theorem proves \(\omega=0\).
At even parity, work modulo the extendable \(t\) and put
\(g_A=b-Az\). Then \(2g_A=z\); its order is twelve in
sector zero and six in the full exotic quotient.
\end{theorem}
\begin{proof}
The triangular change \eqref{p1b:eq:rt-37-1} is an automorphism of the finite
boundary group. Because the coefficient of \(z\) in the sphere
relation is odd, it absorbs exactly the two secondary sphere terms.
Now
\[
 2g_A=-(z_A+2A)\widehat z=-5\widehat z=\widehat z.
\]
Substitute \(\widehat z=2g_A\) and
\(b=(1+2A)g_A\) in the presentation. Before finite reduction
this gives \(C_{12}\oplus C_2^2\), with no relation among
\(e,t\). The exotic determinant-one image kills \(t\).
Its full image also kills \eqref{p1b:eq:rt-37-2}, giving
\(6g_A=\omega e\). Every element reduces uniquely to the displayed
twelve representatives; the only carry in their sum is this relation.
The even assertion follows by the same substitution in the transported
presentations of Subsection~\ref{p1b:sec:rtproof31}, where \(w_e=0\) and \(t\) has already
been killed. The extra even exotic relation is \(3z=6g_A=0\).
\end{proof}

Write \(h_{12}=\gcd(12,A)\), \(h_6=\gcd(6,A)\), including
\(\gcd(m,0)=m\). All groups below act on fibres; the fibres are
torsors until a lift is chosen. The word ``special'' refers to the
intrinsic even exotic row in which the ordinary top insertion survives.

\begin{theorem}[The groups acting simply transitively on fixed-row fibres]
\label{p1b:rt:theorem37-kernels}
Every nonempty ordinary forgetful fibre within a fixed free-region row
is a torsor for the following group:
\[
\begin{array}{c|c|c}
 \text{parity and sector}&\text{row}&\text{fibre group}\\\hline
 A\equiv2\pmod4,\ 0&\delta=0,1&C_{12/h_{12}}\\
 4\mid A,\ 0&\delta=0&C_{12/h_{12}}\\
 4\mid A,\ 0&\delta=1&C_{12/h_{12}}\oplus C_2\\
 A\text{ even},\ 1&\text{special}&C_{6/h_6}\\
 A\text{ even},\ 1&\text{other}&C_{6/h_6}\oplus C_2\\
 A\text{ odd},\ 0&\text{each row}&C_{12/h_{12}}\oplus C_2^2\\
 A\text{ odd},\ 1&\omega=1&C_{12/h_6}\\
 A\text{ odd},\ 1&\omega=0&C_{6/h_6}\oplus C_2
\end{array}
 \tag{RT.37.7}\label{p1b:eq:rt-37-7}
\]
Here \(C_1\) is the trivial group. In the odd sector-zero case an
ordinary fibre contains two disjoint row torsors of this size, because
the two free-region rows have the same ordinary image. No single group
structure on their union is selected. The exotic rows have disjoint
ordinary images.
\end{theorem}
\begin{proof}
The retained central six-cell and top seven-cell translations act
transitively on each ordinary image. Their action descends through the
actual ordinary gauge quotient: the affine carry is zero and the
determinant commutator on the central coefficient is
\(A_*\nu'=\eta_3E\nu'=0\). This is the actual homogeneous-fibre
lemma of Subsection~\ref{p1b:sec:rtproof33}. For a chosen ordinary class, its preimage is therefore
a torsor for the kernel of this translation action.

The raw primary coefficient of \(g_A\) is \(1-2A\), and hence
is one modulo \(h_{12}\) or \(h_6\), respectively. This holds
also with \(\widehat z\), whose primary cover coefficient is two.
The natural \(e\) is the ordinary top operation; \(t\) has zero
ordinary cover. In an even row the mixed group is
\(C_m\{g_A\}\oplus C_2\{e\}\), with \(m=12\) or six.
If the top operation survives, projection of the kernel to the cyclic
factor is an isomorphism onto its subgroup generated by \(hg_A\):
for each cyclic coefficient divisible by \(h\), there is exactly one
\(e\) correction that fixes the ordinary class. If the top operation
is killed, both \(e\) values fix it and the kernel is
\(C_{m/h}\oplus C_2\). This argument does not assume that an
ordinary image has a chosen product chart; a surviving top coordinate
may have a primary carry. The retained survival rows in Subsection~\ref{p1b:sec:rtproof35} give
the first five lines of \eqref{p1b:eq:rt-37-7}.

At odd parity the ordinary top operation is killed in every row.
Sector zero has \(C_{12}\oplus C_2^2\), giving the sixth line.
The full exotic group is cyclic of order twelve when \(\omega=1\),
with \(e=6g_A\), or is \(C_6\oplus C_2\) when \(\omega=0\).
In either case the ordinary primary is \(g_A\bmod h_6\) and
there is no surviving ordinary top coordinate. This gives the final
two lines. The retained free-row incidence determines which row torsors
are disjoint and which contribute to the same ordinary fibre.
\end{proof}

\begin{corollary}[A concrete remaining type test]
\label{p1b:rt:corollary37-test}
In the odd exotic auxiliary quotient,
\(6g_A=3\widehat z\) is a specified nonzero order-two element.
It is killed by the full odd-determinant reduction if and only if
\(\omega=0\). The full quotient has
\[
 \operatorname{ord}(3b)=2\ (\omega=0),\qquad
 \operatorname{ord}(3b)=4\ (\omega=1).
 \tag{RT.37.8}\label{p1b:eq:rt-37-8}
\]
This is one common test for all odd \(A\), both fifth-bit mates,
and their corresponding exotic rows. The actual coincident-circle calculation and split-type theorem below select \(\omega=0\).
\end{corollary}
\begin{proof}
Modulo the determinant-one image the odd-determinant translation is
\(6g_A+\omega e\). It kills \(6g_A\) exactly when
\(\omega=0\), since \(e\) is independent before that reduction.
Since \(b=(1+2A)g_A\) and \(1+2A\) is odd, \(3b\) has
the stated orders in the two presentations.
\end{proof}

The complete source-bound natural evaluations are given in Theorem~\ref{p1b:thm:natural-rank-two-completion}; the evaluated counts, quotient groups and ordinary fibres are summarized in Theorem~\ref{p1b:thm:complete-mixed-rank-two}.



\subsection{The terminal relative obstruction maps}
\label{p1b:sec:rtproof38}

Retain the compact oriented eight-manifold \(W\), its connected
boundary \(R=(S^1\times S^3)\times\mathbb {RP}^3\), and the
actual boundary bundle \(\mathbf1\oplus\lambda_L\). Put
\(C=C_{12}=\pi_6SU(2)\), with coefficient generator \(\nu'\).
Use \(a_F,u_F,v_L\) for the integral circle, sphere, and oriented
lens-top generators, and \(\tau_L=\beta\chi_L\). Their reductions
on \(R\) are denoted \(x,u,\chi\), with
\(H^*(R;C_2)=C_2[x,u,\chi]/(x^2,u^2,\chi^4)\).

The retained integral boundary map is
\[
 H_1(R;\mathbb Z)=\mathbb Z\oplus C_2
 \longrightarrow H_1(W;\mathbb Z)=\mathbb Z\oplus C_2,
 \qquad(m,\epsilon)\longmapsto(2m,\epsilon).
 \tag{RT.38.1}\label{p1b:eq:rt-38-1}
\]
One free generator is a loop moving one particle around its circle
without exchange; the torsion generator is a local exchange loop.
The cover character is zero on the former and one on the latter.

\begin{theorem}[Exact degree-six and degree-seven coefficient sequence]
\label{p1b:rt:theorem38-sequence}
Put \(U=u_Fv_L\), and let \(B=a_Fu_F\tau_L\) denote its
coefficient-reduced class. Then
\[
 H^6(R;C)=C_{12}\{U\}\oplus C_2\{B\},\qquad
 H^7(W,R;C)=C_{12}\{\Gamma_h\}\oplus C_2\{\Gamma_\chi\}.
 \tag{RT.38.2}\label{p1b:eq:rt-38-2}
\]
In duality generators the terminal exact sequence is
\[
 0\longrightarrow C_2
 \xrightarrow{s\mapsto(6s,0)} C_{12}\oplus C_2
 \xrightarrow{\delta_6:(r,s)\mapsto(2r,s)} C_{12}\oplus C_2
 \xrightarrow{(r,s)\mapsto r\bmod2} C_2
 \longrightarrow0.
 \tag{RT.38.3}\label{p1b:eq:rt-38-3}
\]
Its terms are \(H^6(W;C)\), \(H^6(R;C)\),
\(H^7(W,R;C)\), and \(H^7(W;C)\), respectively. Signs of
free duality generators may be fixed to give the displayed positive two.
The natural top map is also an isomorphism
\[
 \delta_7:H^7(R;C_2)\xrightarrow{\cong}H^8(W,R;C_2)=C_2.
 \tag{RT.38.4}\label{p1b:eq:rt-38-4}
\]
\end{theorem}
\begin{proof}
Integral \(H_1(W)=\mathbb Z\oplus C_2\) and connectedness give
\(H_1(W;C)=C_{12}\oplus C_2\), by universal coefficients;
there is no degree-zero Tor term. The same calculation on \(R\),
together with Poincare--Lefschetz duality, gives \eqref{p1b:eq:rt-38-2}. The
classes \(U,B\) are the boundary duals of its free circle and
lens torsion loop. The compatibility of duality with the boundary
map identifies \(\delta_6\) with the inclusion map \eqref{p1b:eq:rt-38-1}.
Thus \(\delta_6(U)=2\Gamma_h\) and
\(\delta_6(B)=\Gamma_\chi\).

The retained \(H^6(W;C)=C_2\) generator restricts to \(6U\),
with no torsion component. Its image is precisely the two-element
kernel of \(\delta_6\). The cokernel of \(\delta_6\) is the
free-coordinate parity, which gives \(H^7(W;C)=C_2\).
Connected oriented boundary gives zero restriction in top boundary
degree, and \(H^8(W;C)=0\). This proves the entire exact segment.
The same connected-boundary duality with \(C_2\) coefficients
identifies the map in \eqref{p1b:eq:rt-38-4} with the isomorphism on \(H_0\).
In particular the absolute equality \(H^8(W;C_2)=0\) does not
make the relative final obstruction group zero.
\end{proof}

\begin{lemma}[The natural finite boundary insertions at their first stage]
\label{p1b:rt:lemma38-insertions}
For a boundary gauge trivial through degree five, the natural insertions
\(z,e,t\) have first degree-six differences \(U,B,0\),
respectively, in the coefficient group \(C\). Changing the supported
representative of \(z\) by its killed coefficient six changes this
cohomology difference by \(6U\), whose relative image is zero.
Consequently their degree-seven relative differences are unambiguously
\[
 z\longmapsto2\Gamma_h,\qquad
 e\longmapsto\Gamma_\chi,\qquad t\longmapsto0.
 \tag{RT.38.5}\label{p1b:eq:rt-38-5}
\]
The natural \(t\) has degree-seven boundary difference the nonzero
\(\pi_7SU(2)=C_2\) coefficient on the boundary fundamental cell.
\end{lemma}
\begin{proof}
The natural \(z\) is the lens-three-cell insertion of \(\nu'\),
parametrized by \(S_F^3\); its six-cell coefficient is \(U\).
The actual map kills coefficient six after lower homotopy, and the
same coefficient is in the kernel of \eqref{p1b:eq:rt-38-3}.

The natural \(e=J\nu'\) is the pullback of the pointed sphere
gauge through \(L\to S^2\). Its four-dimensional parameter is
the smash of the circle and sphere of \(F\). The pullback of the
integral two-class of \(S^2\) is \(\tau_L\). Its first six-cell
coefficient is therefore \(a_Fu_F\tau_L=B\), whose mod-two
reduction is \(xu\chi^2\). The actual pointed calculation allows
the representative to be trivial below lens filtration two; there is
no degree-five factor or evaluation contribution. A change of extended
pinch can add the lens-top \(t\) without changing this first coefficient.
Finally \(t=j_7\mu\) is supported on a seven-disc, giving the
stated last difference and zero in degree six. Apply \eqref{p1b:eq:rt-38-3}.
\end{proof}

\begin{theorem}[Complete seven-skeleton test for pure finite boundary gauges]
\label{p1b:rt:theorem38-pure}
Let \(\gamma(z_0z+e_0e+t_0t)\), with \(z_0\in C_6\),
\(e_0,t_0\in C_2\), be a specified pure finite boundary gauge:
its determinant, circle, and primary data are zero. For every compatible
actual \(V\), in either parity or sector, its essential degree-seven
relative obstruction is
\[
 \boxed{O_7=2z_0\Gamma_h+e_0\Gamma_\chi.}
 \tag{RT.38.6}\label{p1b:eq:rt-38-6}
\]
It extends as an actual gauge through \(W^{(7)}\cup R\) if and
only if \(z_0=e_0=0\). In particular the natural \(e\) fails
already at this stage, while \(t\) extends through this stage.
\end{theorem}
\begin{proof}
Injectivity of determinant restriction and the actual primary
restriction \(nq+p\eta_W\mapsto2nu_F+p a_F\tau_L\)
force zero determinant and zero primary for any proposed extension.
Its global determinant logarithm can be removed relative to the boundary.
Fix the resulting identity homotopy through degree three.

The fourth restriction kernel is spanned by \(\chi^4,\bar c_4\):
the other fourth class \(\alpha\) restricts to
\(x\chi^3\ne0\), the actual nontrivial circle \(2\kappa\).
Returning primary homotopies introduce no extra term: on the sole
integral two-class the operation \(Sq^2+\bar\tau\) is
\(\bar\tau^2+\bar\tau^2=0\). Both fourth kernel directions
have actual global realizations. Their entire boundary values are
\(u_\chi=0\) and \(u_c=(1+(A\bmod2))t\), respectively;
in particular both are the identity through degree six.
Subtract these gauges to normalize any allowed fourth data without
changing the boundary six-stage comparison.

The actual square \(g_\alpha^2\) realizes the only nonzero fifth
difference and has whole boundary value the identity. Subtract it if
necessary to normalize through degree five. These retained whole-gauge
realizations remove all lower-choice indeterminacy in the six-stage
comparison. They apply as well to proposed partial extensions.

With the identity fixed through five, assign the prescribed six-cell
coefficients on \(R\) and arbitrary coefficients on the remaining
six-cells. The next obstruction is the relative coboundary of this
cochain, namely \(\delta_6(z_0U+e_0B)\). Coefficient systems
are untwisted because all bundle transition actions are conjugations
by the connected group \(U(2)\). Changing interior six-cell
choices changes only its coboundary representative, and changing an
absolute sixth class adds a boundary multiple of \(6U\), already
killed by \(\delta_6\). This proves \eqref{p1b:eq:rt-38-6} with the necessary
lower choices included. The map \((z_0,e_0)\mapsto(2z_0,e_0)\)
is injective, proving the exact test.

For \(z_0=e_0=0\), the seven-stage obstruction is zero; hence
the relative cell extension exists through seven. No eight-stage
conclusion is made at this step.
\end{proof}

These are actual relative coefficient and pure finite extension
calculations. They do not yet assign the obstruction of the remaining
odd determinant boundary class, whose low determinant data are nonzero.



\subsection{The final relative stage and the shared odd type test}
\label{p1b:sec:rtproof39}

\begin{theorem}[The final pure finite extension test]
\label{p1b:rt:theorem39-final}
For a pure finite boundary gauge with \(z_0=e_0=0\), the final
obstruction with lower identity choices is \(t_0\Theta\), where
\(\Theta\) generates \(H^8(W,R;C_2)=C_2\).
The permitted fourth-stage change \(\bar c_4\) changes this
obstruction by \((1+(A\bmod2))\Theta\).
Consequently a pure finite boundary gauge extends over all of \(W\)
exactly as follows:
\[
 \boxed{
 \begin{array}{c|c}
 A\text{ even}&z_0=e_0=0,\quad t_0\text{ arbitrary},\\
 A\text{ odd}&z_0=e_0=t_0=0.
 \end{array}}
 \tag{RT.39.1}\label{p1b:eq:rt-39-1}
\]
Of the twenty-four pure finite classes, twenty-two fail at degree seven
in each parity. The two survivors both extend at even parity; at odd
parity exactly one survives degree eight.
\end{theorem}
\begin{proof}
The coefficient \(\mu\in\pi_7SU(2)=C_2\) of the natural
\(t\) is a boundary top cocycle. Its relative difference is
\(\delta_7(t_0[R]^*)=t_0\Theta\), by \eqref{p1b:eq:rt-38-4}.
Interior seven-cell changes cannot cancel this class: the image of
\(H^7(W;C_2)\) in top boundary cohomology is zero. The base
identity extends, so this is the actual difference of final obstructions.

The preceding proof enumerates the permitted lower choices. The
fourth class \(\chi^4\), the fifth realization \(g_\alpha^2\),
and every stage-five-kernel gauge have identity whole boundary value.
The only lower choice that can change the top comparison is the actual
\(g_c\), with boundary \((1+(A\bmod2))t\), computed in Subsection~\ref{p1b:sec:rtproof30}.
Its multiplication therefore changes the final obstruction by precisely
the displayed amount. At even parity this cancels either top value;
at odd parity it supplies no change. The already evaluated pure finite
image theorem excludes further indeterminacy. There are no cells above
eight. Together with \eqref{p1b:eq:rt-38-6}, this proves \eqref{p1b:eq:rt-39-1} and the counts.
\end{proof}

\subsubsection{The type coefficient is at degree seven}

At odd parity retain the actual matched element
\(M=3\widehat z=a+w_e e+w_t t\),
\(\widehat z=z+w_e e+w_t t\), and put
\[
 \omega=w_e+h_e,\qquad \zeta_{\rm nat}=w_t+h_t.
\]
Here \(\zeta_{\rm nat}\) denotes the matched binary comparison, not the
local circle \(2\kappa\). Consider the four specified boundary
classes from Subsection~\ref{p1b:sec:rtproof37}:
\[
 B_{d,p}=2q_++p\kappa+\gamma(M+d e),\qquad d,p\in C_2.
 \tag{RT.39.2}\label{p1b:eq:rt-39-2}
\]

\begin{theorem}[Separated matched obstructions]
\label{p1b:rt:theorem39-matched}
For either odd exotic fifth-bit bundle, the essential degree-seven
obstruction for \eqref{p1b:eq:rt-39-2} is
\[
 \boxed{O_7(B_{d,p})=(d+\omega)\Gamma_\chi.}
 \tag{RT.39.3}\label{p1b:eq:rt-39-3}
\]
When it is zero, the final obstruction is
\[
 \boxed{O_8(B_{d,p})=(p+\zeta_{\rm nat})\Theta.}
 \tag{RT.39.4}\label{p1b:eq:rt-39-4}
\]
Thus the cyclic-versus-split quotient type is decided at degree seven,
independently of the final top obstruction. The extending member of
the four candidates is \(B_{\omega,\zeta_{\rm nat}}\); the included exact-map proofs give \(\omega=0,\zeta_{\rm nat}=0\).
\end{theorem}
\begin{proof}
Use any actual global gauge of low data \((h,0)\), whose normalized
finite restriction is \(a+h_e e+h_t t\), and the actual
\(\eta_W\) gauge with boundary \(\kappa+t\). The known
fourth gauge adjusts a possible even \(2\kappa\) circle discrepancy,
without changing its normalized finite value. Divide the proposed
boundary class by these actual restrictions. The residual pure finite
class is exactly
\[
 \gamma\bigl((d+w_e+h_e)e+(p+w_t+h_t)t\bigr).
 \tag{RT.39.5}\label{p1b:eq:rt-39-5}
\]
The known \(a\) coefficients cancel. Its low data are zero, so
Theorems~\ref{p1b:rt:theorem38-pure} and
\ref{p1b:rt:theorem39-final} apply. Actual global multiplication
identifies the relative extension problems, including partial
extensions. This gives the two displayed essential obstructions.
\end{proof}

\begin{corollary}[A single mod-two pairing for the remaining type]
\label{p1b:rt:corollary39-pairing}
For the prescribed boundary class \(B_{0,0}\), after the indicated
matching of lower data, the remaining type is the value
\[
 \boxed{\omega=
 \left\langle\chi_W\smile\overline{O_7(B_{0,0})},
                  {}[W,R]_2\right\rangle.}
 \tag{RT.39.6}\label{p1b:eq:rt-39-6}
\]
The detector has coefficient one:
\[
 \langle\chi_W\smile\overline{\delta_6(B)},[W,R]_2\rangle
   =\langle\chi_L\smile xu\chi_L^2,[R]_2\rangle=1,
 \qquad
 \langle\chi_W\smile\overline{\delta_6(U)},[W,R]_2\rangle=0.
 \tag{RT.39.7}\label{p1b:eq:rt-39-7}
\]
\end{corollary}
\begin{proof}
Boundary compatibility of relative cup product and the connecting map
turns the first pairing into \(\langle xu\chi^3,[R]_2\rangle=1\).
For \(U\), its reduction is \(u\chi^3\), whose product with
\(\chi\) is zero since \(\chi^4=0\). Equations \eqref{p1b:eq:rt-38-3} and
\eqref{p1b:eq:rt-39-3} give \eqref{p1b:eq:rt-39-6}. The detector and its coefficient are
evaluated; the geometric cochain \(O_7(B_{0,0})\) is not.
\end{proof}

The subsequent coincident-circle calculation evaluates this degree-seven cocycle as zero, independently of the final top coefficient. The split type is \(C_6\oplus C_2\). The final coherent natural comparison evaluates \(\zeta_{\rm nat}=0\).



\subsection{The coincident-circle detector}
\label{p1b:sec:rtproof40}

Retain \(W,R,F,L,V\) and the natural boundary generators of the current
rank-two owners, together with the actual \(g_h,g_2,g_c\) of Subsections~\ref{p1b:sec:rtproof30}--\ref{p1b:sec:rtproof39}. Put \(Q=\mathbb {RP}^2\subset L\). The inclusion uses the
sphere-tangent directions in the product framing
\(T(S^1\times S^3)=\mathbb R\oplus\operatorname {Im}\mathbb H\).
Write
\[
 H_Q=\mathcal G_{U(2)}(\mathbf1\oplus\lambda_Q),\qquad
 R_{\mathcal C}=F\times Q.
 \tag{RT.40.1}\label{p1b:eq:rt-40-1}
\]
This is a restriction of the actual bundle and its gauges, not an
untwisted replacement for the lens calculation.

\subsubsection{The Moore boundary and frame independence}

\begin{lemma}[The boundary information retained on \(Q\)]
\label{p1b:rt:lemma40-Q}
The group \(H_Q\) is connected, and
\[
 \pi_1H_Q=\mathbb Z\{q_+\}\oplus C_2\{\kappa_Q\},\qquad
 \pi_2H_Q=0,\qquad
 \pi_3H_Q=\mathbb Z\{b_Q\}\oplus C_2.
 \tag{RT.40.2}\label{p1b:eq:rt-40-2}
\]
Here \(\kappa_Q=\operatorname{res}_Q\kappa\),
\(b_Q=\operatorname{res}_Q b\), and sphere evaluation of \(b_Q\)
has degree two. There is an injected order-two subgroup
\[
 C_{12}/2C_{12}\lhook\joinrel\longrightarrow\pi_4H_Q,\qquad
 {}[\nu']\longmapsto e_Q=\operatorname{res}_Q e\ne0.
 \tag{RT.40.3}\label{p1b:eq:rt-40-3}
\]
Restriction kills \(z,t,2\kappa\). We do not need to select the
extension type of the whole \(\pi_4H_Q\).
\end{lemma}
\begin{proof}
Let \(S_Q\) be the pointed gauge group. Kill the circle one-cell in
the Moore model \(Q=S^1\cup_2e^2\). The remaining cone sphere has
coaction \(c\mapsto c+2u\). The same component translation as in the
current circle-cofibre proof, now with no lens three-cell, gives the
pointed coefficient segments
\[
 0\longrightarrow
 \operatorname{coker}(2:\pi_{r+2}G\to\pi_{r+2}G)
 \longrightarrow\pi_rS_Q
 \longrightarrow\ker(2:\pi_{r+1}G\to\pi_{r+1}G)
 \longrightarrow0,\quad G=SU(2).
 \tag{RT.40.4}\label{p1b:eq:rt-40-4}
\]
In particular \(\pi_1S_Q=\pi_2S_Q=C_2\), and
\(\pi_3S_Q\) has order four, with an injected bottom
\(C_2\{\eta_3\eta_4\}\).

The evaluation connecting maps are the sphere-clutching commutator
followed by the Moore pinch. The retained actual operation
\(A_*\xi=\eta_3E\xi\) gives
\[
 \delta_3(\iota_3)=[\eta_3]\ne0,\quad
 \delta_4(\eta_3)=[\eta_3\eta_4]\ne0,\quad
 \delta_5(\eta_3\eta_4)=[6\nu']=0
       \text{ in }C_{12}/2C_{12}.
 \tag{RT.40.5}\label{p1b:eq:rt-40-5}
\]
The first map forces even sphere evaluation. The second kills the
bottom order-two subgroup of \(\pi_3S_Q\), leaving a kernel of order
two in \(\pi_3H_Q\); its quotient is the free \(2\mathbb Z\).
The restriction of the current \(b\), of evaluation degree two,
generates this free quotient. The free quotient splits as a group.
The third map leaves the bottom subgroup of \(\pi_4S_Q\) injected
in the unpointed group, proving \eqref{p1b:eq:rt-40-3}.

The component sequence has zero cokernel in degree two and
\(\ker(2:\pi_1U(2)\to\pi_1U(2))=0\), so \(S_Q\), and then
\(H_Q\), are connected. Evaluation injects \(\pi_1S_Q\), and the
local loop \(q_+\) splits determinant winding. The determinant-zero
loop \(\kappa_Q\) has the nonzero clutched primary
\(a_F\tau_Q\), and hence generates this \(C_2\). This proves the
circle assertion and kills \(2\kappa\). The supported lens-three-cell
representatives of \(z,t\) restrict to the identity on \(Q\).
The natural sphere pullback \(J\nu'\) restricts to the bottom Moore
pinch of coefficient \(\nu'\), whose residue modulo two is nonzero.
\end{proof}

\begin{lemma}[The two doubled classes are independent of the boundary frame]
\label{p1b:rt:lemma40-central}
In \(\pi_0\operatorname {Map}(F,H_Q)\), both \(2q_+\) and \(2b_Q\)
are central. Thus conjugating either by any change of the ordinary
boundary bundle frame preserves its component.
\end{lemma}
\begin{proof}
The scalar circle is pointwise central and has class
\(s=2q_+-\kappa_Q\). The loop \(\kappa_Q\) admits a pointed
representative trivial on the lens one-skeleton: its pointed class
comes from the bottom \(\pi_3G/2\pi_3G\) on the lens two-cell.
An evaluation-zero torsion element of \(\pi_3H_Q\) can be represented
trivial at the lens zero-cell. Its commutator with \(\kappa_Q\)
therefore has lens filtration at least \(2+1=3\), outside \(Q\).
This statement follows from a cellular diagonal; conjugation
transition maps preserve the relative identity conditions.
For the free generator \(b_Q\), the commutator with \(\kappa_Q\)
is the restriction of the already proved
\(\langle\kappa,b\rangle=0\) in the actual lens group.
Consequently \(\langle 2q_+,v\rangle=0\) for every
\(v\in\pi_3H_Q\).

Every circle class is a combination of \(q_+\) and \(\kappa_Q\).
The commutator of \(2b_Q\) with \(q_+\) is the restriction of
\(2\langle q_+,b\rangle\), which is zero because
\(\pi_4H=C_2^2\). Its commutator with \(\kappa_Q\) is zero by
the preceding actual lens calculation.

A map from \(F=S^1\times S^3\) to the connected topological group
\(H_Q\) is represented by a product of circle, sphere, and parameter
four-cell factors. The only commutator that can occur on \(F\) is
the circle--sphere Samelson product just computed. Two sphere factors
have parameter smash dimension six; a four-cell factor with a
positive-dimensional factor has dimension at least five. These
exceed the dimension of \(F\). There is no normal component of
\(H_Q\) to add a further conjugation. This proves centrality for
arbitrary boundary frame changes, not only a chosen split frame.
\end{proof}

\subsubsection{A four-dimensional model inside the actual free region}

Let \(\mathcal C\subset W\) be the compactified unordered
configurations whose two circle coordinates coincide:
\[
 \mathcal C=
 \bigl\{\{(s,y),(s,y')\}:y\ne y'\text{ in }S^3\bigr\},
 \qquad\partial\mathcal C=R_{\mathcal C}.
 \tag{RT.40.6}\label{p1b:eq:rt-40-6}
\]
It is a proper seven-manifold; a homotopy model of it has dimension
four. A compatible radial tubular neighborhood may be used for the
diagonal excision, without changing \(W,R\) or the restriction maps.

\begin{lemma}[The coincident-circle retraction and transferred classes]
\label{p1b:rt:lemma40-retraction}
There is an exchange-equivariant retraction of the ordered
coincident-circle configuration onto antipodal sphere pairs. Hence
\[
 \mathcal C\simeq X=S^1\times\mathbb {RP}^3.
 \tag{RT.40.7}\label{p1b:eq:rt-40-7}
\]
On its boundary the retraction is
\[
 \Phi(s,y,[n])=(s,[yn]),\qquad
 n\in S^2\subset\operatorname {Im}\mathbb H.
 \tag{RT.40.8}\label{p1b:eq:rt-40-8}
\]
The cover on \(X\) is \(S^1\times S^3\to S^1\times\mathbb {RP}^3\).
Let \(v_{\mathcal C}\) be the integral projective-three-space top
class, whose cover is \(2u\). The retained transfers restrict as
\[
 h|_{\mathcal C}=2a_F,\qquad
 q|_{\mathcal C}=v_{\mathcal C},\qquad
 c_4|_{\mathcal C}=a_Fv_{\mathcal C}.
 \tag{RT.40.9}\label{p1b:eq:rt-40-9}
\]
In particular \(\bar c_4|_{\mathcal C}=\bar a_F\chi^3\ne0\)
generates \(H^4(X;C_2)=C_2\).
Every compatible actual \(V|_{\mathcal C}\) is ordinarily isomorphic
to \(\mathbf1\oplus\lambda_{\mathcal C}\).
\end{lemma}
\begin{proof}
Regard the sphere as unit vectors in \(\mathbb R^4\). Put
\(m=(y+y')/2\), \(d=(y-y')/2\); then \(m\cdot d=0\),
\(|m|^2+|d|^2=1\), and \(d\ne0\). For \(0\le t\le1\) set
\[
 m_t=(1-t)m,\qquad
 d_t=\sqrt{1-|m_t|^2}\,d/|d|,\qquad
 (y_t,y'_t)=(m_t+d_t,m_t-d_t).
 \tag{RT.40.10}\label{p1b:eq:rt-40-10}
\]
Both points remain on \(S^3\), their separation does not decrease,
and they end at the antipodal pair \((d/|d|,-d/|d|)\).
Exchange sends \(d_t\) to \(-d_t\). The formula therefore descends
to the stated retraction and respects diagonal excision.
Boundary sphere points have form \(y\exp(\pm rn)\), so their
normalized difference is \(yn\). This proves \eqref{p1b:eq:rt-40-8}, including
the actual boundary cover.

On this ordered slice \(a_t=a_E=a_F\). At the antipodal retract the
two sphere orientations agree, since the antipodal map of \(S^3\)
has degree one. The product-cover transfer gives
\(\operatorname {tr}(u)=v_{\mathcal C}\) and obeys the projection
formula. Using the retained definitions
\(h=\operatorname {tr}(a_t)\),
\(q=\operatorname {tr}(u_x)\), and
\(c_4=\operatorname {tr}(a_tu_x)\) proves \eqref{p1b:eq:rt-40-9}.
There is no undetected \(a_F\tau\) term in \(q|_{\mathcal C}\):
this transfer is pulled back from the projective-space factor,
independently of the common circle.

On \(X\), the sign-line Chern class is pulled back from
\(\mathbb {RP}^3\), so its square is zero. Thus the actual bundle
has \(c_1=\tau_{\mathcal C}\), \(c_2=0\).
A complex rank-two bundle on a four-dimensional CW complex is
classified by these first two Chern classes: the four-stage of
\(BU(2)\) has groups \(\pi_2=\mathbb Z,\pi_4=\mathbb Z\),
with no intervening group or five-dimensional \(k\)-invariant.
The split bundle has exactly these data. Homotopy invariance transfers
this ordinary isomorphism from the retract to \(\mathcal C\).
It need not agree with the prescribed boundary frame; that issue is
handled by Lemma~\ref{p1b:rt:lemma40-central}.
\end{proof}

\begin{lemma}[The whole boundary image at zero primary on the slice]
\label{p1b:rt:lemma40-kernel}
Every determinant-one gauge on \(V|_{\mathcal C}\) of primary zero
has identity boundary component on \(R_{\mathcal C}\).
\end{lemma}
\begin{proof}
The retract \(X\) has CW dimension four. Once the primary difference
of a determinant-one section is fixed to zero, the remaining section
classes are a quotient of \(H^4(X;\pi_4G)=C_2\); there are no
higher choices on this homotopy model. Returning lower homotopies
can identify these choices but cannot add any.
The actual \(g_c|_{\mathcal C}\) realizes the nonzero fourth
direction, by \eqref{p1b:eq:rt-40-9}. Thus the identity and this restriction
exhaust the section classes with primary zero.

The retained whole boundary value is
\(g_c|_R=(1+(A\bmod2))t\). Both values restrict to the identity on
\(R_{\mathcal C}\), since \(t\) is supported in the lens three-cell.
The exhaustive section argument therefore proves the claimed whole
boundary image. The conclusion concerns actual section components,
not just their characteristic classes or their six-stage reductions.
\end{proof}

\subsubsection{Evaluation of both natural \(e\) coefficients}

\begin{theorem}[The natural odd \(e\) coefficients vanish]
\label{p1b:rt:theorem40-e}
In the retained natural boundary chart, both shared odd coefficients
are evaluated:
\[
 \boxed{h_e=0,\qquad w_e=0,\qquad \omega=w_e+h_e=0.}
 \tag{RT.40.11}\label{p1b:eq:rt-40-11}
\]
This applies to every odd \(A\), both exotic fifth-bit mates, and the
matched sphere relation in either restriction sector.
The remaining natural odd top coefficients \(h_t,w_t\) are not
evaluated by this detector.
\end{theorem}
\begin{proof}
First work in a split frame of \(V|_{\mathcal C}\) supplied by
Lemma~\ref{p1b:rt:lemma40-retraction}.
The diagonal gauge
\[
 D_2(s,\ell)=\operatorname {diag}(e^{2is},1)
 \tag{RT.40.12}\label{p1b:eq:rt-40-12}
\]
has determinant \(2a_F\), primary \(2a_F\tau_{\mathcal C}=0\),
and boundary \(2q_+\). The primary assertion follows by clutching
the winding-two first eigenline against the fixed second eigenline.

Choose a determinant-one gauge \(K\) on the projective-space factor,
supported in its open three-cell, with primary \(v_{\mathcal C}\).
Its ordered cover has degree two. Pull it back by the retraction.
The boundary family \(\Phi^*K\), independent of the common circle,
has sphere evaluation degree two: at any fixed \(n\), the map
\(y\mapsto[yn]\) has degree two. Its class in \(\pi_3H_Q\)
therefore differs from \(b_Q\) by a torsion element of order two.
Consequently \(K^2\) has primary \(2v_{\mathcal C}\) and exact
boundary component \(2b_Q\), with no parameter-four-cell term.

An arbitrary isomorphism to the split frame changes these boundary
classes by conjugation. Lemma~\ref{p1b:rt:lemma40-central} fixes
both \(2q_+\) and \(2b_Q\). Hence these constructions give actual
gauges on the actual \(V|_{\mathcal C}\) with the indicated low
data and boundary components, independently of the split frame.

Now restrict the actual global \(g_h\). Its determinant becomes
\(2a_F\) and its primary is zero. Its boundary on \(R_{\mathcal C}\)
is \(2q_++h_e e_Q\): the \(a=3z\) and \(t\) terms disappear, and
any even \(2\kappa\) correction also disappears. Divide by \(D_2\).
The determinant can be made exactly one using its zero-winding
logarithm, relative to the prescribed boundary. The resulting gauge
has primary zero and boundary \(h_e e_Q\).
Lemma~\ref{p1b:rt:lemma40-kernel} and \(e_Q\ne0\) force \(h_e=0\).

The actual \(g_2\) has determinant one and primary \(2q\), which
restricts to \(2v_{\mathcal C}\). Its boundary on \(R_{\mathcal C}\)
is \(2b_Q+w_e e_Q\); the entire \(z_Az+w_tt\) correction, except
the displayed \(e\) coefficient, disappears on \(Q\).
Divide by \(K^2\). The resulting primary-zero gauge has boundary
\(w_e e_Q\), which is the identity by the same whole-image lemma.
Thus \(w_e=0\). This proof applies to any compatible actual \(V\);
the already established parity and fifth-bit transports give the
stated common odd values. It also agrees with the evaluated even
\(e\)-coefficients.
\end{proof}

The detector has now evaluated the essential matched degree-seven
class as zero. It has not constructed a complete relative cell
cocycle or supplied the final degree-eight filling. Restriction to
\(Q\) kills the natural \(t\), so it cannot assign \(h_t,w_t\).
The structural quotient consequences are in the following MAIN
supplement; the distinction between this slice extension and an
extension over all of \(W\) is preserved there.



\subsection{The selected odd quotient and the final top frontier}
\label{p1b:sec:rtproof41}

Theorem~\ref{p1b:rt:theorem40-e} gives the actual natural values
\(h_e=w_e=0\). Retain \(z_A=5-2A\bmod6\), \(a=3z\), and the
already evaluated \(j_\eta=t\) in the odd exotic sector. Thus
\[
 r_A=2b+z_Az+w_tt,\qquad j_h=a+h_tt,\qquad j_\eta=t.
 \tag{RT.41.1}\label{p1b:eq:rt-41-1}
\]
The following quotient type is independent of these two top coefficients, because the exotic determinant-one image already kills \(t\). The complete natural comparison later proves \(h_t=w_t=1\).

\begin{theorem}[The odd exotic full quotient is split]
\label{p1b:rt:theorem41-type}
For every odd \(A\) and each exotic fifth-bit bundle, put
\(g_A=b-Az\) in the quotient modulo \(t\). Then
\[
 \boxed{\mathscr D^{\,\det=1}(V,A)
       \text{ is a torsor for }C_{12}\{g_A\}\oplus C_2\{e\},}
 \tag{RT.41.2}\label{p1b:eq:rt-41-2}
\]
and
\[
 \boxed{\mathscr D(V,A)
       \text{ is a torsor for }C_6\{g_A\}\oplus C_2\{e\}.}
 \tag{RT.41.3}\label{p1b:eq:rt-41-3}
\]
The odd-determinant involution on the first torsor is translation by
\(6g_A\). In the full quotient \(3b\) has order two. The alternative
full type \(C_{12}\) is excluded by the geometric calculation.
\end{theorem}
\begin{proof}
Modulo the already extendable \(t\), the actual sphere relation is
\(2b+z_Az=0\), and the auxiliary finite image has no other
generator. Since \(z_A+2A=5\bmod6\),
\[
 2g_A=-(z_A+2A)z=-5z=z,\qquad 12g_A=0.
 \tag{RT.41.4}\label{p1b:eq:rt-41-4}
\]
The triangular presentation gives independent factors \(C_{12}\)
and \(C_2\{e\}\), exactly as in the matched-generator argument of
Subsection~\ref{p1b:sec:rtproof37}. The actual determinant column modulo \(t\) is now
\(a=3z=6g_A\); it has no \(e\) component. Its translation is a
free involution, and quotienting imposes \(6g_A=0\) without
identifying \(e\) with a cyclic carry. This gives \eqref{p1b:eq:rt-41-3}.
Since \(b=(1+2A)g_A\) and \(1+2A\) is odd,
\(3b=3g_A\ne0\) in \(C_6\), of order two.
\end{proof}

After an origin is chosen the twelve exotic full representatives are
\(rg_A+se\), \(0\le r<6\), \(s\in C_2\), with law
\[
 (r,s)+(r',s')=((r+r')\bmod6,s+s').
 \tag{RT.41.5}\label{p1b:eq:rt-41-5}
\]
All previously evaluated counts are retained:
\[
 \begin{array}{c|c|c|c|c}
 &\text{full, sector 0}&\text{full, sector 1}&
   \text{auxiliary, sector 0}&\text{auxiliary, sector 1}\\ \hline
 A\text{ even}&C_{12}\oplus C_2&C_6\oplus C_2&
                   C_{12}\oplus C_2&C_{12}\oplus C_2\\
 A\text{ odd}&C_{12}\oplus C_2^2&C_6\oplus C_2&
                   C_{12}\oplus C_2^2&C_{12}\oplus C_2
 \end{array}
 \tag{RT.41.6}\label{p1b:eq:rt-41-6}
\]
Each entry acts on one fixed compatible-bundle row. There are two
such rows in each restriction sector. Their full totals remain
72 even and 120 odd; the auxiliary totals remain 96 even and 144 odd.
This is the evaluated table of additive quotient types.

\begin{corollary}[All odd exotic fixed-row fibre groups are determined]
\label{p1b:rt:corollary41-fibres}
Write \(h_6=\gcd(6,A)\). Every nonempty ordinary forgetful fibre in
a fixed odd exotic row is a torsor for
\[
 \boxed{C_{6/h_6}\oplus C_2.}
 \tag{RT.41.7}\label{p1b:eq:rt-41-7}
\]
For odd \(A\) divisible by three, this is \(C_2\oplus C_2\),
not \(C_4\). For odd \(A\) not divisible by three it is
\(C_6\oplus C_2\). All other fixed-row fibre groups remain exactly
those proved in Subsection~\ref{p1b:sec:rtproof37}; no unresolved quotient-type branch remains.
\end{corollary}
\begin{proof}
The ordinary top operation is killed in every odd row. The primary
of \(g_A\) is \(1-2A\), hence one modulo \(h_6\).
The homogeneous forgetful action therefore has kernel generated by
\(h_6g_A\), of order \(6/h_6\), and the independent \(e\).
The retained fibre-action proof makes each fibre a torsor for this
kernel. Since \(A\) is odd, \(h_6\) is either one or three.
The odd sector-zero ordinary fibre still comprises two disjoint
row torsors; no group on their union is asserted.
\end{proof}

\subsubsection{What remains at degree eight}

Define \(\zeta_{\rm nat}=w_t+h_t\), the matched top comparison rather than
the local circle \(2\kappa\), and
\(M=3z+w_tt=3\widehat z\), where
\(\widehat z=z+w_tt\). The matched boundary tests are
\[
 B_{d,p}=2q_++p\kappa+\gamma(M+d e),\qquad d,p\in C_2.
 \tag{RT.41.8}\label{p1b:eq:rt-41-8}
\]

\begin{theorem}[The completed degree-seven test and remaining top tests]
\label{p1b:rt:theorem41-stages}
For odd exotic \(V\), the essential matched degree-seven obstruction
is now evaluated:
\[
 \boxed{O_7(B_{d,p})=d\Gamma_\chi.}
 \tag{RT.41.9}\label{p1b:eq:rt-41-9}
\]
Thus \(B_{0,0}\) and \(B_{0,1}\) extend through
\(W^{(7)}\cup R\), while both \(d=1\) candidates fail there.
For \(d=0\) the final obstruction is
\[
 O_8(B_{0,p})=(p+\zeta_{\rm nat})\Theta,\qquad
 H^8(W,R;C_2)=C_2\{\Theta\}.
 \tag{RT.41.10}\label{p1b:eq:rt-41-10}
\]
Exactly one of these two candidates extends over all of \(W\);
the final natural comparison proves \(\zeta_{\rm nat}=0\), so the extending candidate is \(B_{0,0}\).

For the natural determinant and sphere candidates
\[
 D_{e_0,t_0}=2q_++\gamma(a+e_0e+t_0t),\qquad
 S_{e_0,t_0}=\gamma(2b+z_Az+e_0e+t_0t),
 \tag{RT.41.11}\label{p1b:eq:rt-41-11}
\]
the corresponding stages are
\[
 \begin{array}{c|c|c}
 &O_7&O_8\text{ when }e_0=0\\ \hline
 D_{e_0,t_0}&e_0\Gamma_\chi&(t_0+h_t)\Theta\\
 S_{e_0,t_0}&e_0\Gamma_\chi&(t_0+w_t)\Theta
 \end{array}
 \tag{RT.41.12}\label{p1b:eq:rt-41-12}
\]
The natural values \(h_t,w_t\) are two shared odd bits. Neither is
assigned by \eqref{p1b:eq:rt-41-12}.
\end{theorem}
\begin{proof}
Substitute the actual value \(\omega=0\) into the separated relative
obstructions of Subsection~\ref{p1b:sec:rtproof39}. For the natural tests divide by the actual
global \(g_h\) or \(g_2\), respectively, adjusting any even local
\(2\kappa\) discrepancy by the known fourth gauge. The residual pure
finite class has \(e\) coefficient \(e_0\) and the displayed top
coefficient. Subsections~\ref{p1b:sec:rtproof38}--\ref{p1b:sec:rtproof39} give its stage-seven and final obstructions.
All odd pure finite lower-choice top corrections vanish, so there
is no omitted final indeterminacy.
\end{proof}

The \(e_0=1\) natural determinant candidates fail at degree seven. The complete original natural comparison below gives \(h_t=w_t=1\) and \(\zeta_{\rm nat}=0\), deciding both degree-eight alternatives. The resulting inherited ledger is 4 YES/28 NO and the natural target ledger 4 YES/10 NO. The absolute equality \(H^8(W;C_2)=0\) does not replace the relative final group in \eqref{p1b:eq:rt-41-10}.



\subsection{The framed proper-band Thom map}
\label{p1b:sec:rtproof42}

Retain the actual free region \(W\), its boundary
\(R=F\times L\), \(F=S^1\times S^3\), and \(L=\mathbb {RP}^3\).
Let \(C\) be the compactified unordered configuration band of two
circle points, with both sphere coordinates fixed at \(y_0\).
Thus \(C\) is a M\"obius band. Its boundary is the fixed common circle
at \(y_0\) and at the circle-axis point of the normal lens.
Let \(\gamma_C\) be the real line of its ordering double cover.
The boundary frame in this subsection is the prescribed product frame:
the sphere tangent frame at \(y_0\), and the positive-circle-axis chart
in the normal lens. A change of this relative framing is not allowed.

\begin{lemma}[The normal bundle with its boundary frame]
\label{p1b:rt:lemma42-normal}
The actual normal bundle and its specified boundary trivialization are
\[
 \boxed{\nu_C=\mathbb R^3\oplus3\gamma_C.}
 \tag{RT.42.1}\label{p1b:eq:rt-42-1}
\]
Capping \(C\) by a disc using this trivialization gives
\[
 \widehat C=\mathbb {RP}^2,\qquad
 \widehat\nu=\mathbb R^3\oplus3\gamma,\qquad
 w_1(\widehat\nu)=\chi,\quad w_2(\widehat\nu)=\chi^2\ne0.
 \tag{RT.42.2}\label{p1b:eq:rt-42-2}
\]
Although \(\nu_C\) is isomorphic to
\(\mathbb R^5\oplus\gamma_C\) as an absolute bundle on \(C\),
this replacement does not preserve the prescribed boundary frame.
\end{lemma}
\begin{proof}
Before dividing by ordering, the six normal coordinates are the
sphere displacements \((v_1,v_2)\) in
\(T_{y_0}S^3\oplus T_{y_0}S^3\).
Their sum and difference give three exchange-fixed coordinates and
three exchange-negative coordinates. This proves \eqref{p1b:eq:rt-42-1} on the
actual configuration band.

Near its boundary write the circle points as \(m+r,m-r\), choosing
the first point by the positive circle separation. The sum coordinate
is the diagonal sphere displacement. The difference coordinate,
divided by the positive separation, is the sphere coordinate in the
positive-circle-axis lens chart. Positive scaling changes no framing
class. Hence the displayed normal splitting induces exactly the
prescribed product boundary frame.

The ordering line extends over the capped band as the nontrivial
line \(\gamma\) on \(\mathbb {RP}^2\). In the displayed trivialization
the three sum coordinates extend as constant real lines and each of
the three difference coordinates extends as \(\gamma\). Thus the
capped normal bundle is the one in \eqref{p1b:eq:rt-42-2}. The Whitney formula gives
\[
 w(\widehat\nu)=(1+\chi)^3=1+\chi+\chi^2
       \quad\text{in }H^*(\mathbb {RP}^2;C_2).
 \tag{RT.42.3}\label{p1b:eq:rt-42-3}
\]
The capped Thom space also has the concrete stunted-projective model
\(\operatorname {Th}(\widehat\nu)
  \simeq\Sigma^3(\mathbb {RP}^5/\mathbb {RP}^2)\).
The capped \(\mathbb R^5\oplus\gamma\) has second class zero.
An isomorphism of the two bundles on \(C\) preserving the indicated
boundary frames would extend across the cap and contradict this
second-class calculation.
\end{proof}

Write \(P^7(2)=S^6\cup_2 e^7\) and let \(i:S^6\to P^7(2)\)
be its bottom inclusion. Use the product boundary splitting
\[
 \operatorname {Th}(\nu_C|_{\partial C})
       =S^6\wedge S^1_+\simeq S^6\vee S^7.
 \tag{RT.42.4}\label{p1b:eq:rt-42-4}
\]
The top summand is chosen to have zero projection under collapse of
the base circle to a point; this specifies the splitting with its
framing. Its generator is not subsequently shifted by a bottom
\(\eta_6\).

\begin{theorem}[The actual boundary top map]
\label{p1b:rt:theorem42-top}
In this splitting the inclusion of the boundary Thom space has
the evaluated values
\[
 \boxed{k|_{S^6}=i,\qquad k|_{S^7}=i\eta_6\ne0.}
 \tag{RT.42.5}\label{p1b:eq:rt-42-5}
\]
Consequently, for every sixth-stage surgery map
\(f:P^7(2)\to BU(2)\) with \(fi=\alpha_6\),
\[
 \boxed{fk=(\alpha_6,\alpha_6\eta_6),\qquad
       \Omega(\alpha_6\eta_6)
          =\eta_3\eta_4\eta_5=6\nu'\ne0\in C_{12}.}
 \tag{RT.42.6}\label{p1b:eq:rt-42-6}
\]
These are boundary bundle insertions, before their boundary
nullhomotopy. They are not a nonzero class of the completed
compatible boundary bundle.
\end{theorem}
\begin{proof}
The absolute Thom space over the band has the two-cell model
\(P^7(2)\): the orientation of the rank-six normal bundle reverses
around its core, giving the degree-two attachment.
Capping adds one eight-cell. In the product framing its attaching
map is the image of the top summand in \eqref{p1b:eq:rt-42-4}.
Indeed the projection to the fibre \(S^6\) extends over the cap
disc before this projection is collapsed, so the cap uses the
specified top summand without an additional bottom Hopf shift.

Let \(U\) be the degree-six mod-two Thom class of
\(\widehat\nu\). The Thom identity and \eqref{p1b:eq:rt-42-2} give
\[
 Sq^2 U=\chi^2U\ne0\in
       H^8(\operatorname {Th}(\widehat\nu);C_2).
 \tag{RT.42.7}\label{p1b:eq:rt-42-7}
\]
The Moore space is five-connected, so its degree-seven homotopy
group is in the stable range. The degree-two cofibration gives
\(\pi_7P^7(2)=C_2\{i\eta_6\}\).
If the eight-cell attachment were zero, the capped Thom space
would be \(P^7(2)\vee S^8\) and \(Sq^2U\) would be zero.
Thus the attachment is the sole nonzero element \(i\eta_6\).
The fibre inclusion supplies the bottom value in \eqref{p1b:eq:rt-42-5}.

Now apply the actual surgery map \(f\). Its restriction to the
bottom is \(\alpha_6\), so its top restriction is necessarily
\(\alpha_6\eta_6\), independently of the choice of extension
of \(\alpha_6\) across \(P^7(2)\).
The retained loop adjoint of \(\alpha_6\) is
\(\eta_3\eta_4\); composition with the suspended \(\eta_5\)
gives \(6\nu'\) as in the current adjoint-operation owner.
This coefficient is nonzero in the actual \(C_{12}\).
The boundary compatibility nullhomotopy is the separate construction
in the following MAIN supplement.
\end{proof}

\begin{corollary}[The local relative source retains an eight-cell]
\label{p1b:rt:corollary42-relative}
The Thom space relative to the boundary of the band is
\[
 \operatorname {Th}(\nu_C,\nu_C|_{\partial C})
       \simeq S^7\vee S^8.
 \tag{RT.42.8}\label{p1b:eq:rt-42-8}
\]
This describes the local relative source. It does not assert
that the surgery bundle map has a nullhomotopy inside the
boundary normal tube.
The relative Thom collapse \(c:W/R\to
\operatorname {Th}(\nu_C,\nu_C|_{\partial C})\) pulls its
mod-two top generator \(U_{\mathrm{rel}}\) back to
\[
 \boxed{c^*U_{\mathrm{rel}}=\Theta\ne0.}
 \tag{RT.42.9}\label{p1b:eq:rt-42-9}
\]
\end{corollary}
\begin{proof}
Cap the band in the prescribed frame and then collapse the
Thom space over the cap disc. By excision this is the relative
Thom space in the statement. The cap-disc Thom space retracts
to the bottom fibre \(S^6\). The capped Thom space has cells
six, seven, eight, with attachments \(2\) and \(i\eta_6\).
Collapsing its bottom \(S^6\) kills both attachments, leaving
the stated wedge.
For the last assertion the Thom isomorphism identifies
\(U_{\mathrm{rel}}\) with the mod-two relative fundamental
cohomology class of the proper band. The collapse has local
degree one on the interior of its normal tube.
Its pairing with \([W,R]_2\) is therefore one.
The retained \(H^8(W,R;C_2)=C_2\{\Theta\}\) then gives
\eqref{p1b:eq:rt-42-9}.
\end{proof}

The positive result sharpens the zero-or-Hopf alternative in
Lemma~\ref{p1b:rt:lemma34-surgery}. It retains the same surgery and
the same bundles. The relative framing defect, including an explicit
rotation representative, is recorded in
Section~\ref{p1b:app:rt43-framing}.
Neither natural degree-eight gauge coefficient is evaluated by
\eqref{p1b:eq:rt-42-6}.



\subsection{The top boundary comparison required by the surgery}
\label{p1b:sec:rtproof43}

Retain the framed surgery of
Theorem~\ref{p1b:rt:theorem42-top} and the actual circle-cofibre
nullhomotopies of Lemma~\ref{p1b:rt:lemma34-surgery}.
Put \(x=\bar a_F\), \(u=\bar u_F\).
The mod-two boundary ring is
\[
 H^*(R;C_2)=C_2[x,u,\chi]/(x^2,u^2,\chi^4),
 \quad |x|=1,\quad |u|=3,\quad|\chi|=1.
 \tag{RT.43.1}\label{p1b:eq:rt-43-1}
\]
A comparison cochain here belongs to a homotopy between two
boundary bundle classifying maps. It is not an independent class
of an actual boundary gauge. Its later extensions must be retained.

\begin{theorem}[A specified null-comparison of both insertions]
\label{p1b:rt:theorem43-comparison}
The two forced insertions in \eqref{p1b:eq:rt-42-6} admit simultaneous
determinant-preserving boundary null-comparisons.
Choose the bottom comparison independent of the common circle,
and use the natural parameter-top pinch for the additional top
comparison. Their first nonzero comparison inputs are
\[
 \boxed{\beta_4=u\chi\quad(\eta_3),\qquad
        \beta_5^{\mathrm{top}}=xu\chi
                                  \quad(\eta_3\eta_4).}
 \tag{RT.43.2}\label{p1b:eq:rt-43-2}
\]
The second formula specifies the additional comparison for the
forced parameter-top insertion. It does not select all degree-five
or higher cochains of the bottom nullhomotopy.
\end{theorem}
\begin{proof}
In the normal chart around the boundary circle of \(C\), the
six normal coordinates are the three sphere-parameter coordinates
and the three lens-chart coordinates. Collapsing that chart gives
the bottom \(S^6=S_F^3\wedge S_L^3\); adjoining the common-circle
top cell gives \(S^7=S_F^4\wedge S_L^3\).
Thus the bottom and top components of \eqref{p1b:eq:rt-42-6} have respectively
parameter degrees three and four.

Use the actual pointed circle-cofibre in the
Chern-class-one boundary component. Its successive coefficient
values, with their chosen input nullhomotopies, are
\[
 \eta_3\longmapsto(0,\eta_3\eta_4),\qquad
 \eta_3\eta_4\longmapsto(0,6\nu').
 \tag{RT.43.3}\label{p1b:eq:rt-43-3}
\]
The first input supplies the bottom bundle null-comparison:
on \(S_F^3\times L^{(1)}\) its value is \(\eta_3\).
On the lens two-cell its two copies paste by \(2\eta_3=0\);
the actual lens top coaction supplies \(\eta_3\eta_4\).
This is precisely the retained comparison \(\beta_4=u\chi\).
It can be constructed on the sphere parameter and pulled back
independently of the common circle.

Theorem~\ref{p1b:rt:theorem42-top} now forces the extra
parameter-four bundle component with loop value \(6\nu'\).
The second input of \eqref{p1b:eq:rt-43-3} supplies its null-comparison.
On \(S_F^4\times L^{(1)}\) the value is
\(\eta_3\eta_4\); its two copies paste by
\(2\eta_3\eta_4=0\), and the actual lens top coaction gives
exactly \(6\nu'\). Pull this comparison back by the natural
pinch \(F\to F/(S^1\vee S^3)=S_F^4\).
Its first cochain is therefore \(xu\chi\), as asserted.
It is trivial on the lower parameter cells and does not change
the chosen \(\beta_4\).

The bottom comparison may also contain a degree-five term
on \(S_F^3\times L^{(2)}\), with no common-circle factor,
as well as higher comparison choices. No value for those
terms is used or asserted here. The two comparisons just
constructed give a complete boundary compatibility homotopy
by the actual circle-cofibre construction, not by vanishing
of Chern classes alone.
\end{proof}

\begin{corollary}[The determinant-one transport remains exact]
\label{p1b:rt:corollary43-transport}
The evaluated top attachment adds no omitted term to the
determinant-one transport formula:
\[
 \boxed{nq+p\eta_W\quad\longmapsto\quad
        \text{the original normalized boundary class}+p\,t.}
 \tag{RT.43.4}\label{p1b:eq:rt-43-4}
\]
In particular it preserves the already matched sphere-two
relation between the two restriction sectors.
\end{corollary}
\begin{proof}
For a determinant-one gauge the first positive-dimensional
difference is the degree-three primary. Its boundary mod-two value
is \(p x\chi^2\), since \(q|_R=2u_F\).
The fourth comparison gives
\[
 (u\chi)(p x\chi^2)=p xu\chi^3,
 \tag{RT.43.5}\label{p1b:eq:rt-43-5}
\]
which inserts \(p\) copies of
\(\langle\eta_3,\iota_3\rangle=\mu\), hence \(p\,t\).
The new fifth comparison can first interact with that gauge
in degree \(5+3=8\), beyond the seven-dimensional boundary.
Its specific displayed product is also zero because \(x^2=0\).
All later comparisons are still higher-dimensional.
This proves \eqref{p1b:eq:rt-43-4} with the forced top insertion retained.
\end{proof}

\subsubsection{The remaining whole-region test}

The new computations select the proper-band top map and its
additional boundary null-comparison. They do not give a gauge
with prescribed boundary \(D_{0,0}\) or \(S_{0,0}\) on all of \(W\).
For the odd natural chart the intermediate whole-region tests are
\[
 \begin{aligned}
 D_{0,t_0}&=2q_++\gamma(a+t_0t),&
 O_8(D_{0,t_0})&=(t_0+h_t)\Theta,\\
 S_{0,t_0}&=\gamma(2b+z_Az+t_0t),&
 O_8(S_{0,t_0})&=(t_0+w_t)\Theta.
 \end{aligned}
 \tag{RT.43.6}\label{p1b:eq:rt-43-6}
\]
Here \(z_A=5-2A\bmod6\) and
\(H^8(W,R;C_2)=C_2\{\Theta\}\).
Both candidates with \(t_0=0\) already pass degree seven,
by the evaluated \(h_e=w_e=0\).
The complete natural calculation below evaluates their actual relative
eight-cell gauge filling with these boundary representatives, carrying
the gauge comparison as well as the boundary bundle insertion.

The complete source-bound natural evaluations are given in Theorem~\ref{p1b:thm:natural-rank-two-completion}; the evaluated counts, quotient groups and ordinary fibres are summarized in Theorem~\ref{p1b:thm:complete-mixed-rank-two}.



\subsection{The boundary normal-frame loop and the finite generators}
\label{p1b:sec:rtproof44}

Retain the prescribed product coordinates
\(R=F\times L\), \(F=S^1_s\times S^3_y\),
\(L=\mathbb {RP}^3\), and the actual boundary bundle
\(E_\partial=\mathbb C\oplus\lambda_L\).
The line \(\lambda_L\) is the complex ordering line, with its
canonical sign-cover description. The boundary gauge group used
below is the section group of this bundle, not an untwisted
mapping group substituted for it.
Write \(Q=\mathbb {RP}^2\subset L\) for the zero-circle-component
hyperplane, so the complement of \(Q\) is the positive-circle-axis
affine three-cell.

Let \(R_s\in SO(4)\) fix the circle-axis normal and one sphere
normal, and rotate the other two sphere normals by \(s\).
This is a based loop, \(0\le s\le2\pi\).
In quaternion coordinates it is
\(n\mapsto u_snu_s^{-1}\), \(u_s=\exp(si/2)\).
Define
\[
 \Phi(s,y,[n])=(s,y,[R_sn]).
 \tag{RT.44.1}\label{p1b:eq:rt-44-1}
\]
The eigenline-preserving bundle identification covering this map
is fixed, rather than chosen up to a gauge: on the ordering cover
it is \((n,v)\mapsto(R_sn,v)\).
It commutes with \((n,v)\sim(-n,-v)\); its endpoint is the identity.

\begin{lemma}[The actual bundle pullback and the retained parameter filtration]
\label{p1b:rt:lemma44-lift}
With this identification, \(\Phi^*\) is an automorphism of the
actual boundary gauge homotopy group. It fixes
\(q_+,q_-\), and \(\kappa=q_+-q_-\).
On the normalized parameter-top subgroup
\(\pi_4H=C_2\{e\}\oplus C_2\{t\}\) it is the identity.
It fixes \(2b\), without requiring a choice of its value on \(b\).
\end{lemma}
\begin{proof}
The sign-cover lift proves the bundle statement and preserves both
eigenlines. The diagonal eigenline phases defining \(q_+,q_-\)
are therefore fixed as actual gauges.

Here \(H\) is the retained identity-component gauge space on \(L\).
For each \(s\), the rotation and its specified lift give a based
automorphism \(a(s)\) of \(H\); \(a(0)=a(2\pi)=1\).
A parameter-top class is represented by a based map
\(S^1\wedge S^3\to H\).
The difference between applying \(a(s)\) and doing nothing
factors through
\[
 S^1\wedge S^3\longrightarrow
 S^1\wedge S^1\wedge S^3,
\]
using the reduced diagonal of the same parameter circle.
That diagonal \(S^1\to S^1\wedge S^1=S^2\) is null.
This argument takes place in the actual section group: the
difference \(a(s)(g)g^{-1}\) is identity when either input is
based. Thus both \(e\) and \(t\) are fixed.

For a based sphere representative of \(b\), the original and
pulled-back gauges agree on the parameter sphere at \(s=0\)
and on the parameter circle at the sphere basepoint.
Their difference is consequently a parameter-top class in
\(\pi_4H\). This group has exponent two, so doubling gives
\(\Phi^*(2b)=2b\).
This does not calculate the individual \(e,t\) coefficients
of \(\Phi^*b-b\).
\end{proof}

\begin{theorem}[The evaluated finite shear]
\label{p1b:rt:theorem44-shear}
For the natural generators of the retained boundary chart,
\[
 \boxed{\Phi^*z=z+t,\qquad
        \Phi^*e=e,\qquad \Phi^*t=t,\qquad \Phi^*(2b)=2b.}
 \tag{RT.44.2}\label{p1b:eq:rt-44-2}
\]
In particular, on
\(\mathscr A_{\mathrm{even}}
 =2\mathbb Z\{b\}\oplus C_6\{z\}\oplus C_2\{e,t\}\),
the complete action is
\[
 \boxed{\Phi^*(2nb+kz+\epsilon e+\tau t)
  =2nb+kz+\epsilon e+(\tau+k\bmod2)t.}
 \tag{RT.44.3}\label{p1b:eq:rt-44-3}
\]
It is an involution, and sends \(a=3z\) to \(a+t\).
\end{theorem}
\begin{proof}
The natural \(z=j_6\nu'\) is represented by the lens-three-cell
insertion, with sphere parameter \(S^3_y\), and is independent
of \(s\). The gauge is identity off that cell.
Trivialize \(\lambda_L\) there using the positive-circle-axis
ordering representative. The lift in \eqref{p1b:eq:rt-44-1} carries this
trivialization to itself, so it introduces no additional
conjugation or sign.

On that affine cell, the loop acts by a full plane rotation
of \(\mathbb R^3\). The top difference between its action on
the one-point compactification and the constant identity
family is the unstable \(J\)-class
\[
 J([\text{full plane rotation}])=\eta_3
       \in\pi_4S^3.
 \tag{RT.44.4}\label{p1b:eq:rt-44-4}
\]
This is the disc-action definition of \(J\). The rotating
two-plane clutching has Hopf invariant one in dimension two;
adding the fixed third coordinate suspends it to \(\eta_3\).
Smashing with the sphere parameter suspends three more times,
giving \(\eta_6:S^7\to S^6\). Postcomposition with the retained
coefficient \(\nu'\) therefore gives
\[
 j_7(\nu'\eta_6)=j_7\mu=t\ne0.
 \tag{RT.44.5}\label{p1b:eq:rt-44-5}
\]
The insertion is lens-top supported, so it has no \(e\)
component. These constructions use the same cell and local
ordering trivialization as the natural \(j_6,j_7\) insertions.
They prove \(\Phi^*z=z+t\) in the actual twisted gauge space.

Lemma~\ref{p1b:rt:lemma44-lift} supplies the other values.
The retained even-circle subgroup is central, so the
normalized additive chart has no extra circle carry in this
calculation. The killed relation \(6z=0\) is respected because
\(6t=0\). Equation \eqref{p1b:eq:rt-44-3} follows by addition.
Applying it twice adds \(2kt=0\).
Geometrically the double loop is also null in \(\pi_1SO(3)=C_2\).
Finally \(3t=t\), proving the assertion about \(a\).
\end{proof}

\begin{corollary}[The framed-band defect realizes this boundary operation]
\label{p1b:rt:corollary44-band}
The boundary restriction of the explicit absolute normal
isomorphism in Section~\ref{p1b:app:rt43-framing}, extended
using the prescribed product sphere frame, is the operation
\eqref{p1b:eq:rt-44-1} on lens directions.
\end{corollary}
\begin{proof}
That isomorphism is
\(A(\theta)=\operatorname{diag}(I_3,R(\theta/2),1)\)
in diagonal-sphere and difference-sphere coordinates.
The boundary covers the core twice. Hence \(s=\theta/2\)
runs from zero to \(2\pi\): the diagonal sphere coordinates
are fixed and two difference-sphere coordinates make a full
rotation. The circle-axis normal is fixed.
Dividing the difference coordinates by the positive circle
separation identifies precisely the affine lens action in
\eqref{p1b:eq:rt-44-1}. The product sphere frame extends this representative
over \(y\).
\end{proof}

This is a transport of boundary normal coordinates. It does
not assert that \(\Phi\) extends over \(W\) with the same
boundary inclusion. The restriction image must be transported
along with the coordinates. The consequence for the actual
relation image is the next MAIN supplement.



\subsection{A geometric matched finite generator and the invariant top test}
\label{p1b:sec:rtproof45}

Work at odd \(A\), retaining the actual natural values
\(h_e=w_e=0\) and
\[
 j_h=3z+h_tt,\qquad j_\eta=t,\qquad
 r_A=2b+z_Az+w_tt,\qquad z_A=5-2A\bmod6.
 \tag{RT.45.1}\label{p1b:eq:rt-45-1}
\]
The determinant column is the odd exotic column.
The sphere coefficient is shared by the retained odd rows.
Neither natural bit \(h_t,w_t\) is assigned here.
Put \(\zeta_{\rm nat}=h_t+w_t\).

For \(\sigma\in C_2\), define the transported restriction map
\[
 \operatorname{res}_\sigma
   =(\Phi^*)^\sigma\circ\operatorname{res},\qquad
 I_\sigma=(\Phi^*)^\sigma I.
 \tag{RT.45.2}\label{p1b:eq:rt-45-2}
\]
Thus membership of a boundary class \(c\) in \(I\) is equivalent
to membership of \((\Phi^*)^\sigma c\) in \(I_\sigma\).
This is an actual reparametrization of the restriction problem,
not an extension theorem for \(c\) with the original inclusion.

\begin{theorem}[Simultaneous top transport]
\label{p1b:rt:theorem45-bits}
The exact transported columns in fixed additive coordinates are
\[
 \boxed{j_h^\sigma=3z+(h_t+\sigma)t,\qquad
        j_\eta^\sigma=t,\qquad
        r_A^\sigma=2b+z_Az+(w_t+\sigma)t.}
 \tag{RT.45.3}\label{p1b:eq:rt-45-3}
\]
Consequently
\[
 \boxed{h_t^\sigma=h_t+\sigma,\quad
        w_t^\sigma=w_t+\sigma,\quad
        \zeta_{\rm nat}^\sigma=\zeta_{\rm nat}.}
 \tag{RT.45.4}\label{p1b:eq:rt-45-4}
\]
The matched finite generator
\(\widehat z=z+w_tt\) is geometrically realized as
\((\Phi^*)^{w_t}z\).
On the even-sphere subgroup this realizes the matched
relation chart used in Subsection~\ref{p1b:sec:rtproof37}.
\end{theorem}
\begin{proof}
The coefficient three in \(j_h\) and every \(z_A\) are odd.
Apply Theorem~\ref{p1b:rt:theorem44-shear} to \eqref{p1b:eq:rt-45-1}.
The \(2b\) term is fixed; each odd \(z\) coefficient acquires
one \(t\). The two added bits cancel in their sum.

The shear realizes \(z+t\) by a prescribed normal rotation,
and realizes \(z\) by the product frame. Thus either possible
\(\widehat z\) is an admissible finite generator from a
geometric frame. Setting \(\sigma=w_t\) in \eqref{p1b:eq:rt-45-3} gives
\[
 r_A^\sigma=2b+z_Az,\qquad
 j_h^\sigma=3z+\zeta_{\rm nat} t,\qquad j_\eta^\sigma=t.
 \tag{RT.45.5}\label{p1b:eq:rt-45-5}
\]
Equivalently, within the original problem one uses
\(\widehat z=z+w_tt\), obtaining
\(r_A=2b+z_A\widehat z\) and
\(j_h=3\widehat z+\zeta_{\rm nat} t\).
The choice \(\sigma=w_t\) is an existence and normalization
statement. It does not tell which frame is sphere-adapted
until the natural \(w_t\) is evaluated.
No assertion that \(\Phi^*b=b\) is needed or made.
\end{proof}

\begin{corollary}[Two conditional framing orbits]
\label{p1b:rt:corollary45-orbits}
The four conditional bit pairs have two orbits under the
actual boundary frame operation:
\[
 \begin{array}{c|c}
 \zeta_{\rm nat}&\{(h_t,w_t)\}\\\hline
 0&\{(0,0),(1,1)\}\\
 1&\{(0,1),(1,0)\}.
 \end{array}
 \tag{RT.45.6}\label{p1b:eq:rt-45-6}
\]
There is exactly one representative with \(w_t=0\) in each
orbit, namely \((\zeta_{\rm nat},0)\).
\end{corollary}
\begin{proof}
Use \eqref{p1b:eq:rt-45-4}. The operation has no fixed bit pair, preserves
the sum, and flips the second coordinate. The displayed
orbits and unique representative follow.
\end{proof}

The four pairs in this corollary are descriptions of possible presentation coordinates, not four additional compatible bundles. Framing-orbit passage alone does not evaluate their sum; the complete original-map proof below gives \(h_t=w_t=1\) and \(\zeta_{\rm nat}=0\).

\begin{theorem}[Equivariance of the exact extension profiles]
\label{p1b:rt:theorem45-profiles}
Retain the specified natural boundary candidates
\[
 D_{\epsilon,\tau}=2q_++\gamma(3z+\epsilon e+\tau t),\qquad
 S_{\epsilon,\tau}=\gamma(2b+z_Az+\epsilon e+\tau t).
 \tag{RT.45.7}\label{p1b:eq:rt-45-7}
\]
In the transported restriction problem their profiles are
\[
 \begin{array}{c|c|c}
 &O_7&O_8\text{ if }\epsilon=0\\\hline
 D_{\epsilon,\tau}&\epsilon\Gamma_\chi&
             (\tau+h_t+\sigma)\Theta\\
 S_{\epsilon,\tau}&\epsilon\Gamma_\chi&
             (\tau+w_t+\sigma)\Theta.
 \end{array}
 \tag{RT.45.8}\label{p1b:eq:rt-45-8}
\]
The original candidate \(D_{\epsilon,\tau}\), respectively
\(S_{\epsilon,\tau}\), is transported to the candidate with
top index \(\tau+\sigma\). Its obstruction coefficients
are consequently unchanged by transport.

Put \(M^\sigma=3z+(w_t+\sigma)t\), and
\[
 B_{d,p}^\sigma=2q_++p\kappa+\gamma(M^\sigma+d e).
 \tag{RT.45.9}\label{p1b:eq:rt-45-9}
\]
Then the matched extension profiles are
\[
 \boxed{O_7(B_{d,p}^\sigma)=d\Gamma_\chi,\qquad
 O_8(B_{0,p}^\sigma)=(p+\zeta_{\rm nat})\Theta.}
 \tag{RT.45.10}\label{p1b:eq:rt-45-10}
\]
The admissible normal-frame loop cannot select the missing
matched \(p\) value.
\end{theorem}
\begin{proof}
Use the exact retained profiles in
Theorem~\ref{p1b:rt:theorem41-stages} and the definition
\eqref{p1b:eq:rt-45-2}. The shear sends each odd \(z\) coefficient's top
index to \(\tau+\sigma\), fixes \(e\), and fixes \(2q_+,\kappa\).
Thus a fixed-coordinate natural profile in \(I_\sigma\)
has the shifted coefficients in \eqref{p1b:eq:rt-45-8}; substituting the
transported top index cancels that shift.
This is also a way to state the obstructions by conjugating
the boundary-coordinate description of the original relative
problem. No extension of \(\Phi\) into \(W\) is used.

For the matched tests substitute \(M^\sigma\) into the
natural determinant profile and add the actual torsion
column \(p j_\eta=p t\).
The final coefficient is
\((w_t+\sigma)+(h_t+\sigma)+p=p+\zeta_{\rm nat}\).
The degree-seven coefficient remains \(d\).
\end{proof}

The retained relative top group is still
\(H^8(W,R;C_2)=C_2\{\Theta\}\).
The new normal-frame action is nontrivial on the natural
boundary coordinates and invisible on the matched top bit.
Evaluating \(\zeta_{\rm nat}\), or either natural \(h_t,w_t\), still
requires the retained actual relative gauge filling.
The quotient and fixed-row fibre group types and counts are
unchanged; no new natural ledger membership follows here.



\subsection{A geometric pushout for the free region}
\label{p1b:sec:rtproof46}

Retain \(K=S^1\times S^3\), the actual free region \(W\),
its boundary \(R=K\times\mathbb {RP}^3\), and its ordering
cover. Write \(W^\circ=\operatorname{Conf}_2(K)/C_2\)
for the unexcised free configuration space. Radial diagonal
excision gives the retained homotopy equivalence \(W\simeq W^\circ\).
This notation does not replace the boundary isotropy or bundle.

Define three free quotient manifolds:
\[
\begin{split}
 Y&=(S^1_a\times S^3_x\times S^3_p)/
             ((a,x,p)\sim(-a,p,x)),\\
 Z_s&=(T^2_{z_1,z_2}\times S^3_w)/
             ((z_1,z_2,w)\sim(z_2,z_1,-w)),\\
 J&=(S^1_a\times S^3_w)/
             ((a,w)\sim(-a,-w)).
\end{split}
\tag{RT.46.1}\label{p1b:eq:rt-46-1}
\]
Here \(Y\) is the retained exchange mapping torus of the sphere
product. The new symbol \(J\) denotes the joining manifold,
not the proper configuration band of Subsection~\ref{p1b:sec:rtproof42}.
There are actual inclusions
\[
 j_Y[a,w]=[a,w,-w],\qquad
 j_Z[a,w]=[a,-a,w].
 \tag{RT.46.2}\label{p1b:eq:rt-46-2}
\]
Use the double mapping cylinder
\[
 \mathcal P=Y\cup_{j_Y}(J\times[0,1])\cup_{j_Z}Z_s.
 \tag{RT.46.3}\label{p1b:eq:rt-46-3}
\]

\begin{theorem}[The actual configuration pushout]
\label{p1b:rt:theorem46-pushout}
There is a homotopy equivalence
\[
 \boxed{W\simeq\mathcal P,\qquad
 \dim Y=7,\quad\dim Z_s=5,\quad\dim J=4.}
 \tag{RT.46.4}\label{p1b:eq:rt-46-4}
\]
It respects the ordering covers. The canonical map
\(\mathcal P\to W^\circ\) sends the two pieces to their
indicated antipodal configurations and is constant along
each joining cylinder.
\end{theorem}
\begin{proof}
The ordered free space has the exchange-invariant open cover
\[
\begin{split}
 A&=\operatorname{Conf}_2(S^1)\times(S^3)^2,\\
 B&=(S^1)^2\times\operatorname{Conf}_2(S^3).
\end{split}
\tag{RT.46.5}\label{p1b:eq:rt-46-5}
\]
Its intersection has both pairs distinct. Every pair of
distinct product points lies in this cover.

For distinct unit vectors \(x,p\in S^n\), put
\(m=(x+p)/2\), \(d=(x-p)/2\). The deformation
\[
 m_\lambda=(1-\lambda)m,\qquad
 d_\lambda=\sqrt{1-|m_\lambda|^2}\,d/|d|,\qquad
 (x_\lambda,p_\lambda)=(m_\lambda+d_\lambda,m_\lambda-d_\lambda)
 \tag{RT.46.6}\label{p1b:eq:rt-46-6}
\]
retracts onto antipodal pairs. Indeed \(m\perp d\),
the resulting vectors are unit, \(d_\lambda\ne0\), and
exchange negates \(d_\lambda\). Antipodal pairs remain fixed.
Use this for both \(n=1\) and \(n=3\).

First include the span obtained by retracting the circle
configuration in \(A\) and in \(A\cap B\), while leaving
\(B\) unchanged. Each vertex inclusion is an equivariant
homotopy equivalence. Next include the span obtained by
retracting the sphere configuration in its middle vertex
and in \(B\), while leaving its left vertex unchanged.
Again the inclusions form a strict map of spans and are
vertexwise equivalences. The final span, after the free
quotients, is precisely \(Y\leftarrow J\to Z_s\).

An open-cover double mapping cylinder has the homotopy
type of its union, and vertexwise equivalences of spans
preserve that homotopy type. Applying these facts to the
two strict inclusions proves \eqref{p1b:eq:rt-46-4}.
All deformations and inclusions are equivariant, so the
ordering covers are retained. On the two pieces the
canonical map sends
\[
 {}[a,x,p]\mapsto\{(a,x),(-a,p)\},\quad
 {}[z_1,z_2,w]\mapsto\{(z_1,w),(z_2,-w)\}.
\]
Both joining endpoints give \(\{(a,w),(-a,-w)\}\).
Collapsing the joining interval therefore defines the
stated canonical equivalence.
\end{proof}

The joining manifold has a concrete product chart:
\[
 \boxed{J\longrightarrow S^1\times S^3,\qquad
 {}[a,w]\longmapsto(a^2,a^{-1}w).}
 \tag{RT.46.7}\label{p1b:eq:rt-46-7}
\]
Use the fixed complex circle inside the unit quaternions
for \(a^{-1}w\). Both choices of square root of \(a^2\)
give the same class, proving that this is a homeomorphism.
Let \(h_J,u_J,c_J=h_Ju_J\) be its integral generators.

\begin{lemma}[Integral groups and joining maps]
\label{p1b:rt:lemma46-matrices}
The positive-degree nonzero integral groups of the pieces are
\[
\begin{array}{c|ccccc}
 &H^1&H^2&H^3&H^4&H^5\ \text{or}\ H^7\\\hline
 Y&\mathbb Z&0&\mathbb Z&\mathbb Z&H^7=C_2\\
 Z_s&\mathbb Z&C_2&\mathbb Z\oplus C_2&\mathbb Z&H^5=C_2\\
 J&\mathbb Z&0&\mathbb Z&\mathbb Z&0.
\end{array}
\tag{RT.46.8}\label{p1b:eq:rt-46-8}
\]
With compatible orientations, the Mayer--Vietoris joining
matrices on the free summands are
\[
 \boxed{
 H^1:\ (1,-1),\qquad
 H^3:\ (2,-2),\qquad
 H^4:\ (1,-1).}
 \tag{RT.46.9}\label{p1b:eq:rt-46-9}
\]
The torsion summand of \(H^3(Z_s;\mathbb Z)\) maps to zero.
\end{lemma}
\begin{proof}
The circle mapping-torus cochain complex for \(Y\) has
monodromy \(1\), the swap matrix
\(P=\left(\begin{smallmatrix}0&1\\1&0\end{smallmatrix}\right)\),
and \(-1\) in sphere-product degrees \(0,3,6\).
Its sole differential in each such row is \(T-1\).
This gives the \(Y\) groups in \eqref{p1b:eq:rt-46-8}.

For \(Z_s\), regard projection to \([w]\) as a flat
torus bundle over \(\mathbb {RP}^3\).
The cellular torus cochains have zero differential.
The monodromies in fibre degrees \(0,1,2\) are \(1,P,-1\).
In each row the projective-space cochain differentials are
\[
 T-1,\quad 1+T,\quad T-1
 \qquad\text{in base degrees }0,1,2.
 \tag{RT.46.10}\label{p1b:eq:rt-46-10}
\]
These are also a direct cellular calculation: use the
antipodal cellular cover of \(\mathbb {RP}^3\), multiply
by the product CW structure on the torus, and take the
diagonal quotient. There is no fibre cellular differential.
The three row cohomologies are, respectively,
\[
 (\mathbb Z,0,C_2,\mathbb Z),\quad
 (\mathbb Z,0,0,\mathbb Z),\quad
 (0,C_2,0,C_2).
 \tag{RT.46.11}\label{p1b:eq:rt-46-11}
\]
Shift by fibre degree and add them to get the \(Z_s\) groups.
The product chart \eqref{p1b:eq:rt-46-7} gives the \(J\) groups.

Let \(a_0\) be the circle class on the ordered cover of \(J\).
Its product generators pull back to
\(2a_0,u_w,2a_0u_w\).
On the cover of \(Y\), choose free generators whose pullbacks
are \(2a_0,u_x+u_p,a_0(u_x+u_p)\).
The degree-four generator is the primitive mapping-torus
coinvariant, not the product of its degree-one and degree-three
generators. Restriction to the antipodal sphere pairs sends
both \(u_x,u_p\) to \(u_w\).
This gives the coefficients \(1,2,1\) for \(j_Y^*\).

For \(Z_s\), the product-circle function \(z_1z_2\) gives
the primitive degree-one generator; the projective-three-space
top class gives the free degree-three generator.
The primitive degree-four transfer has cover
\((a_1+a_2)u_w\); the cellular coinvariant calculation in
\eqref{p1b:eq:rt-46-10} fixes its primitiveness.
On \(j_Z\), \(a_1=a_2=a_0\), so these restrictions have
the same coefficients \(1,2,1\).
Integral torsion maps to zero in the free group \(H^3(J)\).
Taking the difference of the two restrictions proves
\eqref{p1b:eq:rt-46-9}.
\end{proof}

\begin{corollary}[The torsion four-class is a joining class]
\label{p1b:rt:corollary46-tau}
Let \(\delta_{\rm MV}:H^3(J;\mathbb Z)\to
H^4(W;\mathbb Z)\) be the connecting map for this pushout.
Then
\[
 \boxed{\ker\delta_{\rm MV}=2\mathbb Z\{u_J\},\qquad
        \delta_{\rm MV}(u_J)=\tau^2.}
 \tag{RT.46.12}\label{p1b:eq:rt-46-12}
\]
The same calculation gives
\[
 H^3(W;\mathbb Z)=\mathbb Z\oplus C_2,\quad
 H^4(W;\mathbb Z)=\mathbb Z\oplus C_2,\quad
 H^5(W;\mathbb Z)\xrightarrow{\cong}H^5(Z_s;\mathbb Z)=C_2.
 \tag{RT.46.13}\label{p1b:eq:rt-46-13}
\]
\end{corollary}
\begin{proof}
Since \(H^2(J)=0\), the degree-three group is the kernel
of the \((2,-2)\) matrix together with the torsion summand.
Its image in \(H^3(J)\) is exactly \(2\mathbb Z\).
The resulting cokernel injects into \(H^4(W)\) with order two.
The free part of \(H^4(W)\) is the kernel of \((1,-1)\),
so the integral group is \(\mathbb Z\oplus C_2\).
The retained current owner identifies its unique nonzero
torsion class as \(\tau^2\); hence the nonzero connecting
value is that class. Finally the degree-four joining matrix
is surjective and \(H^5(Y)=H^5(J)=0\), which gives the
last isomorphism. This agrees with, and does not replace,
the retained integral fifth-class calculation.
\end{proof}

The model \eqref{p1b:eq:rt-46-3} is now used with its actual boundary map
and its actual bundle gluing, rather than only its cohomology.
These are specified in the following MAIN supplement.



\subsection{The boundary map and the actual gauge gluing source}
\label{p1b:sec:rtproof47}

Use \(Y,Z_s,J,\mathcal P\) from
Theorem~\ref{p1b:rt:theorem46-pushout}.
Keep the prescribed boundary product frame and the actual
bundle \(E_\partial=\mathbf1\oplus\lambda_L\).
Write a lens direction using its positive-circle-axis hemisphere:
\[
 n=(\sqrt{1-\rho^2},\rho u),\qquad
 0\le\rho\le1,\quad u\in S^2\subset\operatorname{Im}\mathbb H.
 \tag{RT.47.1}\label{p1b:eq:rt-47-1}
\]
The centre collapses \(u\); on the boundary \(\rho=1\),
\(u\sim-u\). This is the closed three-cell model of \(L\).
Put
\[
 a=i e^{is},\qquad
 \alpha(\rho)=3\pi\rho/2,\qquad
 \beta(\rho)=3\pi(1-\rho)/2.
\]

\begin{theorem}[An explicit three-stage boundary map]
\label{p1b:rt:theorem47-boundary}
Under \(W\simeq\mathcal P\), the actual boundary inclusion
is represented by the following map \(b_R:R\to\mathcal P\):
\[
 b_R(s,y,[n])=
 \begin{cases}
 {}[a,y\exp(\alpha(\rho)u),y\exp(-\alpha(\rho)u)]\in Y,
                  &0\le\rho\le1/3,\\
 ([a,yu],\,3\rho-1)\in J\times[0,1],
                  &1/3\le\rho\le2/3,\\
 {}[e^{i(s+\beta(\rho))},e^{i(s-\beta(\rho))},yu]\in Z_s,
                  &2/3\le\rho\le1.
 \end{cases}
 \tag{RT.47.2}\label{p1b:eq:rt-47-2}
\]
It is continuous at both joining endpoints, the lens centre,
and the identified outer sphere. In particular its restriction
to \(F\times Q\), \(Q=\mathbb {RP}^2\), is exactly
\[
 (s,y,[u])\longmapsto[e^{is},e^{is},yu]\in Z_s.
 \tag{RT.47.3}\label{p1b:eq:rt-47-3}
\]
This is the retained coincident-circle boundary map, now
embedded in the global pushout.
\end{theorem}
\begin{proof}
At \(\rho=1/3\) the sphere points are \((yu,-yu)\),
so the first expression is \(j_Y[a,yu]\).
At \(\rho=2/3\) the circle points are \((a,-a)\),
so the last expression is \(j_Z[a,yu]\).
At \(\rho=0\) the sphere points are both \(y\);
the apparent \(u\) dependence disappears continuously.
At \(\rho=1\) the circle points coincide. Changing \(u\)
to \(-u\) then applies the defining exchange of \(Z_s\),
so the outer projective identification is respected.

Here is an explicit check against the actual diagonal link.
Compose \eqref{p1b:eq:rt-47-2} with the canonical map to \(W^\circ\).
Its two half-separation angles are
\[
 d_c(\rho)=
 \begin{cases}\pi/2,&\rho\le2/3,\\
 \beta(\rho),&\rho\ge2/3,\end{cases}
 \qquad
 d_s(\rho)=
 \begin{cases}\alpha(\rho),&\rho\le1/3,\\
 \pi/2,&\rho\ge1/3.\end{cases}
 \tag{RT.47.4}\label{p1b:eq:rt-47-4}
\]
The resulting unordered points are
\((e^{i(s+d_c)},y\exp(d_su))\) and
\((e^{i(s-d_c)},y\exp(-d_su))\).
For a small \(\varepsilon>0\), interpolate the angles linearly
from \((\varepsilon\sqrt{1-\rho^2},\varepsilon\rho)\)
to \((d_c,d_s)\). They never both vanish.
They remain in \([0,\pi/2]\), so the product points
never coincide. The centre and outer projective identification
are preserved throughout. The initial map is the prescribed
small normal link in the product tangent frame.
This gives a homotopy in the actual free configuration space.
Radial excision transports it to the retained \(W\).
It proves the asserted boundary map and preserves the
ordering cover, including the line \(\lambda_L\).
\end{proof}

\begin{theorem}[The actual Chern clutching parity]
\label{p1b:rt:theorem47-clutch}
For every retained compatible actual \(V\), its restriction
to \(J\) is an ordinary trivial rank-two complex bundle.
Fix its actual pieces on \(Y,Z_s\) and a reference gluing.
Changing the gluing by a determinant-one map \(f:J\to SU(2)\)
with primary \(d u_J\) changes its second Chern class by
\[
 \boxed{\Delta c_2=\delta_{\rm MV}(d u_J)
                        =(d\bmod2)\tau^2.}
 \tag{RT.47.5}\label{p1b:eq:rt-47-5}
\]
After choosing a reference gluing with \(c_2=0\), the actual
input \(c_2(V)=(A\bmod2)\tau^2\) therefore imposes
\[
 \boxed{d\equiv A\pmod2.}
 \tag{RT.47.6}\label{p1b:eq:rt-47-6}
\]
An explicit primitive joining map is
\[
 f_1[a,w]=a^{-1}w\in SU(2),\qquad f_d=f_1^d.
 \tag{RT.47.6a}\label{p1b:eq:rt-47-6a}
\]
It has primary \(u_J\); its powers have primary \(d u_J\).
This is a primary clutching constraint, not a classification
of the higher bundle or gauge data. The zero-Chern reference
is an ordinary gluing reference; its boundary compatibility
is a separate condition.
\end{theorem}
\begin{proof}
The product \(J\) has \(H^2(J)=0\), and restriction of the
torsion \(c_2(V)\) to its free \(H^4(J)\) is zero.
Rank-two complex bundles on a four-dimensional CW complex
are classified by \(c_1,c_2\), so \(V|_J\) is trivial.

In a fixed trivialization on \(J\), a determinant-one
clutching change transgresses the sphere primary to \(c_2\).
Indeed the universal clutching of a generator
\(S^3\to SU(2)\) gives the unit second Chern class on
\(S^4\). Naturality with the Mayer--Vietoris connecting
map gives \(\Delta c_2=\delta_{\rm MV}(d u_J)\).
Any fourth homotopy coordinate of \(f\) does not change
this characteristic class. Apply \eqref{p1b:eq:rt-46-12}.
A degree-one sphere map is supplied by the second coordinate
of \eqref{p1b:eq:rt-46-7}, giving \eqref{p1b:eq:rt-47-6a}. Multiplication in \(SU(2)\)
pulls the integral three-class back to the sum of the two
factor classes, so its pointwise powers have primary \(d u_J\).
Thus either Chern parity can be changed to zero to choose
the stated ordinary reference. Relative to it, the retained input gives
\eqref{p1b:eq:rt-47-6}. No identification of all even-degree gluing
changes as ordinary bundle isomorphisms is inferred.
\end{proof}

\begin{theorem}[Exact source for actual global gauges]
\label{p1b:rt:theorem47-gauges}
Pull the actual bundle \(V\) to \(\mathcal P\) by the
equivalence above, retaining its two pieces and their
joining identification. Up to the corresponding homotopy
equivalence of section groups, an actual global gauge is
exactly a triple
\[
 \boxed{(g_Y,g_Z,H),\qquad
 H:J\times[0,1]\to\operatorname{Aut}(V|_J),}
 \tag{RT.47.7}\label{p1b:eq:rt-47-7}
\]
where \(g_Y,g_Z\) are actual gauges on the two pieces and
\[
 H(-,0)=j_Y^*g_Y,\qquad
 H(-,1)=j_Z^*g_Z
 \tag{RT.47.8}\label{p1b:eq:rt-47-8}
\]
after the fixed joining bundle identification.
In separate frames this last endpoint is conjugated by
the actual clutching matrix; that matrix must be retained.
Restriction to \(R\) is obtained by applying \(g_Y,H,g_Z\)
to the three explicit stages of \eqref{p1b:eq:rt-47-2}.
\end{theorem}
\begin{proof}
Trivialize only the interval direction of the actual
bundle on the joining cylinder. Its endpoint bundle
identifications are fixed by the pulled-back \(V\).
A section of its automorphism bundle is then a path of
sections on \(J\). A section on all of \(\mathcal P\)
is precisely two endpoint sections and such a path with
the equalities \eqref{p1b:eq:rt-47-8}. Conversely those equalities
paste the three actual sections continuously.
Products are pointwise, giving the section-group version
of the homotopy pullback, not merely a matching of
cohomology classes.

Theorem~\ref{p1b:rt:theorem47-boundary} supplies the restriction
map. To compare it with the prescribed Born boundary frame,
retain the explicit link homotopy in that theorem.
Parallel transport for a Hermitian connection on the
actual bundle identifies its endpoint pullbacks.
The space of such connections is affine, so different
connections give homotopic identifications and the same
boundary gauge component. This retains the actual
twisted bundle and does not choose a new product frame
for its ordering line.
\end{proof}

\begin{corollary}[A specified seven-sphere source for the remaining bits]
\label{p1b:rt:corollary47-top}
Let \(F_\vee=(S^1\times\{y_0\})\cup(\{s_0\}\times S^3)\), and
\[
 K_R=(F\times Q)\cup(F_\vee\times L).
 \tag{RT.47.9}\label{p1b:eq:rt-47-9}
\]
Then
\[
 \boxed{R/K_R=(F/F_\vee)\wedge(L/Q)\cong S^4\wedge S^3=S^7.}
 \tag{RT.47.10}\label{p1b:eq:rt-47-10}
\]
After the retained low, six-stage, and \(e\)-comparisons
have been made, the difference between a constructed
normalized boundary gauge of low data \((h,0)\) or
\((0,2q)\) and its natural candidate is respectively
\[
 h_t\,t\quad\text{or}\quad w_t\,t.
 \tag{RT.47.11}\label{p1b:eq:rt-47-11}
\]
It can be represented as the natural supported map from
\eqref{p1b:eq:rt-47-10} with coefficient in
\(\pi_7SU(2)=C_2\{\mu\}\).
Thus \eqref{p1b:eq:rt-47-2}, the actual triple \eqref{p1b:eq:rt-47-7}, and the
retained comparisons specify the source on which to
evaluate the two missing coefficients.
\end{corollary}
\begin{proof}
Both quotient factors in \eqref{p1b:eq:rt-47-10} are the stated
spheres in their product and projective cell structures.
The quotient of their product by the union in \eqref{p1b:eq:rt-47-9}
is their smash. The natural \(t=j_7\mu\) is supported
on the resulting circle--sphere--lens top cell.
The retained evaluated values \(h_e=w_e=0\), the known
six-stage columns, and the exact low data leave only
the displayed \(t\) differences. Since the retained
top insertion \(j_7:C_2\to\pi_0\mathcal G(E_\partial)\)
is nonzero, its coefficient is unique.
The following global-triple construction and complete comparison maps
evaluate this coefficient with the quotient source and its frame retained.
\end{proof}

This construction fixes the global geometric source and its boundary map.
The following whole-map evaluations give \(h_t=w_t=1\) and
\(\zeta_{\rm nat}=0\). The retained relative top class is
\(\Theta\), with final obstructions \(h_t\Theta,w_t\Theta\).
Supplementary checks and the limits of the Chern-parity
test are recorded in Section~\ref{p1b:app:rt47-pushout}.



\subsection{The exact component calculation on the joining manifold}
\label{p1b:sec:rtproof48}

Keep the actual bundle pieces and joining identification of
Theorem~\ref{p1b:rt:theorem47-gauges}. Only the joining bundle is
trivialized. In the product chart of \eqref{p1b:eq:rt-46-7}, write
\[
 J=S^1_\xi\times S^3_v,\qquad (\xi,v)=(a^2,a^{-1}w),\qquad
 \Gamma_J=\operatorname{Map}(J,U(2)).
 \tag{RT.48.1}\label{p1b:eq:rt-48-1}
\]
The product chart uses the fixed complex circle in the unit
quaternions. Set \(q(\xi)=\operatorname{diag}(\xi,1)\) and
\(p_n(v)=v^n\). Let \(T:J\to SU(2)\) be the map supported on
\(J/(S^1\vee S^3)=S^4\) with coefficient
\(\eta_3\in\pi_4SU(2)=C_2\).
The secondary coordinate introduced here is a joining coordinate;
it is not either of the natural boundary coefficients \(h_t,w_t\).

\begin{theorem}[Joining components with their multiplication]
\label{p1b:rt:theorem48-components}
Every component of \(\Gamma_J\) has a unique expression
\[
 (k,n,\epsilon)=[q(\xi)^k p_n(v)T(\xi,v)^\epsilon],\qquad
 k,n\in\mathbb Z,\quad\epsilon\in\mathbb F_2.
 \tag{RT.48.2}\label{p1b:eq:rt-48-2}
\]
Here \(k\) is determinant winding and \(n\) is the integral
sphere primary. In these coordinates,
\[
 \boxed{(k,n,\epsilon)(l,m,\delta)
  =(k+l,n+m,\epsilon+\delta+ln),}
 \tag{RT.48.3}\label{p1b:eq:rt-48-3}
\]
where the third coordinate is reduced modulo two. Consequently
\[
 (k,n,\epsilon)^{-1}=(-k,-n,\epsilon+kn),\qquad
 {}[(k,n,\epsilon),(l,m,\delta)]=(0,0,km+ln).
 \tag{RT.48.4}\label{p1b:eq:rt-48-4}
\]
\end{theorem}
\begin{proof}
The map \(U(2)\to S^1\times SU(2)\) given by
\(g\mapsto(\det g,q(\det g)^{-1}g)\) is a homeomorphism
of spaces. Maps to its first factor have the unique integer
\(k\). For the second factor, put
\(M=\operatorname{Map}(S^3,SU(2))\).
Its components are the degree integers. Evaluation at the
sphere basepoint splits as a product of spaces, giving
\(\pi_1M=\pi_1SU(2)\oplus\pi_4SU(2)=C_2\).
Translation by \(p_n\) identifies the components of \(M\).
Free homotopy classes of loops in a component of \(M\) are
therefore the two elements of this abelian fundamental group.
Adjunction identifies the nonzero one with \(T\).
This proves existence and uniqueness in \eqref{p1b:eq:rt-48-2}.

In determinant-one maps on \(J\), the supported \(T\) commutes
up to homotopy with every sphere map: their commutator factors
through \(S^4\wedge S^3=S^7\), above the dimension of \(J\).
The chosen powers satisfy \(p_np_m=p_{n+m}\) pointwise.
Conjugating \(p_n\) by \(q^l\) adds \(lnT\), by the retained
Samelson value \(\langle q,\iota_3\rangle=\eta_3\)
and its bilinearity. Conjugation of \(T\) by \(q^l\)
adds nothing: its commutator factors through \(S^1\wedge S^4\),
again above dimension four. Moving \(q^l\) to the left in
\(q^kp_nT^\epsilon q^lp_mT^\delta\) gives \eqref{p1b:eq:rt-48-3}.
The sign from inverse conjugation has no effect in \(C_2\).
The inverse and commutator formulas follow by multiplication.
The Samelson input is the same actual operation as
Lemma~\ref{p1b:rt:lemma104}, not a stable replacement of it.
\end{proof}

\begin{theorem}[The complete joining test for actual piece gauges]
\label{p1b:rt:theorem48-matching}
Choose separate joining frames on the two actual pieces, with
convention that the \(Z_s\) endpoint is \(c g_Zc^{-1}\) in the
\(Y\) joining frame. Write
\[
 {}[c]=(\ell,d,\epsilon_c),\quad
 {}[j_Y^*g_Y]=(k_Y,n_Y,\epsilon_Y),\quad
 {}[j_Z^*g_Z]=(k_Z,n_Z,\epsilon_Z).
 \tag{RT.48.5}\label{p1b:eq:rt-48-5}
\]
There is a matching path, and hence the actual global triple
of \eqref{p1b:eq:rt-47-7}, if and only if
\[
 \boxed{k_Y=k_Z,\qquad n_Y=n_Z,\qquad
  \epsilon_Y=\epsilon_Z+\ell n_Z+k_Zd.}
 \tag{RT.48.6}\label{p1b:eq:rt-48-6}
\]
The secondary clutching coordinate \(\epsilon_c\) cancels
from this component test. It will still act on matching paths.
\end{theorem}
\begin{proof}
Apply \eqref{p1b:eq:rt-48-3}--\eqref{p1b:eq:rt-48-4} to conjugation:
\[
 \boxed{\operatorname{Ad}_{(\ell,d,\epsilon_c)}
 (k,n,\epsilon)=(k,n,\epsilon+\ell n+kd).}
 \tag{RT.48.7}\label{p1b:eq:rt-48-7}
\]
Two actual maps on \(J\) are joined by a path precisely when
they belong to the same component. The component classification
in Theorem~\ref{p1b:rt:theorem48-components} therefore proves
both necessity and sufficiency. Pasting the two gauges with
that path is the actual section construction of
Theorem~\ref{p1b:rt:theorem47-gauges}.
This assertion presupposes the two actual piece gauges;
it does not substitute three coordinates for their construction.
\end{proof}

Choose the global free generator \(q\) in the orientation of
\eqref{p1b:eq:rt-46-9}. Its joining restriction is \(2u_J\).
The determinant generator \(h\) restricts to \(h_J\), and
the torsion primary \(\eta_W\) restricts to zero in
\(H^3(J;\mathbb Z)\). Thus actual gauges with the retained
low data have the following joining test:
\[
\begin{array}{c|cc|c}
 \text{global low data}&k&n&\epsilon_Y-\epsilon_Z\\\hline
 (h,0)&1&0&d\bmod2\\
 (0,\eta_W)&0&0&0\\
 (0,2q)&0&4&0.
\end{array}
 \tag{RT.48.8}\label{p1b:eq:rt-48-8}
\]
The fourth coefficient in the last row is important: the
joining map already doubles the global free three-primary.
For determinant-one clutching relative to the ordinary
zero-Chern reference of Theorem~\ref{p1b:rt:theorem47-clutch},
\(\ell=0\) and \(d\equiv A\pmod2\). In that reference,
the determinant row requires opposite joining secondary
coordinates when \(A\) is odd. The two determinant-one
rows require equal joining secondary coordinates.
These statements retain the actual piece bundles, their
torsion data, and their frames. They do not evaluate
\(\epsilon_Y,\epsilon_Z\) for the required gauges.



\subsection{Matching paths and actual gauges supported on the joining cylinder}
\label{p1b:sec:rtproof49}

Use the coordinate convention of Theorem~\ref{p1b:rt:theorem48-components}.
Let \(L_0\) be the loop of constant maps
\(q(e^{2\pi it})\) on \(J\). Let \(\alpha\) be the loop
supported on \(S^1_t\wedge S^3_v=S^4\) with coefficient
\(\eta_3\), independent of \(\xi\). Let \(\beta\) be the
loop supported on
\(S^1_t\wedge S^1_\xi\wedge S^3_v=S^5\) with coefficient
\(\eta_3\eta_4\). All three loops have endpoints the identity.

\begin{theorem}[The matching-path group]
\label{p1b:rt:theorem49-loops}
At the identity, and by translation at every component,
\[
 \boxed{\pi_1\Gamma_J=\mathbb Z\{L_0\}
          \oplus C_2\{\alpha\}\oplus C_2\{\beta\}.}
 \tag{RT.49.1}\label{p1b:eq:rt-49-1}
\]
In coordinates \((r,x,y)\), the conjugation action of the
actual clutching component \((\ell,d,\epsilon_c)\) is
\[
 \boxed{\operatorname{Ad}_c(r,x,y)=
 (r,x+dr,y+\ell x+(\epsilon_c+\ell d)r).}
 \tag{RT.49.2}\label{p1b:eq:rt-49-2}
\]
Only the last two coordinates are reduced modulo two.
In particular the secondary clutching bit does act on the
constant determinant loop, even though it cancels from \eqref{p1b:eq:rt-48-6}.
\end{theorem}
\begin{proof}
The space splitting \(U(2)\cong S^1\times SU(2)\) gives
the determinant \(\mathbb Z\) summand on fundamental groups.
For the second factor set \(M=\operatorname{Map}(S^3,SU(2))\).
Then
\[
 \pi_1M=C_2\{\eta_3\},\qquad
 \pi_2M=C_2\{\eta_3\eta_4\}.
\]
The free loop space \(\operatorname{Map}(S^1,M)\) splits
as a space into \(M\) and its based loop space, by dividing
a loop by its value at the circle basepoint. The constant-loop
section splits its evaluation sequence. On fundamental groups
this gives \(\pi_1M\oplus\pi_2M\), with both summands of
order two. This rules out a hidden \(C_4\) extension.
It gives precisely the representatives \(\alpha,\beta\)
above. Equivalently their supported sources appear in the
product suspension splitting
\(\Sigma J\simeq S^2\vee S^4\vee S^5\);
the \(S^2\) summand maps trivially to \(SU(2)\).

The retained operation
\(A_*x=\eta_3Ex\) gives
\(A_*\iota_3=\eta_3\) and
\(A_*\eta_3=\eta_3\eta_4\).
Apply it to the three chosen clutching factors. Their actions are
\[
\begin{array}{c|ccc}
 &L_0&\alpha&\beta\\\hline
 q(\xi)^\ell&L_0&\alpha+\ell\beta&\beta\\
 p_d(v)&L_0+d\alpha&\alpha&\beta\\
 T^{\epsilon_c}&L_0+\epsilon_c\beta&\alpha&\beta.
\end{array}
 \tag{RT.49.3}\label{p1b:eq:rt-49-3}
\]
For the second row, the determinant-loop commutator is
\(d\eta_3\) on \(S^1_t\wedge S^3_v\). For the third it is
\(\epsilon_c\eta_3E\eta_3\) on the five-sphere.
For the first row acting on \(\alpha\), the circle \(\xi\)
and the loop parameter \(t\) are independent, giving the
same five-sphere composition. All the other commutators with
\(\alpha\) or \(\beta\) either have smash dimension above
five or repeat the same sphere or circle coordinate, whose
reduced diagonal is null. The two determinant loops in the
first entry commute pointwise. There is consequently no
further term in this dimension.

The clutching convention is \(c=q^\ell p_dT^{\epsilon_c}\).
Compose the three rows of \eqref{p1b:eq:rt-49-3} in that order. The result
on \(L_0\) is \(L_0+d\alpha+(\epsilon_c+\ell d)\beta\),
and the result on \(\alpha\) is \(\alpha+\ell\beta\).
This proves \eqref{p1b:eq:rt-49-2}. The extra \(\ell d\) term is required
by the noncommutative component law \eqref{p1b:eq:rt-48-3}.
\end{proof}

\begin{theorem}[A finite bound for paths after endpoint homotopies]
\label{p1b:rt:theorem49-paths}
Fix a pair of actual piece gauge components that satisfies
\eqref{p1b:eq:rt-48-6}, and one matching path. Let \(H_{\rm end}\) be
the subgroup of \(\pi_1\Gamma_J\) generated by the actual
endpoint loop images, with the \(Z_s\) image transported by
\eqref{p1b:eq:rt-49-2}. The set of global gauge components over this
fixed pair is a torsor for
\[
 \boxed{\pi_1\Gamma_J/H_{\rm end}.}
 \tag{RT.49.4}\label{p1b:eq:rt-49-4}
\]
Moreover \(2L_0\in H_{\rm end}\). Thus this torsor is
finite, and its cardinality belongs to
\[
 \boxed{\{1,2,4,8\}.}
 \tag{RT.49.5}\label{p1b:eq:rt-49-5}
\]
The actual subgroup \(H_{\rm end}\), and hence the selected
cardinality, requires the endpoint loop images on the actual
piece bundles.
\end{theorem}
\begin{proof}
For fixed endpoints, homotopy classes of paths between them
form a torsor for the fundamental group of their component
of \(\Gamma_J\). Homotopies of either piece gauge act through
their endpoint loops. This is the path description of the
homotopy pullback in Theorem~\ref{p1b:rt:theorem47-gauges}.
Its target fundamental group is abelian, so these actions
give the quotient by their sum, as in \eqref{p1b:eq:rt-49-4}.

On every actual rank-two bundle, the loop of central gauges
\(e^{2\pi it}I\) exists. Its restriction to \(J\) has
determinant winding two. Dividing by \(q(e^{4\pi it})\)
leaves a constant-map loop in \(SU(2)\), which is null since
\(\pi_1SU(2)=0\). Its joining class is exactly \(2L_0\).
Conjugation by the clutching matrix fixes this central loop.
It is in both endpoint images. The quotient in \eqref{p1b:eq:rt-49-4}
therefore factors through
\(\pi_1\Gamma_J/\langle2L_0\rangle=C_2^3\).
Every quotient of this group has one of the orders \eqref{p1b:eq:rt-49-5}.
\end{proof}

\begin{theorem}[Actual joining gauges and odd-parity path independence]
\label{p1b:rt:theorem49-boundary}
For each of \(L_0,\alpha,\beta\), take identity gauges
on \(Y,Z_s\) and the indicated loop on the joining cylinder.
This constructs actual global gauges \(K_{L_0},K_\alpha,K_\beta\)
on every retained \(V\). They have zero determinant winding
and zero integral sphere primary globally. The determinant-one
gauge \(K_\alpha\) has fourth-stage class
\[
 \boxed{\delta_{\rm MV}(\bar u_J)=\overline{\tau^2}=\chi^4.}
 \tag{RT.49.6}\label{p1b:eq:rt-49-6}
\]
Its normalized boundary restriction is zero by the retained
\(u_{\chi,V}=0\) theorem. At odd parity the normalized
boundary restrictions of all three joining gauges are zero.
More generally, the boundary component of a global gauge
at odd parity is independent of the choice of matching path
for a fixed pair of actual piece gauges.
\end{theorem}
\begin{proof}
All three loops equal the identity at both interval ends,
so they paste with the identity piece gauges under the actual
bundle identification. This is an explicit instance of the
global triple \eqref{p1b:eq:rt-47-7}, not only a cohomology realization.

The determinant and integral three-primary restrict to zero
on both pieces. The Mayer--Vietoris restriction maps in
degrees one and three are injective: the former uses
surjectivity in degree zero and the latter uses \(H^2(J)=0\).
Both global low data consequently vanish. For \(K_\alpha\),
the first supported section coordinate is \(\eta_3\) on
the suspended sphere; its fourth-stage joining cochain is
\(\delta_{\rm MV}(\bar u_J)\). Reduction of \eqref{p1b:eq:rt-46-12}
gives \eqref{p1b:eq:rt-49-6}. The retained
Theorem~\ref{p1b:rt:theorem129} evaluates its boundary value
as \(u_\chi=0\).

At odd parity the retained pure finite kernel for zero low
data is zero, because \(u_\chi=u_c=0\). Hence every actual
joining loop gauge has zero normalized boundary value,
including \(K_{L_0}\) and \(K_\beta\).
For any two paths with the same actual endpoints, the pointwise
quotient of their two global triples is identity on the pieces
and a loop on the joining cylinder. It is therefore a zero-low
joining gauge of this type and has zero odd boundary value.
This proves path independence. At even parity the retained
kernel only bounds the possible difference by \(\langle t\rangle\).
\end{proof}

The matching component test, path group and clutching action are proved here. The following included piece-gauge constructions and complete seven-sphere comparisons retain the prescribed boundary transport; their full natural-map evaluation assigns \(h_t=w_t=1\). Supplementary finite endpoint checks are recorded in Section~\ref{p1b:app:rt49-joining}.



\subsection{Gauge representatives on the sphere-antipodal piece}
\label{p1b:sec:rtproof50}

Retain the actual piece
\[
 Z_s=(T^2\times S^3)/((z_1,z_2,w)\sim(z_2,z_1,-w)),
 \qquad \pi_Z:Z_s\longrightarrow\mathbb {RP}^3,
 \tag{RT.50.1}\label{p1b:eq:rt-50-1}
\]
and its actual restricted rank-two bundle. Write \(\gamma\)
for the real ordering line, \(\chi=w_1(\gamma)\), and
\(\lambda=\gamma\otimes\mathbb C\), with \(c_1(\lambda)=\tau_Z\).
The retained integral groups imply
\(c_1(V|_{Z_s})=\tau_Z\) and \(c_2(V|_{Z_s})=0\):
the global second Chern class is torsion, while \(H^4(Z_s;\mathbb Z)\)
is free. In particular \(\bar\tau_Z=\chi^2\).

\begin{lemma}[The evaluated secondary bundle operation]
\label{p1b:rt:lemma50-operation}
The tangent bundle and the secondary operation satisfy
\[
 TZ_s\cong\mathbb R^4\oplus\gamma,\qquad
 \boxed{(Sq^2+\chi^2):H^3(Z_s;\mathbb F_2)
                  \longrightarrow H^5(Z_s;\mathbb F_2)\text{ is zero}.}
 \tag{RT.50.2}\label{p1b:eq:rt-50-2}
\]
More explicitly, there is a class \(r_Z\in H^2(Z_s;\mathbb F_2)\)
such that
\[
 H^3(Z_s;\mathbb F_2)=\langle\chi^3,\chi r_Z\rangle,
 \quad H^5(Z_s;\mathbb F_2)=\langle\chi^3r_Z\rangle,
 \quad Sq^1r_Z=\chi r_Z,\quad Sq^2r_Z=r_Z^2=0.
 \tag{RT.50.3}\label{p1b:eq:rt-50-3}
\]
In these bases both \(Sq^2\) and multiplication by \(\chi^2\)
have the matrix \((0,1)\); their sum has matrix \((0,0)\).
\end{lemma}
\begin{proof}
The projection in \eqref{p1b:eq:rt-50-1} is a flat torus bundle. Its vertical
translation frame transforms by the two-coordinate permutation,
so the vertical tangent bundle is \(\mathbb R\oplus\gamma\).
The base \(\mathbb {RP}^3\) is parallelizable. A bundle connection
splits the tangent sequence and gives the first assertion.

There is a section \(i:\mathbb {RP}^3\to Z_s\) at
\((z_1,z_2)=(1,1)\). Its normal bundle is
\(\mathbb R\oplus\gamma\). Define \(r_Z=i_!(1)\) with
mod-two coefficients. The Thom formula gives
\(Sq^1r_Z=i_!w_1(\nu_i)=\chi r_Z\) and
\(Sq^2r_Z=i_!w_2(\nu_i)=0\). The self-intersection formula
also gives \(r_Z^2=i_!w_2(\nu_i)=0\).
The projection formula gives
\(\langle\chi^3r_Z,[Z_s]_2\rangle
=\langle\chi^3,[\mathbb {RP}^3]_2\rangle=1\).
Also \(i^*r_Z=w_2(\nu_i)=0\), whereas \(i^*\chi^3\ne0\).
Thus \(\chi^3,\chi r_Z\) are independent. The integral groups
of Lemma~\ref{p1b:rt:lemma46-matrices} and the universal
coefficient theorem show that the dimensions are exactly two
and one, proving the basis statements.

The class \(\chi\) is pulled back from the three-dimensional base,
so \(\chi^4=0\). Cartan's formula now gives
\[
 Sq^2\chi^3=\chi^5=0,\qquad
 Sq^2(\chi r_Z)=\chi^2Sq^1r_Z+\chi Sq^2r_Z=\chi^3r_Z.
\]
Multiplication by \(\chi^2\) gives the same two values.
This proves the zero operation. As a check, the second Wu
class from the tangent splitting is \(w_1^2+w_2=\chi^2\),
which gives the same equality in top degree.
\end{proof}

\begin{theorem}[The two actual zero-second-Chern bundles on \(Z_s\)]
\label{p1b:rt:theorem50-bundles}
There are exactly two ordinary rank-two complex bundle classes on
\(Z_s\) with \(c_1=\tau_Z,c_2=0\). One is
\(E_Z^0=\mathbf1\oplus\lambda\). The other can be constructed
by replacing its gluing on a small five-ball, disjoint from
\(j_Z(J)\), by
\(\eta_3:S^4\to SU(2)\). Denote it by \(E_Z^1\).
This is a local description of the two piece classes; it does
not identify their labels with a chosen global fifth-bit origin.
\end{theorem}
\begin{proof}
The first relevant Postnikov invariant of \(BU(2)\) is
\[
 k_6=Sq^2\bar c_2+\bar c_1\bar c_2.
 \tag{RT.50.4}\label{p1b:eq:rt-50-4}
\]
Use the invariant and loop slant proved in
Section~\ref{p1b:app:g20-determinant-forgetting}.
On a five-dimensional CW complex, once \(c_1,c_2\) are fixed,
the lifts form an \(H^5(-;\mathbb F_2)\) torsor before base
homotopies are identified. A loop of the four-stage with data
\((h,v)\in H^1(-;\mathbb Z)\oplus H^3(-;\mathbb Z)\)
acts by
\[
 (Sq^2+\bar c_1)\bar v+\bar h\bar c_2.
 \tag{RT.50.5}\label{p1b:eq:rt-50-5}
\]
This includes ordinary determinant homotopies, not only those
fixing a determinant trivialization. With the stated Chern data,
both terms vanish by Lemma~\ref{p1b:rt:lemma50-operation}.
The split bundle supplies one lift, and
\(H^5(Z_s;\mathbb F_2)=C_2\) supplies exactly two distinct
classes. No higher stage is visible on a five-dimensional source.

The pinch of a five-ball to \(S^5\) detects the nonzero
mod-two top class. Inserting the generator
\(\pi_5BU(2)=\pi_4U(2)=C_2\{\eta_3\}\) on that ball
therefore toggles the two lifts. A ball off the closed
four-dimensional joining hypersurface can be used, so its
modification does not change the bundle on \(J\).
\end{proof}

\begin{theorem}[Actual piece gauges and their joining components]
\label{p1b:rt:theorem50-gauges}
On each actual \(V|_{Z_s}\), there are gauge representatives
\(D_Z,S_Z\) with low data respectively
\((h_Z,0)\) and \((0,2v_Z)\), where \(h_Z\) is the primitive
product-circle class and \(v_Z=\pi_Z^*[\mathbb {RP}^3]\).
In a joining frame induced from their split-or-patched model,
\[
 \boxed{[j_Z^*D_Z]=(1,0,0),\qquad
        {}[j_Z^*S_Z]=(0,4,0).}
 \tag{RT.50.6}\label{p1b:eq:rt-50-6}
\]
The sphere component \((0,4,0)\) is unchanged by every
ordinary joining frame change.
\end{theorem}
\begin{proof}
First work on \(E_Z^0\). Put
\[
 D_Z[z_1,z_2,w]=\operatorname{diag}(z_1z_2,1).
 \tag{RT.50.7}\label{p1b:eq:rt-50-7}
\]
It is exchange-invariant and has determinant \(h_Z\).
Its three-primary is \(h_Z\tau_Z\), as follows by slanting
the second Chern class of the two line summands of its gauge
mapping torus. This product is zero. Indeed
\(h_Z=\operatorname{tr}(a_1)\) for the actual ordered cover;
its cover is \(a_1+a_2\). The projection formula gives
\(\tau_Zh_Z=\operatorname{tr}(p^*\tau_Z\,a_1)=0\), since
the pulled-back ordering line is trivial.

On \(\mathbf1\oplus\lambda\) over \(\mathbb {RP}^3\), choose
a determinant-one gauge \(K\) supported in its open three-cell
with primary the positively oriented top integral class.
This is an actual local gauge: trivialize over a smaller ball,
insert the degree-one \(S^3\to SU(2)\) map and take identity
off the ball. Define
\[
 S_Z=(\pi_Z^*K)^2.
 \tag{RT.50.8}\label{p1b:eq:rt-50-8}
\]
It has primary \(2v_Z\) and no torsion three-primary.

On the joining manifold the sign line is ordinarily trivial;
a trivializing section on its cover is supplied by the circle
coordinate \(a\). In this joining frame \eqref{p1b:eq:rt-50-7} becomes
\(\operatorname{diag}(-\xi,1)\), of component \((1,0,0)\).
The joining restriction of \(\pi_Z^*K\) has component
\((0,2,\epsilon_K)\) for some \(\epsilon_K\in\mathbb F_2\),
because the projective-space class pulls back to \(2u_J\).
The exact multiplication \eqref{p1b:eq:rt-48-3} evaluates its square as
\[
 (0,2,\epsilon_K)^2=(0,4,0),
 \tag{RT.50.9}\label{p1b:eq:rt-50-9}
\]
for every value of \(\epsilon_K\).

For \(E_Z^1\), homotope either gauge to identity on a slightly
larger trivializing ball around the fifth-class patch, with
the homotopy supported off \(J\). Contracting the domain ball
and then its constant value provides this homotopy, in \(SU(2)\)
for \(S_Z\). The resulting gauge pastes through the new
\(\eta_3\) gluing because it is identity near the seam.
Its low data and its exact joining restriction are preserved.
Theorem~\ref{p1b:rt:theorem50-bundles} transports these gauges
to every actual piece by an actual bundle isomorphism, keeping
its induced joining frame. Finally, \eqref{p1b:eq:rt-48-7} sends
\((0,4,0)\) to itself for every joining frame component;
this proves the sphere assertion without selecting a frame.
\end{proof}

The sphere piece is now supplied with actual gauge representatives
for both target columns. The second piece and their joining
homotopies are constructed in the following MAIN supplement.



\subsection{A construction of the two target gauges from the actual pieces}
\label{p1b:sec:rtproof51}

Let \(P=S_x^3\times S_p^3\) and \(B(x,p)=(p,x)\).
Cut the actual \(Y\) at one value of its primitive circle.
Its actual bundle then consists of a bundle \(E\to P\) and
an exchange isomorphism \(m:B^*E\to E\) closing the interval.
Use the convention that a family of gauges closes when
\[
 G(1)=\alpha_m(G(0)),\qquad
 \alpha_m(g)=m B^*g m^{-1}.
 \tag{RT.51.1}\label{p1b:eq:rt-51-1}
\]
The bundle \(E\) has actual G15 sector \(\varepsilon\in\{0,1\}\).
Its determinant is trivialized only for the fibre gauge calculation.
Neither this nor the interval cut trivializes \(E\) itself.

Use the Born-compatible Heisenberg chart supplied by G20--G21
for the actual determinant-one gauge group \(\mathcal G_{SU(2)}(E)\):
\[
 \pi_0\mathcal G_{SU(2)}(E)=
 \{(a,b;z):a,b\in\mathbb Z,\ z\in C_{12}\},\qquad
 (a,b;z)(a',b';z')=(a+a',b+b';z+z'+ab').
 \tag{RT.51.2}\label{p1b:eq:rt-51-2}
\]
Here a gauge component is an actual component of sections,
including in the exotic sector. The reference exchange
isomorphism \(\Phi\) has action
\(\sigma(a,b;z)=(b,a;ab-z)\), in the chart of
Section~\ref{p1b:app:g21}. The retained determinant boundary
also gives
\[
 {}[-I]=(0,0;6\varepsilon)
 \tag{RT.51.3}\label{p1b:eq:rt-51-3}
\]
by Lemma~\ref{p1b:lem:v2170-determinant-boundary}.

\begin{lemma}[The actual monodromy primary]
\label{p1b:rt:lemma51-monodromy}
Up to its scalar determinant factor, write \(m=u\Phi\) with
\(u\) an actual determinant-one gauge on \(E\). Then
\[
 {}[u]=(a,-a;z_m),\qquad
 \alpha_m=\operatorname{Ad}_u\sigma
 \quad\text{on actual gauge components}.
 \tag{RT.51.4}\label{p1b:eq:rt-51-4}
\]
The actual central monodromy coordinate \(z_m\) is retained
but does not enter the component formulas below.
\end{lemma}
\begin{proof}
Compare the actual closing isomorphism with the chosen reference.
Their difference is a unitary gauge on \(E\).
Its determinant has a logarithm since \(H^1(P;\mathbb Z)=0\).
Dividing by its scalar half makes the difference determinant one
without changing its conjugation action. A change of logarithm
only multiplies it by the central \(-I\) and does not change
its primary coordinates.

The reference exchange bundle has \(c_2=0\) on \(Y\):
on the ordered cover it is pulled back from the Chern-zero
G15 bundle, and \(H^4(Y;\mathbb Z)\) is free. If the difference
gauge has primary \((a,b)\), its second-Chern transgression in
the mapping torus is the coinvariant \(a+b\), up to the chosen
overall circle orientation. Equivalently the double-circle
ordered clutching has primary \((a+b,a+b)\), and the primitive
four-class has cover \(a_0(u_x+u_p)\), as in Subsection~\ref{p1b:sec:rtproof46}.
The actual \(c_2(V|_Y)\) is zero, since it is the restriction
of a torsion class to a free integral group. Thus \(a+b=0\).
The equality for \(\alpha_m\) follows directly from composition.
\end{proof}

\begin{theorem}[Two target gauges on the actual mapping torus]
\label{p1b:rt:theorem51-Ygauges}
There are actual gauges \(D_Y^0,S_Y^0\) on \(V|_Y\) of low
data \((h_Y,0)\) and \((0,2q_Y)\), respectively. They can
be constructed with the following fibre components:
\[
 \boxed{[F_h]=(0,0;3\varepsilon),\qquad
        {}[F_s]=(2,2;2+2a).}
 \tag{RT.51.5}\label{p1b:eq:rt-51-5}
\]
The alternative central choices add six to the displayed values.
Their construction uses actual representatives and actual gauge
homotopies; it does not replace the exotic bundle by maps into
an untwisted \(SU(2)\) fibre.
\end{theorem}
\begin{proof}
For a central component \((0,0;z)\), the reference reflection
and actual monodromy send it to \((0,0;-z)\).
The determinant-family closing equation will be
\[
 F_h(1)=-\alpha_m(F_h(0)).
 \tag{RT.51.6}\label{p1b:eq:rt-51-6}
\]
By \eqref{p1b:eq:rt-51-3}, its endpoints have the same actual component
precisely when \(2z=6\varepsilon\pmod{12}\), whose two
solutions are \(3\varepsilon\) and \(3\varepsilon+6\).
Choose an actual gauge in the first component and an actual
path in \(\mathcal G_{SU(2)}(E)\) to its endpoint in \eqref{p1b:eq:rt-51-6}.
The family
\[
 D_Y^0(t)=e^{\pi it}F_h(t)
 \tag{RT.51.7}\label{p1b:eq:rt-51-7}
\]
has determinant \(e^{2\pi it}\) and satisfies \eqref{p1b:eq:rt-51-1}.
The scalar half-path is defined on every actual bundle.
Its fibre three-primary is zero, hence its global three-primary
on \(Y\) is zero. This constructs the determinant gauge.

For a component \((2,2;z)\), \eqref{p1b:eq:rt-51-4} and the Heisenberg
commutator give
\[
 \alpha_m(2,2;z)=(2,2;4-z+4a).
 \tag{RT.51.8}\label{p1b:eq:rt-51-8}
\]
Thus its fixed-component equation is
\(2z=4+4a\pmod{12}\), with solutions
\(z=2+2a\) and \(z=8+2a\).
Choose an actual representative \(F_s\) of the first component
and an actual gauge homotopy from it to \(\alpha_m(F_s)\).
The resulting family \(S_Y^0(t)\) satisfies \eqref{p1b:eq:rt-51-1}, so
it is a determinant-one gauge on the whole actual mapping torus.
Its fibre degrees \((2,2)\) give primary \(2q_Y\).
Equality of components in the actual gauge group is exactly
the criterion for these two paths to exist. This is the
homotopy-class construction, with both path endpoints specified.
\end{proof}

\begin{lemma}[An actual joining-secondary toggle on \(Y\)]
\label{p1b:rt:lemma51-toggle}
There is a determinant-one gauge \(T_Y\) on the actual \(V|_Y\)
with zero integral three-primary and joining component
\[
 \boxed{[j_Y^*T_Y]=(0,0,1).}
 \tag{RT.51.9}\label{p1b:eq:rt-51-9}
\]
\end{lemma}
\begin{proof}
The actual fibre gauge group has a loop \(\alpha_x\) whose
four-stage on the suspended first sphere is \(\eta_3\), as
in the G20 mapping two-type, with no second-sphere coefficient.
Take this loop as the gauge family on the cut mapping torus.
Both endpoints are identity, so it closes under every actual
monodromy in \eqref{p1b:eq:rt-51-1}. It is determinant one and has zero
integral three-primary on every fibre. The degree-three
restriction from \(Y\) to its fibre is injective, so the
global primary is zero.

Its fourth cochain is the primitive mapping-torus coinvariant
represented by \(u_x\) modulo sphere exchange, reduced modulo
two. The joining map restricts that primitive four-class with
coefficient one, by \eqref{p1b:eq:rt-46-9}. On the four-dimensional \(J\)
this is exactly the nonzero supported \(\eta_3\) component.
There is no higher cell in \(J\) to alter it. Frame conjugation
fixes \((0,0,1)\) by \eqref{p1b:eq:rt-48-7}, proving \eqref{p1b:eq:rt-51-9}.
\end{proof}

\begin{theorem}[Target global gauge triples by actual matching]
\label{p1b:rt:theorem51-targets}
Retain the actual joining matrix in the frames induced by the
piece constructions, with component \((\ell,d,\epsilon_c)\).
Write the initial joining components of Theorem~\ref{p1b:rt:theorem51-Ygauges}
as \((1,0,e_h)\) and \((0,4,e_s)\). Define actual piece gauges
\[
 D_Y=D_Y^0T_Y^{e_h+d},\qquad
 S_Y=S_Y^0T_Y^{e_s},
 \tag{RT.51.10}\label{p1b:eq:rt-51-10}
\]
where exponents are reduced modulo two. Then
\[
 \boxed{[j_Y^*D_Y]=(1,0,d),\qquad
        {}[j_Y^*S_Y]=(0,4,0).}
 \tag{RT.51.11}\label{p1b:eq:rt-51-11}
\]
They match \(D_Z,S_Z\) of Theorem~\ref{p1b:rt:theorem50-gauges}
under the actual clutching. Choosing paths between these actual
endpoints constructs global gauge triples \(g_h,g_2\) with
the precise retained low data \((h,0)\), \((0,2q)\).
\end{theorem}
\begin{proof}
The joining-secondary toggle has zero low data and changes
exactly the last joining coordinate. The multiplication in
\eqref{p1b:eq:rt-48-3} therefore proves \eqref{p1b:eq:rt-51-11}.
Conjugation of the \(Z_s\) endpoints by the actual clutching is
\[
 (1,0,0)\longmapsto(1,0,d),\qquad
 (0,4,0)\longmapsto(0,4,0).
\]
Thus the exact criterion of Theorem~\ref{p1b:rt:theorem48-matching}
gives paths \(H_h,H_s\) between the two pairs of actual gauges.
Paste \((D_Y,D_Z,H_h)\) and \((S_Y,S_Z,H_s)\) using
Theorem~\ref{p1b:rt:theorem47-gauges}.

Their determinant and integral three-primary on both pieces
are those of \(h\) and \(2q\). The global restriction maps
in degrees one and three are injective for this pushout.
Consequently the low data are exactly the stated global data,
including the zero torsion-primary coordinate. This construction
uses representatives and homotopies in the actual piece bundles;
it is valid in both restriction sectors and both piece classes.
\end{proof}

The complete homotopy-class recipe is now specified: choose the
actual fibre representatives in \eqref{p1b:eq:rt-51-5}, their closing homotopies,
apply the joining toggle, and choose the joining paths. The subsequent closing and paired comparison constructions supply the actual homotopies needed for the natural evaluation.
The actual boundary value is obtained by this triple along the
three stages of \eqref{p1b:eq:rt-47-2}, with its prescribed parallel transport
to the Born boundary frame. At odd parity its normalized component
is independent of these construction choices at fixed low data,
by the retained zero pure finite kernel. The full natural comparison on \(R/K_R=S^7\), with all these operands retained, is evaluated in the final included natural chain.
Section~\ref{p1b:app:rt51-construction} records the finite residue
checks and the distinction between this gauge construction and
evaluation of the final natural boundary maps.



\subsection{Removing the actual zero-Chern mapping-torus monodromy}
\label{p1b:sec:rtproof52}

Retain the actual cut bundle \(E\to P=S^3\times S^3\), its
G15 sector \(\varepsilon\), and the reference exchange
\(\Phi:B^*E\to E\) used in Subsection~\ref{p1b:sec:rtproof51}. The current
G20--G21 chart of actual determinant-one gauge components has law
\[
 (p,q;r)(a,b;z)=(p+a,q+b;r+z+pb),\qquad z,r\in C_{12},
 \quad \sigma(p,q;r)=(q,p;pq-r).
 \tag{RT.52.1}\label{p1b:eq:rt-52-1}
\]
By Lemma~\ref{p1b:rt:lemma51-monodromy}, after removing a scalar
closing factor the actual closing map is \(m=u\Phi\), with
\([u]=(a,-a;z_m)\). Removing that scalar factor is an ordinary
bundle isomorphism: its logarithm gives the required scalar
path on the cut interval. No trivialization of the exotic \(E\)
is made.

\begin{lemma}[The exact change of the closing component]
\label{p1b:rt:lemma52-twisted}
Changing the fibre frame by an actual determinant-one gauge of
component \([v]=(p,q;r)\) replaces \(u\) by
\(u'=v u\sigma(v)^{-1}\), with component
\[
 \boxed{[u']=(a+p-q,-a-p+q;
                     z_m+2r-p^2-2ap).}
 \tag{RT.52.2}\label{p1b:eq:rt-52-2}
\]
In particular put
\[
 t_0=(z_m+a^2)\bmod2\in\{0,1\},\qquad
 r_0=(t_0^2-z_m-a^2)/2,
 \quad [v]=(t_0-a,t_0;r_0).
 \tag{RT.52.3}\label{p1b:eq:rt-52-3}
\]
Using any integral representative of \(z_m\), the number \(r_0\)
is an integer. This actual frame change sends \([u']\) to
the identity component \((0,0;0)\).
\end{lemma}
\begin{proof}
The closing map in the changed fibre frame is
\(v m B^*v^{-1}\). Its difference from \(\Phi\) is the
displayed twisted product. The inverse of \(\sigma(v)\) is
\((-q,-p;r)\), by \eqref{p1b:eq:rt-52-1}. Multiplying gives \eqref{p1b:eq:rt-52-2}.
For \(p=t_0-a,q=t_0\), both primary entries become zero and
the last entry is \(z_m+a^2-t_0^2+2r_0=0\).
The parity choice ensures integrality of \(r_0\).
Replacing the representative of \(z_m\) by twelve changes
\(r_0\) by minus six and leaves the conclusion unchanged.
Every component used here has an actual gauge representative
in the retained actual gauge group.
\end{proof}

\begin{theorem}[The actual ordinary \(Y\)-bundle normal form]
\label{p1b:rt:theorem52-Ynormal}
For each actual G15 sector \(\varepsilon\), every ordinary
rank-two bundle on \(Y\) with \(c_1=c_2=0\) and that fibre
sector is isomorphic to the reference exchange mapping-torus
bundle closed by \(\Phi\). Thus there are exactly two such
ordinary bundle classes in total, distinguished by the fibre sector.
\end{theorem}
\begin{proof}
Cutting any such bundle gives an actual bundle of the indicated
sector on \(P\). Identify it with the reference fibre bundle.
Its closing difference has a determinant logarithm, since
\(H^1(P;\mathbb Z)=0\); remove the scalar factor as above.
Zero second Chern class forces primary \((a,-a)\), by the
actual transgression in Lemma~\ref{p1b:rt:lemma51-monodromy}.
Lemma~\ref{p1b:rt:lemma52-twisted} then gives an actual frame
change with closing difference in the identity gauge component.
Choose an actual gauge path from identity to its inverse.
Using that path as an interval frame change identifies the
changed closing bundle with the one closed by \(\Phi\).
This is an ordinary bundle isomorphism, not merely a relation
between Chern classes. Restriction to \(P\) distinguishes
the two sectors, so the two reference classes cannot be isomorphic.
\end{proof}

\begin{corollary}[Fixed fibre components for both target constructions]
\label{p1b:rt:corollary52-targets}
After this actual piece isomorphism, the target fibre components
and their closing equations are
\[
 \boxed{
 \begin{array}{c|c|c}
  \text{target}&\text{closing equation}&\text{fibre components}\\\hline
  (h_Y,0)&2z=6\varepsilon\pmod{12}
             &(0,0;3\varepsilon),\ (0,0;3\varepsilon+6)\\
  (0,2q_Y)&2z=4\pmod{12}&(2,2;2),\ (2,2;8).
 \end{array}}
 \tag{RT.52.4}\label{p1b:eq:rt-52-4}
\]
The unknown monodromy primary and central coordinates are no
longer needed as inputs to these reference-piece equations.
\end{corollary}
\begin{proof}
The reference action is now exactly \(\sigma\).
The determinant construction has seam \(-\sigma\) and
\([-I]=(0,0;6\varepsilon)\); the sphere construction has
seam \(\sigma\). Substitute \(a=0\) in the component
calculations of Theorem~\ref{p1b:rt:theorem51-Ygauges}.
Ordinary bundle isomorphisms transport actual gauges and their
low data, so this is a permissible replacement of the piece
construction with the induced joining identification retained.
\end{proof}

The normal form is on the actual piece bundle. Its joining
frame is the one induced by the actual isomorphism; the actual
clutching matrix and the identification with the prescribed Born
boundary frame must be transported with it. No identification
with that boundary frame, and no value of \(h_t,w_t\), follows
from \eqref{p1b:eq:rt-52-3}. The next supplement computes the remaining
closing-family choices in this reference piece.



\subsection{The finite closing-family source for the target comparison}
\label{p1b:sec:rtproof53}

Work on the actual reference \(Y\)-piece of
Theorem~\ref{p1b:rt:theorem52-Ynormal} and keep its actual
fibre gauge group \(\mathcal G=\mathcal G_{SU(2)}(E)\).
G20--G21 supply
\[
 \pi_1\mathcal G=C_2\{\alpha_x\}\oplus
 C_2\{\alpha_p\}\oplus C_2\{\beta\},\qquad
 \sigma_*(r,s,e)=(s,r,e).
 \tag{RT.53.1}\label{p1b:eq:rt-53-1}
\]
Here \(\beta\) has its supported parameter--fibre top coefficient
\(\mu\); it is not the natural boundary coefficient \(t\).
For the determinant target fix the exact determinant function
\(e^{2\pi it}\) on the cut primitive circle, and write the family
as \(e^{\pi it}F(t)\). Its reduced family has seam
\(F(1)=-\sigma F(0)\). For the sphere target fix determinant
identically one; its seam is \(F(1)=\sigma F(0)\).
All homotopies in the following component count preserve those
exact determinant functions. These are not counts of ordinary
gauge classes after allowing arbitrary determinant homotopies.

\begin{lemma}[Closing paths modulo moving their initial representative]
\label{p1b:rt:lemma53-paths}
For each permitted fibre component in \eqref{p1b:eq:rt-52-4}, the reduced
closing-family components form a torsor for
\[
 \boxed{Q_Y=\pi_1\mathcal G/(1-\sigma_*)\pi_1\mathcal G
       =C_2\{\bar\alpha\}\oplus C_2\{\bar\beta\},}
 \quad \bar\alpha=[\alpha_x]=[\alpha_p].
 \tag{RT.53.2}\label{p1b:eq:rt-53-2}
\]
Consequently each target has exactly eight determinant-fixed
closing-family components: two permitted fibre components,
and four path components over each.
\end{lemma}
\begin{proof}
Fix an actual representative \(g\) in one permitted component
and an actual path from it to its prescribed seam endpoint.
Paths between those fixed endpoints form a torsor for
\(\pi_1\mathcal G\). A homotopy of closing families may move
the initial representative around a loop \(L\) in its component.
The endpoint then moves around \(\sigma L\), with the same
action for the determinant seam: multiplication by the central
\(-I\) is a translation of its endpoint component.
The resulting path difference is \((1-\sigma_*)[L]\).
Conversely such a loop supplies exactly that homotopy by the
path-space lifting construction, so these are all identifications.

For clarity, use left translation to identify fundamental groups
of gauge components with \(\pi_1\mathcal G\) at identity.
A gauge-group identity component is an H-space with trivial
fundamental-group action on itself. Thus transport along the
chosen endpoint path agrees on fundamental groups with left
translation, and introduces no additional conjugation by \(g\).
The quotient is precisely \eqref{p1b:eq:rt-53-2}, even when the target
fibre component is not central. Its denominator is the one line
\(\langle\alpha_x+\alpha_p\rangle\).
Finally the two fibre components differ by six and cannot be
joined by a determinant-fixed family homotopy evaluated at
\(t=0\). Their four-element torsors are disjoint, giving eight.
\end{proof}

\begin{theorem}[Evaluated joining map of the closing-path torsor]
\label{p1b:rt:theorem53-joining}
Changing a closing family by a class \((x,y)\in Q_Y\), after
choosing any one actual closing path as origin, changes its
joining secondary coefficient by exactly \(x\):
\[
 \boxed{Q_Y\longrightarrow C_2\{\epsilon_J\},\qquad
        (x,y)\longmapsto x.}
 \tag{RT.53.3}\label{p1b:eq:rt-53-3}
\]
In particular the \(\bar\beta\) change has identity joining
component, and the \(\bar\alpha\) change toggles that component.
\end{theorem}
\begin{proof}
The actual loop \(\alpha_x\), taken as a cut gauge family
with identity endpoints, constructs \(T_Y\) of
Lemma~\ref{p1b:rt:lemma51-toggle}. Its exact joining component
is \((0,0,1)\). The same holds for \(\alpha_p\): its
fourth coinvariant is the same primitive class under exchange.
Their sum is zero on \(J\), so restriction descends to \eqref{p1b:eq:rt-53-2}.

The retained \(\beta\) loop has identity lower section data
through the parameter--fibre six-skeleton and first coefficient
on the seven-cell. Taking it as a cut family therefore gives
an actual zero-low \(Y\)-gauge with no fourth difference.
Restriction to the four-dimensional \(J\) is homotopic to
identity. There are no fifth or higher joining cells to contribute.
This argument uses the actual supported loop in the actual
bundle, not a substituted untwisted sphere bundle. The joining
group multiplication and its conjugation fix \((0,0,1)\) and
the identity, so the conclusions hold in every induced joining
frame. Applying them to the difference from the chosen closing
origin proves \eqref{p1b:eq:rt-53-3}.
\end{proof}

\begin{corollary}[Exactly four matching \(Y\)-classes per target]
\label{p1b:rt:corollary53-matched}
Fix \(D_Z,S_Z\) from Subsection~\ref{p1b:sec:rtproof50} and the actual clutching
\((\ell,d,\epsilon_c)\) in the induced piece frames.
For each target there are exactly four determinant-fixed
\(Y\) closing-family components satisfying the joining test:
two over each permitted fibre component. If its chosen origin
has joining secondary value \(e_z\), the matching paths are
\[
 \boxed{
 \begin{array}{c|c|c}
 \text{target}&\bar\alpha\text{ coordinate}&
                         \bar\beta\text{ coordinate}\\\hline
 D_Y&e_z+d\bmod2&\text{arbitrary}\\
 S_Y&e_z&\text{arbitrary}.
 \end{array}}
 \tag{RT.53.4}\label{p1b:eq:rt-53-4}
\]
After matching, all these choices, and all joining paths for
them, give the same normalized boundary component at odd parity
for the corresponding fixed global low data.
\end{corollary}
\begin{proof}
The clutch-conjugated \(Z_s\) endpoints are \((1,0,d)\) and
\((0,4,0)\), by the actual joining-component law in Subsection~\ref{p1b:sec:rtproof48}. The joining map
\eqref{p1b:eq:rt-53-3} is onto with two-element kernel
\(\langle\bar\beta\rangle\). This gives exactly two matching
classes over either permitted fibre component, hence four.
The displayed affine origin \(e_z\) can differ between those
fibre components and is not assigned numerically.

For any two resulting complete global triples, their quotient
is an actual zero-low gauge on the same actual global bundle.
The retained odd pure finite boundary image is zero, so their
normalized boundary values agree. This includes the difference
between the two fibre-component roots after matching, as well
as the \(\bar\beta\) and joining-path choices. The actual
\(\beta\)-change with identity joining component can also be
pasted to the identity \(Z_s\) gauge; it is such a zero-low
global gauge directly. This is boundary independence, not a
calculation of that common boundary component.
\end{proof}

\begin{lemma}[The unsquared sphere-piece frame gauge]
\label{p1b:rt:lemma53-Zroot}
On either actual \(Z_s\) piece class there is an ordinary
determinant-one gauge \(K_Z\) of primary \(v_Z\) and
joining component \((0,2,\epsilon_K)\) for some
\(\epsilon_K\in C_2\).
\end{lemma}
\begin{proof}
Use the unsquared \(\pi_Z^*K\) in the proof of
Theorem~\ref{p1b:rt:theorem50-gauges}. On the second piece
class, homotope it to identity near the five-ball patch and
paste it through the new gluing exactly as in that proof.
This preserves its low data and joining restriction. An
ordinary joining frame change fixes its component because
its determinant coordinate is zero and its sphere coordinate
is two. Its secondary value need not be assigned.
\end{proof}

\begin{theorem}[An ordinary clutching normal form on the entire free region]
\label{p1b:rt:theorem53-Wnormal}
Fix the actual reference \(Y\) piece in sector \(\varepsilon\)
and either \(Z_s\) piece class \(E_Z^j\), \(j=0,1\).
Choose an ordinary zero-second-Chern reference gluing and
joining frames in which its gluing is identity. In these
reference frames every actual gluing is equivalent, by
actual piece bundle automorphisms and a clutching homotopy,
to
\[
 \boxed{c_a=f_1^a,\qquad [c_a]=(0,a,0),\qquad
        a\in\{0,1\},\quad f_1[a_0,w]=a_0^{-1}w.}
 \tag{RT.53.5}\label{p1b:eq:rt-53-5}
\]
Here the exponent \(a\) is the second-Chern parity, so
\(a=A\bmod2\). It is distinct from the monodromy integer
of Subsection~\ref{p1b:sec:rtproof51}. There is exactly one ordinary gluing class
for each such parity and piece pair. Thus the ordinary
bundles on \(W\) with \(c_1=\tau,c_2=a\tau^2\) have
exactly four classes, represented by the two \(Y\) sectors
and the two \(Z_s\) piece classes with gluing \eqref{p1b:eq:rt-53-5}.
\end{theorem}
\begin{proof}
The zero-Chern reference exists by the actual transgression
of Theorem~\ref{p1b:rt:theorem47-clutch}, for every fixed
piece pair: an odd \(f_1\) change toggles \(c_2\).
Adapt ordinary joining frames to that reference only for
this bundle-classification calculation.

The \(Y\) piece has an actual gauge of joining component
\((1,0,0)\): adjust its determinant gauge by \(T_Y\) if
needed. It also has \(T_Y\) of component \((0,0,1)\).
The \(Z_s\) piece has \(K_Z\) from
Lemma~\ref{p1b:rt:lemma53-Zroot}. Write an arbitrary
joining component as \((\ell,d,\epsilon_c)\), set
\(a=d\bmod2\in\{0,1\}\), and put \(r=(a-d)/2\).
Right multiplication by \(K_Z^r\) changes it to
\((\ell,a,\epsilon_c+r\epsilon_K)\).
Left multiplication by the \(Y\) determinant gauge to
power \(-\ell\), and then by the requisite \(T_Y\),
changes it to \((0,a,0)\). The exact joining-component law of Subsection~\ref{p1b:sec:rtproof48} shows
these equalities, including negative powers. Equality of
joining components gives an actual clutching homotopy,
so the reduction is an ordinary global bundle isomorphism.

Every ordinary piece gauge has even joining sphere degree:
its integral degree-three primary restricts with coefficient
two on either piece, by Lemma~\ref{p1b:rt:lemma46-matrices}.
Therefore left and right piece automorphisms cannot change
the parity of \(d\). This proves precisely one gluing
class per parity, with \(c_2=a\tau^2\).
By Theorem~\ref{p1b:rt:theorem52-Ynormal} and Subsection~\ref{p1b:sec:rtproof50},
there are two choices for each piece. Restriction to the
pieces distinguishes their four pairs. Conversely any
ordinary \(W\) bundle with the stated Chern classes has
those piece Chern classes and hence belongs to one of
these pairs, so the list is exhaustive. In conjunction
with the retained four-compatible-bundle theorem it also
represents the four compatible classes at any fixed \(A\).
It does not supply their prescribed boundary isomorphisms.
\end{proof}

The actual piece-monodromy coordinates have now been removed
from the reference equations, and the determinant-fixed
closing-path quotient and its joining map are evaluated.
Choose the first root for each target in \eqref{p1b:eq:rt-52-4}, either
matching representative of its two-element path torsor, and
a joining path. These choices supply the same odd natural
boundary answer. The final natural proof evaluates that answer using
the actual comparison along \eqref{p1b:eq:rt-47-2}, carrying the induced frame
to the prescribed Born boundary frame on \(R/K_R=S^7\).
Alternatively begin with the ordinary global representatives
\eqref{p1b:eq:rt-53-5}, transporting all their induced frame identifications.
The complete retained comparison below proves \(h_t=w_t=1\) and
\(\zeta_{\rm nat}=h_t+w_t=0\).



\subsection{Ordinary boundary frames preserve the remaining top coefficients}
\label{p1b:sec:rtproof54}

Keep the base coordinates, lens cells, and product normal frame of
\(R=F\times L\), \(F=S^1_s\times S^3_y\), \(L=\mathbb {RP}^3\).
The boundary bundle is the actual
\(E_\partial=\mathbf1\oplus\lambda_L\).
An ordinary boundary frame change below is an automorphism of this
bundle covering the identity of \(R\). It does not change the base
normal coordinates. Put
\(H=\mathcal G_{U(2)}(\mathbf1\oplus\lambda_L)\), so its
boundary frame group is \(\operatorname {Map}(F,H)\).
Retain \(\pi_4H=C_2\{e,t\}\), the central even circle subgroup,
and the natural generators \(b,z\) from the current boundary owner.

\begin{lemma}[The normal component and the only needed circle commutators]
\label{p1b:rt:lemma54-factors}
Every ordinary boundary frame has, at the component level, a
factorization into a normal representative \(h_N\), a circle
factor, a sphere factor, and a parameter-four-cell factor.
The normal integer \(N\in\mathbb Z\) is the component of its
evaluation at the basepoint of \(F\). All the factors except
\(h_N\) conjugate the following subgroup trivially:
\[
 \mathscr B=
 2\mathbb Z\{q_+\}\oplus C_4\{\kappa\}\oplus
 2\mathbb Z\{b\}\oplus C_6\{z\}\oplus C_2\{e,t\}.
 \tag{RT.54.1}\label{p1b:eq:rt-54-1}
\]
In particular, this assertion does not require evaluation of
the possibly nonzero \(\langle q_+,b\rangle\).
\end{lemma}
\begin{proof}
A determinant on \(L\) has a logarithm because
\(H^1(L;\mathbb Z)=0\). Removing its scalar half reduces a
normal gauge to a determinant-one section. On this
three-dimensional base, its component is exactly its
\(H^3(L;\pi_3SU(2))=\mathbb Z\) difference from identity:
the fibre is two-connected and conjugation preserves its
orientation. The determinant-loop ambiguity introduces no
quotient, since the globally defined eigenline loop \(q_+\)
lifts the generator of \(\pi_1\operatorname {Map}(L,U(1))\).
A supported three-cell representative of difference \(N\)
has ordered-cover degree \(2N\). Thus these are exactly the
normal components used in Lemma~\ref{p1b:rt:lemma36-conjugation}.

After division by \(h_N\), a frame takes values in the identity
component of \(H\). Restrict it to
\((S^1\times\{y_0\})\cup(\{s_0\}\times S^3)\).
Products of its based circle and sphere restrictions extend
over \(F\). Homotopy extension makes the residual factor
identity on this wedge, hence a factor from
\(F/F_\vee=S^4\). Only existence of this factorization is
needed; no new choice of a canonical four-cell origin is made.

The only positive-parameter commutator that can survive on
\(F\) is the circle--sphere Samelson product in \(\pi_4H\).
Its value on \(2b\) is zero: bilinearity doubles an element
of the exponent-two group \(\pi_4H\). The \(\kappa\)
circle factor is already central. For the other sphere
generator the exact calculation is
\[
 \boxed{\langle q_+,z\rangle
       =j_7(A_*\nu')=0.}
 \tag{RT.54.2}\label{p1b:eq:rt-54-2}
\]
Indeed \(z\) is identity off the lens three-cell and has
coefficient \(\nu'\) on the smash of that cell with
\(S^3_y\). In its ordering-line frame the \(q_+\)
commutator is the actual winding-one \(U(2)\) adjoint
operation on this coefficient. Lemma~\ref{p1b:rt:lemma104}
gives \(A_*\nu'=\eta_3E\nu'=0\), rather than
\(\nu'\eta_6=\mu\). There are no lower lens cells in this
supported commutator to add another term.

Sphere--sphere factors require parameter dimension six,
and a four-cell factor commuted with a positive-dimensional
factor requires at least five. These exceed \(\dim F=4\).
Circle--circle factors commute in \(\pi_1H\).
The even circle subgroup is central, also across normal
components, by the retained boundary calculation. These facts
prove the assertion for arbitrary circle, sphere, and top
factors, including their products.
\end{proof}

\begin{theorem}[Complete ordinary frame action on the target subgroup]
\label{p1b:rt:theorem54-transport}
Let \(r\) be any ordinary boundary bundle frame, with normal
component \(N\). With conjugation convention
\(\operatorname {Ad}_r(g)=rgr^{-1}\), its exact action is
\[
 \boxed{\begin{split}
 \operatorname {Ad}_r\bigl(2kq_++p\kappa+
       \gamma(2mb+cz+\epsilon e+\tau t)\bigr)
  ={}&2kq_++p\kappa+\\
    &\gamma(2mb+(c+4Nm)z+\epsilon e+\tau t).
 \end{split}}
 \tag{RT.54.3}\label{p1b:eq:rt-54-3}
\]
Here \(k,m\in\mathbb Z\), \(p\in C_4\), \(c\in C_6\),
and \(\epsilon,\tau\in C_2\). Inverse conjugation replaces
\(N\) by \(-N\). In particular all ordinary frame changes
preserve the \(e,t\) coefficients of both target columns.
\end{theorem}
\begin{proof}
All identity-component parameter factors act trivially on
\(\mathscr B\), by Lemma~\ref{p1b:rt:lemma54-factors}.
For the remaining normal factor, the actual retained law is
\(b\mapsto b+2Nz\), with \(z,e,t,2q_+,\kappa\) fixed,
by Lemma~\ref{p1b:rt:lemma36-conjugation}.
Apply it to \(2mb\). This gives the coefficient \(4Nm\).
The even circle subgroup has no multiplication carry, giving
exactly \eqref{p1b:eq:rt-54-3} on the actual component group.
\end{proof}

\begin{corollary}[A convenient ordinary boundary identification suffices]
\label{p1b:rt:corollary54-top}
At odd parity the exotic determinant and shared sphere columns,
after the known even local-circle correction, are
\[
 D=2q_++\gamma(3z+h_tt),\qquad
 S_A=\gamma(2b+z_Az+w_tt),\quad z_A=5-2A\pmod6.
 \tag{RT.54.4}\label{p1b:eq:rt-54-4}
\]
Any two ordinary identifications of an actual model boundary
bundle with \(E_\partial\), covering the fixed \(R\), give
the same two natural top coefficients \(h_t,w_t\).
The sphere coefficient changes only by \(z_A\mapsto z_A+4N\),
equivalently by \(A\mapsto A-2N\) in this conjugation convention.
Consequently an explicit coordinate formula for the ordinary
boundary isomorphism induced by Subsections~\ref{p1b:sec:rtproof52}--\ref{p1b:sec:rtproof53} is not needed
to extract the two top coefficients. An actual ordinary
identification still has to be used; the twisted bundle and
the base normal coordinates remain fixed.
\end{corollary}
\begin{proof}
Two complete identifications differ by an actual boundary
automorphism \(r\). Apply \eqref{p1b:eq:rt-54-3}, with \(m=0\) for
\(D\) and \(m=1\) for \(S_A\). Their top coefficients
are unchanged for every normal integer and every parameter
factor of \(r\). Since
\(5-2(A-2N)=z_A+4N\), only the stated already-known
sixth coordinate needs to be transported.
\end{proof}

This removes the coordinate formula for an ordinary bundle
frame identification from the remaining top calculation.
It does not evaluate either top coefficient. Base normal-frame
rotations still act by the distinct shear of Subsection~\ref{p1b:sec:rtproof44}; that
operation is not conjugation by a bundle frame over identity.
The next MAIN supplement gives a single combined target for
their invariant sum, with no finite sixth-stage term at \(A=1\).



\subsection{A combined target with cancelled sixth-stage finite term}
\label{p1b:sec:rtproof55}

Work on one actual odd exotic compatible free-region bundle.
Use the target global gauges of Subsections~\ref{p1b:sec:rtproof50}--\ref{p1b:sec:rtproof53}, or any other
actual gauges of the same low data. Correct their even
\(2\kappa\) discrepancy by the known fourth gauge, as in
Theorem~\ref{p1b:rt:theorem41-stages}, so their whole boundary
components are \eqref{p1b:eq:rt-54-4}. This correction has zero normalized
finite value at odd parity. Define the actual composite gauge
\[
 C=g_hg_2^{-1},\qquad \text{low data }(h,-2q).
 \tag{RT.55.1}\label{p1b:eq:rt-55-1}
\]
Both factors here are on the same actual bundle.
Put \(\zeta_{\rm nat}=h_t+w_t\); it is not the local circle \(2\kappa\).

\begin{theorem}[The combined boundary column and its cancellation]
\label{p1b:rt:theorem55-cancel}
The exact boundary component of \eqref{p1b:eq:rt-55-1} is
\[
 \boxed{C|_R=2q_++\gamma\bigl(-2b+(2A-2)z+\zeta_{\rm nat} t\bigr).}
 \tag{RT.55.2}\label{p1b:eq:rt-55-2}
\]
At the single representative \(A=1\) it is
\[
 \boxed{C|_R=2q_++\gamma(-2b+\zeta_{\rm nat} t).}
 \tag{RT.55.3}\label{p1b:eq:rt-55-3}
\]
Thus the finite sixth-stage and natural \(e\) terms have
cancelled. The value of \(\zeta_{\rm nat}\) is unchanged by all ordinary
boundary bundle frame changes and by the normal-coordinate
rotation of Subsection~\ref{p1b:sec:rtproof44}.
\end{theorem}
\begin{proof}
The two factors in \eqref{p1b:eq:rt-54-4} lie in the additive target subgroup
with central even circle. Hence their quotient has finite sixth
coefficient \(3-z_A=2A-2\), no \(e\) coefficient, and top
coefficient \(h_t-w_t=h_t+w_t\). This proves \eqref{p1b:eq:rt-55-2}.
The shared odd coefficients of Subsection~\ref{p1b:sec:rtproof36} permit \(A=1\).
Alternatively for \(A=1+2M\), an ordinary frame of normal
component \(M\) changes \(-2b\) by \(-4Mz\), cancelling
the displayed \(4Mz\) term without changing \(\zeta_{\rm nat}\).
The known sixth coefficient must be tracked in an arbitrary
frame; its vanishing is asserted in the representative \eqref{p1b:eq:rt-55-3}.

Theorem~\ref{p1b:rt:theorem54-transport} proves ordinary-frame
invariance of the top coefficient. Subsection~\ref{p1b:sec:rtproof45} changes both
\(h_t,w_t\) by the same bit under its base normal-coordinate
rotation, so their sum is also invariant under that operation.
\end{proof}

\begin{theorem}[A direct actual closing-family construction for the composite]
\label{p1b:rt:theorem55-source}
On the exotic reference \(Y\) piece the composite target has
exact determinant \(e^{2\pi it}\) on the primitive cut circle.
Writing it \(e^{\pi it}F(t)\), its fibre and seam equations are
\[
 \boxed{[F(0)]=(-2,-2;z),\quad
 F(1)=-\sigma F(0),\quad 2z=10\pmod{12},\quad z=5,11.}
 \tag{RT.55.4}\label{p1b:eq:rt-55-4}
\]
For each root its determinant-fixed closing-family components
form the same \(Q_Y=C_2\{\bar\alpha,\bar\beta\}\) torsor
as in Subsection~\ref{p1b:sec:rtproof53}. There are eight such components in total.

In the actual sphere-piece frames of Subsection~\ref{p1b:sec:rtproof50} use
\(C_Z=D_ZS_Z^{-1}\), with joining component \((1,-4,0)\).
For the actual clutching \(c=(\ell,d,\epsilon_c)\), its
conjugated endpoint is \((1,-4,d\bmod2)\).
If a chosen \(Y\) closing origin has joining coefficient
\(e_z\), precisely the choices
\[
 \boxed{\bar\alpha=e_z+d\pmod2,\qquad
        \bar\beta\text{ arbitrary}}
 \tag{RT.55.5}\label{p1b:eq:rt-55-5}
\]
match it. There are four matching \(Y\) components, two per
root. A joining homotopy pastes an actual global composite
target. At odd parity all completed choices have the same
normalized boundary value for low data \((h,-2q)\).
\end{theorem}
\begin{proof}
The retained component action is
\(\sigma(-2,-2;z)=(-2,-2;4-z)\), and
\([-I]=(0,0;6)\) in the exotic fibre. Thus the determinant
seam equation gives \(z=10-z\), with exactly the stated roots.
They also follow directly by multiplying an actual determinant
root \((0,0;3)\) by the inverse of the sphere root
\((2,2;2)\): the latter inverse is \((-2,-2;2)\), so the
product is \((-2,-2;5)\). Changing either root by six gives
the other combined root. Every root has an actual representative,
and equality of its seam components gives an actual closing path.

The path-space argument of Lemma~\ref{p1b:rt:lemma53-paths}
applies to this same central-minus seam, even though the fibre
primary is now negative. Its denominator is still
\((1-\sigma_*)\pi_1\mathcal G\), so it has four path
components per root. The joining change is \((x,y)\mapsto x\)
by Theorem~\ref{p1b:rt:theorem53-joining}.

The exact joining law gives
\((1,0,0)(0,-4,0)=(1,-4,0)\). Clutch conjugation adds
\(\ell(-4)+1\cdot d=d\bmod2\). This proves \eqref{p1b:eq:rt-55-5}
and its two-element kernel per root. The joining component
equalities supply an actual endpoint homotopy and hence an
actual global gauge triple by Theorem~\ref{p1b:rt:theorem47-gauges}.
These are the induced target frames; they are not silently
identified with the separate zero-Chern-reference frames used
for the ordinary clutching normal form in Subsection~\ref{p1b:sec:rtproof53}.

For any two complete composite targets on the same bundle,
their quotient has zero global low data. Its odd pure finite
boundary image is zero. Thus the normalized boundary component
is independent of the root, matched path, and joining choices.
It agrees with \eqref{p1b:eq:rt-55-2} obtained by the product construction.
\end{proof}

\begin{corollary}[One seven-sphere test for the invariant sum]
\label{p1b:rt:corollary55-test}
At \(A=1\) choose the specified natural comparison candidates
\[
 C_\tau=2q_++\gamma(-2b+\tau t),\qquad \tau\in C_2.
 \tag{RT.55.6}\label{p1b:eq:rt-55-6}
\]
Both extend through \(W^{(7)}\cup R\). Their exact last
obstruction and ordinary extension test are
\[
 \boxed{O_7(C_\tau)=0,\qquad
        O_8(C_\tau)=(\tau+\zeta_{\rm nat})\Theta,\qquad
        C_\tau\text{ extends over }W\ \Longleftrightarrow\ \tau=\zeta_{\rm nat}.}
 \tag{RT.55.7}\label{p1b:eq:rt-55-7}
\]
Equivalently restrict one directly constructed composite triple
along the explicit map \eqref{p1b:eq:rt-47-2}, in any actual ordinary
boundary identification covering the fixed base, and divide by
\(C_0\). Remove its known low comparison. The resulting
supported map is
\[
 R/K_R=S^7\longrightarrow SU(2),\qquad [\text{map}]=\zeta_{\rm nat}\mu.
 \tag{RT.55.8}\label{p1b:eq:rt-55-8}
\]
No independent finite sixth-stage insertion is needed in this
combined test. Evaluating this one map gives \(\zeta_{\rm nat}\);
evaluating one additional natural target map then gives both
\(h_t,w_t\). The number of independent missing bits is still two.
\end{corollary}
\begin{proof}
Divide \(C_\tau\) by the actual target of \eqref{p1b:eq:rt-55-3}.
The residual boundary component is exactly
\((\tau+\zeta_{\rm nat})t\). The retained pure finite calculation
has zero seventh-stage obstruction on \(t\), and last
obstruction \(\Theta\) on \(t\) at odd parity, with no
lower-choice correction. This proves \eqref{p1b:eq:rt-55-7}.

The quotient \(R/K_R=S^7\) and its natural supported
insertion \(j_7\mu=t\ne0\) are those of
Corollary~\ref{p1b:rt:corollary47-top}. The actual low comparison
and the zero finite sixth and \(e\) comparisons leave only
the stated supported coefficient. The injection on this
\(C_2\) coefficient makes it unique. An arbitrary ordinary
boundary frame can first be corrected for its known sixth
shear using Theorem~\ref{p1b:rt:theorem54-transport}; this
does not change the coefficient being tested.
Finally \((h_t,w_t)\mapsto(h_t,h_t+w_t)\) is an invertible
linear change of two binary inputs. It gives no value for
either input by itself.
\end{proof}

The composite triple is constructed in the actual homotopy
category, with its comparison source specified without an uncancelled
finite sixth term. The complete original-map proof below evaluates
this supported top map and gives \(h_t=w_t=1\) and
\(\zeta_{\rm nat}=0\). The resulting natural target ledger is 4 YES/10
NO and the inherited candidate ledger 4 YES/28 NO, with all entries
decided by the retained obstruction formulas.



\subsection{Stabilization and an exact mod-two detector for the last bit}
\label{p1b:sec:rtproof56}

Keep the actual compatible rank-two bundle, the boundary base coordinates,
and the natural insertions of Lemma~\ref{p1b:rt:lemma104}.
Write \(\mu=\nu'\eta_6\) for the nonzero element of
\(\pi_7SU(2)=C_2\). Adding an auxiliary trivial line below is a
calculation on \(V\oplus\mathbf1\); it does not change the rank
of the object being classified.

\begin{lemma}[The parity boundary in the three-frame fibration]
\label{p1b:rt:lemma56-parity}
Let \(i:SU(2)\to SU(3)\) be \(g\mapsto\operatorname{diag}(g,1)\)
and let \(p:SU(3)\to S^5\) take the last column. Its exact
homotopy sequence in the relevant degrees can be written
\[
 C_{12}\xrightarrow{\,r\mapsto2r\,}
 C_{24}\{\nu_5\}\xrightarrow{\,n\mapsto(n\bmod2)\mu\,}
 C_2\{\mu\}\longrightarrow0.
 \tag{RT.56.1}\label{p1b:eq:rt-56-1}
\]
In particular \(i_*:\pi_7SU(2)\to\pi_7SU(3)\) is zero.
\end{lemma}
\begin{proof}
The canonical fibration is \(SU(2)\to SU(3)\to S^5\).
The source groups \(\pi_8S^5=C_{24}\{\nu_5\}\),
\(E\nu_4=\nu_5\), and \(\pi_7SU(2)=C_2\{\mu\}\)
are retained in Lemma~\ref{p1b:rt:lemma104}.
Maruyama--Oshima, \textit{Homotopy groups of the spaces of
self-maps of Lie groups}, arXiv:0712.2080, Table 1, gives
\(\pi_8SU(3)=C_4\oplus C_3=C_{12}\) and
\(\pi_7SU(3)=0\). Exactness makes the boundary onto \(C_2\).
The only onto homomorphism \(C_{24}\to C_2\) is parity.
Its kernel is the even subgroup, of order twelve.
The image of \(p_*\) is that subgroup, so \(p_*\) is injective
on \(C_{12}\); choose its generator to project to \(2\nu_5\).
This proves \eqref{p1b:eq:rt-56-1} without imposing a choice of odd sign.
\end{proof}

\begin{theorem}[A mapping-cone square detects exactly the rank-two bit]
\label{p1b:rt:theorem56-detector}
Let \(\varphi:S^7\to SU(2)\) be an actual based supported
comparison map. Choose any extension
\(H:D^8\to SU(3)\) with \(H|_{S^7}=i\varphi\).
Projection is constant on the boundary and gives
\[
 g_H:D^8/S^7=S^8\longrightarrow S^5,\qquad
 X_H=S^5\cup_{g_H}e^9 .
 \tag{RT.56.2}\label{p1b:eq:rt-56-2}
\]
For the bottom and top mod-two generators \(u_5,v_9\) of
\(X_H\), define \(\epsilon(H)\) by
\[
 \operatorname{Sq}^4u_5=\epsilon(H)v_9.
 \tag{RT.56.3}\label{p1b:eq:rt-56-3}
\]
Then
\[
 \boxed{[\varphi]=\epsilon(H)\mu.}
 \tag{RT.56.4}\label{p1b:eq:rt-56-4}
\]
The extension \(H\) always exists. The bit \(\epsilon(H)\)
is independent of that extension.
\end{theorem}
\begin{proof}
Existence follows from \(\pi_7SU(3)=0\).
If \([g_H]=n\nu_5\), the definition of the fibration boundary
and Lemma~\ref{p1b:rt:lemma56-parity} give
\([\varphi]=(n\bmod2)\mu\). Another extension glued to \(H\)
along its boundary gives an element of \(\pi_8SU(3)\).
Its projection changes \(n\) by an even element of \(C_{24}\).
Thus its parity is fixed.

The cone of the quaternionic Hopf map \(\nu_4:S^7\to S^4\)
is \(\mathbb H P^2\), with bottom generator \(u_4\) satisfying
\(\operatorname{Sq}^4u_4=u_4^2=v_8\).
Since \(\nu_5=E\nu_4\), its cone is homotopy equivalent to
\(\Sigma\mathbb H P^2\). Suspension stability gives
\(\operatorname{Sq}^4u_5=v_9\) on this cone.
These standard facts also appear in Schwede,
\textit{Cohomology operations and the Steenrod algebra},
Examples 5.1--5.2.

Represent \(n\nu_5\) by \(\nu_5\) precomposed with the
degree-\(n\) map of \(S^8\). The induced map from its cone
to the cone of \(\nu_5\) is the identity on \(S^5\) and
has degree \(n\) on the top quotient \(S^9\).
Naturality of \(\operatorname{Sq}^4\) therefore gives
\(\operatorname{Sq}^4u_5=(n\bmod2)v_9\).
Homotopic attaching maps give the same relation.
Combining this with the boundary calculation proves \eqref{p1b:eq:rt-56-4}
and its stated independence.
\end{proof}

\begin{theorem}[A cellwise square formula for the relative obstruction]
\label{p1b:rt:theorem56-relative}
Fix an actual boundary candidate already extended as a \(U(2)\)
gauge through \(W^{(7)}\cup R\). Its block stabilization
extends as a \(U(3)\) gauge over \(W\).

On each relative eight-cell choose a trivialization of \(V\)
and remove the determinant of its attaching gauge by a local
logarithm. This leaves \(\varphi_e:S^7\to SU(2)\).
Normalize a rank-three extension in the same determinant
convention and form \(X_{H_e}\) as in \eqref{p1b:eq:rt-56-2}.
Then the cochain
\[
 \epsilon_e=\text{coefficient of }v_9
       \text{ in }\operatorname{Sq}^4u_5
       \text{ on }X_{H_e}
 \tag{RT.56.5}\label{p1b:eq:rt-56-5}
\]
represents the original \(U(2)\) obstruction
\(O_8\in H^8(W,R;C_2)\). In particular, at the odd exotic
representative \(A=1\), the zero-top candidates satisfy
\[
 \boxed{
 \langle\epsilon_D,[W,R]_2\rangle=h_t,\quad
 \langle\epsilon_S,[W,R]_2\rangle=w_t,\quad
 \langle\epsilon_C,[W,R]_2\rangle=\zeta_{\rm nat}=h_t+w_t .
 }
 \tag{RT.56.6}\label{p1b:eq:rt-56-6}
\]
Here \(D_0,S_0\) are the retained natural determinant and
sphere candidates, and
\(C_0=2q_++\gamma(-2b)\) is the candidate of
Corollary~\ref{p1b:rt:corollary55-test}.
\end{theorem}
\begin{proof}
After a fixed seven-skeleton extension, the only possible
remaining obstruction to the stabilized gauge on the
eight-dimensional relative pair lies in \(\pi_7U(3)\).
This group equals \(\pi_7SU(3)=0\), so each cell fills.

On a characteristic disk, the determinant of the attaching
gauge has a logarithm because \(S^7\) is simply connected.
Extend that logarithm across the disk. Multiplying by its
inverse in one local line converts the boundary \(U(2)\)
matrix to an \(SU(2)\) matrix and the extending \(U(3)\)
matrix to an \(SU(3)\) matrix, with boundary \(i\varphi_e\).
The rank-two obstruction on this cell is
\([\varphi_e]\in\pi_7SU(2)\). By
Theorem~\ref{p1b:rt:theorem56-detector} it is exactly
\(\epsilon_e\mu\). This identifies the cellwise cochain with
the usual obstruction cochain. Coefficients are untwisted:
the action on \(C_2\) is necessarily trivial.
Changing the rank-three filling leaves each cell parity
fixed; allowed changes of the seven-skeleton filling give
the usual obstruction indeterminacy. For these odd targets
that essential obstruction has zero lower-choice correction,
as in the retained relative calculation.

The natural top connecting map is the isomorphism of
Theorem~\ref{p1b:rt:theorem38-sequence}; the generator
\(\Theta\) pairs to one on \([W,R]_2\).
The already proved profiles are
\(O_8(D_\tau)=(\tau+h_t)\Theta\),
\(O_8(S_\tau)=(\tau+w_t)\Theta\), and
\(O_8(C_\tau)=(\tau+\zeta_{\rm nat})\Theta\).
Taking \(\tau=0\) proves \eqref{p1b:eq:rt-56-6}.
\end{proof}

The same detector may be applied directly to the single supported
boundary comparison \eqref{p1b:eq:rt-55-8}, instead of summing relative cells.
Neither construction evaluates its square without the actual
comparison map and a geometric model of the projected filling.
In particular, a rank-three extension supplies no rank-two
extension test by itself. The final actual-map construction below supplies this detector input and evaluates \(h_t=w_t=1,\zeta_{\rm nat}=0\).



\subsection{A finite cochain procedure for the top comparison}
\label{p1b:sec:rtproof57}

Use Theorem~\ref{p1b:rt:theorem56-detector} on an actual
comparison map. Replace its projected mapping cone by a finite
ordered simplicial complex \(K\) of the same homotopy type,
with its bottom and top classes identified. Such a model can
be obtained by simplicial approximation of \(g_H\) followed
by triangulation of its mapping cylinder and cone.

\begin{theorem}[The finite cup-one test]
\label{p1b:rt:theorem57-cupone}
Let \(a\in C^5(K;C_2)\) be a cocycle representing \(u_5\).
Let \(c\in Z_9(K;C_2)\) represent the top class, normalized
by \(\langle v_9,c\rangle=1\).
On any ordered nine-simplex write
\[
 (a\smile_1a)(v_0,\ldots,v_9)=
 \sum_{j=0}^{4}
 a(v_0,\ldots,v_j,v_{j+5},\ldots,v_9)\,
 a(v_j,\ldots,v_{j+5}).
 \tag{RT.57.1}\label{p1b:eq:rt-57-1}
\]
All sums are in \(C_2\). The original supported comparison
has the exact value
\[
 \boxed{[\varphi]=\langle a\smile_1a,c\rangle\mu.}
 \tag{RT.57.2}\label{p1b:eq:rt-57-2}
\]
This parity is independent of the representative of \(u_5\),
the representative of the top cycle, and the vertex ordering.
\end{theorem}
\begin{proof}
Steenrod's cup-one formula is obtained by dividing the vertex
string into three intervals \([0,j]\), \([j,k]\), and
\([k,9]\), using the first and third intervals in the
first factor and the second interval in the second factor.
Each factor must be five-dimensional. Thus \(k=j+5\)
and \(0\le j\le4\), giving exactly \eqref{p1b:eq:rt-57-1}.
The general interval formula and the square convention appear
in Brumfiel--Medina-Mardones--Morgan,
\textit{A cochain level proof of Adem relations in the mod 2
Steenrod algebra}, arXiv:2006.09354, Sections 2.1--2.2.
For a degree-five cocycle,
\(\operatorname{Sq}^4[a]=[a\smile_{5-4}a]=[a\smile_1a]\).
Pairing \eqref{p1b:eq:rt-56-3} with \(c\) and using
Theorem~\ref{p1b:rt:theorem56-detector} proves \eqref{p1b:eq:rt-57-2}.
The remaining assertions follow because the square is a
well-defined cohomology operation on the realized complex.
\end{proof}

The executable procedure first generates all faces of \(K\).
It checks \(\delta a=0\) on every six-simplex,
\(\partial c=0\), and a supplied bottom cycle \(b_5\) with
\(\partial b_5=0\), \(\langle a,b_5\rangle=1\).
It then computes the five terms of \eqref{p1b:eq:rt-57-1} on the support
of \(c\). Once the face sets and supports are available,
there are exactly five membership products per top-chain
simplex. This statement does not bound the cost of constructing
the geometric cone or of obtaining a cocycle by linear algebra.

\begin{proposition}[A nonzero calibration using an actual triangulation]
\label{p1b:rt:proposition57-calibration}
Let \(K_0\) be Gaifullin's fifteen-vertex triangulation
\(\mathbb H P^2_{15}(C_6\times C_2)\).
In the frozen vertex labels \(0,\ldots,14\) its face vector is
\[
 (15,105,455,1365,3003,4515,4230,2205,490).
 \tag{RT.57.3}\label{p1b:eq:rt-57-3}
\]
The finite procedure on its suspension \(K=\Sigma K_0\)
returns one for the normalized bottom and top generators.
On \(S^5\vee S^9\), triangulated as two simplex boundaries
with one common vertex, it returns zero.
\end{proposition}
\begin{proof}
Gaifullin's author-supplied data give all 490 maximal simplices.
The identification with \(\mathbb H P^2\) is Proposition 9.4
of \textit{New examples and partial classification of
15-vertex triangulations of the quaternionic projective plane},
arXiv:2311.11309. The face vector is counted from those data.

Take the missing five-simplex with vertices
\((0,1,2,3,9,10)\) and its boundary \(c_4\).
Solve \(\delta a_4=0\), \(\langle a_4,c_4\rangle=1\)
over \(C_2\) with the deterministic elimination convention
in the accompanying script. The solution has support on
196 four-simplices; all cocycle and cycle equations are
checked directly. Counting the ordinary cup product over the
490 top simplices gives
\(\langle a_4\smile a_4,[K_0]_2\rangle=1\).

Give the suspension north vertex label \(-1\), south vertex
label 15, and retain the base vertex order. Define \(a_5\)
to have value \(a_4(f)\) on north cones on four-simplices
and zero on all other five-simplices.
It is a cocycle since \(\delta a_4=0\).
The bottom cycle is the sum of both cones on \(c_4\);
the top cycle is the sum of all 980 suspended facets.
Their boundaries vanish and the bottom pairing is one.
On a north nine-simplex only \(j=0\) contributes to
\eqref{p1b:eq:rt-57-1}, and it is precisely the old cup product term.
All south terms vanish. The total is consequently one.
This also follows from suspension stability.

In the wedge example the bottom cocycle can be supported
on one facet of the \(S^5\) boundary and zero on the \(S^9\)
boundary. Every cup-one term on the latter is zero.
The summed top boundary cycle therefore gives zero.
\end{proof}

The executable calibration evaluates both examples with the same
general face formula. It also transports cochains and cycles under
vertex relabelings, and changes \(a_5\) by coboundaries.
The calibration complexes are external examples.
They are not geometric models for the framework's natural
determinant, sphere, or combined comparison cones.
Before assigning a natural bit from an executable output,
identify the submitted complex and its classes with one of
those actual cones, or with the actual cellwise cones of
Theorem~\ref{p1b:rt:theorem56-relative}.



\subsection{Exotic coordinates and a small-normal boundary frame}
\label{p1b:sec:rtproof58}

The following construction makes the coherent exotic reference
explicit. It supplies the local geometric data for a stabilized
comparison; it does not supply the odd global clutching map.
All unit spheres carry the indicated Euclidean or Hermitian metrics.
Write a based unit quaternion as \(x=(x_0,\mathbf x)\), with
basepoint \(1=(1,0)\), and similarly \(p=(p_0,\mathbf p)\).
Set \(J=\operatorname{diag}(1,1,-1)\) on \(\mathbb C^3\).

\begin{proposition}[Coordinate realization of the exotic section]
\label{p1b:rt:proposition58-section}
Away from the based wedge put
\[
 a_x=1-x_0,\quad a_p=1-p_0,\quad
 K=2a_x+2a_p-a_xa_p,
\]
\[
 u_0=1-\frac{2a_xa_p}{K},\qquad
 \mathbf u_x=\frac{2a_p\mathbf x}{K},\qquad
 \mathbf u_p=\frac{2a_x\mathbf p}{K},\qquad
 \mathbf u_a=\frac{\mathbf u_x+\mathbf u_p}{\sqrt2},\quad
 \mathbf u_b=\frac{\mathbf u_x-\mathbf u_p}{\sqrt2}.
 \tag{RT.58.1}\label{p1b:eq:rt-58-1}
\]
On the wedge use \((u_0,\mathbf u_a,\mathbf u_b)=(1,0,0)\).
Define
\[
 v=(u_{a1},u_{a2},u_{b1}),\quad
 z_0=u_0+i u_{b2},\quad z_1=u_{a3}+i u_{b3},\quad
 r=(|z_0|^2+|z_1|^2)^{1/2}.
\]
For \(r>0\) put
\[
 (h_1,h_2,h_3)=
 \frac{(2\operatorname{Re}(z_0\overline z_1),
         2\operatorname{Im}(z_0\overline z_1),
         |z_0|^2-|z_1|^2)}{r};
\]
at \(r=0\) put \(h=0\). Then
\[
 s(x,p)=(h_3+i h_1,\ u_{a1}+i u_{a2},\ u_{b1}+i h_2)
       \in S(\mathbb C^3),\qquad
 s(p,x)=Js(x,p).
 \tag{RT.58.2}\label{p1b:eq:rt-58-2}
\]
This realizes the supplied equivariant construction
\(\eta_5^{C_2}q_\wedge\) after fixed orthogonal identifications.
Its perpendicular-plane bundle is the coherent exotic reference:
\[
 E_{\mathrm{ex},(x,p)}=s(x,p)^\perp,\qquad
 P_{\mathrm{ex}}=I-ss^*,\qquad
 E_{\mathrm{ex}}\oplus\mathbb C s=\mathbb C^3.
 \tag{RT.58.3}\label{p1b:eq:rt-58-3}
\]
On the diagonal \(s(y,y)=(a_1,a_2,0)\), and its fixed-sphere
map has degree \(+2\) with the orientations specified below.
\end{proposition}
\begin{proof}
Stereographic coordinates
\(X=\mathbf x/(1-x_0)\), \(P=\mathbf p/(1-p_0)\)
identify the complement of the based wedge with \(\mathbb R^6\).
Inverse stereographic projection is
\(((|X|^2+|P|^2-1)/(1+|X|^2+|P|^2),
2X/(1+|X|^2+|P|^2),2P/(1+|X|^2+|P|^2))\).
Using \(|X|^2=2/a_x-1\) gives \eqref{p1b:eq:rt-58-1}.
This is the one-point compactification of the smash product;
the wedge goes to its positive pole. Balanced sum and difference
coordinates make exchange act as
\((u_0,\mathbf u_a,\mathbf u_b)\mapsto
(u_0,\mathbf u_a,-\mathbf u_b)\).

The quadratic Hopf numerator has length \(r^2\), so
\(|h|=r\). Thus \(|v|^2+|h|^2=1\), including \(r=0\).
The estimate \(|h|=r\) proves continuity at the suspended join.
Exchange conjugates \(z_0,z_1\), fixes \(h_1,h_3\), and
negates \(h_2,u_{b1}\), proving \eqref{p1b:eq:rt-58-2}.
Before the three real suspension coordinates \(v\) are added,
this is the conjugation-equivariant Hopf map on \(\mathbb C^2\).
Those coordinates have representation \(2+\sigma\).
The domain and target therefore have the representations
\(3+3\sigma\) and \(3+2\sigma\) in the supplied construction.
Pulling back the unit-vector perpendicular-plane bundle gives
\eqref{p1b:eq:rt-58-3} and the coherent fibre lift \(w\mapsto Jw\).

On the fixed diagonal the stereographic coordinate becomes
\(\sqrt2X\), a degree \(+1\) map of its compactified \(\mathbb R^3\).
The remaining fixed-sphere map is
\[
 (u_0,\mathbf u_a)\longmapsto
 \left(\frac{(u_0+i u_{a3})^2}
                  {(u_0^2+u_{a3}^2)^{1/2}},
        u_{a1}+i u_{a2}\right).
\]
At zero denominator its first coordinate is zero.
This is the twice-suspended circle squaring map, with its
other two real coordinates fixed. Orient source and target
by these suspension identifications. Circle squaring has
degree \(+2\), proving the claimed degree.
\end{proof}

\begin{lemma}[A uniform small-normal chart]
\label{p1b:rt:lemma58-chart}
For \(y\in S^3\), \(n=(n_0,\mathbf n)\in S^3\), and
\(0<\epsilon\le1/12\), let
\[
 x=y\exp(\epsilon\mathbf n),\qquad
 p=y\exp(-\epsilon\mathbf n),\qquad
 a=s(y,y),\quad b=s(x,p).
\]
The scalar-circle normal coordinate is \(\epsilon n_0\);
the section is independent of that coordinate. Then
\[
 \|b-a\|\le3\sqrt2\epsilon|\mathbf n|
 \le\frac{\sqrt2}{4}<1,\qquad
 \operatorname{Re}(a^*b)>0.
 \tag{RT.58.4}\label{p1b:eq:rt-58-4}
\]
\end{lemma}
\begin{proof}
Let \(Q\) denote the smash map before balanced coordinates.
With \(D=1+|X|^2+|P|^2\), conformality of stereographic
projection gives
\[
 \|dQ\|^2=
 \frac{(1+|X|^2)^2\|dx\|^2+
       (1+|P|^2)^2\|dp\|^2}{D^2}
 \le\|dx\|^2+\|dp\|^2.
\]
The estimate extends to paths meeting the wedge by continuity
and integration on the complementary intervals.
For the Hopf numerator \(H\), quaternion multiplication gives
\(|H(z)-H(w)|\le(|z|+|w|)|z-w|\).
If \(|z|\ge|w|>0\), writing the difference after division
by their radii shows
\[
 \left|\frac{H(z)}{|z|}-\frac{H(w)}{|w|}\right|
 \le\left(1+\frac{|w|}{|z|}
              +\frac{|w|}{|z|}\right)|z-w|
 \le3|z-w|.
\]
The bound also holds at zero radius. The suspended map,
which retains \(v\), is consequently \(3\)-Lipschitz.
The product path from \((y,y)\) to \((x,p)\) has length
\(\sqrt2\epsilon|\mathbf n|\). These bounds prove \eqref{p1b:eq:rt-58-4}.
Finally \(\operatorname{Re}(a^*b)=1-\|a-b\|^2/2>0\).
\end{proof}

\begin{theorem}[An explicit equivariant \(SU(3)\) boundary frame]
\label{p1b:rt:theorem58-frame}
For unit \(a,b\in\mathbb C^3\) with \(a^*b\ne-1\), set
\[
 c=a^*b,\quad w=b-ca,\quad
 \alpha=\frac{1+c}{1+\overline c},\quad
 d=\frac{1+\overline c}{|1+c|},
\]
\[
 U=I+(c-1)aa^*+wa^*-\alpha aw^*
                   -\frac{ww^*}{1+\overline c},\qquad
 T(a,b)=(bb^*+d(I-bb^*))U.
 \tag{RT.58.5}\label{p1b:eq:rt-58-5}
\]
Then \(T(a,b)\in SU(3)\), \(T(a,b)a=b\), and
\(T(a,a)=I\). It is continuous on this chart and
\(T(Ra,Rb)=RT(a,b)R^*\) for every unitary \(R\).
For \(a,b\) of Lemma~\ref{p1b:rt:lemma58-chart} put
\[
 f_+=(-\overline a_2,\overline a_1,0),\qquad f_-=e_3,\qquad
 F_\epsilon(y,n)=[T(a,b)f_+,\,T(a,b)f_-,\,b].
 \tag{RT.58.6}\label{p1b:eq:rt-58-6}
\]
This is an \(SU(3)\) frame of the actual coherent reference
stabilization on the small normal boundary, satisfying
\[
 F_\epsilon(y,-n)=J F_\epsilon(y,n)\,
                         \operatorname{diag}(1,-1,1).
 \tag{RT.58.7}\label{p1b:eq:rt-58-7}
\]
Its first two columns identify the boundary reference bundle
with \(1\oplus\lambda\), and its third column is the specified
trivial stabilizing line.
\end{theorem}
\begin{proof}
If \(w\ne0\), use the orthonormal pair \(a,w/\|w\|\).
The matrix of \(U\) on this plane is
\[
 \begin{pmatrix}c&-\alpha\|w\|\\
                 \|w\|&\alpha\overline c\end{pmatrix};
\]
on its orthogonal complement it is the identity.
Since \(|c|^2+\|w\|^2=1\) and \(|\alpha|=1\),
this is unitary, sends \(a\) to \(b\), and has determinant
\(\alpha\). If \(w=0\), the same identities follow directly
from \eqref{p1b:eq:rt-58-5}, without dividing by \(\|w\|\).
The left factor fixes \(b\), is unitary, and has determinant
\(d^2=\alpha^{-1}\). This proves the assertions about \(T\);
continuity and covariance follow from its formula.

The diagonal section is fixed by \(J\).
The frame \(F_0=[f_+,e_3,a]\) is unitary with determinant one,
including at the based pole. Lemma~\ref{p1b:rt:lemma58-chart}
places every indicated \(a,b\) inside the chart of \eqref{p1b:eq:rt-58-5}.
Hence \eqref{p1b:eq:rt-58-6} is \(T(a,b)F_0\in SU(3)\).
Normal reversal exchanges \(x,p\), so it sends \(b\) to \(Jb\).
Covariance, \(Ja=a\), and
\(JF_0=F_0\operatorname{diag}(1,-1,1)\) give \eqref{p1b:eq:rt-58-7}.
Descending the three columns gives respectively a positive
line, the negative line \(\lambda\), and a positive line.
\end{proof}

\begin{proposition}[Global Chern scope of this reference]
\label{p1b:rt:proposition58-scope}
Extend \eqref{p1b:eq:rt-58-2}--\eqref{p1b:eq:rt-58-3} as a coherent exchange model
over the full free configuration space \(W\), keeping the
constant ambient lift \(J\). This reference has \(c_2=0\).
Its coordinate construction alone does not identify the odd
compatible bundle or its natural determinant and sphere gauges.
\end{proposition}
\begin{proof}
The ambient exchange bundle decomposes as
\(\underline{\mathbb C}^{\,2}\oplus\lambda_W\).
The equivariant section \(s\) gives a nowhere-zero section of
the line \(\mathbb C s\) in \eqref{p1b:eq:rt-58-3}; that line is trivial.
Thus the Whitney formula gives
\(c(E_{\mathrm{ex}})=c(\underline{\mathbb C}^{\,2}\oplus\lambda_W)\),
whose second Chern class is zero.
In the retained pushout \(W=Y\cup_J Z_s\), the odd case instead
requires its actual joining clutch with odd degree \(d\);
the connecting homomorphism sends the joining generator to
\(\tau^2\). This information cannot be discarded by replacing
that clutch with the coherent full reference.

The formula is still usable on the exotic \(Y\) piece, whose
second Chern class is zero, and on any actual boundary bundle
after specifying a bundle isomorphism over the fixed boundary
base. The ordinary frame invariance proved in Subsection~\ref{p1b:sec:rtproof54}
preserves the natural top bits under such an isomorphism.
The subsequent actual odd clutching, closing paths and whole paired comparison retain this ordinary-frame invariance and evaluate all natural top coefficients.
\end{proof}



\subsection{From a specified comparison map to its finite square}
\label{p1b:sec:rtproof59}

\begin{proposition}[The last column suffices]
\label{p1b:rt:proposition59-lastcolumn}
On an actual comparison eight-cell choose a unitary frame
\(Q_e\) of the actual rank-three stabilization whose last
column is its trivial stabilizing section \(\sigma\).
Let \(G_e:D^8\to U(3)\) represent an actual stabilized gauge
filling of the prescribed partial \(U(2)\) gauge.
Then
\[
 u_e=Q_e^*G_eQ_e e_3:D^8\longrightarrow S^5
 \tag{RT.59.1}\label{p1b:eq:rt-59-1}
\]
is \(e_3\) on the boundary. Its descended map \(S^8\to S^5\)
has square parity equal to the original rank-two top
obstruction. Determinant normalization of \(G_e\) is optional
for this calculation.
\end{proposition}
\begin{proof}
The boundary gauge preserves \(\sigma\) pointwise, so its
matrix is \(\operatorname{diag}(\varphi,1)\).
Use the full last-column fibration
\(U(2)\to U(3)\to S^5\).
In degrees at least two the inclusions \(SU(n)\to U(n)\)
induce isomorphisms on homotopy groups. Its exact segment is
therefore the segment used in
Lemma~\ref{p1b:rt:lemma56-parity}. The connecting homomorphism
sends the projected filling to \([\varphi]\).
Theorem~\ref{p1b:rt:theorem56-detector} gives the square parity.
Alternatively multiply \(Q_e^*G_eQ_e\) on the right by
\(\operatorname{diag}(\det(Q_e^*G_eQ_e)^{-1},1,1)\).
This makes its determinant one and leaves its last column
unchanged. Its normalized boundary represents the same
class in \(\pi_7U(2)\).
\end{proof}

The coordinates of Theorem~\ref{p1b:rt:theorem58-frame}
give explicit adapted boundary frames for the coherent reference.
They do not supply the \(G_e\) in \eqref{p1b:eq:rt-59-1} for the odd natural
comparison. The following finite procedure begins only after
the projected map itself is specified.

\begin{theorem}[A finite cone and its normalized cochains]
\label{p1b:rt:theorem59-cone}
Let \(A,T\) be finite ordered simplicial complexes, with
\(A\simeq S^8\), \(T\simeq S^5\), and
\(f:A\to T\) a specified simplicial map.
Supply a target five-cocycle \(a\), a target five-cycle
\(b\) with \(\langle a,b\rangle=1\), a domain fundamental
eight-cycle \(c\), and a target nine-chain \(q\) satisfying
\(\partial q=f_*c\). All chains and cochains in this theorem
have coefficients in \(C_2\). Degenerate image simplices
contribute zero. Choose a four-cochain \(\beta\) on \(A\) with
\[
 \delta\beta=f^*a.
 \tag{RT.59.2}\label{p1b:eq:rt-59-2}
\]
There is an explicit finite complex \(X\simeq C_f\), a
five-cocycle \(\widetilde a\) restricting to \(a\) on \(T\),
and a normalized top nine-cycle \(C\). Their exact pairing
\(\langle\widetilde a\smile_1\widetilde a,C\rangle\)
is the parity detected by the cone of \(f\).
\end{theorem}
\begin{proof}
Take disjoint copies of the vertices of \(A\) and \(T\).
Let \(M\) have the target simplices and the generating
simplices \(\sigma_A\cup f(\sigma)_T\) for maximal
\(\sigma\) in \(A\), with all their faces.
The domain copy is a subcomplex. Define the retraction
\(r:M\to T\) by \(r(v_A)=f(v)_T\), with \(r\) the identity
on target vertices. It is simplicial. For each simplex of
\(M\), its vertices together with their retracted images lie
in one of its generating simplices. Hence the identity and
the inclusion followed by \(r\) are contiguous; their linear
homotopy fixes \(T\). Thus \(M\) deformation retracts onto \(T\),
and the composite \(A\to M\to T\) is \(f\).

Add an apex \(*\) and attach the genuine cone on the embedded
domain copy:
\[
 X=M\cup_A(*A).
 \tag{RT.59.3}\label{p1b:eq:rt-59-3}
\]
Both this construction and the usual mapping cone are
homotopy pushouts of \(M\leftarrow A\to *\) and
\(T\leftarrow A\to *\), respectively. The retraction is a
homotopy equivalence commuting with these maps, so
\(X\simeq C_f\). Equivalently use the map induced by \(r\)
and the identity on the attached domain cone.
This construction preserves the domain when \(f\) identifies
vertices; simply coning its image would not do so.

Since \(H^5(A;C_2)=0\), \eqref{p1b:eq:rt-59-2} has a solution.
Put \(\widetilde a=r^*a\) on \(M\), and put
\(\widetilde a(*\sigma)=\beta(\sigma)\) on cone
five-simplices. The cocycle equation on each cone
six-simplex is \(f^*a+\delta\beta=0\).
On the other six-simplices it follows from \(\delta a=0\).
The target cycle \(b\) still pairs to one.

For an ordered eight-simplex
\(\sigma=(v_0,\ldots,v_8)\), put
\[
 P\sigma=\sum_{j=0}^{8}
 {}[v_{0,A},\ldots,v_{j,A},
   f(v_j)_T,\ldots,f(v_8)_T].
 \tag{RT.59.4}\label{p1b:eq:rt-59-4}
\]
Terms with repeated vertices are zero, and surviving terms
are reordered over \(C_2\). Every term belongs to \(M\).
The prism identity is
\(\partial P+P\partial=i_A+i_Tf_*\):
in the boundary expansion adjacent subdivisions cancel,
leaving the two endpoints and the prisms of the boundary.
Consequently
\[
 C=*c+P(c)+i_Tq,\qquad \partial C=0.
 \tag{RT.59.5}\label{p1b:eq:rt-59-5}
\]
Under \(X\to X/M\cong\Sigma A\), this cycle maps to the
suspension of \(c\). The exact sequence and \(T\simeq S^5\)
identify it with the normalized top class of \(C_f\).
The five-cocycle is its normalized bottom class.
Theorem~\ref{p1b:rt:theorem57-cupone} now gives the claimed
pairing and its invariance under cochain and cycle choices.

If \(f\) is injective on vertices, it embeds \(A\) as a
subcomplex of \(T\). One may directly use
\(X=T\cup_{f(A)}*f(A)\), with
\(C=*f_*c+q\); the cone cochain is transported from \(\beta\).
This removes the thick cylinder without changing the proof.
\end{proof}

For context, Hatcher's \textit{Algebraic Topology}, Section 2C,
discusses simplicial mapping cylinders and cones; its chain
homotopy discussion in Section 2.1 gives the prism identity.
The particular thick complex, cochain extension, and explicit
normalized cycle above provide the data used by the executable.
Its dimension can exceed nine. This causes no difficulty:
the square is still evaluated on nine-simplices, and the
program checks its cocycle equation on ten-simplices when
they exist.

\begin{proposition}[A specified positive attaching map]
\label{p1b:rt:proposition59-positive}
Let \(K=\Sigma\mathbb H P^2_{15}(C_6\times C_2)\), with the
fixed triangulation of Subsection~\ref{p1b:sec:rtproof57} and cocycle \(a\).
Delete the maximal nine-simplex
\[
 \Delta=(-1,3,4,5,6,7,9,11,12,14)
 \tag{RT.59.6}\label{p1b:eq:rt-59-6}
\]
while retaining its boundary; call the resulting complex \(T\).
Let \(A=\partial\Delta\), and let \(f:A\hookrightarrow T\)
be the specified vertex inclusion. Then \(A\cong S^8\),
\(T\simeq S^5\), and the cone procedure returns one.
It returns zero for both the constant map
\(\partial\Delta^9\to\partial\Delta^6\) and the vertex map
\(j\mapsto j\bmod6\), \(0\le j\le9\), into a single
five-simplex of \(\partial\Delta^6\).
\end{proposition}
\begin{proof}
The frozen triangulation is PL homeomorphic to the indicated
projective plane before suspension, by Gaifullin's
Proposition 9.4, as used in
Proposition~\ref{p1b:rt:proposition57-calibration}.
Thus \(K\) is four-connected and its integral reduced
homology is \(\mathbb Z\) in degrees five and nine.
Its oriented fundamental suspended chain has coefficient
\(\pm1\) on \(\Delta\); therefore the map
\(H_9(K;\mathbb Z)\to H_9(K,T;\mathbb Z)\) is an isomorphism.
The pair is a single nine-cell attachment. The exact sequence
gives reduced homology \(\mathbb Z\) only in degree five for
\(T\), and cellular approximation preserves its four-connectivity.
By Hurewicz there is a map \(S^5\to T\) generating \(H_5\).
It is an integral homology isomorphism between simply
connected CW complexes, hence a homotopy equivalence
by relative Hurewicz and Whitehead (Hatcher, Corollary 4.33).

The target nine-chain \(q\) is the sum of the remaining
979 old top simplices. Its boundary is \(f_*c\), where
\(c\) is the ten-face boundary of \(\Delta\).
The old bottom cycle lies in \(T\) and retains pairing one.
The direct cone adds ten nine-simplices, replacing \(\Delta\)
by the cone on its boundary. This is a stellar subdivision
of that simplex, so the resulting complex is homeomorphic
to \(K\). The cone cocycle and top cycle therefore have
square pairing one, agreeing with the exact finite calculation.
The four-cochain solve has rank 126 and gives 18 supported
four-simplices; the resulting five-cocycle has 214 terms
and the top cycle has 989 terms.

The two other maps have image in a contractible simplex.
They are nullhomotopic, so their cones have the homotopy
type \(S^5\vee S^9\). The normalized square pairs to zero.
The program independently constructs and evaluates these
cones, including a dimension-fourteen thick cylinder in
the second case. These are specified external examples;
none is asserted to be a natural framework comparison.
\end{proof}

To apply Theorem~\ref{p1b:rt:theorem59-cone} to the framework,
identify \(f\) with the descended last-column map of
\eqref{p1b:eq:rt-59-1}, and identify its sphere generators with the
actual relative eight-cell generators. The missing data
are those actual odd-clutch fillings or equivalent
specified comparison maps, rather than further abstract
knowledge of \(\pi_8S^5\) or the cup-one formula.



\subsection{An explicit joining frame retaining the odd clutch}
\label{p1b:sec:rtproof60}

Keep the actual pushout \(Y\cup_JZ_s\) of
Theorem~\ref{p1b:rt:theorem46-pushout}, the exotic \(Y\)
reference of Subsection~\ref{p1b:sec:rtproof58}, and the split sphere-piece reference
\(E_Z^0=1\oplus\lambda\). Write the ambient exchange matrix
as \(\mathsf J=\operatorname{diag}(1,1,-1)\), to distinguish
it from the joining manifold \(J\). Its product chart is
\[
 {}[a,w]\longmapsto(\xi,u)=(a^2,a^{-1}w)\in S^1\times S^3.
 \tag{RT.60.1}\label{p1b:eq:rt-60-1}
\]
Embed the circle in the fixed quaternionic complex subfield.
The section on the antipodal sphere pairs is
\(b(w)=s(w,-w)\), with the explicit \(s\) of Subsection~\ref{p1b:sec:rtproof58}.

\begin{lemma}[A global antipodal transport chart]
\label{p1b:rt:lemma60-chart}
Writing \(w=(w_0,\mathbf w)\), the smash coordinates on
the joining locus simplify to
\[
 u_0=\frac{1+3w_0^2}{3+w_0^2},\quad
 \mathbf u_a=\frac{2\sqrt2\,w_0\mathbf w}{3+w_0^2},\quad
 \mathbf u_b=\frac{2\sqrt2\,\mathbf w}{3+w_0^2}.
 \tag{RT.60.2}\label{p1b:eq:rt-60-2}
\]
The section satisfies \(b(-w)=\mathsf Jb(w)\) and
\[
 \operatorname{Re}(e_1^*b(w))\ge-7/9,\qquad
 |1+e_1^*b(w)|\ge2/9.
 \tag{RT.60.3}\label{p1b:eq:rt-60-3}
\]
Thus the \(SU(3)\) transport \(T(e_1,b(w))\) of
Theorem~\ref{p1b:rt:theorem58-frame} is defined on the
entire joining sphere, with a uniform denominator bound.
\end{lemma}
\begin{proof}
Substitute \(x=w,p=-w\) in \eqref{p1b:eq:rt-58-1}; its denominator is
\(3+w_0^2\). This gives \eqref{p1b:eq:rt-60-2} and in particular
\(u_0\ge1/3\). Equivariance follows from \eqref{p1b:eq:rt-58-2}.
In the suspended Hopf formula,
\(\operatorname{Re}(e_1^*b)=h_3
 =(2|z_0|^2-r^2)/r\), where
\(|z_0|\ge u_0\ge1/3\) and \(r\le1\).
Therefore \(h_3\ge2/9-1=-7/9\).
The real part of \(1+e_1^*b\) is at least \(2/9\);
this proves \eqref{p1b:eq:rt-60-3}.
\end{proof}

\begin{theorem}[Concrete ordinary joining frames]
\label{p1b:rt:theorem60-frames}
Let \(T_w=T(e_1,b(w))\), and use the full frames
\[
 Q_Z(a)=[e_2,a e_3,e_1],\qquad
 Q_Y(a,w)=T_wQ_Z(a).
 \tag{RT.60.4}\label{p1b:eq:rt-60-4}
\]
Their first two columns are ordinary frames of the split
and exotic joining bundles, respectively. Their last
columns are the specified trivial stabilizing sections
\(e_1\) and \(b(w)\). Both determinants are \(a\), and
\[
 Q_Z(-a)=\mathsf JQ_Z(a),\qquad
 Q_Y(-a,-w)=\mathsf JQ_Y(a,w).
 \tag{RT.60.5}\label{p1b:eq:rt-60-5}
\]
Thus both first-two-column frames descend to the quotient
joining manifold, with no choice of a square root of \(\xi\).
The coherent zero-second-Chern reference gluing is the
identity in these ordinary frames.
\end{theorem}
\begin{proof}
The first two columns of \(Q_Z\) span the complement of
\(e_1\); they are the positive and negative exchange lines.
Multiplying the negative column by \(a\), which changes
sign with the ordering deck transformation, makes it an
ordinary descended joining section. The determinant of
\([e_2,a e_3,e_1]\) is \(a\).
Transport is unitary of determinant one and sends \(e_1\)
to \(b\), so \(Q_Y\) has the asserted column and determinant
properties.

Covariance of \(T\) and \(\mathsf Je_1=e_1\) give
\(T_{-w}=\mathsf JT_w\mathsf J\).
Using the first equality of \eqref{p1b:eq:rt-60-5} proves the second.
This is the descent condition for both full frames.

On all of \(Z_s\), the exotic coherent reference section
on sphere-antipodal pairs depends only on \(w\).
The same transport \(T_w\) therefore identifies that
reference with the split sphere-piece bundle, equivariantly.
At the joining endpoint it is precisely \eqref{p1b:eq:rt-60-4}.
Hence identity matching in these frames is the gluing
induced from the full coherent reference. Subsection~\ref{p1b:sec:rtproof58} proves
that this full reference has \(c_2=0\).
\end{proof}

\begin{theorem}[An explicit odd exotic representative]
\label{p1b:rt:theorem60-clutch}
Represent a unit quaternion by its standard \(SU(2)\)
matrix, and put \(f_d(\xi,u)=u^d\) for \(d\in\mathbb Z\).
In ambient stabilized coordinates define
\[
 \widetilde c_d[a,w]
   =Q_Y(a,w)\operatorname{diag}(f_d(a^{-1}w),1)Q_Z(a)^*.
 \tag{RT.60.6}\label{p1b:eq:rt-60-6}
\]
This is a continuous descended joining isomorphism,
maps \(e_1\) to \(b(w)\), has determinant one, and restricts
to a rank-two clutch of component \((0,d,0)\) in the
ordinary frames \eqref{p1b:eq:rt-60-4}. Gluing the exotic \(Y\) piece
to \(E_Z^0\) using this clutch gives
\[
 c_1=\tau,\qquad c_2=(d\bmod2)\tau^2.
 \tag{RT.60.7}\label{p1b:eq:rt-60-7}
\]
In particular \(d=1\) gives an actual ordinary odd exotic
representative, with its joining frame and matrix specified.
\end{theorem}
\begin{proof}
The function \(a^{-1}w\) is unchanged by
\((a,w)\mapsto(-a,-w)\), so it is the actual quaternion
coordinate of the quotient chart.
The ordinary middle matrix acts on the first two columns
and fixes the stabilizing column. Its determinant is one.
The factors \(a\) and \(\overline a\) from the two full
frames cancel. Equation \eqref{p1b:eq:rt-60-5} gives the ambient descent
law
\(\widetilde c_d[-a,-w]
 =\mathsf J\widetilde c_d[a,w]\mathsf J\).
It consequently glues the actual bundle pieces and their
stabilizations, rather than replacing the odd bundle by
the ambient coherent reference.

For \(d=0\), the matrix is \(T_w\), and the gluing is the
zero-Chern reference of
Theorem~\ref{p1b:rt:theorem60-frames}.
Its determinant line is the ordering line.
The difference for arbitrary \(d\) is the sphere clutch
of degree \(d\) on \(J=S^1\times S^3\).
By Theorem~\ref{p1b:rt:theorem47-clutch} and
\(\delta_{\mathrm{MV}}u_J=\tau^2\), this changes the second
Chern class by \(d\tau^2\), proving \eqref{p1b:eq:rt-60-7}.
\end{proof}

The second \(Z_s\) class is obtained by the retained
\(\eta_3\) patch on a five-ball off the joining locus.
It has the same displayed joining data and the same Chern
parity, by Subsection~\ref{p1b:sec:rtproof50}. The ordinary classification of Subsection~\ref{p1b:sec:rtproof53}
identifies these piece-and-parity models with the corresponding
actual compatible ordinary bundle classes.

The odd clutch is now an explicit matrix. It cannot be
discarded after adding the trivial line: its second Chern
class remains \(\tau^2\). This construction specifies a
joining frame, retained by the following complete natural
target-gauge and boundary constructions. Any actual ordinary
boundary bundle identification over the fixed \(R\) can be
used for the eventual top calculation, by
Corollary~\ref{p1b:rt:corollary54-top}; the complete comparison carries
that identification throughout.



\subsection{A joining homotopy with an explicit top-generator filling}
\label{p1b:sec:rtproof61}

Work in the actual ordinary joining frames of Subsection~\ref{p1b:sec:rtproof60}.
Use \(\lambda=-\xi\), so the sphere-piece determinant
gauge is \(q(\lambda)=\operatorname{diag}(\lambda,1)\).
Take the basepoint \(\lambda=1,u=1\).
For \(k,d\in\mathbb Z\), let
\[
 q_k=\operatorname{diag}(\lambda^k,1,1),\qquad
 c_d=\operatorname{diag}(u^d,1)\in SU(3).
\]
These are stabilized matrices in the ordered column basis
of the joining frames, not a new ambient trivialization of
the global odd bundle.

\begin{theorem}[A based determinant-preserving joining homotopy]
\label{p1b:rt:theorem61-homotopy}
Write the first column of \(u^d\in SU(2)\) as
\((z_0,z_1)\). For \(s\in[0,1]\) set
\[
 \theta=\frac{\pi}{2}(1-s),\qquad
 h_s(u)=
 (\sin^2\theta+\cos^2\theta\,z_0,\,
  \cos\theta\,z_1,\,
  \sin\theta\cos\theta(1-z_0)).
 \tag{RT.61.1}\label{p1b:eq:rt-61-1}
\]
It is a unit vector with \(h_0=e_1\),
\(h_1=(z_0,z_1,0)\), and \(h_s(1)=e_1\).
The family
\[
 L_s=I+(\lambda^k-1)h_s h_s^*,\qquad K_s=L_s q_k^{-1}
 \tag{RT.61.2}\label{p1b:eq:rt-61-2}
\]
is a concrete \(U(3)\) joining homotopy
\[
 L_0=q_k,\quad L_1=c_dq_kc_d^{-1},\quad
 \det L_s=\lambda^k.
 \tag{RT.61.3}\label{p1b:eq:rt-61-3}
\]
Moreover \(K_s\in SU(3)\), \(K_0=I\), and
\[
 K_1=c_dq_kc_d^{-1}q_k^{-1}.
 \tag{RT.61.4}\label{p1b:eq:rt-61-4}
\]
The whole \(K_s\) is identity when \(\lambda=1\) or \(u=1\).
It therefore nullhomotopes the actual supported rank-two
joining commutator after stabilization, relative to its
based wedge.
\end{theorem}
\begin{proof}
Start with
\(\cos\theta(z_0,z_1,0)+\sin\theta e_3\), whose two terms
are orthogonal unit vectors. Rotate by \(-\theta\) in the
real \(e_1,e_3\) plane, leaving \(e_2\) fixed.
This gives exactly \eqref{p1b:eq:rt-61-1}, and proves unit length.
The same rotation fixes its corrected basepoint \(e_1\)
when \(u=1\). The endpoint assertions follow directly.

For any unit \(h\), \(I+(\rho-1)hh^*\) acts by \(\rho\)
on \(\mathbb Ch\) and by identity on its orthogonal
complement. It is unitary with determinant \(\rho\).
At the final endpoint its line is \(c_d\mathbb Ce_1\);
this proves \eqref{p1b:eq:rt-61-3}. Multiplying by \(q_k^{-1}\)
gives determinant one and proves \eqref{p1b:eq:rt-61-4}.
The based-wedge assertion follows from \(h_s(1)=e_1\)
and the factor \(\lambda^k-1\).
\end{proof}

In the retained rank-two joining law, \eqref{p1b:eq:rt-61-4} has
secondary coefficient \(kd\bmod2\):
conjugation by \((0,d,0)\) sends \((k,0,0)\) to
\((k,0,kd)\). Thus for \(k=d=1\) this endpoint is an
actual representative of \(\eta_3:S^4\to SU(2)\).
It has no rank-two nullhomotopy, despite the explicit
rank-three homotopy above.

\begin{theorem}[The projected joining map has unit absolute degree]
\label{p1b:rt:theorem61-degree}
The last column of \(K_s\) is
\[
 g_{k,d}(s,\lambda,u)=
 e_3+(\lambda^k-1)\overline{h_{s3}}\,h_s.
 \tag{RT.61.5}\label{p1b:eq:rt-61-5}
\]
It is \(e_3\) at \(s=0,1\) and on the based wedge,
so it descends to
\[
 g_{k,d}:\Sigma(S^1\wedge S^3)\cong S^5\longrightarrow S^5.
\]
The map \(g_{1,1}\) has degree of absolute value one,
and
\[
 \deg g_{k,d}=kd\,\deg g_{1,1}.
 \tag{RT.61.6}\label{p1b:eq:rt-61-6}
\]
All signs below use the same fixed sphere identifications.
\end{theorem}
\begin{proof}
Since \(q_k^{-1}e_3=e_3\), \eqref{p1b:eq:rt-61-5} follows from \eqref{p1b:eq:rt-61-2}.
The third coordinate of \(h_s\) vanishes at both endpoints
and at \(u=1\); the \(\lambda=1\) assertion is immediate.

For \(k=d=1\), consider the target \(-e_3\).
Its third-coordinate equation forces
\(\lambda=-1\) and \(|h_{s3}|=1\):
the real part of \((\lambda-1)|h_{s3}|^2\) is \(-2\)
only in this case. Unit length then forces
the first two coordinates of \(h_s\) to vanish.
Their formula gives uniquely
\[
 s=1/2,\qquad \lambda=-1,\qquad u=-1.
 \tag{RT.61.7}\label{p1b:eq:rt-61-7}
\]
Use local variables
\((\delta s,\phi,x_1,x_2,x_3)\), with
\(\lambda=-e^{i\phi}\) and
\(u=(-\sqrt{1-|x|^2},x_1,x_2,x_3)\).
Use the Hamilton \(SU(2)\) convention with first column
\((u_0+i u_1,-u_2+i u_3)\).
At this point the differential of the five target
coordinates
\((\operatorname{Re}g_1,\operatorname{Im}g_1,
  \operatorname{Re}g_2,\operatorname{Im}g_2,
  \operatorname{Im}g_3)\) is
\[
 \begin{pmatrix}
 2\pi&0&0&0&0\\
 0&0&-1&0&0\\
 0&0&0&\sqrt2&0\\
 0&0&0&0&-\sqrt2\\
 0&-1&0&0&0
 \end{pmatrix}.
 \tag{RT.61.8}\label{p1b:eq:rt-61-8}
\]
Its determinant is \(4\pi\ne0\).
There is exactly one regular preimage, so the degree has
absolute value one. This is a local degree proof, not
an estimate from sampled Jacobians.

The general map is \(g_{1,1}\) precomposed with the suspension
of \((\lambda,u)\mapsto(\lambda^k,u^d)\) on the smash.
Circle powering has degree \(k\).
Quaternion powering has degree \(d\), including negative
integers: the integral degree-three class of \(S^3\) is
primitive under multiplication, so its pullback under a
product of \(d\) copies is \(d\) times that class, and
inversion gives minus one.
Degrees multiply on the smash and are preserved by
suspension. This proves \eqref{p1b:eq:rt-61-6}.
\end{proof}

\begin{theorem}[An explicit filling of the top generator]
\label{p1b:rt:theorem61-top}
Let \(\nu_4:S^7\to S^4\) be the quaternionic Hopf map,
and identify \(S^4\) with \(S^1\wedge S^3\) as specified
below. Composing \eqref{p1b:eq:rt-61-4} with \(\nu_4\) gives the
supported comparison
\[
 \varphi_{k,d}:S^7\longrightarrow SU(2).
 \tag{RT.61.9}\label{p1b:eq:rt-61-9}
\]
The explicit radial filling
\[
 G_{k,d}(r,x)=K_r(\nu_4(x)),\qquad 0\le r\le1,
 \tag{RT.61.10}\label{p1b:eq:rt-61-10}
\]
is a continuous map \(D^8\to SU(3)\) whose boundary is
the stabilization of \(\varphi_{k,d}\).
Its last column descends to
\[
 g_{k,d}\circ\Sigma\nu_4:S^8\longrightarrow S^5.
 \tag{RT.61.11}\label{p1b:eq:rt-61-11}
\]
In the retained top group and square detector,
\[
 \boxed{[\varphi_{k,d}]=(kd\bmod2)\mu,\qquad
   \langle\operatorname{Sq}^4a_5,C_9\rangle=kd\bmod2.}
 \tag{RT.61.12}\label{p1b:eq:rt-61-12}
\]
In particular the actual displayed map with \(k=d=1\)
has top coefficient one.
\end{theorem}
\begin{proof}
The radial formula is continuous at the cone point because
\(K_0=I\) for every argument. At radius one, the endpoint
is block \(SU(2)\) and fixes \(e_3\). Its last column is
also constant at the cone point and on the collapsed
based line, so the quotient is \eqref{p1b:eq:rt-61-11}.
By Theorem~\ref{p1b:rt:theorem61-degree}, the first factor
has degree \(\pm kd\).
Since \(\pi_8S^5\) is in the stable range, its action by a
degree map is multiplication by that degree.
Thus \eqref{p1b:eq:rt-61-11} has class \(\pm kd\,\nu_5\) in \(C_{24}\),
where \(\nu_5=\Sigma\nu_4\).
Lemma~\ref{p1b:rt:lemma56-parity} sends its class to
\((kd\bmod2)\mu\), and
Theorem~\ref{p1b:rt:theorem56-detector} gives its square value.
This proves \eqref{p1b:eq:rt-61-12} without assigning a sign to the
retained \(C_{24}\) generator.
\end{proof}

For executable coordinates, write a quaternionic unit pair
\((q_0,q_1)\in S^7\) and put
\[
 \nu_4(q_0,q_1)
   =(|q_0|^2-|q_1|^2,\,2q_0\overline q_1).
 \tag{RT.61.13}\label{p1b:eq:rt-61-13}
\]
Away from its based pole, stereographic coordinates give
\[
 X=\frac{2q_0\overline q_1}{2|q_1|^2}=(t,\mathbf y)\in\mathbb R^4,
 \quad
 \lambda=\frac{t^2-1+2it}{1+t^2},\quad
 u=\frac{(|\mathbf y|^2-1,\,2\mathbf y)}{1+|\mathbf y|^2}.
 \tag{RT.61.14}\label{p1b:eq:rt-61-14}
\]
At the based pole use the smash basepoint.
The compactification
\((\mathbb R\times\mathbb R^3)^+
 \cong(\mathbb R^+)\wedge((\mathbb R^3)^+)\)
proves continuity of this identification. Because \(K_s\)
is identity on the based wedge for every \(s\), these
coordinates also give continuous formulas for the filling.

For \(k=d=1\), the endpoint identity is the classical
\(\eta_3\nu_4=\nu'\eta_6=\mu\), now with a concrete
stabilized filling. The same relation appears in Inoue,
\textit{Groups of self-equivalences of suspended real
projective spaces}, Kyushu J. Math. 59 (2005), p. 340.
His \(\nu'\) denotes a two-primary generator; the retained
full \(C_{12}\) generator has the same nonzero composition
with \(\eta_6\). The argument above uses the retained
fibration and the proved degree, and does not rely on
silently identifying these generator conventions.

The ordinary joining matrix of the odd bundle and a
determinant-fixed stable joining path are now explicit.
The Hopf-composed comparison \eqref{p1b:eq:rt-61-9} is also explicit and
evaluated. It is a specified primitive comparison. The complete natural determinant, sphere and combined boundary evaluations below retain the actual \(Y\) fibre representative, closing family, joining toggle and restriction along the actual boundary map of Subsection~\ref{p1b:sec:rtproof47}; those whole-map comparisons identify the detector input before assigning its coefficient.



\subsection{Matrix words for the exotic initial fibre gauges}
\label{p1b:sec:rtproof62}

Keep the actual exotic reference \(E=s^\perp\) over
\(P=S_x^3\times S_p^3\) from
Proposition~\ref{p1b:rt:proposition58-section}.
Its exchange matrix is \(\mathsf J=\operatorname{diag}(1,1,-1)\).
All matrices below act on the actual stabilization
\(E\oplus\mathbb C s=\mathbb C^3\), fix \(s\), and restrict
to determinant-one gauges of \(E\). Write
\[
 \sigma G(x,p)=\mathsf JG(p,x)\mathsf J,
 \qquad N(x,p)=2s(x,p)s(x,p)^*-I.
 \tag{RT.62.1}\label{p1b:eq:rt-62-1}
\]
Thus \(N=-I_E\oplus1\), \(N^2=I\), and \(\sigma N=N\).
In the retained exotic gauge chart, \([N]=(0,0;6)\), by
Lemma~\ref{p1b:lem:v2170-determinant-boundary}.

\begin{lemma}[A uniformly framed sphere cap]
\label{p1b:rt:lemma62-cap}
Put \(\delta=1/12\) and let \(C_\delta\) be the closed
geodesic cap of radius \(\delta\) about \(+1\in S^3_x\).
For \(x\in C_\delta\), uniformly in \(p\),
\[
 \|s(x,p)-e_1\|\le3\delta=1/4,
 \quad \operatorname{Re}(e_1^*s(x,p))\ge31/32,
 \quad |1+e_1^*s(x,p)|\ge63/32.
 \tag{RT.62.2}\label{p1b:eq:rt-62-2}
\]
Consequently the full frame
\[
 Q(x,p)=T(e_1,s(x,p))[e_2,e_3,e_1]\in SU(3)
 \tag{RT.62.3}\label{p1b:eq:rt-62-3}
\]
is defined on \(C_\delta\times S_p^3\), with last column \(s\).
\end{lemma}
\begin{proof}
At \(x=1\), the smash map is its based pole and \(s(1,p)=e_1\).
Along the geodesic in the first sphere with the second sphere
fixed, the smash map is one-Lipschitz and its suspended Hopf
map is three-Lipschitz, as proved in Subsection~\ref{p1b:sec:rtproof58}.
This gives the first bound. For two unit vectors,
\(\operatorname{Re}(e_1^*s)=1-\|s-e_1\|^2/2\), giving the
second and third bounds. The transporter of
Theorem~\ref{p1b:rt:theorem58-frame} is therefore defined.
The displayed constant frame has determinant one, proving \eqref{p1b:eq:rt-62-3}.
\end{proof}

Write \(x=\exp(v)\), \(v\in\operatorname{Im}\mathbb H\),
\(|v|<\delta\), and set \(g=\delta^2-|v|^2\). Define
\[
 q_\delta(x)=
 \frac{(|v|^2-g^2,\,-2g v)}{|v|^2+g^2}
 \quad (|v|<\delta),\qquad
 q_\delta(x)=1\quad (|v|\ge\delta).
 \tag{RT.62.4}\label{p1b:eq:rt-62-4}
\]
The quaternion at the cap centre is \(-1\).
At the cap boundary the expression converges to \(+1\).
This is inverse stereographic projection of \(-v/g\),
so it descends from the oriented cap quotient to a degree
\(+1\) map \(S^3\to S^3\). Use the usual boundary orientation
of the quaternion sphere: the imaginary coordinate chart at
\(-1\) has negative orientation. The derivative of the imaginary
part at \(v=0\) is \(-2\delta^{-2}I_3\); its determinant is
negative, confirming the stated sign.

Using the Hamilton quaternion-to-\(SU(2)\) identification of
Subsection~\ref{p1b:sec:rtproof61}, define
\[
 X(x,p)=Q(x,p)\operatorname{diag}(q_\delta(x),1)Q(x,p)^*
       \quad (x\in C_\delta),\qquad X=I\quad\text{off the cap},
 \qquad P_0=\sigma X.
 \tag{RT.62.5}\label{p1b:eq:rt-62-5}
\]
The notation \(P_0\) here denotes a gauge, not the sphere-product
base \(P\) or the perpendicular-plane projector.
The identity boundary value in \eqref{p1b:eq:rt-62-4} pastes these actual gauges.

\begin{theorem}[Actual initial representatives with no central calibration]
\label{p1b:rt:theorem62-roots}
Form the pointwise gauge words
\[
 A=XP_0,\qquad B=P_0X,\qquad
 H=XP_0X^{-1}P_0^{-1},\qquad V=AB=XP_0^2X.
 \tag{RT.62.6}\label{p1b:eq:rt-62-6}
\]
Then the following actual sections represent precisely the
six permitted initial fibre roots:
\[
 \boxed{
 \begin{array}{c|c|c}
 \text{target}&\text{first actual word}&\text{alternative actual word}\\\hline
 D&H^3\ (0,0;3)&NH^3\ (0,0;9)\\
 S&V\ (2,2;2)&NV\ (2,2;8)\\
 C&H^3V^{-1}\ (-2,-2;5)&NH^3V^{-1}\ (-2,-2;11).
 \end{array}}
 \tag{RT.62.7}\label{p1b:eq:rt-62-7}
\]
In particular this supplies matrices on the actual exotic
bundle, rather than substituting an untwisted \(SU(2)\) mapping
space. The exact pointwise identities are
\[
 \sigma X=P_0,\quad \sigma P_0=X,\quad
 \sigma A=B,\quad \sigma H=H^{-1},\quad \sigma V=BA.
 \tag{RT.62.8}\label{p1b:eq:rt-62-8}
\]
\end{theorem}
\begin{proof}
By the cap degree, \(X\) has primary \((1,0)\).
Write its unknown retained central coordinate as \(c\):
\([X]=(1,0;c)\). It is unnecessary to determine \(c\).
The actual exchange action \eqref{p1b:eq:rt-52-1} gives
\([P_0]=(0,1;-c)\).
In the actual Heisenberg law, their commutator is
\([H]=(0,0;1)\), so \([H^3]=(0,0;3)\).
The central coordinates of the two first generators cancel:
\([A]=(1,1;1)\) and \([B]=(1,1;0)\).
Their product is \([V]=(2,2;2)\).
The inverse of this component is \((-2,-2;2)\), hence
\([H^3V^{-1}]=(-2,-2;5)\).
Multiplication by the actual \(N\) adds six to each central
coordinate, proving all entries of \eqref{p1b:eq:rt-62-7}.
The pointwise identities follow from the involution and the
matrix words, without replacing commutators by central matrices.
The commutator is central only in the component group.
\end{proof}

\begin{corollary}[Explicit determinant seam endpoints]
\label{p1b:rt:corollary62-seam}
For the first determinant root the reduced closing endpoints
are exactly
\[
 F_h(0)=H^3,\qquad F_h(1)=NH^{-3}.
 \tag{RT.62.9}\label{p1b:eq:rt-62-9}
\]
They have the same actual gauge component, and for the alternative
root they are \(NH^3\) and \(H^{-3}\).
The combined seam likewise has explicit endpoints
\(H^3V^{-1}\) and \(N\sigma(H^3V^{-1})\).
\end{corollary}
\begin{proof}
The reduced determinant seam is \(F(1)=-\sigma F(0)\),
where minus acts on the actual rank-two fibre and fixes the
stabilization; this minus is precisely \(N\).
Use \eqref{p1b:eq:rt-62-8} and \(N^2=I\).
The endpoint components are the roots already proved in \eqref{p1b:eq:rt-62-7}.
\end{proof}

These determinant and combined endpoint formulas do not constitute
a closing homotopy. In particular \(H^3\) is not a pointwise square
root of \(N\), and \(H^6=N\) is not asserted as a matrix identity.
The next supplement supplies an actual coordinate closing homotopy
for the sphere root, for which a different conjugation identity applies.



\subsection{A coordinate sphere closing family on the exotic piece}
\label{p1b:sec:rtproof63}

Use the actual gauges \(X,P_0,A,B,V\) of
Theorem~\ref{p1b:rt:theorem62-roots}, on the reference piece of
Theorem~\ref{p1b:rt:theorem52-Ynormal}.
On the rank-two fibre, \(B=X^{-1}AX\) and \(V=AB\).
The intended seam endpoint is \(\sigma V=BA\).

\begin{theorem}[A removable principal-power branch in the actual closing path]
\label{p1b:rt:theorem63-closing}
For \(A\ne-I_E\) let \(A^t\) be its principal \(SU(2)\)
power. The formula
\[
 \boxed{F_s(t)=A^{-t}(AB)A^t,\quad 0\le t\le1,}
 \qquad F_s(t)=I_E\quad\text{when }A=-I_E,
 \tag{RT.63.1}\label{p1b:eq:rt-63-1}
\]
defines a continuous family of actual determinant-one gauges
on the exotic bundle, with
\[
 F_s(0)=V,
 \qquad F_s(1)=A^{-1}(AB)A=BA=\sigma V.
 \tag{RT.63.2}\label{p1b:eq:rt-63-2}
\]
Thus it descends to an actual \(S_Y\) of low data \((0,2q_Y)\).
Multiplying the family by \(N\) supplies the second permitted
root \((2,2;8)\) with the same seam and low data.
\end{theorem}
\begin{proof}
On a rank-two determinant-one fibre write
\(A=\cos\theta\,I_E+\sin\theta\,U\),
where \(U^*=-U\), \(U^2=-I_E\), and \(0<\theta<\pi\).
Its principal power is \(\cos(t\theta)I_E+\sin(t\theta)U\).
This expression is intrinsic under actual bundle-frame changes.
At \(A=I_E\) the power is identity.
At \(A=-I_E\) a direction \(U\) cannot in general be chosen
continuously. However \(B=X^{-1}AX=-I_E\), so \(AB=I_E\).
For the operator norm, uniformly in \(t\),
\[
 \|A^{-t}(AB)A^t-I_E\|=\|AB-I_E\|
 \le\|A+I_E\|+\|B+I_E\|=2\|A+I_E\|.
 \tag{RT.63.3}\label{p1b:eq:rt-63-3}
\]
The inequality follows by expanding
\(AB-I_E=(A+I_E)B-(B+I_E)\).
This proves continuity of \eqref{p1b:eq:rt-63-1} at its defined extension,
jointly in the base and in \(t\). It does not assert a global
continuous principal power of \(A\) at \(-I_E\).
Every value is unitary, determinant one, and fixes the specified
stabilizing line. The endpoint identities are matrix equalities.
They are exactly the reference mapping-torus seam, so this
constructs the actual family, beyond the component-existence
argument in Subsection~\ref{p1b:sec:rtproof51}. Its fibre primary \((2,2)\) gives
\(2q_Y\) and its determinant is identically one.
Finally \(N\) commutes with every fibre gauge and is exchange
invariant, so \(NF_s\) satisfies the same seam.
\end{proof}

For direct ambient evaluation put \(P_s=ss^*\), \(P_E=I-P_s\),
\(W=(A-A^*)/2\), and
\[
 c=\tfrac12(\operatorname{Tr}A-1),\quad
 h=\|W\|_F/\sqrt2,\quad\theta=\operatorname{atan2}(h,c).
\]
Away from \(h=0\), the actual stabilized power is
\[
 R_t=P_s+\cos(t\theta)P_E+
             \frac{\sin(t\theta)}hW,\qquad
 F_s(t)=R_t^*VR_t.
 \tag{RT.63.4}\label{p1b:eq:rt-63-4}
\]
The formula uses the actual last-column section \(s\), not a
global trivialization of its perpendicular bundle.
At the scalar endpoints \(V=I\) and the family is identity
before the optional multiplication by \(N\).

On the union of the two cap supports, the same transporter
\(T(e_1,s)\) is defined by \eqref{p1b:eq:rt-62-2} and its exchanged version.
In the frame \(T(e_1,s)[e_2,e_3,e_1]\), the two gauges are
the quaternions \(q_\delta(x)\) and
\((q_0,q_1,-q_2,-q_3)\) for \(q=q_\delta(p)\).
This is the conjugation by \(\operatorname{diag}(1,-1)\)
induced by \(\mathsf J\). The executable formula uses these
local quaternions to evaluate \eqref{p1b:eq:rt-63-4}; outside their support
the first family is identity. It uses no numerical branch cutoff.

\begin{theorem}[An exactly matched sphere joining component]
\label{p1b:rt:theorem63-joining}
On the actual joining locus \((x,p)=(w,-w)\), at least one of
\(X,P_0\) is identity. Thus \(A=B\), \(H=I\), and
the first sphere closing family is independent of \(t\):
\[
 F_s(t,w,-w)=A(w,-w)^2.
 \tag{RT.63.5}\label{p1b:eq:rt-63-5}
\]
In the concrete ordinary joining frames of
Theorem~\ref{p1b:rt:theorem60-frames},
\[
 \boxed{[j_Y^*S_Y]=(0,4,0).}
 \tag{RT.63.6}\label{p1b:eq:rt-63-6}
\]
This family already passes the sphere joining-secondary test,
with no \(T_Y\) adjustment.
\end{theorem}
\begin{proof}
Both cap supports would require \(w\) and \(-w\) to lie
within \(\delta\) of \(+1\), which is impossible for
\(2\delta<\pi\). The matrix equalities follow.
Then \(V=A^2\) commutes with every locally defined power
of \(A\), and at its scalar branch the same equality holds
by continuity.

Although \(A\) need not close on all of \(Y\), it does on
\(J\): \(\sigma A=B=A\) there.
Consequently it defines an actual joining gauge \(K_J\).
On a sphere fibre its primary is \(1+1=2\): both cap
gauges contribute degree one, and the antipodal sphere map
has degree \(+1\). Its determinant coordinate is zero.
Write \([K_J]=(0,2,\epsilon_K)\), without assigning its
secondary value. Equation \eqref{p1b:eq:rt-63-5} is the pointwise square
of this descended gauge. The exact joining multiplication
law \eqref{p1b:eq:rt-48-3} gives
\((0,2,\epsilon_K)^2=(0,4,0)\), proving \eqref{p1b:eq:rt-63-6}.
For the alternative root the additional ordinary joining
matrix is the constant \(-I_2\), which has identity component
in the ordinary joining gauge group, so its component is the same.
\end{proof}

\begin{corollary}[A concrete compatible sphere-piece endpoint]
\label{p1b:rt:corollary63-Z}
Let \(T_w=T(e_1,s(w,-w))\) and put on the split \(Z_s\) reference
\[
 K_Z(w)=T_w^*A(w,-w)T_w,\qquad S_Z=K_Z^2.
 \tag{RT.63.7}\label{p1b:eq:rt-63-7}
\]
These are actual determinant-one gauges fixing \(e_1\), with
\(K_Z(-w)=\mathsf JK_Z(w)\mathsf J\).
They are independent of the torus variables and descend to
\(Z_s\). The primary of \(K_Z\) is \(v_Z\), and the primary
of \(S_Z\) is \(2v_Z\). Their joining components are
\((0,2,\epsilon_K)\), \((0,4,0)\).
In particular under the explicit odd clutch \(c_1\) of
Theorem~\ref{p1b:rt:theorem60-clutch}, the actual sphere joining endpoints have identical
components \((0,4,0)\).
\end{corollary}
\begin{proof}
The global chart for \(T_w\) is proved in Subsection~\ref{p1b:sec:rtproof60}.
Equivariance of \(A\) on antipodal pairs and of \(T_w\)
proves descent. Its cover-sphere primary two is the pullback
of the projective-space generator one. The joining matrices
are exactly those of the preceding theorem in its transported
frames. Clutch conjugation by \((0,d,0)\) changes the last
coordinate by \(0\cdot4+0\cdot d=0\).
\end{proof}

For the zero-Chern coherent gluing \(c_0=T_w\), equations
\eqref{p1b:eq:rt-63-5}--\eqref{p1b:eq:rt-63-7} actually match pointwise. They therefore
give a complete coordinate sphere target on that coherent
reference. For \(c_1\), the supplied coordinate endpoints are
\(S_Y|_J\) and \(c_1S_Zc_1^*\); a rank-two joining homotopy
between these endpoints still has to be provided for a complete
coordinate odd target. The second \(Z_s\) piece class uses
the off-joining patch construction of Subsection~\ref{p1b:sec:rtproof50}, whose coordinate
patch family is also not supplied here.

\begin{lemma}[A coordinate joining toggle on the actual exotic piece]
\label{p1b:rt:lemma63-toggle}
In the cap frame \eqref{p1b:eq:rt-62-3}, an identity-endpoint reduced family
representing the first-sphere fourth direction can be written
using the suspended complex Hopf map. Put
\[
 r=(4(t-1/2),v/\delta)\in\mathbb R^4,
 \qquad g=1-|r|^2,
 \quad (\tau,z)=
 \frac{(|r|^2-g^2,\,2g r)}{|r|^2+g^2}
 \quad (|r|<1),
\]
where \(z\in\mathbb R^4=\mathbb C^2\).
Define the unit quaternion
\[
 q_\eta=(\tau,H(z)/|z|)\in S^3,
 \qquad H(z)=(|z_0|^2-|z_1|^2,
       2\operatorname{Re}(z_0\overline z_1),
       2\operatorname{Im}(z_0\overline z_1)).
 \tag{RT.63.8}\label{p1b:eq:rt-63-8}
\]
At \(z=0\) use zero imaginary part; off \(|r|<1\) use
\(q_\eta=1\). The gauge
\(Q\operatorname{diag}(q_\eta,1)Q^*\), extended by identity,
is a concrete \(Y\) joining toggle \(T_Y\) with
\[
 {}[j_Y^*T_Y]=(0,0,1).
 \tag{RT.63.9}\label{p1b:eq:rt-63-9}
\]
\end{lemma}
\begin{proof}
The radial compactification of the open four-ball followed
by inverse stereographic projection is the supported cap
quotient \(S^4\). Formula \eqref{p1b:eq:rt-63-8} is \(\Sigma\eta_2=\eta_3\).
Its imaginary numerator has length \(|z|^2\), so the map is
unit and continuous at \(z=0\); at the support boundary it
converges to \(+1\). Its support lies away from the cut
endpoints, so both endpoints are exactly identity and it closes
on the actual bundle. It has zero integral three-primary and
first-sphere fourth coefficient one. The retained joining
restriction of Theorem~\ref{p1b:rt:theorem53-joining} sends this fourth direction to its
secondary coefficient one, independently of any supported
seventh coefficient of the chosen frame. The latter coefficient
is not assigned by this construction.
\end{proof}

The executable coordinate module also evaluates \eqref{p1b:eq:rt-63-1}
on the \(Y\) stage of the actual boundary map \eqref{p1b:eq:rt-47-2}, and
\eqref{p1b:eq:rt-63-7} on its \(Z_s\) stage. It returns the actual odd
joining endpoints separately. The subsequent complete odd boundary comparison retains their joining homotopy and actual boundary bundle identification. Its whole supported seven-sphere map, rather than these endpoints alone, gives \(h_t=w_t=1,\zeta_{\rm nat}=0\) and the complete14/32 extension ledgers.





\section{Natural rank-two boundaries and extension decisions}
\label{p1c:sec:natural-rank-two-full-proofs}
The coefficients and YES/NO test ledger are determined by the exact
boundary and gauge calculations below, not by sampled finite diagnostics.

\subsection{Natural rank-two boundary coefficients: complete source and normalization}
\label{p1b:sec:natural-rank-two-completion}

The following proof chain retains the original odd exotic determinant
gauge, actual odd clutch, ordinary boundary frame over the fixed base,
both primitive circle cuts, both actual lower divisions, and the entire
matched joining. Matched changes of root, closing path and joining at
odd parity and ordinary frame changes over this fixed base preserve
the already normalized coefficients; they do not authorize a rotation
of the boundary base. The original determinant bit is \(h_t\), the
sphere bit is \(w_t\), and the rebased scalar bit is
\begin{equation}
 \zeta_{\rm nat}=\zeta_{\rm helper}=h_t+w_t=r_J\in\mathbb F_2.
 \label{p1b:eq:natural-scalar-notation}
\end{equation}
The scalar \(\zeta_{\rm nat}\) is distinct from the local circle
class \(\zeta=2\kappa\) used in the mixed-gauge construction.
Intermediate equations below retain symbolic coefficients and exact
auxiliary filling criteria. Their original status paragraphs are
superseded by the complete evaluation
\eqref{p1b:eq:natural-aj-4}--\eqref{p1b:eq:natural-aj-7}; no
auxiliary source homotopy or separate primitive is substituted for
the actual original whole-map calculation.


\subsection{The actual odd sphere comparison on the retained seven-cell}
\label{p1b:sec:natural-sphere-comparison}

This calculation retains the accumulated counts, ordinary bundle types,
natural boundary generators, and closures. Its sole classification target
is the pair \(h_t,w_t\), equivalently \(h_t,\zeta_{\rm nat}=h_t+w_t\).
The dependencies are: complete actual piece gauges and their matching;
an actual boundary bundle identification over the retained base;
removal of the known lower terms; evaluation of the resulting actual
supported maps by the existing detector. The sphere prerequisites are
applied below. No new detector or calibration is used.

Use the exotic reference \(E_s=s(x,p)^\perp\) of
Proposition~\ref{p1b:rt:proposition58-section}, with its actual exchange
\(\mathsf J=\operatorname{diag}(1,1,-1)\). Use \(\delta=1/12\)
and the actual cap gauges \(X,P_0\), words \(A=XP_0\),
\(B=P_0X\), \(V=AB\), and \(H=[X,P_0]\) of
Theorem~\ref{p1b:rt:theorem62-roots}. In particular the sphere
root is the actual \(V\) in component \((2,2;2)\), not a
replacement map into an untwisted fibre. Its closing family is
\[
 F_s(t,x,p)=A^{-t}(AB)A^t,\qquad
 F_s(0)=V,\quad F_s(1)=BA=\sigma V.                 \tag{NC.1}\label{p1b:eq:natural-nc-1}
\]
At \(A=-I_E\) its value is \(I_E\), with the removable
branch justified in Theorem~\ref{p1b:rt:theorem63-closing}.
Let \(S_Y\) be its descended gauge, and retain the actual
\(S_Z=T_w^*A(w,-w)^2T_w\) of
Corollary~\ref{p1b:rt:corollary63-Z}.

\begin{lemma}[A complete actual odd sphere triple]
\label{p1b:lemma:natural-sphere-joining}
Let \(j\) denote the quaternion \((0,0,1,0)\), and let
\(q_\delta\) be the degree-one cap map of Subsection~\ref{p1b:sec:rtproof62}.
In the retained joining coordinates \(\xi=a^2,u=a^{-1}w\), put
\[
 f_j(u)=q_\delta(j^{-1}u),\qquad
 c_j[a,w]=Q_Y(a,w)\operatorname{diag}(f_j(a^{-1}w),1)Q_Z(a)^*.
                                                        \tag{NC.2}\label{p1b:eq:natural-nc-2}
\]
This is an actual odd clutch in ordinary component \((0,1,0)\).
The triple \((S_Y,S_Z,H_s)\), with
\(H_s(\lambda,a,w)=S_Y|_J(a,w)\) constant in \(\lambda\),
is a complete actual global sphere target of low data \((0,2q)\).
Its matching toggle is zero.
\end{lemma}
\begin{proof}
The cap of \(f_j\) is centered at \(j\) and avoids the joining
basepoint \(u=1\). Left translation preserves degree, so the
based map has degree one and is independent of \(\xi\);
its ordinary joining component is exactly \((0,1,0)\).
The joining frame formulas of
Theorem~\ref{p1b:rt:theorem60-frames} show its descent under
\((a,w)\mapsto(-a,-w)\). The ordinary normal form of
Theorem~\ref{p1b:rt:theorem53-Wnormal} permits this representative
of the same odd bundle and the same exotic/split piece pair.
This is a choice within the established ordinary class, not a
change of the boundary base or of its product normal coordinates.

On the support of \(f_j(a^{-1}w)\), the quaternion \(w\)
is within \(\delta\) of \(aj\). Since \(a\) is complex
unit, \(aj\) is purely imaginary and its distance from either
\(+1\) or \(-1\) is \(\pi/2\). Thus this support misses
the two supports of \(A(w,-w)\), since \(2\delta<\pi/2\).
There the sphere gauge is identity. Off the clutch support,
\(c_j=T_w\). Consequently, as actual matrices everywhere,
\[
 c_jS_Zc_j^*=A(w,-w)^2=S_Y|_J.                       \tag{NC.3}\label{p1b:eq:natural-nc-3}
\]
This supplies the required joining homotopy immediately, including
both endpoints. Subsection~\ref{p1b:sec:rtproof63} already evaluates the sphere joining
component as \((0,4,0)\), so its \(T_Y\) exponent is zero.
The actual triple therefore pastes by
Theorem~\ref{p1b:rt:theorem47-gauges}. Its fixed low data and
ordinary piece pair are retained. Any transport to another actual
representative is an ordinary bundle identification over the same
base; its boundary frame action preserves \(w_t\) by
Corollary~\ref{p1b:rt:corollary54-top}.
\end{proof}

Write the retained boundary as \(R=S^1_s\times S^3_y\times L\),
\(L=\mathbb{RP}^3\), with \(a=i e^{is}\), and use the exact
three-stage boundary map \eqref{p1b:eq:rt-47-2}. On its inner stage,
\[
 \alpha=3\pi\rho/2,\quad 0\le\rho\le1/3,\qquad
 x=y e^{\alpha u},\quad p=y e^{-\alpha u},\quad u\in S^2.
                                                        \tag{NC.4}\label{p1b:eq:natural-nc-4}
\]
At the neck \(w=yu\). No base normal-frame rotation is made.

\begin{lemma}[A coherent reference frame along the complete inner stage]
\label{p1b:lemma:natural-coherent-frame}
Let \(s_\alpha=s(y e^{\alpha u},y e^{-\alpha u})\),
\(a_y=s(y,y)\), and set
\[
 F_0(y)=[(-\overline{a_{y,2}},\overline{a_{y,1}},0),e_3,a_y].
                                                        \tag{NC.5}\label{p1b:eq:natural-nc-5}
\]
Subdivide \([0,\pi/2]\) into sixteen equal subintervals.
Starting from \(F_0\), multiply the successive exact transporters
\(T(s_{\alpha_k},s_{\alpha_{k+1}})\) on the left, using
\(T(s_{\alpha_k},s_\alpha)\) on the final partial subinterval.
This defines a continuous actual \(SU(3)\) frame \(F_Y(y,u,\alpha)\)
whose last column is \(s_\alpha\). Put
\[
 F_b=F_Y(y,u,\pi/2),\quad
 F_Z=T_w^*F_b,\quad Q_0=[e_2,e_3,e_1],\quad C_0=Q_0^*F_Z.
                                                        \tag{NC.6}\label{p1b:eq:natural-nc-6}
\]
The matrix \(C_0\) fixes \(e_3\) and has its first block in
\(SU(2)\). Together with the constant cylinder frame \(F_b\)
and outer frame \(F_Z\), these give the coherent zero-Chern
boundary identification with \(\mathbf1\oplus\lambda_L\).
\end{lemma}
\begin{proof}
The retained section is three-Lipschitz. Each transport step has
section distance at most \(3\sqrt2\pi/32<1\), so its
\(\langle s_\alpha,s_{\alpha'}\rangle\ne-1\) chart is valid
uniformly. The transporter of Subsection~\ref{p1b:sec:rtproof58} is an exact \(SU(3)\)
matrix, not a numerical integration rule, and sends the first
section to the second. Products match at every subdivision point.
The initial frame is unitary of determinant one since \(a_y\)
lies in the first two coordinate axes. The assertions about
last columns and determinants follow. Covariance gives
\[
 F_Y(y,-u,\alpha)=\mathsf J F_Y(y,u,\alpha)
                    \operatorname{diag}(1,-1,1).        \tag{NC.7}\label{p1b:eq:natural-nc-7}
\]
At \(\alpha=\pi/2\) the last column is \(s(w,-w)\).
The outer last column is \(e_1\), so \eqref{p1b:eq:natural-nc-6} fixes \(e_3\).
Equation \eqref{p1b:eq:natural-nc-7} and \(T_{-w}=\mathsf JT_w\mathsf J\)
give precisely the ordering-line transition on the outer
\(u\sim-u\) boundary. The positive and negative axes at the
inner origin are the retained ones, and the normal coordinates
are those of \eqref{p1b:eq:natural-nc-4}.
\end{proof}

\begin{lemma}[An explicit actual identification for the odd boundary bundle]
\label{p1b:lemma:natural-odd-boundary-frame}
Let \(D_a=\operatorname{diag}(1,a,1)\). In the coherent frames
the odd neck transition is
\[
 f_\partial=C_0(y,u)^*D_a
       \operatorname{diag}(f_j(a^{-1}yu),1)D_a^*C_0(y,u).
                                                        \tag{NC.8}\label{p1b:eq:natural-nc-8}
\]
It has the following continuous \(SU(2)\) extension \(R\) over
the entire inner normal disk, with polar radius \(r=3\rho\).
Set \(q_r=e^{\pi r u/2}\) and \(x=yq_r\). On the active
clutch tube write \(j^{-1}a^{-1}x=e^v\), \(|v|<\delta\).
For \(r\ge1/2\), put \(h=2(1-r)\) and
\[
 \mathcal A_h(a)=(\cos(\pi h/2)a+\sin(\pi h/2)j)
                 (\cos(\pi h/2)+\sin(\pi h/2)j)^{-1},
\]
\[
 \widetilde y_h=\mathcal A_h(a)e^{(1-h)\pi j/2}
                         e^{(1-h)v}e^{-(1-h)\pi u/2},
 \qquad C_r=C_0(\widetilde y_h,u).                       \tag{NC.9}\label{p1b:eq:natural-nc-9}
\]
For \(r\le1/2\), put
\[
 C_r=\operatorname{diag}(e^{2ir\ell(u)},e^{-2ir\ell(u)},1),
                                                        \tag{NC.10}\label{p1b:eq:natural-nc-10}
\]
where the following is a globally continuous real logarithm:
\[
 \ell(u)=\sum_{k=0}^{15}
  \arg(1+\overline{\langle s_{k\pi/32}(1,u),
                              s_{(k+1)\pi/32}(1,u)\rangle})
       -\arg(1+\overline{s_{\pi/2,1}(1,u)}).             \tag{NC.11}\label{p1b:eq:natural-nc-11}
\]
Here the notation in the last term means the first coordinate
of the section at \(y=1,\alpha=\pi/2\).
Now define, on that tube,
\[
 R=C_r^*D_a\operatorname{diag}(f_j(a^{-1}x),1)D_a^*C_r,
                                                        \tag{NC.12}\label{p1b:eq:natural-nc-12}
\]
and define \(R=I\) off it. Then \(R|_{r=1}=f_\partial\),
and \(R|_{r=0}\) is independent of \(u\).
The inner actual frame \(F_YR\), cylinder frame \(F_bf_\partial\),
and outer frame \(F_Z\) are an actual ordinary boundary bundle
identification over the fixed \(R\).
\end{lemma}
\begin{proof}
The joining frames are \(Q_Y=T_wQ_0D_a\), \(Q_Z=Q_0D_a\).
Substitution of \eqref{p1b:eq:natural-nc-2} proves \eqref{p1b:eq:natural-nc-8} exactly.
For \(y=1\), all sections in the transport path lie in the
\((e_1,e_3)\) plane. Each transport acts on \(e_2\) by
\((1+\bar c)/|1+c|\), with \(c\) the adjacent section
overlap. The neck transporter acts by the same formula using
\(c=s_{\pi/2,1}\). Since all first arguments in \eqref{p1b:eq:natural-nc-11}
have positive real part, and the neck bound is
\(\operatorname{Re}s_{\pi/2,1}\ge-7/9\), their principal
arguments are continuous. Thus
\(C_0(1,u)=\operatorname{diag}(e^{i\ell(u)},e^{-i\ell(u)},1)\).

The actual periodic based contraction \(\mathcal A_h\) is
the retained quaternion formula of Subsection~\ref{p1b:sec:rtproof32}; its two endpoints
are \(a\) and \(1\). At \(r=1\), the active-tube equation
gives \(y=aj e^v(-u)\), so \(\widetilde y_0=y\).
At \(r=1/2\), \(\widetilde y_1=1\), proving equality
of \eqref{p1b:eq:natural-nc-9} and \eqref{p1b:eq:natural-nc-10}. At \(r=0\), \(C_r=I\) and
\(q_r=1\), so \eqref{p1b:eq:natural-nc-12} is independent of the polar direction.
The cap map tends to identity at the tube boundary. Since all
conjugators are unitary, \eqref{p1b:eq:natural-nc-12} pastes continuously to identity
there, even though \(C_r\) is needed only on the tube.
All formulas fix the last coordinate and are determinant one.

The displayed actual frames agree at the inner seam. At the
outer cylinder seam, \(c_jF_Z=F_bf_\partial\) by \eqref{p1b:eq:natural-nc-8}.
At the projective outer boundary their positive/negative
transition is \eqref{p1b:eq:natural-nc-7}; thus they give a full identification
with the prescribed \(\mathbf1\oplus\lambda_L\), including
the joining data. This is over the identity of the boundary
base. In particular it is not the base rotation of Subsection~\ref{p1b:sec:rtproof44}.
\end{proof}

\begin{theorem}[The actual sphere input for the retained detector]
\label{p1b:theorem:natural-sphere-input}
The full sphere triple restricts in the identification above to
\[
 \widetilde S(s,y,\rho,u)=
 R^*F_Y^* S_Y(a,x,p)F_YR\quad(\rho\le1/3),            \tag{NC.13}\label{p1b:eq:natural-nc-13}
\]
and to \(F_Z^*S_Z(yu)F_Z\) on both outer stages.
The actual difference
\[
 \Phi_S(s,y,\rho,u)=\widetilde S(s,y,\rho,u)
                         \widetilde S(0,y,\rho,u)^{-1}  \tag{NC.14}\label{p1b:eq:natural-nc-14}
\]
is identity at \(s=0\), on \(\rho\ge1/3\), and at \(y=-1\).
It therefore defines an actual supported map
\[
 S^1\wedge(S^3/\{-1\})\wedge(D^3/\partial D^3)
           =S^7\longrightarrow SU(2),\qquad
                     [\Phi_S]=w_t\mu.                  \tag{NC.15}\label{p1b:eq:natural-nc-15}
\]
For the retained \(y\)-basepoint \(+1\), use
\(\Phi_S(s,-y,\rho,u)\). The explicit isotopy
\(y\mapsto e^{i\pi\tau}y\), with identity on the other base
coordinates and on \(\lambda_L\), justifies this rebase.
It changes no natural top coefficient.
\end{theorem}
\begin{proof}
Equations \eqref{p1b:eq:natural-nc-3} and \eqref{p1b:eq:natural-nc-8} show that the inner conjugation
agrees with the outer value at the neck: the support of
\(f_\partial\) misses the joining sphere support, so it
commutes with that gauge. The complete constant joining
homotopy therefore has exactly the outer value in \eqref{p1b:eq:natural-nc-13}.
That outer value is independent of \(s\), proving the lens
collapse in \eqref{p1b:eq:natural-nc-14}. The circle basepoint collapse is immediate.
For \(y=-1\), both inner quaternions \(-e^{\pm\alpha u}\)
have distance at least \(\pi/2\) from \(+1\); hence
\(X=P_0=I\) there. The joining quaternion \(-u\) is also
purely imaginary. Thus the complete sphere gauge is identity
at this \(y\), including its conjugation by \(R\).
The polar origin is legitimate by Lemma~\ref{p1b:lemma:natural-odd-boundary-frame}.
The upper two-dimensional lens skeleton is outside this inner
three-cell. Hence the map really is supported on the natural
seven-cell, rather than on a replaced cover sphere.

To identify its coefficient, let \(B_0(y,n)=\widetilde S(0,y,n)\).
This is an actual gauge on the fixed boundary bundle, independent
of \(s\). In \(\operatorname{Map}(S^1\times S^3,H)\) it
comes from \(\operatorname{Map}(S^3,H)\); it has neither a
parameter circle nor a parameter four-cell component. In particular
its natural \(e,t\) components vanish. Its sphere component
is the actual known lower part of \(\widetilde S\): primary
\(2b\) and the six-stage coefficient carried by this ordinary
frame. An ordinary frame may shear that coefficient by \(4N\),
exactly as in Theorem~\ref{p1b:rt:theorem54-transport}, but cannot
alter \(w_t\). Division by the actual \(B_0\) removes those
lower terms with an exact comparison equation, not a guessed
calibration. The remaining top factor is central and \eqref{p1b:eq:natural-nc-14}
is exactly that factor, proving \eqref{p1b:eq:natural-nc-15}. No independent
identification of a primitive benchmark with the target is used.
The stated isotopy lifts because \(\lambda_L\) depends only
on the unchanged normal coordinate; composition along it is
an actual gauge homotopy. It is independent of \(s\), and
does not make a parameter-dependent normal-frame rotation.
\end{proof}

\subsection{An evaluated filling for the actual clutch contribution}
\label{p1b:sec:natural-clutch-filling}

The following removes a named component of the actual map \eqref{p1b:eq:natural-nc-14}.
It does not construct another detector. Retain the reference frame
\(F_Y\), the prescribed normal coordinates, and the actual transition
extension \(R\) above. On the inner stage put
\[
\begin{gathered}
 G_0=F_Y^*S_YF_Y,\qquad \Phi_0=G_0(s)G_0(0)^{-1},\qquad C_R(s)=R(s)^{-1}G_0(s)R(s)G_0(s)^{-1}, \\
 \Psi_C=C_R(s)C_R(0)^{-1}.
\end{gathered}
 \tag{NC.19}\label{p1b:eq:natural-nc-19}
\]
Define these comparisons to be identity on the outer stages where
their values are circle independent or commute, respectively.
The formulas use the unrebased \(y=-1\) basepoint of \eqref{p1b:eq:natural-nc-15}.
Precomposing every comparison, homotopy and filling with \(y\mapsto-y\)
gives the retained \(+1\) basepoint, with exactly the previously
specified isotopy and no normal-frame change.

\begin{lemma}[An exact split into two actual supported comparisons]
\label{p1b:lemma:natural-sphere-split}
Both \(\Phi_0\) and \(\Psi_C\) are based maps on the retained
seven-cell. Their supports are disjoint and
\[
             \Phi_S=\Phi_0\Psi_C                     \tag{NC.20}\label{p1b:eq:natural-nc-20}
\]
as actual matrices. The first support requires both \(x,p\)
to lie in their \(+1\) caps. On the second support \(X=I\)
and the sphere gauge is the actual \(P_0^2\), independent
of the circle.
\end{lemma}
\begin{proof}
With either cap gauge identity, \(A=B\) and the entire closing
family is constant. Thus \(\Phi_0\) is supported only where
both caps occur. The support of \(R\) has \(x\) within
\(\delta\) of the purely imaginary \(aj\); it misses the
\(+1\) cap. There \(X=I\), so the same closing family
is exactly \(P_0^2\). On the neck \(R\) commutes with the
joining gauge, so \(\Psi_C=I\). At \(y=-1\) both underlying
gauges are identity. The circle basepoint is also identity
by the definitions, so both comparisons descend to \(S^7\).
Algebraically \(\Phi_S=C_R(s)\Phi_0C_R(0)^{-1}\).
At every point with \(C_R(s)\) or \(C_R(0)\) nonidentity,
\(\Phi_0=I\). This proves the actual equality \eqref{p1b:eq:natural-nc-20}.
\end{proof}

\begin{lemma}[Actual homotopies to a fixed product-cap clutch word]
\label{p1b:lemma:natural-clutch-reduction}
The actual \(\Psi_C\) is based homotopic, through rank-two
comparisons fixing the stabilizing line, to
\[
 \Psi_{
 \mathrm{red}}=C_sC_0^{-1},\qquad
 C_s=f_s^{-1}Q^2 f_s Q^{-2},\quad
 f_s=D_a fD_a^*,\quad Q=\operatorname{diag}(q,1).       \tag{NC.21}\label{p1b:eq:natural-nc-21}
\]
Here \(f=q_\delta(i^{-1}x)\) is supported in the fixed cap
about the pure quaternion \(i\), and
\(q=(q_0,q_1,-q_2,-q_3)\) for \(q_\delta(p)=(q_0,q_1,q_2,q_3)\).
The source chart is the actual product of those two caps.
Both cap maps have degree one.
\end{lemma}
\begin{proof}
First flatten the auxiliary \(C_r\) in \eqref{p1b:eq:natural-nc-12} only in
an interior radial band. Let \(\chi_1\) be zero outside
\([3/8,5/8]\), one on \([7/16,9/16]\), and linear
on its two intervening intervals. On the active tube, for
\(r\ge1/2\), replace \(C_r\), at homotopy time \(v\), by
\[
 C_0\bigl(\exp((1-v\chi_1)\log\widetilde y_h),u\bigr)
 C_0(1,u)^*
 \operatorname{diag}(e^{i(1-v\chi_1)\ell(u)},
                    e^{-i(1-v\chi_1)\ell(u)},1).       \tag{NC.22}\label{p1b:eq:natural-nc-22}
\]
For \(r\le1/2\), replace it by
\(\operatorname{diag}(e^{2ir(1-v\chi_1)\ell},
e^{-2ir(1-v\chi_1)\ell},1)\).
These are exact homotopies of the actual matrix formula.
They match at \(r=1/2\) and leave the center and neck
unchanged. The logarithm in \eqref{p1b:eq:natural-nc-22} is legitimate: the retained
contraction satisfies \(d(\mathcal A_h(a),1)\le\pi(1-h)\),
so on this band
\[
 d(\widetilde y_h,1)\le(2\pi+\delta)(2r-1)
                      \le(2\pi+\delta)/4<\pi.
\]
This is a proved chart bound, rather than a numerical branch choice.
At the final time \(C_r=I\) on the plateau.

On an interaction of the clutch cap with the sphere gauge,
\(p\) lies in the \(+1\) cap and \(x\) lies in a cap
about a pure imaginary unit. Consequently
\[
 |r-1/2|\le2\delta/\pi<1/18.                          \tag{NC.23}\label{p1b:eq:natural-nc-23}
\]
This follows from \(p^{-1}x=e^{\pi r u}\) and the spherical
triangle inequality. All such interactions lie inside the
flattened plateau. Now let \(\chi_2\) be zero outside
\([7/16,9/16]\), one on \([4/9,5/9]\), and linear on
the two intervening intervals. Move the interior support center by
\[
 c_v(a,r)=\cos(\pi v\chi_2(r)/2)\,aj
                  +\sin(\pi v\chi_2(r)/2)\,i,\qquad
 R_v=D_a\operatorname{diag}(q_\delta(c_v^{-1}x),1)D_a^*.
                                                        \tag{NC.24}\label{p1b:eq:natural-nc-24}
\]
The two summands are orthogonal pure imaginary quaternions,
so \(c_v\) is unit and always has distance \(\pi/2\)
from \(\pm1\). This formula is used only on the flattened
band and is pasted to the already flattened \(R\) outside it.
It leaves the actual neck transition \eqref{p1b:eq:natural-nc-8} unchanged.
It is an ordinary boundary homotopy; it does not replace the
clutch on the global joining quotient by a non-descending map.
Bound \eqref{p1b:eq:natural-nc-23} applies throughout the motion. At its final time
every nonidentity clutch--sphere interaction therefore has
the fixed center \(i\) and \(R=D_a\operatorname{diag}(f,1)D_a^*\).

In this product-cap chart, the actual pair-coordinate inverse is
\[
 \theta=|\log(p^{-1}x)|,\quad u=\log(p^{-1}x)/\theta,
 \quad y=xe^{-\theta u/2},\quad \rho=\theta/(3\pi).
                                                        \tag{NC.25}\label{p1b:eq:natural-nc-25}
\]
Its chart avoids antipodal pairs. In the common local frame
\(Q_s=T(e_1,s)[e_2,e_3,e_1]\), the actual sphere word
is exactly \(Q^2\). The residual boundary frame is the
actual \(U(x,p)=F_Y^*Q_s\), a block \(SU(2)\) matrix
independent of \(s\). Write \(x=i e^{v_x}\), \(p=e^{v_p}\).
The explicit path
\[
 U(x,p)\longrightarrow
 U(i e^{(1-v)v_x},e^{(1-v)v_p})
 \longrightarrow U(i,1)^{1-v}                         \tag{NC.26}\label{p1b:eq:natural-nc-26}
\]
contracts it to identity. The second segment is a path of one
fixed matrix, so an axis at a scalar endpoint may be fixed once;
no base-dependent principal-power singularity is involved.
Apply \eqref{p1b:eq:natural-nc-26} inside the commutator \eqref{p1b:eq:natural-nc-19}. At either cap
boundary one of its two factors is identity, so the homotopy
pastes by identity and retains every based collapse.
This yields \eqref{p1b:eq:natural-nc-21} by actual homotopies of the comparison,
including the frame terms and the interior support motion.
\end{proof}

\begin{theorem}[The actual clutch contribution has detector parity zero]
\label{p1b:theorem:natural-clutch-parity}
There is an explicit filling \(G_C:D^8\to SU(3)\) with
\(G_C|_{S^7}=\operatorname{diag}(\Psi_C,1)\). For its
last-column map \(g_C:S^8\to S^5\), the retained square
detector evaluates to zero. Thus
\[
                    [\Psi_C]=0\in\pi_7SU(2).          \tag{NC.27}\label{p1b:eq:natural-nc-27}
\]
This is an evaluation of the actual clutch contribution, not
an assignment of the full natural sphere bit.
\end{theorem}
\begin{proof}
In \eqref{p1b:eq:natural-nc-21} set \(f_0=D_i fD_i^*\),
\(\chi=f_sf_0^{-1}\), and use the actual normalized identity
\[
 \Psi_{\mathrm{red}}
      =f_0^{-1}\chi^{-1}Q^2\chi Q^{-2}f_0.            \tag{NC.28}\label{p1b:eq:natural-nc-28}
\]
Let \(K_t(\lambda,f)\) be exactly the rank-one-phase
\(SU(3)\) path \eqref{p1b:eq:rt-61-1}--\eqref{p1b:eq:rt-61-4}, with
\(\lambda=e^{-is}\), \(K_0=I\), \(K_1=[f,
\operatorname{diag}(\lambda,1,1)]\).
Define
\[
 \widehat\chi_t=D_iK_t^{-1}D_i^*,\qquad
 \widehat G_t=f_0^{-1}\widehat\chi_t^{-1}
                      Q^2\widehat\chi_t Q^{-2}f_0.    \tag{NC.29}\label{p1b:eq:natural-nc-29}
\]
On the first two columns, \(\operatorname{diag}(e^{-is},1)\)
has the same conjugation action as
\(\operatorname{diag}(1,e^{is})\). Hence
\(\widehat\chi_1=\chi\) exactly, so
\(\widehat G_0=I\), \(\widehat G_1=\Psi_{\mathrm{red}}\).
It is identity on the based wedge for all \(t\), giving a
concrete based stabilized filling of \eqref{p1b:eq:natural-nc-28}.

Precede its boundary by the reverse of the three actual
rank-two homotopies \eqref{p1b:eq:natural-nc-22}, \eqref{p1b:eq:natural-nc-24}, and \eqref{p1b:eq:natural-nc-26} on an
outer radial collar. More explicitly use \eqref{p1b:eq:natural-nc-29} with
\(t=2r\) for filling radius \(0\le r\le1/2\), and
use the reversed concatenated homotopies for \(1/2\le r\le1\).
The seam is exactly \(\Psi_{\mathrm{red}}\); the outer boundary
is exactly \(\Psi_C\). This is the stated actual \(G_C\).
Every matrix on its collar is rank two and fixes \(e_3\),
so the collar projects constantly and adds no term to \(g_C\).

In \eqref{p1b:eq:natural-nc-29} the sphere variable \(q\) occurs only as \(q^2\).
The last-column map therefore factors exactly through
\[
 S^8=\Sigma(S^1\wedge S^3_f\wedge S^3_q)
   \xrightarrow{\,\Sigma(1\wedge1\wedge(q\mapsto q^2))\,}
 S^8\longrightarrow S^5.                             \tag{NC.30}\label{p1b:eq:natural-nc-30}
\]
The first map has degree two. The final map need not be
calibrated: if its class is any \(n\nu_5\in C_{24}\),
then the class of the actual \(g_C\) is \(2n\nu_5\).
The two cap maps used to obtain this chart have degree one,
and the intervening homotopies project constantly on the collar;
there is no additional unidentified summand. Apply the actual
parity boundary of Lemma~\ref{p1b:rt:lemma56-parity} and
Theorem~\ref{p1b:rt:theorem56-detector}: the square coefficient
is \(2n\bmod2=0\), proving \eqref{p1b:eq:natural-nc-27}.
This same cancellation agrees with the retained action
\(\rho(0,2;0)(r,0,0)=(r,0,0)\) in the G20 mapping two-type,
but the displayed filling proves it directly for the actual input.
\end{proof}

\begin{theorem}[The reduced actual input for the natural sphere bit]
\label{p1b:theorem:natural-closing-input}
For unit quaternions \(X,P\), put \(A=XP\), \(B=PX\),
\(V=AB\), and define
\[
 F(\theta)=
 \begin{cases}
 A^{-\theta/\pi} V A^{\theta/\pi},&0\le\theta\le\pi,\\
 B^{-(\theta/\pi-1)}(BA)B^{\theta/\pi-1},&\pi\le\theta\le2\pi.
 \end{cases}
\]
\[
 \boxed{L(s,X,P)=F((\pi/2+s)\bmod2\pi)F(\pi/2)^{-1},
                   \qquad [L]=w_t\mu.}                \tag{NC.31}\label{p1b:eq:natural-nc-31}
\]
At the scalar branches use the removable value proved for
the actual sphere family. The domain is the actual supported
\(S^1\wedge S^3_X\wedge S^3_P=S^7\).
\end{theorem}
\begin{proof}
The common local frame on the two \(+1\) caps is again
\(Q_s=T(e_1,s)Q_0\). The actual cap gauges in it are
exactly \(q_\delta(x)\) and the rotation of
\(q_\delta(p)\) used above, by Subsection~\ref{p1b:sec:rtproof63}. Thus the actual
\(\Phi_0\) is the conjugate of \eqref{p1b:eq:natural-nc-31} by
\(U=F_Y^*Q_s\). Contract this frame by the explicit path
\(U(e^{(1-v)\log x},e^{(1-v)\log p})\); at the final
point \(x=p=1\), it is identity. The pair chart \eqref{p1b:eq:natural-nc-25}
is valid throughout, with its continuous value at \(x=p\).
The comparison is identity when either cap variable is identity,
so this homotopy pastes by identity at both cap boundaries.
The cap product is an actual coordinate neighborhood of
\((y,n)=(1,0)\) in the inner six-cell. Its quotient has
degree \(\pm1\); each cap map has degree one, and the rotation
of the second quaternion has degree one. Such a sign is irrelevant
in \(\pi_7SU(2)=C_2\). This proves the identification of
the actual \(\Phi_0\) with \eqref{p1b:eq:natural-nc-31} in the retained top group.
Finally \eqref{p1b:eq:natural-nc-20} and the evaluated actual cancellation \eqref{p1b:eq:natural-nc-27}
give \([\Phi_S]=[\Phi_0]=[L]\). This retains the prescribed
baseline \(a=i\), the complete joining, and the actual frame
homotopies; no benchmark map is substituted.
\end{proof}

\paragraph{Evaluation of the full sphere equation.}
The actual odd-clutch filling has been constructed and evaluated;
its parity is zero. The remaining specific \eqref{p1b:eq:natural-nc-31}, implemented by
\texttt{closing\_comparison}, is now evaluated in
Theorem~\ref{p1b:theorem:natural-sphere-bit}. Its whole-map comparison,
including all lower-cell and splitting terms, constructs the filling
\[
 \mathcal F_L:D^8\to SU(3),\qquad
 \mathcal F_L|_{S^7}=\operatorname{diag}(L,1),\qquad
 \operatorname{Sq}^4u_5=v_9
       \text{ on }C_{[\mathcal F_Le_3]}.               \tag{NC.32}\label{p1b:eq:natural-nc-32}
\]
The proof uses the complete actual branch, its global framed comparison,
and the actual source top component. With the already evaluated clutch
term it gives \(w_t=1\). The retained one-point diagnostic in the
appendix is not used as a parity argument.

\paragraph{The retained determinant endpoint equations.}

For the determinant bit, the actual initial representative is
the retained \(H^3\) of component \((0,0;3)\).
Since \(\sigma H=H^{-1}\) exactly, the rank-two closing family
constructed in \eqref{p1b:eq:natural-nd-1}--\eqref{p1b:eq:natural-oa-4} has the exact endpoint equations
\[
 F_h:[0,1]\to\mathcal G_{SU(2)}(E_s),\qquad
 F_h(0,x,p)=H(x,p)^3,\quad
 F_h(1,x,p)=N(x,p)H(x,p)^{-3}.                          \tag{NC.17}\label{p1b:eq:natural-nc-17}
\]
Every value is an actual section of the gauge bundle and must
fix \(s(x,p)\).
The prescribed determinant gauge is \(e^{\pi it}F_h(t)\).
The chosen local joining origin is now measured as \(e_h=0\), and
the toggle exponent \(e_h+1=1\) is applied in \eqref{p1b:eq:natural-nd-5}. The constructed
joining path has the exact endpoints, in the actual ambient frames,
\[
 H_h(0,a,w)=D_Y|_J(a,w),\qquad
 H_h(1,a,w)=c_j(a,w)\operatorname{diag}(1,-a^2,1)c_j(a,w)^*.
                                                        \tag{NC.18}\label{p1b:eq:natural-nc-18}
\]
The right endpoint is the actual split \(D_Z\) of \eqref{p1b:eq:rt-50-7}
on \((z_1,z_2)=(a,-a)\), with positive/negative axes
\((e_2,e_3)\). The local closing, the measured joining origin, the applied matching toggle
and the joining homotopy are now supplied in \eqref{p1b:eq:natural-nd-1}--\eqref{p1b:eq:natural-nd-10}. The global
closing patch is \eqref{p1b:eq:natural-nd-11}, constructed in \eqref{p1b:eq:natural-ny-1}--\eqref{p1b:eq:natural-oa-4}. Equality
of its endpoint components alone does
not calculate the determinant boundary comparison or \(h_t\).

The complete original comparison is evaluated in
\eqref{p1b:eq:natural-aj-4}--\eqref{p1b:eq:natural-aj-7}:
\(h_t=w_t=1\), \(r_J=\zeta_{\rm nat}=0\), with every retained
extension test decided. This component is used with its full frames,
lower divisions and joining; its isolated value is not a replacement
for that complete comparison. In particular \(S_{0,0}\) fails and \(S_{0,1}\) extends. 


\subsection{An actual determinant closing chart and its matched joining}
\label{p1b:sec:natural-determinant-joining}

Keep the actual section \(s(x,p)\), the stabilizing line \(\mathbb Cs\),
the actual root \(H^3=[X,P_0]^3\), and \(N=2ss^*-I\).
The following supplies the closing family on a neighborhood of the entire
joining locus, measures its joining secondary origin, applies the required
toggle, and supplies a joining homotopy with the retained odd clutch.
It does not supply the remaining global closing patch or evaluate \(h_t\).

\begin{proposition}[Actual local closing and its exact singular set]
\label{p1b:proposition:natural-determinant-local}
Put \(P_s=I-ss^*\), \(v_s=P_se_2\), and, where \(v_s\ne0\),
\[
 \ell_s=\frac{v_sv_s^*}{\|v_s\|^2},\qquad
 K_+(t)=ss^*+\cos(\pi t)P_s
                 +i\sin(\pi t)(2\ell_s-P_s).           \tag{ND.1}\label{p1b:eq:natural-nd-1}
\]
It is an actual determinant-one rank-two gauge, fixes \(s\), and
has endpoints \(I,N\). On the complement of the two-cap product
\(B=\overline{B_\delta(1)}\times\overline{B_\delta(1)}\),
where \(H=I\), it solves the prescribed closing endpoints
\(H^3,NH^{-3}\) exactly. Its descended determinant gauge is
\[
 D_+(a,x,p)=ss^*+a^2\ell_s+(P_s-\ell_s).                \tag{ND.2}\label{p1b:eq:natural-nd-2}
\]
In the cut circle \(a=e^{\pi it}\), its reduced family is
exactly \eqref{p1b:eq:natural-nd-1}, and its determinant is \(a^2=e^{2\pi it}\).
It obeys \(D_+(-a,p,x)=\mathsf J D_+(a,x,p)\mathsf J\).

The entire singular set of this local chart is the explicit circle
\[
 \Sigma_+=\{x=p=y(r):\quad
 y(r)=(-1/3,\ (2\sqrt2/3)\cos r,
                 (2\sqrt2/3)\sin r,\ 0)\}.             \tag{ND.3}\label{p1b:eq:natural-nd-3}
\]
It is disjoint from the two-cap product and from the joining locus.
The latter has \(\|v_s\|^2\ge5/6\), so \eqref{p1b:eq:natural-nd-1}--\eqref{p1b:eq:natural-nd-2} give
a uniformly defined closing chart on its whole neighborhood.
\end{proposition}
\begin{proof}
The three orthogonal projections \(ss^*,\ell_s,P_s-\ell_s\)
sum to identity. Formula \eqref{p1b:eq:natural-nd-1} acts on them by \(1,e^{\pi it},
e^{-\pi it}\), proving every rank-two, determinant and endpoint assertion.
Multiplication by the scalar half-phase on the actual rank-two bundle
gives \eqref{p1b:eq:natural-nd-2}, while leaving the specified stabilizing line fixed.
Since \(e_2\) is fixed by \(\mathsf J\), the projections transform
covariantly and the displayed descent is exact. The root commutator
is identity when either of its two cap gauges is identity, hence off \(B\).

The singular condition is precisely \(s\in\mathbb Ce_2\).
Using the actual suspended Hopf section, this requires
\(u_{a1}^2+u_{a2}^2=1\) and all other source coordinates zero.
The balanced smash inverse gives \(X=P=(\cos r,\sin r,0)/\sqrt2\)
in the original stereographic chart. Inverting it gives \eqref{p1b:eq:natural-nd-3},
including all points, not just a sampled zero.

On \(x=w,p=-w\), \eqref{p1b:eq:rt-60-2} gives
\[
 |s_2|^2\le |\mathbf u_a|^2
   =\frac{8w_0^2(1-w_0^2)}{(3+w_0^2)^2}\le\frac16.
                                                        \tag{ND.4}\label{p1b:eq:natural-nd-4}
\]
The maximum occurs at \(w_0^2=3/7\). Thus \(\|v_s\|^2
=1-|s_2|^2\ge5/6\), proving all chart assertions.
\end{proof}

\begin{proposition}[Measured joining origin and the applied odd toggle]
\label{p1b:proposition:natural-determinant-toggle}
In the \emph{actual} ordinary frames \(Q_Y(a,w)=T_wQ_Z(a)\),
the joining value of \eqref{p1b:eq:natural-nd-2} has component \((1,0,0)\).
Thus this prescribed local closing origin has \(e_h=0\).
For the actual retained odd clutch \(c_j\), the required matching
exponent is \(e_h+1=1\). Apply the already constructed global
\(T_Y\) once, and retain the actual joining value
\[
 D_{J}(a,w)=D_+(a,w,-w)T_Y(a,w,-w).                     \tag{ND.5}\label{p1b:eq:natural-nd-5}
\]
It has component \((1,0,1)\), exactly the component of
\(c_j\operatorname{diag}(1,-a^2,1)c_j^*\).
\end{proposition}
\begin{proof}
Here is the complete line comparison, retaining the actual frame.
Write \(b=T_we_2\), \(\widehat v=v_s/\|v_s\|\),
and \(c=e_1^*s\). The actual transporter gives
\[
 e_2^*T_we_2=d\left(1-\frac{|s_2|^2}{1+\overline c}\right).
                                                        \tag{ND.6}\label{p1b:eq:natural-nd-6}
\]
The retained bound \(|1+c|\ge2/9\) and \eqref{p1b:eq:natural-nd-4} show that the
fraction has norm at most \(3/4\); this overlap never vanishes.
Set \(\gamma=\overline{\langle\widehat v,b\rangle}/
|\langle\widehat v,b\rangle|\), and define
\[
 r_\tau=\frac{(1-\tau)\widehat v+\tau\gamma b}
                  {\|(1-\tau)\widehat v+\tau\gamma b\|}.
                                                        \tag{ND.7}\label{p1b:eq:natural-nd-7}
\]
The aligned overlap is positive, so the denominator is nonzero.
This is a continuous equivariant line homotopy on the whole joining
locus, including its deck identification. Replace \(\ell_s\) by
\(r_\tau r_\tau^*\) in \eqref{p1b:eq:natural-nd-2}. At its end the ordinary frame
value is precisely \(\operatorname{diag}(a^2,1,1)\), independent
of the ordinary sphere coordinate. It has determinant-circle
component one and zero sphere and joining-secondary components.
This measures \(e_h=0\) without selecting an unmeasured global
closing origin. The explicit \(T_Y\) has joining component
\((0,0,1)\); multiplying it immediately gives \eqref{p1b:eq:natural-nd-5}.
The actual clutch has component \((0,1,0)\), and its conjugation
of the split determinant endpoint gives \((1,0,1)\), proving matching.
\end{proof}

\begin{proposition}[Constructed rank-two joining with the actual clutch]
\label{p1b:proposition:natural-determinant-joining}
There is a constructed ordinary rank-two joining homotopy with endpoints
\[
 H_J(0,a,w)=D_J(a,w),\qquad
 H_J(1,a,w)=c_j(a,w)\operatorname{diag}(1,-a^2,1)c_j(a,w)^*.
                                                        \tag{ND.8}\label{p1b:eq:natural-nd-8}
\]
It retains the actual \(Q_Y,Q_Z\), \(f_j(a^{-1}w)\), and their
deck descent. Its cylinder restriction uses exactly \(a=i e^{is}\),
\(w=yu\) in the retained boundary map. No inner \(Y\)-stage
restriction is supplied by this joining construction alone.
\end{proposition}
\begin{proof}
Use three consecutive comparisons. First apply \eqref{p1b:eq:natural-nd-7} to \(D_+\)
while retaining the right factor \(T_Y\). Next rotate the positive-line
phase from \(a^2\) to \(-a^2\). Their explicit ordinary endpoints
are \(q_1 T_J\), where
\[
 q_1=\operatorname{diag}(-a^2,1,1),\qquad
 T_J=Q_Y^*T_YQ_Y.
\]
The third comparison is an ordinary \(SU(2)\) homotopy from \(T_J\)
to the \emph{actual} difference
\[
 \tau_J=q_1^{-1}\operatorname{diag}(f_j,1)q_1
                                    \operatorname{diag}(f_j^{-1},1).
                                                        \tag{ND.9}\label{p1b:eq:natural-nd-9}
\]
Here are the lower comparison and construction of that homotopy.
Use the product coordinates \(\xi=a^2,u=a^{-1}w\), already retained.
Put \(q_0=\operatorname{diag}(-1,1)\),
\(D(\xi)=\operatorname{diag}(\xi,1)\), and \(f'=q_0^{-1}f_jq_0\).
The circle-independent value at \(\xi=1\) is
\(\tau_0=f'f_j^{-1}\), with the explicit nullhomotopy obtained
by rotating the constant positive-line phase from \(1\) to \(-1\).
Dividing by this actual lower value gives the exact based equality
\[
 \tau_J\tau_0^{-1}=D(\xi)^{-1}f'D(\xi)f'^{-1}.         \tag{ND.10}\label{p1b:eq:natural-nd-10}
\]
It is a based \(S^4=S^1\wedge S^3\) map. The circle degree is
\(-1\), the localized actual cap has degree \(+1\), and the fixed
inner rotation does not change it. Hence it is the actual adjoint
conjugation representative of \(\eta_3\), with the retained composition
order. This is a named joining input, not a new top calibration.

For the other endpoint, the retained explicit toggle is supported
in \(1/4<t<3/4\) and the first cap in a cut circle. On the joining
locus its first support has \(a\) near \(i\), \(w\) near \(1\);
its deck mate is the same supported four-cell in the quotient.
In the actual ordinary frame its value is
\(D_a^*\eta_3(4(t-1/2),\log w/\delta)D_a\), with
\(D_a=\operatorname{diag}(1,a)\). Contract \(a\) to \(i\) inside
this support; the map is identity on the support boundary, so this
is a genuine rank-two frame homotopy. The source coordinates
\((t,\log w)\mapsto(\xi,u)\) are the retained invertible four-cell
chart. Thus \(T_J\) is a degree-one supported suspension-Hopf
representative of the same \(\eta_3\).

Construct the comparison by the Pontryagin tube collapse on these two
\emph{actual} maps after the stated frame and lower homotopies.
The suspension-Hopf regular fibre is one Hopf circle with its three-normal
framing; the conjugation regular fibre is \(\xi=-1\),
\(f'_0=f'_1=0\), also one circle. Each is an unknotted circle
in its indicated supported four-cell, with the nontrivial class in
\(\pi_1SO(3)=C_2\). Move the four-cell charts through the ordinary
product \(S^1\times S^3\), extend the circle isotopy, and use the
framing homotopy in this same class. Collapse the complements and undo
the collapses, exactly as in the complete comparison used for \eqref{p1b:eq:natural-ne-8}.
This supplies an ordinary rank-two homotopy of the full four-sphere maps.
Undo the explicit lower normalization \(\tau_0\) to end at \eqref{p1b:eq:natural-nd-9}.
Every step is in the descended ordinary joining frame and hence fixes
the specified stabilizing line. Multiplication by \(q_1\) and transport
by \(Q_Y\) gives \eqref{p1b:eq:natural-nd-8}. The equality
\(c_j=Q_Y\operatorname{diag}(f_j,1)Q_Z^*\) verifies its exact right
endpoint: the ambient split matrix is \(\operatorname{diag}(1,-a^2,1)\),
whereas in the ordered \(Q_Z=[e_2,a e_3,e_1]\) frame it is the
\(q_1\) above. Thus the stabilizing line and both joining axes are retained.
\end{proof}

\paragraph{The prescribed global determinant map.}
The full closing patch, constructed by \eqref{p1b:eq:natural-ny-1}--\eqref{p1b:eq:natural-oa-4}, retains
the prescribed joining origin and the exact equation
\[
 \begin{aligned}
 \mathcal K:&[0,1]\times P\longrightarrow SU(3),\\
 \mathcal K(t,x,p)s(x,p)&=s(x,p),\\
 \mathcal K(0,x,p)&=H(x,p)^6,\qquad
 \mathcal K(1,x,p)=N(x,p),\\
 \mathcal K(t,w,-w)&=K_+(t,w,-w).
 \end{aligned}                                             \tag{ND.11}\label{p1b:eq:natural-nd-11}
\]
It must be a rank-two determinant-one gauge for every parameter. Then
\(F_h=\mathcal K H^{-3}\) has exactly the retained endpoints
\(H^3,NH^{-3}\), agrees with \eqref{p1b:eq:natural-nd-1} on the joining locus, and its
toggle and joining are the already constructed \eqref{p1b:eq:natural-nd-5}--\eqref{p1b:eq:natural-nd-8}.
Equality of the global endpoint components does not evaluate this map.

There is no obstruction to prescribing this joining origin at the
homotopy-class level: comparing any existing closing path on the
antipodal \(S^3\) with \eqref{p1b:eq:natural-nd-1} gives an element of
\(\pi_1\operatorname{Map}(S^3,SU(2))=C_2\); the retained
\(\alpha_x\) restriction is onto this group. Section homotopy
extension then fixes that chosen origin. This proves that the local
and joining data are legitimate target data, without claiming to have
computed the needed global patch in \eqref{p1b:eq:natural-nd-11}.

The local formula cannot simply be patched on the two-cap product
alone: after contracting its actual frame, that boundary-cylinder
obstruction is \(2[H^3]=6\nu'\ne0\). The compensating central-path
defect is at the actual circle \eqref{p1b:eq:natural-nd-3}. A global comparison must retain
both contributions; this is a failure of that separate-patch ansatz,
not an impossibility claim about the target. The complete compensation, inner boundary comparison and evaluation
are retained in \eqref{p1b:eq:natural-ny-1}--\eqref{p1b:eq:natural-oa-4}, \eqref{p1b:eq:natural-ob-1}--\eqref{p1b:eq:natural-ob-20} and \eqref{p1b:eq:natural-ci-8},
\eqref{p1b:eq:natural-aj-4}--\eqref{p1b:eq:natural-aj-5}, with \(h_t=w_t=1\) and \(\zeta_{\rm nat}=0\).



\subsection{The complete branch tube of the actual closing comparison}
\label{p1b:sec:natural-closing-branch-tube}

The retained complete triple and the evaluated clutch filling reduce the
actual sphere bit to \eqref{p1b:eq:natural-nc-31}. This subsection constructs its full antipodal
locus, its normal comparison, and its actual preimages under the retained
cap and boundary maps. It does not evaluate the parity of the whole map.

\begin{lemma}[An exact projective parameterization of the actual map]
\label{p1b:lemma:natural-closing-projective}
Write \(A=XP\), \(B=X^{-1}AX\), and retain the coordinate change
\(P=X^{-1}A\) explicitly. Then
\[
 \widehat L(s,A,X)=L(s,X,X^{-1}A),\qquad
 \widehat L(s,A,-X)=\widehat L(s,A,X).                 \tag{NC.33}\label{p1b:eq:natural-nc-33}
\]
Thus this particular parameterization descends in \(X\) to
\(\mathbb {RP}^3\). The original comparison is still exactly
\(L(s,X,P)=\widehat L(s,XP,X)\); the original based collapse and
the boundary chart are not replaced by an unproved top-cell identification.
\end{lemma}
\begin{proof}
Both conjugate quaternions in \eqref{p1b:eq:natural-nc-31} depend on \(X\) only through
\(X^{-1}AX\). Changing the sign of \(X\) cancels there. Their
specified powers and the prescribed midpoint value are therefore unchanged.
The coordinate shear is an actual invertible coordinate change on the
product, with the displayed inverse. A degree-two cover of a projective
space does not alone evaluate the natural bit: no lower-cell term has been
removed from this projective parameterization.
\end{proof}

\begin{theorem}[The entire antipodal locus and its actual normal map]
\label{p1b:theorem:natural-closing-full-branch}
For the original actual \(L\), every point with \(L=-1\) is
uniquely described by
\[
 s=\pi\pmod{2\pi},\qquad
 A=-\tfrac12+\tfrac{\sqrt3}{2}u,\qquad
 u\cdot X^{-1}uX=\tfrac13,\qquad P=X^{-1}A.           \tag{NC.34}\label{p1b:eq:natural-nc-34}
\]
Here \(u\) is a unit imaginary quaternion. Its original-source locus is
the compact smooth manifold \(M=S^1\times SO(3)\), parameterized by
\[
 X=\sqrt{\tfrac23}(\cos t+\sin t\,u)
                  +\sqrt{\tfrac13}w,\qquad
 |u|=|w|=1,\quad u\cdot w=0.                          \tag{NC.35}\label{p1b:eq:natural-nc-35}
\]
The sign quotient in \eqref{p1b:eq:natural-nc-33} is not taken in this description of \(M\).
There is a global normal chart obtained by replacing \(s=\pi\),
\(\phi=2\pi/3\), \(d=1/3\) with nearby values in
\[
 A=\cos\phi+\sin\phi\,u,\quad
 X=\sqrt{\tfrac{1+d}{2}}(\cos t+\sin t\,u)
                 +\sqrt{\tfrac{1-d}{2}}w,\quad P=X^{-1}A. \tag{NC.36}\label{p1b:eq:natural-nc-36}
\]
Put \(w'=\sin t\,w+\cos t\,(u\times w)\), so that
\(X^{-1}uX=d u+\sqrt{1-d^2}w'\). In the output basis
\((u,w',u\times w')\), the actual normal differential is
\[
 \left.\frac{\partial\operatorname {Im}L}
                  {\partial(s,\phi,d)}\right|_M
 =\begin{pmatrix}
 8/9&4/3&-\sqrt3/2\\
 4\sqrt2/9&-4\sqrt2/3&\sqrt6/2\\
 0&2\sqrt6/3&3\sqrt2/2
 \end{pmatrix},\qquad \det=-8.                       \tag{NC.37}\label{p1b:eq:natural-nc-37}
\]
Thus \(-1\) is a regular value of this actual input, with its normal
comparison explicitly retained, rather than merely one sampled antipode.
\end{theorem}
\begin{proof}
The scalar branches \(A=\pm1\) or \(B=\pm1\) give \(L=1\)
and cannot contribute. Otherwise write
\(A=c+\sqrt{1-c^2}u\), \(B=c+\sqrt{1-c^2}v\),
\(d=u\cdot v\). Conjugation keeps the real part of the closing family
fixed, equal to \(c^2-(1-c^2)d\). If \(L=-1\), its value
must be the negative of the prescribed midpoint value, so necessarily
\[
 c^2-(1-c^2)d=0.                                    \tag{NC.38}\label{p1b:eq:natural-nc-38}
\]
In the first half of the circle the displacement from the midpoint
is rotation about \(u\) by an angle of magnitude at most \(\phi<\pi\).
It cannot send a nonzero vector to its negative. Thus the antipode lies
in the second half. The midpoint formula, obtained by direct multiplication,
is
\[
 F(\pi/2)=c^2-(1-c^2)d+
     \sqrt{1-c^2}\{v+[c(1+d)-d]u\}.                  \tag{NC.39}\label{p1b:eq:natural-nc-39}
\]
The projection of the second half onto its rotation axis \(v\) is
\(c\sqrt{1-c^2}(1+d)\). Matching it to the negative projection
of \eqref{p1b:eq:natural-nc-39} gives
\[
 (1+d)\{c(1+d)+1-d\}=0.                             \tag{NC.40}\label{p1b:eq:natural-nc-40}
\]
The case \(d=-1\) has \(AB=1\) and fails \eqref{p1b:eq:natural-nc-38}.
Combining the other case with \eqref{p1b:eq:natural-nc-38} yields \(c=2c^2-1\).
Its scalar solution \(c=1\) is excluded; hence \(c=-1/2\)
and \(d=1/3\). At those values direct multiplication gives
\(F(3\pi/2)=-F(\pi/2)\). The second arc rotates by total angle
\(4\pi/3<2\pi\), so this occurrence is unique. With the retained
baseline, its circle coordinate is exactly \(s=\pi\).
This proves necessity and sufficiency of \eqref{p1b:eq:natural-nc-34}, not just existence.

Write \(X=\alpha+\beta u+q_\perp\). Quaternion conjugation gives
\(u\cdot X^{-1}uX=1-2|q_\perp|^2\). On \eqref{p1b:eq:natural-nc-34},
\(|q_\perp|^2=1/3\), \(\alpha^2+\beta^2=2/3\).
Thus \(\alpha,\beta\) determine a unique \(t\) modulo \(2\pi\),
and \(q_\perp\) determines a unique unit tangent vector \(w\).
The oriented pair \((u,w)\) is exactly a point of \(SO(3)\).
This proves \eqref{p1b:eq:natural-nc-35} and the product description of \(M\).
The same calculation gives \eqref{p1b:eq:natural-nc-36}, a genuine global normal chart.

For \eqref{p1b:eq:natural-nc-37} use quaternion covariance to take \(u=i\), \(w'=j\).
Differentiate the exact second-half word
\(B^{1-\tau}AB^\tau\), \(\tau=s/\pi-1/2\), and divide
by \(A^{1/2}BA^{1/2}\). At
\((s,\phi,d)=(\pi,2\pi/3,1/3)\) the real derivative is zero;
the three imaginary derivatives are exactly the displayed columns.
Their determinant is \(-8\). Covariance gives the asserted differential
at every \((t,u,w)\), including its actual rotating output basis.
No parity of a mapping cone follows merely from this regular-value data.
\end{proof}

\begin{proposition}[Actual cap and boundary preimages of the full branch]
\label{p1b:proposition:natural-closing-branch-preimages}
All branch points \eqref{p1b:eq:natural-nc-34} have explicit preimages under the retained
two degree-one cap maps and the actual boundary restriction \eqref{p1b:eq:natural-nc-4}.
For a target quaternion \(Q\ne1\), put
\[
 k=\frac{|\operatorname {Im}Q|}{1-\operatorname {Re}Q},\qquad
 r_Q=\frac{2k\delta^2}{1+\sqrt{1+4k^2\delta^2}},\qquad
 b(Q)=\exp\left(-r_Q\frac{\operatorname {Im}Q}
                                  {|\operatorname {Im}Q|}\right). \tag{NC.41}\label{p1b:eq:natural-nc-41}
\]
Use \(b(-1)=1\). For \eqref{p1b:eq:natural-nc-34}--\eqref{p1b:eq:natural-nc-36}, set
\[
\begin{gathered}
 x=b(X),\qquad p=b(P_0,P_1,-P_2,-P_3),\qquad v_0=\log(p^{-1}x),\qquad y=x e^{-v_0/2},\\
 \rho=|v_0|/(3\pi), \\
 u_0=v_0/|v_0|.
\end{gathered}
\tag{NC.42}\label{p1b:eq:natural-nc-42}
\]
At \(v_0=0\) use the actual center of the normal disk. Then
\(\Phi_S(s,y,\rho,u_0)=-I\) at every branch point, in the
unrebased convention. For the retained basepoint use \(y_{\rm ret}=-y\)
and precisely the retained isotopy of \eqref{p1b:eq:natural-nc-15}.
The actual target normal differential is \eqref{p1b:eq:natural-nc-37} composed with
the two cap and pair-chart differentials and conjugated by
\(U=F_Y^*Q_s\). The homotopy contracting this actual \(U\)
is the one already proved in Theorem~\ref{p1b:theorem:natural-closing-input}.
\end{proposition}
\begin{proof}
The actual rational cap has
\(Q_0=(r^2-g^2)/(r^2+g^2)\),
\(\operatorname {Im}Q=-2g\,v/(r^2+g^2)\),
where \(g=\delta^2-r^2\) and \(r=|v|\).
Thus \(k=r/(\delta^2-r^2)\); solving the positive quadratic
gives \eqref{p1b:eq:natural-nc-41}, including the required negative direction.
The rotation on the second cap is its own inverse, which gives
the second formula of \eqref{p1b:eq:natural-nc-42}. The branch cannot have \(X=\pm1\)
or \(P=\pm1\), since either would force \(A=B\) and \(L=1\).
These inverses are therefore continuous on the full compact branch.
Both framework points are inside their \(+1\) caps and their
pair chart avoids the antipode. Equation \eqref{p1b:eq:natural-nc-42} is exactly the inverse
of \eqref{p1b:eq:natural-nc-4}; no base or normal coordinate is changed.

The odd clutch support misses this product of caps. Here the complete
actual sphere comparison is therefore \(U L U^{-1}\). At
\(L=-1\) that is exactly \(-I\). Its normal derivative is
\(\operatorname {Ad}_U DL\), since all derivatives of the frame
cancel at a scalar value. This retains the actual bundle identification
and its already constructed homotopy, as well as the source differentials.
The coordinate implementation applies \eqref{p1b:eq:natural-nc-41}--\eqref{p1b:eq:natural-nc-42} immediately and
checks the value on the complete original comparison, not on a benchmark.
\end{proof}

\paragraph{Evaluation after the branch construction.}
The global framed comparison of this entire branch, together with the
actual equal-trace factorization and all lower-cell cancellations, is now
constructed in Theorem~\ref{p1b:theorem:natural-sphere-bit}. It evaluates
\[
 \mathcal F_L:D^8\longrightarrow SU(3),\qquad
 \mathcal F_L|_{S^7}=\operatorname{diag}(L,1),\qquad
 \operatorname{Sq}^4u_5=v_9
       \text{ on }C_{[\mathcal F_Le_3]}.             \tag{NC.43}\label{p1b:eq:natural-nc-43}
\]
The construction uses the proved full-map identification and a constant
last-column outer collar; it does not infer parity from this branch alone.
The complete original comparison is evaluated in
\eqref{p1b:eq:natural-aj-4}--\eqref{p1b:eq:natural-aj-7}:
\(h_t=w_t=1\), \(r_J=\zeta_{\rm nat}=0\), with every retained
extension test decided. This component is used with its full frames,
lower divisions and joining; its isolated value is not a replacement
for that complete comparison.


\subsection{Evaluation of the retained natural sphere comparison}
\label{p1b:sec:natural-sphere-evaluation}

Retain the actual triple, boundary chart, cap maps, bundle identifications
and frame homotopies of \eqref{p1b:eq:natural-nc-1}--\eqref{p1b:eq:natural-nc-31}, and the complete branch calculation
\eqref{p1b:eq:natural-nc-33}--\eqref{p1b:eq:natural-nc-42}. The already filled odd-clutch comparison has value zero.
The remaining actual map is the prescribed \(L\) in \eqref{p1b:eq:natural-nc-31}.
The following proof evaluates that map, including its lower-cell terms.

\paragraph{Executed dependencies.}
Factor this specific \(L\) through the equal-trace pair quotient; construct
a degree-one top comparison and compare its \emph{entire} framed inverse
image with the quaternion commutator; compute the actual source map's top
component; remove the lower component with the retained composition law;
apply the retained stabilized parity detector to this identified input.
Each construction and comparison is made below.

\begin{lemma}[The actual factorization and source shear]
\label{p1b:lemma:natural-sphere-factor}
Let \(\Delta\subset S^2\times S^2\) be the diagonal and put
\[
 Q=(S^2\times S^2)/\Delta,\qquad Y=\Sigma Q,\qquad E=\Sigma Y.
                                                        \tag{NE.1}\label{p1b:eq:natural-ne-1}
\]
For \(A=\cos\phi+\sin\phi\,u\), \(0\le\phi\le\pi\),
define the actual based map
\[
 k_1(A,X)=[\phi;u,X^{-1}uX]\in Y,\qquad
 k=\Sigma k_1:S^7\longrightarrow E.                     \tag{NE.2}\label{p1b:eq:natural-ne-2}
\]
Both polar endpoints and the diagonal pair are the basepoint. Thus
\(k_1\) is continuous, is constant when \(A=1\) or \(X=1\),
and descends from \(S^3\times S^3\) to \(S^6\).

Define \(G:E\to SU(2)\) by the \emph{retained} closing word, using
\(A=\cos\phi+\sin\phi\,u\),
\(B=\cos\phi+\sin\phi\,v\), and precisely the retained midpoint:
\[
 G[s,\phi;u,v]=F_{A,B}((\pi/2+s)\bmod2\pi)
                         F_{A,B}(\pi/2)^{-1}.            \tag{NE.3}\label{p1b:eq:natural-ne-3}
\]
This is a map on \(E\): it is identity on \(s=0\), on
\(\phi=0,\pi\), and on \(u=v\). The scalar branches are the
removable branches already proved in the actual closing construction.
Then, as \emph{actual functions},
\[
 \widehat L(s,A,X)=Gk(s,A,X),\qquad
 L(s,X,P)=Gk(s,XP,X).                                   \tag{NE.4}\label{p1b:eq:natural-ne-4}
\]
The last source shear preserves the extracted top class up to sign.
It introduces no unremoved lower comparison.
\end{lemma}
\begin{proof}
In the retained map \(B=PX=X^{-1}(XP)X\), so \eqref{p1b:eq:natural-ne-4} follows by
substitution, with no change of closing equation or circle baseline.
The map in \eqref{p1b:eq:natural-ne-3} is constant for commuting equal-trace pairs, in
particular on the diagonal. Its endpoint behavior proves continuity
on the displayed quotients.

The shear \(T(X,P)=(XP,X)\) is an actual based diffeomorphism of
the product, although it is not asserted to preserve its coordinate
wedge. Use the suspended product quotient instead: for a top section
\(j:S^7\to\Sigma(S^3\times S^3)\), the composite
\(\Sigma q\,\Sigma T\,j:S^7\to S^7\) has degree \(-1\).
Indeed \(T^*a=a_x+a_p\), \(T^*b=a_x\), and hence
\(T^*(ab)=-a_xa_p\). All maps from the lower \(S^4\) summands
to \(S^7\) are null. Both actual closing comparisons restrict
to the identity on their original coordinate wedges. Consequently
their extracted top classes differ only by that degree sign; choices
of the suspended product section add lower maps on which the actual
comparison is null. This justifies the source shear without declaring
it a map of the original smash quotients. The sign has no effect in
\(\pi_7SU(2)=C_2\).
\end{proof}

\begin{lemma}[The two-cell space and every splitting ambiguity]
\label{p1b:lemma:natural-sphere-two-cell}
The quotient \(Q\) has the two-cell homotopy type
\(S^2\cup_{\pm2\eta_2}e^4\). Thus
\[
 Y\simeq S^3\vee S^5,\qquad E\simeq S^4\vee S^6,       \tag{NE.5}\label{p1b:eq:natural-ne-5}
\]
where the second splitting can be chosen as a suspension of the first.
Write \(j_4:S^4\to E\) for the bottom inclusion and choose a top
section \(h_5:S^5\to Y\) using any based nullhomotopy of the
actual attaching map \(\pm2\eta_3\); retain its suspension \(Eh_5\).
The homotopy exists since \(2\eta_3=0\). The calculation below
accounts for every choice of this homotopy, rather than assigning
zero to a choice-dependent lower term.

Every degree-one top map \(H:S^6\to E\) differs from
\(Eh_5\), up to a common orientation sign, by
\(j_4\epsilon\), where
\(\epsilon\in\pi_6S^4=C_2\{\eta_4\eta_5\}\).
Postcomposition with \(G\) changes \([GH]\) by either zero
or \(6\nu'\). Its coefficient of \(\nu'\) modulo two is
therefore independent of the splitting.
\end{lemma}
\begin{proof}
The complement of the diagonal has the actual tangent-bundle chart
\[
 (u,v)\longmapsto\left(u,
             \frac{v-(u\cdot v)u}{1-u\cdot v}\right)\in TS^2.
                                                        \tag{NE.6}\label{p1b:eq:natural-ne-6}
\]
It is a diffeomorphism with inverse given by inverse stereographic
projection in each tangent plane. Approaching the removed diagonal
is going to infinity. Hence \(Q\) is the one-point compactification
of \(TS^2\), its Thom space. Its bottom and top cells have dimensions
two and four, and the top attaching map has Hopf invariant of absolute
value two, the Euler number of \(TS^2\). This proves the claimed
attaching map. Suspending gives \(\pm2\eta_3=0\), proving \eqref{p1b:eq:natural-ne-5}.
A nullhomotopy of this actual attaching map constructs the top section
by the usual mapping-cone extension; suspending that construction
retains a suspension splitting.

There are no mixed Whitehead terms in \(\pi_6(S^4\vee S^6)\):
their first possible dimension is nine. Its degree-one top lifts
therefore differ by precisely the stated bottom group. The bottom
map \(Gj_4\) is either zero or \(\eta_3\). The retained composition
law \(\eta_3\eta_4\eta_5=6\nu'\) proves the ambiguity statement.
No particular nullhomotopy coefficient needs to be guessed.
\end{proof}

\begin{theorem}[A complete framed comparison of the actual top input]
\label{p1b:theorem:natural-sphere-top-comparison}
For independent quaternions
\(a'=\cos\alpha+\sin\alpha\,u\) and
\(b'=\cos\beta+\sin\beta\,v\), construct
\[
 H(a',b')=[2\alpha,\,g(\beta);u,v]\in E,
 \qquad g(\beta)=\beta+\frac\pi6\sin\beta.             \tag{NE.7}\label{p1b:eq:natural-ne-7}
\]
The map is constant on both coordinate wedges and both polar endpoint
sets, so it descends to \(S^6=S^3\wedge S^3\). It has top degree
of absolute value one. For this \emph{actual} comparison map,
\[
 {}[GH]=\pm\nu'\in\pi_6SU(2),                           \tag{NE.8}\label{p1b:eq:natural-ne-8}
\]
where \(\nu'\) is the retained quaternion-commutator generator.
In particular its coefficient is odd. This identifies the whole map,
not a sampled locus, and includes its complementary region.
\end{theorem}
\begin{proof}
The function \(g\) is a diffeomorphism of \([0,\pi]\), since
\(g'=1+(\pi/6)\cos\beta>0\); it has
\(g(\pi/2)=2\pi/3\), \(g'(\pi/2)=1\).
The polar endpoint conditions prove that \eqref{p1b:eq:natural-ne-7} is continuous on
the actual quotients. A generic point of the top open cell of \(E\)
has exactly one preimage, with nonzero Jacobian, proving its top degree.

Apply the \emph{complete} antipodal equations \eqref{p1b:eq:natural-nc-38}--\eqref{p1b:eq:natural-nc-40} to
the independent equal-trace pair in \eqref{p1b:eq:natural-ne-3}. They use only \(A,B\)
and their equal trace, not a chosen conjugating \(X\). Thus
\[
 (GH)^{-1}(-1)=\{\alpha=\beta=\pi/2,\ u\cdot v=1/3\}.
                                                        \tag{NE.9}\label{p1b:eq:natural-ne-9}
\]
There are no other inverse images. Parameterize
\(v=d u+\sqrt{1-d^2}w\), with \((u,w)\) an orthonormal
imaginary pair. In the output frame \(R=(u,w,u\times w)\),
the differential of \(GH\) on the three normal coordinates
\((\alpha,\beta,d)\) is exactly
\[
 R D_*\operatorname{diag}(2,1,1),\qquad\det=-16,       \tag{NE.10}\label{p1b:eq:natural-ne-10}
\]
where \(D_*\) is the \emph{actual} matrix \eqref{p1b:eq:natural-nc-37}.
Its dependence on the entire \(SO(3)\) locus through \(R\) is retained.

Here is a global ambient isotopy of this locus. Choose a smooth cutoff
\(\chi\) equal to one on a neighborhood of \(\pi/2\)
and zero near both polar endpoints. Put \(c=t\chi(\alpha)\chi(\beta)/3\),
\(0\le t\le1\), and replace the imaginary direction of \(b'\) by
\[
 v_t=\frac{(d-c)u+\sqrt{1-c^2}(v-du)}{1-cd}.
                                                        \tag{NE.11}\label{p1b:eq:natural-ne-11}
\]
For each \(u\) this is a smooth sphere diffeomorphism, with
\(d_t=(d-c)/(1-cd)\), fixing its two axial points. Its inverse
uses \(-c\). The two cutoffs make it the identity near the poles
of both \(a'\) and \(b'\), so it extends smoothly there. It fixes both
coordinate wedges and therefore gives an isotopy of the actual
smash source. At \(t=1\), \eqref{p1b:eq:natural-ne-9} is exactly
\(\alpha=\beta=\pi/2,d_t=0\).
The perpendicular unit vector \(w\) is unchanged throughout.
In the retained global normal coordinates, transport multiplies the
third coordinate by a positive constant along the locus; cutoff
derivatives vanish there. Tangential normal-coordinate shears can
be removed continuously and do not change the framing.

The standard commutator \(C(a',b')=a'b'a'^{-1}b'^{-1}\)
has \(-1\) inverse image exactly the orthogonal imaginary pairs,
the locus just reached. Direct differentiation in the same frame
\(R\) and coordinates gives
\[
 DC|_{-1}=R\operatorname{diag}(-2,2,-2),\qquad
 D(C^{-1})|_{-1}=R\operatorname{diag}(2,-2,2).          \tag{NE.12}\label{p1b:eq:natural-ne-12}
\]
The transported normal matrix in \eqref{p1b:eq:natural-ne-10} and the second matrix
in \eqref{p1b:eq:natural-ne-12} lie in the same connected component of \(GL(3,\mathbb R)\).
A path between their \emph{constant} matrices gives a global framing
homotopy \(nD_t^{-1}R^{-1}\) over the entire locus. In particular
the nonconstant frame \(R\) is present at both ends and along the
homotopy; it has not been discarded as a constant frame.

For completeness, the passage from this framing comparison to maps
includes the complementary region. Choose a small target ball about
\(-1\), disjoint from the based value \(+1\), and postcompose each
map with a degree-one radial collapse supported on this ball. Such a
collapse is based homotopic to the identity. Its pullback is the
Pontryagin tube collapse with the induced framing, and its complement
is the based constant map. Smooth approximation can be made relative
to these compact regular tubes, keeping the complementary image away
from \(-1\). The isotopy \eqref{p1b:eq:natural-ne-11} and the global framing homotopy
therefore give a based homotopy of the \emph{full} collapsed maps.
Undoing the target collapse proves \eqref{p1b:eq:natural-ne-8}, in fact the commutator
inverse for compatible source orientations. This is the Pontryagin
construction for actual framed submanifolds (Milnor, \S7, Theorem B).
It accounts for the full inverse image and the full complement.
\end{proof}

\begin{lemma}[The actual source top component]
\label{p1b:lemma:natural-sphere-source-top}
Let \(q_5:Y\to S^5\) be the top-cell quotient. Then
\[
 {}[q_5 k_1]=\eta_5,\qquad
 k\simeq j_4 E(m\nu')+(Eh_5)\eta_6
 \quad\text{for some }m\in\mathbb Z/12.                \tag{NE.13}\label{p1b:eq:natural-ne-13}
\]
The unknown integer is a lower component whose entire contribution
is evaluated zero in the next theorem; it is not a conditional natural bit.
\end{lemma}
\begin{proof}
Represent the quotient on \(Q\) by collapsing the complement of a
small four-ball about one generic pair \((u_0,v_0)\), chosen away
from the diagonal and from a bottom-cell representative. This has
top degree one and is homotopic to the cellular top quotient.
Its suspension represents \(q_5\). Choose a regular interior polar
coordinate \(\phi_0\). Its \emph{entire} inverse image under the
actual \(k_1\) is
\[
 A=\cos\phi_0+\sin\phi_0 u_0,\qquad
 X^{-1}u_0X=v_0.                                      \tag{NE.14}\label{p1b:eq:natural-ne-14}
\]
It is one Hopf circle in the \(X\)-sphere, with all three \(A\)
coordinates fixed. The differential is block triangular: the first
three normals are the fixed \(A\)-chart normals, and the remaining
two are precisely the standard Hopf differential. Removing the
off-diagonal block is a framing homotopy, constant on the added three
normal coordinates. The Pontryagin collapse is therefore the third
suspension of the Hopf map, \(\eta_5\). This proves the first equation
using the actual source circle and its complete framing.

In the chosen \(Y\simeq S^3\vee S^5\) there is no mixed Whitehead
term in dimension six, since its first dimension is seven. Thus
\([k_1]\) has bottom component \(m\nu'\) and top component
\(\eta_5\). Suspending this equality gives \eqref{p1b:eq:natural-ne-13}. All these
homotopies are on the based quotients of the displayed actual maps.
\end{proof}

\begin{theorem}[The original sphere bit is evaluated]
\label{p1b:theorem:natural-sphere-bit}
For the retained natural normalization and actual odd reference triple,
\[
 \boxed{[L]=\nu'\eta_6=\mu,\qquad w_t=1.}             \tag{NE.15}\label{p1b:eq:natural-ne-15}
\]
This assigns one of the original two bits. It does not assign \(h_t\).
\end{theorem}
\begin{proof}
Theorem~\ref{p1b:theorem:natural-sphere-top-comparison} and
Lemma~\ref{p1b:lemma:natural-sphere-two-cell} give
\([G(Eh_5)]=n\nu'\), where \(n\) is odd. By \eqref{p1b:eq:natural-ne-13},
\[
 {}[Gk]=(Gj_4)E(m\nu')+n\nu'\eta_6.                    \tag{NE.16}\label{p1b:eq:natural-ne-16}
\]
The first summand is zero for \emph{every} \(m\): the bottom map is
zero or \(\eta_3\), and the retained, proved law is
\(\eta_3E\nu'=0\), not the different nonzero law
\(\nu'\eta_6=\mu\). Every choice of the splitting nullhomotopy
changes \(n\) by \(0\) or \(6\), so it also contributes zero after
\(\eta_6\). Thus all possible lower-cell and splitting cancellations
are explicitly removed. The source shear of \eqref{p1b:eq:natural-ne-4} preserves this
order-two value. The two actual cap degrees and the actual residual
bundle-frame contractions were retained in \eqref{p1b:eq:natural-nc-31}, so this is the
original closing-map coefficient.

Apply the \emph{existing} stabilized detector to this actual input.
The homotopies just constructed identify it with the retained generator
\(\mu=\nu'\eta_6=\eta_3\nu_4\). Subsection~\ref{p1b:sec:rtproof61} supplies its explicit
filling \(K_r(\nu_4)\), whose projected cone has
\(\operatorname{Sq}^4u_5=v_9\). To form a filling of \(L\), put
that filling on an inner ball and the stabilization of the above based
homotopies on an outer collar. Each collar map is block \(SU(2)\)
and fixes \(e_3\); consequently its last-column projection is constant.
The projected filling therefore has the same parity, one. Choices
of the lower nullhomotopies do not alter it, by \eqref{p1b:eq:natural-ne-16}.
This constructs the detector input from the identified actual map;
no primitive value is substituted without its comparison proof.

Finally the actual full sphere comparison is \(\Phi_0\Psi_C\),
with disjoint supports. The already constructed filling of \(\Psi_C\)
has parity zero, and the contraction of the actual closing frame
identifies \(\Phi_0\) with \(L\). Adding that evaluated clutch term
gives exactly \(w_t=1\), including the prescribed zero matching toggle
and constant joining of the complete triple.
\end{proof}

The complete original comparison is evaluated in
\eqref{p1b:eq:natural-aj-4}--\eqref{p1b:eq:natural-aj-7}:
\(h_t=w_t=1\), \(r_J=\zeta_{\rm nat}=0\), with every retained
extension test decided. This component is used with its full frames,
lower divisions and joining; its isolated value is not a replacement
for that complete comparison.


\subsection{An actual exterior closing and its six-sphere patch input}
\label{p1b:sec:natural-determinant-patch}

Retain the actual section \(s\), root \(H^3\), exchange
\(\mathsf J\), odd clutch, toggle and joining of \eqref{p1b:eq:natural-nc-17}--\eqref{p1b:eq:natural-nc-18}
and \eqref{p1b:eq:natural-nd-1}--\eqref{p1b:eq:natural-nd-11}. The retained I-B source owners and all later comparison
prerequisites are included in the present geometry proof chain. In particular \(w_t=1\) remains proved.
The construction below supplies the exterior part of \eqref{p1b:eq:natural-nd-11}, with
its \emph{exact} prescribed joining origin, and computes the class of
its actual remaining patch input. It does not assign \(h_t\).

\paragraph{Applied dependencies.}
Move the actual singular circle into the commutator cap by an explicit
based sphere isotopy; transport the actual bundle along that isotopy;
form a global exterior positive line agreeing with the original line
on the joining locus; construct the full cylinder-boundary map in the
actual cap frame; calculate its two contributions and cancel them.
All five operations are applied below.

\begin{lemma}[An actual shifted section and actual bundle transport]
\label{p1b:lemma:natural-determinant-shift}
In stereographic coordinates \(X=\mathbf x/(1-x_0)\), let
\(\phi_\mu\) be the smooth based sphere diffeomorphism
\(X\mapsto X-\mu e_i\), \(0\le\mu\le32\). It fixes the pole
\(+1\), and the two-factor isotopy commutes with exchange. Put
\[
 s_\mu(x,p)=s(\phi_{-\mu}x,\phi_{-\mu}p).                 \tag{NT.1}\label{p1b:eq:natural-nt-1}
\]
Use intervals of length \(1/16\) in \(\mu\). Multiply the exact
transporters \(T(s_{\mu_k},s_{\mu_{k+1}})\) from \eqref{p1b:eq:rt-58-5},
with a final partial interval, to define \(U_\mu\). Then
\[
 U_\mu\in SU(3),\quad U_\mu s=s_\mu,\quad U_0=I,
 \quad U_\mu(p,x)=\mathsf J U_\mu(x,p)\mathsf J.          \tag{NT.2}\label{p1b:eq:natural-nt-2}
\]
Every transporter chart is valid uniformly. These are actual products,
not an unspecified parallel-transport choice or numerical integration.
\end{lemma}
\begin{proof}
Stereographic translation has sphere speed at most two. The two-factor
speed is at most \(2\sqrt2\); the retained section is three-Lipschitz.
The section distance on one step is therefore at most
\(6\sqrt2/16<1\). This excludes the antipodal chart failure.
The final partial transporter tends to identity at each subdivision
point, proving continuity in \(\mu\). The transporter sends its first
section to its second, giving the column equation. Its covariance gives
the exchange equation. At the based wedge all sections are \(e_1\),
so every product there is identity.
\end{proof}

Let \(\mathbb B\) be the round six-ball in the actual pair chart
\[
 x=\exp z_x,\quad p=\exp z_p,\qquad
 |z_x|^2+|z_p|^2\le R_B^2,\qquad R_B=1/6.               \tag{NT.3}\label{p1b:eq:natural-nt-3}
\]
It contains the entire two-cap commutator support, whose radius is
at most \(\sqrt2/12<R_B\). Thus \(H=I\) outside \(\mathbb B\).
Write \(d(x,p)\) for quaternion sphere distance. Choose the following
specific continuous cutoff:
\[
 \mu(x,p)=
 \begin{cases}
 32,&d\le1/3,\\
 32(1-3v^2+2v^3),\ v=3d-1,&1/3\le d\le2/3,\\
 0,&d\ge2/3.
 \end{cases}                                           \tag{NT.4}\label{p1b:eq:natural-nt-4}
\]
It is constant on a neighborhood of the diagonal and on the whole
joining locus. It is also 32 on \(\mathbb B\), since
\(d(x,p)\le|z_x|+|z_p|\le\sqrt2/6<1/3\) there.

\begin{proposition}[Complete exterior closing with retained joining]
\label{p1b:proposition:natural-determinant-exterior}
Put \(v_\mu=(I-s_\mu s_\mu^*)e_2\). Where it is nonzero let
\[
 \mathsf V_\mu
  =U_\mu^*[\widehat v_\mu, s_\mu\mathbin\times\widehat v_\mu, s_\mu],
 \qquad \widehat v_\mu=v_\mu/\|v_\mu\|,                 \tag{NT.5}\label{p1b:eq:natural-nt-5}
\]
using the complex cross product fixed by the ambient \(SU(3)\) volume.
Set \(\mathsf V^\#=\mathsf V_{\mu(x,p)}\), and let
\(\ell^\#\) be the projection onto its first column. Its entire
singular set is
\[
 x=p=\phi_{32}y(r),\qquad
 y(r)=(-1/3,(2\sqrt2/3)\cos r,(2\sqrt2/3)\sin r,0).    \tag{NT.6}\label{p1b:eq:natural-nt-6}
\]
This set lies strictly inside \(\mathbb B\). Consequently
\[
 K^\#(t,x,p)=\mathsf V^\#
       \operatorname{diag}(e^{\pi it},e^{-\pi it},1)
          (\mathsf V^\#)^*                             \tag{NT.7}\label{p1b:eq:natural-nt-7}
\]
is an actual, everywhere defined exterior rank-two path from \(I\)
to \(N\), fixes \(s\), and satisfies
\[
 K^\#(t,w,-w)=K_+(t,w,-w).                             \tag{NT.8}\label{p1b:eq:natural-nt-8}
\]
Its descended determinant gauge is
\(D_Y^\#=ss^*+a^2\ell^\#+P_s-\ell^\#\).
Thus the measured origin \(e_h=0\), applied \(T_Y^1\), and the
actual joining \eqref{p1b:eq:natural-nd-8} are retained exactly.
\end{proposition}
\begin{proof}
The complex cross product satisfies
\([s_\mu,\widehat v_\mu,s_\mu\times\widehat v_\mu]\in SU(3)\).
Cyclically permuting its three columns and multiplying by \(U_\mu^*\)
give an actual \(SU(3)\) frame with last column \(s\).

By the complete singular equation \eqref{p1b:eq:natural-nd-3}, \(v_\mu=0\) precisely
when \(x=p=\phi_\mu y(r)\). At such a pair \eqref{p1b:eq:natural-nt-4} sets
\(\mu=32\), proving the entire set \eqref{p1b:eq:natural-nt-6}, including the variable
cutoff. The old circle has stereographic radius \(1/\sqrt2\).
After translation its radius is at least \(32-1/\sqrt2\); hence
the round pair radius is at most
\[
 2\sqrt2\arctan\frac1{32-1/\sqrt2}<1/6.                \tag{NT.9}\label{p1b:eq:natural-nt-9}
\]
This proves the required strict containment. Equation \eqref{p1b:eq:natural-nt-7} acts
on the two rank-two lines by inverse phases and on \(s\) by identity.
It therefore has determinant one and endpoints \(I,N\).
On the joining locus \(d(w,-w)=\pi\), so \(\mu=0\) and \(U_0=I\);
its first projection is exactly the old \(\ell_s\). This proves
\eqref{p1b:eq:natural-nt-8}, without changing joining coordinates or their positive axis.
Exchange fixes the first projection covariantly; the negative cross
column changes sign, which does not change its projection. Thus
the descended determinant gauge has the exact required deck law.
\end{proof}

\begin{theorem}[The constructed actual cap-boundary comparison]
\label{p1b:theorem:natural-determinant-patch-input}
The actual frame
\(Q_s=T(e_1,s)[e_2,e_3,e_1]\) is defined on the whole ball \eqref{p1b:eq:natural-nt-3}.
In its rank-two block put \(g=Q_s^*\mathsf V^\#|_{\partial\mathbb B}\)
and \(\Pi=g\operatorname{diag}(1,0)g^*\). The following is the
\emph{complete} actual map on the boundary of the seven-dimensional
cylinder, which is a six-sphere:
\[
 \Omega_B(t,z)=
 \begin{cases}
  (Q_s^*H^6Q_s)_{2\times2},&t=0,\\
  -I_2,&t=1,\\
  \cos(\pi t)I_2+i\sin(\pi t)(2\Pi(z)-I_2),
                                      &|z|=R_B.
 \end{cases}                                           \tag{NT.10}\label{p1b:eq:natural-nt-10}
\]
All three faces and their seams are retained. Its actual class is
\[
 \boxed{[\Omega_B]=6\nu'+\eta_3\eta_4\eta_5=0
                 \quad\text{in }\pi_6SU(2).}           \tag{NT.11}\label{p1b:eq:natural-nt-11}
\]
This is the cap-extension obstruction, not the natural bit \(h_t\).
\end{theorem}
\begin{proof}
Comparing \((x,p)\) with \((1,p)\) gives
\(\|s-e_1\|\le3|z_x|\le1/2\), so
\(\operatorname{Re}(e_1^*s)\ge7/8\). The cap frame is therefore
valid everywhere on this larger round ball. The bottom map is identity
on its boundary because the commutator support lies strictly inside.
The side endpoints are \(I_2,-I_2\), proving both seam equalities.

First retain the entire bottom contribution. In \(Q_s\) the two actual
cap quaternions are \(q_\delta(x)\) and the fixed degree-one rotation
of \(q_\delta(p)\), exactly as in Theorem~\ref{p1b:rt:theorem62-roots}.
Collapsing the ball boundary gives the degree-one cap product, with
the actual frame contraction at \((1,1)\). Thus its commutator is
the retained \(\nu'\), and its sixth power is \(6\nu'\).

Next retain the entire side contribution. Frame \eqref{p1b:eq:natural-nt-5} is defined
on the whole exterior and agrees with \([e_2,e_3,e_1]\) on the based
wedge. Consequently \(g:S^5\to SU(2)\) is the actual clutch of
\(E_s\) between that exterior frame and the actual cap frame.
Its class is the nonzero \(\eta_3\eta_4\), by the retained exotic
sector and Theorem~\ref{p1b:thm:g15-unstable}.
There is no unaccounted change from the exterior framing: the punctured
product retracts to \(S^3\vee S^3\), and every map of that wedge
into \(SU(2)\) extends over its product by group multiplication.
Its restriction to the attaching five-sphere is therefore null.

The line projection \(g\mapsto g\mathbb Ce_1\) is the actual
complex Hopf map \(\eta_2\) in these ordered two-column frames.
The side formula, capped by its two constant poles, is its suspension
composed with \(Eg\). It has precisely the class
\(\eta_3 E(\eta_3\eta_4)=\eta_3\eta_4\eta_5=6\nu'\),
by Lemma~\ref{p1b:rt:lemma104}. Pinching the compact bottom support
away from the side gives the sum of these two full maps. All constants,
seams and complements are included in that pinch. Orientation signs
do not affect either term, which has order two. Their sum is zero.
This proves \eqref{p1b:eq:natural-nt-11} for the displayed actual input, rather than
for an independently calibrated model.
\end{proof}

\paragraph{The prescribed construction and boundary restriction.}
The interior filling must satisfy the following equation:
\[
 \mathcal X:[0,1]\times\mathbb B\longrightarrow SU(2),
 \qquad \mathcal X|_{\partial([0,1]\times\mathbb B)}
                              =\Omega_B.                \tag{NT.12}\label{p1b:eq:natural-nt-12}
\]
Equation \eqref{p1b:eq:natural-nt-11} proves existence; it is not a formula for this filling
and does not calculate its natural boundary comparison. Once an actual
\(\mathcal X\) is constructed, set
\(\mathcal K=Q_s\operatorname{diag}(\mathcal X,1)Q_s^*\)
inside the ball and \(\mathcal K=K^\#\) outside. The seams in
\eqref{p1b:eq:natural-nt-10} then give exactly \eqref{p1b:eq:natural-nd-11}, and
\(F_h=\mathcal K H^{-3}\) has the exact original endpoints.

The exterior restriction has been applied along the retained map,
with no normal-frame rotation:
\[
 a=i e^{is},\quad x=y e^{3\pi\rho u/2},\quad
 p=y e^{-3\pi\rho u/2},\qquad
 \widetilde D^\#=R^*F_Y^*D_Y^\# T_YF_YR.                \tag{NT.13}\label{p1b:eq:natural-nt-13}
\]
This is defined on the entire inner stage outside the preimage of
\(\mathbb B\). On the cylinder the original constructed joining is
retained; on the outer stage its exact matrix is
\(F_Z^*\operatorname{diag}(1,e^{2is},1)F_Z\).
The ordinary boundary frames \(F_YR,F_Z\), actual clutch extension
\(R\), stabilizing columns and all base coordinates remain those of
\eqref{p1b:eq:natural-nc-4}--\eqref{p1b:eq:natural-nc-12}. The preimage of the cap lies strictly inside the
inner stage; it has not been replaced by a different boundary sphere.

The interior restriction uses the retained \eqref{p1b:eq:natural-nt-12}, supplied by
\eqref{p1b:eq:natural-ny-1}--\eqref{p1b:eq:natural-oa-4}. Its lower-stage
comparisons, including the known local determinant and sixth-stage
term \(3z\), must then be carried to the retained seven-cell before
using the existing detector. The complete original comparison is evaluated in
\eqref{p1b:eq:natural-aj-4}--\eqref{p1b:eq:natural-aj-7}:
\(h_t=w_t=1\), \(r_J=\zeta_{\rm nat}=0\), with every retained
extension test decided. This component is used with its full frames,
lower divisions and joining; its isolated value is not a replacement
for that complete comparison.


\subsection{An applied bottom homotopy and the actual side filling}
\label{p1b:sec:natural-determinant-cap-comparison}

Retain the actual cylinder, cap frame, exterior closing, matching toggle
and joining of \eqref{p1b:eq:natural-nt-1}--\eqref{p1b:eq:natural-nt-13}. The actual sphere value \(w_t=1\),
all counts and all group types remain fixed. This continuation constructs
two named components of \eqref{p1b:eq:natural-nt-12}: a homotopy of its actual bottom and
a filling of its actual side and top. It does not evaluate \(h_t\).

\paragraph{Dependencies executed.}
Normalize the displayed cylinder boundary by an extendible rank-two
path; fill its actual side and top radially; apply a based characteristic-map
homotopy to the actual bottom cap quaternions; retain the seams; write
the single remaining relative cancellation and its exact assembly equation.
No calibration value is substituted for a natural map.

\begin{proposition}[Exact normalization and an actual side/top filling]
\label{p1b:proposition:natural-determinant-side-fill}
Write \(z=R_B r v\), \(0<r\le1\), \(v\in S^5\), and use
the actual \(g(v)=(Q_s^*\mathsf V^\#)_{2\times2}\) of \eqref{p1b:eq:natural-nt-10}.
Put
\[
 A(t)=\operatorname{diag}(e^{\pi it},e^{-\pi it}),\qquad
 S_g(t,v)=g(v)A(t)g(v)^{-1}A(t)^{-1}.                 \tag{NU.1}\label{p1b:eq:natural-nu-1}
\]
The exact normalization \(M_B=\Omega_B A(t)^{-1}\) has faces
\[
 M_B(0,z)=(Q_s^*H^6Q_s)_2,\quad M_B(1,z)=I_2,
 \quad M_B(t,R_Bv)=S_g(t,v).                           \tag{NU.2}\label{p1b:eq:natural-nu-2}
\]
Its actual side and top are filled by
\[
 G_0(t,z)=S_g(1-(1-t)r,v),\quad G_0(t,0)=I_2,
 \qquad E_g(z)=G_0(0,z).                               \tag{NU.3}\label{p1b:eq:natural-nu-3}
\]
The map \(G_0 A(t)\) fixes the stabilizing column in the actual
cap frame, has final value \(-I_2\), and agrees with the exact
exterior closing on the whole side. Its initial value is \(E_g\),
not the required \(H^6\). Thus it is a constructed background
for the actual patch, not a claimed filling of \eqref{p1b:eq:natural-nt-12}.
\end{proposition}
\begin{proof}
The line projection in \eqref{p1b:eq:natural-nt-10} gives its side exactly as
\(gA(t)g^{-1}\). Right multiplication proves \eqref{p1b:eq:natural-nu-2}, including
both seams. The normalizing factor extends over the entire actual
cylinder. More explicitly \(\Omega_B A(\lambda t)^{-1}\),
\(0\le\lambda\le1\), is a homotopy of the complete boundary;
it is fixed on the bottom and respects both seams for every \(\lambda\).
Its changing top is a constant matrix over that face.

Uniformly in \(v\), the operator norm satisfies
\[
 \|S_g(q,v)-I\|\le2\|A(q)+I\|
        \le2\pi(1-q),\qquad 0\le q\le1.               \tag{NU.4}\label{p1b:eq:natural-nu-4}
\]
This follows by subtracting the central matrix \(-I\) in the
commutator difference and using unitarity. Substitution gives
\(\|G_0(t,z)-I\|\le2\pi(1-t)r\), proving continuity at the
entire cap center, including its polar direction collapse.
At \(r=1\) it is exactly \(S_g(t,v)\); at \(t=1\) it is
identity. Since \(S_g(0,v)=S_g(1,v)=I\), the bottom \(E_g\)
is identity on \(\partial\mathbb B\) and at its center.

The based direction \(v_*=(1,0,0,0,0,0)\) has \(p=1\).
The actual section, shifted sections, and their transporters on
this wedge are \(e_1,e_1,I\). Hence both ordered frames are
\([e_2,e_3,e_1]\) and \(g(v_*)=I\), exactly. No constant
frame rotation or unmeasured basepoint adjustment is used.

On the side, \(G_0A=S_gA=gAg^{-1}\), the precise side of
\(\Omega_B\). At the top it is \(-I\). Conjugating its block
by the unchanged actual \(Q_s\) fixes \(s\), ends at \(N\),
and pastes exactly to \(K^\#\). The initial value assertion follows
from \eqref{p1b:eq:natural-nu-3}. In particular the known lower-stage terms have not
been silently removed by this normalization.
\end{proof}

\begin{theorem}[A based homotopy of the actual bottom, with its cap frame]
\label{p1b:theorem:natural-determinant-bottom-homotopy}
In the actual \(Q_s\) frame put
\[
 a=q_\delta(x),\qquad
 b=(b_0,b_1,-b_2,-b_3),\quad
 q_\delta(p)=(b_0,b_1,b_2,b_3).                         \tag{NU.5}\label{p1b:eq:natural-nu-5}
\]
Then \((Q_s^*H^6Q_s)_2=[a,b]^6\) exactly. There is the
following applied based homotopy \(B_\lambda(z)\) from this
entire actual bottom to \(D(\mu(a,b))^6\). It is identity
outside the original two-cap product for every \(\lambda\).

Away from \(a=1\) or \(b=1\), define
\(p_D=\pi^{-1}\log(-b)\in\operatorname{Im}\mathbb H\),
using \(|p_D|<1\), and
\(u_D=a\sqrt{1-|p_D|^2}\). For \(R=\operatorname{Re}u_D\)
define \(R_\ell=\min\{1,R+\ell(1+R)\}\) and
\[
 f_\ell(p_D,u_D)=
 \left(c_\ell p_D,\ R_\ell+c_\ell\operatorname{Im}u_D\right),
 \quad c_\ell=\sqrt{\frac{1-R_\ell^2}{1-R^2}}.         \tag{NU.6}\label{p1b:eq:natural-nu-6}
\]
At \(R=\pm1\) use the fixed pole value. The collapsed
three-disk \(\operatorname{Im}u_D=0,R\ge0\) is preserved
and is sent to \((0,1)\) by \(f_1\). Thus
\(\mu(a,b)=f_1(p_D,u_D)\) extends over the whole based smash,
with its wedge sent to \((0,1)\).

For \(a_D=u_D/|u_D|\), use the based formula
\[
 T_\ell(p_D,u_D)=a_De^{\pi p_D}a_D^{-1}
       e^{-(1-\ell)\pi p_D+\ell\pi i}e^{-\ell\pi i},
 \qquad D(p_D,u_D)=a_De^{\pi p_D}a_D^{-1}.              \tag{NU.7}\label{p1b:eq:natural-nu-7}
\]
At \(u_D=0\), the first factor is \(-1\); this defines
the continuous value independently of its absent direction.
Set, writing \((p_\ell,u_\ell)=f_\ell(p_D,u_D)\),
\[
 B_\lambda(z)=
 \begin{cases}
 {}[u_{2\lambda}/|u_{2\lambda}|,
                     -e^{\pi p_{2\lambda}}]^6,
                                      &0\le\lambda\le1/2,\\
 T_{2\lambda-1}(\mu(a,b))^6,&1/2\le\lambda\le1.
 \end{cases}                                           \tag{NU.8}\label{p1b:eq:natural-nu-8}
\]
The commutator value at \(u_{2\lambda}=0\) is identity.
On the based wedge set \(B_\lambda=I\). All formulas are in
the actual ordered rank-two block and retain its stabilizing column.
\end{theorem}
\begin{proof}
Both retained cap gauges use \(T(e_1,s)\). Exchange of the second
cap conjugates its block by \(\operatorname{diag}(1,-1)\), giving
precisely the fixed rotation in \eqref{p1b:eq:natural-nu-5}. Their commutator in the actual
\(Q_s\) frame is therefore exactly \([a,b]\); taking its sixth
power proves the initial equation without a frame contraction.

The source collapse \eqref{p1b:eq:natural-nu-6} and characteristic-map deformation in
\eqref{p1b:eq:natural-nu-7} are the concrete source of P\"uttmann's Lemmas 6.2--6.4
and Corollary 6.3. They are applied here to the actual cap quaternions,
with the displayed based normalization before taking the sixth power.
For completeness, the needed continuity and endpoint checks follow
directly. Formula \eqref{p1b:eq:natural-nu-6} stays on the unit six-sphere, preserves the
collapsed disk, and equals its positive pole where \(R_\ell=1\).
The initial quotient expression is exactly \((a,b)\).
Whenever \(u_\ell=0\), its second quaternion is \(+1\), so
the commutator tends to identity independently of the first axis.
At the pole endpoints the shrinking transverse coordinates prove
continuity of \(f_\ell\).

At \(\ell=0\), \(T_0=[a_D,-e^{\pi p_D}]\), proving
the intermediate seam in \eqref{p1b:eq:natural-nu-8}. At \(\ell=1\), the two
last factors cancel and \(T_1=D\). When \(u_D\to0\),
\(e^{\pi p_D}\to-1\) uniformly in the missing axis, proving
continuity of the first factor and hence of the entire deformation.
At \(\mu(a,b)=(0,1)\), its value is
\(e^{\ell\pi i}e^{-\ell\pi i}=1\) for every \(\ell\).
This right factor is essential: the unnormalized published deformation
would move the based value and would not paste to our fixed side seam.

The characteristic map with the usual Dur\'an sign is \(-D\).
Its sixth power is exactly \(D^6\), so no sign is being replaced
at the endpoint. When either actual cap quaternion is identity,
all of \eqref{p1b:eq:natural-nu-8} is identity by the based extension just proved. Thus
the homotopy is fixed on a neighborhood of \(\partial\mathbb B\).
It applies to the full bottom while keeping \(Q_s\), all side data,
the matching toggle and the joining data unchanged.
\end{proof}

\paragraph{The complete applied boundary homotopy.}
After \eqref{p1b:eq:natural-nu-2}, replace its bottom by \(B_\lambda\), keep its top
identity and its entire side \(S_g(t,v)\) unchanged:
\[
 M_{B,\lambda}|_{\mathrm{bottom}}=B_\lambda,
 \quad M_{B,\lambda}|_{\mathrm{top}}=I,
 \quad M_{B,\lambda}|_{\mathrm{side}}=S_g.             \tag{NU.9}\label{p1b:eq:natural-nu-9}
\]
The just-proved boundary support gives both seam equalities for all
\(\lambda\). This is an actual homotopy of the complete cylinder
boundary, applied to the retained target; it is not merely an equality
of homotopy-group coefficients.

\paragraph{The retained relative comparison and its exact role.}
The endpoint cancellation uses the based homotopy, supplied by
\eqref{p1b:eq:natural-ny-1}--\eqref{p1b:eq:natural-nz-20}, between these two constructed entire maps:
\[
 \begin{aligned}
 \mathcal R:&[0,1]\times(\mathbb B/\partial\mathbb B)\to SU(2),\\
 \mathcal R(0,z)&=D(\mu(q_\delta(x),\operatorname{rot}q_\delta(p)))^6,\\
 \mathcal R(1,z)&=E_g(z),\qquad
 \mathcal R(\ell,\partial\mathbb B)=I.
 \end{aligned}                                            \tag{NU.10}\label{p1b:eq:natural-nu-10}
\]
Both endpoint classes are \(6\nu'\); \eqref{p1b:eq:natural-nt-11} proves their
homotopy existence. Neither \eqref{p1b:eq:natural-nu-10} nor the natural bit is being
declared calculated from that equality. P\"uttmann's later
characteristic-map deformation leads to a particular suspended
map \(\tau\), with its coordinate and fibre transport. Its
identification with the actual \(g\), after the displayed \(\mu\),
is supplied through the intrinsic bridges \eqref{p1b:eq:natural-ny-1}--\eqref{p1b:eq:natural-nz-20}. A
twelfth-power nullhomotopy alone does not perform \eqref{p1b:eq:natural-nu-10}.

The precise assembly removes any separate seam or endpoint ambiguity.
Concatenate the constructed \(B_\lambda\) with the required
\(\mathcal R\), calling the full path \(\mathcal B_\ell\), and
put \(\chi(t)=\min\{4t,1\}\). The complete desired ordinary cap
map would be exactly
\[
 \mathcal X(t,z)=\mathcal B_{\chi(t)}(z)E_g(z)^{-1}
                      G_0(t,z)A(t).                    \tag{NU.11}\label{p1b:eq:natural-nu-11}
\]
Indeed its bottom is \(H^6\), its top is \(-I\), and its side
is \(S_gA=gAg^{-1}\). These are exactly all three faces of
\eqref{p1b:eq:natural-nt-10}. Conjugating its stabilized block by \(Q_s\), pasting to
\(K^\#\), and setting \(F_h=\mathcal K H^{-3}\) give the
original closing equation, the exact measured origin, applied toggle
and joining. Formula \eqref{p1b:eq:natural-nu-11} identifies the role of \(\mathcal R\). The
full intrinsic closing is supplied in \eqref{p1b:eq:natural-ny-1}--\eqref{p1b:eq:natural-oa-4}, with these
exact prescribed faces.

The complete original comparison is evaluated in
\eqref{p1b:eq:natural-aj-4}--\eqref{p1b:eq:natural-aj-7}:
\(h_t=w_t=1\), \(r_J=\zeta_{\rm nat}=0\), with every retained
extension test decided. This component is used with its full frames,
lower divisions and joining; its isolated value is not a replacement
for that complete comparison.

\paragraph{Primary source.}
T. P\"uttmann, \emph{On the commutator of unit quaternions and the
numbers 12 and 24}, arXiv:1101.5147v4 (2012), \S6, Lemmas 6.2--6.4
and Corollary 6.3; see also \S2 for the subsequent construction.
The application above retains the actual target's source maps and adds
its necessary based normalization.



\subsection{An applied characteristic-map comparison with its fibre transport}
\label{p1b:sec:natural-determinant-stiefel-comparison}

Retain the actual cap, section, cap frame, side map, matching toggle and
joining of \eqref{p1b:eq:natural-nt-1}--\eqref{p1b:eq:natural-nu-11}. The original target still has one evaluated
bit, \(w_t=1\). This continuation constructs an actual based homotopy
\(D^6\mu\to J_\tau\), including its source and rank-two fibre
identification. It does not identify \(J_\tau\) with the actual \(E_g\)
or evaluate \(h_t\).

\paragraph{Applied dependencies.}
Solve the two-column horizontal deformation, including its lower-hemisphere
correction; assemble the full Stiefel map; apply it to the actual cap source
\(\mu\); identify its initial fibre map with the actual \(D^6\); normalize
the moving fibre at the actual based direction; replace the entire bottom
while retaining the actual side, top and both seams. All six constructions
are performed below. No primitive parity or calibration map is substituted
for the natural target.

\subsubsection{A complete horizontal deformation and its explicit endpoint}

For \(z=(z_0,z_1)\in\mathbb C^2\), put \(L_z=|z|\) and
use the following homogeneous real three-vectors, with value zero at \(z=0\):
\[
\begin{gathered}
h_1(z)=\frac1{L_z}\bigl(|z_0|^2-|z_1|^2,
                   \operatorname{Re}(2\bar z_0z_1),
                   \operatorname{Im}(2\bar z_0z_1)\bigr),\\
 h_2(z)=\frac1{L_z}\bigl(-2\operatorname{Re}(z_0z_1),
                   \operatorname{Re}(\bar z_0^2-z_1^2),
                   \operatorname{Im}(\bar z_0^2-z_1^2)\bigr).
\end{gathered}
\tag{NV.1}\label{p1b:eq:natural-nv-1}
\]
Direct multiplication gives \(|h_1|=|h_2|=L_z\) and
\(h_1\cdot h_2=0\). For \(x\in S^5\subset\mathbb R^6\), define
\[
\begin{gathered}
\widehat h(x)=\bigl((x_0,x_1,h_1(x_4+ix_5,x_2+ix_3)),
                    (x_1,-x_0,h_2(x_4+ix_5,x_2+ix_3))\bigr),
 \\
 P_2(x)=(2x_0^2-1,2x_0x_1,\ldots,2x_0x_5).
\end{gathered}
\tag{NV.2}\label{p1b:eq:natural-nv-2}
\]
The columns are an ordered orthonormal pair. The initial pair needed here
is precisely \(\widehat h(P_2x)\), rather than an assumed pointwise
commutation of squaring and homogeneous suspension.

\begin{proposition}[A closed horizontal lift of the actual initial pair]
\label{p1b:proposition:natural-determinant-V5-lift}
There is a specified three-stage horizontal homotopy
\(\mathcal H^5:S^5\times I\to V_{5,2}\) with
\[
 \mathcal H^5_0=\widehat hP_2,\qquad
 \mathcal H^5_1=(e_0,(0,\tau)),\qquad \tau:S^5\to S^3.
                                                        \tag{NV.3}\label{p1b:eq:natural-nv-3}
\]
Its second column is the following explicit parallel transport, including
all pole and seam limits.
\end{proposition}
\begin{proof}
Write \(x=(\cos t,\sin t\,v)\), \(0\le t\le\pi\), \(v\in S^4\),
and set \(c=\cos2t\), \(s=\sin2t\), \(k=|s|\), \(d=sv_0\).
The initial columns are
\(\gamma=(c,d,k h_1(z_v))\),
\(\delta=(d,-c,k h_2(z_v))\), where
\(z_v=(v_3+iv_4,v_1+iv_2)\).
For \(t\le\pi/2\) leave them fixed in the first stage.
For \(t>\pi/2\), put \(R=k|z_v|\), \(L=(c^2+d^2)^{1/2}\),
\(e_j=h_j(z_v)/|z_v|\), and, for \(0\le\theta\le\pi\), put
\[
 \begin{aligned}
 e_\theta&=\cos\theta\,e_1+\sin\theta\,e_2,\quad
 T_\theta=-\sin\theta\,e_1+\cos\theta\,e_2,\\
 \gamma_\theta&=(c,d,R e_\theta),\\
 \delta_\theta&=\left(d-R^2c\frac{\sin L\theta}{L},
             -c-R^2d\frac{\sin L\theta}{L},
             R\cos L\theta\,T_\theta+
                         RL\sin L\theta\,e_\theta\right).
 \end{aligned}                                            \tag{NV.4}\label{p1b:eq:natural-nv-4}
\]
Use \(\sin(L\theta)/L=\theta\) at \(L=0\); at \(R=0\) the
constant initial pair is the continuous value. This is the Hopf sign
correction induced by
\(z\mapsto\cos(\theta/2)z+\sin(\theta/2)(-\bar z_1,\bar z_0)\).
Its first column ends at \((c,s\rho)\), where
\(\rho=(v_0,h_1(z_v))\in S^3\).

Here is a direct verification of the full lift. When \(L\ne0\), set
\(m=(-d,c)/L\), \(N=(c,d)/L\),
\(B_\theta=RN-Le_\theta\). The initial second column is
\(-Lm+RT_0\). The tangent connection along \(\gamma_\theta\)
satisfies \(\nabla T=L B\), \(\nabla B=-L T\).
Solving the two-dimensional rotation gives
\(-Lm+R(\cos L\theta\,T_\theta-\sin L\theta\,B_\theta)\),
which expands to \eqref{p1b:eq:natural-nv-4}. It satisfies exactly
\(\dot\delta=-\langle\delta,\dot\gamma\rangle\gamma\).
The expansion, not the divided moving frame, proves its \(L=0\) limit.
The factors \(R\) prove its \(|z_v|=0\) limit.
At \(t=\pi/2\), \(R=0\), so the upper and lower pieces agree.

For the second stage use \(J\rho=(-\rho_1,\rho_0,-\rho_3,\rho_2)\).
Keep the upper hemisphere fixed, and on the lower hemisphere rotate
\(e_\theta=\cos\theta\rho+\sin\theta J\rho\), \(0\le\theta\le\pi\).
Given its actual initial \(\delta_0\), put
\[
 \begin{aligned}
 T_0&=(0,J\rho),\quad B_0=(s,-c\rho),\quad
 a_0=\delta_0\cdot T_0,\quad b_0=\delta_0\cdot B_0,\\
 W&=\delta_0-a_0T_0-b_0B_0,\quad
 T_\theta=(0,-\sin\theta\rho+\cos\theta J\rho),\quad
 B_\theta=(s,-c e_\theta),\\
 a_\theta&=a_0\cos c\theta+b_0\sin c\theta,\quad
 b_\theta=b_0\cos c\theta-a_0\sin c\theta,\\
 \gamma_\theta&=(c,s e_\theta),\qquad
 \delta_\theta=a_\theta T_\theta+b_\theta B_\theta+W.
 \end{aligned}                                            \tag{NV.5}\label{p1b:eq:natural-nv-5}
\]
The connection equations are \(\nabla T=cB\), \(\nabla B=-cT\);
\(W\) is a fixed vector perpendicular to this moving plane.
Thus \eqref{p1b:eq:natural-nv-5} is again the exact horizontal solution.
At \(s=0\), \(c=\pm1\), the frame and coefficient rotations cancel,
so the second column remains constant. This proves the hemisphere seam
and both source pole limits. The first column now equals \((c,k\rho)\)
on both halves.

For the final stage put \(c_\ell=\ell+(1-\ell)c\),
\(k_\ell=\sqrt{1-c_\ell^2}\). With the actual second-stage endpoint
\(\delta_0\), set
\[
 T_0=(-k,c\rho),\quad T_\ell=(-k_\ell,c_\ell\rho),\quad
 a=\delta_0\cdot T_0,\quad
 \gamma_\ell=(c_\ell,k_\ell\rho),\quad
 \delta_\ell=\delta_0+a(T_\ell-T_0).
                                                        \tag{NV.6}\label{p1b:eq:natural-nv-6}
\]
The meridian tangent is parallel; every perpendicular tangent is fixed.
This proves horizontality, unit length and orthogonality. At the end
\(\gamma=e_0\) and the first coordinate of \(\delta\) is zero. Define
\[
                   \tau(x)=(\delta_1(x))_{1,\ldots,4}.
                                                        \tag{NV.7}\label{p1b:eq:natural-nv-7}
\]
This is a whole map, including its polar limits. The preceding direct
formulae prove continuity at the seams and absent axes. At
\(x=\pm e_0\) it is \((-1,0,0,0)\). This value is retained in the
subsequent fibre normalization; it is not silently reset by a frame choice.
\end{proof}

\subsubsection{The full Stiefel deformation, in specified coordinates}

Use real eight-coordinates
\(x=(\operatorname{Re}z_0,\operatorname{Im}z_0,
\operatorname{Re}z_1,\operatorname{Im}z_1,a,b,c,d)\), with
\(z_2=a+id,z_3=c+ib\). Keep the target product on
\(\mathbb O=\mathbb C\times\mathbb C^3\) from \S3 of P\"uttmann;
use the explicit source isometry \(F(z_0,z)=(\bar z_0,iz)\).
Thus \(\kappa(x)=(\overline{F(x)}e_1F(x),
\overline{F(x)}e_2F(x))\) is expressed in the fixed target frame,
where \(e_1=(i,0)\), \(e_2=(0,1,0,0)\).
In target real seven-coordinates the order is
\((\operatorname{Im}z_0,\operatorname{Re}z_1,
\operatorname{Im}z_1,a,b,c,d)\).
Direct multiplication gives
\[
 \kappa_1=(i(|z_0|^2-|z_1|^2-|z_2|^2-|z_3|^2),
              2\bar z_0z_1,2\bar z_0z_2,2\bar z_0z_3),
\]
\[
 \kappa_2=(-2i\operatorname{Re}(z_0z_1),
      \bar z_0^2-z_1^2-|z_2|^2-|z_3|^2,
      -2i\operatorname{Im}z_1 z_2+2i\operatorname{Re}z_0\bar z_3,
      -2i\operatorname{Im}z_1 z_3-2i\operatorname{Re}z_0\bar z_2).
                                                        \tag{NV.8}\label{p1b:eq:natural-nv-8}
\]
The minus sign on \(z_1^2\) follows from the actual product.

For completeness the first deformation is specified as follows.
Put \(C=\cos\theta,S=\sin\theta\), \(0\le\theta\le\pi/2\),
and use raw columns
\[
 u_1=\left(\begin{gathered}
 i(|z_0|^2-|z_1|^2-(|z_2|^2+|z_3|^2)C^2),\quad
 2\bar z_0z_1,\\
 2z_2(S^2+\bar z_0C^2)-2\bar z_3(1-\bar z_0)CS,\\
 2z_3(S^2+\bar z_0C^2)+2\bar z_2(1-\bar z_0)CS
 \end{gathered}\right),
\]
\[
 \begin{aligned}
 u_2=\bigl(&-2i\operatorname{Re}(z_0z_1),
 \bar z_0^2-z_1^2-(|z_2|^2+|z_3|^2)C^2,\\
 &2z_2\operatorname{Im}z_1(CS-iC^2)
       -2\bar z_3((1-\operatorname{Re}z_0)CS-i(S^2+\operatorname{Re}z_0C^2)),\\
 &2z_3\operatorname{Im}z_1(CS-iC^2)
       +2\bar z_2((1-\operatorname{Re}z_0)CS-i(S^2+\operatorname{Re}z_0C^2))\bigr).
 \end{aligned}                                            \tag{NV.9}\label{p1b:eq:natural-nv-9}
\]
First Gram--Schmidt \((u_1(\theta),u_2(0))\), then
\((u_1(\pi/2),u_2(\theta))\). The columns are independent:
the real and imaginary equations in the two orthogonal last complex
directions give no solution for proportionality, as in \S7.1 of the
primary source. They never vanish. At the last endpoint rescale the first
Hopf block by \(L/(1-r+rL)\) and the last four coordinates by \(2-r\),
\(0\le r\le1\), then Gram--Schmidt, where
\(L=(|z_0|^2+|z_1|^2)^{1/2}\). The result is the pair
\[
 \widetilde h(x)=\bigl((h_1(z_0,z_1),a,b,c,d),
                    (h_2(z_0,z_1),b,-a,-d,c)\bigr).
\]
The \(L=0\) block has value zero and a uniform continuous limit.

There is also a complete source homotopy from \(P_2:S^7\to S^7\)
to its double lower suspension. For \(x=(v,w)\in\mathbb R^6\times\mathbb R^2\),
normalize the following nonzero vector for \(0\le\theta\le\pi/2\):
\[
\begin{gathered}
\left(\begin{gathered}
 x_0^2-\sum_{j=1}^5x_j^2-|w|^2\cos^2\theta,
 2x_0x_1,\ldots,2x_0x_5,\\
 2(w_0 A-w_1 B),2(w_1 A+w_0 B)
\end{gathered}\right),\\
 A=\sin^2\theta+x_0\cos^2\theta,\qquad
 B=(1-x_0)\cos\theta\sin\theta.
\end{gathered}
\tag{NV.10}\label{p1b:eq:natural-nv-10}
\]
Its length is \(1+|w|^2\sin^2\theta\). At the endpoint it is
\((|v|^2P_2(v/|v|),2w)\) before normalization. Complete the homotopy
by normalizing
\((( (1-r)|v|^2+r|v|)P_2(v/|v|),(2-r)w)\), \(0\le r\le1\).
The \(v=0\) block is zero; no direction there is selected.

Use actual rotations from identity to the following two isometries:
\[
 f_5(x)=(x_4,x_5,x_2,x_3,x_0,x_1),\qquad
 f_6(y)=(y_3,y_4,y_0,y_1,y_2).                         \tag{NV.11}\label{p1b:eq:natural-nv-11}
\]
For \(f_5\), rotate by \(\pi r\) in the plane spanned by
\((e_0-e_4)/\sqrt2,(e_1-e_5)/\sqrt2\).
For \(f_6\), use the two Fourier planes
\(c_k(j)=\sqrt{2/5}\cos(2\pi kj/5)\),
\(s_k(j)=\sqrt{2/5}\sin(2\pi kj/5)\), \(k=1,2\), rotating them
by \(4\pi kr/5\). These fix the constant Fourier axis and give
exactly the permutation at \(r=1\). Apply these rotations after the
squaring map in the first six-coordinate block, rather than assuming
that \(f_5\) commutes with squaring.

Double-lift \eqref{p1b:eq:natural-nv-3} by
\((\gamma,\delta)\mapsto((|v|\gamma,c,d),(|v|\delta,-d,c))\).
The two added coordinates preserve unit length and orthogonality and have
unambiguous limits at \(v=0\). At its endpoint contract the first column
\((|v|e_0,c,d)\) along its meridian to \(e_0\). Its second column is
perpendicular to that entire meridian, so stays fixed. Then rotate its
last two coordinates by \(-\pi/2\), leaving the first column fixed.

\begin{theorem}[An actual whole Stiefel map]
\label{p1b:theorem:natural-determinant-V7-map}
The ordered concatenation just specified gives a whole map
\[
 \alpha:S^7\times I\longrightarrow V_{7,2},\qquad
 \alpha_0=\kappa P_2,\quad
 \alpha_1(x)=(e_1,\Sigma^2_{\rm low}\tau(x)),
                                                        \tag{NV.12}\label{p1b:eq:natural-nv-12}
\]
where the second endpoint is precisely
\((|v|\tau(v/|v|),w)\in S^5\) in the retained real target order.
All rotations, source rescalings, hemisphere corrections and pole values
are part of this map. Neither source nor fibre endpoints are inferred
solely from a homotopy-group class.
\end{theorem}
\begin{proof}
First use \eqref{p1b:eq:natural-nv-9} at \(P_2x\), then \(\widetilde h\) along \eqref{p1b:eq:natural-nv-10}.
Next double-lift the rotations \eqref{p1b:eq:natural-nv-11}, then \eqref{p1b:eq:natural-nv-3}, and finally apply
the described meridian contraction and last-plane rotation.
The endpoint of each stage is the actual initial pair of the next.
Equations \eqref{p1b:eq:natural-nv-1}--\eqref{p1b:eq:natural-nv-11} prove all column and polar assertions.
In particular the lower-hemisphere correction precedes the nullhomotopy;
omitting it changes the initial map. The first column is \(e_1\) at
the final stage and the second is exactly the displayed lower suspension.
This proves \eqref{p1b:eq:natural-nv-12}, not merely the existence of some unspecified \(\tau\).
\end{proof}

\subsubsection{The actual fibre identification and based cap homotopy}

The identification of the \(Sp(2)\) spinor coordinates with the target
octonion coordinates is
\[
 \mathsf J(r,u)=(r_0+ir_1,r_2+ir_3,u_0-iu_1,u_2+iu_3).
\]
It is an isometry of \(\mathbb R^8\), distinct from the usual
Cayley--Dickson product identification. In the actual bottom, with
\((p_D,u_D)=\mu(q_\delta(x),\operatorname{rot}q_\delta(p))\), use
the resulting specific unit direction in the source real seven-order:
\[
 C(p_D,u_D)=(-p_{D,1},p_{D,3},-p_{D,2},
             -u_{D,1},-u_{D,2},u_{D,3},-u_{D,0}),\qquad
 v_*=(0,0,0,0,0,0,-1).                                \tag{NV.13}\label{p1b:eq:natural-nv-13}
\]
It follows by applying \(F^{-1}\mathsf J\) to \((p_D,u_D)\).
The actual collapsed wedge gives exactly \(v_*\). No independent
degree-one replacement of \(\mu\) is made.

Here is a finite exact specification of the spinor/fibre identification.
Let \(E_{ab}\) send target basis vector \(e_a\) to \(e_b\) and
\(e_b\) to \(-e_a\), with \(e_0=1\), indices \(0,\ldots,7\) in
the target real order. For the \(Sp(2)\) Lie generators
\(\operatorname{diag}(q,0)\), \(\operatorname{diag}(0,q)\), and
\(\left(\begin{smallmatrix}0&q\\-\bar q&0\end{smallmatrix}\right)\),
their first triality component \(a\) is respectively
\[
\begin{array}{c|ccc}
 &q=i&q=j&q=k\\\hline
 \operatorname{diag}(q,0)&E_{56}-E_{47}&E_{46}+E_{57}&-E_{45}+E_{67}\\
 \operatorname{diag}(0,q)&-E_{56}-E_{47}&E_{46}-E_{57}&E_{45}+E_{67}
\end{array}
\]
\[
\begin{array}{c|cccc}
 q&1&i&j&k\\\hline
 \left(\begin{smallmatrix}0&q\\-\bar q&0\end{smallmatrix}\right)
     &-2E_{35}&-2E_{36}&-2E_{37}&-2E_{34}
\end{array}.                                           \tag{NV.14}\label{p1b:eq:natural-nv-14}
\]
Specifically \(b=\mathsf J Q\mathsf J^{-1}\) at the Lie-algebra
level, and \(c=a-L_{b1}\). For each generator and each of the eight
unit coordinate octonions \(x\), multiplication gives exactly
\([a,L_x]=L_{bx}-L_xL_{b1}\), with \(c\) skew. This is the
infinitesimal triality equation \(a(xy)=b(x)y+x c(y)\).
The table also preserves Lie brackets and the trace inner product.
These are finite integer identities; thus the table integrates, using
simple connectivity of \(Sp(2)\), to its specified embedding in
\(Spin(7)\) fixing \(e_1,e_2\) in the first component.

Its \(Sp(1)\) endpoint is the following actual fixed-fibre representation:
\[
 \varrho(q)\mathsf J(0,u)=\mathsf J(0,qu),\qquad
 \varrho(q)e_j=e_j\quad(1\le j\le3).
                                                        \tag{NV.15}\label{p1b:eq:natural-nv-15}
\]
In complex target coordinates this block is the standard retained
quaternion block with the fixed rotation
\((q_0,q_1,q_2,q_3)\mapsto(q_0,-q_1,-q_2,q_3)\).
Using \(\varrho^{-1}\) below retains the original quaternion and the
actual ordered \(Q_s\) cap block; it does not leave this rotation unmeasured.

Specify the connection on \(G_2\to V_{7,2}\) by the trace metric.
For each \(\alpha_\lambda\), lift its north frame along \(\lambda\)
starting at identity, giving \(p_\lambda\). For
\(\gamma_v(t)=(\cos t,\sin t\,v)\), let \(L_\lambda(t,v)\) be
the following unique horizontal solution starting at \(p_\lambda\):
at the pair \((\gamma,\delta)=\alpha_\lambda\gamma_v(t)\), take
the minimum trace-norm \(A\in\mathfrak g_2\) satisfying
\(A\gamma=\dot\gamma,A\delta=\dot\delta\), and solve
\(\dot L=AL\). The evaluation has constant rank eleven because
\(G_2\) acts transitively on \(V_{7,2}\); its kernel has dimension
three. Thus this specifies the full transport, rather than just an
endpoint class. Use the piecewise smooth curves and their continuous
polar limits from \eqref{p1b:eq:natural-nv-12}. Equivalently the minimum is computed by
the fixed-rank linear system in the accompanying implementation.
Define
\[
 b_\lambda(v)=L_\lambda(\pi,v_*)^{-1}L_\lambda(\pi,v).
 \qquad b_\lambda(v_*)=I,\quad
 b_\lambda(v)\in G_{2,(e_1,e_2)}=\varrho(SU(2)).
                                                        \tag{NV.16}\label{p1b:eq:natural-nv-16}
\]

\begin{theorem}[A transported homotopy of the actual determinant bottom]
\label{p1b:theorem:natural-determinant-transported-bottom}
In the actual cap frame the following is a complete based homotopy:
\[
 \mathcal H_\lambda(z)=\varrho^{-1}
       b_\lambda\bigl(C(\mu(q_\delta(x),\operatorname{rot}q_\delta(p)))\bigr),
 \quad \mathcal H_0=D(\mu)^6,\quad
 \mathcal H_1=J_\tau,\quad
 \mathcal H_\lambda|_{\partial\mathbb B}=I.
                                                        \tag{NV.17}\label{p1b:eq:natural-nv-17}
\]
Here \(J_\tau\) denotes the displayed entire endpoint at \(\lambda=1\).
All maps fix the prescribed stabilizing column when inserted into \(Q_s\).
\end{theorem}
\begin{proof}
The north and south of the meridian do not depend on \(v\). The based
division in \eqref{p1b:eq:natural-nv-16} consequently lands in one fixed stabilizer throughout,
even while both endpoint frames of \(\alpha\) move. It is exactly identity
at \(v_*\). Path lifting for the specified connection gives continuous
maps on all the pieces and their seams. The formulae have bounded meridian
derivatives on the finitely many pieces, with continuous shrinking polar
limits; the ODE solutions therefore have the same limits. The north lift
moves only during the explicit rotation stage. Hence no unchosen north
lift or constant fibre factor is hidden in \eqref{p1b:eq:natural-nv-16}.

It remains essential to identify the initial map, not just its class.
Apply the horizontal-difference construction to the embedding \eqref{p1b:eq:natural-nv-14}.
Let \(Q(t)\) be the explicit horizontal \(Sp(2)\) lift of
\(\mathsf J(\cos t+\sin t\,p_D,\sin t\,u_D)\).
Its endpoint at \(6\pi\) is
\(\operatorname{diag}(1,D(p_D,u_D)^6)\): the lower block is
\(\widehat u e^{6\pi p_D}\widehat u^{-1}\), exactly the sixth
power of \eqref{p1b:eq:natural-nu-7}, including its \(u_D=0\) limit.
The \(Spin(7)\) horizontal lift has first component
\(C_{\mathsf J(\gamma(2t))}\) at parameter \(6t\), and endpoint
identity at \(6\pi\). Their difference is horizontal in \(G_2\).
Its two columns are precisely
\(\overline{\mathsf J(\gamma(2t))}e_i\mathsf J(\gamma(2t))\),
\(i=1,2\), which equal \(\kappa P_2\gamma_{C(p_D,u_D)}(t)\)
by \eqref{p1b:eq:natural-nv-8} and \eqref{p1b:eq:natural-nv-13}. The trace metric in \eqref{p1b:eq:natural-nv-14} gives the same
connection. Uniqueness of horizontal lift now proves
\(L_0(\pi,C(p_D,u_D))=\varrho(D(p_D,u_D)^6)\), pointwise.
At \(v_*\) the value is identity, so the based division does not change
this equation. This establishes the initial equality in \eqref{p1b:eq:natural-nv-17}.
It is the coordinate application of the horizontal-difference construction
in \S4.2 and \S5.5 of the cited primary source, with the actual embedding
and based direction supplied here.

The actual \(\mu\) is continuous and based by \eqref{p1b:eq:natural-nu-5}--\eqref{p1b:eq:natural-nu-8}.
Its wedge and the complement of the two-cap product give \(v_*\) exactly.
Thus \eqref{p1b:eq:natural-nv-17} is identity there for all parameters, including a neighborhood
of \(\partial\mathbb B\). Inserting the retained quaternion block in
the actual \(Q_s\) frame fixes \(s\) and has determinant one.
The construction is therefore a homotopy of the actual whole bottom,
not an independently calibrated sixth-power map.
\end{proof}

\paragraph{Immediate application and the retained relative comparison.}
After \eqref{p1b:eq:natural-nu-9}, replace its entire bottom by \eqref{p1b:eq:natural-nv-17}, while keeping top
identity and side \(S_g\). The complete boundary family now has
bottom \(J_\tau\), top \(I\), and side the original actual
\(S_g(t,v)=g(v)A(t)g(v)^{-1}A(t)^{-1}\). All seams are fixed.
The original endpoint cancellation has this exact equation:
\[
 \begin{aligned}
 \mathcal R_\tau:&I\times(\mathbb B/\partial\mathbb B)\to SU(2),\\
 \mathcal R_\tau(0,z)&=\varrho^{-1}b_1(C\mu(z))=J_\tau(z),\\
 \mathcal R_\tau(1,z)&=S_g(1-|z|/R_B,z/|z|)=E_g(z),\\
 \mathcal R_\tau(\ell,\partial\mathbb B)&=I,
 \end{aligned}                                            \tag{NV.18}\label{p1b:eq:natural-nv-18}
\]
with the retained center value of \(E_g\). Both maps in \eqref{p1b:eq:natural-nv-18} have
class \(6\nu'\), by \eqref{p1b:eq:natural-nv-17} and \eqref{p1b:eq:natural-nt-11}. This proves existence,
while its constructive intrinsic bridge is supplied in
\eqref{p1b:eq:natural-ny-1}--\eqref{p1b:eq:natural-nz-20}, \eqref{p1b:eq:natural-oa-1}--\eqref{p1b:eq:natural-oa-4}.
In particular the suspended map \(\tau\), its actual transported
endpoint, and the actual clutch \(g\) are not equated merely because
their relevant homotopy classes agree.

If \eqref{p1b:eq:natural-nv-18} is supplied, concatenate \(B\) of \eqref{p1b:eq:natural-nu-8}, the completed
\(\mathcal H\) of \eqref{p1b:eq:natural-nv-17}, and \(\mathcal R_\tau\), denoting it
\(\mathcal B\). The exact complete closing assembly remains
\[
 \mathcal X(t,z)=\mathcal B_{\min(4t,1)}(z)E_g(z)^{-1}
                                G_0(t,z)A(t).            \tag{NV.19}\label{p1b:eq:natural-nv-19}
\]
Its bottom, top and side are exactly \(H^6,-I,gAg^{-1}\), respectively.
The \eqref{p1b:eq:natural-nv-18} endpoint bridge is supplied by the adopted intrinsic
comparisons; \eqref{p1b:eq:natural-nv-17} supplies the retained preceding leg.
The complete inner cap restriction uses
\(a=ie^{is},x=y e^{3\pi\rho u/2},p=y e^{-3\pi\rho u/2}\),
the retained \(F_YR,F_Z\), matching toggle one, and the actual joining.
The complete local determinant and \(3z\) divisions are retained
in its natural \(S^7\to SU(2)\) input, identified and evaluated
by (OA), (OB), (CI) and (AJ). No input to that
detector or value of \(h_t\) is claimed from the cap homotopy alone.

The complete original comparison is evaluated in
\eqref{p1b:eq:natural-aj-4}--\eqref{p1b:eq:natural-aj-7}:
\(h_t=w_t=1\), \(r_J=\zeta_{\rm nat}=0\), with every retained
extension test decided. This component is used with its full frames,
lower divisions and joining; its isolated value is not a replacement
for that complete comparison.

\paragraph{Primary source.}
T. P\"uttmann, \emph{On the commutator of unit quaternions and the numbers
12 and 24}, arXiv:1101.5147v4 (2012), \S\S4.2, 5.2, 5.4--5.5, 7.1--7.3.
The explicit horizontal second column, source correction, actual spinor
identification and its target application above are supplied here.



\subsection{Two actual characteristic-map legs of the determinant patch}
\label{p1b:sec:natural-determinant-characteristic-comparison}

Retain the actual maps, cap frame and source of \eqref{p1b:eq:natural-nu-1}--\eqref{p1b:eq:natural-nv-19}.
The sphere value is \(w_t=1\); the complete determinant value is \eqref{p1b:eq:natural-aj-5}.
The calculation below supplies actual based homotopies from the two
endpoints in \eqref{p1b:eq:natural-nv-18} to one explicit characteristic map applied to two
specified \(S^4\)-valued maps. It supplies neither their intervening
homotopy nor a natural determinant parity.

\paragraph{Dependencies executed.}
Rotate the actual lower suspension into the retained meridian convention;
transport its actual north frame and normalize at the actual cap based
direction; evaluate its characteristic endpoint, retaining the constant
fibre identification; move its based characteristic value to a pole;
perform the same characteristic comparison for the actual exterior map;
apply both legs to the complete bottom and side/top background. These
constructions are executed immediately below. There is no substituted
primitive comparison or new detector.

\subsubsection{The actual source rotation and its transported endpoint}

Use the real source and target order from \eqref{p1b:eq:natural-nv-8}--\eqref{p1b:eq:natural-nv-13}. On \(S^7\)
let \(R_\theta\), \(0\le\theta\le\pi/2\), rotate the
\((x_0,x_7)\) plane by
\[
 (x_0,x_7)\longmapsto
 (\cos\theta\,x_0-\sin\theta\,x_7,
  \sin\theta\,x_0+\cos\theta\,x_7).
\]
On the fixed target \(\mathbb C^3=e_1^\perp\) use the actual
special unitary matrices
\[
 G_\theta=
 \begin{pmatrix}
 \cos\theta&-i\sin\theta&0\\
 -i\sin\theta&\cos\theta&0\\
 0&0&1
 \end{pmatrix},\qquad
 \alpha^\circ_\theta(x)=
 (e_1,G_\theta(\Sigma^2_{\rm low}\tau)(R_\theta x)).
                                                        \tag{NW.1}\label{p1b:eq:natural-nw-1}
\]
The action of \(G_\theta\) is the retained \(SU(3)\subset G_2\)
action, rather than an arbitrary real target rotation. This whole homotopy
starts exactly at \(\alpha_1\) of \eqref{p1b:eq:natural-nv-12}.

For \(v=(v_1,\ldots,v_7)\in S^6\), put
\[
 r_v=(v_1^2+\cdots+v_5^2+v_7^2)^{1/2},\qquad
 n_v=(-v_7,v_1,v_2,v_3,v_4,v_5)/r_v.
\]
Use a zero \(r_v\tau(n_v)\) block at \(r_v=0\).
The actual endpoint along the retained meridian is
\[
 \alpha^\circ_{\pi/2}(\cos t,\sin t\,v)=
 \left(e_1,
   (\cos t+i\sin t\,y_v,\sin t\,z_{v0},\sin t\,z_{v1})\right),
\]
\[
 \beta_\tau(v)=(y_v,z_{v0},z_{v1})=
 \left(-r_v\tau_2(n_v),
 r_v(\tau_1(n_v)-i\tau_0(n_v)),
 v_6+i r_v\tau_3(n_v)\right)\in S^4.
                                                        \tag{NW.2}\label{p1b:eq:natural-nw-2}
\]
These are pointwise equations: at the endpoint the old lower-suspension
coordinates are \((-x_7,x_1,\ldots,x_5,x_6,x_0)\), and multiplication
by \(G_{\pi/2}\) sends their complex second column to
\((x_0-i r_v\tau_2,r_v\tau_1-i r_v\tau_0,v_6+i r_v\tau_3)\).
Unit length and the uniform factor \(r_v\) prove the missing-axis limit.

Let \(p_1\) be the actual north lift in \eqref{p1b:eq:natural-nv-16}, kept with its exact
value. The new north and south second columns are, for all \(\theta\),
\[
 \alpha^\circ_\theta(N)_2=(-\cos2\theta,i\sin2\theta,0),
 \qquad \alpha^\circ_\theta(S)_2=-e_2.
\]
Use the explicit north lift \(p_\theta=G_{2\theta}p_1\).
It sends the standard ordered pair to this exact north frame. Lift
\(\alpha^\circ_\theta\gamma_v\) by the same trace connection as
\eqref{p1b:eq:natural-nv-16}, starting at \(p_\theta\), and call the endpoint
\(L^\circ_\theta(\pi,v)\). Normalize in the original fixed fibre by
\[
 b^\circ_\theta(v)=
  L^\circ_\theta(\pi,v_*)^{-1}L^\circ_\theta(\pi,v),
 \quad v_*=(0,0,0,0,0,0,-1).
                                                        \tag{NW.3}\label{p1b:eq:natural-nw-3}
\]
Every value fixes the original ordered pair, so lies in
\(\varrho(SU(2))\), and the based value is exactly identity.
At \(\theta=0\), it is the actual \(b_1\) in \eqref{p1b:eq:natural-nv-16}.
Thus substitution of the actual \(v=C\mu(z)\) gives a based homotopy
starting exactly at \(J_\tau(z)\). No meridian or north-frame convention
is silently changed. The actual moving \(\tau\)-hemisphere seam is
\(\cos\theta\cos t-\sin\theta v_7\sin t=0\); the replay retains
this seam when differentiating the piecewise formula.

\subsubsection{The characteristic endpoint and its retained fibre factor}

For \(b=(y,z)\in S^4\subset\mathbb R\times\mathbb C^2\), define
\[
 \vartheta(b)=\frac\pi2(y+1),\qquad
 \phi(b)=\cos\vartheta\,I_2+
 i\sin\vartheta\left(2\frac{zz^*}{|z|^2}-I_2\right).
                                                        \tag{NW.4}\label{p1b:eq:natural-nw-4}
\]
Its two pole values are \(\phi(-1,0)=I_2\) and
\(\phi(1,0)=-I_2\). This specifies a whole continuous map; at a pole
\(|\sin\vartheta|\) tends to zero uniformly in the absent direction.
Use the actual fibre conversion from \eqref{p1b:eq:natural-nv-15}, written on complex
two-blocks as
\[
 M=\begin{pmatrix}0&i\\i&0\end{pmatrix},\qquad
 \psi(b)=M^{-1}\phi(b)M,
 \quad
 \varrho(q)_{\mathbb C^2}=M\operatorname{block}_{SU(2)}(q)M^{-1}.
                                                        \tag{NW.5}\label{p1b:eq:natural-nw-5}
\]
Thus \(\psi\), inserted in the actual \(Q_s\) block, is in the
original cap frame. In particular this conversion cannot be omitted
when the exterior frame is compared with the Spin(5) fibre.

Set \(c=G_\pi p_1\). Its exact first complex column is \(e_2\),
so \(c\in\varrho(SU(2))\). The restriction of the trace connection
to a constant first column is precisely the trace connection of
\(SU(2)\to SU(3)\to S^5\): a generator fixing \(e_1\) lies in
\(\mathfrak{su}(3)\), and the metric differs only by a scalar factor.
The explicit geodesic lift of \eqref{p1b:eq:natural-nw-2}, evaluated at \(t=\pi\), is
\(\operatorname{diag}(-1,i\phi(\beta_\tau(v)))\).
At the actual based direction,
\[
 \beta_\tau(v_*)=(0,i,0)=:\beta_*,\qquad
 \phi(\beta_*)=\operatorname{diag}(i,-i),\qquad
 \widetilde\gamma_{\beta_*}(\pi)=G_\pi.
\]
Uniqueness and right equivariance of horizontal lift now give the
pointwise endpoint identity
\[
 b^\circ_{\pi/2}(v)=
 c^{-1}\operatorname{diag}\bigl(1,
            \phi(\beta_*)^{-1}\phi(\beta_\tau(v))\bigr)c,
                                                        \tag{NW.6}\label{p1b:eq:natural-nw-6}
\]
where the displayed \(SU(3)\) matrix is interpreted in the retained
\(G_2\) representation. This proves the characteristic identity for
the actual input and actual north frame, rather than a class identity.

Here is an explicit removal of the constant conjugation. Write
\(q_c=\varrho^{-1}c\). Select its sign so
\(\operatorname{Re}q_c\ge0\), put
\(q_s=\exp((1-s)\log q_c)\), and conjugate the right-hand characteristic
factor in \eqref{p1b:eq:natural-nw-6} by \(\varrho(q_s)\). The logarithm is defined on this
closed positive-real hemisphere; at the real unit quaternion the
path is constant. A sign change in \(q_c\) leaves the initial
conjugation unchanged. This constructs the based homotopy removing
the actual \(c\); it does not assume \(c=I\).

Finally rotate the real \((y,\operatorname{Im}z_0)\) plane of \(S^4\)
by \(B_\theta(y,\operatorname{Im}z_0)=
(\cos\theta\,y-\sin\theta\,\operatorname{Im}z_0,
 \sin\theta\,y+\cos\theta\,\operatorname{Im}z_0)\).
Use the entire based family
\[
 M^{-1}\phi(B_\theta\beta_*)^{-1}
                  \phi(B_\theta\beta_\tau(C\mu(z)))M.
                                                        \tag{NW.7}\label{p1b:eq:natural-nw-7}
\]
At \(\theta=0\) it is exactly the previous endpoint in the original
cap block. At \(\theta=\pi/2\), its based point is the south pole.
Consequently \eqref{p1b:eq:natural-nw-1}--\eqref{p1b:eq:natural-nw-7} construct and apply a complete based leg
\[
 \mathcal L_\tau:J_\tau\longrightarrow\psi(U_\tau),\qquad
 U_\tau(z)=
 \left(r_v\tau_0(n_v),
 r_v(\tau_1(n_v)-i\tau_2(n_v)),
 v_6+i r_v\tau_3(n_v)\right),\quad v=C\mu(z).
                                                        \tag{NW.8}\label{p1b:eq:natural-nw-8}
\]
At the actual collapsed cap wedge it has value \((-1,0,0)\),
so the entire leg is identity on \(\partial\mathbb B\).
The source is the original \(\mu\), not a replacement selected only
for its degree.

\subsubsection{The actual exterior leg and its immediate application}

Keep the actual \(g(v)=(Q_s^*\mathsf V^\#)_2\) on
\(\partial\mathbb B\), including all transporters and their coordinates.
For \(z=R_Brv\), define
\[
 U_g(z)=\left(1-2r,\,2\sqrt{r(1-r)}\,Mg(v)e_1\right)\in S^4,
 \quad U_g(0)=(1,0),\quad U_g|_{\partial\mathbb B}=(-1,0).
                                                        \tag{NW.9}\label{p1b:eq:natural-nw-9}
\]
The uniform factor proves both polar limits. The left factor \(M\)
is essential: these are the original exterior columns converted to
the target fibre of \eqref{p1b:eq:natural-nw-4}. Set \(A(q)=\operatorname{diag}(e^{\pi iq},
e^{-\pi iq})\). Direct multiplication gives
\[
 \phi(U_g)=MgA(1-r)g^{-1}M^{-1},\qquad
 MA(q)M^{-1}=A(q)^{-1}.
                                                        \tag{NW.10}\label{p1b:eq:natural-nw-10}
\]
Hence the following is an actual based homotopy in the original cap block:
\[
 E_\ell(z)=M^{-1}\phi(U_g(z))A((1-\ell)(1-r))M,
 \qquad E_0=E_g,\quad E_1=\psi(U_g).
                                                        \tag{NW.11}\label{p1b:eq:natural-nw-11}
\]
It is identity on the entire cap boundary for every \(\ell\).
At the center its value is
\(M^{-1}(-I)A(1-\ell)M\), a specified continuous constant.
No fixed-center requirement is imposed by the based cap quotient.

Apply this leg to the actual side/top background immediately:
\[
 \mathcal Q_g(t,z)=E_{1-t}(z)E_g(z)^{-1}G_0(t,z)A(t).
                                                        \tag{NW.12}\label{p1b:eq:natural-nw-12}
\]
Its actual three faces are
\[
 \mathcal Q_g(0,z)=\psi(U_g(z)),\quad
 \mathcal Q_g(1,z)=-I_2,\quad
 \mathcal Q_g(t,R_Bv)=g(v)A(t)g(v)^{-1}.
                                                        \tag{NW.13}\label{p1b:eq:natural-nw-13}
\]
All equations follow from \eqref{p1b:eq:natural-nu-3} and the boundary values of \eqref{p1b:eq:natural-nw-11}.
Thus this is a constructed rank-two side/top background with exactly
the retained exterior seam. It starts at \(\psi(U_g)\), not at \(H^6\).

Likewise concatenate the completed \(H^6\to D^6\mu\to J_\tau\)
with \(\mathcal L_\tau\). Replace the bottom of the complete normalized
cylinder boundary by this actual path, retaining its top identity,
side \(S_g\), both seams, \(Q_s\) and stabilizing column. This ends
with actual bottom \(\psi(U_\tau)\). The matching toggle \(T_Y^1\)
and the actual joining outside the cap are unchanged.

\paragraph{The auxiliary comparison criterion.}
The two entire source maps \eqref{p1b:eq:natural-nw-8}--\eqref{p1b:eq:natural-nw-9} are explicit. A separate
source-level bridge for this auxiliary route would be the based map
\[
 \mathcal C:[0,1]\times(\mathbb B/\partial\mathbb B)\longrightarrow S^4,
 \quad \mathcal C(0,z)=U_\tau(z),\quad
 \mathcal C(1,z)=U_g(z),\quad
 \mathcal C(\ell,\partial\mathbb B)=(-1,0,0).
                                                        \tag{NW.14}\label{p1b:eq:natural-nw-14}
\]
Equality of their classes, or the primitive value of \(\phi\), is not
a construction of \eqref{p1b:eq:natural-nw-14}. In particular the displayed two legs do
not by themselves complete \eqref{p1b:eq:natural-nv-18}.

Their exact assembly is nevertheless fixed. A supplied \(\mathcal C\)
would produce \eqref{p1b:eq:natural-nv-18} by concatenating
\(\mathcal L_\tau\), \(\psi\mathcal C\), and the reverse of
\eqref{p1b:eq:natural-nw-11}. Equivalently concatenate the actual bottom path already
constructed with \(\psi\mathcal C\), calling it \(\mathcal B\), and use
\[
 \mathcal X(t,z)=\mathcal B_{\min(4t,1)}(z)
           \psi(U_g(z))^{-1}\mathcal Q_g(t,z).
                                                        \tag{NW.15}\label{p1b:eq:natural-nw-15}
\]
Its prescribed faces would be exactly \(H^6,-I,gA(t)g^{-1}\).
The only unsupplied part of \eqref{p1b:eq:natural-nw-15} is \eqref{p1b:eq:natural-nw-14}; the formula is not
reported as a completed cap closing.

The complete original comparison is evaluated in
\eqref{p1b:eq:natural-aj-4}--\eqref{p1b:eq:natural-aj-7}:
\(h_t=w_t=1\), \(r_J=\zeta_{\rm nat}=0\), with every retained
extension test decided. This component is used with its full frames,
lower divisions and joining; its isolated value is not a replacement
for that complete comparison.

\paragraph{Primary source and actual application.}
T. P\"uttmann, \emph{On the commutator of unit quaternions and the
numbers 12 and 24}, arXiv:1101.5147v4 (2012), \S4.2,
Lemma 4.5, and \S5.1, Lemma 5.1 and Theorem 5.2, give the meridian
lifting identity and the displayed special-unitary characteristic map.
The source rotation, actual north frame, based division, fixed fibre
conversion and both cap-source substitutions above are part of the
present actual-map application.



\subsection{The complete stabilized determinant target and its actual column defect}
\label{p1b:sec:natural-determinant-stabilized-closing}

Keep every actual source, frame, root, exterior, toggle and joining of
\eqref{p1b:eq:natural-nd-1}--\eqref{p1b:eq:natural-nw-15}. The original target remains \(h_t,w_t\), with
\(w_t=1\) already evaluated. The construction here completes a specific
stabilized target, applies its whole boundary restriction, and identifies
its entire moving-column map. It does not replace the unsupplied
rank-two comparison \eqref{p1b:eq:natural-nw-14} by a primitive value.

\paragraph{Dependencies executed.}
Construct the based special-unitary filling of the retained characteristic
map in the exact cap fibre; apply it to both actual sources to bridge the
bottoms; fill all prescribed cap faces and paste the exterior; form the
stabilized determinant gauge with the retained root and matching toggle;
restrict that complete triple, and its sphere quotient, along the retained
boundary map; identify the whole projected column and the precise repair
it requires. Every construction is applied below.

\subsubsection{An explicit filling in the actual stabilized cap}

For \(b=(y,z)\in S^4\subset\mathbb R\times\mathbb C^2\),
\(0\le\theta\le\pi\), put \(P_z=zz^*/|z|^2\) when \(z\ne0\),
and define
\[
 L_\theta(b)=
 \begin{pmatrix}
  \cos\theta+iy\sin\theta&
               -\sin\theta\,e^{iy\theta/2}z^*\\
  \sin\theta\,z&
  e^{-iy\theta/2}(I-P_z)+
          e^{iy\theta/2}(\cos\theta-iy\sin\theta)P_z
 \end{pmatrix}.                                         \tag{NX.1}\label{p1b:eq:natural-nx-1}
\]
At either pole use
\(L_\theta(\pm1,0)=\operatorname{diag}
(e^{\pm i\theta},e^{\mp i\theta/2},e^{\mp i\theta/2})\).
This is the entire horizontal geodesic lift in the fixed
\(SU(2)\to SU(3)\to S^5\) convention, not a new detector.
Its bottom-right direction-dependent coefficient is
\[
 e^{iy\theta/2}(\cos\theta-iy\sin\theta)-e^{-iy\theta/2}.
\]
It vanishes at \(y=\pm1\). Division by \(1-y^2\) therefore
has finite one-sided limits, uniformly for \(\theta\in[0,\pi]\).
Multiplication by \(zz^*\) gives the claimed continuous, in fact analytic,
polar extension. Direct multiplication gives \(L_\theta\in SU(3)\),
its first column \((\cos\theta+iy\sin\theta,\sin\theta z)\),
and the exact endpoint
\(L_\pi(b)=\operatorname{diag}(-1,i\phi(b))\).

Use the actual \(M\) from \eqref{p1b:eq:natural-nw-5}. Let \(P\) be the determinant-one
cyclic permutation with columns \((e_3,e_1,e_2)\), and set
\(V=P\operatorname{diag}(1,M^{-1})\). In the retained cap order
(two fibre axes, then the prescribed stabilizing axis), define
\[
 \Lambda_\theta(b)=
 V L_\theta(-1,0)^{-1}L_\theta(b)V^{-1}.                \tag{NX.2}\label{p1b:eq:natural-nx-2}
\]
Every factor is specified. It obeys
\[
 \Lambda_0(b)=I_3,\quad
 \Lambda_\pi(b)=\operatorname{diag}(\psi(b),1),\quad
 \Lambda_\theta(-1,0)=I_3.                             \tag{NX.3}\label{p1b:eq:natural-nx-3}
\]
Thus it is a based stabilized filling of the \emph{actual} characteristic
map with its original fibre conversion. This construction does not
assert a filling in \(SU(2)\).

Let \(\mathcal B^{\rm old}:H^6\to\psi(U_\tau)\) denote the
complete actual bottom homotopy already constructed in (NU)--(NW).
Apply \eqref{p1b:eq:natural-nx-2} to its endpoint and to the actual exterior source:
\[
 \mathcal B^{\rm st}_r(z)=
 \begin{cases}
  \mathcal B^{\rm old}_{3r}(z),&0\le r\le1/3,\\
  \Lambda_{\pi(2-3r)}(U_\tau(z)),&1/3\le r\le2/3,\\
  \Lambda_{\pi(3r-2)}(U_g(z)),&2/3\le r\le1.
 \end{cases}                                           \tag{NX.4}\label{p1b:eq:natural-nx-4}
\]
Its endpoints are exactly \(H^6\) and \(\operatorname{diag}
(\psi(U_g),1)\). Its two seams are respectively
\(\operatorname{diag}(\psi(U_\tau),1)\) and \(I_3\).
Because both actual sources equal the south pole on \(\partial\mathbb B\),
this entire path is identity there. The comparison has been constructed
on the actual source, rather than selected only for its class.

Apply it to the already constructed \(\mathcal Q_g\) of \eqref{p1b:eq:natural-nw-12}:
\[
 \mathcal X^{\rm st}(t,z)=
 \mathcal B^{\rm st}_{\min(4t,1)}(z)
 \operatorname{diag}(\psi(U_g(z))^{-1},1)
 \operatorname{diag}(\mathcal Q_g(t,z),1).              \tag{NX.5}\label{p1b:eq:natural-nx-5}
\]
Here and below \(\mathcal Q_g\) denotes its two-block when inserted
in a displayed diagonal. The complete three faces are
\[
\begin{gathered}
 \mathcal X^{\rm st}(0,z)=(Q_s^*H^6Q_s)(z),\qquad \mathcal X^{\rm st}(1,z)=\operatorname{diag}(-I_2,1), \\
 \mathcal X^{\rm st}(t,R_Bv)=
       \operatorname{diag}(g(v)A(t)g(v)^{-1},1).
\end{gathered}
 \tag{NX.6}\label{p1b:eq:natural-nx-6}
\]
This is an actual \(SU(3)\) filling of the complete actual cap-boundary
map \eqref{p1b:eq:natural-nt-10}. The specified northern frame, both fibre axes and every side
coordinate have been retained. In particular there is no guessed
\(U_\tau\to U_g\) leg in \eqref{p1b:eq:natural-nx-4}.

\subsubsection{The actual stabilized triple and its complete restriction}

Paste immediately in the original ambient bundle:
\[
 \mathcal K^{\rm st}(t,x,p)=
 \begin{cases}
  Q_s\mathcal X^{\rm st}(t,z)Q_s^*,&(x,p)\in\mathbb B,\\
  K^\#(t,x,p),&(x,p)\notin\mathbb B.
 \end{cases}                                           \tag{NX.7}\label{p1b:eq:natural-nx-7}
\]
Equations \eqref{p1b:eq:natural-nx-6} give the exact common seam. It has endpoints
\(H^6,N\) on all of \(P\), and still equals \(K_+\) on the
joining locus, which is disjoint from \(\mathbb B\).
It acts on the actual stabilized bundle \(E_s\oplus\mathbb Cs\).
It need not fix \(s\) in the cap interior.

On the first primitive cut \(a=e^{\pi it}\), define the actual gauge
\[
 D_Y^{\rm st}(a,x,p)=
 (ss^*+aP_s)\mathcal K^{\rm st}(t,x,p)H^{-3}(x,p).
                                                        \tag{NX.8}\label{p1b:eq:natural-nx-8}
\]
It has determinant exactly \(a^2\); its first and last reduced
representatives are \(H^3,NH^{-3}\). The full gauge values at the
cut endpoints are \(H^3,H^{-3}\), since
\((ss^*-P_s)N=I\). Therefore extend to the second half by
\[
 D_Y^{\rm st}(-a,p,x)=\mathsf J
                      D_Y^{\rm st}(a,x,p)\mathsf J.     \tag{NX.9}\label{p1b:eq:natural-nx-9}
\]
The established equation \(\mathsf J H(p,x)\mathsf J=H(x,p)^{-1}\)
proves both cut seams exactly. No symmetry of the interior path at
intermediate times is assumed. This constructs a full stabilized gauge
on the actual mapping torus. Apply the already prescribed matching toggle:
\(\widehat D_Y^{\rm st}=D_Y^{\rm st}T_Y^1\).

On the joining locus \eqref{p1b:eq:natural-nx-7} equals the original \(K_+\), so
\(\widehat D_Y^{\rm st}=D_J\) there. Consequently the original
constructed joining \(H_J\) of \eqref{p1b:eq:natural-nd-8}--\eqref{p1b:eq:natural-nd-10}, with its actual
odd clutch, pastes this gauge to the actual
\(D_Z=\operatorname{diag}(1,-a^2,1)\).
The result is a complete stabilized triple on the actual bundle; there
is no unfilled cap in it. Its rank-two sphere gauge is precisely the
already constructed \(S_Y,S_Z,H_J^S\), with \(w_t=1\).
Their stabilized composite is likewise an actual complete triple,
\(C^{\rm st}=D^{\rm st}S^{-1}\), with the existing joining replaced
by \(H_J(H_J^S)^{-1}\).

Here is its complete boundary restriction, with the fixed base coordinates
and all actual ordinary identifications. Put \(a=ie^{is}\),
\(w=yu\), \(\alpha=3\pi\rho/2\), and retain the \(F_Y,R,F_Z\)
of \eqref{p1b:eq:natural-nc-4}--\eqref{p1b:eq:natural-nc-12}. Define \(F_N=F_Y(y,u,\pi/2)R(s,y,1,u)\).
Then
\[
\widetilde D^{\rm st}(s,y,\rho,u)=
 \begin{cases}
 \begin{gathered}R(s,y,3\rho,u)^*F_Y(y,u,\alpha)^*
       \widehat D_Y^{\rm st}(a,x,p)\\
       F_Y(y,u,\alpha)R(s,y,3\rho,u)\end{gathered},&0\le\rho\le1/3,\\
 F_N^*H_J(3\rho-1,a,w)F_N,&1/3\le\rho\le2/3,\\
 F_Z(y,u)^*D_Z(a)F_Z(y,u),&2/3\le\rho\le1,
 \end{cases}
\tag{NX.10}\label{p1b:eq:natural-nx-10}
\]
where \(x=ye^{\alpha u},p=ye^{-\alpha u}\) in the inner stage.
The first expression is now defined also on the full actual preimage
of \(\mathbb B\). At the first seam it is \(F_N^*D_JF_N\).
At the second use the actual equality \(F_N=c_jF_Z\), giving
exactly the outer expression. At \(\rho=0\) the original inner
extension is independent of the absent polar direction; at \(\rho=1\)
the retained ordering-line identification supplies the \(u\sim-u\)
descent. These are the old base coordinates, without a normal rotation.
The joining phase changes the determinant from \(a^2=-e^{2is}\)
to \(-a^2=e^{2is}\). Its actual normalization is constructed next.

The actual sphere restriction \(\widetilde S\) is unchanged. Thus
\[
 \widetilde C^{\rm st}=\widetilde D^{\rm st}
                              (\widetilde S)^{-1}        \tag{NX.11}\label{p1b:eq:natural-nx-11}
\]
is also specified on all three stages. This is a stabilized restriction.

\subsubsection{The applied literal positive-eigenline circle comparison}

Parameterize the three already constructed \eqref{p1b:eq:natural-nd-8} joining stages by equal
thirds. This only fixes the time parametrization of that actual homotopy.
Its first stage has determinant \(a^2\), its second stage has determinant
\(e^{\pi i(3j-1)}a^2\), and its third stage has determinant
\(-a^2\), with \(j=3\rho-1\). Thus, on the entire retained boundary,
\[
 \det\widetilde D^{\rm st}
     =-e^{2is}e^{\pi i f(\rho)},\qquad
 f(\rho)=\min(1,\max(0,9\rho-4)).                     \tag{NX.12}\label{p1b:eq:natural-nx-12}
\]
The sphere determinant is one, so this is also the composite determinant.
In the actual \(\mathbf1\oplus\lambda_L\oplus\mathbf1\) order set
\[
 B_\rho=\operatorname{diag}
        (e^{\pi i(1-f(\rho))},1,1),\qquad
 Q_2(s)=\operatorname{diag}(e^{2is},1,1).               \tag{NX.13}\label{p1b:eq:natural-nx-13}
\]
Both matrices are actual automorphisms in the retained boundary bundle:
the first axis is the positive trivial eigenline and is fixed by the
ordering-line descent. The first matrix has the displayed global real
logarithm; multiplying its exponent by \(\ell\in[0,1]\) is an
explicit ordinary bundle homotopy from identity to \(B_\rho\).
There is no change of base normal coordinates or bundle identification.
The second matrix is literally the natural positive-eigenline insertion
\(2q_+\). It is not \(e^{is}I_2\), whose retained component
is \(2q_+-\kappa\).

Apply these comparisons immediately to the complete restrictions:
\[
 M_D^{\rm st}=Q_2^{-1}B_\rho\widetilde D^{\rm st},
 \qquad M_C^{\rm st}=Q_2^{-1}B_\rho\widetilde C^{\rm st}.
                                                        \tag{NX.14}\label{p1b:eq:natural-nx-14}
\]
Equations \eqref{p1b:eq:natural-nx-12}--\eqref{p1b:eq:natural-nx-13} prove that both determinants are exactly one
on every inner, middle and outer point, with both seams retained.
This supplies the actual determinant-function normalization and the
literal \(2q_+\) comparison on the complete stabilized targets.
The retained fourth-stage correction and the actual \(3z\) or
\(-2b\) divisions are carried through the complete (OA), (OB),
(CI), (AJ) comparison; determinant normalization alone supplies
no value for those maps.

\subsubsection{The whole column defect and its relative repair criterion}

In the retained cap order the last column of \eqref{p1b:eq:natural-nx-2} is the explicit map
\[
 \xi_\theta(y,z)=\left(
 \sin\theta\,e^{-i\theta/2}M^{-1}z,
 e^{i\theta}(\cos\theta+iy\sin\theta)\right)\in S^5.
                                                        \tag{NX.15}\label{p1b:eq:natural-nx-15}
\]
It is \(e_3\) at \(\theta=0,\pi\), and at \(b=(-1,0)\)
for every \(\theta\). The entire projection of the actual cap
filling is
\[
 \Xi(t,z)=\mathcal X^{\rm st}(t,z)e_3=
 \begin{cases}
 e_3,&t\le1/12\text{ or }t\ge1/4,\\
 \xi_{\pi(2-12t)}(U_\tau(z)),&1/12\le t\le1/6,\\
 \xi_{\pi(12t-2)}(U_g(z)),&1/6\le t\le1/4.
 \end{cases}                                           \tag{NX.16}\label{p1b:eq:natural-nx-16}
\]
The right two factors of \eqref{p1b:eq:natural-nx-5} fix the last column; the old bottom
path does also. This proves the equation for the full map, including
its constant complement. Its boundary is exactly constant, so it is
an actual map \(\partial\)-quotiented \(I\times\mathbb B\simeq S^7\to S^5\).

This column has also been applied to the \emph{full} actual boundary
restriction. On the first cut, inside the actual cap preimage,
\[
 \widetilde D^{\rm st}e_3=\widetilde C^{\rm st}e_3
   =(F_YR)^*Q_s\operatorname{diag}(a,a,1)\Xi(t,z).
                                                        \tag{NX.17}\label{p1b:eq:natural-nx-17}
\]
The root and the matching toggle on the right fix the actual section;
the left half phase has exactly the displayed action in \(Q_s\).
On the second cut use \eqref{p1b:eq:natural-nx-9} on the swapped pair and apply the actual
\(\mathsf J\), retaining that pair's \(Q_s\). Outside the cap
preimage, on the entire joining cylinder, and on the outer stage the
column is exactly \(e_3\). After the completed positive-line circle
comparison its value is \(Q_2^{-1}B_\rho\) times \eqref{p1b:eq:natural-nx-17}.
The sphere quotient fixes the prescribed column, so both normalized
targets have this identical complete column formula. In particular no
basepoint frame, half phase or deck factor is discarded on passing from
the cap projection to the actual boundary map.

The map \(\xi\) itself has suspension degree of absolute value one.
Indeed its entire inverse image of \(-e_3\) is
\(\theta=\pi/2,b=(1,0)\): its first two components require either
a polar value or an endpoint, the latter giving \(+e_3\); the south
pole gives \(+e_3\), and the north pole gives \(e^{2i\theta}e_3\).
In the north \(S^4\) chart the four \(z\)-normals have differential
\(e^{-i\pi/4}M^{-1}\), an orientation-preserving real isometry,
and the remaining \(\theta\)-normal has derivative \(-2i\).
This identifies the whole degree-one quotient map, not only a
benchmark value. Consequently \eqref{p1b:eq:natural-nx-16} is the actual projected
input of class
\[
 {}[\Xi]=\pm\bigl(-E[U_\tau]+E[U_g]\bigr)=0
                         \quad\text{in }\pi_7S^5.      \tag{NX.18}\label{p1b:eq:natural-nx-18}
\]
The endpoint class equality is the established exotic clutch comparison;
its use here proves the class of this \emph{constructed} projection.
It is not a construction of a nullhomotopy of \(\Xi\).

This filling really moves the prescribed column. At \(z=0\),
\(U_g=(1,0)\). Thus at \(t=5/24\) the ascending leg has
\(\theta=\pi/2\) and \(\Xi=-e_3\). In particular
\(\mathcal K^{\rm st}s\ne s\), so it cannot already serve as
the rank-two closing in \eqref{p1b:eq:natural-nd-11}.

The auxiliary column repair criterion is the relative homotopy
\[
 \mathcal A:[0,1]\times I\times\mathbb B\to S^5,
 \quad \mathcal A(0,t,z)=\Xi(t,z),\quad
 \mathcal A(1,t,z)=e_3,\quad
 \mathcal A(\ell,\partial(I\times\mathbb B))=e_3.      \tag{NX.19}\label{p1b:eq:natural-nx-19}
\]
There is no additional untracked axis choice: a supplied actual \eqref{p1b:eq:natural-nw-14}
gives the following precise construction of this repair. For
\(0\le\ell\le1/2\), replace \(U_\tau\) by
\(\mathcal C(2\ell,z)\) in the descending leg of \eqref{p1b:eq:natural-nx-16},
keeping the ascending \(U_g\) leg fixed. At \(\ell=1/2\)
the two legs are exact reverse paths on the same source. For
\(1/2\le\ell\le1\), replace either leg's angle \(\theta(t)\)
by
\[
 \theta_\ell(t)=\pi+(2-2\ell)(\theta(t)-\pi).          \tag{NX.20}\label{p1b:eq:natural-nx-20}
\]
This contracts the retracing pair to \(\theta=\pi\); all three
seams and all boundary points stay fixed. This is an adapter for a supplied \(\mathcal C\), not an assumption
that such a source-level comparison has been constructed. The
intrinsic bridge actually used is \eqref{p1b:eq:natural-ny-1}--\eqref{p1b:eq:natural-nz-20}.

Given such a repair, lift its \(\ell\)-curves horizontally in
\(SU(3)\to S^5\), starting at the actual frame
\(\mathcal X^{\rm st}(t,z)\), whose last column is \(\Xi\).
Call the lifted frames \(Z_\ell(t,z)\), and set
\(T_\ell=Z_\ell(\mathcal X^{\rm st})^{-1}\). Then
\(T_0=I_3\) and \(T_\ell\Xi=\mathcal A(\ell,t,z)\).
The constant boundary column has constant horizontal lift starting at
its actual boundary frame, so \(T_\ell\) is identity on the whole
cap boundary. Then \(T_1\mathcal X^{\rm st}\) fixes
\(e_3\) and has exactly the faces \eqref{p1b:eq:natural-nx-6}; pasting it in \eqref{p1b:eq:natural-nx-7}
would supply the original rank-two closing. This states the required
repair equation and its exact application to the actual completed
stabilized filling. It does not assert that \eqref{p1b:eq:natural-nx-19} has been evaluated.

The complete original comparison is evaluated in
\eqref{p1b:eq:natural-aj-4}--\eqref{p1b:eq:natural-aj-7}:
\(h_t=w_t=1\), \(r_J=\zeta_{\rm nat}=0\), with every retained
extension test decided. This component is used with its full frames,
lower divisions and joining; its isolated value is not a replacement
for that complete comparison. The seven-sphere column \eqref{p1b:eq:natural-nx-16} and the eight-sphere detector input have different domains; a zero class for the former alone supplies no natural coefficient. The rank-two bridge actually used is \eqref{p1b:eq:natural-ny-1}--\eqref{p1b:eq:natural-nz-20}, followed by \eqref{p1b:eq:natural-oa-1}--\eqref{p1b:eq:natural-oa-4}, rather than an assumed separate \eqref{p1b:eq:natural-nw-14} source homotopy.

\paragraph{Preserved primary reference.}
T. P\"uttmann, \emph{On the commutator of unit quaternions and the
numbers 12 and 24}, arXiv:1101.5147v4 (2012), \S5.1,
Lemma 5.1 and Theorem 5.2. Formula \eqref{p1b:eq:natural-nx-1} is its horizontal lift
in dimension five. The actual south normalization, retained cap fibre
conversion, two actual-source applications, pasted target and entire
projected restriction are the present calculations.



\subsection{An applied quaternion endpoint comparison for the determinant bottom}
\label{p1b:sec:natural-quaternion-repair}

Retain the actual cap, its frame \(Q_s\), its source \(C_\mu\), and
the already constructed bottom homotopy \(Q_s^*H^6Q_s\to D^6C_\mu\).
The following replaces the later Stiefel path by a specified path with an
explicit endpoint comparison. It does not replace the actual cap source.
Write a point of \(S^5\) as \(x=(a,b,x_2,x_3,x_4,x_5)\), and set
\[
 L=\|(x_2,x_3,x_4,x_5)\|,\quad
 q=L^{-1}(x_4,-x_5,-x_3,x_2),\quad
 e_1=qiq^{-1},\ e_2=qjq^{-1},\quad u=b+Le_1.             \tag{NY.1}\label{p1b:eq:natural-ny-1}
\]
Missing \(q\) at \(L=0\) is used only with a vanishing coefficient.
The retained homogeneous perpendicular Hopf pair is \((Le_1,Le_2)\).

\begin{proposition}[A complete replacement of the five-dimensional frame path]
\label{p1b:proposition:natural-quaternion-frame}
Put \(c=2a^2-1,d=2ab,r=2|a|L\). For \(a<0\), rotate the Hopf pair
through \(\theta=\pi\lambda\), while for \(a\ge0\) put \(\theta=0\):
\[
 \gamma=(c,d,r(\cos\theta e_1+\sin\theta e_2)),\qquad
 \delta=(d,-c,r(-\sin\theta e_1+\cos\theta e_2)).         \tag{NY.2}\label{p1b:eq:natural-ny-2}
\]
This starts at the actual \(\widehat hP_2\) frame. It ends at
\(\gamma_1=(2a^2-1,2au)\) and
\(\delta_1=(2ab,1-2a^2+2aLe_2)\).

For a unit quaternionic column \((r_0,U)\) with \(\operatorname{Re}r_0\ge0\),
the following matrix is unitary and has that first column:
\[
 R(r_0,U)=
 \begin{pmatrix}
 r_0&-(1+r_0)(1+\overline r_0)^{-1}\overline U\\
 U&1-U(1+\overline r_0)^{-1}\overline U
 \end{pmatrix}.                                         \tag{NY.3}\label{p1b:eq:natural-ny-3}
\]
Let \(e=(i,0)\), \(E=\operatorname{diag}(i,1)\), and
\[
 v_t=\cos t\,(a,\overline u)+\sin t\,e,\quad
 A_t=E R(E^*v_t)E^*,\qquad 0\le t\le\pi/2.              \tag{NY.4}\label{p1b:eq:natural-ny-4}
\]
Then \(A_te=v_t\), \(A_{\pi/2}=I\). Encode a vector \((c,w)\in\mathbb R^5\)
by \(\Gamma(c,w)=\left(\begin{smallmatrix}c&w\\\bar w&-c\end{smallmatrix}\right)\).
The pair
\[
 \Gamma(\gamma_t)=2v_tv_t^*-I,\qquad
 \Gamma(\delta_t)=A_tA_0^*\Gamma(\delta_1)A_0A_t^*         \tag{NY.5}\label{p1b:eq:natural-ny-5}
\]
contracts the first vector to north. Finally multiply the quaternionic
tail of the second vector by \(\exp(ti)\), \(0\le t\le\pi/2\).
These three specified legs give a full \(V_{5,2}\) homotopy with endpoint
\((e_0,(0,\tau))\), where
\[
 q_2={2a-\overline uiu\over1+a^2},\qquad
 k=2L^2-1+2Lq(-bi+aj)\bar q,\qquad
 \boxed{\tau=-kq_2^{-1}.}                                \tag{NY.6}\label{p1b:eq:natural-ny-6}
\]
Its pole values are \(\tau(e_0)=1\), \(\tau(-e_0)=-1\).
\end{proposition}
\begin{proof}
The two Hopf vectors are perpendicular unit imaginary quaternions.
Thus \eqref{p1b:eq:natural-ny-2} has unit perpendicular columns because
\(c^2+d^2+r^2=1\). At \(a=0\) the entire rotating tail vanishes,
which gives the required hemisphere seam. Direct multiplication proves
the unitary identity in \eqref{p1b:eq:natural-ny-3}; its denominator has norm at least one.
The two spinors in \eqref{p1b:eq:natural-ny-4} are real-orthogonal, so \(v_t\) is unit.
Quaternionic conjugation in \eqref{p1b:eq:natural-ny-5} preserves the trace form on the
Hermitian matrices and gives the claimed frame path.

The endpoint is identified by the following pointwise identities, with
\(F=\left(\begin{smallmatrix}a&-u\\\bar u&a\end{smallmatrix}\right)\):
\[
 A_0^*F=\operatorname{diag}(i,q_2),\qquad
 F^*\Gamma(\delta_1)F=\Gamma(0,k).                        \tag{NY.7}\label{p1b:eq:natural-ny-7}
\]
Consequently the tail before the final rotation is \(ikq_2^{-1}\);
multiplying by \(i\) gives \eqref{p1b:eq:natural-ny-6}. This is an identity of the actual
endpoint, not a homotopy-group identification by its order.
Substitution at both poles gives the stated values.
\end{proof}

\begin{proposition}[An explicit nullhomotopy of the actual extra factor]
\label{p1b:proposition:natural-quaternion-factor}
The factor \(k\) in \eqref{p1b:eq:natural-ny-6} has a specified homotopy to \(-1\), fixing
both poles. Put \(R=\sqrt{a^2+b^2}\). First replace
\((-bi+aj)/R\) by
\[
 \cos\theta\,(-bi+aj)/R+\sin\theta\,k,\quad
 0\le\theta\le\pi/2,                                    \tag{NY.8}\label{p1b:eq:natural-ny-8}
\]
in \(2L^2-1+2LRq(\cdot)\bar q\). Then replace \(L\) by
\((1-\lambda)L\) in
\[
 2L^2-1+2L\sqrt{1-L^2}\,qk\bar q.                        \tag{NY.9}\label{p1b:eq:natural-ny-9}
\]
The resulting actual endpoint comparison is
\[
 \tau_\lambda=-k_\lambda q_2^{-1},\qquad
 \tau_0=\tau,\quad \tau_1=q_2^{-1}.                       \tag{NY.10}\label{p1b:eq:natural-ny-10}
\]
\end{proposition}
\begin{proof}
The directions in \eqref{p1b:eq:natural-ny-8} are orthogonal unit imaginary quaternions.
The norm identity follows from \(L^2+R^2=1\).
At \(R=0\) the first rotating coefficient vanishes; at \(L=0\)
every missing \(q\) coefficient vanishes. The second formula is a map
of the actual four-disc \(Lq\), contracted radially to zero.
Its scalar endpoint is \(-1\). This proves all seams, pole conditions
and endpoint equations, and immediately supplies \eqref{p1b:eq:natural-ny-10}.
\end{proof}

\paragraph{Application to the actual bottom and its moving lifts.}
Use the first three legs of the retained \(V_{7,2}\) path unchanged.
Use the double lift of \eqref{p1b:eq:natural-ny-2}--\eqref{p1b:eq:natural-ny-5} and its final tail rotation for the
fourth leg, and the retained last-two-coordinate contraction for the last
leg. Thus the complete path starts at the same actual \(\kappa P_2\)
and ends at \((e_1,\Sigma_{\rm low}^2\tau)\). Apply it immediately to
the retained meridians
\[
 x(t,z)=(\cos t,\sin t\,C_\mu(z)),\qquad
 C_\mu(z)=(-p_1,p_3,-p_2,-u_1,-u_2,u_3,-u_0).             \tag{NY.11}\label{p1b:eq:natural-ny-11}
\]
Use the same exact minimum-trace horizontal \(G_2\) connection as \eqref{p1b:eq:natural-nv-16}.
The north lift is transported through the whole new path, including its
new fourth leg. Normalize each boundary lift by the lift at
\(v_*=(0,\ldots,0,-1)\). The initial based boundary is the retained
\(\rho(D^6C_\mu)\), by the unchanged initial triality identification.
This constructs an actual bottom homotopy to its new endpoint \(J_{\rm new}\).

Rotate source coordinates \((x_0,x_7)\) and apply the retained target
\(SU(3)\) rotation from \eqref{p1b:eq:natural-nw-1}--\eqref{p1b:eq:natural-nw-2}. Because the two new \(\tau\)
pole values are \((1,-1)\), the north and south frames stay fixed throughout
this rotation. Its north lift is the actual \(c=p_{\rm new}\in\rho SU(2)\),
not identity. At its end the meridian is the actual geodesic with direction
\[
 \beta(v)=(-r\tau_2,\ r(\tau_1-i\tau_0),\ v_6+ir\tau_3),\quad
 r^2=v_7^2+\sum_{j=1}^5v_j^2,\quad
 \tau=\tau((-v_7,v_1,\ldots,v_5)/r).                     \tag{NY.12}\label{p1b:eq:natural-ny-12}
\]
Here \(\beta(v_*)=(0,-i,0)\). The geodesic formula gives the based
boundary exactly as
\[
 c^{-1}\operatorname{diag}(1,\phi(\beta_*)^{-1}\phi(\beta(v)))c.
                                                               \tag{NY.13}\label{p1b:eq:natural-ny-13}
\]
Remove this actual constant conjugation by its quaternionic logarithm
path, retaining its sign; choose the \(i\) axis if the constant is \(-1\).
Rotate \((y,\operatorname{Im}z_0)\) through \(-\pi/2\), carrying the
base value to the south pole. In the retained cap fibre this constructs
\(J_{\rm new}\to\psi(U_{\rm new})\), where
\[
 U_{\rm new}(v)=(-r\tau_0,\ r(\tau_1+i\tau_2),\ v_6+ir\tau_3),
 \qquad \psi=M^{-1}\phi M.                               \tag{NY.14}\label{p1b:eq:natural-ny-14}
\]
Apply \eqref{p1b:eq:natural-ny-10} in \eqref{p1b:eq:natural-ny-14}. Next rotate the source pair \((v_6,v_7)\)
and the target pair \((y,\operatorname{Re}z_1)\) through \(\pi/2\).
The combined rotation fixes the south base value throughout. Its endpoint is
\[
\begin{gathered}
 U_H(v)=\bigl(v_7,R\Theta(q_2^{-1}(n))\bigr),\qquad R=\sqrt{1-v_7^2},\qquad n=(-v_6,v_1,\ldots,v_5)/R, \\
 \Theta(q)=(q_1+iq_2,-q_0+iq_3).
\end{gathered}
 \tag{NY.15}\label{p1b:eq:natural-ny-15}
\]
Every leg is applied to \(v=C_\mu(z)\); \(\partial B\) remains south.
The actual applied bottom path is therefore
\[
 \boxed{Q_s^*H^6Q_s\longrightarrow D^6C_\mu
 \longrightarrow J_{\rm new}\longrightarrow\psi(U_H(C_\mu)).}
                                                               \tag{NY.16}\label{p1b:eq:natural-ny-16}
\]
The preceding equations retain the full lifts, base coordinates, pole
normalization and fibre matrix. No bit is inferred from this endpoint.




\subsection{Actual inverse embeddings, bundle transport and the exterior clutch}
\label{p1b:sec:natural-bundle-comparison}

\paragraph{The retained cap source is compared as a map.}
Let \(B=\{z=(z_x,z_p):|z|\le r_B\}\), \(r_B=1/6\), \(\delta=1/12\), and
\[
 A(z_x,z_p)=(-z_{p1},-z_{p3},z_{p2},z_{x1},z_{x2},-z_{x3}),\qquad
 \nu_A(z)=(\sin(\pi|z|/r_B)Az/|z|,\cos(\pi|z|/r_B)).      \tag{NZ.1}\label{p1b:eq:natural-nz-1}
\]
The center is north and \(\partial B\) is south. The actual \(C_\mu\)
also has a unique north preimage, the cap center. Its inverse on the
north geodesic ball is as follows. For \(y\in\mathbb R^6\), put
\(\theta=|y|\), \(v=\operatorname{sinc}(\theta)y\), \(c=\cos(\theta/2)\), and
\[
 \begin{aligned}
 p_D&={\sqrt{1+c^2}\over2c}(-v_1,-v_3,v_2),\\
 u_D&=\left(-c^2,{\sqrt{1+c^2}\over2c}(-v_4,-v_5,v_6)\right),\\
 a&=u_D/\sqrt{1-|p_D|^2},\qquad
 b=-\exp(\pi p_D),\qquad b'=(b_0,b_1,-b_2,-b_3),\\
 I_\delta(q)&={2\delta^2X\over1+\sqrt{1+4\delta^2|X|^2}},
 \quad X=-\operatorname{Im}q/(1-\operatorname{Re}q),\\
 e_\mu(y)&=(I_\delta(a),I_\delta(b')).
 \end{aligned}                                             \tag{NZ.2}\label{p1b:eq:natural-nz-2}
\]
The zero values of the sinc ratios have their continuous extensions.
These formulas invert the actual cap quaternion, the actual second-variable
rotation, \(\iota\), and \(f_1\), in that order. In particular
\[
 De_\mu(0)=E_\mu
 =A^T\operatorname{diag}\left(
 {\pi\delta^2\over2\sqrt2}I_3,{\delta^2\over2\sqrt2}I_3\right),
 \qquad De_{\nu_A}(0)=E_\nu={r_B\over\pi}A^T.             \tag{NZ.3}\label{p1b:eq:natural-nz-3}
\]
The unique-north assertion also follows directly: \(f_1\) takes the
real quaternion coordinate to \(-1\) only at \(p_D=0,u_D=-1\), hence
\(a=b=-1\), which inverts to the two actual cap centers.

\begin{proposition}[Quantified actual source homotopies]
\label{p1b:proposition:natural-inverse-source}
Set \(\epsilon=2^{-12}\). On the real box \(|y_j|\le1/128\), the actual
inverse embeddings \eqref{p1b:eq:natural-nz-2} and \eqref{p1b:eq:natural-nz-9} below have the certified bounds
\[
 \|D^2e_\mu\|\le1/32,\qquad \|D^2e_\kappa\|\le32.        \tag{NZ.4}\label{p1b:eq:natural-nz-4}
\]
Consequently their straight interpolations to their indicated linear
derivatives are embeddings on \(|y|\le2\epsilon\). All their images
lie inside \(B\). These interpolations, followed by the specified linear
framing paths, construct the source homotopies used here.
\end{proposition}
\begin{proof}
The derivative bounds are finite interval estimates of the displayed
inverse formulas. The standalone supplied arithmetic certificate uses
50-decimal directed rounding for addition, multiplication and division,
and an extra outward unit for square roots. It bounds \(\pi\) by the
alternating rational series for
\(16\arctan(1/5)-4\arctan(1/239)\).
Every cosine, sinc and inverse-sine quotient is evaluated as a series
in its squared argument through order 32, with derivative tails included.
For \(|s|\le1/4\), the absolute coefficient bounds by one give tails
\(2(1/4)^{33}\), \(66(1/4)^{32}\),
and \(2178(1/4)^{31}\) in orders zero, one and two.
The chain rule is propagated on values, six first derivatives and all
36 second derivatives, with interval operations throughout. Every
denominator and square root is verified positive on the whole box.
The resulting second-derivative entry bounds are \(1/512\) and \(2\).
For a six-output Hessian, \(6\sqrt6<16\), proving \eqref{p1b:eq:natural-nz-4}.

On the \(2\epsilon\) ball the derivative differences from the two
linear maps are at most \(1/65536\) and \(1/64\).
The first minimum singular value exceeds \(\delta^2/3=1/432\);
the second equals \(4/3\). Hence the interpolated inverse embeddings
are injective, with inverse contraction ratio below \(1/2\).
The Taylor bounds
\(2\delta^2|y|+|y|^2/64\) and
\((4/3)|y|+16|y|^2\) are below \(r_B\) on this ball.
This proves both image and inverse claims with demonstrated bounds.
\end{proof}

Use the based target compression whose north polar angle is
\[
 \chi(\theta)=
 \begin{cases}
 (\pi/2)(3\theta/\epsilon-(\theta/\epsilon)^3),&\theta\le\epsilon,\\
 \pi,&\theta\ge\epsilon.
 \end{cases}                                             \tag{NZ.5}\label{p1b:eq:natural-nz-5}
\]
Interpolate the angle linearly from \(\theta\) to \(\chi(\theta)\).
The compressed map equals south off the inverse north ball.
For an interpolated inverse embedding \(e_t\), define its map on
\(e_t(D_\epsilon^6)\) by the compressed north exponential of
\(e_t^{-1}(z)\), and define it to be south elsewhere.
The boundary agrees because \(\chi(\epsilon)=\pi\).
The inverse is the unique contraction solution of
\[
 y=E^{-1}\bigl(z-(1-t)(e(y)-Ey)\bigr),                   \tag{NZ.6}\label{p1b:eq:natural-nz-6}
\]
with the indicated actual \(e,E\). Thus this is a specified map and
homotopy, rather than an appeal to equality of degrees.
For \(C_\mu\), compress, interpolate \(e_\mu\to E_\mu\), interpolate
the positive diagonal coefficients in \eqref{p1b:eq:natural-nz-3} to \(E_\nu\), and undo
the compression on \(\nu_A\). It is relative to \(\partial B\).
Apply \(U_H\) from \eqref{p1b:eq:natural-ny-15} at every step. Its endpoint is
\(U_H\nu_A\), with all preceding actual bottom comparisons retained.

\paragraph{A whole actual base homotopy for the bundle comparison.}
Write the retained balanced smash coordinates as \((u_0,\mathbf u_a,\mathbf u_b)\).
Put
\[
 \kappa(x,p)=(u_{b1},u_{a3},u_{b3},u_0,u_{b2},u_{a2},-u_{a1}),\qquad
 F_0(v)=(-v_7+iv_6,H_k+iH_j,H_i+iv_1),                  \tag{NZ.7}\label{p1b:eq:natural-nz-7}
\]
where \(H\) is the actual homogeneous Hopf vector of
\((v_4+iv_5,v_2+iv_3)\). The pointwise equality is
\(F_0\kappa=A_6s\), where in real target coordinates
\[
 A_6(o_1,o_2,o_3,o_4,o_5,o_6)=(o_3,o_4,-o_6,o_2,o_1,o_5).
                                                               \tag{NZ.8}\label{p1b:eq:natural-nz-8}
\]
First rotate the actual Hopf direction from \(qi\bar q\) to
\(-qi\bar q\) through \(q(i\cos t+j\sin t)\bar q\).
Then apply the prescribed \(SO(6)\) path from identity to
\(A_6\operatorname{diag}(-1,-1,1,1,1,-1)\).
This constructs \(s\to A_6s=F_0\kappa\).

Let \(c_0=(-1/3,-2\sqrt2/3,0,0)\), and move the actual source by
\((x,p)\mapsto(c_\ell x,c_\ell p)\),
\(c_\ell=\exp(-\ell i\arccos(-1/3))\).
The endpoint \(\kappa_1\) has its unique north preimage at \((1,1)\).
Its inverse is explicit. For \(v=\exp_N(y)\), recover
\(u_0=v_4\), \(\mathbf u_a=(-v_7,v_6,v_2)\),
\(\mathbf u_b=(v_1,v_5,v_3)\); set
\[
 X={\mathbf u_a+\mathbf u_b\over\sqrt2(1-u_0)},\quad
 P={\mathbf u_a-\mathbf u_b\over\sqrt2(1-u_0)},\quad
 e_\kappa(y)=\bigl(\log(c_0^{-1}\operatorname{st}^{-1}X),
                   \log(c_0^{-1}\operatorname{st}^{-1}P)\bigr).
                                                               \tag{NZ.9}\label{p1b:eq:natural-nz-9}
\]
These are logarithms about \(+1\); no antipodal logarithm is used.
Their derivative is \(E_\kappa=(4/3)O_\kappa\), with
\(O_\kappa=(R_i\oplus R_i)B_6\),
\[
 R_i=\begin{pmatrix}-1&0&0\\0&-1/3&-2\sqrt2/3\\0&2\sqrt2/3&-1/3\end{pmatrix},
 \quad B_6y=2^{-1/2}(y_1-y_4,y_6+y_5,y_2+y_3,
                         -y_1-y_4,y_6-y_5,y_2-y_3).      \tag{NZ.10}\label{p1b:eq:natural-nz-10}
\]
Both \(O_\kappa\) and \(A^T\) have determinant minus one.
Use \eqref{p1b:eq:natural-nz-5}--\eqref{p1b:eq:natural-nz-6} to linearize the actual inverse \eqref{p1b:eq:natural-nz-9}, then the
linear path
\[
 E_t=((1-t)4/3+t r_B/\pi)O_\kappa R_t,\quad
 R_0=I,\ R_1=O_\kappa^TA^T,                             \tag{NZ.11}\label{p1b:eq:natural-nz-11}
\]
and undo the target compression on \(\nu_A\). The two \(SO(6)\)
paths just used are specified by the same finite Givens construction:
eliminate each below-diagonal entry from the last row upwards using
\(\operatorname{atan2}(M_{ij},M_{jj})\), in increasing column order,
then multiply the inverses of those rotations with angles scaled by \(t\).
Their endpoint is the indicated actual matrix. This retains the local
framing; a degree calculation alone is not substituted.
Apply \(F_0\) immediately to this whole source homotopy.

For a unit cap direction \(v=(v_x,v_p)\), put
\(n=(v_{x3},-v_{p1},-v_{p3},v_{p2},v_{x1},v_{x2})\), and use
\((a,u)\) from \eqref{p1b:eq:natural-ny-1}. The final actual section direction is
\[
 \beta(v)=\left(-{4\arctan a\over\pi},\quad
 M\operatorname{block}(\bar u)e_1\,
 {\sqrt{1-(4\arctan a/\pi)^2}\over|u|}\right).            \tag{NZ.12}\label{p1b:eq:natural-nz-12}
\]
At \(u=0\) its complex vector is zero. Interpolate its real coordinate
from \(-a\) to \(-4\arctan(a)/\pi\), retaining this unit direction formula.
Starting at \(F_0\nu_A\), this ends at the actual section
\[
 s_{\rm can}(z)=\gamma_{\pi(1-|z|/r_B)}(\beta(z/|z|))
 \quad(z\in B),\qquad s_{\rm can}=e_1\quad\text{off }B.  \tag{NZ.13}\label{p1b:eq:natural-nz-13}
\]
Its center is \(-e_1\). The section path \(s_\ell:P\to S^5\) assembled
above is joint-continuous and piecewise differentiable in time.

\begin{proposition}[Actual bundle identification and pointwise frame factors]
\label{p1b:proposition:natural-exterior-factors}
Use the specified connection
\[
 \dot T=\mathcal A T,\quad T_0=I,\qquad
 \mathcal A=\dot s s^*-s\dot s^*
       -{s^*\dot s\over2}(I+ss^*).                      \tag{NZ.14}\label{p1b:eq:natural-nz-14}
\]
It gives \(T_1s=s_{\rm can}\), with \(T_1\in SU(3)\).
Let \(L_q=(1-i-j+k)/2\), \(R_q=(-1+i-j-k)/2\), and let
\(\mathcal L_\theta\) be the retained actual geodesic lift \eqref{p1b:eq:natural-nx-1}.
The following actual frames have last column \(s_{\rm can}\):
\[
 \begin{aligned}
 Q_{\rm cap}&=\mathcal L_\theta(\beta)
       \mathcal L_\pi(\beta)^{-1}\operatorname{diag}(-1,i,i)Q_0
       \operatorname{diag}(M\operatorname{block}(L_q^{-1}),1),\\
 Q_{\rm ext}&=Q_0\operatorname{diag}(M\operatorname{block}(R_q),1).
 \end{aligned}                                             \tag{NZ.15}\label{p1b:eq:natural-nz-15}
\]
At the center the first frame has its displayed constant extension.
Their actual clutch is
\[
 g_{\rm can}(v)=Q_{\rm cap}^*Q_{\rm ext}
 =\operatorname{block}(L_qq_2^{-1}(n)R_q),\qquad
 Mg_{\rm can}e_1=\Theta(q_2^{-1}(n)).                    \tag{NZ.16}\label{p1b:eq:natural-nz-16}
\]
In the original frames define
\[
 k_B=Q_{\rm cap}^*T_1Q_s,\quad
 h_E=Q_{\rm ext}^*T_1V^\#,\qquad
 \boxed{g=k_B^{-1}g_{\rm can}h_E\quad\text{on }\partial B.}
                                                               \tag{NZ.17}\label{p1b:eq:natural-nz-17}
\]
All these are actual \(SU(2)\) factors in the retained fibre axes.
\end{proposition}
\begin{proof}
The generator in \eqref{p1b:eq:natural-nz-14} is trace free and skew Hermitian;
\(\mathcal A s=\dot s\). Uniqueness of the transport ODE gives the
column and determinant assertions. Pointwise substitution in \eqref{p1b:eq:natural-nz-12}
gives \(M^{-1}\phi(\beta)M=\operatorname{block}(q_2^{-1})\).
The endpoint of the geodesic lift is
\(\operatorname{diag}(-1,i\phi)\). Hence on the cap boundary
\(Q_{\rm cap}=Q_0\operatorname{diag}(\phi^{-1}ML_q^{-1},1)\),
which proves the first identity in \eqref{p1b:eq:natural-nz-16}. Quaternion multiplication
gives the second. The factors in \eqref{p1b:eq:natural-nz-17} compare frames of the same
actual bundle, so they fix the last column and have determinant one.
Cancelling \(T_1\) proves the boxed equality, without changing base
coordinates or discarding a frame term.
\end{proof}

\paragraph{Both frame terms are removed by specified homotopies.}
Contract \(k_B|_{\partial B}\) radially inside its actual ball to \(k_B(0)\).
For \(E=P\setminus\operatorname{int}B\), define a strong retraction to
the antipodal wedge \(W_-=(S^3\times\{-1\})\cup(\{-1\}\times S^3)\).
If \(\alpha_x,\alpha_p\) are the angles about \(+1\), multiply both
angles by
\[
 1-t+t\pi/\max(\alpha_x,\alpha_p).                        \tag{NZ.18}\label{p1b:eq:natural-nz-18}
\]
An antipodal coordinate stays antipodal; its direction is not chosen.
The sine vanishes uniformly at that endpoint, proving continuity there.
The radial norm never decreases, so the retraction remains in \(E\).
Set
\[
 H(x,p)=h_E(x,-1)h_E(-1,-1)^{-1}h_E(-1,p).                \tag{NZ.19}\label{p1b:eq:natural-nz-19}
\]
It extends the actual wedge restrictions to all of \(P\).
On \(\partial B\) concatenate
\(h_E\to h_Er_1=Hr_1\to H\to H(1,1)\), using \eqref{p1b:eq:natural-nz-18} and then
the actual cap radial contraction. This gives an explicit nullhomotopy
up to the displayed constant, not a cancellation of its class by fiat.
Apply these two homotopies immediately in \eqref{p1b:eq:natural-nz-17}, then remove the two
actual constants by their specified quaternionic logarithm paths.
This constructs the actual homotopy \(g\to g_{\rm can}\).

\paragraph{Immediate closing application.}
The endpoint of the actual bottom path \eqref{p1b:eq:natural-ny-16}, after the applied
\(C_\mu\to\nu_A\) homotopy, is \(\psi(U_H\nu_A)\).
Change its radial real coordinate from \(\cos(\pi r)\) to \(1-2r\),
preserving both poles. Reverse the now constructed \(g\to g_{\rm can}\)
homotopy in its complex column. This gives the full relative bridge
\[
 \boxed{\mathcal B_\lambda:B\to SU(2),\quad
 \mathcal B_0=Q_s^*H^6Q_s,\quad
 \mathcal B_1=\psi(U_g),\quad \mathcal B_\lambda|_{\partial B}=I.}
                                                               \tag{NZ.20}\label{p1b:eq:natural-nz-20}
\]
Every source, frame and constant has been constructed in the preceding
equations. This supplies an intrinsic bridge directly. A comparison of
the old and new \(S^4\) endpoint sources is not assumed and is unnecessary
for this route. The bridge is applied to the complete target in the next
supplement, before any natural-bit inference.



\subsection{The actual intrinsic target and its parameter comparison}
\label{p1b:sec:natural-intrinsic-closing}

The short dependency chain is now applied: the actual Stiefel endpoint
comparison \eqref{p1b:eq:natural-ny-1}--\eqref{p1b:eq:natural-ny-16}, the inverse-source and retained-frame
comparison \eqref{p1b:eq:natural-nz-1}--\eqref{p1b:eq:natural-nz-20}, the intrinsic closing below, and its actual
boundary normalization. The sphere comparison remains the proved
\(w_t=1\). No primitive value is assigned to the determinant target.

\begin{theorem}[The actual prescribed rank-two closing]
\label{p1b:theorem:natural-intrinsic-closing}
Let \(Q_g(t,z)\) be the actual background of \eqref{p1b:eq:natural-nw-15}, with bottom
\(\psi(U_g)\), top \(-I_2\), and side \(gA(t)g^{-1}\).
Apply the bridge \eqref{p1b:eq:natural-nz-20} immediately:
\[
 X(t,z)=\mathcal B_{\min(4t,1)}(z)\psi(U_g(z))^{-1}Q_g(t,z).
                                                               \tag{OA.1}\label{p1b:eq:natural-oa-1}
\]
All its values fix the third cap column and have determinant one.
It has exactly the three prescribed cap faces
\[
 X(0,z)=Q_s^*H^6Q_s,\quad X(1,z)=\operatorname{diag}(-I_2,1),
 \quad X(t,r_Bv)=\operatorname{diag}(g(v)A(t)g(v)^{-1},1).
                                                               \tag{OA.2}\label{p1b:eq:natural-oa-2}
\]
Paste \(Q_sXQ_s^*\) to the original exterior \(K^\#\). This gives
\[
 K(t)s=s,\quad K(0)=H^6,\quad K(1)=N,\quad K|_J=K_+,
 \qquad F_h=KH^{-3},\quad F_h(1)=-\sigma F_h(0).          \tag{OA.3}\label{p1b:eq:natural-oa-3}
\]
Here \(-\sigma\) is the central minus on the actual rank-two bundle;
in the ambient three-frame it is represented by \(N\sigma\).
The original fibre representative is the actual ambient \(H^3\) of
component \((0,0;3)\), without an extra cap-frame conjugation.
\end{theorem}
\begin{proof}
Every factor of \eqref{p1b:eq:natural-oa-1} belongs to the same actual intrinsic \(SU(2)\).
At the bottom its two last factors cancel; at the top the bridge cancels
\(\psi(U_g)^{-1}\); on the side both are identity. This proves \eqref{p1b:eq:natural-oa-2}.
The side is exactly the retained exterior seam, so the pasting is a
continuous intrinsic closing. The exterior shift is zero on the joining
locus, hence its first eigenline there is exactly the measured \(K_+\).
The retained ambient root satisfies \(\sigma H=H^{-1}\).
Thus \(F_h(0)=H^3\), \(F_h(1)=NH^{-3}\), proving the original closing
equation, including its central sign.
\end{proof}

On the primitive cut \(a=e^{\pi it}\) define the actual determinant gauge
\[
 D_Y^0(a,x,p)=(ss^*+aP_s)F_h(t,x,p),\qquad
 D_Y(-a,p,x)=\mathsf J D_Y(a,x,p)\mathsf J,\qquad
 D_Y=D_Y^0T_Y.                                          \tag{OA.4}\label{p1b:eq:natural-oa-4}
\]
Its endpoints are \(H^3,H^{-3}\), so both exchange seams agree exactly.
Its determinant is \(a^2\). The measured joining origin is still
\(e_h=0\); the matching toggle is still \(T_Y^1\).
Retain the already constructed actual joining \(H_h\) of \eqref{p1b:eq:natural-nd-8}, with
the same localized odd clutch, ordinary frames and fixed four-cell
comparison. Its endpoints are the joining value of \eqref{p1b:eq:natural-oa-4} and the actual
split \(D_Z=\operatorname{diag}(1,z_1z_2,1)\), respectively.
This completes the intrinsic determinant triple \((D_Y,H_h,D_Z)\).
Multiply it by the inverse of the unchanged actual sphere triple:
\[
 C_Y=D_YS_Y^{-1},\qquad H_C=H_hH_s^{-1},\qquad C_Z=D_ZS_Z^{-1}.
                                                               \tag{OA.5}\label{p1b:eq:natural-oa-5}
\]
All factors are gauges of the same actual bundle.

\paragraph{The complete retained boundary restriction.}
Keep the original \(s,y,\rho,u\) coordinates and all the normal-link
and ordinary-frame identifications. On the inner stage put
\[
 a=ie^{is},\quad x=y\exp(3\pi\rho u/2),\quad
 p=y\exp(-3\pi\rho u/2),\quad F=F_YR.                    \tag{OA.6}\label{p1b:eq:natural-oa-6}
\]
On the middle stage put \(w=yu\), \(t_J=3\rho-1\), and
\(F_N=F_Y^{\rm neck}R_{\rm neck}=c_jF_Z\).
On the outer stage put \(\beta=3\pi(1-\rho)/2\),
\(z_1=e^{i(s+\beta)}\), \(z_2=e^{i(s-\beta)}\), \(w=yu\).
The full restriction of the determinant triple is the actual gauge
\[
 D_\partial=
 \begin{cases}
 F^*D_Y(a,x,p)F,&0\le\rho\le1/3,\\
 F_N^*H_h(t_J,a,w)F_N,&1/3\le\rho\le2/3,\\
 F_Z^*D_Z(z_1,z_2,w)F_Z,&2/3\le\rho\le1.
 \end{cases}                                             \tag{OA.7}\label{p1b:eq:natural-oa-7}
\]
Its endpoints use precisely the retained joining equation; the middle
map is the fixed actual \eqref{p1b:eq:natural-nd-8}, not a substitute stabilized joining.
The complete composite is
\(C_\partial=D_\partial S_\partial^{-1}\), with the original sphere
restriction \eqref{p1b:eq:natural-nc-7}. Every value of \eqref{p1b:eq:natural-oa-7}, including the cap preimage,
fixes the prescribed last column. The ordering line and both base-coordinate
maps have been retained.

\paragraph{Actual low and sphere-slice removal.}
Use the same phase comparison as \eqref{p1b:eq:natural-nx-12}--\eqref{p1b:eq:natural-nx-14}:
\[
\begin{gathered}
 f(\rho)=\min(1,\max(0,9\rho-4)),\qquad B_\rho=\operatorname{diag}(e^{i\pi(1-f(\rho))},1,1), \\
 Q_2(s)=\operatorname{diag}(e^{2is},1,1),\qquad M_D=Q_2^{-1}B_\rho D_\partial.
\end{gathered}
 \tag{OA.8}\label{p1b:eq:natural-oa-8}
\]
The real logarithm contracts the first phase; \(Q_2\) is the literal
\(2q_+\) representative. This makes determinant exactly one on all
three stages. Put \(M_C=Q_2^{-1}B_\rho C_\partial\).
No scalar-circle representative is substituted.

Make the same explicit sphere-base isotopy as in the retained sphere
proof, \(y\mapsto\exp(\pi\ell i)y\). Its endpoint is \(y\mapsto-y\),
and its bundle lift is identity on \(\mathbf1\oplus\lambda_L\) in
these prescribed split axes. It fixes the normal coordinates and is
homotopic to identity; it does not perform the normal-frame shear.
For either actual normalized gauge \(M\), define
\[
 \begin{aligned}
 G(s,y,n)&=M(s,-y,n)M(0,-y,n)^{-1},\\
 A_0(s,n)&=M(s,-1,n)M(0,-1,n)^{-1},\\
 E(s,y,n)&=G(s,y,n)A_0(s,n)^{-1}.
 \end{aligned}                                             \tag{OA.9}\label{p1b:eq:natural-oa-9}
\]
These are sequential divisions of actual maps on the same bundle.
The resulting \(E\) is exactly identity at \(s=0\) and \(y=1\).
It is therefore an actual parameter-smash comparison
\(S^1\wedge S^3\to H\), where \(H\) is the retained lens gauge space.

\begin{proposition}[The known lower terms are removed in the retained chart]
\label{p1b:proposition:natural-parameter-lower}
The determinant comparison \eqref{p1b:eq:natural-oa-9} has class
\(h_t t\in\pi_4H\); the composite comparison has class
\(\zeta_{\rm nat} t\in\pi_4H\), with the retained \(\zeta_{\rm nat}=h_t+1\).
Neither statement evaluates its coefficient.
\end{proposition}
\begin{proof}
In the retained product chart, the natural \(3z\) is independent of
\(s\), as is the sphere-slice \(-2b\) in the combined target.
Their actual representatives in the present normalization are the
\(s=0\) restrictions used in the first line of \eqref{p1b:eq:natural-oa-9}.
Thus that division removes the sphere slice, including its sixth term.
The restriction to the sphere basepoint carries precisely the remaining
circle slice, including the known even \(2\kappa\) discrepancy.
The second division removes this actual slice. The even-circle subgroup
is central in the retained target chart, so it contributes no commutator
carry to \(\pi_4H\). The resulting parameter-top coordinates are
\(h_e e+h_t t\), or the corresponding composite coordinates.
The retained evaluated \(h_e=w_e=0\) leaves the two stated classes.
The sphere-base isotopy has degree one and retains the ordering-line
lift, so it changes no natural generator or top coefficient.
This uses the actual lower restrictions; it does not identify a raw
primitive seven-sphere map with the natural target.
\end{proof}

\paragraph{The retained projective comparison and detector input.}
Let \(Q=\mathbb {RP}^2\subset L\) be the normal zero-circle-component
hyperplane. At \(Q\) the outer stage has \(\rho=1\).
Write the actual retained coherent frame as
\(F_Z=Q_0C_0(y,u)\), and put \(a_s=\operatorname{diag}(e^{is},e^{-is})\).
The actual determinant parameter comparison restricts to the explicit
equation
\[
 \begin{aligned}
 E_Q(s,y,u)&=\left(a_s^{-1}C_0(-y,u)^{-1}a_sC_0(-y,u)\right)\\
 &\hspace{6mm}\cdot
 \left(a_s^{-1}C_0(-1,u)^{-1}a_sC_0(-1,u)\right)^{-1}.
 \end{aligned}                                             \tag{OA.10}\label{p1b:eq:natural-oa-10}
\]
All \(C_0\) here are the actual finite transporter products of \eqref{p1b:eq:natural-nc-3},
not a degree-calibrated sphere map. On the ordering cover
\(C_0(y,-u)=dC_0(y,u)d\), \(d=\operatorname{diag}(1,-1)\), so this
equation retains the actual ordering-line gauge transformation.
The specific projective nullhomotopy is
\[
\boxed{\begin{gathered}
 \mathcal N_Q:[0,1]\times S^1\times S^3\times Q
     \to SU(\mathbf1\oplus\lambda_Q),\\
 \mathcal N_Q(0)=E_Q,\qquad \mathcal N_Q(1)=I,\\
 \mathcal N_Q|_{s=0\ \mathrm{or}\ y=1}=I.
\end{gathered}}
\tag{OA.11}\label{p1b:eq:natural-oa-11}
\]
The retained vanishing \(h_e=0\) proves that such a comparison exists.
Its actual map is constructed in \eqref{p1b:eq:natural-ob-1}--\eqref{p1b:eq:natural-ob-16}. In particular
\eqref{p1b:eq:natural-oa-10} is not replaced by identity merely because its component is zero.
Extend this specified homotopy over the normal affine-cell collar by
homotopy extension. That would give the actual supported sphere map
\[
 \varphi_D:R/K_R=S^7\to SU(2),\qquad [\varphi_D]=h_t\mu.   \tag{OA.12}\label{p1b:eq:natural-oa-12}
\]
The existing stabilized detector would then require the actual input
\[
\begin{gathered}
 \mathcal F_D:D^8\to SU(3),\qquad \mathcal F_D|_{S^7}=\operatorname{diag}(\varphi_D,1), \\
 g_D=[\mathcal F_De_3]:S^8\to S^5,\qquad \operatorname{Sq}^4u_5=h_t v_9\text{ on }C_{g_D}.
\end{gathered}
 \tag{OA.13}\label{p1b:eq:natural-oa-13}
\]
No such natural input or parity is claimed from the constructed cap
closing. The retained top insertion is injective; different genuine
relative \(Q\) nullhomotopies cannot change its \(C_2\) coefficient
by the cofiber exact sequence. The adopted original nullhomotopy is \eqref{p1b:eq:natural-ob-16}, and the complete
resulting coefficient is evaluated in \eqref{p1b:eq:natural-ci-8}, \eqref{p1b:eq:natural-aj-4}--\eqref{p1b:eq:natural-aj-5}.

The complete original comparison is evaluated in
\eqref{p1b:eq:natural-aj-4}--\eqref{p1b:eq:natural-aj-7}:
\(h_t=w_t=1\), \(r_J=\zeta_{\rm nat}=0\), with every retained
extension test decided. This component is used with its full frames,
lower divisions and joining; its isolated value is not a replacement
for that complete comparison.

The older named intrinsic endpoint comparisons can also be completed
immediately: reverse the retained \(H^6\to D^6C_\mu\), concatenate
\eqref{p1b:eq:natural-nz-20}, and reverse \(E_g\to\psi(U_g)\). This gives
\(D^6C_\mu\to E_g\); reversing the retained \(D^6C_\mu\to J_\tau\)
first gives \(J_\tau\to E_g\). Reversing the retained whole
\(H^6\to\psi(U_{\tau,\rm old})\) and concatenating \eqref{p1b:eq:natural-nz-20} gives an
actual \(SU(2)\) bridge between the old and new characteristic endpoints.
Every concatenation fixes \(\partial B\). No \(S^4\) source homotopy
between the old endpoint sources is assumed in these applications.



\subsection{The actual projective-plane comparison and natural sphere input}
\label{p1b:sec:natural-projective-comparison}

Retain the actual intrinsic target of \eqref{p1b:eq:natural-oa-1}--\eqref{p1b:eq:natural-oa-9}, the actual
\(C_0\), all original base coordinates and the ordering line. This
construction completes the specific missing equation \eqref{p1b:eq:natural-oa-11}. It is
applied immediately to construct \eqref{p1b:eq:natural-oa-12}. The coefficient \(h_t\) is
not assigned without the actual filling and evaluation in \eqref{p1b:eq:natural-oa-13}.

\paragraph{Executed dependencies.}
Normalize the actual diagonal base frame; construct the one-cell and
ordered-disc nullhomotopies; compare the actual neck frame in its same
intrinsic fibre; retain both actual inverse-image caps; correct and
contract the complete two-cell seam; paste the relative projective
nullhomotopy and apply it to the complete determinant comparison.

\subsubsection{Actual one-cell and ordered-disc comparisons}

At \(y=-1\), every section on the coherent path has second complex
coordinate zero. Its first frame column stays in \(\mathbb Ce_2\).
After the actual neck transport, the resulting \(C_b(u)=C_0(-1,u)\)
is diagonal in its intrinsic two-block. Thus it commutes with \(a_s\).
Put
\[
 c(y,u)=C_b(u)^{-1}C_0(-y,u),\qquad
 E_Q(s,y,u)=a_s^{-1}c(y,u)^{-1}a_sc(y,u),\quad c(1,u)=I.
                                                               \tag{OB.1}\label{p1b:eq:natural-ob-1}
\]
This is an exact rewrite of \eqref{p1b:eq:natural-oa-10}, not a replacement of its frame.
The actual base constant has been retained. On the ordering cover,
\(c(y,-u)=dc(y,u)d\), where \(d=\operatorname{diag}(1,-1)\).

Represent \(Q\) by the hemisphere \(u_2\ge0\), with
\[
 u(r,\theta)=\bigl(\sin(\pi r/2)\sin\theta,
                    \cos(\pi r/2),
                    \sin(\pi r/2)\cos\theta\bigr),\quad
 u_\theta=u(1,\theta),\quad 0\le\theta\le2\pi.
                                                               \tag{OB.2}\label{p1b:eq:natural-ob-2}
\]
Its boundary is the double attaching cover of the projective one-cell,
whose parameter is \(\theta\bmod\pi\). Let
\(j_\theta=\cos\theta\,j+\sin\theta\,k\), and use the retained
based B32 circle contraction, now with this specified ordering-line axis:
\[
 b_\tau(s,\theta)=
 \bigl(\cos(\pi\tau/2)e^{is}+\sin(\pi\tau/2)j_\theta\bigr)
 \bigl(\cos(\pi\tau/2)-\sin(\pi\tau/2)j_\theta\bigr).
                                                               \tag{OB.3}\label{p1b:eq:natural-ob-3}
\]
These are unit quaternions in the same retained intrinsic block.
They satisfy \(b_0=a_s\), \(b_1=I\), \(b_\tau(0,\theta)=I\),
and \(b_\tau(s,\theta+\pi)=db_\tau(s,\theta)d\).
For \(\mathcal C(B,c)=B^{-1}c^{-1}Bc\), define actual nullhomotopies
\[
 N_1(\tau,s,y,\theta)=\mathcal C(b_\tau(s,\theta),c(y,u_\theta)),
 \quad N_D(\tau,s,y,u)=\mathcal C(b_\tau(s,0),c(y,u)).
                                                               \tag{OB.4}\label{p1b:eq:natural-ob-4}
\]
The first descends to the projective one-cell. The second is defined on
the entire ordered hemisphere. Each starts at the actual \(E_Q\),
ends at identity and is identity when \(s=0\) or \(y=1\).
The second is not presumed to descend on the hemisphere boundary.

Set \(D(\tau,s,\theta)=b_\tau(s,\theta)b_\tau(s,0)^{-1}\).
It is identity on both \(\tau\) faces, on \(s=0\), and on
\(\theta=0,2\pi\). Direct multiplication gives the entire uncorrected
seam, with no omitted conjugation:
\[
 N_1N_D^{-1}
  =b_\tau(s,0)^{-1}\mathcal C(D(\tau,s,\theta),c(y,u_\theta))
                          b_\tau(s,0).
                                                               \tag{OB.5}\label{p1b:eq:natural-ob-5}
\]
It is a based six-sphere map. We contract a corrected version of this
actual map below, rather than inferring a nullhomotopy from \(h_e=0\).

\subsubsection{An applied intrinsic homotopy of the actual neck frame}

Only the evaluation \(u=k\) is needed after the explicitly separated
circle diagonal below. Write \(F=F_Y(y,k,\pi/2)\),
\(s_N=s(yk,-yk)\), and \(Q_0=[e_2,e_3,e_1]\).
An actual three-frame path starts at \(C_0(y,k)\) and has three legs:
\[
 \begin{array}{ll}
 Q_0^*T(r_\ell,s_N)^*F,
   &r_\ell=((1-\ell)e_1+\ell s_N)/\|(1-\ell)e_1+\ell s_N\|,\\
 Q_0^*F_Y(y,k,\alpha),&\alpha:\pi/2\longrightarrow0,\\
 R_\vartheta (Q_0^*F_Y(y,k,0))R_\vartheta^{-1},
   &R_\vartheta=
      \begin{pmatrix}1&0&0\\0&\cos\vartheta&\sin\vartheta\\
                         0&-\sin\vartheta&\cos\vartheta\end{pmatrix},
      \quad\vartheta:0\longrightarrow\pi/2.
 \end{array}                                                    \tag{OB.6}\label{p1b:eq:natural-ob-6}
\]
Here \(T\) is the original actual special-unitary transporter. The first
leg starts at the actual inverse neck transport, not a chosen frame.
All last columns of \eqref{p1b:eq:natural-ob-6} avoid \(+e_2\). Apply, at every value,
\[
 S(v)=T(-e_2,v)R_{-\pi/2},\qquad S(e_3)=I,
 \quad \widehat C=S(Ce_3)^{-1}C.
                                                               \tag{OB.7}\label{p1b:eq:natural-ob-7}
\]
This immediately converts \eqref{p1b:eq:natural-ob-6} to an intrinsic \(SU(2)\) path.
Divide on the left by its actual moving value at \(y=-1\), and
substitute \(y\mapsto-y\). The resulting based homotopy
\(H_3(\ell,y)\) starts at \(c(y,k)\).

Here are the chart exclusions used in \eqref{p1b:eq:natural-ob-7}. For the interior
\(0\le\alpha<\pi/2\) of the second leg, the two real coordinates
\((u_{a1},u_{a2})\) equal a nonzero complex scalar times
\((y_1,y_2)\). Its real scalar is proportional to
\(\cos\alpha(1-y_0\cos\alpha)>0\), away from the removable based
pole. If \(s(y,k,\alpha)=e_3\), these coordinates vanish and force
\(y_1=y_2=0\), which also forces \(u_{b1}=0\), contradicting the
real third coordinate one. At the neck \(u_0\ge1/3\), so the Hopf
radius is nonzero; \(s=e_3\) would force all its three Hopf coordinates
to vanish. This is again impossible. For the first leg, if \(s_N\)
has second coordinate zero and real nonzero third coordinate, its first
coordinate is positive real: when \(w_0\ne0\) the zero second
coordinate forces the real third coordinate zero; when \(w_0=0\),
zero imaginary third coordinate forces \(u_{b3}=0\), and then
\(h_3=\sqrt{u_0^2+u_{b2}^2}>0\). Its positive linear interpolation
with \(e_1\) cannot equal \(e_3\). In the third leg the real part of
the second last-column coordinate is
\(\sin\vartheta\cos\vartheta(\operatorname{Re}\bar\alpha-1)\le0\),
where \(\alpha\) is the first diagonal entry of the initial two-block.
It therefore cannot equal \(+e_2\). These exclusions prove every
transporter and column section in the actual path is defined.

The endpoint is a pointwise specified map, not just a degree comparison.
Put
\[
\begin{gathered}
 (A,B,C,D)=M(-y),\qquad M(q)=\left(\frac{1+3q_0}{3+q_0},
                  \frac{2\sqrt2\,\mathbf q}{3+q_0}\right), \\
 \alpha=\frac{(A-iD)^2}{\sqrt{A^2+D^2}},\qquad \beta=B+iC, \\
 C_2(y)=\begin{pmatrix}\alpha&\beta\\-\bar\beta&\bar\alpha\end{pmatrix}.
\end{gathered}
 \tag{OB.8}\label{p1b:eq:natural-ob-8}
\]
Use \(\alpha=0\) at its removable radius-zero locus.
Then \(H_3(1,y)=C_2(y)\) exactly and \(C_2(1)=I\).
This endpoint follows directly from the actual diagonal section and
the last rotation of \eqref{p1b:eq:natural-ob-6}. No homotopy of the actual bundle or its
normal framing has been discarded; \eqref{p1b:eq:natural-ob-6}--\eqref{p1b:eq:natural-ob-7} were applied in the
one fixed fibre evaluation needed for the seam.

\subsubsection{The two actual caps and their relative source isotopy}

For a quaternion \(q=\cos\varphi+\sin\varphi\,n\), let
\(\epsilon=1/128\),
\(\varphi_c=\pi\max(0,(\varphi-\pi+\epsilon)/\epsilon)\), and
\(p_\ell(q)=\cos((1-\ell)\varphi+\ell\varphi_c)
+\sin((1-\ell)\varphi+\ell\varphi_c)n\).
This is a based actual target compression. The map \(p_1C_2\) is
supported in precisely two disjoint inverse-image balls. They are
labelled by the sign of \(D\) in \eqref{p1b:eq:natural-ob-8}. Write their separate maps
as \(Q_+,Q_-\), identity off their respective balls. Then
\[
 p_1C_2=Q_+Q_-,\qquad
 \mathcal C(D,p_1C_2)=\mathcal C(D,Q_+)\mathcal C(D,Q_-).
                                                               \tag{OB.9}\label{p1b:eq:natural-ob-9}
\]
The second equation is pointwise exact because the supports are disjoint.

Their original source embeddings are explicit. For a target
\((\alpha,\beta)\) near \((-1,0)\), choose the local phase
\(\arg\alpha=\pi+\arg(-\alpha)\), take
\(z_\pm=\pm|\alpha|e^{i\arg\alpha/2}\), and put
\[
\begin{gathered}
 (A,B,C,D)=(\operatorname{Re}z_\pm,
              \operatorname{Re}\beta,\operatorname{Im}\beta,
                           -\operatorname{Im}z_\pm), \\
 y=-M^{-1}(A,B,C,D), \\
 M^{-1}(v)=\left(\frac{3v_0-1}{3-v_0},
                       \frac{2\sqrt2\,\mathbf v}{3-v_0}\right).
\end{gathered}
 \tag{OB.10}\label{p1b:eq:natural-ob-10}
\]
The positive-\(D\) cap uses \(z_-\), and the negative-\(D\) cap uses
\(z_+\). Both original inverse-image balls and their framings are retained.

Let \(\Phi_t\) be the flow in the \((A,D)\) plane
\[
\begin{gathered}
 \dot A=\chi(A)D,\qquad \dot D=-\chi(A)A,\qquad \dot B=\dot C=0,\qquad \chi(A)=v^2(3-2v), \\
 v=\min(1,\max(0,4(A+1/2))).
\end{gathered}
 \tag{OB.11}\label{p1b:eq:natural-ob-11}
\]
This is a globally defined \(C^1\) tangent vector field. It fixes
\((-1,0,0,0)\), the actual sphere basepoint in these coordinates.
On the entire positive cap and its path to the negative cap it equals
the clockwise rotation: \(|A|<\epsilon\), \(|D|>1-2\epsilon\)
at the endpoints, and its intervening semicircle stays in \(A>-1/4\).
Thus \(\Phi_\pi\) maps the positive cap to the negative cap, agrees
there exactly with \((A,D)\mapsto(-A,-D)\), and retains the output
of \eqref{p1b:eq:natural-ob-8}. In particular
\[
 Q_{-,\ell}(y)=Q_+\bigl(-M^{-1}\Phi_{-\pi(1-\ell)}M(-y)\bigr),
 \qquad Q_{-,0}=Q_-,\quad Q_{-,1}=Q_+.
                                                               \tag{OB.12}\label{p1b:eq:natural-ob-12}
\]
This is a relative homotopy of the two actual caps. It fixes \(y=1\).

\subsubsection{Complete contraction of the corrected seam}

Use the actual maps, with no calibration substitution:
\[
\begin{gathered}
K(\tau,s,y,\theta)=\mathcal C(D(\tau,s,\theta),Q_+(y)),\qquad
 R(\theta)=
 \begin{cases}
 K(2\theta)^{-1},&0\le\theta\le\pi,\\
 dK(2\theta-2\pi)^{-1}d,&\pi\le\theta\le2\pi.
 \end{cases}\\
N_1'=N_1R.
\end{gathered}
\tag{OB.13}\label{p1b:eq:natural-ob-13}
\]
The suppressed arguments of \(R\) are \((\tau,s,y)\).
It descends on the projective one-cell and is identity on both
\(\tau\) faces, \(s=0\), \(y=1\), and \(\theta=0,\pi\).
It therefore immediately corrects the actual one-cell nullhomotopy
without changing either of its endpoints.

Here is a complete specified homotopy \(\Pi:X\to I\) of the actual
corrected seam \(X=N_1'N_D^{-1}\). Parameterize its nine steps by
equal intervals. Put
\[
 \theta_1^\ell=(1-\ell)\theta+\ell\min(2\theta,2\pi),\qquad
 \theta_2^\ell=(1-\ell)\theta+\ell\max(2\theta-2\pi,0).
                                                               \tag{OB.14}\label{p1b:eq:natural-ob-14}
\]
This is a based homotopy of the circle diagonal to its two successive
coordinate circles. The steps, including the actual factors, are:
\begin{enumerate}
\item In \(X=(N_1N_D^{-1})(N_DRN_D^{-1})\), replace the conjugating
\(N_D(\tau)\) in the second factor by
\(N_D((1-\ell)\tau+\ell)\). This ends at \((N_1N_D^{-1})R\).
\item In \eqref{p1b:eq:natural-ob-5} replace the conjugating \(b_\tau(s,0)\) by
\(b_{(1-\ell)\tau+\ell}(s,0)\). This ends at
\(\mathcal C(D(\theta),c(y,u_\theta))R\).
\item Substitute \(\theta_1^\ell\) in \(D\) and
\(\theta_2^\ell\) in \(c(y,u_\theta)\). At its end the first half
circle has \(c(y,k)\); the second half has \(D=I\).
\item Expand the compressed first-half \(D\) to its full circle by
replacing \(\min(2\theta,2\pi)\) by its linear interpolation to
\(\theta\), while retaining \(c(y,k)\).
\item Apply the actual intrinsic \(H_3(\ell,y)\) of
\eqref{p1b:eq:natural-ob-6}--\eqref{p1b:eq:natural-ob-8} inside \(\mathcal C(D,H_3)R\).
\item Apply the actual target compression \(p_\ell C_2\).
\item Use the exact disjoint-support equality \eqref{p1b:eq:natural-ob-9} and replace
\(Q_-\) in its second commutator by \eqref{p1b:eq:natural-ob-12}. This ends at \(K^2R\).
\item Remove the constant deck conjugation in the second half of \(R\)
using the matrix
\[
 \operatorname{diag}(e^{i(1-\ell)\pi/2},e^{-i(1-\ell)\pi/2}).
\]
Its support-boundary values are identity.
\item Replace the now unconjugated two-circle concatenation by the
pointwise product, using \eqref{p1b:eq:natural-ob-14} backwards. Its exact endpoint is
\(K^2K^{-1}K^{-1}=I\).
\end{enumerate}
Every step is a homotopy of the whole actual six-sphere map, with the
prescribed composition order. It is identity on both \(\tau\) faces,
\(s=0\), \(y=1\), and \(\theta=0,2\pi\). This supplies
\[
 \Pi(0,\tau,s,y,\theta)=N_1'N_D^{-1},\qquad
 \Pi(1,\tau,s,y,\theta)=I.
                                                               \tag{OB.15}\label{p1b:eq:natural-ob-15}
\]
In particular no primitive value, Moore residue or equality of degrees
has been used to replace the actual seam by identity. The actual seam
is homotoped to two actual commutators and both are cancelled.

\begin{theorem}[The specific relative projective-plane nullhomotopy]
\label{p1b:theorem:natural-projective-null}
The missing actual map in \eqref{p1b:eq:natural-oa-11} is
\[
 \mathcal N_Q(\tau,s,y,[u(r,\theta)])=
 \begin{cases}
 N_D(\tau,s,y,u(r,\theta)),&0\le r\le1/2,\\
 \Pi(2(1-r),\tau,s,y,\theta)N_D(\tau,s,y,u(r,\theta)),
                                      &1/2\le r\le1.
 \end{cases}                                                   \tag{OB.16}\label{p1b:eq:natural-ob-16}
\]
It satisfies all of \eqref{p1b:eq:natural-oa-11}, in the actual ordering-line gauge bundle.
\end{theorem}
\begin{proof}
At \(r=1/2\), \eqref{p1b:eq:natural-ob-15} makes the two expressions identical. At
\(r=1\), their product is \(N_1'\), which has the correct antipodal
deck descent. At the hemisphere centre the expression is independent
of the absent \(\theta\). On the initial and terminal \(\tau\) faces,
\(\Pi=I\), so the value is respectively the actual \(E_Q\) and
identity. At \(s=0\) or \(y=1\) every factor is identity.
Every factor lies in the same intrinsic \(SU(2)\); its third column
stays fixed. These equalities prove the full relative statement,
including the boundary identification. This construction uses the
actual map \eqref{p1b:eq:natural-ob-1}, not merely its proved zero component.
\end{proof}

\subsubsection{Immediate application to the supported natural target}

Let \(E_D(s,y,\rho,u)\) be the complete actual determinant parameter
comparison \eqref{p1b:eq:natural-oa-7}--\eqref{p1b:eq:natural-oa-9}, including the original \eqref{p1b:eq:natural-nd-8} middle map.
Use the same normal-ball coordinate \(\rho\), the same bundle axes,
and the same basepoint \(y=1\). Apply \eqref{p1b:eq:natural-ob-16} immediately:
\[
 \varphi_D(s,y,\rho,u)=
 \begin{cases}
 E_D(s,y,2\rho,u),&0\le\rho\le1/2,\\
 \mathcal N_Q(2\rho-1,s,y,[u]),&1/2\le\rho\le1.
 \end{cases}                                                   \tag{OB.17}\label{p1b:eq:natural-ob-17}
\]
The seam is exactly \eqref{p1b:eq:natural-oa-10}. It is identity on \(\rho=1\),
\(s=0\) and \(y=1\), and is independent of \(u\) at \(\rho=0\).
Consequently this is an actual based map
\[
 \varphi_D:S^1\wedge S^3\wedge(D^3/S^2)=S^7\longrightarrow SU(2),
 \qquad [\varphi_D]=h_t\mu.
                                                               \tag{OB.18}\label{p1b:eq:natural-ob-18}
\]
Compressing the original normal ball and using \eqref{p1b:eq:natural-ob-16} on its collar
is its explicit homotopy-extension comparison. The retained top
insertion is injective, so genuine choices of the relative comparison
cannot alter this coefficient. No rotation or shear of the normal
product framing has been made. The complete \(Y\) cap, original root,
closing equation, matching toggle, middle joining and \(Z\) sphere-piece
gauge are precisely those of \eqref{p1b:eq:natural-oa-1}--\eqref{p1b:eq:natural-oa-9}.

On \(Q\) the actual sphere restriction is independent of \(s\).
In the first line of \eqref{p1b:eq:natural-oa-9} that same sphere factor cancels pointwise
from \(M_C(s)M_C(0)^{-1}\), and likewise in the base circle slice.
Thus \(E_C|_Q=E_D|_Q\) exactly. Applying the same \eqref{p1b:eq:natural-ob-16} in \eqref{p1b:eq:natural-ob-17},
with the complete actual \(E_C\) in its inner part, also constructs
\(\varphi_C:S^7\to SU(2)\) with \([\varphi_C]=\zeta_{\rm nat}\mu\).
No separate projective comparison or conditional bit model is needed.

This constructs the actual boundaries of the retained stabilized detector.
The stabilized detector has the exact filling boundary
\[
 \mathcal F_D:D^8\longrightarrow SU(3),\qquad
 \mathcal F_D|_{S^7}=\operatorname{diag}(\varphi_D,1).
                                                               \tag{OB.19}\label{p1b:eq:natural-ob-19}
\]
Its full projected map and evaluation would be
\[
 g_D=[\mathcal F_De_3]:S^8\longrightarrow S^5,\qquad
 \operatorname{Sq}^4u_5=h_tv_9\text{ on }C_{g_D}.
                                                               \tag{OB.20}\label{p1b:eq:natural-ob-20}
\]
The older actual cap filling \eqref{p1b:eq:natural-nx-5} is not \eqref{p1b:eq:natural-ob-19}: it fills different
faces, before the complete boundary joining and the actual lower and
projective comparisons. No identification with \eqref{p1b:eq:natural-ob-19} is asserted.
Thus its primitive degree or column value cannot assign \(h_t\).

The complete original comparison is evaluated in
\eqref{p1b:eq:natural-aj-4}--\eqref{p1b:eq:natural-aj-7}:
\(h_t=w_t=1\), \(r_J=\zeta_{\rm nat}=0\), with every retained
extension test decided. This component is used with its full frames,
lower divisions and joining; its isolated value is not a replacement
for that complete comparison.


\subsection{The actual matching-toggle contraction and natural filling annulus}
\label{p1b:sec:natural-toggle-filling}

Keep the complete intrinsic determinant triple \eqref{p1b:eq:natural-oa-1}--\eqref{p1b:eq:natural-oa-9} and its
actual supported boundary \eqref{p1b:eq:natural-ob-17}. The original two-bit target is still
\(h_t,w_t\), with \(w_t=1\). This construction supplies a particular
outer annulus of the missing filling \eqref{p1b:eq:natural-ob-19}, with its complete projected
column. It does not assign a bit from a characteristic-map calibration.

\paragraph{Executed dependencies.}
Identify the actual supported toggle with the retained characteristic
map by an explicit intrinsic homotopy; apply its existing three-frame
filling in the actual cap frame; apply the resulting contraction to the
complete target, retaining the original joining through an endpoint
track; move the determinant logarithm with that track; perform the
actual parameter divisions and projective collar; paste the resulting
annulus into the required natural filling.

\subsubsection{A contraction of the entire actual toggle}

Let \(v\in D^4\) be precisely the coordinates of the retained toggle,
and put \(g=1-|v|^2\). Its actual suspended-Hopf source is
\[
 y={|v|^2-g^2\over |v|^2+g^2},\qquad
 z={2gv\over |v|^2+g^2}\in\mathbb R^4=\mathbb C^2,
 \qquad y^2+|z|^2=1.                                  \tag{OC.1}\label{p1b:eq:natural-oc-1}
\]
The support boundary is \((y,z)=(1,0)\). Write
\(h(z)=(|z_0|^2-|z_1|^2,2\operatorname{Re}(z_0\bar z_1),
2\operatorname{Im}(z_0\bar z_1))/|z|^2\) when \(z\ne0\).
The actual toggle block is
\(q(v)=y+\sqrt{1-y^2}\,h(z)\) in the original Hamilton axes.
Set
\[
\begin{gathered}
 \theta_\ell=(1-\ell)\arccos y+\ell\pi(1-y)/2,\qquad r_\ell=\cos(\pi\ell/2)
        +{j+k\over\sqrt2}\sin(\pi\ell/2), \\
 H_\ell(v)=r_\ell(\cos\theta_\ell+\sin\theta_\ell h(z))r_\ell^{-1}.
\end{gathered}
 \tag{OC.2}\label{p1b:eq:natural-oc-2}
\]
At the two poles the value is the appropriate real quaternion;
the sine tends uniformly to zero there. Thus this is a whole based
intrinsic homotopy, including its support boundary.

Let \(\psi\) be exactly the retained cap conversion of the characteristic
map in \eqref{p1b:eq:natural-nx-2}--\eqref{p1b:eq:natural-nx-3}, and put \(b(v)=(-y,z_0,z_1)\).
Direct multiplication in those original axes gives
\[
 H_0(v)=q(v),\qquad H_1(v)=\psi(b(v)).                  \tag{OC.3}\label{p1b:eq:natural-oc-3}
\]
Indeed the characteristic matrix is
\(\cos\theta I+i\sin\theta(2P_z-I)\),
\(\theta=\pi(1-y)/2\). Its retained fibre conversion sends the
Hopf vector to \((-h_1,h_3,h_2)\). This is exactly conjugation
by \((j+k)/\sqrt2\). Equation \eqref{p1b:eq:natural-oc-3} identifies the actual
map pointwise; no equality of homotopy classes replaces it.

Use the already supplied \(\Lambda_\theta(b)\) of \eqref{p1b:eq:natural-nx-2}.
The specific nullhomotopy of this actual toggle block is
\[
 U_\tau(v)=
 \begin{cases}
  \operatorname{diag}(H_{2\tau}(v),1),&0\le\tau\le1/2,\\
  \Lambda_{2\pi(1-\tau)}(b(v)),&1/2\le\tau\le1.
 \end{cases}                                                   \tag{OC.4}\label{p1b:eq:natural-oc-4}
\]
Here the first notation uses the same Hamilton two-block as the retained
toggle. Equations \eqref{p1b:eq:natural-oc-3} and \eqref{p1b:eq:natural-nx-3} give the exact seam.
It starts at the actual toggle and ends at identity in \(SU(3)\).
It is identity on the support boundary because there \(b=(-1,0)\).

Apply \eqref{p1b:eq:natural-oc-4} in the actual cap frame, rather than a benchmark frame:
\[
 v=(4(t-1/2),\log(x)/\delta),\quad
 \mathsf U_\tau(a,x,p)=Q_s(x,p)U_\tau(v)Q_s(x,p)^*,
 \quad a=e^{\pi it},\quad\delta=1/12.                  \tag{OC.5}\label{p1b:eq:natural-oc-5}
\]
Set it to identity off the actual support. On the second primitive cut
define it by
\(\mathsf U_\tau(-a,p,x)=\mathsf J\mathsf U_\tau(a,x,p)\mathsf J\).
The support misses both cut endpoints, so both seams agree for every
\(\tau\). All values have determinant one. The entire actual cap frame,
including its dependence on \(p\), is retained. Consequently
\[
 \mathsf U_0=T_Y^1,\qquad \mathsf U_1=I,
 \qquad \mathsf U_\tau\text{ is an actual stabilized gauge on }Y.
                                                               \tag{OC.6}\label{p1b:eq:natural-oc-6}
\]
For \(\tau>1/2\) it can move the prescribed stabilizing line.
At \(v=0,\tau=3/4\), its cap-frame last column is \(-e_3\).
Thus its contribution cannot be dropped from the projected filling.

\subsubsection{Immediate application to the complete natural target}

Keep the original root \(H^3\), intrinsic closing \(K\), original closing
equation and prescribed joining \(H_J\) of \eqref{p1b:eq:natural-nd-8}. Define
\[
 D_Y^\tau=D_Y^0\mathsf U_\tau,\qquad D_Z^\tau=D_Z,
 \quad D_Y^0=(ss^*+aP_s)KH^{-3}.                        \tag{OC.7}\label{p1b:eq:natural-oc-7}
\]
This applies the contraction on the same right side as the actual
matching toggle; it is not a commutation past the root or closing.
The initial triple will be exactly the retained rank-two target.

Let \(b_\tau=\tau/(1+\tau)\), and let
\(D_+(a,w,-w)\) denote the actual untoggled joining endpoint \eqref{p1b:eq:natural-nd-2}.
Paste the reversed actual endpoint track to the unchanged joining:
\[
 H_J^\tau(t,a,w)=
 \begin{cases}
 D_+(a,w,-w)\mathsf U_{\tau-(1+\tau)t}(a,w,-w),&t\le b_\tau,\\
 H_J((1+\tau)t-\tau,a,w),&t\ge b_\tau.
 \end{cases}                                                   \tag{OC.8}\label{p1b:eq:natural-oc-8}
\]
At \(\tau=0\) the first interval has zero length and the map is exactly
\(H_J\). At the seam both sides are \(D_+T_Y\). At the initial
endpoint it is the restriction of \eqref{p1b:eq:natural-oc-7}; its terminal endpoint is
the original clutch-transported \(D_Z\). The shrinking interval is
joint-continuous at \(\tau=0\), by continuity of the actual endpoint
track and original joining. Thus \eqref{p1b:eq:natural-oc-7}--\eqref{p1b:eq:natural-oc-8} form an actual
homotopy of complete stabilized triples, retaining every original
root, frame, toggle endpoint, odd clutch and joining factor.

Restrict this whole family along \eqref{p1b:eq:natural-oa-6}--\eqref{p1b:eq:natural-oa-7}; call the restriction
\(D_\partial^\tau\). The inserted prefix changes the location of the
joining phase. Its actual determinant logarithm must move with it:
\[
\begin{gathered}
f_\tau(\rho)=
 \begin{cases}
 0,&\rho\le1/3,\\
 \min(1,\max(0,3\max(0,(1+\tau)(3\rho-1)-\tau)-1)),
          &1/3\le\rho\le2/3,\\
 1,&\rho\ge2/3,
 \end{cases}\\
B_\rho^\tau=\operatorname{diag}(e^{i\pi(1-f_\tau(\rho))},1,1).
\end{gathered}
\tag{OC.9}\label{p1b:eq:natural-oc-9}
\]
It agrees with \eqref{p1b:eq:natural-oa-8} at \(\tau=0\). On the inserted track the
determinant remains \(a^2\); on the retained joining part its phase
is the original phase at its reparameterized time. Hence
\[
 M_\tau=Q_2^{-1}B_\rho^\tau D_\partial^\tau\in SU(3),
 \qquad Q_2=\operatorname{diag}(e^{2is},1,1).            \tag{OC.10}\label{p1b:eq:natural-oc-10}
\]
This retains the actual real-logarithm comparison and literal \(2q_+\).
Keeping a fixed old logarithm on \eqref{p1b:eq:natural-oc-8} would not give this equation.

Apply the actual sequential divisions, in their original order:
\[
 G_\tau=M_\tau(s,-y,n)M_\tau(0,-y,n)^{-1},\quad
 A_\tau=M_\tau(s,-1,n)M_\tau(0,-1,n)^{-1},\quad
 E_\tau=G_\tau A_\tau^{-1}.                            \tag{OC.11}\label{p1b:eq:natural-oc-11}
\]
The circle and sphere base faces are exactly identity. On \(Q\) the
outer map is unchanged, so \(E_\tau|_Q=E_Q\) for every \(\tau\).
Apply the already completed \(\mathcal N_Q\) immediately:
\[
 \Phi_\tau(s,y,\rho,u)=
 \begin{cases}
 E_\tau(s,y,2\rho,u),&\rho\le1/2,\\
 \mathcal N_Q(2\rho-1,s,y,[u]),&\rho\ge1/2.
 \end{cases}
 \quad \Phi_\tau:S^7\longrightarrow SU(3),\quad
 \Phi_0=\operatorname{diag}(\varphi_D,1).                \tag{OC.12}\label{p1b:eq:natural-oc-12}
\]
All boundary identifications and both base faces are retained. The
normal product framing has not been rotated. At \(\tau=1\), the inner
\(Y\) gauge has its toggle removed, but the full joining includes the
actual reverse contraction \eqref{p1b:eq:natural-oc-8}. Its column can still move.

\subsubsection{The actual annulus and its entire last column}

Use \(r\) for the eight-ball radial coordinate and \(\xi\in S^7\)
for the supported boundary coordinates in \eqref{p1b:eq:natural-oc-12}. The constructed
part of the specific natural filling is
\[
 \mathcal A_D(r,\xi)=\Phi_{2(1-r)}(\xi),\qquad
 1/2\le r\le1,\quad
 \mathcal A_D(1)=\operatorname{diag}(\varphi_D,1),\quad
 \mathcal A_D(1/2)=\Phi_1.                              \tag{OC.13}\label{p1b:eq:natural-oc-13}
\]
The exact unsupplied core and the immediate pasting equation are
\[
 \mathcal G_D:D^8\to SU(3),\quad\mathcal G_D|_{S^7}=\Phi_1,
 \qquad
 \mathcal F_D(r,\xi)=
 \begin{cases}
 \mathcal G_D(2r,\xi),&r\le1/2,\\
 \mathcal A_D(r,\xi),&r\ge1/2.
 \end{cases}                                                   \tag{OC.14}\label{p1b:eq:natural-oc-14}
\]
This identifies a completed actual component of \eqref{p1b:eq:natural-ob-19}, with its exact
remaining boundary; it does not assert an actual core from existence.

There is also a full column formula that does not require a numerical
replacement of the original geometric joining. On the circle-base slice
\(y=1\), the original inner pairs have real coordinate
\(-\cos\alpha\le0\), \(0\le\alpha\le\pi/2\); at the neck
\(w=-u\) is purely imaginary. Both cuts miss the actual toggle cap.
Thus \(A_\tau=A_0\) remains intrinsic and fixes \(e_3\), giving
\[
 \Phi_\tau e_3=M_\tau(s,-y,n)M_\tau(0,-y,n)^*e_3
                         \quad\text{on its inner normal half}. \tag{OC.15}\label{p1b:eq:natural-oc-15}
\]
On the inner \(Y\) stage let \(F_s=F_YR(s)\),
\(F_0=F_YR(0)\), \(s_*\) be the actual section at its actual pair,
and \(U_s=\mathsf U_\tau(a,x,p)\),
\(U_0=\mathsf U_\tau(i,x,p)\). Expanding \eqref{p1b:eq:natural-oc-15} and using only
the exact equality \(D_Y^0(i)s_*=s_*\) yields
\[
 \Phi_\tau e_3=
 Q_2^{-1}B_\rho^\tau F_s^*D_Y^0(a)U_s
                          F_sF_0^*U_0^*s_* .           \tag{OC.16}\label{p1b:eq:natural-oc-16}
\]
Here \(\rho\) is the original, uncompressed normal radius.
In particular the root in \(D_Y^0(a)\), moving cap frame, actual
baseline toggle and frame comparison \(F_sF_0^*\) remain present.
None is discarded as a static factor.

On the inserted joining prefix the same equation holds with
\[
\begin{gathered}
 F_s=F_N(s),\qquad F_0=F_N(0),\qquad D_Y^0(a)=D_+(a,w,-w),\qquad U_s=\mathsf U_{\tau-(1+\tau)t_J}(a,w,-w), \\
 U_0=\mathsf U_{\tau-(1+\tau)t_J}(i,w,-w).
\end{gathered}
 \tag{OC.17}\label{p1b:eq:natural-oc-17}
\]
Off that prefix, the retained original joining is intrinsic; on the
outer stage and projective collar it is intrinsic as well. Their columns
are therefore exactly \(e_3\). The column is also exactly \(e_3\)
for \(\tau\le1/2\). Equations \eqref{p1b:eq:natural-oc-15}--\eqref{p1b:eq:natural-oc-17}, with these constant
complements, construct the whole projected annulus from actual inputs.

The remaining natural detector must use the entire pasted map:
\[
 g_D=[\mathcal F_De_3]:S^8\to S^5,
 \qquad \operatorname{Sq}^4u_5=h_tv_9\text{ on }C_{g_D}.
                                                               \tag{OC.18}\label{p1b:eq:natural-oc-18}
\]
The annulus cannot be assigned a separate natural parity: its inner
face \(\Phi_1e_3\) is not constant. Its contribution must be retained
along with the actual core in \eqref{p1b:eq:natural-oc-14}. This is the actual remaining
comparison equation, not a conditional enumeration of bits.

The complete original comparison is evaluated in
\eqref{p1b:eq:natural-aj-4}--\eqref{p1b:eq:natural-aj-7}:
\(h_t=w_t=1\), \(r_J=\zeta_{\rm nat}=0\), with every retained
extension test decided. This component is used with its full frames,
lower divisions and joining; its isolated value is not a replacement
for that complete comparison.


\subsection{The actual core neck-frame comparison and its second filling annulus}
\label{p1b:sec:natural-core-frame}

Keep the actual supported \(\Phi_1\) of \eqref{p1b:eq:natural-oc-12}, not a calibrated
seven-sphere map. Its entire moving column is on the inserted joining
prefix, where \(w=yu\) is in one of the two actual toggle caps.
The following constructs a homotopy of its actual neck frame on both
caps, applies it to the whole \(\Phi_1\), and pastes the resulting
second annulus into the same natural filling \eqref{p1b:eq:natural-ob-19}. It does not
assign \(h_t\) from either annulus.

\paragraph{The actual local frame is contracted as a map.}
Let \(\epsilon=\pm1\), \(d(w,\epsilon)<2\delta\),
\(\delta=1/12\), and \(y=w(-u)\). First contract the actual cap
coordinate by
\[
 w_\ell=\epsilon\exp((1-\ell)\log(\epsilon w)),\qquad
 y_\ell=w_\ell(-u),\qquad
 C_0(y,u)\longrightarrow C_0(-\epsilon u,u).            \tag{OD.1}\label{p1b:eq:natural-od-1}
\]
Every value here is the original finite-product \(C_0\) in the same
intrinsic fibre. No nontrivial frame term is erased by a degree argument.

At this actual pole evaluation \(y_*=-\epsilon u\), the neck section
is \(e_1\), so its actual neck transporter is identity. Undo the
original coherent path, then rotate the original diagonal-section frame:
\[
 Q_0^*F_Y(y_*,u,\alpha),\quad\alpha:\pi/2\longrightarrow0;
 \qquad R_\vartheta(Q_0^*F_Y(y_*,u,0))R_\vartheta^{-1},
       \quad\vartheta:0\longrightarrow\pi/2.            \tag{OD.2}\label{p1b:eq:natural-od-2}
\]
Apply the actual section \eqref{p1b:eq:natural-ob-7} immediately at every value of this path,
\(\widehat C=S(Ce_3)^{-1}C\). Its endpoint is intrinsic.

Here is a chart exclusion on the entire two-sphere family, rather than
only its fixed-axis evaluation. Put \(c=\cos\alpha\),
\(s=\sin\alpha\). The original balanced coordinates on this actual
path are
\[
 u_0={1+3s^2\over3+s^2},\qquad
 \mathbf u_a={-2\sqrt2\epsilon c\over3+s^2}u,
 \qquad \mathbf u_b={-2\sqrt2sc\over3+s^2}u.            \tag{OD.3}\label{p1b:eq:natural-od-3}
\]
For \(\alpha<\pi/2\), vanishing of the second complex section
coordinate forces \(u_1=u_2=0\); its third complex coordinate is then
purely imaginary. It cannot be real one. At the neck the section is
\(e_1\). Thus no section in the first leg is \(+e_3\), equivalently
no last column in \eqref{p1b:eq:natural-od-2} is \(+e_2\). The Hopf radius is nonzero
because \(u_0\ge1/3\). In the second leg the real part of the
second last-column coordinate is
\(\sin\vartheta\cos\vartheta(\operatorname{Re}\bar s_1-1)\le0\).
This also excludes \(+e_2\). Hence the intrinsic correction is
defined on the entire actual family, including both poles.

The endpoint can be read directly from the original diagonal section:
\[
\begin{gathered}
 (A,B,C,D)=\left(1/3,{2\sqrt2\over3}\operatorname{Im}y_*\right),\qquad \alpha_*={(A-iD)^2\over\sqrt{A^2+D^2}}, \\
 \beta_*=B+iC, \\
 C_*(u,\epsilon)=
 \begin{pmatrix}\alpha_*&\beta_*\\-\bar\beta_*&\bar\alpha_*\end{pmatrix}.
\end{gathered}
 \tag{OD.4}\label{p1b:eq:natural-od-4}
\]
This is the exact endpoint, with no divided base constant: none is
needed on this unbased local two-sphere family. It avoids \(-I\).
Indeed \(\beta_*=0\) forces \(u_1=u_2=0\); then \(D\ne0\),
and \(\operatorname{Im}\alpha_*=-2AD/\sqrt{A^2+D^2}\ne0\).
Its quaternionic principal logarithm is therefore continuous over the
whole family. Contract it by \(\exp((1-\ell)\log C_*)\).

Concatenate \eqref{p1b:eq:natural-od-1}, the corrected \eqref{p1b:eq:natural-od-2}, and this logarithm contraction
in equal intervals. This constructs an actual specified homotopy
\[
 \mathcal C(\lambda,w,u)\in SU(2),\quad
 \mathcal C(0,w,u)=C_0(w(-u),u),\qquad
 \mathcal C(1,w,u)=I.                                  \tag{OD.5}\label{p1b:eq:natural-od-5}
\]
Both actual pole caps and all original frame constants are retained.

\paragraph{Apply it to the whole actual core boundary.}
Use \(\rho\) for the compressed normal radius of \eqref{p1b:eq:natural-oc-12}.
Let \(\chi_N\) be zero outside \([1/8,1/4]\), one on
\([1/6,5/24]\), and linear on the two intervening intervals.
Let \(\chi_P(\theta)\) be one for \(\theta\le\delta\), zero for
\(\theta\ge2\delta\), and \(2-\theta/\delta\) between.
Set \(w=(-y)u\), \(\theta=d(w,\{1,-1\})\), and
\(\chi=\chi_N(\rho)\chi_P(\theta)\). Define
\[
 L_\ell=Q_2^{-1}B_{2\rho}^{1}
          \mathcal C(1-\ell\chi,w,u)
                       (B_{2\rho}^{1})^{-1}Q_2,        \tag{OD.6}\label{p1b:eq:natural-od-6}
\]
and set it to identity off the two cap neighborhoods or the normal
support. It is continuous there because \(\mathcal C(1)=I\).
The normal cutoff makes it independent of the absent polar direction
at \(\rho=0\) and identity on the whole projective collar.
The two caps are disjoint, so no pole-sign choice is made on their
common complement. Every value fixes \(e_3\) and has determinant one.

Apply it immediately as a homotopy of the whole actual matrix map:
\[
 \Xi_\ell=L_\ell\Phi_1L_\ell^{-1}:S^7\to SU(3),
 \quad\Xi_0=\Phi_1,\qquad
 \Xi_\ell e_3=L_\ell\Phi_1e_3.                         \tag{OD.7}\label{p1b:eq:natural-od-7}
\]
All relative faces of \(\Phi_1\) remain identity under conjugation.
This includes the actual changes to its intrinsic complementary region;
only the column there stays constant. Thus this is a whole-map comparison,
not a replacement of the moving-column support alone.

On that support the cutoff is one. The original odd transition is
identity: its cap requires \(w\) near \(aj\), whose distance from
either real pole is \(\pi/2\), while the actual toggle support has
\(d(w,\{1,-1\})<\delta\). Hence \(F_N(s)=F_N(0)=F_Y^{\rm neck}\)
there. With the retained cap frame \(Q_{\rm cap}=T_wQ_0\),
one has \(F_Y^{\rm neck}=Q_{\rm cap}C_0\). Formula \eqref{p1b:eq:natural-oc-17} therefore
gives the exact endpoint column
\[
 \Xi_1e_3=Q_2^{-1}B_{2\rho}^{1}
        D_+^{\rm ord}(a,w)
        U_s^{\rm ord}(U_0^{\rm ord})^*e_3,              \tag{OD.8}\label{p1b:eq:natural-od-8}
\]
where every ordinary matrix is conjugated by the same actual
\(Q_{\rm cap}\), and the toggle parameter remains
\(1-2t_J\). The source is still the original \(w=(-y)u\),
with both actual cuts and the actual \(a=i\) baseline.
All complementary columns are exactly \(e_3\).
The actual neck frame has been removed by \eqref{p1b:eq:natural-od-5}--\eqref{p1b:eq:natural-od-7}, not by
silently resetting it. Its entire homotopy has been retained in the filling.

\paragraph{Second annulus in the original natural filling.}
Paste this comparison immediately inside \eqref{p1b:eq:natural-oc-13}. The actual remaining
core and full assembly equation are
\[
 \mathcal J_D:D^8\to SU(3),\quad\mathcal J_D|_{S^7}=\Xi_1,
 \qquad
 \mathcal F_D(r,\xi)=
 \begin{cases}
 \mathcal J_D(4r,\xi),&0\le r\le1/4,\\
 \Xi_{2-4r}(\xi),&1/4\le r\le1/2,\\
 \Phi_{2(1-r)}(\xi),&1/2\le r\le1.
 \end{cases}                                                   \tag{OD.9}\label{p1b:eq:natural-od-9}
\]
Both annular seams are exact. This completes an actual region of the
specific filling \eqref{p1b:eq:natural-ob-19}, with its original boundary and all parameter,
bundle, lower-comparison and projective identifications. The fractions
are coordinate choices; they are not progress percentages or a bound
on the remaining calculation.

The entire projected region follows directly from \eqref{p1b:eq:natural-oc-15}--\eqref{p1b:eq:natural-oc-17}
and \eqref{p1b:eq:natural-od-7}. The resulting complete detector input is
\[
 \mathcal J_D|_{S^7}=\Xi_1\text{ in \eqref{p1b:eq:natural-od-7}},\qquad
 g_D=[\mathcal F_De_3]:S^8\to S^5,\qquad
 \operatorname{Sq}^4u_5=h_tv_9\text{ on }C_{g_D}.        \tag{OD.10}\label{p1b:eq:natural-od-10}
\]
No actual \(\mathcal J_D\), finite mapping-cone input, or natural
parity is supplied by this frame comparison. The nonconstant inner
column is retained and cannot be omitted from the final pasting.
The complete original comparison is evaluated in
\eqref{p1b:eq:natural-aj-4}--\eqref{p1b:eq:natural-aj-7}:
\(h_t=w_t=1\), \(r_J=\zeta_{\rm nat}=0\), with every retained
extension test decided. This component is used with its full frames,
lower divisions and joining; its isolated value is not a replacement
for that complete comparison.


\subsection{The actual core column contraction and its intrinsic remainder}
\label{p1b:sec:natural-core-retraction}

Keep the complete actual \(\Xi_1\) of \eqref{p1b:eq:natural-od-7}, including its intrinsic
complement. This construction contracts its entire column, relative to
the retained base faces, and applies that contraction to the whole map.
It supplies a specified intrinsic endpoint; it does not presume that
endpoint is null in \(SU(2)\) or assign the determinant bit.

\paragraph{Executed dependencies.}
Construct a full ordinary lift of the exact \eqref{p1b:eq:natural-od-8} column, retaining both
toggles and their baseline; contract the actual normal direction with
the required moving sphere-base correction; apply this homotopy to the
whole \(\Xi_1\); contract the resulting radial column, again on the
whole map; paste this comparison immediately into \eqref{p1b:eq:natural-od-9}.

For every \(w\in S^3\) use the retained joining frame
\(Q_w=T(e_1,s(w,-w))Q_0\). It is globally defined by the original
joining chart bound \(|1+s_1|\ge2/9\). Let
\(D_s(w)=Q_w^*D_+(ie^{is},w,-w)Q_w\) and let
\(U_s^\lambda(w)=Q_w^*\mathsf U_\lambda(ie^{is},w,-w)Q_w\).
All factors here are the actual \eqref{p1b:eq:natural-nd-2} and \eqref{p1b:eq:natural-oc-5}, with both cuts,
not just their homotopy classes. Set
\[
\begin{gathered}
 \lambda(\nu)=\min(1,\max(0,3-12\nu)),\qquad B=\operatorname{diag}(-1,1,1), \\
 P(s,w,\nu)=Q_2^{-1}B D_s(w)U_s^{\lambda(\nu)}(w)
       U_0^{\lambda(\nu)}(w)^*D_0(w)^*B^*.
\end{gathered}
 \tag{OE.1}\label{p1b:eq:natural-oe-1}
\]
Here \(\nu\) is the compressed normal radius in \eqref{p1b:eq:natural-oc-12}.
The extension at \(\lambda=0,1\) is exactly the displayed intrinsic
matrix; it is not replaced by identity. Directly,
\[
 P\in SU(3),\quad P(0,w,\nu)=I,\qquad
 Pe_3=e_3\quad\text{if }\nu\notin(1/6,5/24).
                                                        \tag{OE.2}\label{p1b:eq:natural-oe-2}
\]
Indeed the determinant ratio of the two positive joining-line factors
is \(e^{2is}\), canceled by literal \(Q_2\). The toggle factors have
determinant one. For \(\lambda\le1/2\), or \(\lambda=1\), they
are intrinsic. The right factors \(D_0^*B^*\) fix the last column.
Thus on its entire moving support \(Pe_3\) is exactly \eqref{p1b:eq:natural-od-8}.

Use the already specified \(\chi_N\) of \eqref{p1b:eq:natural-od-6}, including its ramps.
For \(0\le t\le1\) define the following actual unit quaternion and
based matrix; do not omit the second factor:
\[
\begin{gathered}
 \begin{aligned}
 q_t(\nu,u)&=\cos\frac\pi2(1-t\chi_N(\nu))\\
  &\quad+\sin\frac\pi2(1-t\chi_N(\nu))u,
 \end{aligned}\\
 P_t(s,y,\nu,u)=P(s,(-y)q_t,\nu)P(s,-q_t,\nu)^*.
\end{gathered}
\tag{OE.3}\label{p1b:eq:natural-oe-3}
\]
At \(t=0\), the baseline \(-u\) is purely imaginary and misses both
real toggle caps. Consequently its full matrix is intrinsic, and
\[
 P_0e_3=\Xi_1e_3,\qquad P_t|_{s=0\text{ or }y=1}=I.
                                                        \tag{OE.4}\label{p1b:eq:natural-oe-4}
\]
This identifies the entire column, including its constant complement.
It uses the actual \(w=(-y)u\), not a source class or a primitive map.

Apply the comparison immediately to the full original core:
\[
 Z_t=P_tP_0^*\Xi_1,\quad Z_0=\Xi_1,\quad Z_te_3=P_te_3.
                                                        \tag{OE.5}\label{p1b:eq:natural-oe-5}
\]
Off the support of \(\chi_N\), the multiplier is defined to be
identity. This agrees with \eqref{p1b:eq:natural-oe-3} and avoids using a polar direction
where that direction is absent. At every support edge its value tends
to identity, uniformly in \(u\), by continuity on compact sources.
All original base faces, the outer collar, and the original matrix map
on the complementary region are retained. In the two ramps the new
values are intrinsic. This is a homotopy of the entire actual \(S^7\)
map; it is not an unrecorded alteration of its column alone.

The endpoint column is independent of the normal direction. More
precisely let
\[
 A(s,y,\nu)=P(s,-y,\nu)P(s,-1,\nu)^*,\qquad
 Z_1e_3=A(s,y,\nu)e_3.                                \tag{OE.6}\label{p1b:eq:natural-oe-6}
\]
On the moving support \(\chi_N=1\), so \(q_1=1\) and this is
pointwise \eqref{p1b:eq:natural-oe-3}. Outside that support both sides fix \(e_3\) by
\eqref{p1b:eq:natural-oe-2}. Thus \eqref{p1b:eq:natural-oe-6} is an equation for the entire map, not merely its
regular part. The moving baseline in \eqref{p1b:eq:natural-oe-3} is indispensable: the
uncorrected substitution \(u\to1\) would fail the \(y=1\) face.

Now contract the radial dependence, retaining the actual matrix lift:
\[
 \nu_b=(1-b)\nu+b,\qquad
 V_b=A(s,y,\nu_b)A(s,y,\nu)^*Z_1,\qquad 0\le b\le1.
                                                        \tag{OE.7}\label{p1b:eq:natural-oe-7}
\]
Its entire column is \(A(s,y,\nu_b)e_3\). At \(\nu=1\) the
multiplier is identity for every \(b\), so the collapsed normal
boundary stays identity. At \(\nu=0\) the multiplier is independent
of the absent direction, and \(Z_1\) is already independent of it.
Both parameter-base faces are identity, since the two factors in
\(A\) agree there. At the endpoint \(A(s,y,1)\) is intrinsic by
\eqref{p1b:eq:natural-oe-2}, proving that \(V_1\) fixes \(e_3\) pointwise.

Concatenate these two whole-map homotopies in equal intervals:
\[
 \mathcal H_t=
 \begin{cases}Z_{2t},&0\le t\le1/2,\\
 V_{2t-1},&1/2\le t\le1,
 \end{cases}
 \quad \mathcal H_0=\Xi_1,\quad
 \mathcal H_1=\operatorname{diag}(\psi_D,1).
                                                        \tag{OE.8}\label{p1b:eq:natural-oe-8}
\]
This constructs a specific relative column nullhomotopy and applies it
to the actual core, using multiplication by specified \(SU(3)\) matrices.
No existence-only horizontal-lift or framed-cobordism step is required.
The full intrinsic endpoint is the actual equation
\[
 \operatorname{diag}(\psi_D,1)=
 A(s,y,1)A(s,y,\nu)^*P_1(s,y,\nu,u)P_0(s,y,\nu,u)^*
           \Xi_1(s,y,\nu,u).                           \tag{OE.9}\label{p1b:eq:natural-oe-9}
\]
At the polar faces use the identity convention in \eqref{p1b:eq:natural-oe-5}.
Every old complementary matrix, root, joining, phase and lower
comparison is still in this expression through the complete \(\Xi_1\).

\paragraph{Immediate pasting and the retained assembly equation.}
The remaining core now has the specified intrinsic boundary
\[
 \mathcal K_D:D^8\to SU(3),\qquad
 \mathcal K_D|_{S^7}=\operatorname{diag}(\psi_D,1).
                                                        \tag{OE.10}\label{p1b:eq:natural-oe-10}
\]
In particular this has not been replaced by an intrinsic filling:
no value of \([\psi_D]\in\pi_7SU(2)\) has been calculated.
The third actual annulus and exact full assembly are
\[
 \mathcal A_D^{(3)}(r,\xi)=\mathcal H_{2-8r}(\xi),
                         \quad1/8\le r\le1/4,         \tag{OE.11}\label{p1b:eq:natural-oe-11}
\]
\[
 \mathcal F_D(r,\xi)=
 \begin{cases}
  \mathcal K_D(8r,\xi),&0\le r\le1/8,\\
  \mathcal H_{2-8r}(\xi),&1/8\le r\le1/4,\\
  \Xi_{2-4r}(\xi),&1/4\le r\le1/2,\\
  \Phi_{2(1-r)}(\xi),&1/2\le r\le1.
 \end{cases}                                          \tag{OE.12}\label{p1b:eq:natural-oe-12}
\]
All three seams and the original outer boundary are exact. Every
projected column on \(1/8\le r\le1\) is now constructed, and both
ends of this annular region have constant column. The radial fractions
are just coordinates, not a completion percentage or a batch bound.

The closed annular projection is now a specific reduced-suspension
map \(g_{\rm ann}:S^8\to S^5\). It is not the whole natural detector
input. For any filling with the actual \eqref{p1b:eq:natural-oe-10} boundary, collapsing its constant-column
boundary gives \(g_{\mathcal K}:S^8\to S^5\). With compatible source
orientations the pasting gives
\([g_D]=[g_{\rm ann}]+[g_{\mathcal K}]\); signs have no effect on the
retained mod-two detector. Its natural coefficient uses the entire
\eqref{p1b:eq:natural-oe-12}, as evaluated by the complete (CI), (AJ) comparison, and
\(\operatorname{Sq}^4u_5=h_tv_9\) on its mapping cone.
In particular a zero column at the new inner face does not justify a
constant-column core or omission of its contribution.

The complete original comparison is evaluated in
\eqref{p1b:eq:natural-aj-4}--\eqref{p1b:eq:natural-aj-7}:
\(h_t=w_t=1\), \(r_J=\zeta_{\rm nat}=0\), with every retained
extension test decided. This component is used with its full frames,
lower divisions and joining; its isolated value is not a replacement
for that complete comparison.


\subsection{The actual root comparison on both primitive cuts}
\label{p1b:sec:natural-root-cut-comparison}

Retain \eqref{p1b:eq:natural-oe-9}, including the whole original \(\Xi_1\). On its inner
stage \(\Phi_1\) has the actual toggle removed by \eqref{p1b:eq:natural-oc-7}.
Here is an exact simplification of that full stage, with both cut seams
retained. It is applied immediately in \eqref{p1b:eq:natural-oe-9}; it does not discard the
second-cut root contribution or assign its top coefficient.

For the actual pair \(x=(-y)e^{\alpha u}\),
\(p=(-y)e^{-\alpha u}\), \(\alpha=3\pi\rho/2\), write
\[
\begin{gathered}
 F_s=F_YR_s,\qquad F_0=F_YR_0,\qquad
 \Gamma_s=F_sF_0^*,\\
 \Pi=ss^*,\qquad P_s=I-\Pi,\qquad
 K_i=K(1/2,x,p),\qquad a=ie^{is}.
\end{gathered}
\tag{OF.1}\label{p1b:eq:natural-of-1}
\]
The \(s\) in the projection notation is the original section, as in
\eqref{p1b:eq:natural-oa-4}; the subscript on \(F_s\) is the circle parameter.
The actual root commutes with \(\Gamma_s\):
\[
 H^3\Gamma_s=\Gamma_sH^3.                              \tag{OF.2}\label{p1b:eq:natural-of-2}
\]
Indeed \(H\) is identity unless both \(x,p\) belong to the original
\(+1\) caps. On that entire double cap, the actual odd extension
requires \(x\) within \(\delta\) of \(aj\), while
\(d(aj,1)=\pi/2\). Since \(2\delta<\pi/2\), both \(R_s\) and
\(R_0\) are identity there. Thus \(\Gamma_s=I\) wherever the root
is nontrivial. This proves the pointwise equation for the complete inner
stage, using the actual support and bundle frames.

Let \(\theta=\arg a\in[0,2\pi]\). On the first primitive cut put
\(K_a=K(\theta/\pi,x,p)\). Expanding the actual first division of
\eqref{p1b:eq:natural-oc-11}, with \(\tau=1\), and using \eqref{p1b:eq:natural-of-2}, gives
\[
 G_1=Q_2^{-1}B_\rho^1 F_s^*(\Pi+aP_s)K_a
       \Gamma_sK_i^*(\Pi-iP_s)F_0(B_\rho^1)^*,
                            \qquad0\le\theta\le\pi.
                                                        \tag{OF.3}\label{p1b:eq:natural-of-3}
\]
The two actual powers \(H^{-3}\) and \(H^3\) from the \(i\)
baseline cancel in this equation. This is a simplification of the
actual divided map, not a replacement of the original fibre root.

On the second primitive cut put
\(K_a^\sigma=\mathsf J K(\theta/\pi-1,p,x)\mathsf J\).
The original exchange equation sends \(H(p,x)^{-3}\) to
\(H(x,p)^3\), so the same expansion gives instead
\[
 G_1=Q_2^{-1}B_\rho^1 F_s^*(\Pi-aP_s)K_a^\sigma
       \Gamma_sH^6K_i^*(\Pi-iP_s)F_0(B_\rho^1)^*,
                            \qquad\pi\le\theta\le2\pi.
                                                        \tag{OF.4}\label{p1b:eq:natural-of-4}
\]
The retained baseline is still \(a=i\) on the first cut.
Consequently the second-cut power is \(H^3H^3=H^6\), not identity.
Moving it through \(\Gamma_s\) uses \eqref{p1b:eq:natural-of-2} and no commutation with
\(K_i\) or with the closing comparison.

Both actual cut seams are exact. At \(a=-1\), the first expression
has \((\Pi-P_s)K(1)=N^2=I\). The second has
\[
 K_0^\sigma=H^{-6},\qquad
 K_0^\sigma\Gamma_sH^6=\Gamma_s.                       \tag{OF.5}\label{p1b:eq:natural-of-5}
\]
At \(a=1\), the first has \(K(0)\Gamma_s=H^6\Gamma_s\),
and the second has \((\Pi-P_s)K_1^\sigma=N^2=I\), followed by
\(\Gamma_sH^6\). Equation \eqref{p1b:eq:natural-of-2} identifies these endpoints too.
Thus neither cut is dropped and the original central-minus closing
equation remains unchanged.

Apply the second actual lower comparison immediately:
\[
 E_1(s,y,\rho,u)=G_1(s,y,\rho,u)
                          G_1(s,1,\rho,u)^*.           \tag{OF.6}\label{p1b:eq:natural-of-6}
\]
Use this exact two-cut formula in the inner stage of \eqref{p1b:eq:natural-oc-12}, then
the same actual \(L_1\) of \eqref{p1b:eq:natural-od-6} and the multipliers in \eqref{p1b:eq:natural-oe-9}.
This gives a pointwise formula for the inner stage of the remaining
intrinsic \(\psi_D\), with no numerical replacement of the original
joining. The root in its closing still enters through \(K_a\) and
the second-cut \(H^6\); no top cancellation is inferred from \eqref{p1b:eq:natural-of-3}.

The complete original comparison is evaluated in
\eqref{p1b:eq:natural-aj-4}--\eqref{p1b:eq:natural-aj-7}:
\(h_t=w_t=1\), \(r_J=\zeta_{\rm nat}=0\), with every retained
extension test decided. This component is used with its full frames,
lower divisions and joining; its isolated value is not a replacement
for that complete comparison.


\subsection{A cut-adapted closed annular comparison of the natural target}
\label{p1b:sec:natural-cut-annular-input}

The original target remains \(h_t,w_t\), with the proved \(w_t=1\).
The executed dependencies here are: rebase the actual circle at \(a=1\),
including its actual projective null; apply the retained neck-frame
comparison; contract the normal direction separately on the two actual
toggle arcs, with its moving sphere-base division; apply the resulting
radial contraction to the whole matrix. The next MAIN proof evaluates
this particular closed annular input. Its intrinsic core is retained.

\paragraph{Rebase the actual complete target, rather than its class.}
Keep \(M_\tau\) in \eqref{p1b:eq:natural-oc-10}, with precisely its actual root, closing,
matching-toggle track, original \eqref{p1b:eq:natural-nd-8}, phase logarithm, frames and bundle
identifications. Put \(\sigma_\ell=-\pi\ell/2\) and define
\[
 \begin{aligned}
 G_\tau^\ell(s,y,n)&=M_\tau(s+\sigma_\ell,-y,n)
                              M_\tau(\sigma_\ell,-y,n)^*,\\
 A_\tau^\ell(s,n)&=G_\tau^\ell(s,1,n),\qquad
 E_\tau^\ell=G_\tau^\ell(A_\tau^\ell)^*.
 \end{aligned}                                                   \tag{OG.1}\label{p1b:eq:natural-og-1}
\]
These are actual sequential divisions on the same bundle. They fix
\(s=0\) and \(y=1\) throughout. At \(Q\), the actual circle-base
matrix is identity: \(C_0(-1,u)\) is diagonal. Thus
\[
 E_\tau^\ell|_Q=E_Q(s+\sigma_\ell,y,u)E_Q(\sigma_\ell,y,u)^*,
 \quad
 \mathcal N_Q^\ell(\kappa)=
 \mathcal N_Q(\kappa,s+\sigma_\ell,y,u)
 \mathcal N_Q(\kappa,\sigma_\ell,y,u)^* .               \tag{OG.2}\label{p1b:eq:natural-og-2}
\]
This pastes on the very same collar as \eqref{p1b:eq:natural-ob-17}, with exact initial,
terminal and parameter-base faces. Let \(\Phi_\tau^\ell\) denote
the supported whole matrix. At \(\tau=0\), every value of the entire
\(\ell\)-homotopy is intrinsic. In particular
\(\Phi_0^0=\operatorname{diag}(\varphi_D,1)\) and
\(\Phi_0^1=\operatorname{diag}(\varphi'_D,1)\) are intrinsically
homotopic, with the same natural coefficient \(h_t\). Its projected
column is constant, so the rebase contributes no detector carry.
Write \(\Phi'_\tau=\Phi_\tau^1\); its actual parameter is now
\(a=e^{is}\), with circle base \(a=1\).

At this base the whole actual toggle is identity, including both
primitive seams, at every \(\tau\). The \(y=1\) slice also misses
both actual toggle caps, just as in \eqref{p1b:eq:natural-oc-15}. Both divided baselines in
\eqref{p1b:eq:natural-og-1} consequently fix \(e_3\). On the inner and inserted-prefix
stages this gives the complete column equation
\[
 \Phi'_\tau e_3=M_\tau(s-\pi/2,-y,n)e_3.               \tag{OG.3}\label{p1b:eq:natural-og-3}
\]
The actual original joining, outer stage and \eqref{p1b:eq:natural-og-2} have column
\(e_3\). No factor of the full matrix is dropped in \eqref{p1b:eq:natural-og-1}.

Apply \eqref{p1b:eq:natural-od-5}--\eqref{p1b:eq:natural-od-7}, with the circle argument shifted by \(-\pi/2\),
immediately to \(\Phi'_1\). Denote its whole-map homotopy by
\(\Xi'_\ell=L'_\ell\Phi'_1(L'_\ell)^*\). On the moving support,
\(Q_2(s-\pi/2)^{-1}B=\operatorname{diag}(e^{-2is},1,1)\).
In the actual ordinary frame \(Q_w=T(e_1,s(w,-w))Q_0\), write
\(D_+(a,w)\) and \(U^\lambda(a,w)\) for the actual positive-line
endpoint and actual toggle track. Define on every \(w\in S^3\)
\[
 \begin{aligned}
 \lambda(\nu)&=\operatorname{clamp}(3-12\nu,0,1),\\
 P'(s,w,\nu)&=\operatorname{diag}(e^{-2is},1,1)
             D_+(e^{is},w)U^{\lambda(\nu)}(e^{is},w).
 \end{aligned}                                                   \tag{OG.4}\label{p1b:eq:natural-og-4}
\]
This is a specified full \(SU(3)\) matrix, with \(P'(0,w,\nu)=I\).
Its column is the exact \(\Xi'_1\) column at \(w=(-y)u\) on the
moving prefix. It is intrinsic everywhere else. The ordinary frame is
defined globally by the retained joining-chart bound; neither the
positive line nor the actual second-cut toggle is replaced.

\paragraph{A direction comparison adapted to the actual two cuts.}
The first toggle arc lies in \((\pi/4,3\pi/4)\), and the second in
\((5\pi/4,7\pi/4)\). Choose the following explicit continuous,
periodic function, with its indicated linear connecting intervals:
\[
 \theta(s)=
 \begin{cases}
 \pi/2-2s,&0\le s\le\pi/4,\\
 0,&\pi/4\le s\le3\pi/4,\\
 2(s-3\pi/4),&3\pi/4\le s\le5\pi/4,\\
 \pi,&5\pi/4\le s\le7\pi/4,\\
 \pi-2(s-7\pi/4),&7\pi/4\le s\le2\pi.
 \end{cases}                                                   \tag{OG.5}\label{p1b:eq:natural-og-5}
\]
Keep the original \(\chi_N\) of \eqref{p1b:eq:natural-od-6}. Put
\[
 \begin{aligned}
 \vartheta_t&=(1-t\chi_N(\nu))\pi/2+t\chi_N(\nu)\theta(s),\\
 q_t&=\cos\vartheta_t+\sin\vartheta_t u,\\
 P'_t&=P'(s,(-y)q_t,\nu)P'(s,-q_t,\nu)^*.
 \end{aligned}                                                   \tag{OG.6}\label{p1b:eq:natural-og-6}
\]
The moving baseline \(P'(s,-q_t,\nu)\) is intrinsic for the entire
homotopy. On the first active arc \(\vartheta_t\le\pi/2\), so
\(\operatorname{Re}(-q_t)\le0\), excluding the actual first \(+1\)
cap. On the second active arc \(\vartheta_t\ge\pi/2\), so
\(\operatorname{Re}(-q_t)\ge0\), excluding the actual second \(-1\)
cap. On the connecting arcs the actual toggle is identity. The baseline
remains in the full formula; its column being constant does not remove
its intrinsic complementary matrix.

The actual initial-column identification is
\(P'_0e_3=\Xi'_1e_3\). Apply it to the whole target immediately:
\[
 Z'_t=P'_t(P'_0)^*\Xi'_1,                              \tag{OG.7}\label{p1b:eq:natural-og-7}
\]
with identity multiplier off \(\chi_N\). This is jointly continuous
on the original normal cell, fixes both parameter-base faces, retains
every complementary matrix, and has column \(P'_te_3\). At its moving
endpoint \(q_1=1\) on the first cut and \(q_1=-1\) on the second.

For a global direction-independent lift use the fixed original axis
\(j\), without changing the source normal framing:
\[
 \widehat q(s)=\cos\theta(s)+\sin\theta(s)j,
 \qquad
 A'(s,y,\nu)=P'(s,(-y)\widehat q(s),\nu)
                           P'(s,-\widehat q(s),\nu)^*.
                                                               \tag{OG.8}\label{p1b:eq:natural-og-8}
\]
On the moving endpoint \(\widehat q=q_1=\pm1\); on its complement
both columns are \(e_3\). Hence \(Z'_1e_3=A'e_3\) everywhere.
This is an exact column identification, not a homotopy assertion about
the complementary intrinsic matrix. Apply the whole radial comparison
\[
 V'_b=A'(s,y,(1-b)\nu+b)A'(s,y,\nu)^*Z'_1.             \tag{OG.9}\label{p1b:eq:natural-og-9}
\]
It is based, fixes the collapsed normal boundary, and is independent
of the absent direction at \(\nu=0\). Its endpoint is intrinsic.
In particular the complete remaining boundary is the actual map
\[
 \boxed{\operatorname{diag}(\psi',1)=
 A'(s,y,1)A'(s,y,\nu)^*P'_1(P'_0)^*\Xi'_1.}            \tag{OG.10}\label{p1b:eq:natural-og-10}
\]
The entire original \eqref{p1b:eq:natural-nd-8} still occurs in \(\Xi'_1\). The closing
and root have not been removed on its inner stage. The known two-cut
root formulas \eqref{p1b:eq:natural-of-3}--\eqref{p1b:eq:natural-of-6} remain valid before rebase; using \eqref{p1b:eq:natural-og-1}
changes the actual baseline, and does not license deleting a root power.

Concatenate \eqref{p1b:eq:natural-og-7} and \eqref{p1b:eq:natural-og-9} in equal time intervals to form
\(\mathcal H'_0=\Xi'_1\),
\(\mathcal H'_1=\operatorname{diag}(\psi',1)\). Paste immediately:
\[
 \mathcal F'_D(r,\xi)=
 \begin{cases}
 \mathcal K'(8r,\xi),&r\le1/8,\\
 \mathcal H'_{2-8r}(\xi),&1/8\le r\le1/4,\\
 \Xi'_{2-4r}(\xi),&1/4\le r\le1/2,\\
 \Phi'_{2(1-r)}(\xi),&1/2\le r\le1,
 \end{cases}
 \quad \mathcal K'|_{S^7}=\operatorname{diag}(\psi',1).
                                                               \tag{OG.11}\label{p1b:eq:natural-og-11}
\]
All displayed seams are exact. The annular projection on
\(1/8\le r\le1\) has constant column on both ends and on all
relative faces. Collapsing those faces constructs an actual map
\[
 g'_{\rm ann}:S^8\longrightarrow S^5.                 \tag{OG.12}\label{p1b:eq:natural-og-12}
\]
No core filling is inferred from this construction. The annular map
is a named, now constructed component of the existing natural detector.
Its complete regular-fibre evaluation follows, including the actual
frame and closing contributions. It uses neither a new general detector
nor a calibration value to assign a natural bit.



\subsection{Parity of the actual closed annular comparison}
\label{p1b:sec:natural-annular-parity}

This calculation evaluates \eqref{p1b:eq:natural-og-12}, retaining both primitive cuts,
the actual coherent frames, the actual prescribed closing and the
odd transition. It does not set the parity of the unfilled \eqref{p1b:eq:natural-og-10}.

The retained characteristic filling in \eqref{p1b:eq:natural-nx-2}, used pointwise in
\eqref{p1b:eq:natural-oc-4}, has last column \(-e_3\) exactly at its characteristic pole
\(v=0\), \(\lambda=3/4\). This follows directly from its displayed
two-plane rotation: its column is \(e_3\) plus a multiple of
\((z_0,z_1,y)\); a zero first pair forces \(z=0\), and only the
positive pole at angle \(\pi/2\) gives \(-e_3\). Its derivative in
\((\lambda-3/4,v)\in\mathbb R\oplus\mathbb R^4\) is invertible.
The actual cap conversion and intrinsic half of \eqref{p1b:eq:natural-oc-4} retain this
statement. Every left factor in the complete column fixes \(e_3\),
so it also fixes \(-e_3\), and cannot create any other inverse image.

Put \(\epsilon=1\) on the first cut, \(\epsilon=-1\) on the second,
\(s_\epsilon=\pi/2,3\pi/2\), and \(\nu_*=3/16\). Solving the
actual equations on all stages of \eqref{p1b:eq:natural-og-11} gives these entire fibres:
\[
 \begin{array}{c|c|c|c}
 \text{stage}&\text{annular parameter}&\nu&y\\ \hline
 \text{inner}&\tau=3/4&\alpha/(3\pi)&
            -\cos\alpha+\epsilon\sin\alpha\,u,
                         \quad0\le\alpha\le\pi/2\\
 \text{prefix}&3/4\le\tau\le1&
       {1+(\tau-3/4)/(1+\tau)\over6}&\epsilon u\\
 \text{frame}&0\le\ell\le1&\nu_*&\epsilon u\\
 \text{direction}&0\le t\le1&\nu_*&-\epsilon q_t^{-1}\\
 \text{radial}&b={\nu_*-\nu\over1-\nu}&0\le\nu\le\nu_*&-1.
 \end{array}                                                   \tag{OH.1}\label{p1b:eq:natural-oh-1}
\]
Here \(s=s_\epsilon\) throughout. In the inner row the actual first
pair has \(x=1,p=e^{-2\alpha u}\); the second has
\(p=1,x=e^{2\alpha u}\). All adjacent rows agree literally at
their seams. The absent \(u\) at \(\alpha=0\) and \(\nu=0\)
caps the two ends. Each fibre is therefore a complete embedded
\(S^3\), comprising two three-balls and the intervening collars.
At the two parameter bases, normal boundary, complementary original
joining, outer stage and projective collar there are no inverse images.

On each row five transverse coordinates are the actual toggle
\((\lambda-3/4,v)\); on the prefix, frame and direction rows
\(\lambda=3-12\nu\), and on the radial row replace \(\nu\) by
\((1-b)\nu+b\). Away from cap origins these have transverse rank
five. At the inner cap use the original Cartesian pair coordinates.
At the radial cap the explicit triangular coordinates are
\(((1-b)|n|+b-\nu_*,n)\), \(n=\nu u\). These give a product tubular
chart although the unsmoothed radius is not differentiable at zero.
To apply smooth Pontryagin--Thom, replace \(\nu\) in the first
coordinate by a monotone \(\rho_\delta(\nu)\), equal to
\(\nu^2/(2\delta)\) near zero and \(\nu\) for \(\nu\ge2\delta\).
Use a cutoff \(\beta(b)\), one near \(b=\nu_*\) and zero near
\(b=0,1\); in the whole radial matrix \eqref{p1b:eq:natural-og-9} replace its numerator's
radius by
\((1-b)((1-t\beta(b))\nu+t\beta(b)\rho_\delta(\nu))+b\).
This applies the smoothing immediately to the whole comparison, fixes
its seams and endpoint, and is a homotopy of the actual \eqref{p1b:eq:natural-og-12}.
Its cap is smooth in \(n\); the same five normal coordinates now
have rank five there. The other piecewise seams can be smoothed in
their product collars, carrying these normal coordinates and framings.
This identifies the full Pontryagin--Thom input of \eqref{p1b:eq:natural-og-12},
with its actual complement, rather than isolated preimages of a model.

\paragraph{Construct the coherent-frame comparison on the whole fibre.}
On the inner ball put \(y_0=e^{-\epsilon\alpha u}\). Its actual
endpoint section is \(e_1\), so
\(C_Y=Q_0^*F_Y(y_0,u,\alpha)\) is intrinsic. Retain this matrix
on the first third of the comparison; on the second third reverse
the actual coherent path \(\beta:\alpha\to0\), then perform the
same actual rotation of the diagonal-section frame as in \eqref{p1b:eq:natural-od-2}.
At every intermediate value apply the original \eqref{p1b:eq:natural-ob-7} column section.
The endpoint is exactly
\[
 \begin{aligned}
 (A,B,C,D)&=\left({1+3\cos\alpha\over3+\cos\alpha},
          {2\sqrt2\over3+\cos\alpha}\operatorname{Im}y_0\right),\\
 C_*(\alpha,u,\epsilon)&=
 \begin{pmatrix}
 (A-iD)^2/\sqrt{A^2+D^2}&B+iC\\
 -B+iC&(A+iD)^2/\sqrt{A^2+D^2}
 \end{pmatrix}.
 \end{aligned}                                                   \tag{OH.2}\label{p1b:eq:natural-oh-2}
\]
It avoids \(-I\), since \(A\ge1/3\): if \(B=C=0\) and
\(D\ne0\), its first entry is nonreal; if \(D=0\) it is positive.
Use its actual principal quaternionic logarithm on the last third.
At \(\alpha=\pi/2\) the entire comparison is exactly \eqref{p1b:eq:natural-od-5} at
\(w=\epsilon\), including its first constant radial leg. At
\(\alpha=0\) it is identity, independent of the absent direction.

For completeness the section chart is valid on this entire three-ball.
The actual intermediate pair is
\(x=e^{-\epsilon(\alpha-\beta)u}\),
\(p=e^{-\epsilon(\alpha+\beta)u}\) for \(\epsilon=1\), with the
order exchanged for \(\epsilon=-1\); here \(0\le\beta\le\alpha\).
Away from its actual basepoint wedge its stereographic lengths are
\(X=\cot((\alpha-\beta)/2)\) and
\(P=\cot((\alpha+\beta)/2)\). Their product is at least one for
\(\alpha\le\pi/2\), so the retained balanced coordinate satisfies
\[
 u_0={X^2+P^2-1\over X^2+P^2+1}\ge1/3.               \tag{OH.3}\label{p1b:eq:natural-oh-3}
\]
For \(\beta<\alpha\), the coefficient of \(u\) in
\(\mathbf u_a\) is nonzero. A vanishing second complex section
coordinate therefore forces \(u_1=u_2=0\), when the third is purely
imaginary and cannot equal real one. At the wedge the section is
\(e_1\). The rotation leg has the same nonpositive real second-column
bound as \eqref{p1b:eq:natural-od-3}. Thus the correction never meets its excluded column.
There is an endpoint-track comparison to account for at the neck.
The actual conjugation \eqref{p1b:eq:natural-od-7} gives the framing path
\(\mathcal C(1-\ell)C_0^*\), whereas the inverse of the ball
contraction gives \(\mathcal C(\ell)^*\). They are related by the
specified relative triangle homotopy
\[
 \mathcal C(1-(1-t)\ell)\mathcal C(t\ell)^*,
       \qquad 0\le t,\ell\le1.                        \tag{OH.4}\label{p1b:eq:natural-oh-4}
\]
At \(\ell=0\) it is \(C_0^*\), at \(\ell=1\) it is identity,
for all \(t\). Its two \(t\)-faces are exactly the two paths.
Apply this actual homotopy to the neck framing, keeping the inner,
prefix, direction and radial pieces fixed. Write \(C_Y(v,h)\) for
the just constructed whole ball contraction. The ensuing explicit
whole-fibre nullhomotopy, with parameter \(h\), is
\(C_Y(v,h)^*\) on the inner ball,
\(\mathcal C(h)^*\) on the prefix,
\(\mathcal C(h+(1-h)\ell)^*\) on the neck-frame track, and identity
on the direction and radial pieces. At every seam these equations
agree; at \(h=0\) they are exactly the track-corrected actual framing,
and at \(h=1\) they are identity. Thus the complete coherent-frame
contribution on both capped regular fibres is null by a specified map.

\paragraph{Retain and compare the actual closing restriction.}
On either free fibre the actual quaternionic commutator root is identity:
one argument is \(1\), whose actual cap value is the central \(-1\).
Thus \(H(1,p)=H(x,1)=I\), by multiplication in the retained word.
For the first fibre the normalized intrinsic restriction is the actual
family
\[
 L_t(p)=\operatorname{diag}(e^{-2\pi it},1)
                     e^{\pi it}\bigl(Q_0^*K(t,1,p)Q_0\bigr)_{12}.
                                                               \tag{OH.5}\label{p1b:eq:natural-oh-5}
\]
At \(t=0\) it is identity, because the prescribed closing has
bottom \(H^6=I\) on this slice. At \(p=-1\), the actual joining
equation gives \(K=\operatorname{diag}(e^{\pi it},e^{-\pi it})\),
so \(L_t(-1)=I\) throughout. This is an actual based homotopy of
the entire free \(S^3\) restriction to identity. Its restriction
at \(t=1/2\) is precisely the closing factor in the first regular
fibre's framing, after the fixed phase is included. On the second
fibre use its prescribed exchanged closing, the same equation with
free \(x\), and the constant retained cap exchange. It gives the
same based nullhomotopy. Thus both actual closing-frame contributions
extend over their full capped fibres and are null. This local comparison
does not delete the second-cut \(H^6\) in the whole \eqref{p1b:eq:natural-og-10}.

\paragraph{Compare the actual odd transition, with its local conjugator.}
The first fibre has \(x=1\); its actual cap centred at \(aj\) is
disjoint from \(1\). Hence its \(R\) is identity on the whole
fibre. On the second fibre the free variable is the entire \(x\in S^3\),
\(a=-i\), and its actual transition is
\[
 R(x)=C_r(x)^*D_a\,\operatorname{diag}(f_j(a^{-1}x),1)
                  D_a^*C_r(x),\qquad D_a=\operatorname{diag}(1,a,1).
                                                               \tag{OH.6}\label{p1b:eq:natural-oh-6}
\]
Here \(r=\arccos(\operatorname{Re}x)/\pi\), with the original
coherent auxiliary \(C_r\) of \eqref{p1b:eq:natural-nc-5}, not a new clutch frame.
The support is the actual three-ball about \(aj=-k\), strictly inside
the inner free ball and away from all its joining seams. If
\(x=aj\exp v\) lies in this ball, apply the actual homotopy
\[
 C_r(x)\longrightarrow C_r(aj\exp((1-\ell)v)),
 \quad f_j(a^{-1}x)\text{ held fixed in \eqref{p1b:eq:natural-oh-6}}.        \tag{OH.7}\label{p1b:eq:natural-oh-7}
\]
The original auxiliary is defined on this entire ball, including the
\(r=1/2\) seam; its original formula gives identical values there.
Off the ball paste identity, since the held cap map is identity on its
boundary. This is an actual intrinsic homotopy of the complete transition
restriction, ending at a fixed conjugate of
\(q_{\rm cap}((aj)^{-1}x)\). The retained cap map is the actual
degree-one three-ball collapse: its radial latitude runs once from
\(-1\) to \(1\) and its imaginary direction is the original sphere
direction. Consequently the complete second-fibre framing transition
has degree \(\pm1\); the first has degree zero. Only parity is used.

\paragraph{Account for cancellation and identify the detector input.}
At a regular fibre, the remaining left matrix is intrinsic and fixes
\(e_3\), so it fixes the fifth transverse coordinate \(ie_3\).
It acts on the other four as the real fundamental \(SU(2)\)
representation. Derivatives of that left matrix give no extra normal
term because it fixes \(-e_3\) identically. The sphere-base division
and right intrinsic matrices also fix \(e_3\) identically in the
full column equation. Thus \eqref{p1b:eq:natural-oh-2}--\eqref{p1b:eq:natural-oh-7} are comparisons of the
actual five-dimensional normal framing, not merely scalar degrees.
Constant phase conjugations retain these null comparisons; their
determinants cancel against the actual determinant factor before
the four-plane framing is read.

After applying the two actual null comparisons, remove the sole odd
transition temporarily to define the two bare framed fibres. Their
piecewise coordinate equations are carried into each other by
\[
 (s,y,\nu,u)\longmapsto(s+\pi,y,\nu,-u).               \tag{OH.8}\label{p1b:eq:natural-oh-8}
\]
For the direction row \(q_t\) changes to \(-q_t\); for the inner
row the two actual pairs are exchanged. The actual toggle conversion
on the second cut is the first conversion conjugated by the retained
constant cap exchange. Its five transverse coordinates therefore
agree up to constant invertible changes and an orientation sign.
The collars and two Cartesian caps agree as well. These local tubular
identifications compare the full bare framed \(S^3\)'s; their tubes
avoid all parameter-base faces. Their stable classes are equal up to
sign. They need not individually bound. Their sum has even coefficient
in \(\pi_8S^5=C_{24}\), and hence zero detector parity.

Restore the actual second-fibre transition by reversing \eqref{p1b:eq:natural-oh-7}. The
change of its four-dimensional normal framing is the displayed
degree-one map into \(SU(2)\subset SO(4)\), with the fifth normal
line fixed. The Pontryagin--Thom difference of these two actual tubular
maps is the attaching map of the rank-five Thom bundle over \(S^4\)
clutched by that very transition. To see the identification, retain the
fibre's \(D^3\times D^5\) tube on its support ball, use the homotopy
on its boundary, and glue the before and after transverse disks there;
their gluing is exactly the normal-framing transition of \eqref{p1b:eq:natural-oh-6}.
The complement is unchanged throughout this local comparison. This
constructs the relevant clutch input, rather than selecting a benchmark.

For the standard clutch-to-Thom-cell construction compare Vera Christina
Gahlen, \emph{Quaternionic line bundles on spin 7-manifolds} (2022),
Lemma 5.1 and its proof on page 38, citing Adams, Lemma 10.1.
The clutch used here is the explicitly identified \eqref{p1b:eq:natural-oh-6}--\eqref{p1b:eq:natural-oh-7}.

That bundle is the real quaternionic line clutched by \eqref{p1b:eq:natural-oh-7}, plus its
fixed trivial line. Its fourth Stiefel--Whitney number is its Euler
number modulo two, namely the degree of that actual clutch, one.
The Thom identity therefore evaluates this actual difference:
\[
 \operatorname{Sq}^4u_5=w_4(\mathbb H\oplus\mathbb R)u_5=v_9.
                                                               \tag{OH.9}\label{p1b:eq:natural-oh-9}
\]
Equivalently it is the suspension of the attaching map of this
quaternionic Thom complex. Combining its parity with the proved even
bare sum gives the value of the whole actual annular input
\[
 \boxed{\operatorname{Sq}^4u_5=v_9\text{ on }C_{g'_{\rm ann}};
              \quad\epsilon(g'_{\rm ann})=1.}          \tag{OH.10}\label{p1b:eq:natural-oh-10}
\]
This is an evaluation of \eqref{p1b:eq:natural-og-12}. Both complete fibres, their actual
frames, the actual closing and the complement have been accounted for.
Coordinate replays check the map equations and all named fibre pieces;
they are not the proof of the topological value.

\paragraph{Apply the evaluated carry to the original two-bit target.}
Keep the exact \eqref{p1b:eq:natural-og-10} intrinsic boundary \(\psi'\). For any actual
core \(\mathcal K'\) satisfying \eqref{p1b:eq:natural-og-11}, the existing detector and
the actual annular pasting give
\[
 \boxed{h_t=1+\epsilon([\mathcal K'e_3]),\qquad
        \zeta_{\rm nat}=\epsilon([\mathcal K'e_3]),\qquad
        {}[\psi']=\zeta_{\rm nat}\mu.}                            \tag{OH.11}\label{p1b:eq:natural-oh-11}
\]
The last equation uses the original fibration's connecting map and
the proved \(w_t=1\), not an assigned value of a primitive map.
The complete original comparison is evaluated in
\eqref{p1b:eq:natural-aj-4}--\eqref{p1b:eq:natural-aj-7}:
\(h_t=w_t=1\), \(r_J=\zeta_{\rm nat}=0\), with every retained
extension test decided. This component is used with its full frames,
lower divisions and joining; its isolated value is not a replacement
for that complete comparison. Changing a core filling adds an even projected coefficient. This annular value alone does not decide the complete coefficient.


\subsection{The actual rebased root terms in the intrinsic core}
\label{p1b:sec:natural-rebased-root}

Apply \eqref{p1b:eq:natural-og-1} to the whole inner-stage root calculation, with
\(a=e^{is}\). The actual baseline is now \(a=1\). Write
\(F_s=F_YR(s-\pi/2)\), \(F_0=F_YR(-\pi/2)\),
\(\Gamma_s=F_sF_0^*\), \(\Pi=ss^*\), \(P=I-\Pi\), using the
actual section at the pair in \eqref{p1b:eq:natural-oa-6}. On the actual root support both
odd transition factors are identity by the same cap separation as
\eqref{p1b:eq:natural-of-2}. Hence \(H^3\Gamma_s=\Gamma_sH^3\). At \(a=1\) the
prescribed target is \(H^3\), rather than the former \(i\)-slice.
The inverse baseline is therefore \(H^{-3}\), on its actual section.

For \(0\le\theta=\arg a\le\pi\), the full first division is
\[
 G'_1=\operatorname{diag}(e^{-2is},1,1)
 F_s^*(\Pi+aP)K(\theta/\pi,x,p)\Gamma_s H^{-6}F_0.
                                                               \tag{OI.1}\label{p1b:eq:natural-oi-1}
\]
The two actual inverse root factors multiply to \(H^{-6}\);
the first-cut carry is not zero. For \(\pi\le\theta\le2\pi\),
exchange sends the numerator root to \(H^3\), which cancels the
inverse actual baseline. Thus the full second division numerator is
\[
 G'_1=\operatorname{diag}(e^{-2is},1,1)
 F_s^*(\Pi-aP)K^\sigma(\theta/\pi-1,x,p)\Gamma_s F_0,
 \quad K^\sigma(t,x,p)=\mathsf J K(t,p,x)\mathsf J.
                                                               \tag{OI.2}\label{p1b:eq:natural-oi-2}
\]
Both equations retain the same actual central-minus closing and its
original bundle frames. At \(a=-1\) the first expression has
\(N K(1)H^{-6}=H^{-6}\), while the second begins at
\(K^\sigma(0)=H^{-6}\). At \(a=1\), the first begins at
\(K(0)H^{-6}=I\), and the second ends at \(N K^\sigma(1)=I\).
The actual frame difference commutes with the retained root at both
seams, so these are exact whole-map seam equations:
\[
 G'_1(-1)^-=G'_1(-1)^+,\qquad
 G'_1(1)^-=G'_1(1)^+=I.                                \tag{OI.3}\label{p1b:eq:natural-oi-3}
\]

Apply the actual second lower division immediately:
\[
 E'_1(s,y,\rho,u)=G'_1(s,y,\rho,u)G'_1(s,1,\rho,u)^*.
                                                               \tag{OI.4}\label{p1b:eq:natural-oi-4}
\]
Use this expression in the inner stage of \(\Phi'_1\), then apply
the actual \eqref{p1b:eq:natural-od-5} rebase conjugation and the two complete multipliers
\eqref{p1b:eq:natural-og-7}--\eqref{p1b:eq:natural-og-9}. The whole intrinsic inner restriction is therefore
\[
 \operatorname{diag}(\psi'_{\rm inner},1)=
 A'(s,y,1)A'(s,y,\nu)^*P'_1(P'_0)^*
 L'_1E'_1(s,y,2\nu,u)(L'_1)^*,\qquad\nu\le1/6.
                                                               \tag{OI.5}\label{p1b:eq:natural-oi-5}
\]
This applies the new comparison to its actual target. In particular,
the local root vanishing on the two annular regular fibres does not
remove the first-cut \(H^{-6}\) from this remaining complete core.
The original v1\_12 formulas with the \(i\)-baseline remain valid;
\eqref{p1b:eq:natural-oi-1}--\eqref{p1b:eq:natural-oi-5} are their specified rebase, not a correction to them.

The remaining equation is exactly
\[
 \mathcal K':D^8\to SU(3),\quad
 \mathcal K'|_{S^7}=\operatorname{diag}(\psi',1),\quad
 \operatorname{Sq}^4u_5=\zeta_{\rm nat} v_9\text{ on }C_{[\mathcal K'e_3]}.
                                                               \tag{OI.6}\label{p1b:eq:natural-oi-6}
\]
Here the whole \(\psi'\) is \eqref{p1b:eq:natural-og-10}, with the original joining
and both complete lower divisions; \eqref{p1b:eq:natural-oi-5} specifies its actual inner
part. The annular carry has been evaluated in \eqref{p1b:eq:natural-oh-10}. The complete original comparison is evaluated in
\eqref{p1b:eq:natural-aj-4}--\eqref{p1b:eq:natural-aj-7}:
\(h_t=w_t=1\), \(r_J=\zeta_{\rm nat}=0\), with every retained
extension test decided. This component is used with its full frames,
lower divisions and joining; its isolated value is not a replacement
for that complete comparison.


\subsection{Intrinsic localization of the actual determinant core}
\label{p1b:sec:natural-exterior-localization}

Keep the whole actual \(\psi'\) of \eqref{p1b:eq:natural-og-10}, the evaluated annular
carry \eqref{p1b:eq:natural-oh-10}, and the unchanged sphere value \(w_t=1\).
The executed dependencies are: identify both exterior frames under
exchange; contract the actual full-circle exterior comparison; extend it
across a collar disjoint from the root support; contract the actual
intrinsic inner direction multiplier; retain and contract the radial
endpoint's actual two-cap prefactor; paste the resulting moving neck
difference into the complete core. The next two MAIN supplements remove
the cap's actual lower baseline and frame, and simplify its prescribed
background relative to every face. None of these comparisons assigns
a core parity from an isolated component. Its complete value is \eqref{p1b:eq:natural-aj-5}.

\begin{proposition}[The exact exterior frame and both actual cuts]
\label{p1b:proposition:natural-exterior-exchange}
Let \(V(x,p)=V^\#(x,p)\) be the actual finite-shift frame of
\eqref{p1b:eq:natural-nt-4}--\eqref{p1b:eq:natural-nt-6}, and let \(d_3=\operatorname{diag}(1,-1,1)\).
Then the following is a pointwise frame equation:
\[
 V(p,x)=\mathsf J V(x,p)d_3,\qquad
 K^{\#\sigma}(t,x,p)=K^\#(t,x,p).
                                                        \tag{OJ.1}\label{p1b:eq:natural-oj-1}
\]
Consequently the whole untoggled exterior determinant, on both primitive
cuts with \(a=e^{is}\), is
\[
 D_Y^0(a,x,p)=V(x,p)\operatorname{diag}(a^2,1,1)V(x,p)^*.
                                                        \tag{OJ.2}\label{p1b:eq:natural-oj-2}
\]
\end{proposition}
\begin{proof}
The actual shift is symmetric in \(x,p\) and applies the same
stereographic translation to both. The exotic section satisfies
\(s(p,x)=\mathsf J s(x,p)\) at every finite shift. Each retained
transporter therefore satisfies the same conjugation equation, and so
does their entire product. The positive vector is the projection of
\(e_2\), which \(\mathsf J\) fixes. The complex cross-product
column transforms with the additional sign \(\det\mathsf J=-1\).
The ordered columns are thus exactly those of \eqref{p1b:eq:natural-oj-1}. Its last constant
factor commutes with the actual diagonal closing path. This proves the
closing equation. Outside the actual ball the root is identity; multiply
the actual half-phase by that closing, on each cut. Exchange and the
unchanged squared phase give precisely \eqref{p1b:eq:natural-oj-2}. No symmetry of the cap
interior closing is asserted.
\end{proof}

Write \(\nu=\rho/2\) in the supported normal cell, and put
\(b(s)=\operatorname{diag}(e^{is},e^{-is},1)\).
Let \(F_s=F_YR(s-\pi/2)\) be the complete actual inner frame at
\(x=(-y)e^{3\pi\nu u},p=(-y)e^{-3\pi\nu u}\).
The actual circle base is \(a=1\). On the exterior \eqref{p1b:eq:natural-oj-2} is identity
there, so the first divided numerator, in its prescribed axes, is
\[
\begin{gathered}
 U(s,y,\nu,u)=b(s)^*C(s,y,\nu,u)^*b(s)C(s,y,\nu,u), \\
 C=V^*F_s, \\
 E'_1=U(s,y,\nu,u)U(s,1,\nu,u)^*.
\end{gathered}
 \tag{OJ.3}\label{p1b:eq:natural-oj-3}
\]
The last column of \(C\) is \(e_3\), so every displayed commutator
is intrinsic. The actual odd transition remains in \(C\), including
its circle dependence. The second lower division is retained literally.

Use the retained based circle disk, in the same fixed two-block,
\(b_\lambda(s)=(c_\lambda e^{is}+s_\lambda j)
(c_\lambda-s_\lambda j)\), where
\(c_\lambda=\cos(\pi\lambda/2)\),
\(s_\lambda=\sin(\pi\lambda/2)\).
Replace both occurrences of \(b\) in \eqref{p1b:eq:natural-oj-3} by \(b_\lambda\),
holding \(C(s,y,\nu,u)\) pointwise fixed; denote the results
\(U_\lambda,E_\lambda^{\rm ext}\). Then
\[
 E_0^{\rm ext}=E'_1,\qquad E_1^{\rm ext}=I,
 \qquad E_\lambda^{\rm ext}|_{s=0\ \mathrm{or}\ y=1}=I.
                                                        \tag{OJ.4}\label{p1b:eq:natural-oj-4}
\]
This is an intrinsic comparison of the actual exterior map. Its circle
periodicity and both primitive-cut seams follow from \eqref{p1b:eq:natural-oj-1}--\eqref{p1b:eq:natural-oj-3};
a partial-circle contraction has not been used.

\paragraph{An applied cap-collar extension.}
Keep the actual round log ball \(B\), \(r_B=1/6\).
Its retained pair-coordinate preimage is an embedded six-ball in the
inner stage. The triangle inequality in the quaternion sphere gives
\[
 \nu\le\frac{\sqrt2r_B}{6\pi}<\frac1{8},\qquad
 \operatorname{dist}(y,-1)\le r_B+\frac{\sqrt2r_B}{2}<\frac12.
                                                        \tag{OJ.5}\label{p1b:eq:natural-oj-5}
\]
In particular this whole preimage misses \(y=1\); the original lower
baseline is exterior. The odd cap centers \(aj\) and their exchanged
mates have distance at least \(\pi/2-r_B\) from these near-real
coordinates, also on the lower baseline. The prescribed odd transitions
are therefore identity on this actual cap and its lower slice, for every
circle value. Here \(\chi_N=0\), \(L'_1=I\), \(\lambda(\nu)=1\),
and the direction multiplier is identity. The radial endpoint still
contains the intrinsic two-cap prefactor \(B_0(s,y)\) defined below.
The cap's actual core is \(B_0E'_1\) until that prefactor is removed
by the explicit whole-map comparison below. The lower division remains
present throughout.

For \(z\in B\), let \(\bar z=r_Bz/|z|\) on its outer collar,
and use the actual inverse pair coordinates of each of these points.
Set \(r_0=3r_B/4=1/8\),
\(\chi(r)=\operatorname{clamp}((r-r_0)/(r_B-r_0),0,1)\).
Define on the entire inner stage
\[
 \widetilde E_\lambda=
 \begin{cases}
 E_\lambda^{\rm ext},&\text{outside }B,\\
 E^{\rm ext}_{\lambda\chi(|z|)}(s,\bar z)
 E^{\rm ext}_0(s,\bar z)^*E'_1(s,z),&z\in B,\ |z|>r_0,\\
 E'_1(s,z),&|z|\le r_0.
 \end{cases}                                           \tag{OJ.6}\label{p1b:eq:natural-oj-6}
\]
In this equation the notation \((s,z)\) always means the retained
pair-coordinate map, with its actual lower slice; no base coordinates
are identified by a degree argument. The entire original root support
lies in \(|z|<\sqrt2/12<r_0\). Thus this extension fixes that
support, pastes exactly at \(r_0,r_B\), and is based throughout.
At \(\lambda=1\) it is identity on the complete inner exterior.
Every matrix is in the same prescribed intrinsic axes. Conjugating
back by the actual ordinary frame gives the corresponding bundle map;
there is no dropped bundle identification.

\paragraph{Remove the actual inner multiplier and paste its neck track.}
For \(\nu\le\nu_0=1/6\), the actual toggle parameter in \eqref{p1b:eq:natural-og-4}
is one, hence its toggle is identity. The radial endpoint has parameter
zero, however, and restores the original intrinsic toggle. Retain
\(B_0(s,y)=A'(s,y,1)A'(s,y,0)^*\); this is independent of
\(\nu,u\) on the entire inner stage. Put
\(g_w=Q_w^*V(w,-w)\), using the original joining frame and positive
line. Its ordinary positive-line prefix is
\(p(s,w)=b(s)^*g_wb(s)g_w^*\).
Replace \(b\) by its actual \(b_\lambda\), call the result
\(p_\lambda\), and retain both based direction prefixes:
\[
 \begin{aligned}
 p^{\rm bas}_\lambda(s,y,q)&=p_\lambda(s,(-y)q)
                                      p_\lambda(s,-q)^*,\\
 T_\lambda&=p^{\rm bas}_\lambda(s,y,q_1)
                                      p^{\rm bas}_\lambda(s,y,q_0)^*,\\
 T_0&=P'_1(P'_0)^*,\qquad T_1=I,
 \quad q_0=u,\quad q_1=q_1(s,\nu,u)\text{ of \eqref{p1b:eq:natural-og-6}}.
 \end{aligned}                                             \tag{OJ.7}\label{p1b:eq:natural-oj-7}
\]
All values are intrinsic. Use the literal complete \(L=L'_1\) and
construct the prefactor's actual nullhomotopy before applying it.

On its first supported four-ball use the actual coordinates
\(v=(4(s-\pi/2)/\pi,\log(-y)/\delta)\), \(|v|<1\),
with \(\delta=1/12\). Both ordinary frames of the actual toggle
there are exactly \(Q_w=T(e_1,s(w,-w))Q_0\), \(w=-y\).
The based lower prefix misses that toggle. Thus its first value is
exactly \(p(s,-y)\eta(v)p(s,-y)^*\). On the second ball,
centered at \(s=3\pi/2\) with the same sphere coordinates, the value
is \(d_3p(s-\pi,-y)\eta(v)p(s-\pi,-y)^*d_3\).
Here \(\eta\) is the \emph{actual} retained supported suspension-Hopf
formula, and the exact second identity follows from \eqref{p1b:eq:natural-oj-1} on the
joining pair. Off these two balls \(B_0=I\).

In three stages define \(B_h\) as follows. First, hold \(\eta(v)\)
fixed and contract each actual conjugator by evaluating it at
\(s=\pi/2+(1-h)\pi v_0/4\),
\(w=\exp((1-h)\delta(v_1,v_2,v_3))\). Its center value is
exactly identity. The second ball retains its constant \(d_3\)
conjugation. Next write \(z=(v_0+iv_1,v_2+iv_3)\),
\(Jz=(-\bar z_1,\bar z_0)\), and replace that second source by
\(\cos(\pi h/2)z+\sin(\pi h/2)Jz\). This preserves its
norm and fixes its based support boundary. The actual Hopf vector
satisfies \(H(Jz)=-H(z)\), so its endpoint is exactly
\(\eta(v)^{-1}\). Simultaneously remove the constant conjugation
using \(\operatorname{diag}(e^{\pi i(1-h)/2},
e^{-\pi i(1-h)/2},1)\), whose initial conjugation is that of
\(d_3\). Finally slide the second, now inverse, circle support
center from \(3\pi/2\) to \(\pi/2\), keeping its actual
\(\log(-y)\) coordinates. Each support has half-width at most
\(\pi/4\), so this slide misses the circle base. At the endpoint
the full product is pointwise \(\eta(v)\eta(v)^{-1}=I\).
Use equal thirds for these three stages. This constructs
\(B_0\to B_1=I\), based at \(s=0\) and \(y=1\), with no
class-only invocation of \(2\eta_3=0\).

Apply the exterior, direction and prefactor comparisons in three
consecutive equal intervals to the actual core:
\[
 \Theta_t^{\rm in}=
 \begin{cases}
 B_0T_0L\widetilde E_{3t}L^*,&0\le t\le1/3,\\
 B_0T_{3t-1}L\widetilde E_1L^*,&1/3\le t\le2/3,\\
 B_{3t-2}L\widetilde E_1L^*,&2/3\le t\le1.
 \end{cases}
 \quad \Theta_0^{\rm in}=\operatorname{diag}(\psi'_{\rm in},1).
                                                        \tag{OJ.8}\label{p1b:eq:natural-oj-8}
\]
Both time seams are exact. The endpoint is identity on the whole
inner exterior; on \(B\) it is \(\widetilde E_1\).

The neck is disjoint from \(B\). Form its actual intrinsic difference
and its specified cutoff, rather than substituting another joining:
\[
 \delta_t(s,y,u)=\Theta_t^{\rm in}(s,y,\nu_0,u)
                         \Theta_0^{\rm in}(s,y,\nu_0,u)^*,
 \quad \kappa(\nu)=\operatorname{clamp}
       \frac{1/4-\nu}{1/4-\nu_0}.
                                                        \tag{OJ.9}\label{p1b:eq:natural-oj-9}
\]
Apply this to the entire actual \eqref{p1b:eq:natural-og-10} matrix:
\[
 \Theta_t(s,y,\nu,u)=
 \begin{cases}
 \Theta_t^{\rm in},&\nu\le\nu_0,\\
 \delta_{t\kappa(\nu)}(s,y,u)\operatorname{diag}(\psi',1),
                                                &\nu\ge\nu_0.
 \end{cases}                                            \tag{OJ.10}\label{p1b:eq:natural-oj-10}
\]
At \(\nu_0\) cancellation proves the exact seam; at \(1/4\)
the multiplier is identity. The entire original \eqref{p1b:eq:natural-nd-8}, original phase,
matching-toggle endpoint track, ordinary frames, lower divisions and
rebased projective null remain in the second expression. All collapsed
and parameter-base faces stay fixed. This is an actual intrinsic
homotopy of the complete seven-sphere core, not just its column.

Let \(\Gamma_B(s,z)=\widetilde E_1(s,z)\) on \(B\), extended
by identity, and let \(\mathcal R\) be identity for \(\nu\le\nu_0\)
and the second expression of \eqref{p1b:eq:natural-oj-10} at \(t=1\) otherwise.
Both are actual based maps \(S^7\to SU(2)\), with disjoint support,
and the endpoint factorization is pointwise:
\[
 \Theta_1=\operatorname{diag}(\Gamma_B\mathcal R,1),
 \qquad [\psi']=[\Gamma_B]+[\mathcal R].              \tag{OJ.11}\label{p1b:eq:natural-oj-11}
\]
This completes the exterior localization on the actual target. It does
not prove that either remaining map has zero or primitive parity. In
particular the joining remainder is retained explicitly by \eqref{p1b:eq:natural-oj-9}--\eqref{p1b:eq:natural-oj-10}.



\subsection{An applied comparison with the actual paired cap word}
\label{p1b:sec:natural-cap-word-identification}

Start with the actual localized \(\Gamma_B\) of \eqref{p1b:eq:natural-oj-6}, \eqref{p1b:eq:natural-oj-11},
not with a primitive map. Write \(z=(z_x,z_p)\),
\(\sigma z=(z_p,z_x)\), and, in the actual northern cap frame, put
\[
 q(z)=Q_s(z)^*H(x,p)^3Q_s(z),\quad
 k(t,z)=X(t,z)\text{ of \eqref{p1b:eq:natural-oa-1}},\quad
 f^\sigma(z)=d_3f(\sigma z)d_3.
 \quad q^\sigma=q^{-1}.
                                                        \tag{OK.1}\label{p1b:eq:natural-ok-1}
\]
Every two-block below fixes the last cap column. The complete cap frame
has the exact exchange equation \(Q_s(p,x)=\mathsf JQ_s(x,p)d_3\).
Thus the last identity is the original ambient root exchange identity
in these actual cap axes. Set \(A_t=\operatorname{diag}
(e^{\pi it},e^{-\pi it},1)\). The retained rebased root equations
\eqref{p1b:eq:natural-oi-1}--\eqref{p1b:eq:natural-oi-2} give the following entire word, including both seams:
\[
\begin{gathered}
W(s,z)=
 \begin{cases}
 A_t^*k(t,z)q(z)^{-2},&s=\pi t,\quad0\le t\le1,\\
 A_t^*k^\sigma(t,z),&s=\pi(t+1),\quad0\le t\le1.
 \end{cases}\\
W(0)=W(2\pi)=I,\quad W(\pi)=q^{-2}.
\end{gathered}
\tag{OK.2}\label{p1b:eq:natural-ok-2}
\]
The first-cut \(q^{-2}=Q_s^*H^{-6}Q_s\) is essential.
Neither half of this word is a closed based map on its own.

On \(\partial B\), let \(g(v)=Q_s^*V^\#\). The actual frame
equations give \(g(\sigma v)=d_3g(v)d_3\); the closing there is
\(gA_tg^*\), on both cuts. Consequently the entire boundary word,
with \(b(s)=\operatorname{diag}(e^{is},e^{-is},1)\), is
\[
 W(s,r_Bv)=b(s)^*g(v)b(s)g(v)^*,\qquad
 N_\lambda(s,v)=b_\lambda(s)^*g(v)b_\lambda(s)g(v)^*.
                                                        \tag{OK.3}\label{p1b:eq:natural-ok-3}
\]
Here \(N_0=W|_{\partial B}\), \(N_1=I\). Define the specified
capped word immediately, using the same \(\chi\) as in \eqref{p1b:eq:natural-oj-6}:
\[
 \widehat W(s,z)=
 \begin{cases}
 N_{\chi(|z|)}(s,z/|z|)N_0(s,z/|z|)^*W(s,z),&|z|>r_0,\\
 W(s,z),&|z|\le r_0.
 \end{cases}
 \quad\widehat W|_{s=0\ \mathrm{or}\ z\in\partial B}=I.
                                                        \tag{OK.4}\label{p1b:eq:natural-ok-4}
\]
This constructs a particular \(S^1\wedge(B/\partial B)=S^7\to SU(2)\)
input. The following comparison identifies it with the actual localized
cap; constructing a closed word alone does not assign its parity.

\paragraph{Retain the exact ordinary frame and lower baseline first.}
The odd cap exclusion in \eqref{p1b:eq:natural-oj-5} makes the actual frame on \(B\)
circle independent. Put \(c(z)=Q_s(z)^*F_Y(z)\). Its last column
is fixed, it extends over the whole ball, and \(c(0)=I\) follows
from the literal retained ordered frames at \((x,p)=(1,1)\).
Let \(V_\lambda(s,z)\) be the actual lower-baseline commutator
\(U_\lambda(s,1,\nu(z),u(z))\) of \eqref{p1b:eq:natural-oj-3}--\eqref{p1b:eq:natural-oj-4}.
Its normal-origin value has the original direction-independent
extension. Thus it is a specified whole-ball intrinsic nullhomotopy,
not a constant inserted for the lower slice. On either cut the exact
ordinary numerator is
\[
 U_c(s,z)=A_t^*c(z)^*A_tW(s,z)c(z),\quad
 E'_1(s,z)=U_c(s,z)V_0(s,z)^*,\quad
 U_c|_{\partial B}=b^*c^*gbg^*c.
                                                        \tag{OK.5}\label{p1b:eq:natural-ok-5}
\]
The matrix \(A_t\) in the first expression is the local half-circle
parameter of \eqref{p1b:eq:natural-ok-2}. Multiplication by its central sign on the second
cut does not change the expression. Direct substitution of \eqref{p1b:eq:natural-oi-1}--\eqref{p1b:eq:natural-oi-2}
proves \eqref{p1b:eq:natural-ok-5}; the literal diagonal circle has not been commuted past
the actual frame.

For any raw cap map \(F(s,z)\), together with a prescribed based
boundary null \(J_\lambda\) from \(F|_{\partial B}\) to identity,
write its actual collar closure as
\(\mathfrak C(F,J)=J_{\chi}J_0^*F\) on the collar and \(F\)
on the inner ball. This is exactly the construction already applied
in \eqref{p1b:eq:natural-oj-6}; it is notation for those full matrices.
For \(0\le h\le1\), define
\[
 F_h=U_cV_h^*,\qquad
 J_{h,\lambda}=b_\lambda^*c(\bar z)^*g b_\lambda g^*c(\bar z)
                         V_{h+(1-h)\lambda}(s,\bar z)^*.
 \quad\Gamma_h=\mathfrak C(F_h,J_{h,\cdot}).
                                                        \tag{OK.6}\label{p1b:eq:natural-ok-6}
\]
Each \(J_{h,0}\) is the boundary of \(F_h\), and each
\(J_{h,1}=I\). At \(h=0\) this is precisely \eqref{p1b:eq:natural-oj-6}, so
\(\Gamma_0=\Gamma_B\). At \(h=1\) it is the closed \(U_c\),
with its actual boundary commutator null. Equation \eqref{p1b:eq:natural-ok-6} removes the
actual lower baseline while moving the entire boundary null with it.
It fixes the circle base and cap boundary throughout.

\paragraph{Apply the actual frame comparison to the full capped map.}
Set \(c_h(z)=c((1-h)z)\); these are the actual retained frame
differences at the displayed ball coordinates, with \(c_0=c\) and
\(c_1=I\). Replace \(c\) by \(c_h\) in both \(U_c\) of \eqref{p1b:eq:natural-ok-5}
and its boundary null. This gives
\[
 U_h=A_t^*c_h^*A_tWc_h,\qquad
 J^c_{h,\lambda}=b_\lambda^*c_h(\bar z)^*g
                            b_\lambda g^*c_h(\bar z),\qquad
 \Gamma^c_h=\mathfrak C(U_h,J^c_{h,\cdot}).
                                                        \tag{OK.7}\label{p1b:eq:natural-ok-7}
\]
Again both boundary equations are exact. Its initial map is the
endpoint of \eqref{p1b:eq:natural-ok-6}, and its terminal map is \(\widehat W\) in \eqref{p1b:eq:natural-ok-4}.
All frames and their bundle identifications occur literally; the frame
was not deleted on the grounds that its restriction has zero class.
Extend both comparisons by identity off \(B\), keeping the actual
\(\mathcal R\) unchanged. They paste to the complete localized
target immediately, because their whole cap boundaries are identity.
Combining with \eqref{p1b:eq:natural-oj-10} constructs the intrinsic whole-map comparison
\[
 \boxed{\psi'\simeq\widehat W\mathcal R,
       \qquad [\psi']=[\widehat W]+[\mathcal R].}
                                                        \tag{OK.8}\label{p1b:eq:natural-ok-8}
\]
Neither map on the right is assigned a value here. In particular the
exchange of the two cuts does not license cancelling them separately:
their intermediate seam in \eqref{p1b:eq:natural-ok-2} is the actual \(q^{-2}\), not
the based identity. The next comparison simplifies \(k\) relative
to all of its prescribed faces and is applied to this same word.



\subsection{Relative simplification of the actual closing background}
\label{p1b:sec:natural-background-reduction}

Retain \(q,k,g,W,\widehat W,\mathcal R\) of \eqref{p1b:eq:natural-ok-1}--\eqref{p1b:eq:natural-ok-8},
and the complete actual bridge \(\mathcal B\) in \eqref{p1b:eq:natural-nz-20}.
The construction below simplifies the actual \(Q_g\) of \eqref{p1b:eq:natural-nw-12},
then applies that relative comparison immediately to the prescribed
closing, both primitive cuts, the complete core and the capped input.

For \(z=r_Brv\), put \(\Psi_g(z)=\psi(U_g(z))\).
The retained characteristic and side equations are pointwise
\(\Psi_g=gA_{1-r}g^*\),
\(E_g=\Psi_g A_{1-r}^*\),
\(G_0=gA_{1-r+tr}g^*A_{1-r+tr}^*\).
Using also \eqref{p1b:eq:natural-nw-11}, cancellation in the complete \eqref{p1b:eq:natural-nw-12} gives
\[
 Q_g(t,z)=\Psi_g(z)A_{(1-t)(1-r)}
                       g(v)A_{tr}g(v)^*A_{(1-t)(1-r)}^*.
                                                        \tag{OL.1}\label{p1b:eq:natural-ol-1}
\]
This is an equality of the full actual matrices; every constant fibre
conversion in \eqref{p1b:eq:natural-nw-5}, \eqref{p1b:eq:natural-nw-10} has already been used in obtaining
\(\Psi_g=gA_{1-r}g^*\).

Define and apply the actual intrinsic family
\[
 Q_{g,h}(t,z)=\Psi_g(z)A_{(1-h)(1-t)(1-r)}
                        g(v)A_{tr}g(v)^*
                        A_{(1-h)(1-t)(1-r)}^*,
 \quad Q_{g,1}=g(v)A_{1-(1-t)r}g(v)^*.
                                                        \tag{OL.2}\label{p1b:eq:natural-ol-2}
\]
The center value is \(\operatorname{diag}(-I_2,1)\).
For example \(\|\Psi_g+I_2\|\le\pi r\) and
\(\|gA_{tr}g^*-I_2\|\le\pi r\) on the two-block,
uniformly in \(t,h,v\); conjugation by the intervening axis matrix
does not change this bound. Hence \eqref{p1b:eq:natural-ol-2} has that continuous central
extension uniformly over the entire parameter square. It retains
exactly the following faces and symmetry:
\[
 Q_{g,h}(0)=\Psi_g,\quad Q_{g,h}(1)=-I_2,\quad
 Q_{g,h}(t,r_Bv)=g(v)A_tg(v)^*,\quad
 Q_{g,h}^\sigma=Q_{g,h}.
                                                        \tag{OL.3}\label{p1b:eq:natural-ol-3}
\]
The last equation follows from the actual \(g^\sigma=g\) in
\eqref{p1b:eq:natural-ok-3} and from \(d_3A_td_3=A_t\). It is an equation for this
background, not for the original bottom bridge or root.

\paragraph{Apply the comparison to the actual target immediately.}
In the same actual cap axes define
\[
\begin{gathered}
 k_h(t,z)=\mathcal B_{\min(4t,1)}(z)\Psi_g(z)^*Q_{g,h}(t,z).\qquad k_0=k,\qquad k_h(0)=q^2,\qquad k_h(1)=-I_2, \\
 k_h|_{\partial B}=gA_tg^*.
\end{gathered}
 \tag{OL.4}\label{p1b:eq:natural-ol-4}
\]
Paste \(Q_sk_hQ_s^*\) to the unchanged \(K^\#\). This gives
a family of actual intrinsic closings with the original endpoints,
the prescribed \(K|_J=K_+\), and
\(F_h(1)=-\sigma F_h(0)\) in the prescribed central-minus sense.
Their determinant representatives retain the actual right matching
toggle. The original \eqref{p1b:eq:natural-nd-8} joining is unchanged, since its entire
locus is disjoint from \(B\). Thus the original complete triple is
compared through actual triples, including its retained sphere piece.

Replace \(k\) by \(k_h\) in both cuts of \eqref{p1b:eq:natural-ok-2}. Their shared
midpoint remains exactly \(q^{-2}\), their circle endpoints remain
identity, and their actual boundary null \eqref{p1b:eq:natural-ok-3} is unchanged. Apply
the same collar closure \eqref{p1b:eq:natural-ok-4} to obtain \(\widehat W_h\).
The whole comparison \eqref{p1b:eq:natural-ok-8} therefore ends at the explicit input
\[
 \begin{aligned}
 \psi'&\simeq\widehat W_1\mathcal R,\\
 W_1(\pi t,z)&=A_t^*\mathcal B_{\min(4t,1)}(z)
       \Psi_g(z)^*g(v)A_{1-(1-t)r}g(v)^*q(z)^{-2},\\
 W_1(\pi(t+1),z)&=A_t^*\mathcal B_{\min(4t,1)}^\sigma(z)
       \Psi_g(z)^*g(v)A_{1-(1-t)r}g(v)^*.
 \end{aligned}                                             \tag{OL.5}\label{p1b:eq:natural-ol-5}
\]
Use the already proved central extensions at \(r=0\), and the literal
\(\chi\)-collar of \eqref{p1b:eq:natural-ok-4}. The actual source/bundle/frame bridge
\(\mathcal B\) remains in both cuts. In particular no assumption
\(\mathcal B^\sigma=\mathcal B\) is made. For \(t\ge1/4\)
both bridges equal \(\Psi_g\) and cancel their following inverses,
whereas the first-cut \(q^{-2}\) still remains. The terms on the
two cuts are not separate based maps and cannot be cancelled by adding
their primitive coefficients.

For an explicit whole-core implementation before localization, keep
the actual \(c,V_0\) of \eqref{p1b:eq:natural-ok-5}, replace \(W\) by \(W_h\) there,
and set
\[
\psi'_h|_B=B_0(s,y)U_{c,h}V_0^*,\qquad
 \psi'_h|_{S^7\setminus B}=\psi'.
                                                        \tag{OL.6}\label{p1b:eq:natural-ol-6}
\]
Here the actual prefactor \(B_0\) is the retained radial-endpoint
factor of \eqref{p1b:eq:natural-oj-8}, before its constructed nullhomotopy is applied.
The cap boundary is unchanged by \eqref{p1b:eq:natural-ol-3}--\eqref{p1b:eq:natural-ol-4}, so this pastes
exactly and is an intrinsic whole-map homotopy. The cap remains away
from all parameter-base faces; every original lower slice, joining and
projective collar remains present. Equations \eqref{p1b:eq:natural-ol-5}--\eqref{p1b:eq:natural-ol-6} are the
immediate application of this named prerequisite to the actual target.

\paragraph{The complete split detector input.}
The actual input identification \(\psi'\simeq\widehat W_1\mathcal R\)
has now been constructed, including every changed boundary null,
frame factor and lower slice. Their complete sum is evaluated in \eqref{p1b:eq:natural-aj-4}--\eqref{p1b:eq:natural-aj-5}. The equivalent
full filling equation is \eqref{p1b:eq:natural-oi-6}. An equivalent split input for the \emph{existing}
detector is
\[
\begin{gathered}
 K_B,K_R:D^8\longrightarrow SU(3),\\
 K_B|_{S^7}=\operatorname{diag}(\widehat W_1,1),\qquad
 K_R|_{S^7}=\operatorname{diag}(\mathcal R,1),\\
 g_B=[K_Be_3],\quad g_R=[K_Re_3]:S^8\to S^5,\\
 \operatorname{Sq}^4u_5=\epsilon(g_B)v_9\text{ on }C_{g_B},\\
 \operatorname{Sq}^4u_5=\epsilon(g_R)v_9\text{ on }C_{g_R},\\
 \boxed{\begin{gathered}
 \zeta_{\rm nat}=\epsilon(g_B)+\epsilon(g_R),\\
 h_t=1+\epsilon(g_B)+\epsilon(g_R),\quad w_t=1.
 \end{gathered}}
\end{gathered}
\tag{OL.7}\label{p1b:eq:natural-ol-7}
\]
In the first boundary equation \(\widehat W_1\) is inserted by the
actual cap embedding already constructed in \eqref{p1b:eq:natural-oj-5}, \eqref{p1b:eq:natural-oj-11}; both maps
are thereby on the original \(S^7\). The disjoint-support factorization
and its actual whole-map homotopies justify their addition in
\(\pi_7SU(2)\). The existing stabilized detector identifies that
addition with parity addition. Genuine changes of either filling add
the retained even projected image and cannot assign an odd parity by
choice. A pinch construction after compressing the two actual supports
pastes such fillings and the proved intrinsic comparisons to the full
original \eqref{p1b:eq:natural-og-11} input. No new detector is required or constructed.

The complete original comparison is evaluated in
\eqref{p1b:eq:natural-aj-4}--\eqref{p1b:eq:natural-aj-7}:
\(h_t=w_t=1\), \(r_J=\zeta_{\rm nat}=0\), with every retained
extension test decided. This component is used with its full frames,
lower divisions and joining; its isolated value is not a replacement
for that complete comparison.


\subsection{Removal of the complete actual inserted joining prefix}
\label{p1b:sec:natural-prefix-flattening}

Keep the actual core \(\psi'\), the actual closing and its central-minus
equation, the right matching toggle, the original \eqref{p1b:eq:natural-nd-8} joining, and
the completed sphere comparison \(w_t=1\). The dependencies executed
here are: enlarge the existing intrinsic frame and direction cutoffs
where the original column is already fixed; cancel the complete prefix
toggle against its actual inverse ordinary prefix; apply the already
constructed inner localization across the resulting constant prefix.
No joining interior is replaced.

Write \(\nu_0=1/6\), \(\nu_1=1/4\), \(\nu_2=1/3\).
Replace the actual \(\chi_N\) of \eqref{p1b:eq:natural-og-6} by
\[
 \widetilde\chi_N(\nu)=
 \begin{cases}
 \chi_N(\nu),&\nu\le\nu_0,\\
 1,&\nu_0\le\nu\le\nu_1,\\
 (\nu_2-\nu)/(\nu_2-\nu_1),&\nu_1\le\nu\le\nu_2,\\
 0,&\nu\ge\nu_2,
 \end{cases}
 \qquad \chi_{N,h}=(1-h)\chi_N+h\widetilde\chi_N.
                                                        \tag{OM.1}\label{p1b:eq:natural-om-1}
\]
In the actual \eqref{p1b:eq:natural-od-6} conjugator use \(\chi_{N,h}\chi_P\), retaining
the same actual neck contraction and the same pole cutoff. Call the
result \(L_h\). In the actual direction prefix use
\[
\begin{gathered}
 q_{1,h}=\cos\vartheta_h+\sin\vartheta_h u,\qquad \vartheta_h=(1-\chi_{N,h})\pi/2+\chi_{N,h}\theta(s), \\
 P_{1,h}=p(s,(-y)q_{1,h},\nu)p(s,-q_{1,h},\nu)^*.
\end{gathered}
 \tag{OM.2}\label{p1b:eq:natural-om-2}
\]
Here \(p\) is the full actual \eqref{p1b:eq:natural-og-4} prefix, including its actual
\(U^{\lambda(\nu)}\); the actual moving lower value is retained.
Keep \(P_0,A_\nu,A_1\) unchanged and apply the comparison to the
complete supported boundary \(\Phi'_1\):
\[
 \psi'_h=A_1A_\nu^*P_{1,h}P_0^*L_h\Phi'_1L_h^*,
 \qquad \psi'_0=\psi'.                                 \tag{OM.3}\label{p1b:eq:natural-om-3}
\]
On the entire moving-column region \(1/6<\nu<5/24\), the two
cutoffs are already exactly one; \eqref{p1b:eq:natural-om-3} equals the retained intrinsic
core there for every \(h\). On the changed complement all its
prefixes and \(\Phi'_1\) are intrinsic. Consequently \eqref{p1b:eq:natural-om-3} is an
intrinsic homotopy on the whole source. It fixes both parameter-base
faces and the collapsed normal boundary. It is unchanged for
\(\nu\le\nu_0\) and \(\nu\ge\nu_2\). In particular it retains
the actual root cap and all original bundle identifications.

\paragraph{Cancel full matrices, including the lower value.}
On \(\nu_0\le\nu\le\nu_1\), set \(w=(-y)u\), \(w_b=-u\),
\(Q_w=T(e_1,s(w,-w))Q_0\),
\(D_s=Q_w^*D_+(e^{is},w,-w)Q_w\), and let \(U_\lambda\)
denote the actual toggle in these axes. On its entire nonidentity
support the actual odd transition is identity, since \(w\) lies in
one of the real pole caps while the odd clutch centers are purely
imaginary. Also \(L_1=q_i C_0q_i^*\), with
\(q_i=\operatorname{diag}(e^{-2is},1,1)\), exactly there.
The actual normalized \(y\)-base matrix \(M_b\) is intrinsic and
independent of \(\lambda\): its pair \((w_b,-w_b)\) misses both
toggles, and the actual circle base \(a=1\) misses them too.
Thus on this support the complete expressions are
\[
 \begin{aligned}
 \Xi_\lambda&=q_iD_sU_\lambda C_0M_b^*q_i C_0^*q_i^*,\\
 P_0&=q_iD_sU_\lambda D_b^*q_i^*,\\
 P_0^*\Xi_\lambda&=q_iD_b C_0M_b^*q_i C_0^*q_i^*.
 \end{aligned}                                             \tag{OM.4}\label{p1b:eq:natural-om-4}
\]
The actual lower value \(D_b\), the complete \(M_b\), and the actual
\(C_0\) have not been deleted. Off this support \(U_\lambda=I\),
so the same independence follows without changing the retained
partial conjugator. This proves independence of the full product,
not only of its last column.

On both active circle arcs \(q_{1,1}\) is exactly the same real
quaternion as the retained radial prefix. Hence \(A_\nu^*P_{1,1}=I\)
there. Off those arcs every toggle is identity, so that product equals
its actual untoggled positive-line expression. The radial endpoint is
still \(A_1=B_0A_0\), with the exact \(B_0\) of \eqref{p1b:eq:natural-oj-8}, not its
replacement by identity. All remaining untoggled factors and
\(L_1,\Phi'_1\) on this prefix are independent of \(\nu\).
Consequently the following is an exact whole-matrix equation:
\[
 \boxed{\psi'_1(s,y,\nu,u)=\psi'(s,y,\nu_0,u)
       \quad(\nu_0\le\nu\le\nu_1).}                  \tag{OM.5}\label{p1b:eq:natural-om-5}
\]
At \(\nu_1\) use the original \eqref{p1b:eq:natural-nd-8} initial endpoint, not a new
joining path. Its equality to the end of the inserted track proves
the same equation at that seam.

Apply \(\Theta^{\rm in}_t\) and its actual \(\delta_t\) of
\eqref{p1b:eq:natural-oj-8}--\eqref{p1b:eq:natural-oj-9} immediately, this time over the whole flattened prefix:
\[
 \begin{aligned}
 k_J(\nu)&=
 \begin{cases}1,&\nu_0\le\nu\le\nu_1,\\
 (\nu_2-\nu)/(\nu_2-\nu_1),&\nu_1\le\nu\le\nu_2,\\
 0,&\nu\ge\nu_2,
 \end{cases}\\
 T_t&=\Theta^{\rm in}_t\quad(\nu\le\nu_0),\qquad
 T_t=\delta_{t k_J(\nu)}\psi'_1\quad(\nu\ge\nu_0).
 \end{aligned}                                             \tag{OM.6}\label{p1b:eq:natural-om-6}
\]
Equation \eqref{p1b:eq:natural-om-5} proves that \(T_1=I\) on the entire inserted
prefix. Its seams at \(\nu_0,\nu_1,\nu_2\) are exact. The original
three-leg localization, including the actual radial-endpoint prefactor
null, has been applied; no constant-column inference is used.
Define \(\mathcal R_J=I\) for \(\nu\le\nu_1\), and let it be
the second expression of \eqref{p1b:eq:natural-om-6} at \(t=1\) otherwise. Then
\[
 T_1=\Gamma_B\mathcal R_J,\qquad
 \psi'\simeq\widehat W_1\mathcal R_J
       \quad\hbox{intrinsically on the original }S^7.   \tag{OM.7}\label{p1b:eq:natural-om-7}
\]
The actual \eqref{p1b:eq:natural-ok-6}--\eqref{p1b:eq:natural-ok-7} lower/frame comparisons and \eqref{p1b:eq:natural-ol-4} background
reduction are applied to \(\Gamma_B\) in the last equation. Their
cap boundaries are identity, so the new \(\mathcal R_J\) is retained
unchanged. It starts at the original joining, with its actual changed
neck/direction multipliers, and equals the original outer and projective
tail for \(\nu\ge\nu_2\). Its zero class follows from the complete
\eqref{p1b:eq:natural-aj-4}--\eqref{p1b:eq:natural-aj-5} calculation, not from a pointwise identity here.

The retained parity equation of the existing detector is
\[
 {}[\psi']=\zeta_{\rm nat}\mu,
 \quad\zeta_{\rm nat}=\epsilon(K'e_3),\quad h_t=1+\zeta_{\rm nat},\quad w_t=1,
 \qquad K'|_{S^7}=\operatorname{diag}(\psi',1).
                                                        \tag{OM.8}\label{p1b:eq:natural-om-8}
\]
The following constructions complete specific cap-seam and boundary-null
components of this actual filling input. They do not assign either
natural coefficient from \eqref{p1b:eq:natural-om-7} without its complete evaluation
in \eqref{p1b:eq:natural-aj-4}--\eqref{p1b:eq:natural-aj-5}.



\subsection{An actual stabilized root-seam comparison}
\label{p1b:sec:natural-root-seam-stabilization}

Keep the exact \(q=Q_s^*H^3Q_s\), \(k_1\) of \eqref{p1b:eq:natural-ol-4}, its
prescribed faces, and the actual \(\widehat W_1\) in \eqref{p1b:eq:natural-om-7}.
Construct the stabilized null of its actual \(q^2\), then attach it
to both complete cuts. In the same actual cap axes, concatenate
the completed \eqref{p1b:eq:natural-nz-20} bridge with its identified characteristic filling:
\[
 C_h(z)=
 \begin{cases}
 \mathcal B_{2h}(z),&0\le h\le1/2,\\
 \Lambda_{2\pi(1-h)}(U_g(z)),&1/2\le h\le1,
 \end{cases}
 \quad C_0=q^2,\quad C_1=I,\quad C_h|_{\partial B}=I.
                                                        \tag{ON.1}\label{p1b:eq:natural-on-1}
\]
Here \(\Lambda\) is the actual retained \eqref{p1b:eq:natural-nx-2} filling in its
specified fibre conversion, with
\(\Lambda_\pi(U_g)=\psi(U_g)=\mathcal B_1\).
Thus the time seam is a proved pointwise identity. All actual source,
frame and constant identifications in \eqref{p1b:eq:natural-nz-1}--\eqref{p1b:eq:natural-nz-20} remain in
\(\mathcal B\); a degree-calibrated source is not inserted.
The actual \(U_g\) has the prescribed south value on \(\partial B\),
where \(\Lambda\) is identity. This proves the spatial face too.
For \(h>1/2\) this is a stabilized comparison; it is not assumed
intrinsic.

Set \(A_t=\operatorname{diag}(e^{\pi it},e^{-\pi it},1)\).
The full actual two-cut family is
\[
 V_h(s,z)=
 \begin{cases}
 A_t^*C_{h(1-t)}q^{-2}k_1(t,z)C_h^*,&s=\pi t,\\
 A_t^*C_{h(1-t)}^*q^2 k_1^\sigma(t,z),&s=\pi(t+1).
 \end{cases}                                             \tag{ON.2}\label{p1b:eq:natural-on-2}
\]
The first-cut right inverse \(C_h^*\) is essential. Direct full-matrix
multiplication, using \(k_1(0)=q^2\),
\(k_1^\sigma(0)=q^{-2}\), and \(k_1(1)=k_1^\sigma(1)=-I_2\),
gives all three circle equations:
\[
 V_h(0)=V_h(2\pi)=I,\qquad
 V_h(\pi)^-=V_h(\pi)^+=C_h^*,\qquad V_0=W_1.
                                                        \tag{ON.3}\label{p1b:eq:natural-on-3}
\]
This uses \(C_h\) literally. No unproved relation
\(C_h^\sigma=C_h^*\) is imposed on the constructed nullhomotopy.
At the spatial boundary \(C_h=I\), so the boundary word of every
\(V_h\) is exactly the unchanged actual \(N_0(s,v)\) of \eqref{p1b:eq:natural-ok-3}.
Apply its original collar immediately:
\[
 \widehat V_h=N_{\chi(|z|)}N_0^*V_h\quad(|z|>r_0),
 \qquad \widehat V_h=V_h\quad(|z|\le r_0),
 \quad\widehat V_0=\widehat W_1.
                                                        \tag{ON.4}\label{p1b:eq:natural-on-4}
\]
This is a whole \(SU(3)\) homotopy relative to the circle base and
the cap boundary. Extend by identity off the actual cap and multiply
by the unchanged \(\mathcal R_J\). It is immediately a comparison
of the complete original-source input, not of an isolated primitive.
Its complete projected column is
\[
 e_{h}(s,z)=N_{\chi}N_0^*V_h(s,z)e_3
 \quad\hbox{on the collar, and }V_he_3\hbox{ on the inner ball}.
                                                        \tag{ON.5}\label{p1b:eq:natural-on-5}
\]
Every factor of \eqref{p1b:eq:natural-on-2} is retained in this equation. In particular the
first right inverse can move the column and cannot be deleted in a
detector calculation.

Let \(V=V_1\), \(\Omega_B=\widehat V_1\).
Although the raw midpoint is now identity, the whole capped midpoint is
\[
 V(\pi,z)=I,\qquad
 \Omega_B(\pi,z)=N_{\chi(|z|)}(\pi,z/|z|)
                       \quad(|z|>r_0).                 \tag{ON.6}\label{p1b:eq:natural-on-6}
\]
It is generally nonidentity. The exact matrix replay and a concrete
nonidentity collar value are recorded in
the archived evidence file \nolinkurl{Paper1B-NaturalDeterminantPrefixSeamDiagnostics-Appendix-v1_0.tex}.
Thus \eqref{p1b:eq:natural-on-1}--\eqref{p1b:eq:natural-on-5} complete an actual stabilized seam and its whole
annular projection, while \eqref{p1b:eq:natural-on-6} identifies a real remaining collar
term. The next supplement constructs and applies its boundary-disk
comparison. Neither the stabilized seam nor a vanishing raw midpoint
assigns \(h_t\) or \(\zeta_{\rm nat}\).



\subsection{Closed actual cuts, their collar, and the retained filling equation}
\label{p1b:sec:natural-boundary-disk-comparison}

Retain the actual \(g(v)=Q_s^*V^\#\), the full \(V\) of \eqref{p1b:eq:natural-on-2},
and the original \(b_\lambda(s)\) boundary disk. The dependencies
executed here are: construct a stabilized null for each actual half
boundary; compare it with the original full-circle null, including
the centralizer loop; apply that relative comparison to the entire
capped target; retain its exact collar remainder and paste the complete
new annuli into the existing detector input.

Let \(B^{12}_\lambda(s)=b_\lambda(s)\), and use the same retained
B32 formula in the \(23\) block, denoted
\(B^{23}_\lambda(s,\alpha)\) with imaginary axis
\(\cos\alpha\,j+\sin\alpha\,k\). Write
\[
 d(s)=\operatorname{diag}(e^{is},e^{is},e^{-2is}),\quad
 Z_\lambda(s)=B^{12}_\lambda(s)B^{23}_\lambda(2s,0),\quad
 L_\lambda(t)=d(\pi t)B^{23}_\lambda(-2\pi t,0).
                                                        \tag{OO.1}\label{p1b:eq:natural-oo-1}
\]
The latter retains both ends \(L_\lambda(0)=I\),
\(L_\lambda(1)=N=\operatorname{diag}(-1,-1,1)\).
Its initial path is the actual \(A_t\), and its final path is in
the centralizer of the actual \(g\). Hence
\(L_\lambda(t)^*gL_\lambda(t)g^*\) is a specified stabilized null
of the actual half boundary, identity at both half endpoints.
The \(Z\) path starts at \(d\), ends at identity, and is based in
\(s\). The extra segment \(\mathcal C(Z_\lambda,g)\), where
\(\mathcal C(a,g)=a^*gag^*\), has identity endpoints but is retained.

Construct the full based loop lift by pasting
\[
 \widetilde L_\lambda(s)=
 \begin{cases}
 L_\lambda(s/\pi),&0\le s\le\pi,\\
 L_\lambda(s/\pi-1)N,&\pi\le s\le2\pi.
 \end{cases}
 \quad\widetilde L_0=b,\quad\widetilde L_1=d.
                                                        \tag{OO.2}\label{p1b:eq:natural-oo-2}
\]
It is a continuous full circle, even while its \(23\) block moves.
Both endpoints of its commutator on each half remain identity.

\paragraph{A complete relative comparison of the actual disks.}
First define the literal square of based matrix loops
\[
 T(f_0,f_1;s)=B^{12}_{f_0}(s)B^{23}_{f_0}(2s,0)
                                      B^{23}_{f_1}(2s,0)^*.
                                                        \tag{OO.3}\label{p1b:eq:natural-oo-3}
\]
Its diagonal is the old disk \(b_\lambda\). Interpolate its diagonal
path \((\lambda,\lambda)\) to
\((\max(2\lambda-1,0),\min(2\lambda,1))\).
This fixes both disk faces and the circle base. The new first segment is
\(d B^{23}_{2\lambda}(2s,0)^*\); its second is \(Z_{2\lambda-1}\).
For \(0\le r\le1\), the explicit constant conjugation
\[
\begin{gathered}
 E_{\lambda,r}(s)=\gamma_{\lambda r}
 B^{23}_\lambda(-2s,\pi r)\gamma_{\lambda r}^*, \\
 \gamma_{\lambda r}=\exp((\pi\lambda r/2)j)_{23},\qquad E_{\lambda,0}=B^{23}_\lambda(-2s,0), \\
 E_{\lambda,1}=B^{23}_\lambda(2s,0)^*
\end{gathered}
 \tag{OO.4}\label{p1b:eq:natural-oo-4}
\]
compares the inverse disk to the negative-angle disk. Apply it backwards
on the first segment, keeping the second segment fixed. Finally rotate
its \(23\) axis on the second circle half from \(j\) to \(-j\).
At \(s=\pi\) that block is identity for every disk parameter, so
this is a continuous based homotopy relative to both disk faces.
The resulting first segment is precisely \(\widetilde L\) in \eqref{p1b:eq:natural-oo-2}.
Use equal thirds for these three comparisons; call the full lift
\(A_{r,\lambda}(s)\). The accompanying implementation gives their
literal piecewise formula, including every parameter seam.
Consequently
\[
 \begin{aligned}
 J_{r,\lambda}(s,v)&=\mathcal C(A_{r,\lambda}(s),g(v)),\\
 J_{0,\lambda}&=N_\lambda,\quad J_{r,0}=N_0,\quad
 J_{r,1}=I,\quad J_{r,\lambda}(0)=I,\\
 J_{1,\lambda}&=
 \begin{cases}
 \mathcal C(\widetilde L_{2\lambda},g),&\lambda\le1/2,\\
 \mathcal C(Z_{2\lambda-1},g),&\lambda\ge1/2.
 \end{cases}
 \end{aligned}                                             \tag{OO.5}\label{p1b:eq:natural-oo-5}
\]
All equations are pointwise for the actual \(g\), with its retained
bundle/frame identification. This constructs the disk comparison,
rather than replacing one null disk by another without its track.

Apply it immediately to the full actual raw \(V\): on the collar
use \(J_{r,\chi}N_0^*V\), and on the inner ball use \(V\).
Call this \(\Omega_{B,r}\); its initial map is the exact
\(\Omega_B\) of \eqref{p1b:eq:natural-on-6}. Also define the short, half-closed cap by
the same formula with
\(J^S_\lambda=\mathcal C(\widetilde L_\lambda,g)\).
Call it \(\Omega^S_B\). The actual relative equations give
\[
 \Omega^S_B(0,z)=\Omega^S_B(\pi,z)=\Omega^S_B(2\pi,z)=I,
 \qquad \Omega^S_B|_{\partial B}=I.                    \tag{OO.6}\label{p1b:eq:natural-oo-6}
\]
Its two cuts are therefore actual closed maps on the same source:
\(\Omega_+\) equals \(\Omega^S_B\) on the first half and identity
on the second; \(\Omega_-\) is defined conversely. They are
\(SU(3)\) maps; no intrinsic endpoint or primitive value is assumed.
Retain the exact additional collar
\[
 \Delta=J_{1,\chi}(J^S_\chi)^*\quad(|z|>r_0),\qquad
 \Delta=I\quad(|z|\le r_0),\qquad
 \boxed{\Omega_{B,1}=\Delta\Omega_+\Omega_-.}
                                                        \tag{OO.7}\label{p1b:eq:natural-oo-7}
\]
It is identity on the cap boundary and circle base, but generally
nonidentity in the collar. A concrete retained nonidentity value and
the equation replay are recorded in
the archived evidence file \nolinkurl{Paper1B-NaturalDeterminantPrefixSeamDiagnostics-Appendix-v1_0.tex}.
The exact equality in \eqref{p1b:eq:natural-oo-7} retains all possible contributions and
cancellations. Neither the two cuts nor \(\Delta\) is discarded.

\paragraph{Pasting into the existing original-source detector.}
Insert each cap map by the actual \eqref{p1b:eq:natural-oj-5} embedding and keep
\(\mathcal R_J\) of \eqref{p1b:eq:natural-om-6}. The resulting actual inner boundary is
\[
 \Psi_*=\Delta\Omega_+\Omega_-\mathcal R_J:S^7\to SU(3).
                                                        \tag{OO.8}\label{p1b:eq:natural-oo-8}
\]
Let \(\mathfrak I\) concatenate \eqref{p1b:eq:natural-om-3}, \eqref{p1b:eq:natural-om-6}, the actual
\eqref{p1b:eq:natural-ok-6}--\eqref{p1b:eq:natural-ok-7} lower/frame comparisons and the actual \eqref{p1b:eq:natural-ol-4}
background reduction. It is intrinsic, starts at \(\psi'\), and
ends at \(\widehat W_1\mathcal R_J\). Let \(\mathfrak S\) be
the whole-cap \eqref{p1b:eq:natural-on-4} comparison, retaining \(\mathcal R_J\), and
let \(\mathfrak D\) be the whole-cap \eqref{p1b:eq:natural-oo-5} comparison, again
retaining that remainder. All intermediate source and bundle maps
are the original ones. The exact remaining map is now
\[
 \boxed{K_*:D^8\to SU(3),\qquad K_*|_{S^7}=\Psi_*.}    \tag{OO.9}\label{p1b:eq:natural-oo-9}
\]
Given this actual map, its assembly with the constructed comparisons is
\[
 K'(r,\xi)=
 \begin{cases}
 K_*(4r,\xi),&0\le r\le1/4,\\
 \mathfrak D_{2-4r}(\xi),&1/4\le r\le1/2,\\
 \mathfrak S_{3-4r}(\xi),&1/2\le r\le3/4,\\
 \mathfrak I_{4-4r}(\xi),&3/4\le r\le1.
 \end{cases}
 \quad K'|_{S^7}=\operatorname{diag}(\psi',1).
                                                        \tag{OO.10}\label{p1b:eq:natural-oo-10}
\]
Every seam and the full outer boundary are constructed and checked.
The complete matrices and projected columns on \(1/4\le r\le1\)
are supplied. The inner column \(\Psi_*e_3\) is not assumed constant;
it cannot be collapsed to produce a separate projected sphere map.
The existing detector must be applied to the full projection after
\eqref{p1b:eq:natural-oo-9} is constructed:
\(g'=[K'e_3]:S^8\to S^5\),
\(\operatorname{Sq}^4u_5=\zeta_{\rm nat} v_9\) on \(C_{g'}\),
\(h_t=1+\zeta_{\rm nat}\), \(w_t=1\).

The complete original comparison is evaluated in
\eqref{p1b:eq:natural-aj-4}--\eqref{p1b:eq:natural-aj-7}:
\(h_t=w_t=1\), \(r_J=\zeta_{\rm nat}=0\), with every retained
extension test decided. This component is used with its full frames,
lower divisions and joining; its isolated value is not a replacement
for that complete comparison.


\subsection{Cancellation of the identified pair with its root difference retained}
\label{p1b:sec:natural-paired-cut-cancellation}

Keep the entire \eqref{p1b:eq:natural-oo-8} input, including its actual \(\Delta\),
the original joining remainder and all cap identifications.
The dependencies executed here are the pointwise exchange equation for
the actual cuts, an explicit null of its identified pair, and immediate
application to the complete original-source boundary and filling annuli.
No independent parity is assigned to either stabilized cut.

Let \(\sigma(z_x,z_p)=(z_p,z_x)\),
\(D=\operatorname{diag}(1,-1,1)\), and
\(C_h^\sigma(z)=DC_h(\sigma z)D\), for the literal \eqref{p1b:eq:natural-on-1} null.
Write \(P(t,z)=\Omega_B^S(\pi t,z)\). On the actual spatial collar put
\(F(s,z)=J_\chi^S(s,z/|z|)N_0(s,z/|z|)^*\), and put \(F=I\)
on the inner ball. The retained cap and exterior frame equations give
\[
 q^\sigma=q^{-1},\qquad
 g(\sigma v)=Dg(v)D,\qquad
 F(\pi(t+1),z)=DF(\pi t,\sigma z)D.
                                                        \tag{OP.1}\label{p1b:eq:natural-op-1}
\]
For the last equation, the second half of the disk lift is
\(L_\chi(t)N\), where \(N=\operatorname{diag}(-1,-1,1)\)
centralizes \(g\). Its conjugation on the moving \(23\) block is
the same as that of \(D\). Both conjugations commute with \(A_t\)
and \(d(s)\). This proves the equation for the entire collar,
including its unchanged \(N_0\); it is not an assumed symmetry
of the original closing interior.

Substituting the actual complete formulas \eqref{p1b:eq:natural-on-2}, and cancelling the
literal closing only between the two identified expressions, gives
\[
 \begin{aligned}
 \Gamma_C(\pi(t+1),z)
 &=F(\pi(t+1),z)A_t^*C_{1-t}^*
                       (C^\sigma_{1-t})^*A_tF(\pi(t+1),z)^*,\\
 \Gamma_C(s,z)&=I\quad(0\le s\le\pi),\\
 \Omega_B^S(\pi(t+1),z)
 &=\Gamma_C(\pi(t+1),z)\,DP(t,\sigma z)D.
 \end{aligned}                                             \tag{OP.2}\label{p1b:eq:natural-op-2}
\]
The right inverse and every collar factor are retained.
At \(t=0\), \(C_1=I\). At \(t=1\),
\(C_0^*(C_0^\sigma)^*=q^{-2}q^2=I\).
On \(\partial B\), both null factors are identity.
Thus \(\Gamma_C\) is an actual closed second-cut map,
not a deleted root correction.

Define the identified pair \(\mathcal P\) to be \(P(s/\pi,z)\)
on the first half, and \(DP(s/\pi-1,\sigma z)D\) on the second.
Its midpoint and both circle endpoints are identity. The supports of
\(\Omega_+\) and \(\Gamma_C\) are disjoint, so the exact
whole-matrix equality is
\[
 \boxed{\Psi_*=\Delta\Gamma_C\mathcal P\mathcal R_J.}
                                                        \tag{OP.3}\label{p1b:eq:natural-op-3}
\]
There is no commutation of overlapping nonidentity terms in this equation.

\paragraph{An explicit relative null of this actual pair.}
First undo the two-half pinch by interpolating
\((\min(2\theta,1),\max(2\theta-1,0))\) to
\((\theta,\theta)\), where \(\theta=s/(2\pi)\).
This gives \(P(\theta,z)DP(\theta,\sigma z)D\).
Remove its constant conjugation through
\(\operatorname{diag}(e^{i\pi(1-h)},1,e^{-i\pi(1-h)})\).
Its initial conjugation is that of \(-D\), hence of \(D\);
its final value is identity.

In the actual log coordinates put
\(a=(z_x-z_p)/\sqrt2,\ b=(z_x+z_p)/\sqrt2\).
The literal orthogonal family
\[
 \sigma_h(a,b)=
 \bigl(-a_0,\,
       \cos(\pi(1-h))a_1-\sin(\pi(1-h))a_2,\,
       \sin(\pi(1-h))a_1+\cos(\pi(1-h))a_2,\,
       b\bigr)
                                                        \tag{OP.4}\label{p1b:eq:natural-op-4}
\]
starts at the actual swap and ends at the reflection \(a_0\mapsto-a_0\).
Apply it only inside the retained cap map.
This is a homotopy of its source ball fixing the collapsed boundary,
not a change of the original normal framing or generator.

For fixed \(v=(a_1,a_2,b)\), let
\(\ell(v)=\sqrt{r_B^2-|v|^2}\) and
\(u=(1+a_0/\ell)/2\). Write \(P_v(u)=P(\theta,z(u,v))\).
It is identity at \(u=0,1\). The remaining pair is
\(P_v(u)P_v(1-u)\). Interpolate
\[
 \begin{aligned}
 u_1^h&=(1-h)u+h\min(2u,1),\\
 u_2^h&=(1-h)u+h\max(2u-1,0),\\
 P_v(u)P_v(1-u)&\longrightarrow
              P_v(u_1^h)P_v(1-u_2^h)\\
 &\longrightarrow
              P_v\bigl((1-h)\min(2u,2-2u)\bigr)
              \longrightarrow I.
 \end{aligned}                                             \tag{OP.5}\label{p1b:eq:natural-op-5}
\]
The middle arrow is the exact endpoint equation for the two
forward/backtracking halves. If \(\ell=0\), every value is the
original cap-boundary identity. Uniform continuity of the actual
closed \(P\) proves continuity there. The circle base remains fixed.
Use equal fifths for pinch removal, constant conjugation, \eqref{p1b:eq:natural-op-4},
the first comparison in \eqref{p1b:eq:natural-op-5}, and its backtracking contraction.
Call the resulting complete homotopy \(\mathcal P_h\).
It starts at the exact \eqref{p1b:eq:natural-op-3} factor and ends at identity.
In particular \(\mathcal P_{1-r}\) is a constructed \(D^8\) filling
of this named actual factor, without an existence-only appeal
to \(\pi_7SU(3)=0\).

Apply it immediately on the original source:
\[
 \Psi_h=\Delta\Gamma_C\mathcal P_h\mathcal R_J,\qquad
 \Psi_0=\Psi_*,\quad
 \Psi_1=\Delta\Gamma_C\mathcal R_J.
                                                        \tag{OP.6}\label{p1b:eq:natural-op-6}
\]
All bundle and cap insertions are those already retained in \eqref{p1b:eq:natural-oj-5}.
The collar and \(\Gamma_C\) are unchanged. The latter is generally
nonidentity; it cannot be inferred to vanish from the paired cancellation.
The complete factorization replay retains a concrete nonidentity root
word; the archived evidence file \nolinkurl{Paper1B-NaturalDeterminantNormTailDiagnostics-Appendix-v1_0.tex}
records its value and the exact seam checks.

The exact remaining filling at this stage is
\[
 K_{\rm flat}:D^8\to SU(3),\qquad
 K_{\rm flat}|_{S^7}=\Delta\Gamma_C\mathcal R_J.
                                                        \tag{OP.7}\label{p1b:eq:natural-op-7}
\]
Paste it with the full new comparison and the existing three annuli:
\[
 K'(r,\xi)=
 \begin{cases}
 K_{\rm flat}(5r,\xi),&0\le r\le1/5,\\
 \Psi_{2-5r}(\xi),&1/5\le r\le2/5,\\
 \mathfrak D_{3-5r}(\xi),&2/5\le r\le3/5,\\
 \mathfrak S_{4-5r}(\xi),&3/5\le r\le4/5,\\
 \mathfrak I_{5-5r}(\xi),&4/5\le r\le1.
 \end{cases}
 \qquad K'|_{S^7}=\operatorname{diag}(\psi',1).
                                                        \tag{OP.8}\label{p1b:eq:natural-op-8}
\]
Every newly required annular matrix and complete projected column is
constructed. The whole boundary and closed-projection parity are evaluated by
the retained exact-map chain through \eqref{p1b:eq:natural-aj-4}--\eqref{p1b:eq:natural-aj-5}, carrying all
these constructed annuli.
The filled paired factor does not independently assign a natural bit.
The complete original comparison is evaluated in
\eqref{p1b:eq:natural-aj-4}--\eqref{p1b:eq:natural-aj-7}:
\(h_t=w_t=1\), \(r_J=\zeta_{\rm nat}=0\), with every retained
extension test decided. This component is used with its full frames,
lower divisions and joining; its isolated value is not a replacement
for that complete comparison.


\subsection{The literal inverse root norm and the centralizer collar}
\label{p1b:sec:natural-norm-remainder-reduction}

Start with the actual \eqref{p1b:eq:natural-op-7} boundary, not a primitive replacement.
Remove the conjugators of its literal \(\Gamma_C\) through based
relative homotopies, compress its exact two-disk collar difference,
and apply both constructions to the complete original-source input.
These steps retain the root norm and the centralizer-loop term.

Put
\[
 T_C(t,z)=C_{1-t}(z)^*(C^\sigma_{1-t}(z))^*,\qquad
 F_h=J^S_{(1-h)\chi}N_0^*,\qquad
 \Gamma_{C,h}=F_hA_{(1-h)t}^*T_CA_{(1-h)t}F_h^*
                                                        \tag{OQ.1}\label{p1b:eq:natural-oq-1}
\]
on the second cut and collar, with \(F_h=I\) on the inner ball
and the whole word identity on the first cut.
The endpoint \(F_1=I\) is the exact equation
\(J^S_0=N_0\). Therefore
\[
 \Gamma_{C,0}=\Gamma_C,\qquad
 \Gamma_{C,1}=\mathcal T_C,\quad
 \mathcal T_C(\pi(t+1),z)=T_C(t,z),\quad
 \mathcal T_C|_{s\le\pi}=I.
                                                        \tag{OQ.2}\label{p1b:eq:natural-oq-2}
\]
The root norm is identity on \(t=0,1\) and on the entire spatial
boundary, by the actual root equations in \eqref{p1b:eq:natural-op-1}--\eqref{p1b:eq:natural-op-2}.
Hence these homotopies remain based there even though their
conjugators need not be identity there. On the inner collar face
\(\chi=0\), \(F_h=I\) exactly, proving the spatial paste.
Every conjugator is retained during its removal; none is
declared null as a seven-sphere map.

For the actual long and short disks in \eqref{p1b:eq:natural-oo-5}--\eqref{p1b:eq:natural-oo-7}, set
\[
 \ell_h(\chi)=(1-h)\chi+h\min(2\chi,1),\qquad
 \Delta_h=J_{1,\chi}(J^S_{\ell_h(\chi)})^*.
                                                        \tag{OQ.3}\label{p1b:eq:natural-oq-3}
\]
The inner and outer collar faces, and circle base, are fixed.
At \(h=0\) this is exactly \(\Delta\).
At \(h=1\), the first half of the long disk cancels pointwise
against its own short disk, while its second half remains:
\[
 \Delta_1=\mathcal Z_g=
 \begin{cases}
 I,&\chi\le1/2,\\
 Z_{2\chi-1}(s)^*g(v)Z_{2\chi-1}(s)g(v)^*,&\chi\ge1/2.
 \end{cases}
                                                        \tag{OQ.4}\label{p1b:eq:natural-oq-4}
\]
Here \(Z_\lambda\) is the exact \eqref{p1b:eq:natural-oo-1} matrix, not a substituted
null of its component. At \(\lambda=0\), \(Z_0=d\) centralizes
the actual \(g\); at \(\lambda=1\), \(Z_1=I\).
These prove both new collar faces. The nonidentity commutator
between them is retained.

Apply these two relative comparisons immediately:
\[
 \mathfrak Q_h=
 \begin{cases}
 \Delta\Gamma_{C,2h}\mathcal R_J,&0\le h\le1/2,\\
 \Delta_{2h-1}\mathcal T_C\mathcal R_J,&1/2\le h\le1,
 \end{cases}
 \quad
 \mathfrak Q_0=\Delta\Gamma_C\mathcal R_J,\quad
 \mathfrak Q_1=\mathcal Z_g\mathcal T_C\mathcal R_J.
                                                        \tag{OQ.5}\label{p1b:eq:natural-oq-5}
\]
All cap insertions use the original \eqref{p1b:eq:natural-oj-5} coordinates.
The joining and projective tail remain unchanged. The literal
complete matrices and their third columns are supplied.

The characteristic half of \(C\) also has a proved actual
exchange equation. In the retained \(S^4\) coordinates
\(U_g=(y,w)\), the frame equation gives
\(U_g(\sigma z)=(y,MD_2M^*w)\), where
\(D_2=\operatorname{diag}(1,-1)\) and \(M\) is the original
fibre conversion. The \eqref{p1b:eq:natural-nx-2} lift is equivariant under this
constant action; its conversion into the cap axes is exactly
conjugation by \(D\). Thus, on this half only,
\[
 C_h^\sigma=C_h\quad(1/2\le h\le1),\qquad
 T_C(t,z)=\Lambda_{2\pi t}(U_g(z))^{-2}
                              \quad(0\le t\le1/2).
                                                        \tag{OQ.6}\label{p1b:eq:natural-oq-6}
\]
No such equation is assumed for the intrinsic bridge half.
It starts at the actual \(q^2\), whose exchange is \(q^{-2}\).
The next supplement constructs an intrinsic square null at
this tail's seam and retains its full characteristic comparison.

The exact remaining whole-core equation is now
\[
 \boxed{K_{\rm red}:D^8\to SU(3),\qquad
 K_{\rm red}|_{S^7}=\mathcal Z_g\mathcal T_C\mathcal R_J.}
                                                        \tag{OQ.7}\label{p1b:eq:natural-oq-7}
\]
Its full assembly with every constructed annulus is
\[
 K'(r,\xi)=
 \begin{cases}
 K_{\rm red}(6r,\xi),&0\le r\le1/6,\\
 \mathfrak Q_{2-6r}(\xi),&1/6\le r\le2/6,\\
 \Psi_{3-6r}(\xi),&2/6\le r\le3/6,\\
 \mathfrak D_{4-6r}(\xi),&3/6\le r\le4/6,\\
 \mathfrak S_{5-6r}(\xi),&4/6\le r\le5/6,\\
 \mathfrak I_{6-6r}(\xi),&5/6\le r\le1.
 \end{cases}
 \qquad K'|_{S^7}=\operatorname{diag}(\psi',1).
                                                        \tag{OQ.8}\label{p1b:eq:natural-oq-8}
\]
The complete original comparison is evaluated in
\eqref{p1b:eq:natural-aj-4}--\eqref{p1b:eq:natural-aj-7}:
\(h_t=w_t=1\), \(r_J=\zeta_{\rm nat}=0\), with every retained
extension test decided. This component is used with its full frames,
lower divisions and joining; its isolated value is not a replacement
for that complete comparison.



\subsection{An intrinsic tail and its identified characteristic comparison}
\label{p1b:sec:natural-root-norm-tail}

Keep the exact \eqref{p1b:eq:natural-oq-7} target. Construct an intrinsic null of its
characteristic-square seam, apply it to the actual root norm,
retain the entire difference as a specified characteristic comparison,
and evaluate the projection degrees of that comparison and the
actual centralizer collar. These are components of the actual input,
not independently substituted benchmark maps.

For \(b=(y,w_0,w_1)\in S^4\), write \(Jw=(-\overline w_1,\overline w_0)\),
and retain the original \(\psi=M^{-1}\phi M\) in the cap block.
The real orthogonal family and its intrinsic square null are
\[
 B_hb=(y,\cos(\pi h/2)w+\sin(\pi h/2)Jw),\qquad
 \mathcal J_h(b)=\psi(b)^*\psi(B_hb)^*,\qquad
 \mathcal J_0=\psi^{-2},\quad\mathcal J_1=I.
                                                        \tag{OR.1}\label{p1b:eq:natural-or-1}
\]
The real inner product of \(w,Jw\) is zero, so the family stays
on the original \(S^4\). Their Hopf projections are opposite;
therefore \(\psi(y,Jw)=\psi(y,w)^*\), proving the final equation
pointwise. Both poles give identity in the product for every \(h\).
All values lie in the retained intrinsic cap \(SU(2)\).

Apply this immediately with the actual \(U_g(z)\):
\[
 T_{\rm int}(t,z)=
 \begin{cases}
 \mathcal J_{1-2t}(U_g(z)),&0\le t\le1/2,\\
 C_{1-t}^*(C_{1-t}^\sigma)^*,&1/2\le t\le1.
 \end{cases}
                                                        \tag{OR.2}\label{p1b:eq:natural-or-2}
\]
Both seam values are the literal \(\psi(U_g)^{-2}\).
The second part uses the actual intrinsic bridge of \eqref{p1b:eq:natural-on-1},
not a selected substitute. Its end is \(q^{-2}q^2=I\);
the first part starts at \(I\). Both parts are identity on
\(\partial B\). Insert this map on the second cut and identity
on the first, obtaining the actual closed intrinsic
\(\mathcal T_{\rm int}:S^7\to SU(2)\).
Its top class is not assigned.

The comparison retained by this replacement is the exact closed map
\[
 a:([0,\pi]/\partial[0,\pi])\wedge(S^4,\mathrm{south})
                                      \longrightarrow SU(3),
 \qquad
 a(\theta,b)=\Lambda_\theta(b)^{-2}
                          \mathcal J_{1-\theta/\pi}(b)^*.
                                                        \tag{OR.3}\label{p1b:eq:natural-or-3}
\]
Both time endpoints and the entire south arc are identity,
so its domain is the specified \(S^5\). In the original cap insert
\(a(2\pi t,U_g(z))\) for \(0\le t\le1/2\), and identity elsewhere;
call it \(\mathcal A_g\). The complete factorization is pointwise:
\[
 \boxed{\mathcal T_C=\mathcal A_g\mathcal T_{\rm int},\qquad
 \Phi_{\rm red}=\mathcal Z_g\mathcal A_g
                          \mathcal T_{\rm int}\mathcal R_J.}
                                                        \tag{OR.4}\label{p1b:eq:natural-or-4}
\]
In particular the stabilized square path has not been replaced by
its intrinsic endpoint null without retaining its comparison.
The entire projected column of \(a\) is
\(a(\theta,b)e_3=\Lambda_\theta(b)^{-2}e_3\).

\paragraph{Degrees of the identified actual comparison factors.}
Orient the characteristic suspension by
\(d\theta\wedge\operatorname{vol}_{S^4}\), where
the ambient \(S^4\) order is
\((y,\operatorname{Re}w_0,\operatorname{Im}w_0,
       \operatorname{Re}w_1,\operatorname{Im}w_1)\).
Use the original complex target orientation. The regular value
\(-e_3\) of \(ae_3:S^5\to S^5\) has exactly the two preimages
\[
 (\theta,b)=(\pi/4,\mathrm{north}),\
             (3\pi/4,\mathrm{north}),\qquad
 \deg(ae_3)=-2.
                                                        \tag{OR.5}\label{p1b:eq:natural-or-5}
\]
Here is a source-bound regular-value calculation. In the
unconverted geodesic two-plane, the off-axis column vanishes only
at a pole, a time endpoint, or when the two-plane trace vanishes.
In the last case its squared inverse is
\(-e^{-i(1+y)\theta/2}I\). In the open interval and away from
the poles this cannot be \(-I\). At the south pole the value is
\(e_3\); at the north pole its axial value is \(e^{-4i\theta}\).
This proves the exact preimage list. The derivative in the four
north coordinates is complex multiplication, up to the retained
unitary fibre conversion, by
\[
 k(\theta)=-2\sin\theta\cos(3\theta/2)e^{-2i\theta}.
                                                        \tag{OR.6}\label{p1b:eq:natural-or-6}
\]
It is nonzero at both preimages. The axial derivative is \(+4i\)
at the value \(-e_3\), whose induced target chart has the opposite
orientation. Thus each local degree is \(-1\).
The fixed cap conversion preserves the real target orientation.

The literal collar factor also has a universal expression in its
actual variables, with no replacement of its clutch:
\[
 z(s,\lambda,g)=Z_\lambda(s)^*gZ_\lambda(s)g^*,\qquad
 \mathcal Z_g(s,R_Brv)=
 z(s,2\chi(R_Br)-1,g(v))
 \quad(\chi\ge1/2).
                                                        \tag{OR.7}\label{p1b:eq:natural-or-7}
\]
It is identity on both \(\lambda\) faces, the circle base and
\(g=I\), giving a specified \(S^2\wedge SU(2)=S^5\) map.
Orient that source by \(ds\wedge d\lambda\) and the standard
\((\operatorname{Re}a,\operatorname{Im}a,
\operatorname{Re}b,\operatorname{Im}b)\) sphere orientation of
\(\left(\begin{smallmatrix}a&b\\-\bar b&\bar a\end{smallmatrix}\right)\).
The regular value \(e_1\) of \(ze_3\) has two preimages:
\[
 s=\pi/2,\ 3\pi/2,\quad \lambda=1/2,\quad
 g'=B_\lambda^{12}(s)^*gB_\lambda^{12}(s)
      =\begin{pmatrix}0&1\\-1&0\end{pmatrix},
 \qquad \deg(ze_3)=+2.
                                                        \tag{OR.8}\label{p1b:eq:natural-or-8}
\]
Indeed, if \(B^{23}_\lambda(2s)\) has two-block entries
\((\alpha,\beta)\), its column has first coordinate \(b\beta\).
The value \(e_1\) forces \(|b|=|\beta|=1\).
Its literal disk formula gives exactly the two listed points,
with \(\alpha=0,\beta=1,a=0,b=1\).
At either point, writing \(g'\) in the oriented coordinates
\((a_R,a_I,b_I)\), the first-order column is
\((1+i(b_I+\delta s),-\pi\delta\lambda-i\delta s,a_R-i a_I)\).
Its target-chart Jacobian in \((s,\lambda,a_R,a_I,b_I)\)
has determinant \(+\pi\), proving both signs.

These degrees identify actual component maps, not a natural
two-bit assignment. Their original source maps and embeddings remain
in \eqref{p1b:eq:natural-or-3}--\eqref{p1b:eq:natural-or-7}. In particular \(g\mapsto Mg e_1\) in the
actual \(U_g\) suspension has degree \(-1\) in the displayed
orientations: \(g e_1=(a,-\bar b)\) reverses one real coordinate,
while \(M\) preserves orientation. Opposite displayed projection
degrees therefore do not justify deleting the two actual factors.
Any cancellation requires their source and comparison tracks.

The exact remaining full filling is
\[
 \boxed{K_{\rm red}:D^8\to SU(3),\qquad
 K_{\rm red}|_{S^7}=
       \mathcal Z_g\mathcal A_g\mathcal T_{\rm int}\mathcal R_J.}
                                                        \tag{OR.9}\label{p1b:eq:natural-or-9}
\]
The endpoint in \eqref{p1b:eq:natural-or-9} is the exact \eqref{p1b:eq:natural-oq-7} boundary, so every
constructed annulus in \eqref{p1b:eq:natural-oq-8} applies unchanged. Its whole projected
column is retained. No independent sphere projection is produced
by collapsing the generally nonconstant inner boundary column.
For a filling with this retained boundary, the original detector
requires
\[
 g'=[K'e_3]:S^8\to S^5,\qquad
 \operatorname{Sq}^4u_5=\zeta_{\rm nat} v_9\text{ on }C_{g'},\qquad
 h_t=1+\zeta_{\rm nat},\quad w_t=1.
                                                        \tag{OR.10}\label{p1b:eq:natural-or-10}
\]
The complete original comparison is evaluated in
\eqref{p1b:eq:natural-aj-4}--\eqref{p1b:eq:natural-aj-7}:
\(h_t=w_t=1\), \(r_J=\zeta_{\rm nat}=0\), with every retained
extension test decided. This component is used with its full frames,
lower divisions and joining; its isolated value is not a replacement
for that complete comparison. The matrix, source and regular-value checks are recorded in the archived evidence file \nolinkurl{Paper1B-NaturalDeterminantNormTailDiagnostics-Appendix-v1_0.tex}; they corroborate the displayed formulas rather than certify the homotopy coefficient.


\subsection{The actual characteristic and collar on one clutch source}
\label{p1b:sec:natural-common-clutch-source}

Retain the whole \eqref{p1b:eq:natural-or-9} input and all its original source and bundle
identifications. The dependencies executed here are to remove the
actual based-direction arc of its characteristic factor, move the
collar's radial and circle coordinates to the same source, and paste
those full comparisons into the original filling annuli. The actual
intrinsic norm and original joining remainder are retained.

Write \(r=|z|/r_B\), \(g=g(v)\), \(v=z/|z|\), and
\[
\begin{gathered}
 b(r,g)=(1-2r,\,2\sqrt{r(1-r)}\,Mg e_1),\qquad a_*(t,r)=a(\pi t,b(r,I)), \\
 \widetilde a(t,r,g)=a(\pi t,b(r,g))a_*(t,r)^* .
\end{gathered}
 \tag{OS.1}\label{p1b:eq:natural-os-1}
\]
The two pole values are independent of \(g\). In particular
\(\widetilde a=I\) on \(r=0,1\), \(t=0,1\), and \(g=I\).
The original characteristic comparison has not been asserted to
have these based-direction faces: its north arc is nonconstant.

Put \(t(s)=\operatorname{clamp}(2(s/\pi-1),0,1)\) on
\(0\le s\le2\pi\). The following removes that precise arc:
\[
 \mathcal A_h(s,z)=
 \mathcal A_g(s,z)a_*(t(s),r)^*
                  a_*(t(s),h+(1-h)r),\quad
 \mathcal A_0=\mathcal A_g,\quad
 \mathcal A_1=\widetilde a(t(s),r,g(v)).
                                                        \tag{OS.2}\label{p1b:eq:natural-os-2}
\]
At the cap center its value is \(a_*(t(s),h)\), independent of
the absent direction. At the cap boundary all factors are identity.
The circle base and both ends of the active cut are fixed. These
equations and uniform continuity on the compact pole models prove a
jointly continuous relative homotopy of full \(SU(3)\) matrices.
Its last column is retained, rather than declared constant.

The actual collar initially has radial coordinate
\(\lambda_0(r)=\operatorname{clamp}(8r-7,0,1)\). Use the two
consecutive source comparisons
\[
 \lambda_h(r)=(1-h)\lambda_0(r)+hr,\qquad
 \theta_h(s)=(1-h)s+h\,2\pi t(s).
                                                        \tag{OS.3}\label{p1b:eq:natural-os-3}
\]
First replace its radial coordinate by \(\lambda_h\), retaining \(s\).
Then replace \(s\) by \(\theta_h\), retaining \(r\).
Both radial faces give identity for every \(g\); the lifted circle
coordinates start at zero and end at \(2\pi\).
Consequently these are actual relative source homotopies.
At \(r=0\) their value is identity uniformly in direction.
The old interior half-circle seams need not be identity throughout
this comparison: their whole collar values are retained.

Apply \eqref{p1b:eq:natural-os-2} and then the two parts of \eqref{p1b:eq:natural-os-3} to the complete map,
in three equal intervals. Denote the full comparison by
\(\mathfrak C_h\). Its literal endpoint is
\[
 \mathfrak C_0=\mathcal Z_g\mathcal A_g\mathcal T_{\rm int},
 \qquad
 \mathfrak C_1=\mathfrak b(t(s),r,g(v))\mathcal T_{\rm int},
 \quad
 \boxed{\mathfrak b(t,r,g)=
 z(2\pi t,r,g)\widetilde a(t,r,g).}
                                                        \tag{OS.4}\label{p1b:eq:natural-os-4}
\]
Every face of the specified \(S^2\wedge SU(2)=S^5\) domain of
\(\mathfrak b\) is identity. This identifies the complete pair
pointwise on the actual double-suspended clutch source. It does not
substitute a map having only a matching component degree.
On the original source use \(\mathfrak C_h\mathcal R_J\) and keep
the original \eqref{p1b:eq:natural-nd-8} throughout. The cap and joining seams remain
identity, so this is an immediate full-input application.

\paragraph{Orientations and the next comparison equation.}
Use \(dt\wedge dr\wedge\operatorname{vol}_{SU(2)}\).
For the unreduced \(S^4\) suspension in \eqref{p1b:eq:natural-os-1}, its oriented volume
in the product chart is
\(-dy\wedge\operatorname{vol}_{S^3_w}\).
Here \(dy=-2dr\), and \(g\mapsto Mg e_1\) has the actual degree
minus one proved in \eqref{p1b:eq:natural-or-8}. Thus this source conversion has degree
minus one. The full arc homotopy \eqref{p1b:eq:natural-os-2} descends before the
based-direction arc is collapsed; it preserves the source degree.
Together with \eqref{p1b:eq:natural-or-5}, this proves
\[
 \deg(\widetilde a e_3)=2,\qquad
 \deg(z e_3)=2,\qquad
 \deg(\mathfrak b e_3)=4 .
                                                        \tag{OS.5}\label{p1b:eq:natural-os-5}
\]
The last addition is justified for the full maps into \(SU(3)\):
the degree-five cohomology class is primitive under multiplication,
since the positive cohomology degrees below five are three alone.
In particular the actual two factors reinforce in this common
source convention; their previously displayed opposite degrees
do not cancel.

Define the named power-collar map on the same domain:
\[
\mathfrak p(t,r,g)=z(2\pi t,r,g^2),\qquad
 \deg(\mathfrak p e_3)=4,\qquad
 \boxed{\begin{gathered}
 H_5(0,t,r,g)=\mathfrak b(t,r,g),\\
 H_5(1,t,r,g)=\mathfrak p(t,r,g).
 \end{gathered}}
\tag{OS.6}\label{p1b:eq:natural-os-6}
\]
All faces of \(H_5\) must be identity. The degree of the power
map \(g\mapsto g^2\) on this \(S^3\) is two: multiplication
pulls its degree-three class to the sum of the two factor classes.
The degree statement for \(\mathfrak p\) therefore follows from
the identified actual \(z\), rather than from a benchmark.

There is no homotopy-class obstruction to this specific \(H_5\).
For the retained \(SU(2)\to SU(3)\xrightarrow{p}S^5\) fibration,
Bott's stable range gives \(\pi_5SU(3)=\mathbb Z\).
The exact sequence, its nontrivial characteristic map \(\eta_3\),
and the finite \(\pi_5SU(2)\) give
\[
 p_*:\pi_5SU(3)\hookrightarrow\pi_5S^5=\mathbb Z,
 \quad\operatorname{im}p_*=2\mathbb Z,\qquad
 {}[\mathfrak b]=[\mathfrak p].
                                                        \tag{OS.7}\label{p1b:eq:natural-os-7}
\]
This proves the required full-map class identification and existence
of a based comparison. It does not specify its actual matrix family
or evaluate its complete projected track. Those are not reported
as constructed here.

Comparison choices introduce no concealed independent natural bit.
Two \(H_5\) choices differ by an element of \(\pi_6SU(3)\).
The boundary of \(\eta_5\in\pi_6S^5\) is
\(\eta_3\eta_4\), the nonzero generator of \(\pi_5SU(2)=C_2\);
exactness therefore gives
\[
 p_*\pi_6SU(3)=0.
                                                        \tag{OS.8}\label{p1b:eq:natural-os-8}
\]
After insertion on the actual double-suspended clutch, the closed
difference projects as \(p\alpha\) composed with its suspended
source and is null. The right intrinsic norm and joining factors
fix \(e_3\), so they do not change this column equation.
This is a comparison-choice statement, not the absolute parity of
one missing \(H_5\) track.

\paragraph{Immediate paste into the existing detector input.}
Write \(\mathfrak C_h^{\rm full}\) for \eqref{p1b:eq:natural-os-4} with the retained joining.
Using precisely the old \eqref{p1b:eq:natural-oq-8} annulus, the new full assembly is
\[
 K'(R,\xi)=
 \begin{cases}
 K_{\rm common}(7R,\xi),&R\le1/7,\\
 \mathfrak C_{2-7R}^{\rm full}(\xi),&1/7\le R\le2/7,\\
 K'_{\rm OQ}((7R-1)/6,\xi),&2/7\le R\le1,
 \end{cases}
 \quad
 K_{\rm common}|_{S^7}=
       \mathfrak b(t(s),r,g(v))\mathcal T_{\rm int}\mathcal R_J .
                                                        \tag{OS.9}\label{p1b:eq:natural-os-9}
\]
Every whole matrix and projected column for \(R\ge1/7\) is constructed.
The generally nonconstant inner column is not collapsed.
The next proof constructs the actual clutch-square null and its
power-collar filling, a named terminal component of \eqref{p1b:eq:natural-os-6}.
It does not fill \(K_{\rm common}\) without the intervening comparison
and the retained intrinsic norm and original joining remainder.

The complete original comparison is evaluated in
\eqref{p1b:eq:natural-aj-4}--\eqref{p1b:eq:natural-aj-7}:
\(h_t=w_t=1\), \(r_J=\zeta_{\rm nat}=0\), with every retained
extension test decided. This component is used with its full frames,
lower divisions and joining; its isolated value is not a replacement
for that complete comparison. Supplementary coordinate checks are in the archived evidence file \nolinkurl{Paper1B-NaturalDeterminantCommonSourceDiagnostics-Appendix-v1_0.tex}.


\subsection{A filling from the actual clutch square}
\label{p1b:sec:natural-clutch-square-filling}

Start with the precise \(\mathfrak p\) in \eqref{p1b:eq:natural-os-6}.
The dependencies executed here are to base the actual \eqref{p1b:eq:natural-nz-17}--\eqref{p1b:eq:natural-nz-19}
clutch comparison with its moving base value, construct a literal
inverse path for its actual canonical \(q_2\), retain its constant
base-frame conjugation, and apply the resulting square null to
the named power-collar map. No characteristic degree is used to
assign a natural coefficient.

Let \(v_*=(1,0,0,0,0,0)\) be the original direction with \(g(v_*)=I\).
Use the entire previously constructed clutch path \(g_h^{NZ}\), and put
\[
 G_h(v)=g_h^{NZ}(v)g_h^{NZ}(v_*)^*,\qquad
 G_0=g,\quad
 G_1=g_{\rm can}(v)g_{\rm can}(v_*)^*=:f(v).
                                                        \tag{OT.1}\label{p1b:eq:natural-ot-1}
\]
This retains the actual moving base value at every time, including
the connection transport and both actual frame factors.
All values are in the original intrinsic rank-two axes, and
\(G_h(v_*)=I\).

Keep the \eqref{p1b:eq:natural-ny-1}, \eqref{p1b:eq:natural-ny-6}, \eqref{p1b:eq:natural-nz-16} coordinates:
\[
\begin{gathered}
 n(v)=(v_{x3},-v_{p1},-v_{p3},v_{p2},v_{x1},v_{x2}),\qquad q_2(n)={2a-\bar uiu\over1+a^2},\qquad q_0=q_2(n(v_*)), \\
 f(v)=L_q q_2(n(v))^{-1}q_0L_q^{-1}.
\end{gathered}
 \tag{OT.2}\label{p1b:eq:natural-ot-2}
\]
The constant \(R_q\) cancels only by the full based division in
\eqref{p1b:eq:natural-ot-1}. The \(L_q\) and actual \(q_0\) remain. In these exact
coordinates \(q_0=-i\); the construction below also retains it
as its actual computed value.

Rotate the quaternionic target axis, without altering the actual
source or choosing a Hopf replacement:
\[
 q_{2,h}(n)={2a-\bar u(i\cos\pi h+j\sin\pi h)u\over1+a^2},
 \quad q_{2,0}=q_2,\quad q_{2,1}=q_2^{-1}.
                                                        \tag{OT.3}\label{p1b:eq:natural-ot-3}
\]
Indeed \(|u|^2=1-a^2\). Its numerator has squared norm
\(4a^2+(1-a^2)^2=(1+a^2)^2\).
The imaginary numerator changes sign at \(h=1\), proving the
literal inverse equation. The homogeneous vanishing coefficients
of \eqref{p1b:eq:natural-ny-1} prove continuity at its missing Hopf axes.

For the first half of a based inverse path use
\[
 f_h^{(1)}(v)=L_q q_{2,h}(n(v))^{-1}
                         q_{2,h}(n(v_*))L_q^{-1}.
                                                        \tag{OT.4}\label{p1b:eq:natural-ot-4}
\]
It starts at \(f\) and ends at
\(L_q q_2(n(v))q_0^{-1}L_q^{-1}\).
That end is not yet \(f^{-1}\).
Let \(\gamma_k\) be the actual quaternionic logarithm path
from \(I\) to \(q_0^{-1}\), using the same fixed-axis convention
as \eqref{p1b:eq:natural-nz-19}. The second half is the full constant-frame comparison
\[
 f_k^{(2)}(v)=L_q\gamma_k q_2(n(v))q_0^{-1}
                           \gamma_k^{-1}L_q^{-1},
 \quad f_0^{(2)}=f_1^{(1)},\quad f_1^{(2)}=f^{-1}.
                                                        \tag{OT.5}\label{p1b:eq:natural-ot-5}
\]
Both families equal identity at \(v_*\).
Concatenate them in equal intervals and call the path \(F_h\).
Its endpoints are the actual \(f,f^{-1}\), with exact whole-matrix seam.

Now the actual clutch-square null is
\[
 S_h(v)=
 \begin{cases}
 G_{2h}(v)^2,&0\le h\le1/2,\\
 f(v)F_{2h-1}(v),&1/2\le h\le1,
 \end{cases}
 \quad
 \boxed{S_0=g^2,\quad S_1=I,\quad S_h(v_*)=I.}
                                                        \tag{OT.6}\label{p1b:eq:natural-ot-6}
\]
This is an explicit intrinsic null of the original clutch square,
with its original coordinates and both actual frame comparisons.
The seam is exactly \(f^2\). It has not been deduced solely from
the order of \([g]\).

Apply it immediately to the power collar, on the original cap:
\[
 P_h(s,z)=z(2\pi t(s),r,S_h(v)),\qquad
 P_h(s,0)=I,\quad
 P_h|_{r=1}=I,\qquad
 K_{\mathfrak p}(R,s,z)=P_{1-R}(s,z).
                                                        \tag{OT.7}\label{p1b:eq:natural-ot-7}
\]
The source gives a closed \(S^7\) at \(R=1\).
Both radial faces and both circle faces are identity for every
clutch value, and \(z(\theta,\lambda,I)=I\).
Uniform continuity of this universal map proves the cap-center
extension and the \(R=0\) cone extension. Thus \eqref{p1b:eq:natural-ot-7} is an actual
\(D^8\to SU(3)\) filling of the named \(\mathfrak p\) input.
Its complete last column is \(P_{1-R}e_3\).

After a constructed \(H_5\), its immediate full-input application
would be
\[
 H_5(h,t(s),r,g(v))\mathcal T_{\rm int}\mathcal R_J.
\]
The terminal power-collar comparison is already implemented as
\[
 \Phi_h^{\rm power}=P_h\mathcal T_{\rm int}\mathcal R_J,\qquad
 \Phi_0^{\rm power}=\mathfrak p(g)\mathcal T_{\rm int}\mathcal R_J,
 \quad
 \Phi_1^{\rm power}=\mathcal T_{\rm int}\mathcal R_J.
                                                        \tag{OT.8}\label{p1b:eq:natural-ot-8}
\]
All original joining values remain; the cap boundary makes the paste
exact. No value of a missing \(H_5\) is supplied by a fallback.
Consequently \eqref{p1b:eq:natural-ot-8} constructs the specified terminal comparison
component, not an already applied comparison from the full original
input without \eqref{p1b:eq:natural-os-6}.

\paragraph{The exact assembly and original target.}
The auxiliary intervening family \eqref{p1b:eq:natural-os-6}, if used, must have its
full stated boundary. Independently, the intrinsic boundary is the
specified \(\mathcal T_{\rm int}\mathcal R_J\), with its actual
inverse bridge norm and original geometric \eqref{p1b:eq:natural-nd-8}. Its full filling,
together with the complete columns of \eqref{p1b:eq:natural-os-9}, \(H_5\) and \eqref{p1b:eq:natural-ot-8},
must enter the existing detector:
\[
 K'|_{S^7}=\operatorname{diag}(\psi',1),\qquad
 \operatorname{Sq}^4u_5=\zeta_{\rm nat} v_9
       \ \text{on }C_{[K'e_3]},\qquad
 h_t=1+\zeta_{\rm nat},\quad w_t=1 .
                                                        \tag{OT.9}\label{p1b:eq:natural-ot-9}
\]
The complete original comparison is evaluated in
\eqref{p1b:eq:natural-aj-4}--\eqref{p1b:eq:natural-aj-7}:
\(h_t=w_t=1\), \(r_J=\zeta_{\rm nat}=0\), with every retained
extension test decided. This component is used with its full frames,
lower divisions and joining; its isolated value is not a replacement
for that complete comparison.


\subsection{The actual pair-to-power comparison}
\label{p1b:sec:natural-pair-to-power}

Retain the original two-bit target and the complete \eqref{p1b:eq:natural-os-9} input.
The dependencies executed here are: identify the actual characteristic
factor as a commutator of its actual based-direction disk; compare that
disk relatively with the actual centralizer disk; separate the two
copies of the clutch sphere in a specified based homotopy; construct
the full \eqref{p1b:eq:natural-os-6} matrix family and apply it immediately. The last-column
track is retained at every step. No degree calibration assigns a bit.

Put \(A(t,r)=a_*(t,r)\), \(Z(t,r)=Z_r(2\pi t)\), and
\(z=Z^*gZg^*\), with all matrices and fibre axes precisely as in
\eqref{p1b:eq:natural-os-1}, \eqref{p1b:eq:natural-or-1}, \eqref{p1b:eq:natural-nx-2} and \eqref{p1b:eq:natural-oo-1}. Quaternionic \(J\) commutes with
the original \(M\) and every \(g\in SU(2)\). Thus
\(J(Mge_1)=Mge_2\). The entire geodesic lift is equivariant under
this same fibre action, including its south normalization.
Substitution into the full \eqref{p1b:eq:natural-or-3}, not just its last column, proves
\[
 a(\pi t,b(r,g))=gA(t,r)g^*,\qquad
 \widetilde a=gAg^*A^*,\qquad
 \boxed{\mathfrak b=Z^*gZAg^*A^*.}                 \tag{OU.1}\label{p1b:eq:natural-ou-1}
\]
The exact disk boundaries are
\[
\begin{gathered}
 A(t,0)=Z(t,0)=d(2\pi t), \\
 A(t,1)=Z(t,1)=I, \\
 A(0,r)=A(1,r)=Z(0,r)=Z(1,r)=I.
\end{gathered}
 \tag{OU.2}\label{p1b:eq:natural-ou-2}
\]
In particular \(D=AZ^*\) is identity on every boundary edge.
This two-dimensional comparison retains the actual disk; it does
not replace the five-sphere map by one with the same degree.

\paragraph{A specified relative comparison of these two disks.}
Write \(T(v)=T(e_3,v)\) for the original canonical special-unitary
transporter, \(v=De_3\), and \(Q=T(v)^*D\).
Thus \(Q=\operatorname{diag}(Q_2,1)\), \(Q_2\in SU(2)\).
In the coordinates
\[
\begin{gathered}
 Q_2=\begin{pmatrix}q_0+iq_1&q_2+iq_3\\
                    -q_2+iq_3&q_0-iq_1\end{pmatrix},
 \qquad |q|=1,\quad p=(0,0,0,1),\\
 1+\operatorname{Re}D_{33}>\frac1{100},\qquad
 1+q_3>\frac1{100}.
\end{gathered}
\tag{OU.3}\label{p1b:eq:natural-ou-3}
\]
Both strict inequalities are proved on the complete actual disk.
Here are the exact input and enclosure argument, so no sampled
gap is a premise. Set \(r=(1-\cos\rho)/2\), \(0\le\rho\le\pi\),
\(y=\cos\rho\), \(\theta=\pi t\). In the actual cap axes the
based geodesic matrix has its only off-diagonal entries in axes
one and three:
\[
\begin{gathered}
\begin{aligned}
 L_{11}&=e^{-i(1-y)\theta/2}(\cos\theta-iy\sin\theta),&
 L_{22}&=e^{-i(1+y)\theta/2},\\
 L_{13}&=\sin\theta\sin\rho\,e^{-i\theta/2},&
 L_{31}&=-\sin\theta\sin\rho\,e^{i(1+y/2)\theta},\\
 L_{33}&=e^{i\theta}(\cos\theta+iy\sin\theta).
 \end{aligned}\\
A=(L^*)^2V_\alpha\Psi V_\alpha^*\Psi,\quad
 \alpha=\frac\pi2(1-t),\\
 \Psi=\operatorname{diag}
     (e^{i\pi(1-r)},e^{-i\pi(1-r)},1).
\end{gathered}
\tag{OU.4}\label{p1b:eq:natural-ou-4}
\]
Here \(V_\alpha\) is the real determinant-one two-plane rotation.
These are pointwise simplifications of \eqref{p1b:eq:natural-nx-2} and \eqref{p1b:eq:natural-or-1}.
For either of the two original circle disks the quaternion is
\((c^2\cos s+h^2,c^2\sin s,ch(1-\cos s),-ch\sin s)\),
where \(c=\cos(\pi r/2)\), \(h=\sin(\pi r/2)\).
Their product in axes \(12,23\), at angles \(2\pi t,4\pi t\),
is the exact \(Z\).
The special-unitary cofactor identity gives
\[
 (Q_2)_{12}=
       \frac{D_{12}-\overline{D_{21}}}{|1+D_{33}|}.
                                                        \tag{OU.5}\label{p1b:eq:natural-ou-5}
\]
The supplied certificate evaluates exactly these expressions with
outward integer intervals at denominator \(2^{80}\).
Its enclosure of \(\pi\) follows from the Machin identity and the
alternating rational arctangent bounds with 72 terms.
After quadrant reduction, Taylor polynomials of degrees 24 and
25 enclose cosine and sine: the reduced argument has absolute
value at most one, and each omitted remainder is less than
\(2^{-80}\). Wider reduced intervals use the rigorous range
\([-1,1]\). Products, quotients and square roots are rounded
outwards by integer division and integer square root.
The recorded dyadic rectangles cover the whole \((t,\rho/\pi)\)
square, including all edges; four sibling rectangles are replaced
by their parent until the entire unit square remains.
On every leaf both lower bounds strictly exceed \(1/100\).
This proves \eqref{p1b:eq:natural-ou-3} by finite exact inequalities. The certificate
is an essential chart-validity input for these actual disks,
not a natural-bit detector. Supplementary coordinate replays are
in the archived evidence file \nolinkurl{Paper1B-NaturalDeterminantPairToPowerDiagnostics-Appendix-v1_0.tex}.

Let \(v_k=((1-k)v+ke_3)/|(1-k)v+ke_3|\).
Let \(F_k(q)=((1-k)q+kp)/|(1-k)q+kp|\), interpreted in the
displayed quaternion coordinates. The two inequalities exclude
respectively \(v=-e_3\) and \(q=-p\).
Also \(p\ne-(1,0,0,0)\). Consequently every denominator below
is nonzero on its complete domain:
\[
\begin{gathered}
D_h=
 \begin{cases}
 T(v_{2h})Q,&0\le h\le1/2,\\
 \operatorname{diag}\bigl(F_{2h-1}(q)
                         F_{2h-1}(1)^*,1\bigr),&1/2\le h\le1,
 \end{cases}\\
D_0=D,\quad D_1=I,\quad
 \boxed{A_h=D_hZ:\ A\longrightarrow Z.}
\end{gathered}
\tag{OU.6}\label{p1b:eq:natural-ou-6}
\]
The seam is exactly \(Q\). On each disk edge \(D=Q=I\);
the moving base division \(F_k(1)^*\) makes the second leg
identity there as well. Thus \eqref{p1b:eq:natural-ou-6} is a constructed relative
homotopy of full matrices, with its complete column.

\paragraph{Construct the clutch diagonal comparison.}
Use the original sphere coordinates
\(g=(q_0+iq_1,q_2+iq_3)\) in the displayed matrix convention.
For \(g\ne I\) set \(x=(q_1,q_2,q_3)/(1-q_0)\).
For \(\epsilon=\pm1\), put \(x_h^\epsilon=x-2\epsilon h e_1\) and
\[
 c_h^\epsilon(g)=
 \begin{cases}
 \operatorname{st}^{-1}\!
 \left(\dfrac{x_h^\epsilon}{1-h|x_h^\epsilon|^2}\right),
                 &h|x_h^\epsilon|^2<1,\\
 I,&h|x_h^\epsilon|^2\ge1,
 \end{cases}
 \quad c_h^\epsilon(I)=I,\quad c_0^\epsilon(g)=g .
                                                        \tag{OU.7}\label{p1b:eq:natural-ou-7}
\]
Here the inverse stereographic map is
\(\operatorname{st}^{-1}(x)=((|x|^2-1),2x)/(1+|x|^2)\).
As the denominator tends to zero the value tends to \(I\)
uniformly; near the absent stereographic point \(|x|\) tends
to infinity uniformly, while the translations remain bounded.
The norm of the noncollapsed argument is at least
\(|x_h^\epsilon|\). These bounds prove joint continuity also
at \(h=0,g=I\). At \(h=1\) the two nonidentity supports are
the disjoint open unit three-balls with centers \(\pm2e_1\).
Hence the following is an actual based nullhomotopy, with all
circle and radial faces fixed:
\[
 B_h=c_h^+(g)\,z(t,r,c_h^-(g))\,
                  c_h^+(g)^*z(t,r,c_h^-(g))^*,
 \quad B_0=gzg^*z^*,\quad B_1=I.                   \tag{OU.8}\label{p1b:eq:natural-ou-8}
\]
This constructs the comparison of the actual diagonal source.
It does not assume that the clutch itself is null.

The full missing \eqref{p1b:eq:natural-os-6} family is now
\[
\begin{gathered}
\boxed{
 H_5(h,t,r,g)=
 \begin{cases}
 Z^*gZA_{3h}g^*A_{3h}^*,&0\le h\le1/3,\\
 z\,Z(t,r+(3h-1)(1-r))\,z\,
             Z(t,r+(3h-1)(1-r))^*,&1/3\le h\le2/3,\\
 zB_{3-3h}z,&2/3\le h\le1 .
 \end{cases}}\\
 H_5(0)=\mathfrak b,\quad H_5(1)=z(t,r,g^2)=\mathfrak p.
\end{gathered}
\tag{OU.9}\label{p1b:eq:natural-ou-9}
\]
At the first seam the exact value is \(zZzZ^*\);
at the second it is \(z^2\).
At the terminal endpoint the commutator identity gives
\(z(g)\,g z(g)g^*=z(g^2)\).
Every face is identity, including \(g=I\).
Thus the actual matrix family, rather than only the class
identification \eqref{p1b:eq:natural-os-7}, has been constructed. Its last column is
the displayed full product applied to \(e_3\); it is not
replaced by a constant or assigned a parity.




\subsection{Apply the constructed comparison and retain the exact terminal input}
\label{p1b:sec:natural-pair-to-power-assembly}

Apply \eqref{p1b:eq:natural-ou-9} immediately on the original cap, in the actual
\(t(s),r,g(v)\) coordinates of \eqref{p1b:eq:natural-os-1}--\eqref{p1b:eq:natural-os-4}:
\[
 \Phi_h^{H}=H_5(h,t(s),r,g(v))
                         \mathcal T_{\rm int}\mathcal R_J,
 \quad
 \Phi_0^{H}=\mathfrak b(g)\mathcal T_{\rm int}\mathcal R_J,\quad
 \Phi_1^{H}=\mathfrak p(g)\mathcal T_{\rm int}\mathcal R_J .
                                                        \tag{OV.1}\label{p1b:eq:natural-ov-1}
\]
Off the cap retain the original \(\mathcal R_J\). The cap faces
of \(H_5\) are identity, so this pastes exactly, including the
cap center where the absent direction has a uniform extension.
The complete matrices and their columns are retained.
There is no substitution of the original \eqref{p1b:eq:natural-nd-8} interior.

Apply the actual \eqref{p1b:eq:natural-ot-6}--\eqref{p1b:eq:natural-ot-8} clutch-square filling next:
\[
 \Phi_h^{P}=P_h\mathcal T_{\rm int}\mathcal R_J,\qquad
 \Phi_0^P=\Phi_1^H,\quad
 \Phi_1^P=\mathcal T_{\rm int}\mathcal R_J .
                                                        \tag{OV.2}\label{p1b:eq:natural-ov-2}
\]
These are now applied consecutive comparisons from the actual
\eqref{p1b:eq:natural-os-9} input, with a constructed intervening \(H_5\).
The earlier conditional adapter is no longer a missing
component. No independent calibration or degree assignment
is used in this application.

With \(K'_{\rm OS}\) the exact previously assembled outer
annuli in \eqref{p1b:eq:natural-os-9}, the resulting entire assembly is
\[
 \boxed{
 K'(R,\xi)=
 \begin{cases}
 K_{\rm int}(9R,\xi),&0\le R\le1/9,\\
 \Phi^P_{2-9R}(\xi),&1/9\le R\le2/9,\\
 \Phi^H_{3-9R}(\xi),&2/9\le R\le3/9,\\
 K'_{\rm OS}((9R-2)/7,\xi),&3/9\le R\le1,
 \end{cases}\qquad
 K_{\rm int}|_{S^7}=\mathcal T_{\rm int}\mathcal R_J .
 }                                                      \tag{OV.3}\label{p1b:eq:natural-ov-3}
\]
At \(R=1/9\) the exact boundary is
\(\mathcal T_{\rm int}\mathcal R_J\); at \(2/9\) it is
\(\mathfrak p(g)\mathcal T_{\rm int}\mathcal R_J\); at
\(3/9\) it is \(\mathfrak b(g)\mathcal T_{\rm int}\mathcal R_J\).
At \(R=1\) it is the original
\(\operatorname{diag}(\psi',1)\).
Every whole matrix and projected column on the displayed
outer annular region is constructed. The radial fractions
are coordinates, not a completion percentage or a batch bound.

Both end columns of the constructed annulus are exactly \(e_3\).
Thus its complete projection now supplies a specific closed input
to the retained detector, without an unknown filling:
\[
\begin{gathered}
 \gamma_{\rm ann}(u,\xi)=K'((1+8u)/9,\xi)e_3,\qquad 0\le u\le1, \\
 \gamma_{\rm ann}(0,\xi)=\gamma_{\rm ann}(1,\xi)=e_3,\qquad \gamma_{\rm ann}:S^8\longrightarrow S^5 .
\end{gathered}
 \tag{OV.4}\label{p1b:eq:natural-ov-4}
\]
This is a construction of the actual annular input from all the
identified comparisons, not an evaluation of its square.
Every column-moving factor is supported on the original two-cap
ball. On its entire complement the retained matrices are intrinsic,
so the exact projected value is \(e_3\). This supplies the complete
column evaluator without requiring a numerical replacement of the
original geometric joining. The full matrix joining is unchanged.
Its connecting image is the difference between the original
intrinsic boundary and the specified terminal intrinsic boundary.
The full annular and terminal classes are evaluated in
\eqref{p1b:eq:natural-pf-3}--\eqref{p1b:eq:natural-pf-9}, \eqref{p1b:eq:natural-pg-14} and \eqref{p1b:eq:natural-aj-4}--\eqref{p1b:eq:natural-aj-5}. This annular
construction alone supplies neither complete coefficient.

The terminal boundary fixes \(e_3\); that property alone does
not evaluate its class or make
\(K_{\rm int}\) an intrinsic filling.
On the second cut its exact cap factor is
\[
\begin{gathered}
\mathcal T_{\rm int}(\pi(t+1),z)=
 \begin{cases}
 \mathcal J_{1-2t}(U_g(z)),&0\le t\le1/2,\\
 C_{1-t}(z)^*(D C_{1-t}(\sigma z)D)^*,&1/2\le t\le1,
 \end{cases}\\
\mathcal T_{\rm int}|_{s\le\pi}=I,\quad
 D=\operatorname{diag}(1,-1,1).
\end{gathered}
\tag{OV.5}\label{p1b:eq:natural-ov-5}
\]
Here the original intrinsic bridge half of \(C\) and the
original \eqref{p1b:eq:natural-om-6}--\eqref{p1b:eq:natural-om-7} joining remainder are retained.
Neither top coefficient is deleted by their boundary equations.

The retained complete natural filling criterion is
\[
\boxed{\begin{gathered}
 K_{\rm int}:D^8\longrightarrow SU(3),\\
 K_{\rm int}|_{S^7}=\mathcal T_{\rm int}\mathcal R_J,\\
 \operatorname{Sq}^4u_5=\zeta_{\rm nat} v_9
       \text{ on }C_{[K'e_3]},\\
 h_t=1+\zeta_{\rm nat},\quad w_t=1 .
\end{gathered}}
\tag{OV.6}\label{p1b:eq:natural-ov-6}
\]
The complete intrinsic comparison through \eqref{p1b:eq:natural-aj-4}--\eqref{p1b:eq:natural-aj-5} evaluates
this same retained terminal boundary. The complete projected
parity in \eqref{p1b:eq:natural-ov-6}, including the new tracks, is not supplied
by the two component degrees or by the chart certificate.
The complete original comparison is evaluated in
\eqref{p1b:eq:natural-aj-4}--\eqref{p1b:eq:natural-aj-7}:
\(h_t=w_t=1\), \(r_J=\zeta_{\rm nat}=0\), with every retained
extension test decided. This component is used with its full frames,
lower divisions and joining; its isolated value is not a replacement
for that complete comparison.


\subsection{Compare the actual intrinsic bridge norm and evaluate its mirror correction}
\label{p1b:sec:natural-intrinsic-norm-mirror}

Retain the original \(B_\rho\) of \eqref{p1b:eq:natural-nz-20}, from \(q^2=H^6\) to
\(P=\psi(U_g)\), the original spatial cap, \(D=\operatorname{diag}(1,-1,1)\),
and \(\sigma(z_x,z_p)=(z_p,z_x)\). Write
\(B_\rho^\sigma=DB_\rho(\sigma z)D\). Thus
\[
 B_0^\sigma=q^{-2},\qquad B_1^\sigma=P,\qquad
 N_B(\rho,z)=B_\rho^{-1}(B_\rho^\sigma)^{-1},
 \quad N_B(0,z)=I,\quad N_B(1,z)=P^{-2}.       \tag{OW.1}\label{p1b:eq:natural-ow-1}
\]
Every spatial boundary value is identity. The bridge half of the
actual \eqref{p1b:eq:natural-ov-5} terminal is exactly \(N_B(2-2t,z)\).
No root or original joining is changed in this notation.

Construct the following specific alternative bridge and its closed
difference with the original bridge:
\[
\begin{gathered}
B^\#_\rho=
 \begin{cases}
 (B^\sigma_{2\rho})^{-1},&0\le\rho\le1/2,\\
 P\mathcal J_{2\rho-1}(U_g),&1/2\le\rho\le1,
 \end{cases}\\
B^\#_0=q^2,\quad B^\#_{1/2}=P^{-1},\quad B^\#_1=P,\quad
 L_\rho=B^\#_\rho B_\rho^{-1}.
\end{gathered}
\tag{OW.2}\label{p1b:eq:natural-ow-2}
\]
The seam follows from \(\mathcal J_0=P^{-2}\).
In particular \(L_0=L_1=I\) and \(L|_{\partial B}=I\);
it is a specified intrinsic \(S^7\) map, with no coefficient assigned.

\paragraph{A relative intrinsic null of the actual norm difference.}
Set \(C_\rho=B_\rho^\sigma B_\rho\). Direct multiplication gives
\[
 \boxed{
 N_{B^\#}N_B^{-1}
   =(B_\rho^{-1}L_\rho^{-1}B_\rho)
      (C_\rho^{-1}(L_\rho^\sigma)^{-1}C_\rho).
 }                                                     \tag{OW.3}\label{p1b:eq:natural-ow-3}
\]
Both factors are identity at \(\rho=0,1\).
First replace both moving conjugators \(B_\rho,C_\rho\) by
\(B_{(1-h)\rho},C_{(1-h)\rho}\), holding the entire \(L\) fixed.
At the end \(B_0=q^2\) and \(C_0=I\).
This operation is relative to every collapsed face, even though
\(q^2\) is not declared null.

Next remove the remaining \(q^2\) conjugation using the actual
spatial source \(S^6=B/\partial B\). In the south-based
stereographic coordinate \(x\in\mathbb R^6\), use
\[
 x_h^\pm=x\mp2h e_1,\qquad
 c_h^\pm(x)=
 \begin{cases}
 \operatorname{st}^{-1}\bigl(x_h^\pm/(1-h|x_h^\pm|^2)\bigr),
       &h|x_h^\pm|^2<1,\\
 \mathrm{south},&h|x_h^\pm|^2\ge1.
 \end{cases}                                           \tag{OW.4}\label{p1b:eq:natural-ow-4}
\]
These start at the identity, fix the based point, and end with
disjoint open unit six-ball supports centered at \(\pm2e_1\).
The homogeneous inverse stereographic formula gives continuity
at the collapse seam. Near the based point its norm tends to
infinity uniformly, as in \eqref{p1b:eq:natural-ou-7}.
Precompose \(q^2\) by \(c_h^+\) and \(L_\rho\) by \(c_h^-\)
in the first factor of \eqref{p1b:eq:natural-ow-3}. At \(h=1\) their supports are
disjoint, so that factor is literally \(L_\rho(c_1^-)^{-1}\).
Reverse the compression of \(L\). The complete difference has
now been compared with \(L_\rho^{-1}(L_\rho^\sigma)^{-1}\).

Remove the constant \(D\) conjugation by
\(\operatorname{diag}(e^{i\pi(1-h)/2},e^{-i\pi(1-h)/2},1)\),
and apply the actual swap-to-reflection family \eqref{p1b:eq:natural-op-4}.
For fixed \((\rho,a_1,a_2,b)\), the remaining word is
\(L_\rho(u)^{-1}L_\rho(1-u)^{-1}\), with identity at both
ends of the spatial interval. The two exact pinch and
backtracking stages \eqref{p1b:eq:natural-op-5}, applied to \(L^{-1}\), contract it.
All stages are intrinsic; all spatial and \(\rho\) faces are fixed.
Use intervals \([0,.2],[.2,.4],[.4,.6],[.6,.8],[.8,.9],[.9,1]\)
for the six displayed operations. Denote this specified null by
\(\mathcal D_h\), with
\(\mathcal D_0=N_{B^\#}N_B^{-1}\), \(\mathcal D_1=I\).
This is a constructed null of the difference, not a filling of \(N_B\).

Apply it immediately to the complete actual terminal:
\[
\begin{gathered}
\mathcal T_h^\#(\pi(t+1),z)=
 \begin{cases}
 \mathcal J_{1-2t}(U_g(z)),&0\le t\le1/2,\\
 \mathcal D_{1-h}(2-2t,z)N_B(2-2t,z),&1/2\le t\le1,
 \end{cases}\\
\mathcal T_h^\#|_{s\le\pi}=I,\qquad
 \Phi_h^\#=\mathcal T_h^\#\mathcal R_J.
\end{gathered}
\tag{OW.5}\label{p1b:eq:natural-ow-5}
\]
At both bridge-time faces \(\mathcal D=I\), so the midpoint
and circle seams are the original literal seams.
The cap boundary is identity and the original \eqref{p1b:eq:natural-nd-8} joining
remains in \(\mathcal R_J\) unchanged. Thus \eqref{p1b:eq:natural-ow-5} is an
intrinsic comparison on the original \(S^7\), including its
bundle identifications and actual lower values.
It starts at \eqref{p1b:eq:natural-ov-5} and ends at the specified mirrored bridge norm.
Its projected column is identically \(e_3\).

\paragraph{Evaluate the actual square-null mirror correction.}
The characteristic exchange is the fixed real map
\(\widehat\sigma(y,w_0,w_1)=(y,-w_0,w_1)\) in \eqref{p1b:eq:natural-nw-5} axes.
It satisfies
\[
 D\psi(\widehat\sigma b)D=\psi(b),\qquad
 B_h\widehat\sigma=\widehat\sigma B_{-h}.               \tag{OW.6}\label{p1b:eq:natural-ow-6}
\]
Here \(B_h\) in the second equation is the spinor rotation in \eqref{p1b:eq:natural-or-1},
not the intrinsic bridge \(B_\rho\). Its anticommutation follows
directly from \(Jw=(-\bar w_1,\bar w_0)\).
Put \(J_h(b)=\mathcal J_h(b)\) and \(W_h(b)=\psi(b)J_h(b)\).
The entire three-leg correction is
\[
 \mathcal U(\lambda,b)=
 \begin{cases}
 J_{1-3\lambda}(b),&0\le\lambda\le1/3,\\
 W_{2-3\lambda}(b)^{-1}
       (D W_{2-3\lambda}(\widehat\sigma b)D)^{-1},
                                      &1/3\le\lambda\le2/3,\\
 J_{3\lambda-2}(b)^{-1},&2/3\le\lambda\le1.
 \end{cases}                                           \tag{OW.7}\label{p1b:eq:natural-ow-7}
\]
Its successive endpoints are \(I,P^{-2},P^2,I\).
Both poles of \(b\) give identity on every leg.
Thus this is an actual universal \(S^5\to SU(2)\) map.
The exchange difference \(J_h(b)(D J_h(\widehat\sigma b)D)^{-1}\)
is another actual \(S^5\) map with the same based faces.

The required substitution is the original \(U_g\), not a benchmark.
By \eqref{p1b:eq:natural-nw-9}, with \(\kappa(g)=Mg e_1\), it is the literal suspension
of \(\kappa\circ g\) in the radial coordinate \(y=1-2r\).
The longitude map \(\kappa:SU(2)\to S^3\) is a degree-minus-one
homeomorphism; the established original clutch has class
\(\eta_3\eta_4\). Consequently the actual suspended substitution is
\[
 \Sigma U_g=\Sigma^2(\kappa\circ g),\qquad
 {}[\Sigma U_g]=\eta_5\eta_6\in\pi_7S^5.                 \tag{OW.8}\label{p1b:eq:natural-ow-8}
\]
All meridians, the south base, the actual \(M\) and the original
\(g\) have been retained in this identification.
No primitive value of \(\mathcal U\) needs to be selected:
the entire group \(\pi_5S^3=C_2\{\eta_3\eta_4\}\) is annihilated
by this actual precomposition, since the retained \eqref{p1b:eq:rt-23-3} relation gives
\[
 \boxed{
 {}[\mathcal U(\lambda,U_g(z))]
   =0,\qquad
 {}[J_h(U_g)(D J_h(U_g(\sigma z))D)^{-1}]=0,\qquad
 \eta_3\eta_4\eta_5\eta_6=6\nu'\eta_6=6\mu=0.
 }                                                     \tag{OW.9}\label{p1b:eq:natural-ow-9}
\]
These are evaluations of specified comparison maps with their
actual source identifications. They are not assignments of the
remaining natural bit from a primitive value.

Cancelling the common final \(N_B^{-1}\) path in the mirrored
terminal comparison leaves exactly \eqref{p1b:eq:natural-ow-7}.
Hence \([\mathcal T^\#_{\rm int}]+[\mathcal T_{\rm int}]=0\).
Together with the constructed \eqref{p1b:eq:natural-ow-5} this is compatible with
the exponent-two top group. It does not establish
\([\mathcal T_{\rm int}]=0\).
The complete original comparison is evaluated in
\eqref{p1b:eq:natural-aj-4}--\eqref{p1b:eq:natural-aj-7}:
\(h_t=w_t=1\), \(r_J=\zeta_{\rm nat}=0\), with every retained
extension test decided. This component is used with its full frames,
lower divisions and joining; its isolated value is not a replacement
for that complete comparison.


\subsection{The complete actual annular attaching map on Cartesian \(S^8\)}
\label{p1b:sec:natural-closed-annular-sphere}

Use the already constructed \eqref{p1b:eq:natural-ov-3}--\eqref{p1b:eq:natural-ov-4} comparisons.
Their complete columns are constant on the complement of the actual
two-cap log ball; the retained joining there is intrinsic.
This permits a standard-domain input without a numerical substitute
for the original \eqref{p1b:eq:natural-nd-8} full matrix.

The actual source embedding can be written independently of inverse
coordinate routines. For \(z=(z_x,z_p)\in B\), put
\[
 x=\exp z_x,\quad p=\exp z_p,\quad
 d=\tfrac12\log(p^{-1}x),\quad
 y=-p\exp d,\quad n=\frac{d}{3\pi}=\nu u.
                                                        \tag{OX.1}\label{p1b:eq:natural-ox-1}
\]
The logarithm is the original short quaternionic logarithm.
It is unique throughout this cap since
\(\operatorname{dist}(p^{-1}x,1)\le\sqrt2r_B<\pi\).
Then \(x=(-y)e^{3\pi\nu u}\), \(p=(-y)e^{-3\pi\nu u}\)
are exactly the retained boundary coordinates.
At \(d=0\) the Cartesian \(n\) is zero and the absent \(u\)
is immaterial. The image lies in the original inner normal cell by
\eqref{p1b:eq:natural-oj-5}, misses the parameter base, and retains the actual ordinary
bundle frame. At the center,
\[
 d a=(d z_x+d z_p)/2,\qquad
 d n=(d z_x-d z_p)/(6\pi),\qquad y=-\exp a,
                                                        \tag{OX.2}\label{p1b:eq:natural-ox-2}
\]
so this is an embedding with orientation-minus-one in the
displayed \((a,n)\) coordinates. Only this explicit source
identification, not a chosen map of the same degree, is used.

Write \(f_R(s,z)\) for the complete column on this cap.
The recursive pasting \eqref{p1b:eq:natural-ov-3}, \eqref{p1b:eq:natural-os-9}, \eqref{p1b:eq:natural-oq-8} unfolds into
eight equal annular legs:
\[
\begin{array}{c|c}
 R\text{ interval}&\text{actual full cap matrix whose column is }f_R\\ \hline
 {}[1/9,2/9]&P_{2-9R}\mathcal T_{\rm int}\\
 {}[2/9,3/9]&H_5(3-9R,t(s),r,g)\mathcal T_{\rm int}\\
 {}[3/9,4/9]&\text{\eqref{p1b:eq:natural-os-4} at }4-9R\\
 {}[4/9,5/9]&\text{\eqref{p1b:eq:natural-oq-5} cap at }5-9R\\
 {}[5/9,6/9]&\text{\eqref{p1b:eq:natural-op-6} cap at }6-9R\\
 {}[6/9,7/9]&\text{\eqref{p1b:eq:natural-oo-8} cap comparison at }7-9R\\
 {}[7/9,8/9]&\text{\eqref{p1b:eq:natural-on-2} capped comparison at }8-9R\\
 {}[8/9,1]&\text{the original intrinsic input comparison}
\end{array}                                             \tag{OX.3}\label{p1b:eq:natural-ox-3}
\]
The last leg's complete column is literally \(e_3\).
The evaluator retains every factor of the other seven legs,
including the actual moving root-null inverse, original bridge,
cap collar and based arc. It does not replace their matrices
by a projection degree. All seams are the original seams.

Now let \(X=(X_0,\ldots,X_8)\in S^8\subset\mathbb R^9\).
On the open suspension coordinates set
\[
 u=(1-X_0)/2,\qquad
 Y=(X_1,\ldots,X_8)/\sqrt{1-X_0^2}\in S^7,
\]
\[
 s=2\operatorname{atan2}(|Y_\perp|,Y_0),\qquad
 V=Y_\perp/|Y_\perp|\in S^6,\qquad
 z=\frac{r_B}{\pi}
       \operatorname{atan2}(|V_\perp|,V_0)
       \frac{V_\perp}{|V_\perp|}.
                                                        \tag{OX.4}\label{p1b:eq:natural-ox-4}
\]
At \(V=\mathrm{north}\) use \(z=0\).
At \(V=\mathrm{south}\), at both \(Y\) poles and both \(X\)
poles use the exact column \(e_3\).
The complete map is
\[
 \boxed{
 \gamma_{\rm cart}(X)=f_{(1+8u)/9}(s,z):
                  S^8\longrightarrow S^5 .
 }                                                      \tag{OX.5}\label{p1b:eq:natural-ox-5}
\]
These assignments are continuous: \eqref{p1b:eq:natural-ov-4} fixes both annular
ends; the original circle base is fixed; the entire spatial
boundary is fixed; and the cap center has its uniform original
matrix extension. The inverse open chart is explicit:
\[
\begin{gathered}
 V=(\cos(\pi|z|/r_B),\sin(\pi|z|/r_B)z/|z|), \\
 Y=(\cos(s/2),\sin(s/2)V), \\
 X=(1-2u,2\sqrt{u(1-u)}Y).
\end{gathered}
 \tag{OX.6}\label{p1b:eq:natural-ox-6}
\]

Collapse the complement of the actual embedded \eqref{p1b:eq:natural-ox-1} ball
in the original spatial \(S^6\), and retain the two parameter
quotients. This is the precise degree-minus-one source collapse
relating the original \eqref{p1b:eq:natural-ov-4} map to \eqref{p1b:eq:natural-ox-5}, by \eqref{p1b:eq:natural-ox-2}.
Thus the two detector parities agree; reversing this displayed
source orientation changes the stable coefficient's sign and
does not change its parity.
No source, bundle or joining identification is inferred from
a benchmark value.

This constructs the actual attaching map for the existing
cone \(C_{\gamma_{\rm cart}}=S^5\cup_{\gamma_{\rm cart}}D^9\).
The exact whole-map proof \eqref{p1b:eq:natural-pf-3}--\eqref{p1b:eq:natural-pf-9}, rather than a sampled
simplicial approximation, evaluates its square:
\[
 \operatorname{Sq}^4u_5=\epsilon_{\rm ann}v_9,\qquad
 \epsilon_{\rm ann}=0\text{ by \eqref{p1b:eq:natural-pf-3}--\eqref{p1b:eq:natural-pf-9}}.         \tag{OX.7}\label{p1b:eq:natural-ox-7}
\]
The intrinsic annulus \eqref{p1b:eq:natural-ow-5} preserves the whole original filling
criterion. Its retained boundary and detector normalization are
\[
\boxed{\begin{gathered}
 K_{\rm int}:D^8\to SU(3),\\
 K_{\rm int}|_{S^7}=\mathcal T_{\rm int}\mathcal R_J,\\
 \operatorname{Sq}^4u_5=\zeta_{\rm nat} v_9\text{ on }C_{[K'e_3]},\\
 h_t=1+\zeta_{\rm nat},\quad w_t=1.
\end{gathered}}
\tag{OX.8}\label{p1b:eq:natural-ox-8}
\]
The specific correction evaluated in \eqref{p1b:eq:natural-ow-9} contributes zero.
It does not cancel another boundary contribution by itself. The
characteristic and joining terms are evaluated by \eqref{p1b:eq:natural-pg-14}--\eqref{p1b:eq:natural-pg-16}
and \eqref{p1b:eq:natural-ci-8}, \eqref{p1b:eq:natural-aj-4}--\eqref{p1b:eq:natural-aj-5}, with their complete source maps retained.



\subsection{Eliminate the actual bridge from the terminal norm}
\label{p1b:sec:natural-bridge-elimination}

The executed dependencies are the retained intrinsic bridge \(B_\rho\),
the actual swap and constant cap frame, the relative spatial
reflection null, and a literal sweep of the resulting square.
The prescribed root \(H^3\), central-minus closing, matching toggle one,
original \eqref{p1b:eq:natural-nd-8}, sphere gauge and all lower/frame comparisons are fixed.
No independent primitive map is substituted.

Write \(q=Q_s^*H^3Q_s\), \(P=\psi(U_g)\) and
\(B_0=q^2,\ B_1=P\). Retain the exact \eqref{p1b:eq:natural-op-1} equations
\[
 Dq(\sigma z)D=q^{-1},\qquad DP(\sigma z)D=P,\qquad
 B_\rho|_{\partial B}=I.                              \tag{OY.1}\label{p1b:eq:natural-oy-1}
\]
Both roots here are the actual matrices in the original cap frame.

For an arbitrary actual cap map \(f:B\to SU(2)\), define the following
path, retaining its generally nonidentity initial value.
In the first third use the original \eqref{p1b:eq:natural-op-4} source path \(\sigma_h\)
and
\[
 d_h=\operatorname{diag}(e^{i\pi(1-h)/2},
                         e^{-i\pi(1-h)/2},1),\qquad
 Z_f(h/3,z)=f(z)^{-1}d_h f(\sigma_hz)^{-1}d_h^{-1}.
                                                               \tag{OY.2}\label{p1b:eq:natural-oy-2}
\]
It starts at \(f^{-1}(Df(\sigma z)D)^{-1}\) and ends at
\(f(z)^{-1}f(rz)^{-1}\), where \(r\) is the retained single reflection.
For \(a=(z_x-z_p)/\sqrt2,\ b=(z_x+z_p)/\sqrt2\), put
\(\ell=(r_B^2-|a_\perp|^2-|b|^2)^{1/2}\) and
\(u=(1+a_0/\ell)/2\). Let \(z_v(u)\) replace \(a_0\) by
\(\ell(2u-1)\), retaining every other coordinate.
On the next two thirds use
\[
\begin{aligned}
 Z_f((1+h)/3,z)
 &=f(z_v((1-h)u+h\min(2u,1)))^{-1}\\
 &\quad\cdot f(z_v(1-(1-h)u-h\max(2u-1,0)))^{-1},\\
 Z_f((2+h)/3,z)
 &=f(z_v((1-h)\min(2u,2-2u)))^{-1}.
\end{aligned}                                                   \tag{OY.3}\label{p1b:eq:natural-oy-3}
\]
When \(\ell=0\) every value is the actual boundary identity.
Only that zero interval is collapsed.
At the pinch seam one factor is the boundary identity, so the
two displayed formulas agree exactly. Uniform continuity of \(f\)
proves continuity at \(\ell=0\).
Thus \(Z_f(1)=I\); every spatial boundary is fixed.
This is a path from the specified norm to \(I\), not a null of a
closed sphere map unless its initial norm is literally identity.

Apply \eqref{p1b:eq:natural-oy-2}--\eqref{p1b:eq:natural-oy-3} to the entire original bridge:
\[
\begin{gathered}
 F(\rho,\lambda,z)=Z_{B_\rho}(\lambda,z),\qquad F(\rho,0)=N_B(\rho),\qquad F(\rho,1)=I, \\
 F(0,\lambda)=Z_{q^2}(\lambda),\qquad F(1,\lambda)=Z_P(\lambda).
\end{gathered}
 \tag{OY.4}\label{p1b:eq:natural-oy-4}
\]
All values are in the actual intrinsic \(SU(2)\).
The two vertical edges have been constructed using their actual inputs,
not inferred from equality of their homotopy classes.

\begin{theorem}[An applied whole-matrix bridge elimination]
\label{p1b:theorem:natural-bridge-elimination}
There is a specified intrinsic homotopy on the original \(S^7\)
\[
 \mathcal T_{\rm int}\mathcal R_J
       \simeq \mathcal V_P\mathcal Q_{q^2}^{\rm rev}\mathcal R_J,
 \quad
 \mathcal Q_{q^2}(\lambda,z)=Z_{q^2}(\lambda,z),
 \quad
 \mathcal V_P=\overline{\mathcal J}\,*\,Z_P .
                                                               \tag{OY.5}\label{p1b:eq:natural-oy-5}
\]
Here \(\overline{\mathcal J}\) is the original \(I\to P^{-2}\)
square-null path, and \(\mathcal Q^{\rm rev}\) reverses the path,
rather than silently replacing its actual values by pointwise inverses.
The loops are inserted on disjoint portions of the original second cut.
The original \(\mathcal R_J\), including \eqref{p1b:eq:natural-nd-8} and its complete
neck/direction multipliers, is an unchanged right factor.
\end{theorem}
\begin{proof}
Keep the first half \(J_{1-2t}(U_g)\) of \eqref{p1b:eq:natural-ov-5} unchanged.
On its second half put \(u=2t-1\).
First reparameterize the original \(\rho=1-u\) path to
\(\rho=\min(1,\max(0,2-3u))\), by linear interpolation.
This fixes both seam values.
Then, for height \(h\in[0,1]\), evaluate the complete \(F\) on
\[
 (\rho,\lambda)=
 \begin{cases}
 (1,3hu),&u\le1/3,\\
 (2-3u,h),&1/3\le u\le2/3,\\
 (0,3h(1-u)),&u\ge2/3 .
 \end{cases}                                             \tag{OY.6}\label{p1b:eq:natural-oy-6}
\]
At \(h=0\) this is the reparameterized bottom edge.
At \(h=1\) it is the right edge, the constant top edge, and the
reversed left edge of the actual square.
Their joins agree by \eqref{p1b:eq:natural-oy-4}.
This gives the claimed intrinsic homotopy with all values supplied.
More explicitly, its endpoint uses \(J_{1-2t}\) for \(t\le1/2\),
\(Z_P(6t-3)\) for \(1/2\le t\le2/3\), identity up to \(5/6\),
and \(Z_{q^2}(6-6t)\) for \(5/6\le t\le1\).
The first cut and the cap boundary remain identity.
Insert this whole homotopy through the actual \eqref{p1b:eq:natural-ox-1} embedding.
The cap faces allow pasting to the unchanged original joining;
right multiplication by that joining gives \eqref{p1b:eq:natural-oy-5}.
Every new full matrix fixes \(e_3\), so the complete projected
track of this applied homotopy is constant.
\end{proof}

This completes elimination of the original bridge from this target
calculation. It leaves its two specified endpoint comparisons;
it does not prove either comparison null.
In particular the zero in \eqref{p1b:eq:natural-ow-9} evaluates its own mirror correction,
not the newly constructed \(\mathcal V_P\).

The existing detector now has the equivalent, actually constructed
terminal boundary
\[
 \boxed{
 K_{\rm end}:D^8\longrightarrow SU(3),\qquad
 K_{\rm end}|_{S^7}
       =\mathcal V_P\mathcal Q_{q^2}^{\rm rev}\mathcal R_J .
 }                                                        \tag{OY.7}\label{p1b:eq:natural-oy-7}
\]
An actual filling of this boundary, pasted with the reverse \eqref{p1b:eq:natural-oy-5}
annulus and the original \eqref{p1b:eq:natural-ov-3} annulus, supplies the original \(K'\).
The new annulus adds no projected column.
Neither that filling nor the whole projected parity is supplied by
the square sweep itself:
\[
 \operatorname{Sq}^4u_5=\zeta_{\rm nat} v_9
       \text{ on }C_{[K'e_3]},\qquad
 h_t=1+\zeta_{\rm nat},\qquad w_t=1.                               \tag{OY.8}\label{p1b:eq:natural-oy-8}
\]
The complete original comparison is evaluated in
\eqref{p1b:eq:natural-aj-4}--\eqref{p1b:eq:natural-aj-7}:
\(h_t=w_t=1\), \(r_J=\zeta_{\rm nat}=0\), with every retained
extension test decided. This component is used with its full frames,
lower divisions and joining; its isolated value is not a replacement
for that complete comparison.


\subsection{The complete root-power carry and its actual relative source}
\label{p1b:sec:natural-root-power-carry}

Apply the constructed \eqref{p1b:eq:natural-oy-2}--\eqref{p1b:eq:natural-oy-3} path to \(q\) and \(q^2\), with
the same original source and cap frame. The exact exchange \eqref{p1b:eq:natural-oy-1}
makes both initial norms identity. Thus both are specified closed
\(S^7\to SU(2)\) maps, written \(\mathcal Q_q,\mathcal Q_{q^2}\).
Construct their complete root-power carry, rather than assuming that
the norm of a square is the square of the norm:
\[
 \mathcal W_q(\lambda,z)
     =\mathcal Q_{q^2}(\lambda,z)
          \mathcal Q_q(\lambda,z)^{-2},\qquad
 \mathcal Q_{q^2}=\mathcal W_q\mathcal Q_q^2 .
                                                               \tag{OZ.1}\label{p1b:eq:natural-oz-1}
\]
This is a pointwise full-matrix equality on the original source.
It has not replaced the full boundary in \eqref{p1b:eq:natural-oy-7} by a homotopy-class
representative.

On the straightening third of the actual path put
\[
\begin{gathered}
 a=q(z)^{-1},\qquad b=d_hq(\sigma_hz)^{-1}d_h^{-1};\qquad \text{on its pinch third put}\qquad a=q(z_v(u_1))^{-1}, \\
 b=q(z_v(1-u_2))^{-1}.
\end{gathered}
 \tag{OZ.2}\label{p1b:eq:natural-oz-2}
\]
Here \(u_1,u_2\) are exactly the retained \eqref{p1b:eq:natural-oy-3} coordinates.
At every such point \(\mathcal Q_q=ab\) and
\(\mathcal Q_{q^2}=a^2b^2\), so direct multiplication proves
\[
 \boxed{
 \mathcal W_q=\mathcal F(a,b),\qquad
 \mathcal F(a,b)=a^2ba^{-1}b^{-1}a^{-1}
                 =a(aba^{-1}b^{-1})a^{-1}.
 }                                                        \tag{OZ.3}\label{p1b:eq:natural-oz-3}
\]
The conjugator in this actual word has been retained.
On the backtracking third, both complete words are powers of the
same \(q(z_v(\cdot))^{-1}\). Consequently
\(\mathcal W_q=I\) there, pointwise; this specified part of the carry
has been evaluated as zero.

Both collapsed faces of the remaining word matter.
At its first face \eqref{p1b:eq:natural-oy-1} gives \(b=a^{-1}\); the pair need not belong
to the coordinate wedge. At the pinch endpoint one of \(a,b\) is
identity. The entire spatial boundary has both values identity.
The universal word \eqref{p1b:eq:natural-oz-3} is literally identity on all these faces.
It therefore has the exact, forced relative domain
\[
 \mathcal X=
 \frac{SU(2)\times SU(2)}
 {(SU(2)\times\{I\})\cup(\{I\}\times SU(2))
                 \cup\{(a,a^{-1}):a\in SU(2)\}},
 \qquad \overline{\mathcal F}:\mathcal X\longrightarrow SU(2).
                                                               \tag{OZ.4}\label{p1b:eq:natural-oz-4}
\]
This quotient records the faces of one actual comparison.
It is not a new detector or a replacement of the existing
stabilized filling/finite mapping-cone detector.

Use the entire pair \eqref{p1b:eq:natural-oz-2} for \(0\le\lambda\le2/3\), and collapse
its two end faces and its spatial boundary. This gives a specified map
\[
 \boxed{
 \alpha_q:S^7\longrightarrow\mathcal X,\qquad
 \mathcal W_q=\overline{\mathcal F}\alpha_q .
 }                                                        \tag{OZ.5}\label{p1b:eq:natural-oz-5}
\]
All entries use the original \(Q_s^*H^3Q_s\), the original
\(\sigma_h,d_h,z_v,u_1,u_2\). Equations \eqref{p1b:eq:natural-oz-2}--\eqref{p1b:eq:natural-oz-5} construct
the input from the actual comparison, including both cancellations.
No benchmark source of the same apparent degree has been inserted.

For coefficients only, the retained \(\pi_7SU(2)=C_2\) and the exact
\eqref{p1b:eq:natural-oz-1} imply
\[
 {}[\mathcal Q_{q^2}]=[\mathcal W_q]
       =[\overline{\mathcal F}\alpha_q].                 \tag{OZ.6}\label{p1b:eq:natural-oz-6}
\]
This justified group equality does not supply an actual nullhomotopy
of \(\mathcal Q_q^2\). Accordingly the full original matrix
\(\mathcal Q_{q^2}^{\rm rev}\), rather than a discarded square
factor, remains in the complete \eqref{p1b:eq:natural-oy-7} boundary.

The actual relative identification, supplied in \eqref{p1b:eq:natural-pb-1}--\eqref{p1b:eq:natural-pc-7},
is a comparison of
\(\overline{\mathcal F}\alpha_q\) with a source-identified evaluated
seven-sphere map, retaining its antidiagonal null face and conjugator.
In particular the coordinate wedge alone does not define
\(\alpha_q:S^7\to SU(2)\wedge SU(2)=S^6\):
its first face is generally \((q^{-1},q)\), with neither entry identity.
The known class of \(H\), a primitive commutator or a self-interchange
calibration therefore does not evaluate this actual input.
Any application of a secondary track formula must prove that formula's
track is the complete \eqref{p1b:eq:natural-oz-2}, including its two end-face comparisons.
The complete antidiagonal and closing tracks are retained in
\eqref{p1b:eq:natural-pb-1}--\eqref{p1b:eq:natural-pc-7}, rather than replaced by this coordinate wedge.

The original detector criterion retains the actual complete maps:
\[
 \boxed{
 K_{\rm end}|_{S^7}
       =\mathcal V_P\mathcal Q_{q^2}^{\rm rev}\mathcal R_J,\qquad
 \operatorname{Sq}^4u_5=\zeta_{\rm nat} v_9\text{ on }C_{[K'e_3]},
 \qquad h_t=1+\zeta_{\rm nat},\quad w_t=1.
 }                                                        \tag{OZ.7}\label{p1b:eq:natural-oz-7}
\]
The full relative root carry is evaluated in \eqref{p1b:eq:natural-pc-7}, the complete
characteristic loop in \eqref{p1b:eq:natural-pg-14}, and the original joining in
\eqref{p1b:eq:natural-aj-4}--\eqref{p1b:eq:natural-aj-5}. Each uses its retained source comparison.
The complete original comparison is evaluated in
\eqref{p1b:eq:natural-aj-4}--\eqref{p1b:eq:natural-aj-7}:
\(h_t=w_t=1\), \(r_J=\zeta_{\rm nat}=0\), with every retained
extension test decided. This component is used with its full frames,
lower divisions and joining; its isolated value is not a replacement
for that complete comparison.


\subsection{Localize the actual root exchange and retain its closing discrepancy}
\label{p1b:sec:natural-root-exchange-localization}

Use the actual spatial cap compactification
\(V=(\cos(\pi|z|/r_B),\sin(\pi|z|/r_B)z/|z|)\)
and its south-based stereographic coordinate \(x=V_\perp/(1+V_0)\).
In the original six-coordinate order put
\[
\begin{gathered}
 e=(1,0,0,-1,0,0)/\sqrt2,\qquad
 x_h^\pm=x\mp2he,\\
 c_h^\pm(z)=
 \begin{cases}
 \operatorname{cap}\operatorname{st}^{-1}
       (x_h^\pm/(1-h|x_h^\pm|^2)),&h|x_h^\pm|^2<1,\\
 \partial B,&h|x_h^\pm|^2\ge1 .
 \end{cases}
\end{gathered}
\tag{PA.1}\label{p1b:eq:natural-pa-1}
\]
The south pole and \(h=0\) have their original based extensions.
The homogeneous inverse stereographic formula proves joint continuity
at the collapse seam and at the based point, as in \eqref{p1b:eq:natural-ow-4}.
Both maps start at identity and have degree one.
At \(h=1\) their supports are the disjoint open unit six-balls
centered at \(2e,-2e\); the distance between centers is four.
This uses the actual minus direction of the original exchange.

Every \(\sigma_k\) in the retained \eqref{p1b:eq:natural-op-4} family sends \(e\) to \(-e\)
and fixes the south pole. The scalar denominators in \eqref{p1b:eq:natural-pa-1} therefore
give the exact source equations
\[
 c_h^\pm(\sigma_k z)=\sigma_k c_h^\mp(z).
                                                               \tag{PA.2}\label{p1b:eq:natural-pa-2}
\]
In particular this holds for the original \(\sigma\) and final reflection.

Construct and apply the following homotopy of the actual root square:
\[
 f_h(z)=q(c_h^+z)q(c_h^-z),\qquad f_0=q^2,\qquad
 Df_h(\sigma z)D=f_h(z)^{-1}.                            \tag{PA.3}\label{p1b:eq:natural-pa-3}
\]
The last equation follows in the displayed order: exchange sends the
first factor to the inverse of the second and the second to the inverse
of the first. No commutation of overlapping factors is assumed.
All spatial boundary values are identity.
Consequently the already constructed path \(Z_{f_h}\) has both
parameter ends identity for every \(h\), and supplies the complete
intrinsic root-loop homotopy
\[
 \mathcal Q_{q^2}\simeq\mathcal Q_{f_1},\qquad
 \Phi_h=\mathcal V_P\mathcal Q_{f_h}^{\rm rev}\mathcal R_J .
                                                               \tag{PA.4}\label{p1b:eq:natural-pa-4}
\]
This is applied on the same last portion of the actual second cut in
\eqref{p1b:eq:natural-oy-5}. The actual cap faces allow pasting to the unchanged original
joining. All full matrices fix \(e_3\).
This changes a comparison map inside the cap; the prescribed
original \(Y\) fibre representative and its closing equation are retained.

At \(h=1\), on the entire first straightening third,
disjoint supports and \eqref{p1b:eq:natural-pa-2} prove the exact whole-matrix equation
\[
 \mathcal Q_{f_1}(\lambda,z)
  =\mathcal Q_q(\lambda,c_1^+z)\,
                    \mathcal Q_q(\lambda,c_1^-z),
 \qquad 0\le\lambda\le1/3.                              \tag{PA.5}\label{p1b:eq:natural-pa-5}
\]
Here each local factor still uses the original \(\sigma_k,d_k\);
no root normal framing or target-frame path has been reset.
This equation does not assert equality of their later reflection
pinches, which use different complete spatial interval coordinates.

Retain that difference by constructing its full closed matrix map:
\[
 \boxed{
 \mathcal E_q(\lambda,z)
  =\mathcal Q_{f_1}(\lambda,z)
       \left(\mathcal Q_q(\lambda,c_1^+z)
                    \mathcal Q_q(\lambda,c_1^-z)\right)^{-1}
       :S^7\longrightarrow SU(2).
 }                                                        \tag{PA.6}\label{p1b:eq:natural-pa-6}
\]
It is literally identity on the first third, both circle ends and the
spatial boundary. It retains both global and both local closing tracks
on the remaining two thirds.
Equation \eqref{p1b:eq:natural-pa-6} is a constructed map of the actual comparison, including
its cancellations; the entire straightening contribution to this
particular discrepancy has been evaluated as zero.
The later discrepancy is not declared zero.

For its coefficient, both \(c_1^\pm\) have the same degree one and keep
the original based cap. The retained top group therefore gives
\[
 {}[\mathcal Q_{q^2}]=[\mathcal Q_{f_1}]
       =[\mathcal E_q]+2[\mathcal Q_q]=[\mathcal E_q].
                                                               \tag{PA.7}\label{p1b:eq:natural-pa-7}
\]
Every full matrix remains in \eqref{p1b:eq:natural-pa-4}; no intrinsic nullhomotopy of a
square has been supplied by this group equality.
Thus the actual remaining root coefficient is now the coefficient of
the explicit \eqref{p1b:eq:natural-pa-6} closing-track discrepancy.
Assigning the value of a two-copy self-interchange would additionally
require a relative homotopy identifying that entire map with the proposed
interchange, including the global and local pinch matchings, their fixed
cap frames and both backtracking tracks.
That named comparison has not been proved here.

The equivalent original detector input is now
\[
 \boxed{
 K_{\rm loc}:D^8\to SU(3),\qquad
 K_{\rm loc}|_{S^7}
       =\mathcal V_P\mathcal Q_{f_1}^{\rm rev}\mathcal R_J,
 \qquad
 \operatorname{Sq}^4u_5=\zeta_{\rm nat} v_9
       \text{ on the complete }C_{[K'e_3]} .
 }                                                        \tag{PA.8}\label{p1b:eq:natural-pa-8}
\]
Pasting with the actual reverse \eqref{p1b:eq:natural-pa-4}, reverse \eqref{p1b:eq:natural-oy-5}, and retained
\eqref{p1b:eq:natural-ov-3} comparisons recovers precisely the original target.
Both new annuli have the constant projected column \(e_3\).
The complete original comparison is evaluated in
\eqref{p1b:eq:natural-aj-4}--\eqref{p1b:eq:natural-aj-7}:
\(h_t=w_t=1\), \(r_J=\zeta_{\rm nat}=0\), with every retained
extension test decided. This component is used with its full frames,
lower divisions and joining; its isolated value is not a replacement
for that complete comparison.


\subsection{Remove the actual first path from the root-power carry}
\label{p1b:sec:natural-root-path-comparison}

Retain the actual \(q\), its original cap frame, and every later pinch
and backtracking value in \eqref{p1b:eq:natural-oy-3}. Define the entire closing part of the
single-root loop by
\[
 H_v(z)=\mathcal Q_q((1+2v)/3,z),\qquad
 H_0=q^{-1}q_r^{-1},\quad H_1=I,\qquad q_r(z)=q(rz).
                                                               \tag{PB.1}\label{p1b:eq:natural-pb-1}
\]
It is the specified reflection pinch followed by its backtrack.
On the first third, with \(k=3\lambda\), put
\[
 a=q^{-1},\quad b_k=d_kq(\sigma_kz)^{-1}d_k^{-1},\quad
 \widehat b_k=qH_{1-k},\quad
 b_0=\widehat b_0=q,\quad b_1=\widehat b_1=q_r^{-1}.
                                                               \tag{PB.2}\label{p1b:eq:natural-pb-2}
\]
Both paths have the same actual endpoints. The first uses the retained
source and constant-frame path. The second is constructed from the
actual closing part, without substituting a representative of its class.

Use \(\mathcal F(a,\widehat b_k)\) on this first third and exactly
the original \eqref{p1b:eq:natural-oz-3} values on the pinch third; use identity on the
backtracking third. Denote this full closed map by
\(\mathcal W_{\rm can}:S^7\to SU(2)\).
It retains the initial antidiagonal and final wedge faces.
There is no omission of the seam at \(\lambda=1/3\).

For the complete difference on the first third define
\[
 L_k=b_k\widehat b_k^{-1},\qquad
 C_k=\widehat b_k a\widehat b_k^{-1},\qquad
 \boxed{\mathcal W_q\mathcal W_{\rm can}^{-1}
      =a^2L_kC_k^{-1}L_k^{-1}C_ka^{-2}.}
                                                               \tag{PB.3}\label{p1b:eq:natural-pb-3}
\]
Direct multiplication of the two \eqref{p1b:eq:natural-oz-3} words proves the equality.
In particular \(L_0=L_1=I\). All spatial boundary values of
\(L,C,a\) are identity. Outside this first third the difference is
literally identity. The displayed conjugator is retained.

Apply the actual based cap compressions \eqref{p1b:eq:natural-pa-1} separately to these
two named inputs:
\[
\begin{gathered}
\mathcal D_h(\lambda,z)=
 \begin{cases}
 a(z)^2 L_k(c_h^+z)C_k(c_h^-z)^{-1}
                 L_k(c_h^+z)^{-1}C_k(c_h^-z)a(z)^{-2},
                                      &0\le\lambda\le1/3,\\
 I,&1/3\le\lambda\le1 .
 \end{cases}\\
\mathcal D_0=\mathcal W_q\mathcal W_{\rm can}^{-1},
 \quad \mathcal D_1=I.
\end{gathered}
\tag{PB.4}\label{p1b:eq:natural-pb-4}
\]
At the two time faces \(L=I\), independently of the compressed
\(C\). The spatial boundary is fixed as well. At \(h=1\) the two
supports are disjoint, so at every point either \(L\) or \(C\) is
identity. Thus this is a jointly continuous, relative intrinsic null
of the entire actual difference \eqref{p1b:eq:natural-pb-3}.
It evaluates this named comparison as zero, including its moving
conjugator and all collapsed faces.

Apply it immediately to the full root loop:
\[
 \boxed{
 \mathcal Q_h^{\rm can}
   =\mathcal D_h\mathcal W_{\rm can}\mathcal Q_q^2,\qquad
 \mathcal Q_0^{\rm can}=\mathcal Q_{q^2},\quad
 \mathcal Q_1^{\rm can}=\mathcal W_{\rm can}\mathcal Q_q^2 .
 }                                                        \tag{PB.5}\label{p1b:eq:natural-pb-5}
\]
Every displayed equality is an equality of complete matrices.
The square of the actual single-root loop has not been discarded.
Its known coefficient is zero because \(\pi_7SU(2)=C_2\);
that class statement has not supplied a nullhomotopy of the square.
The new homotopy fixes \(e_3\) everywhere.

To start precisely at the latest \eqref{p1b:eq:natural-pa-4} endpoint, first use that
actual localization backwards, then use \eqref{p1b:eq:natural-pb-5}:
\[
 \mathcal Q_h^{\rm appl}=
 \begin{cases}
 \mathcal Q_{f_{1-2h}},&h\le1/2,\\
 \mathcal Q_{2h-1}^{\rm can},&h\ge1/2,
 \end{cases}\qquad
 \Phi_h^{\rm can}
    =\mathcal V_P(\mathcal Q_h^{\rm appl})^{\rm rev}\mathcal R_J .
                                                               \tag{PB.6}\label{p1b:eq:natural-pb-6}
\]
These loops occupy the same last interval of the actual second cut
as in \eqref{p1b:eq:natural-oy-6}. Their seam is literally \(\mathcal Q_{q^2}\).
Insert them through the original \eqref{p1b:eq:natural-ox-1} embedding and retain the
original joining as an unchanged right factor on its entire region.
The constant projected column proves that this applied whole-boundary
comparison adds no projected track.
The original fibre representative, central-minus closing, toggle,
lower comparisons, coordinates and bundle identifications are fixed.

The coefficient now has the justified actual comparison
\[
 {}[\mathcal Q_{q^2}]=[\mathcal W_{\rm can}]+2[\mathcal Q_q]
                     =[\mathcal W_{\rm can}].             \tag{PB.7}\label{p1b:eq:natural-pb-7}
\]
This removes the first source/frame path from the coefficient.
It does not evaluate the canonical pinch carry.
In particular a standard self-interchange or a James--Hopf value
cannot be applied until its complete relative track is identified
with this specified map. The original conjugated commutator,
antidiagonal face and closing cancellation remain in its definition.

The exact original detector has the following equivalent boundary,
constructed from the actual comparisons:
\[
\boxed{\begin{gathered}
 K_{\rm can}:D^8\to SU(3),\\
 K_{\rm can}|_{S^7}
   =\mathcal V_P
        (\mathcal W_{\rm can}\mathcal Q_q^2)^{\rm rev}\mathcal R_J,\\
 \operatorname{Sq}^4u_5=\zeta_{\rm nat} v_9
              \text{ on }C_{[K'e_3]},\\
 h_t=1+\zeta_{\rm nat},\quad w_t=1 .
\end{gathered}}
\tag{PB.8}\label{p1b:eq:natural-pb-8}
\]
Paste reverse \eqref{p1b:eq:natural-pb-6}, reverse \eqref{p1b:eq:natural-pa-4}, reverse \eqref{p1b:eq:natural-oy-5}, and the
retained \eqref{p1b:eq:natural-ov-3} annulus to obtain the original \(K'\).
Every new annulus is intrinsic. The complete annular projected
input remains the actually constructed \eqref{p1b:eq:natural-ox-5}.
Neither \(K_{\rm can}\), the endpoint loop coefficient, the canonical
carry coefficient, the original joining coefficient, nor the full
projected square has been evaluated by \eqref{p1b:eq:natural-pb-4}.
The complete original comparison is evaluated in
\eqref{p1b:eq:natural-aj-4}--\eqref{p1b:eq:natural-aj-7}:
\(h_t=w_t=1\), \(r_J=\zeta_{\rm nat}=0\), with every retained
extension test decided. This component is used with its full frames,
lower divisions and joining; its isolated value is not a replacement
for that complete comparison.


\subsection{Evaluate the complete actual canonical root carry}
\label{p1b:sec:natural-canonical-carry-evaluation}

The dependencies are the actual relative comparison \eqref{p1b:eq:natural-pb-1}--\eqref{p1b:eq:natural-pb-6},
its original reflection coordinate, the signed-word two-cycle of that
specific closing track, and the suspension adjunction of the actual
root. We execute them before applying a known seven-sphere value.
The original root, cap frame, joining, and square factor are retained.

Write \(v=(a_1,a_2,b)\in D^5\),
\(\ell=(r_B^2-|v|^2)^{1/2}\), and use exactly the original coordinate
\[
 \Theta(u,v):\quad a_0=\ell(2u-1),\quad
 (a_1,a_2,b)=v,\qquad
 q_v(u)=q(\Theta(u,v)),\qquad q_v(0)=q_v(1)=I.
                                                               \tag{PC.1}\label{p1b:eq:natural-pc-1}
\]
On \(\partial D^5\) every value is identity. Thus \(\Theta\) induces
the actual based identification \(S^1\wedge S^5\cong B/\partial B\).
The original source axes and every transverse coordinate are retained.
Continuity at \(\ell=0\) is the same uniform-continuity statement
used in \eqref{p1b:eq:natural-oy-3}.

Let \(F(S^1)\) be the free based word group, in its CW word model,
and let \(\gamma(u)\) denote its generator with
\(\gamma(0)=\gamma(1)=1\). For each actual \(v\), the based map
\(q_v\) extends to the continuous homomorphism
\(e_v:F(S^1)\to SU(2)\). This is a jointly continuous family,
trivial on \(\partial D^5\).
Put
\[
\begin{gathered}
 p(u)=\min(2u,1), \\
 m(u)=\max(2u-1,0), \\
 j(u)=p(u)-m(u), \\
 l_h=(1-h)u+hp(u), \\
 b_h=1-(1-h)u-hm(u).
\end{gathered}
 \tag{PC.2}\label{p1b:eq:natural-pc-2}
\]
The actual closing part in \eqref{p1b:eq:natural-pb-1} becomes the following universal word:
\[
 \widehat H_v(u)=
 \begin{cases}
 \gamma(l_{2v})^{-1}\gamma(b_{2v})^{-1},&v\le1/2,\\
 \gamma((2-2v)j(u))^{-1},&v\ge1/2.
 \end{cases}
                                                               \tag{PC.3}\label{p1b:eq:natural-pc-3}
\]
Here the letter \(v\) in \(\widehat H_v\) is the closing parameter;
it does not change the transverse coordinate in \eqref{p1b:eq:natural-pc-1}.
Its seam is an exact word equality: for \(u\le1/2\) the second
pinch letter is based, and for \(u\ge1/2\) the first is based.

Define \(U:I^2/\partial I^2=S^2\to F(S^1)\), using the very same
three intervals as \(\mathcal W_{\rm can}\):
\[
\begin{gathered}
U(\lambda,u)=
 \begin{cases}
 {}[\gamma(u)^{-1},\widehat H_{1-k}(u)],\quad k=3\lambda,
                                                    &\lambda\le1/3,\\
 \gamma(l_h)^{-2}\gamma(b_h)^{-1}
       \gamma(l_h)\gamma(b_h)\gamma(l_h),\quad h=3\lambda-1,
                                                    &1/3\le\lambda\le2/3,\\
 1,                                                &\lambda\ge2/3.
 \end{cases}\\
\boxed{\mathcal W_{\rm can}(\lambda,\Theta(u,v))=e_v(U(\lambda,u)).}
\end{gathered}
\tag{PC.4}\label{p1b:eq:natural-pc-4}
\]
The boxed equality is strict, not a comparison of primitive classes.
Indeed \(\mathcal F(q^{-1},qH)=q^{-1}HqH^{-1}\) on the first
third, and the second word is exactly \eqref{p1b:eq:natural-oz-3}.
At \(k=0\) the commutator is identity. At \(k=1\) it is precisely
the second-third word at \(h=0\). At \(h=1\) it is identity.
All four \((\lambda,u)\) boundary faces are based.
This proves the actual relative source identification, including
the original antidiagonal face and closing cancellation.

\paragraph{An exact two-cycle calculation for this particular word.}
On a signed-word chart
\(\gamma(t_1)^{\epsilon_1}\cdots\gamma(t_r)^{\epsilon_r}\) use
\[
 \omega_r=\sum_{i<j}\epsilon_i\epsilon_j\,dt_i\wedge dt_j,
 \qquad \epsilon_i\in\{1,-1\}.                       \tag{PC.5}\label{p1b:eq:natural-pc-5}
\]
This defines a cellular real two-cocycle in the CW word model.
Here is the needed compatibility argument. Triangulate each signed
word cube by all coordinate-order simplices. Every base-letter face
and every adjacent inverse-letter diagonal is a subcomplex.
On a base-letter face its differential is zero, so deletion gives
exactly the shorter-word form. On an adjacent cancellation diagonal
\(t_i=t_{i+1}\), \(\epsilon_{i+1}=-\epsilon_i\), terms involving
an outside letter cancel in pairs and the internal wedge is zero.
Thus the form pulls back from the reduced word on every identified
face. Iterated reductions have the same property.
Integrate these compatible constant forms on the two-simplices of
the word complex. Stokes on each three-simplex gives zero coboundary;
identified or collapsed faces agree by the just-proved compatibility.
This constructs a cellular cohomology class and justifies evaluating
it on the complete piecewise-smooth cycle \eqref{p1b:eq:natural-pc-4}.
It is not a new detector for the natural target.

For normalization, the positive two-letter torus
\((s,t)\mapsto\gamma(s)\gamma(t)\) has integral one.
Its homology class is the Pontryagin product \(x^2\), where
\(x\in H_1F(S^1;\mathbb Z)\) is the generator loop.
The James/free-group suspension model gives
\(F(S^1)\simeq\Omega S^2\) and
\(H_2F(S^1;\mathbb Z)=\mathbb Z\{x^2\}\).
These standard model facts, also stated in the introduction and
Section 7 of Weber--Ziemia\'nski, apply to the connected based circle.
The normalization fixes the real cocycle's pairing with the integral
generator; no value of a natural map has been inferred from it.

Orient each moving rectangle by its increasing track coordinate
followed by \(u\). Direct substitution of the full signed words into
\eqref{p1b:eq:natural-pc-5} gives the following exact calculation:
\[
\begin{array}{c|c|c}
 \text{actual part of \eqref{p1b:eq:natural-pc-4}}&\text{pullback of }\omega&\text{integral}\\ \hline
 k\in[0,1/2],\ \widehat H_{1-k}=\gamma(2kj)^{-1}
       &-4j(u)\,dk\wedge du&-1\\
 k\in[1/2,1],\ h=2-2k
       &+4j(u)\,dk\wedge du&+1\\
 h\in[0,1],\ u\le1/2
       &-4u\,dh\wedge du&-1/2\\
 h\in[0,1],\ u\ge1/2
       &-4(1-u)\,dh\wedge du&-1/2\\
 \text{final backtrack}&0&0
\end{array}
\qquad
 \boxed{\langle[\omega],U_*[S^2]\rangle=-1.}          \tag{PC.6}\label{p1b:eq:natural-pc-6}
\]
For example the first two words are respectively
\(\gamma(u)^{-1}\gamma(2kj)^{-1}\gamma(u)\gamma(2kj)\)
and
\(\gamma(u)^{-1}\gamma(l_h)^{-1}\gamma(b_h)^{-1}
  \gamma(u)\gamma(b_h)\gamma(l_h)\).
The second-third word has pullback \(2\,dl_h\wedge db_h\).
Since \(l_h+b_h=1+hj(u)\), all integrals in \eqref{p1b:eq:natural-pc-6} are rational
integrals of the displayed actual coordinates.
In particular the first two contributions cancel; retaining only
one self-interchange would not give the calculation of this map.

\paragraph{Suspend the actual source identification before using the value.}
The Hurewicz map \(\pi_2F(S^1)\to H_2F(S^1;\mathbb Z)\) is an
isomorphism. One may see this by looping the Hopf fibration:
\(\Omega S^3\to\Omega S^2\to S^1\). Its two-dimensional edge
map is an isomorphism because \(H_1\Omega S^3=0\), and
\(\Omega S^3\) is simply connected. The Hopf map induces the
generator on both \(\pi_2\) and \(H_2\). Therefore \eqref{p1b:eq:natural-pc-6} says
\([U]=-\operatorname{adj}(\eta_2)\), with a consistent choice
of orientation.

Apply a based homotopy representing that equality through the
\emph{actual} entire family \(e_v\) in \eqref{p1b:eq:natural-pc-4}.
It remains based on \(\partial D^5\), so it is a homotopy of the
original closed seven-sphere map. Under classifying-space adjunction,
\(Be_v:S^2\to BSU(2)\) is the adjoint of the actual \(q_v\).
Consequently the adjoint of the resulting whole map is
\(\widetilde q\circ(\eta_2\wedge 1_{S^5})
 =\widetilde q\eta_7\), where \(\widetilde q\) is the adjoint
of the actual root on the coordinates \eqref{p1b:eq:natural-pc-1}.
Since \(\eta_7=\Sigma\eta_6\), unadjunction gives
\[
 \boxed{
 {}[\mathcal W_{\rm can}]=-[q\eta_6]
       =-3\nu'\eta_6=\mu\quad\text{in }\pi_7SU(2)=C_2.
 }                                                        \tag{PC.7}\label{p1b:eq:natural-pc-7}
\]
Any orientation reversal in the retained transverse chart is harmless
in this order-two group. The equality was obtained from the actual
full word \eqref{p1b:eq:natural-pc-4}, not from a calibration followed by an assumed
identification. Equation \eqref{p1b:eq:natural-pb-7} now also evaluates
\([\mathcal Q_{q^2}]=\mu\) and the \eqref{p1b:eq:natural-pa-6} complete closing
discrepancy as \(\mu\). The already constructed \eqref{p1b:eq:natural-pb-6} application
therefore puts this evaluated contribution into the actual boundary.

For the existing stabilized detector, the retained source-identified
\(\mu\) filling has odd \(\operatorname{Sq}^4\) coefficient.
The homotopy just proved transports its boundary to the actual
\(\mathcal W_{\rm can}\). This evaluates the named root component.
It does not specify a new full matrix replacement: the complete
\(\mathcal Q_q^2\) factor remains in \eqref{p1b:eq:natural-pb-8}, with class zero.

The immediate whole-boundary consequence and the still exact target are
\[
\boxed{\begin{gathered}
 {}[\mathcal T_{\rm int}]=[\mathcal V_P]+\mu,\\
 K_{\rm can}|_{S^7}
    =\mathcal V_P(\mathcal W_{\rm can}\mathcal Q_q^2)^{\rm rev}
       \mathcal R_J,\\
 \operatorname{Sq}^4u_5=\zeta_{\rm nat} v_9\text{ on }C_{[K'e_3]},\\
 h_t=1+\zeta_{\rm nat},\quad w_t=1.
\end{gathered}}
\tag{PC.8}\label{p1b:eq:natural-pc-8}
\]
The first equality is a coefficient statement on the actual boundary;
every full factor in the second equality is retained.
The complete original comparison is evaluated in
\eqref{p1b:eq:natural-aj-4}--\eqref{p1b:eq:natural-aj-7}:
\(h_t=w_t=1\), \(r_J=\zeta_{\rm nat}=0\), with every retained
extension test decided. This component is used with its full frames,
lower divisions and joining; its isolated value is not a replacement
for that complete comparison.


\subsection{Apply the evaluated carry to the actual determinant comparison}
\label{p1b:sec:natural-carry-subtraction}

Use the evaluated full actual carry \eqref{p1b:eq:natural-pc-7}, not a primitive replacement.
Insert it into the same last interval of the second cut and the same
original cap embedding as \eqref{p1b:eq:natural-pb-6}. Explicitly put
\[
 \mathcal C(s,z)=
 \begin{cases}
 \mathcal W_{\rm can}(6-6t,z),\quad t=s/\pi-1,
                                      &11\pi/6\le s\le2\pi,\\
 I,                                   &\text{otherwise},
 \end{cases}\qquad
 \widehat{\mathcal C}:S^7\to SU(2).
                                                               \tag{PD.1}\label{p1b:eq:natural-pd-1}
\]
Extend by identity off the actual \eqref{p1b:eq:natural-ox-1} cap. The two interval ends,
the cap boundary and every original parameter-base face are identity.
The original transverse orientation, ordinary bundle identification,
and \eqref{p1b:eq:natural-ok-6}--\eqref{p1b:eq:natural-ok-7} frame comparisons are retained.
The reversing interval and the actual cap orientation do not alter
the order-two value, so \([\widehat{\mathcal C}]=\mu\).

Apply this specified factor to the complete actual \eqref{p1b:eq:natural-og-10} core:
\[
 \boxed{\Psi_D^\dagger=\psi'\widehat{\mathcal C}^{-1}:S^7\to SU(2),
 \qquad [\Psi_D^\dagger]=[\psi']+\mu=(\zeta_{\rm nat}+1)\mu=h_t\mu.}
                                                               \tag{PD.2}\label{p1b:eq:natural-pd-2}
\]
The last equality is the already proved actual \eqref{p1b:eq:natural-oh-11} comparison,
with the retained \(w_t=1\). No value of \(h_t\) is asserted.
This constructs an equivalent determinant detector boundary directly
from two actual comparisons whose source identifications are proved.
The original \(Y\) representative, central-minus closing, matching
toggle, original \eqref{p1b:eq:natural-nd-8}, sphere gauge and lower comparisons remain
in \(\psi'\). On the entire joining and outer region the new
factor is literally identity.

For each of the already constructed whole comparisons \(\Phi_h\),
apply the same actual right factor:
\[
 \Phi_h^\dagger=\Phi_h\widehat{\mathcal C}^{-1},\qquad
 \Phi_h^\dagger e_3=\Phi_he_3 .                       \tag{PD.3}\label{p1b:eq:natural-pd-3}
\]
The equality holds for full matrices even on moving-column regions,
because \(\widehat{\mathcal C}^{-1}e_3=e_3\).
All pasting equations therefore remain exact.
Apply it to the full \eqref{p1b:eq:natural-ov-3} annulus and to the subsequent reverse
\eqref{p1b:eq:natural-oy-5}, \eqref{p1b:eq:natural-pa-4}, \eqref{p1b:eq:natural-pb-6} comparisons. No new projected track is added.
The closed Cartesian input remains exactly \(\gamma_{\rm cart}\)
of \eqref{p1b:eq:natural-ox-5}, with all seven moving legs and the intrinsic eighth leg.

At the latest canonical endpoint, \(\mathcal V_P\) and
\(\mathcal C\) occupy disjoint time intervals, and the original
\(\mathcal R_J\) is identity on the cap. Thus multiplication gives
the following strict full boundary, with no discarded square:
\[
 \boxed{
 K_{\rm can}^\dagger|_{S^7}
   =\mathcal V_P\,
      \mathcal C(\mathcal Q_q^{\rm rev})^2\mathcal C^{-1}
        \mathcal R_J .
 }                                                        \tag{PD.4}\label{p1b:eq:natural-pd-4}
\]
The middle factor is generally nonidentity. Its coefficient is zero:
conjugation of closed based \(S^7\) maps induces ordinary conjugation
in the abelian group \(\pi_7SU(2)\), and
\(2[\mathcal Q_q]=0\). This statement supplies no formula for an
intrinsic null of that full factor, so it remains in the boundary.
The endpoint and joining coefficients are evaluated in \eqref{p1b:eq:natural-pg-14}
and \eqref{p1b:eq:natural-aj-4}--\eqref{p1b:eq:natural-aj-5}.

The connecting map of the original fibration identifies the parity
of the closed annular projection with the difference of its two
intrinsic boundary classes. Using \eqref{p1b:eq:natural-ov-4}, \eqref{p1b:eq:natural-pc-7}, \eqref{p1b:eq:natural-oh-11}, and then
the actual application \eqref{p1b:eq:natural-pd-3} gives the exact residual equation
\[
 \boxed{
 h_t\mu=\partial[\gamma_{\rm cart}]+[\mathcal V_P]
                       +[\mathcal R_J],\qquad
 \operatorname{Sq}^4u_5=h_tv_9
       \text{ on }C_{[K_D^\dagger e_3]},\qquad
 K_D^\dagger|_{S^7}=\Psi_D^\dagger .
 }                                                        \tag{PD.5}\label{p1b:eq:natural-pd-5}
\]
Here \(K_D^\dagger\) denotes the full paste of a filling
\(K_{\rm can}^\dagger:D^8\to SU(3)\) with the constructed
reverse comparisons and corrected original annulus.
The boundary in \eqref{p1b:eq:natural-pd-4} and every outer annular matrix are supplied.
The first equation is the retained complete three-map sum, evaluated
by \eqref{p1b:eq:natural-pf-3}--\eqref{p1b:eq:natural-pf-9}, \eqref{p1b:eq:natural-pg-14} and \eqref{p1b:eq:natural-aj-4}--\eqref{p1b:eq:natural-aj-5}. This equivalent
filling criterion does not replace the existing detector.

The complete original comparison is evaluated in
\eqref{p1b:eq:natural-aj-4}--\eqref{p1b:eq:natural-aj-7}:
\(h_t=w_t=1\), \(r_J=\zeta_{\rm nat}=0\), with every retained
extension test decided. This component is used with its full frames,
lower divisions and joining; its isolated value is not a replacement
for that complete comparison. The evaluated root \(\mu\) cancels the outer carry when passing from \(\zeta_{\rm nat}\) to \(h_t\); the other complete terms in \eqref{p1b:eq:natural-pd-5} are retained and separately evaluated.


\subsection{Reduce the actual characteristic meridian, retaining its whole closing loop}
\label{p1b:sec:natural-characteristic-meridian-reduction}

Start with the complete corrected boundary \eqref{p1b:eq:natural-pd-4}. The dependencies
executed here are the actual radial axis factor, its endpoint track in
the prescribed square null, the complete characteristic norm, and the
actual meridian value of the resulting closed loop. The original
root, central-minus closing, toggle, joining, sphere gauge, cap
embedding, frames and lower comparisons remain in the full input.

For \(z=r_Brv\), retain \(P=g(v)A_{1-r}g(v)^{-1}\) and put
\[
 P_\epsilon(z)=P(z)A_{\epsilon(1-r)}^{-1},\qquad
 P_\epsilon(0)=A_{1-\epsilon},\qquad
 P_\epsilon|_{\partial B}=I,\qquad
 P_1=gA_{1-r}g^{-1}A_{1-r}^{-1}=:\overline P.
                                                               \tag{PE.1}\label{p1b:eq:natural-pe-1}
\]
This is an actual intrinsic homotopy, not deletion of a factor by
its class. The original center estimate is uniform in \(v\);
multiplication by the displayed radial axis path preserves it.
The exchange equations remain \(DP_\epsilon(\sigma z)D=P_\epsilon\).
At \(\epsilon=1\), both radial poles are identity. On the actual
retained meridian \(v_*= (1,0,0,0,0,0)\), where \(g(v_*)=I\),
\(\overline P(r_Brv_*)=I\) for every \(r\).
Consequently \(\overline P\) has the strict reduced-suspension
factorization \(\overline P=\beta(r,g(v))\),
\(\beta(r,G)=GA_{1-r}G^{-1}A_{1-r}^{-1}\).
All three based faces of this specific \(S^1\wedge SU(2)\) map
are identity. This factorization concerns the cap map only; it
does not yet factor its entire closing loop.

Move the original square null with the actual endpoint path.
With \(\mathcal J_\nu\) exactly as in \eqref{p1b:eq:natural-or-1}, define
\[
 \mathcal J^{(\epsilon)}_\nu(z)=
 \begin{cases}
 P_{\epsilon-(1+\epsilon)\nu}(z)^{-2},
                              &0\le\nu\le\epsilon/(1+\epsilon),\\
 \mathcal J_{(1+\epsilon)\nu-\epsilon}(U_g(z)),
                              &\epsilon/(1+\epsilon)\le\nu\le1.
 \end{cases}
 \quad
 \mathcal J^{(\epsilon)}_0=P_\epsilon^{-2},\quad
 \mathcal J^{(\epsilon)}_1=I,\quad
 \mathcal J^{(0)}=\mathcal J.
                                                               \tag{PE.2}\label{p1b:eq:natural-pe-2}
\]
The seam is the literal \(P_0^{-2}\). The first interval shrinks
continuously at \(\epsilon=0\), by uniform continuity of \eqref{p1b:eq:natural-pe-1}.
Every spatial boundary value is identity. In particular the original
square null is retained, rather than replaced by a null chosen
solely from \(2[P]=0\).

Use the complete original \eqref{p1b:eq:natural-oy-2}--\eqref{p1b:eq:natural-oy-3} on \(P_\epsilon\),
including the actual \(d_h,\sigma_h\), Cartesian reflection pinch
and backtrack. On the same second-cut intervals put
\[
 \mathcal V^{(\epsilon)}(\pi(1+t),z)=
 \begin{cases}
 \mathcal J^{(\epsilon)}_{1-2t}(z),&0\le t\le1/2,\\
 Z_{P_\epsilon}(6t-3,z),&1/2\le t\le2/3,\\
 I,&2/3\le t\le1,
 \end{cases}
 \qquad \mathcal V^{(\epsilon)}|_{s\le\pi}=I.
                                                               \tag{PE.3}\label{p1b:eq:natural-pe-3}
\]
Its seam is exactly \(P_\epsilon^{-2}\), because the exchange
equation in \eqref{p1b:eq:natural-pe-1} holds. The circle faces and cap boundary are
identity for all \(\epsilon\). Thus \eqref{p1b:eq:natural-pe-3} is a specified intrinsic
homotopy of closed original \(S^7\) maps, with
\(\mathcal V^{(0)}=\mathcal V_P\).
It preserves its coefficient without assigning it.

The full loop \(\mathcal V^{(1)}\) need not be based on its meridian
or at its north pole: the endpoint track in \eqref{p1b:eq:natural-pe-2} and the actual
\(g(\sigma_hv_*)\) in \eqref{p1b:eq:natural-oy-2} still occur. Retain, measure and
remove these actual lower values by
\[
\begin{gathered}
 M(s,r)=\mathcal V^{(1)}(s,r_Brv_*),\qquad \mathcal M_h(s,z)=M(s,(1-h)|z|/r_B+h),\qquad \mathcal M_0=M(s,r), \\
 \mathcal M_1=I.
\end{gathered}
 \tag{PE.4}\label{p1b:eq:natural-pe-4}
\]
The displayed radial null is a whole-cap intrinsic homotopy.
It is circle based and cap-boundary based since
\(M(0,r)=M(2\pi,r)=M(s,1)=I\).
It is continuous at the cap center, independent of the absent
longitude, because it depends only on \(|z|\).
This is a justified lower comparison, not an assumed meridian null.
Define
\[
 \mathcal V^{\rm red}(s,z)=\mathcal V^{(1)}(s,z)M(s,r)^{-1}.
 \quad
 {}[\mathcal V^{\rm red}]=[\mathcal V_P],\qquad
 \mathcal V^{\rm red}|_{r=0,\ r=1,\ v=v_*,\ s=0}=I.
                                                               \tag{PE.5}\label{p1b:eq:natural-pe-5}
\]
The homotopy is \(\mathcal V^{(1)}\mathcal M_{1-h}^{-1}\).
All values of \(M\), including its closing, actual frames and
source path, are retained during their removal. The new map therefore
factors the actual reduced source \(S^1_s\wedge\Sigma(S^5,v_*)\).
This factorization alone does not identify a universal five-sphere
loop precomposed with \(\Sigma^2g\). The actual restricted loop
identification in \eqref{p1b:eq:natural-pg-1}--\eqref{p1b:eq:natural-pg-14} retains both multiplication
rectangles and the full seam correction.

Apply \eqref{p1b:eq:natural-pe-3}, followed by \eqref{p1b:eq:natural-pe-5}, immediately to the full boundary.
Replace only its characteristic factor, insert by the original \eqref{p1b:eq:natural-ox-1}
cap, and extend by identity on its complement. All new annular columns
are exactly \(e_3\), so the complete original \(\gamma_{\rm cart}\)
and every previously retained projected leg remain in the detector.
The exact new terminal and original residual comparison are
\[
 \boxed{
 K_{\rm red}^{\dagger}|_{S^7}
  =\mathcal V^{\rm red}\,
       \mathcal C(\mathcal Q_q^{\rm rev})^2\mathcal C^{-1}
       \mathcal R_J,\qquad
 h_t\mu=\partial[\gamma_{\rm cart}]+[\mathcal V^{\rm red}]
                                      +[\mathcal R_J].
 }                                                             \tag{PE.6}\label{p1b:eq:natural-pe-6}
\]
Both applications start at the preceding strict whole boundary;
neither uses a numerical joining substitute. All square factors,
the original \eqref{p1b:eq:natural-nd-8}, and the original bundle/base identifications
are present. The actual reduced characteristic loop is now supplied
with both its axis and meridian lower comparisons.

The complete original comparison is evaluated in
\eqref{p1b:eq:natural-aj-4}--\eqref{p1b:eq:natural-aj-7}:
\(h_t=w_t=1\), \(r_J=\zeta_{\rm nat}=0\), with every retained
extension test decided. This component is used with its full frames,
lower divisions and joining; its isolated value is not a replacement
for that complete comparison. The whole restricted characteristic loop, its multiplication rectangles, Cartesian pinch and seam correction are identified and evaluated in \eqref{p1b:eq:natural-pg-1}--\eqref{p1b:eq:natural-pg-16}. The factorization \eqref{p1b:eq:natural-pe-1} alone does not imply that value.


\subsection{The parity of the complete original Cartesian annulus}
\label{p1b:sec:natural-sewn-annular-parity}

All maps and frames in this argument are the retained actual ones.
Write \(q=Q_s^*H^3Q_s\), \(P=\Psi_g\),
\(\mathcal B_0=q^2\), \(\mathcal B_1=P\), and
\(f^\sigma(z)=Df(\sigma z)D\).
In particular \((q^2)^\sigma=q^{-2}\), \(P^\sigma=P\),
and every \(\mathcal B_\rho\) is identity on \(\partial B\).
By \eqref{p1b:eq:natural-ol-2}--\eqref{p1b:eq:natural-ol-4}, the actual reduced background is
\[
 Q(t,z)=g(v)A_{1-(1-t)r}g(v)^{-1},\qquad
 P(z)=g(v)A_{1-r}g(v)^{-1},\qquad z=r_Brv.
\]
It has bottom \(P\), top \(-I_2\), and side
\(gA_tg^{-1}\); both its central values and those of \(P\)
are the actual direction-independent \(-I_2\).

\begin{theorem}[The original complete annular input has even detector parity]
\label{p1b:theorem:natural-sewn-annular-parity}
For the original complete eight-leg map \eqref{p1b:eq:natural-ox-5},
\[
 \boxed{\partial[\gamma_{\rm cart}]=0\in\pi_7SU(2),\qquad
         \operatorname{Sq}^4u_5=0\text{ on }C_{\gamma_{\rm cart}}.}
                                                               \tag{PF.1}\label{p1b:eq:natural-pf-1}
\]
Consequently the exact original target equation \eqref{p1b:eq:natural-pe-6} reduces to
\[
 \boxed{h_t\mu=[\mathcal V^{\rm red}]+[\mathcal R_J].} \tag{PF.2}\label{p1b:eq:natural-pf-2}
\]
This evaluates the entire original annular contribution. It does not
assign either remaining term in \eqref{p1b:eq:natural-pf-2}.
\end{theorem}
\begin{proof}
The following comparison is of whole capped seven-sphere maps, rather
than of individual projected legs. For \(0\le\rho\le1\), put
\[
 f_\rho(z)=\mathcal B_\rho(z),\qquad
 k_\rho(t,z)=\mathcal B_{\rho+(1-\rho)\min(4t,1)}(z)
                         P(z)^{-1}Q(t,z).
                                                               \tag{PF.3}\label{p1b:eq:natural-pf-3}
\]
Thus \(k_0=k_1\) in the notation of \eqref{p1b:eq:natural-ol-4}, whereas
\(k_{\rho=1}=Q\). For every \(\rho\), the three literal
faces are
\[
 k_\rho(0)=f_\rho,\qquad k_\rho(1)=-I_2,\qquad
 k_\rho(t,r_Bv)=g(v)A_tg(v)^{-1}.
\]
Use the \emph{entire actual} norm contraction
\(Z_f\) of \eqref{p1b:eq:natural-oy-2}--\eqref{p1b:eq:natural-oy-3}, whose faces for any such cap map are
\[
 Z_f(0,z)=f(z)^{-1}(f^\sigma(z))^{-1},\qquad
 Z_f(1,z)=I,\qquad Z_f(\lambda,z)|_{\partial B}=I.
\]
It is jointly continuous in the bridge parameter by the same
uniform-continuity and zero-interval extension used in \eqref{p1b:eq:natural-oy-3}.
Define the following \emph{sewn} two-cut family:
\[
 U_\rho(s,z)=
 \begin{cases}
 A_t^{-1}k_\rho(t,z)f_\rho(z)^{-1},&s=\pi t,\\
 A_t^{-1}Z_{f_\rho}(t,z)k_\rho^\sigma(t,z),
                                      &s=\pi(t+1).
 \end{cases}                                                  \tag{PF.4}\label{p1b:eq:natural-pf-4}
\]
The first cut ends at \(f_\rho^{-1}\). The second starts
at \(f_\rho^{-1}(f_\rho^\sigma)^{-1}f_\rho^\sigma
=f_\rho^{-1}\). Their outer circle endpoints are both identity.
This proves the midpoint equation while retaining its full, generally
nonidentity, matrix. At \(z=r_Bv\), both cuts have the same actual
boundary word \(N_0(s,v)=b(s)^{-1}g(v)b(s)g(v)^{-1}\), because
\(g^\sigma=g\) in the retained cap axes. No inverse-exchange
identity is imposed on the intermediate bridge \(f_\rho\).

Apply the \emph{original} \eqref{p1b:eq:natural-ok-4} collar to every \(U_\rho\):
\[
 \widehat U_\rho=
 \begin{cases}
  N_{\chi(|z|)}(s,v)N_0(s,v)^{-1}U_\rho(s,z),&|z|>r_0,\\
  U_\rho(s,z),&|z|\le r_0.
 \end{cases}                                                  \tag{PF.5}\label{p1b:eq:natural-pf-5}
\]
The inner collar seam is exact, every spatial boundary is identity,
and the circle base remains identity. Thus \eqref{p1b:eq:natural-pf-5} is an actual intrinsic
homotopy of complete original cap seven-sphere maps. It is inserted by
the original \eqref{p1b:eq:natural-ox-1} source and extended by identity; the original
joining factor is untouched throughout.

At \(\rho=0\), its first cut is exactly the retained
\(W_1\) of \eqref{p1b:eq:natural-ol-5}. Its second cut differs by the actual closed
\(Z_{q^2}=\mathcal Q_{q^2}\) loop. The difference is conjugation
of this loop by \(A_t\) and by the original collar factor. These
conjugations introduce no extra top term: \(A_t\) can be contracted
in the conjugation while the loop's time endpoints are identity;
the collar conjugator is a continuous map on \(S^1\times B\),
and its radial contraction is allowed because the conjugated loop is
identity on \(\partial B\). Therefore
\[
 {}[\widehat U_0]=[\widehat W_1]+[\mathcal Q_{q^2}].       \tag{PF.6}\label{p1b:eq:natural-pf-6}
\]

At \(\rho=1\), the first cut is \(A_t^{-1}QP^{-1}\),
and the second is \(A_t^{-1}Z_P(t,z)Q(t,z)\). Define instead
the sewn characteristic word
\[
 U_{\rm char}(s,z)=
 \begin{cases}
  A_t^{-1}Q(t,z)P(z)^{-1},&s=\pi t,\\
  A_t^{-1}\mathcal J_t(U_g(z))Q(t,z),&s=\pi(t+1),
 \end{cases}                                                  \tag{PF.7}\label{p1b:eq:natural-pf-7}
\]
where \(\mathcal J_t\) is \emph{exactly} \eqref{p1b:eq:natural-or-1}, from
\(P^{-2}\) to identity. Its midpoint is \(P^{-1}\) on
both cuts; it has the same circle endpoints and spatial boundary word
as \eqref{p1b:eq:natural-pf-4}. Close it by exactly \eqref{p1b:eq:natural-pf-5}, and denote the closed map by
\(\widehat U_{\rm char}\).
The difference between the two second-cut paths is the actual loop
\(\overline{\mathcal J}*Z_P=\mathcal V_P\) of \eqref{p1b:eq:natural-oy-5}.
The same harmless conjugations just used in \eqref{p1b:eq:natural-pf-6} give
\[
 {}[\widehat U_1]=[\widehat U_{\rm char}]+[\mathcal V_P].  \tag{PF.8}\label{p1b:eq:natural-pf-8}
\]
Equations \eqref{p1b:eq:natural-pf-5}--\eqref{p1b:eq:natural-pf-8} consequently prove the whole-map equality
\[
 {}[\widehat W_1]+[\mathcal Q_{q^2}]
       =[\widehat U_{\rm char}]+[\mathcal V_P].         \tag{PF.9}\label{p1b:eq:natural-pf-9}
\]

We now evaluate the \emph{identified complete} characteristic word.
Every factor in \eqref{p1b:eq:natural-pf-7} and its original collar depends only on
\((s,r,g(v))\): \(U_g=(1-2r,2\sqrt{r(1-r)}Mg(v)e_1)\),
and \(N_\lambda=b_\lambda^{-1}gb_\lambda g^{-1}\).
At \(r=0\), \eqref{p1b:eq:natural-pf-7} is \(A_t^{-1}\) on the first cut and
\(-A_t^{-1}\) on the second, independently of \(g\).
At \(r=1\) its closed value is identity. Retain the actual meridian
\[
 M(s,r)=\widehat U_{\rm char}(s,r,I).
\]
It satisfies \(M(0,r)=M(2\pi,r)=M(s,1)=I\). Hence
\(M(s,h+(1-h)r)\) is a specified radial null, continuous at the
cap center. Right multiplication by the reverse of this null compares
the whole map with
\(\widehat U_{\rm char}(s,r,g)M(s,r)^{-1}\).
This new map is literally identity when \(s=0\), \(r=0\),
\(r=1\), or \(g=I\). It therefore factors through the specific
universal smash source
\[
 S^1_s\wedge S^1_r\wedge SU(2)=S^5.
\]
Write its universal map as \(F:S^5\to SU(2)\). On the original
source the complete normalized map is exactly the precomposition of
\(F\) with \(\Sigma^2g:S^7\to S^5\), up to the harmless
orientation conventions of the displayed smash coordinates.
Both \([F]\) and \([g]\) belong to
\(\pi_5SU(2)=C_2\{\eta_3\eta_4\}\). Thus either one is zero,
or their composite is
\[
 (\eta_3\eta_4)(\eta_5\eta_6)
       =6\nu'\eta_6=6\mu=0.
\]
This evaluates \([\widehat U_{\rm char}]=0\) for the actual
whole word, without selecting either universal coefficient.

Finally the retained bridge elimination \eqref{p1b:eq:natural-oy-5} gives
\([\mathcal T_{\rm int}]=[\mathcal V_P]-[\mathcal Q_{q^2}]\).
Together with \eqref{p1b:eq:natural-pf-9}, this says
\([\widehat W_1]=[\mathcal T_{\rm int}]\).
The complete original annulus \eqref{p1b:eq:natural-ov-3}--\eqref{p1b:eq:natural-ov-4}, equivalently \eqref{p1b:eq:natural-ox-5},
has the literal outer intrinsic end \(\operatorname{diag}(\psi',1)\)
and literal inner intrinsic end
\(\mathcal T_{\rm int}\mathcal R_J\).
The already constructed whole intrinsic comparison \eqref{p1b:eq:natural-om-7}, with
the original insertion and collar, gives
\([\psi']=[\widehat W_1]+[\mathcal R_J]\).
Thus the fibration connecting image of the \emph{original whole annulus}
is precisely
\[
 \partial[\gamma_{\rm cart}]
 =([\widehat W_1]+[\mathcal R_J])
   -([\mathcal T_{\rm int}]+[\mathcal R_J])=0.
\]
The unchanged joining cancels only in this difference of two whole
end classes, proving \(\partial[\gamma_{\rm cart}]=0\).
The retained connecting map \(\pi_8S^5=C_{24}\to\pi_7SU(2)=C_2\)
is coefficient parity, the same parity detected by \(\operatorname{Sq}^4\).
This proves \eqref{p1b:eq:natural-pf-1}, and substitution into \eqref{p1b:eq:natural-pe-6} proves \eqref{p1b:eq:natural-pf-2}.
\end{proof}

\noindent\emph{Scope of the annular comparison.}
The argument evaluates the full annular parity with its original
joining and moving right inverse. It does not declare an isolated
annular leg null or assign $h_t$ from a comparison choice.





\subsection{The actual characteristic norm and its restricted secondary composition}
\label{p1b:sec:natural-characteristic-norm-compatibility}

Retain PE.1--PE.6, the actual original square null, the original cap,
the source swap and Cartesian closing, and the unchanged joining.
Write $X=B/\partial B=S^6$, $r=|z|/r_B$, $v=z/|z|$,
and $\beta(r,G)=GA_{1-r}G^{-1}A_{1-r}^{-1}$.
This reduced map is based at both radial poles and at $G=I$,
so defines the actual characteristic $\alpha:S^4\to S^3$.
Its class is $\eta_3$ in the retained cap axes.
The actual cap is $f(z)=\beta(r,g(v))$, with
$[g]=\eta_3\eta_4$ and $[f]=6\nu'$.

\paragraph{Evaluate the actual first-track base.}
Let $v_*=(1,0,0,0,0,0)$ and $v_p=(0,0,0,1,0,0)$.
The \eqref{p1b:eq:natural-op-4} source path has
$a=(v_x-v_p)/\sqrt2$, $b=(v_x+v_p)/\sqrt2$, rotates
only $(a_1,a_2)$, and changes the sign of $a_0$.
At $v_*$ those rotated coordinates are zero, so
\[
 \sigma_hv_*=v_p\quad(0\le h\le1),\qquad
 g(v_*)=g(v_p)=I.
 \tag{PG.1}\label{p1b:eq:natural-pg-1}
\]
The second assertion is an actual frame calculation. At either
direction one cap quaternion is $1$. The original section and
every shifted section are on the smash wedge, hence are $e_1$.
Each successive special-unitary transporter is identity, the cap
frame is $Q_0$, and the exterior line frame is
$[e_2,e_3,e_1]=Q_0$. Thus their clutch is identity exactly.
Consequently
\[
\begin{gathered}
 H_h(v)=d_hg(\sigma_hv)d_h^{-1},\qquad H_h(v_*)=I,\qquad H_0=g,\qquad H_1=g\circ R, \\
 Z_f(h/3)=\beta(r,g)^{-1}\beta(r,H_h)^{-1},
\end{gathered}
 \tag{PG.2}\label{p1b:eq:natural-pg-2}
\]
where $R$ is the original final single reflection.
The commutation of $d_h$ with the retained radial axis proves
the last identity pointwise. In particular this whole first
track is already based; its PE meridian correction is identity.

\paragraph{The square leg is an actual universal based left null.}
For an arbitrary $G\in SU(2)$ put
\[
 P_\epsilon(r,G)=GA_{1-r}G^{-1}A_{\epsilon(1-r)}^{-1},
 \quad b(r,G)=(1-2r,2\sqrt{r(1-r)}\,MGe_1),
\]
and define the entire moved null by the exact \eqref{p1b:eq:natural-pe-2} formula
\[
 N_\nu(r,G)=
 \begin{cases}
 P_{1-2\nu}(r,G)^{-2},&0\le\nu\le1/2,\\
 \mathcal J_{2\nu-1}(b(r,G)),&1/2\le\nu\le1.
 \end{cases}
 \qquad
 N^{\rm b}_\nu(r,G)=N_\nu(r,G)N_\nu(r,I)^{-1}.
 \tag{PG.3}\label{p1b:eq:natural-pg-3}
\]
This is the original square null, with its actual endpoint path,
not a null inferred from $2[f]=0$. Its seam is the original
$P_0^{-2}$. At $r=0$ its raw value is independent of $G$,
so the displayed quotient is identity; at $r=1$ all factors
are identity. The quotient is identity at $G=I$ as well.
Its two time endpoints are
$N^{\rm b}_0=\beta(r,G)^{-2}$ and $N^{\rm b}_1=I$.
Thus it is literally a based universal square null of $\alpha$.
On the actual cap, \eqref{p1b:eq:natural-pg-3} is exactly the square portion of
\eqref{p1b:eq:natural-pe-5}, because the actual meridian has $g(v_*)=I$.
No portion of the nonuniversal Cartesian closing is absorbed
into this universal identification.

Any change of this universal left null is a map
$S^5\to S^3$. Its actual source substitution is
$\Sigma^2g$, and hence its possible contribution is
\[
 \pi_5S^3\circ[\Sigma^2g]
 =\{0,\eta_3\eta_4\eta_5\eta_6\}=\{0\}.
 \tag{PG.4}\label{p1b:eq:natural-pg-4}
\]
This uses the already proved \eqref{p1b:eq:natural-ow-9} relation $6\mu=0$.

\paragraph{A relative kernel for a change of the whole longitude track.}
If $H^0,H^1$ are based paths having the literal endpoints
in \eqref{p1b:eq:natural-pg-2}, put $L_h=H^0_h(H^1_h)^{-1}$.
It is a based $S^1_h\wedge S^5\to SU(2)$ map.
The exact commutator identity is
\[
 \beta(r,LH)=L\beta(r,H)L^{-1}\beta(r,L).
 \tag{PG.5}\label{p1b:eq:natural-pg-5}
\]
It identifies the first-track difference, apart from the fixed
conjugation by $\beta(r,g)$, with $\beta(r,L)^{-1}$ times
a complete commutator of $L$ and $\beta(r,H^1)$.
That commutator has a relative intrinsic null by separating
their actual longitude supports. In the stereographic
$x\in\mathbb R^5$ coordinate from $v_*$ use
\[
 c_\tau^\pm(v)=
 \begin{cases}
 \operatorname{st}^{-1}
       ((x\mp2\tau e_1)/(1-\tau|x\mp2\tau e_1|^2)),
       &\tau|x\mp2\tau e_1|^2<1,\\
 v_*,&\tau|x\mp2\tau e_1|^2\ge1,
 \end{cases}
 \quad c_\tau^\pm(v_*)=v_*.
 \tag{PG.6}\label{p1b:eq:natural-pg-6}
\]
This is \eqref{p1b:eq:natural-ou-7} in dimension five. Both maps start at identity
and end with disjoint open unit five-ball supports centered at
$\pm2e_1$. Its homogeneous inverse stereographic formula
proves continuity on the collapse seam and at the based point.
Precompose only the two operands in the closed commutator.
It remains identity at both time seams because $L_0=L_1=I$;
it remains based on every spatial and radial face. At the end
one operand is identity at every point. This gives a null of
the entire commutator without changing either norm endpoint.
The fixed-conjugator correction is another closed commutator,
whose null is given by the same two-support construction.

The remaining difference is therefore
$\eta_3E[L]$. The relevant source relation is
\[
 {}[\iota_4,\iota_4]=\pm(2\nu_4-E\nu'),\qquad
 \eta_3\nu_4=\mu,\qquad 2\mu=0,
 \quad\Longrightarrow\quad \eta_3E\pi_6S^3=0.
 \tag{PG.7}\label{p1b:eq:natural-pg-7}
\]
Naturality sends the Whitehead square to $[\eta_3,\eta_3]$,
which vanishes in the topological group $S^3$.
The relation $E\nu'=2\nu_4$ would be false and is not used.
This proves that replacing the actual based longitude track
with another having its exact endpoints adds zero, with every
relative comparison carried.

\paragraph{Universal inverse compatibility.}
The actual norm uses inverse characteristic values. It must not
identify these silently with the characteristic of inverse longitudes.
Use the exact identity
\[
 \beta(r,G^{-1})=G^{-1}\beta(r,G)^{-1}G,
 \quad
 \beta(r,G)^{-1}\beta(r,G^{-1})^{-1}
       =[G,\beta(r,G^{-1})].
 \tag{PG.8}\label{p1b:eq:natural-pg-8}
\]
The last commutator has a specified universal null: independently
compress its two $S^3$ inputs by \eqref{p1b:eq:natural-ou-7} in dimension three.
Its two radial faces stay identity because $\beta$ is identity
there. Its final operands have disjoint support. Multiplication
by $\beta(r,G^{-1})$ supplies a relative comparison
$\beta(r,G)^{-1}\to\beta(r,G^{-1})$.
Apply it to each complete norm operand and transport \eqref{p1b:eq:natural-pg-3}
along its actual changed initial endpoint. This is a homotopy
of the whole closed map. The transported left null is still a
universal null; its choice difference has the zero pullback \eqref{p1b:eq:natural-pg-4}.
The intermediate compression need not preserve the original
$D$ exchange, so the endpoint transport is the following entire
universal path, rather than an assumed square. Let
$a_\kappa(r,G)$ be the specified comparison of \eqref{p1b:eq:natural-pg-8}, and put
\[
 F_\kappa(r,G)=a_\kappa(r,G)D
                       a_\kappa(r,DGD)D.
\]
Then $F_0=\beta(r,G)^{-2}$ and
$F_1=\beta(r,G^{-1})^2$. For an intermediate $\kappa$ use
the actual left path $I\to F_0$, followed by
$F_t$, $0\le t\le\kappa$, and use $a_\kappa$ in both
complete norm operands and in the full closing word. Their
initial seam is literally $F_\kappa$. This transport is based
on both radial faces and at $G=I$ for every $\kappa$.
Thus the complete circle is closed even if the intermediate
inverse comparison itself is not exchange-equivariant. At its
terminal endpoint the inverse characteristic recovers the
original exchange equation, and the universal left leg is
indeed a square null again.
Hence one can work with $g'=g^{-1}$ and positive characteristic
operands. Here $[g']=\eta_3\eta_4$ and
$[f']= [\beta(r,g')]=6\nu'$.

\paragraph{Match the complete Cartesian and radial closing starts.}
Choose a based longitude reflection $R'$ fixing $v_*$ and a
linear orientation-reversing source path $R_t$ from $R$ to $R'$.
It can be chosen to vary only the original $(x_1,p_1)$ plane,
so $R_tv_*$ joins $v_p$ to $v_*$ in that plane.
Explicitly on that plane take
$R_t=\left(\begin{smallmatrix}\cos2\theta_t&\sin2\theta_t\\
 \sin2\theta_t&-\cos2\theta_t\end{smallmatrix}\right)$,
$\theta_t=(1-t)\pi/4$, and retain the identity on the other four
coordinates. Its initial matrix is the actual swap of $x_1,p_1$
and its final matrix fixes $x_1$ and changes the sign of $p_1$.
Put $b_t=g'(R_tv_*)$ and
$K_t(v)=g'(R_tv)b_t^{-1}$. Both endpoints of $b_t$ are identity,
and $K$ has the exact endpoints $g'R,g'R'$.

This lower base loop has a specified null. In that plane the
two source quaternions lie on the actual quaternionic $i$ circle.
The balanced smash has only its first $u_a,u_b$ components;
the exotic section is the real vector $(h_3,u_{a1},u_{b1})$.
Every shifted section and transporter is real, and both ordered
frames are real. Hence $b_t$ is in the real $SO(2)\subset SU(2)$
circle. In quaternion coordinates choose the orthogonal unit $i$,
put $F_\kappa(q)=((1-\kappa)q+\kappa i)/
\sqrt{(1-\kappa)^2+\kappa^2}$, and divide
$F_\kappa(b_t)$ by $F_\kappa(I)$ on the right.
This contracts $b_t$ to identity, fixing both time endpoints,
with a denominator bounded below by $1/\sqrt2$.
Denote this actual contraction by $b_t^\kappa$.
The whole characteristic-basing defect has the relative null
\[
 \delta_\kappa=
 \beta(r,K_t)
 \left(\beta(r,K_tb_t^\kappa)
                 \beta(r,b_t^\kappa)^{-1}\right)^{-1}.
 \tag{PG.9}\label{p1b:eq:natural-pg-9}
\]
It starts at the literal defect between a based longitude
and its raw meridian division and ends at identity.
It is identity on the longitude base, both time endpoints,
and both radial faces. Thus the lower source path is retained
and evaluated, not discarded by a basedness assumption.

The new radial closing first follows $R_t$, with its actual
meridian division, and then uses the standard based longitude
reflection pinch/backtrack for $R'$. Its initial target value is
the \emph{same} $f'(z)f'(Rz)$ as the transformed Cartesian tail.
Its final value is identity. The comparison of \eqref{p1b:eq:natural-pg-9} identifies
its based-$K$ and raw-word forms without changing a coefficient.

For the closing-tail comparison alone, lift every raw spatial
letter to the free topological group $F(X)$, $X=B/\partial B$.
Normalize each \emph{whole} lifted closing by its own actual
meridian word, exactly as \eqref{p1b:eq:natural-pe-5}. Both lifted closings then have
the same literal initial word
\[
 j(z)j(Rz)\bigl(j(r_Brv_*)j(r_BrRv_*)\bigr)^{-1}
 \tag{PG.10}\label{p1b:eq:natural-pg-10}
\]
and terminal word identity. At the cap center each raw word
is its meridian word and cancels literally; the same is true
on the actual meridian. At the cap boundary every letter is
the based letter. At the endpoint of the longitude backtrack
each raw letter is on the same meridian, so the terminal
division cancels literally as well. These are the needed
lower word repairs on the original $X$, not a collapse of the
center that would change $X$ to $S^6\vee S^1$.
Their full closed difference is therefore a based
$S^7\to F(S^6)\simeq\Omega S^7$. Evaluation under $f'$ gives
\[
 \pi_7F(S^6)=\pi_8S^7=C_2
   \longmapsto [f'\eta_6]=6\nu'\eta_6=0.
 \tag{PG.11}\label{p1b:eq:natural-pg-11}
\]
Only the two reflection closings with the common start \eqref{p1b:eq:natural-pg-10}
are compared here. The universal characteristic square null
has not been lifted to $F(S^6)$ or absorbed into this argument.

\paragraph{Retain the multiplication square and its source endpoint.}
After this closing replacement the complete source norm is a
based path from $g'^2$ to identity. Its straightening and pinch
stages have two based operands $A_s,B_s:S^5\to SU(2)$;
at the pinch seam their pointwise disjoint supports allow them
to be written as one folded letter. The remaining backtrack
contracts that folded letter. There is no continuous pair of
individual based operands with one operand becoming identity:
each original operand has the nontrivial clutch class. On the
paired stages retain the entire comparison
\[
 D_\kappa(r,s,v)=
 \beta(r,A_s(c_\kappa^+v))
 \beta(r,B_s(c_\kappa^-v))
 \beta(r,A_s(c_\kappa^+v)B_s(c_\kappa^-v))^{-1}.
 \tag{PG.12}\label{p1b:eq:natural-pg-12}
\]
At $\kappa=0$ it is the literal multiplication defect and
at $\kappa=1$ it is identity. At the pinch seam its
$\kappa=0$ value is identity because the original operands
have disjoint hemisphere supports. Its intermediate
$\kappa$ values need not be identity: the whole seam edge
is retained and filled below before pasting the single-letter
backtrack. No null of either individual clutch operand is used.
At the initial time its track is the actual comparison
$\beta(r,g')^2\leftrightarrow\beta(r,g'^2)$.
That initial track is pasted into the left square null;
it is not deleted or assumed constant. All cap and longitude
base faces are fixed.

\paragraph{Fill the actual pinch-seam edge.}
Write $d(r,U,V)=\beta(r,U)\beta(r,V)\beta(r,UV)^{-1}$.
It is identity on either coordinate wedge face and on both
radial faces. Thus it strictly factors as a map
$\overline d:S^1_r\wedge(S^3\wedge S^3)=S^7\to S^3$.
At the original pinch seam the pair has the form
$\Phi_0(v)=(g'(p_Av),g'(p_Bv))$, where $p_A,p_B$ are the
literal two hemisphere pinch maps, of degrees $1,-1$.
The $\kappa=1$ compressed pair is
$\Phi_1(v)=(g'(p_Ac_1^+v),g'(p_Bc_1^-v))$.
Both pairs are pointwise wedge-valued. Before evaluating $g'$,
they are maps $T_0,T_1:S^5\to S^5\vee S^5$, with the
same two degrees $(1,-1)$. The first mixed Whitehead term
in that source wedge has dimension nine; hence its fifth
group is $\mathbb Z\oplus\mathbb Z$, and these exact maps
admit a based source-wedge comparison. Evaluate it by
$g'\vee g'$ and denote the resulting wedge homotopy by $W$.
This is a comparison of the original hemisphere and compressed
ball support configurations, not a class cancellation of either
operand.

Concatenate the whole actual pair-compression path $\Phi_\kappa$
with $W$ backwards. It gives
$P:S^1\times S^5\to S^3\times S^3$ with initial and final
slice $P_0=\Phi_0$. That slice is not constant, so the pair
map is not silently treated as an $S^6$ map. Translate it
in the product group: $L=P_0^{-1}P$. Then $L$ is literally
identity on both time ends and on the longitude base, and
does define $S^6\to S^3\times S^3$.

The closed defect $d(r,P_0L)$ has an actual relative
comparison with $d(r,L)$. Independently precompose $P_0$
and $L$ by the two longitude compressions \eqref{p1b:eq:natural-pg-6}. At a
time end $L=I$, and the other pair is wedge-valued, so
$d(r,P_0)=I$ throughout this operation. At disjoint support
the result is literally $d(r,L(c_1^-v))$: where $L=I$
the wedge pair has zero defect, and where $P_0=I$ only
$L$ remains. Reverse the compression of $L$ alone.
All radial, longitude and time faces remain fixed. Finally,
for the smash quotient $q:S^3\times S^3\to S^6$,
\[
 q_*\pi_6(S^3\times S^3)
   =q_*(C_{12}\oplus C_{12})=0\subset\mathbb Z,
 \qquad [qL]=0.
 \tag{PG.12a}\label{p1b:eq:natural-pg-12a}
\]
Evaluating $\overline d$ after suspending the radial coordinate
therefore supplies a based null of this entire seam edge.
The part of the pair loop lying in $W$ has defect identity.
The null hence fills precisely the originally omitted
$\kappa$ edge of the paired square. Paste it before the
single folded backtrack; the backtrack has one letter and
introduces no further multiplication defect.
The source closing at the compressed endpoint is retained
as well: after the paired stages append the actual product
of the wedge comparison $W$, from $\Phi_1$ to $\Phi_0$,
and then the original folded backtrack. Along $W$ one
operand is identity at each point, so its characteristic
product is literally the characteristic of its source
product. Thus the source and target seams both paste;
the backtrack is not attached to a different initial map.

\paragraph{Identify the initial edge with its universal counterpart.}
The initial track of \eqref{p1b:eq:natural-pg-12} must also be compared, rather
than declared universal. Its pair-compression path is
\[
 (g'(c_\kappa^+v),g'(c_\kappa^-v)),
\]
whereas a universal null pulled back from the characteristic
longitude uses
\[
 (c_\kappa^{3,+}(g'(v)),c_\kappa^{3,-}(g'(v))).
\]
Here $c^{3,\pm}$ is the same based compression formula on
the universal $S^3$ input. Both paths begin at the literal
diagonal pair $(g',g')$ and end in the coordinate wedge.
Their exact endpoint wedge maps admit a based wedge comparison.
Indeed $[g']=\eta_3\eta_4=\Sigma(\eta_2\eta_3)$, so choose
a based comparison with that suspension representative. Compare
each \eqref{p1b:eq:natural-ou-7} source pinch with the ordinary suspension pinch in
$S^5\vee S^5$ and $S^3\vee S^3$ respectively; their two
projection degrees are $(1,1)$. Suspension pinch naturality
then gives the endpoint wedge comparison, and transport it
back through the actual $g'$ comparison. Its whole values
remain in the wedge, so its defect is identity.

Closing the two pair paths with this wedge comparison gives
$P:S^1\times S^5\to S^3\times S^3$, now with nonwedge
initial slice $P_0=(g',g')$. No claim that this pair map
factors through $S^6$ is made. The map $qP$ nevertheless
has a based null on
$Y=(S^1\times S^5)/(S^1\times\{v_*\})$.
The target is five-connected and $Y$ has dimension six,
so the sole obstruction is its pulled-back degree-six class.
If $a_3,b_3$ are the two product generators and $u_6$ is
the smash generator, then
\[
 q^*u_6=a_3b_3,\qquad H^3(Y;\mathbb Z)=0,
 \qquad (qP)^*u_6=0.
 \tag{PG.12b}\label{p1b:eq:natural-pg-12b}
\]
Choose the resulting whole based null $T_\tau$ of $qP$.
Its initial-time slice need not remain constant or fixed;
retain it by using the correlated target normalization
\[
 \overline d(r,T_\tau(t,v))
       \overline d(r,T_\tau(0,v))^{-1}.
 \tag{PG.12c}\label{p1b:eq:natural-pg-12c}
\]
This is identity on the collapsed time face for every
$\tau$, and on the longitude base and radial faces.
At its initial value it is the complete difference of the
two initial defect-null tracks, normalized by their common
actual diagonal slice. At its terminal value it is identity.
Thus the initial transport is identified with the universal
transport by a supplied zero correction, even though its
initial slice is nonconstant. The actual initial lower
track moves with the whole null and is not dropped.

Together \eqref{p1b:eq:natural-pg-12a}--\eqref{p1b:eq:natural-pg-12c} supply the two missing rectangles
of the multiplication square. Its initial endpoint comparison
is retained, its entire pinch-seam edge is filled, and the
single-letter folded backtrack pastes after that filling.

There is a universal typing comparison for that initial
transport. The universal defect null compares
$\beta(r,G)^2$ with $\beta(r,G^2)$, the latter being
$\alpha\circ\Sigma p_2$ for $p_2(G)=G^2$ on $S^3$.
The power map has degree two because the degree-three
class is primitive under group multiplication.
Both $\Sigma p_2:S^4\to S^4$ and the radial suspension
fold $2\iota_4$ therefore have degree two and admit a
based universal source comparison. Compose that whole
comparison with $\alpha$ and retain it after the universal
initial transport. Two choices differ by a map $S^5\to S^3$,
whose pullback on $\Sigma^2g'$ is zero by \eqref{p1b:eq:natural-pg-4}.
Consequently this retained universal transport is identified
with the independently written radial left-null construction
below, with their whole comparison and its zero ambiguity.
The $\beta(G^2)$ endpoint is not identified silently with
the radial degree-two endpoint.

At the compressed endpoint the entire source right path
starts at the two-copy co-H sum, follows the compressed
paired stages, the just-retained wedge-product seam path,
and the original folded backtrack, and ends at identity.
Prepending the product comparison from $g'^2$ to that
co-H sum gives a based null of $g'^2$. Its difference
from another such null is in
$\pi_6S^3$, whose suspended characteristic image is zero by
\eqref{p1b:eq:natural-pg-7}. Therefore it can be compared with the actual \eqref{p1b:eq:natural-ot-6}
null with its values inverted: if \eqref{p1b:eq:natural-ot-6} is $S_s:g^2\to I$,
put $S'_s(v)=S_s(v)^{-1}$, so $S'_s:g'^2\to I$.
This is pointwise inversion of its values, not reversal of
its time direction.

Let $m_\kappa:g'^2\to
g'(c_1^+)g'(c_1^-)$ be the displayed source product comparison.
The required null from that co-H sum to identity is
$\overline m*S'$, namely the product comparison traversed
backwards, followed by $S'$ in its original forward time.
Suspend this \emph{whole} source path. Thus the right null in
the comparison is an actual suspended clutch-square null with
the correct two endpoints. Changing it contributes only
$\eta_3E\pi_6S^3=0$, not $\eta_3\pi_7S^4$.

There is one source-coordinate comparison in this suspension.
Its longitude pinch $\Delta_L$ and the radial suspension
pinch $\Delta_R$ are specified based maps
$S^6\to S^6\vee S^6$, each of projection degrees $(1,1)$.
The wedge is five-connected and its sixth group is
$\mathbb Z\oplus\mathbb Z$, so these two source maps admit
a based comparison. Two comparison choices differ by
\[
 \pi_7(S^6\vee S^6)=C_2\oplus C_2;
 \quad (f'\vee f')_*\pi_7(S^6\vee S^6)=0.
 \tag{PG.13}\label{p1b:eq:natural-pg-13}
\]
The first mixed Whitehead term occurs in dimension eleven,
so there is no omitted mixed contribution. Each summand is
sent to $f'\eta_6=0$. Crucially this comparison is chosen
on the original $S^6$ source, before applying $\Sigma g'$.
It does not choose an arbitrary $S^4$ null with an uncontrolled
Hopf summand $\nu_4$. The radial pinch commutes literally
with the suspension $\Sigma g'$, giving the ordinary Toda
defining diagram with this specified restricted right null.

The universal pointwise-square to radial-pinch comparison on the
left is also literal. For the characteristic $\alpha(r,G)$ use
\[
 \alpha((1-u)r+u\min(2r,1),G)\,
 \alpha((1-u)r+u\max(2r-1,0),G),\quad 0\le u\le1.
\]
At $u=0$ this is $\alpha(r,G)^2$; at $u=1$ one factor
is a radial based face and this is precisely $\alpha\circ
2\iota_4$ in the suspension pinch model. Both radial poles
and $G=I$ are fixed throughout. Concatenate its inverse with
the actual universal square null \eqref{p1b:eq:natural-pg-3}, transported through
\eqref{p1b:eq:natural-pg-8}. This supplies the universal left null for the ordinary
Toda diagram. Its substitution by $\Sigma g'$ is strict,
so no additional $S^4$ comparison null is introduced here.

\paragraph{Evaluate only after the actual diagram is identified.}
The preceding squares identify the entire normalized loop with
the secondary composition having left map $\eta_3$, middle
map $2\iota_4$, right map $\Sigma g'=\eta_4\eta_5$,
and right null required to be a suspension of the clutch-square
null. Its left choice ambiguity is zero by \eqref{p1b:eq:natural-pg-4}, its right
choice ambiguity is zero by \eqref{p1b:eq:natural-pg-7}, and its source-pinch and
closing-word ambiguities are zero by \eqref{p1b:eq:natural-pg-11}--\eqref{p1b:eq:natural-pg-13}.

Now use the supplied actual class $[g']=\eta_3\eta_4$.
Choosing a based representative comparison and a square null
factored through the inner $\eta_3$ gives, with the same
restricted diagram,
\[
 {}[\mathcal V^{\rm red}]
   =\langle\eta_3,2\iota_4,\eta_4\rangle\eta_6
   =(\text{odd}\,\nu')\eta_6=\mu.
 \tag{PG.14}\label{p1b:eq:natural-pg-14}
\]
The primitive bracket has indeterminacy $6\nu'$; composing
with $\eta_6$ annihilates it. Its classical value is a generator
of the two-primary $C_4$ summand. In the retained integral
$C_{12}$ normalization its lifts are $3\nu'$ or $9\nu'$;
both have odd coefficient and both compose to $\mu$.
No generator of the full $C_{12}$ group is asserted to lie
in that primitive bracket. The unrestricted bracket with right map
$\eta_4\eta_5$ has the larger indeterminacy
$\eta_3\pi_7S^4=C_2$ and is deliberately not treated as a
singleton. The zero kernels above are what justify the
restricted calculation for this actual map.


\paragraph{Bridge-controlled parity, with no separately selected root class.}
\begin{lemma}[Bridge-controlled parity of the actual root-carry word]
\label{p1c:lem:v263-bridge-controlled-carry}
Use the based cap-smash coordinates of \eqref{p1b:eq:natural-pc-1}--\eqref{p1b:eq:natural-pc-4}.
Suppose $q,P:S^6\to SU(2)$ have a specified boundary-preserving
bridge $B_\rho:q^2\simeq P$ as in NZ.20 and
$[P]=6\epsilon\nu'$ in $\pi_6(SU(2))\cong\mathbb Z_{12}$, where
$\epsilon\in\mathbb Z_2$.  Assume the complete signed-word identification
\eqref{p1b:eq:natural-pb-7}, \eqref{p1b:eq:natural-pc-4}--\eqref{p1b:eq:natural-pc-7}
is applied to this same $q$, including its meridian, source orientation,
and both moving rectangles.  Then
\[
  [\mathcal Q_{q^2}]=\epsilon\mu
       \quad\text{in }\pi_7(SU(2))\cong\mathbb Z_2\{\mu\}.
\]
In particular, the root-carry class depends on the bridge class
$[P]$, not on a choice between the two possible square-root classes.
\end{lemma}
\begin{proof}
Write $[q]=k\nu'$ with $k\in\mathbb Z_{12}$.  The supplied
bridge implies $2k\equiv6\epsilon\pmod{12}$, hence
\[
  \begin{array}{c|c|c}
    \epsilon&k\pmod{12}&k\pmod2\\\hline
    0&0,6&0\\
    1&3,9&1
  \end{array}.
\]
The full word calculation \eqref{p1b:eq:natural-pc-4}--\eqref{p1b:eq:natural-pc-7},
including the precise integer winding $-1$ in
\eqref{p1b:eq:natural-pc-6}, gives
$[\mathcal Q_{q^2}]=-[q\eta_6]$; the extra complete
square $\mathcal Q_q^2$ in \eqref{p1b:eq:natural-pb-7}
contributes twice a class in the order-two target and hence zero.
Since $\nu'\eta_6=\mu$ has order two, this is
$-k\mu=(k\bmod2)\mu=\epsilon\mu$.
Equivalently,
$\ker(2:\mathbb Z_{12}\to\mathbb Z_{12})=\{0,6\}$
is contained in the kernel of parity.  This proves root-class
independence for the stated, fully typed cap construction.
\end{proof}

\begin{corollary}[Old-cap cancellation with the original characteristic comparison]
\label{p1c:cor:v263-old-cap-cancellation}
For the actual nontrivial reduced clutch
$[g]=\eta_3\eta_4$, the retained bridge NZ.20 gives
$[P]=\eta_3\eta_4\eta_5=6\nu'$, so the preceding lemma fixes
$[\mathcal Q_{q^2}]=\mu$ without invoking the particular
identity $[q]=3\nu'$.  Keeping the \emph{source-specific}
meridian removal \eqref{p1b:eq:natural-pe-5} and the restricted
characteristic-to-Toda comparison
\eqref{p1b:eq:natural-pg-8}--\eqref{p1b:eq:natural-pg-14},
we have $[\mathcal V_P]=[\mathcal V^{\rm red}]=\mu$.
The actual whole bridge square \eqref{p1b:eq:natural-oy-5}
therefore yields
\[
 \boxed{[\mathcal T_{\rm int}]
   =[\mathcal V_P]-[\mathcal Q_{q^2}]=\mu-\mu=0.}
\]
This is a consequence of the previously specified maps and
restricted nulls; it does not replace the full frame, the original
cap normalization, or either half of the joined boundary.
\end{corollary}

\noindent\emph{Nonconverse and logical scope.}
A vanishing \emph{formal} top difference does not imply the
existence of the intrinsic bridge.  For example, with
$[P]=6\nu'$ and a hypothetical $[q]=\nu'$, both mod-two
seven-dimensional contributions could be $\mu$, although
$2[q]=2\nu'\ne6\nu'=[P]$ prohibits any based bridge
$q^2\simeq P$.  Thus cap triviality is strictly weaker than
bridge compatibility.  The cancellation above uses the actual
NZ.20 and PG constructions; the displayed parity arithmetic
alone does not independently certify the complete PG pinch-seam
transport or the original OK/OL/PF/CI frame comparisons.

\paragraph{Primary source and normalization.}
Brayton Gray, \emph{Periodicity, compositions and EHP sequences},
Contemporary Mathematics 617 (2014), Section 2, p.~191,
doi:10.1090/conm/617/12284, gives the Moore-map construction
of the two-primary $\nu'$ and its composition with $\eta$.
Zhongjian Zhu and Tian Jin, \emph{The relative James construction and its
application to homotopy groups}, arXiv:2402.07072v2 (10 August 2024),
Lemma 4.7, printed p.~20, record the two-local bracket as
$\{\nu',-\nu'\}$. The retained full-group normalization is
the one used in RT.23 and \eqref{p1b:eq:natural-pc-7}. These primary computations
are applied only after the entire actual restricted input,
both multiplication rectangles and their ambiguity kernels
have been identified above.

Every comparison is intrinsic on the original cap and fixes its
boundary, so it can be inserted in \eqref{p1b:eq:natural-pe-6} and extended by identity.
Its complete projected column is unchanged. The proved residual
equation is
\[
 h_t\mu=\partial[\gamma_{\rm cart}]+\mu+[\mathcal R_J].
 \tag{PG.15}\label{p1b:eq:natural-pg-15}
\]
The preceding PF calculation evaluates
$\partial[\gamma_{\rm cart}]=0$. Writing the remaining actual
joining class as $[\mathcal R_J]=r_J\mu$ therefore gives
\[
 \boxed{h_t=1+r_J,\qquad \zeta_{\rm nat}=r_J,\qquad w_t=1.}
 \tag{PG.16}\label{p1b:eq:natural-pg-16}
\]
The complete original comparison is evaluated in
\eqref{p1b:eq:natural-aj-4}--\eqref{p1b:eq:natural-aj-7}:
\(h_t=w_t=1\), \(r_J=\zeta_{\rm nat}=0\), with every retained
extension test decided. This component is used with its full frames,
lower divisions and joining; its isolated value is not a replacement
for that complete comparison.


\subsection{The actual projective seam has a primitive quaternion commutator}
\label{p1b:sec:natural-joining-projective-primitive}

Retain the actual maps, intrinsic quaternion block, parameter directions,
and inverse-image caps of \eqref{p1b:eq:natural-ob-1}--\eqref{p1b:eq:natural-ob-15}.  This calculation evaluates
the particular primitive \(K\) in \eqref{p1b:eq:natural-ob-13}.  It does not replace the
complete joining remainder by that primitive.

\subsubsection{The axis-change map has unit absolute degree}

Put \(c=\cos(\pi\tau/2)\), \(h=\sin(\pi\tau/2)\), and
\(b=b_\tau(s,0)\).  Quaternion multiplication in the original
\((i,j,k)\) block gives
\[
 b=r+iI+jJ+kL,\qquad
 \begin{array}{ll}
 r=h^2+c^2\cos s,&I=c^2\sin s,\\
 J=ch(1-\cos s),&L=-ch\sin s.
 \end{array}                                                    \tag{PJ.1}\label{p1b:eq:natural-pj-1}
\]
Since \(j_\theta=e^{i\theta/2}je^{-i\theta/2}\), the exact map
in \eqref{p1b:eq:natural-ob-5} is
\[
 D(\tau,s,\theta)
   =\operatorname{Ad}_{e^{i\theta/2}}(b)b^{-1}.
                                                               \tag{PJ.2}\label{p1b:eq:natural-pj-2}
\]
It is identity at \(\tau=0,1\), \(s=0,2\pi\), and
\(\theta=0,2\pi\).  Thus the retained based quotient is
\[
 S^1_\tau\wedge S^1_s\wedge S^1_\theta\cong S^3,
 \qquad D:S^3\longrightarrow SU(2).
                                                               \tag{PJ.3}\label{p1b:eq:natural-pj-3}
\]
In particular there is no boundary contribution to the degree
calculation below.

If \(D=-I\), \eqref{p1b:eq:natural-pj-2} says
\(\operatorname{Ad}_{e^{i\theta/2}}(b)=-b\).  The real and
\(i\) coordinates of \(b\) must therefore vanish.  Its remaining
\((j,k)\) coordinate is a unit vector, so its rotation through
\(\theta\) is its negative only when \(\theta=\pi\).
The possibility \(c=0\) gives \(b=I\) and is excluded.
Consequently \(I=0\) implies \(s=0\) or \(\pi\); the first
again gives \(b=I\).  At \(s=\pi\), \(r=0\) is
\(h^2-c^2=0\), hence \(\tau=1/2\).  Conversely these values
give \(b=j\) and \(D=-I\).  There is exactly one preimage:
\[
 D^{-1}(-I)=\{(1/2,\pi,\pi)\}.                              \tag{PJ.4}\label{p1b:eq:natural-pj-4}
\]

Write the displacements from this point as
\((\delta\tau,\delta s,\delta\theta)\).  Differentiating
\eqref{p1b:eq:natural-pj-1}--\eqref{p1b:eq:natural-pj-2}, in the original parameter directions, gives
\[
 D=-I-\delta\theta\,i-2\pi\delta\tau\,j+\delta s\,k
       +O\bigl(\lVert(\delta\tau,\delta s,\delta\theta)\rVert^2\bigr).
                                                               \tag{PJ.5}\label{p1b:eq:natural-pj-5}
\]
Thus the Jacobian in the target imaginary-coordinate chart at
\(-I\), with column order \((\tau,s,\theta)\), is
\[
 \frac{\partial(\operatorname{Im}_iD,
                     \operatorname{Im}_jD,
                     \operatorname{Im}_kD)}
      {\partial(\tau,s,\theta)}
   =\begin{pmatrix}0&0&-1\\-2\pi&0&0\\0&1&0\end{pmatrix},
 \qquad \det=2\pi.
                                                               \tag{PJ.6}\label{p1b:eq:natural-pj-6}
\]
The unique preimage is regular, proving \(\lvert\deg D\rvert=1\).
For explicit sign bookkeeping, orient the source by
\(d\tau\wedge ds\wedge d\theta\), and orient \(SU(2)\) by the
left-invariant orientation whose basis at identity is \((i,j,k)\).
The imaginary-coordinate chart at \(-I\) has the opposite orientation:
left translation carries that basis to \((-i,-j,-k)\).  With these
conventions \eqref{p1b:eq:natural-pj-6} therefore gives \(\deg D=-1\).  The positive
coordinate determinant is not silently identified with a positive
global degree.  Reversing a supplied sphere identification changes
this displayed sign, while retaining unit absolute degree.

\subsubsection{Each original frame cap has unit absolute degree}

Use the actual endpoint \eqref{p1b:eq:natural-ob-8}, writing its fourth source coordinate
as \(D_0\) to distinguish it from \eqref{p1b:eq:natural-pj-2}:
\[
 (A,B,C,D_0)=M(-y),\qquad
 \alpha=\frac{(A-iD_0)^2}{\sqrt{A^2+D_0^2}},\quad
 \beta=B+iC,\qquad
 C_2(y)=\begin{pmatrix}\alpha&\beta\\-\bar\beta&\bar\alpha\end{pmatrix}.
                                                               \tag{PJ.7}\label{p1b:eq:natural-pj-7}
\]
The two preimages of target \(-I\) in these coordinates are exactly
\((0,0,0,1)\) and \((0,0,0,-1)\).  In the tangent chart
\((A,B,C)\) at either point, with \(D_0=\varepsilon\in\{1,-1\}\),
the target imaginary-coordinate differential is
\[
 (\delta A,\delta B,\delta C)
       \longmapsto(-2\varepsilon\delta A,\delta B,\delta C),
 \qquad \det=-2\varepsilon.
                                                               \tag{PJ.8}\label{p1b:eq:natural-pj-8}
\]
Both inverse branches \eqref{p1b:eq:natural-ob-10} are consequently genuine local
degree-one charts in absolute value.  They are the original labelled
caps: \(Q_+\) uses \(\varepsilon=1\) and the \(z_-\) branch,
and \(Q_-\) uses \(\varepsilon=-1\) and the \(z_+\) branch.
No cap has been selected by a parity calibration.

For example, orient the source unit sphere by its boundary orientation
in \((A,B,C,D_0)\).  The chart \((A,B,C)\) has orientation sign
\(-\varepsilon\); the target chart at \(-I\) has sign \(-1\),
as above.  The local degrees in these conventions are both \(-1\).
The map \(M\) is isotopic to identity through the same rational
formula with \(1/3\) replaced by a parameter from zero to \(1/3\);
all denominators stay positive.  The antipodal map on \(S^3\) has
degree \(+1\).  Hence \(y\mapsto M(-y)\) does not change these
local signs when the source uses its usual quaternion orientation.
Equivalently, orienting the source by the twice-suspension coordinates
\((A,-D_0,B,C)\) gives both local degrees \(+1\).  These explicit
coordinate conventions distinguish the conjugated circle in \eqref{p1b:eq:natural-pj-7}
from the circle-squaring orientation of the fixed section in
\eqref{p1b:eq:rt-58-3}.

The target compression \(p_\ell\) in \eqref{p1b:eq:natural-ob-9} is a based homotopy
from identity.  Its restriction to each of the two original inverse
balls gives the corresponding separate cap, equal to identity on the
cap boundary.  Each target \(-I\) still has the one preimage in its
ball.  Thus
\[
       \lvert\deg Q_+\rvert=\lvert\deg Q_-\rvert=1.
                                                               \tag{PJ.9}\label{p1b:eq:natural-pj-9}
\]
The actual relative cap homotopy \eqref{p1b:eq:natural-ob-11}--\eqref{p1b:eq:natural-ob-12} also confirms that
the two signs agree.  Its source flow fixes the basepoint and retains
the original output; it is not used to substitute a new cap.

\subsubsection{The original correction primitive and its Hopf composition}

Now retain the precise commutator order in \eqref{p1b:eq:natural-ob-13}:
\[
 K(\tau,s,y,\theta)
     =D(\tau,s,\theta)^{-1}Q_+(y)^{-1}
                          D(\tau,s,\theta)Q_+(y).
                                                               \tag{PJ.10}\label{p1b:eq:natural-pj-10}
\]
It is identity on the based wedge, and hence is the based map
\(S^3_{\tau,s,\theta}\wedge S^3_y\to SU(2)\) after the specified
regrouping of its source coordinates.  If
\(\mathfrak c(a,b)=aba^{-1}b^{-1}\), the pointwise identity is
\[
                  K=\mathfrak c(D^{-1},Q_+^{-1}).             \tag{PJ.11}\label{p1b:eq:natural-pj-11}
\]
Inversion on \(S^3\) has degree \(-1\).  Both operands in
\eqref{p1b:eq:natural-pj-11} therefore have degree \(\pm1\), by \eqref{p1b:eq:natural-pj-6} and \eqref{p1b:eq:natural-pj-9}.
Based self-maps of \(S^3\) with those degrees are based homotopic
to identity or a reflection.  Their smash is consequently a
degree-\(\pm1\) self-map of \(S^6\).  It follows that \(K\) is
the quaternion commutator generator up to the retained source and
generator signs.  Here the generator fact is the one for the actual
order \(aba^{-1}b^{-1}\) in Thomas P\"uttmann,
\emph{On the commutator of unit quaternions and the numbers 12 and 24},
arXiv:1101.5147v4, Section~1, p.~1
(\texttt{https://arxiv.org/abs/1101.5147}).  No external source is used
to evaluate a natural boundary coefficient.

For further sign tracking, with the product orientation
\((\tau,s,\theta;y)\) and the usual quaternion orientation on
\(y\), the two operand degrees in \eqref{p1b:eq:natural-pj-11} are both \(+1\),
so this convention gives \([K]=[\mathfrak c]\).
Regrouping the written argument order \((\tau,s,y,\theta)\)
to \((\tau,s,\theta,y)\) moves a one-dimensional coordinate past
a three-dimensional factor and introduces a minus sign.  The safe
statement in the retained \(\nu'\) convention, without changing
any supplied sphere identification, is therefore
\[
          {}[K]=\pm\nu'\in\pi_6SU(2)=C_{12}.
                                                               \tag{PJ.12}\label{p1b:eq:natural-pj-12}
\]
In particular \(K\) is primitive, rather than merely an arbitrary
odd multiple of a two-primary generator.  Using the already retained
relation \(\nu'\eta_6=\mu\) from \eqref{p1b:eq:rt-61-12}, or equivalently
the normalization used in \eqref{p1b:eq:natural-ne-15}, gives
\[
                  {}[K\eta_6]=\mu\ne0
                    \quad\hbox{in }\pi_7SU(2)=C_2.             \tag{PJ.13}\label{p1b:eq:natural-pj-13}
\]
The source and generator signs disappear in this order-two group.

Equation \eqref{p1b:eq:natural-pj-13} evaluates this actual projective correction primitive
only.  The whole remainder \(\mathcal R_J\) additionally retains
the original joining path \eqref{p1b:eq:natural-nd-8}, localization and prefix comparison
\eqref{p1b:eq:natural-om-6}, and the complete projective nullhomotopy \(\Pi\) in
\eqref{p1b:eq:natural-ob-14}--\eqref{p1b:eq:natural-ob-15}.  No relative homotopy identifying that whole map
with \(K\eta_6\), or proving that its other contributions cancel,
is supplied by this degree calculation.  Accordingly
\(r_J,h_t,\zeta_{\rm nat}\) are not assigned here.



\subsection{A cut-defined intrinsic conjugation and its complete matching}
\label{p1b:sec:natural-joining-cut-conjugation}

Retain the actual determinant gauge $D_Y^0$, actual root $H^3$,
actual ordinary frames and odd clutch of \eqref{p1b:eq:natural-oa-3}--\eqref{p1b:eq:natural-oa-4} and \eqref{p1b:eq:natural-nc-2},
and the actual extension $R$ of \eqref{p1b:eq:natural-nc-9}--\eqref{p1b:eq:natural-nc-12}.
The construction below supplies a different complete matched intrinsic
triple, with an explicit two-leg joining.  It also retains the exact
second-cut frame comparison.  In particular the auxiliary conjugator
is not asserted to descend as a gauge on $Y$.

\paragraph{A conjugator on each pair space in the first primitive cut.}
For $0\le r\le1$, put
\[
 x=y e^{\pi r u/2},\qquad p=y e^{-\pi r u/2},\qquad
 X_a(x,p)=F_Y(y,u,\pi r/2)\,
                 R(a,y,u,r)\,F_Y(y,u,\pi r/2)^*.
 \tag{PI.1}\label{p1b:eq:natural-pi-1}
\]
Here $R(a,y,u,r)$ is literally \eqref{p1b:eq:natural-nc-12}, not a bare cap quaternion
in a chosen global fibre frame.  The displayed pair coordinates cover
all nonantipodal pairs.  Their inverse uses the unique logarithm of
$p^{-1}x$ of norm less than $\pi$; the polar direction is absent
when $x=p$.
At $r=0$, $F_Y=F_0(y)$ and $C_r=I$, so \eqref{p1b:eq:natural-pi-1} is independent of
that absent direction.  At $r=1$, \eqref{p1b:eq:natural-nc-6} and \eqref{p1b:eq:natural-nc-8} give the exact
limit
\[
 X_a(w,-w)=Q_w D_a\operatorname{diag}(f_j(a^{-1}w),1)
                         D_a^*Q_w^*,\qquad
 Q_w=T_wQ_0,\quad D_a=\operatorname{diag}(1,a,1).
 \tag{PI.2}\label{p1b:eq:natural-pi-2}
\]
It is independent of the redundant choices $w=yu$ at an antipodal
pair.  Uniform continuity on the compact coordinate source therefore
proves continuity at both quotient loci.  Thus $X_a$ is an actual
determinant-one gauge of $E_s$ on $P$ for each first-cut parameter.
It fixes $s$ throughout.  Its support requires $x$ within $\delta$
of the purely imaginary quaternion $aj$.  Consequently it is identity
on the entire original round real cap $B$, including its collar.

\paragraph{Descend the determinant gauge, rather than the conjugator.}
On the first primitive cut define
\[
 \widetilde D_Y=X_aD_Y^0X_a^*,\qquad
 T_{\rm com}=(D_Y^0)^{-1}X_aD_Y^0X_a^*.
 \tag{PI.3}\label{p1b:eq:natural-pi-3}
\]
On the second cut define each of these two gauges by the original
exchange law.  At $a=\pm1$, $D_Y^0=I$ off the root support, whereas
$X_a=I$ on that support.  Hence the exact endpoint equations are
\[
 T_{\rm com}|_{a=\pm1}=I,\qquad
 \widetilde D_Y|_{a=1}=H^3,\qquad
 \widetilde D_Y|_{a=-1}=H^{-3}.
 \tag{PI.4}\label{p1b:eq:natural-pi-4}
\]
The original inverse-exchange identity for $H$ proves both descended
cut seams.  All values of $T_{\rm com}$ are intrinsic and have
determinant one; $\det\widetilde D_Y=a^2$.  The root cap and original
closing seam are retained literally.
The analogous endpoint assertion for $X_a$ itself is not made:
its odd joining sphere degree prevents it from being a descended
$Y$ gauge.  Its use inside \eqref{p1b:eq:natural-pi-3}, with the explicit second-cut
definition and \eqref{p1b:eq:natural-pi-4}, is what gives legitimate descended gauges.

\paragraph{An explicit complete joining.}
On the joining locus retain $X_a$ of \eqref{p1b:eq:natural-pi-2}, and first replace the
positive line in $D_+$ by the actual aligned line $r_\lambda$ of
\eqref{p1b:eq:natural-nd-7}, while holding $X_a$ fixed.  At the end the value is
\[
 Q_Y(a,w)\operatorname{diag}(f_j,1)
       \operatorname{diag}(a^2,1,1)
       \operatorname{diag}(f_j^{-1},1)Q_Y(a,w)^*.
 \tag{PI.5}\label{p1b:eq:natural-pi-5}
\]
Here the argument of $f_j$ remains $a^{-1}w$.
In a second interval rotate the displayed positive-line phase from
$a^2$ to $-a^2$.  Its terminal value is precisely
\[
 Q_Y\operatorname{diag}(f_j,1)
       \operatorname{diag}(-a^2,1,1)
       \operatorname{diag}(f_j^{-1},1)Q_Y^*
       =c_jD_Zc_j^*.
 \tag{PI.6}\label{p1b:eq:natural-pi-6}
\]
These are actual ordinary rank-two paths with the original stabilizing
line and deck descent.  Equal half-intervals specify their concatenation.
There is no comparison of two Pontryagin representatives in this joining.

The measured secondary matching is also retained.  Applying \eqref{p1b:eq:natural-nd-7}
to both determinant factors in $T_{\rm com}|_J$ gives, in $Q_Y$ axes,
\[
 D(a^2)^{-1}f_j(a^{-1}w)D(a^2)f_j(a^{-1}w)^{-1},
 \qquad D(\xi)=\operatorname{diag}(\xi,1).
 \tag{PI.7}\label{p1b:eq:natural-pi-7}
\]
In the original $(\xi,u)=(a^2,a^{-1}w)$ joining coordinates it is
the based adjoint conjugation representative of $\eta_3$, with circle
degree $-1$ and the actual cap degree $+1$.  Thus its joining secondary
component is one, as required for the measured $e_h=0$ origin.
The triple consisting of \eqref{p1b:eq:natural-pi-3}, this two-leg joining, and the original
$D_Z$ is a complete matched determinant triple with the original low
data.  The complete matched-choice invariance of \eqref{p1b:eq:rt-53-4}, including
joining-path invariance \eqref{p1b:eq:rt-49-6}, says its normalized natural $h_t$
equals that of the original triple.  This is not an identification of
its residual exterior map with the original stabilized $\mathcal R_J$.

\paragraph{The exact whole-circle restriction keeps a second-cut word.}
On the first primitive cut the actual odd frame gives the pointwise
cancellation
\[
 (F_YR)^*\widetilde D_Y(F_YR)=F_Y^*D_Y^0F_Y.
 \tag{PI.8}\label{p1b:eq:natural-pi-8}
\]
On the second cut the conjugator is the exchanged first-cut formula,
whereas the retained odd frame remains the actual \eqref{p1b:eq:natural-nc-12}.
Let $d=\operatorname{diag}(1,-1)$ in the intrinsic block.  The exact
remaining frame word is
\[
 G_a(y,u,r)=R(a,y,u,r)^*\,dR(-a,y,-u,r)d,
 \qquad
 (F_YR)^*\widetilde D_Y(F_YR)
     =G_a\,(F_Y^*D_Y^0F_Y)\,G_a^{-1}.
 \tag{PI.9}\label{p1b:eq:natural-pi-9}
\]
This follows directly from
$F_Y(y,-u,\alpha)=\mathsf JF_Y(y,u,\alpha)d_3$.
At the neck $G_a=I$, by \eqref{p1b:eq:natural-nc-7}--\eqref{p1b:eq:natural-nc-8},
$D_{-a}=D_ad_3$, and $(-a)^{-1}(-w)=a^{-1}w$.
It is also identity on the whole original real cap and its collar.
At the polar centre its actual direction-independent value is
\[
 G_a(y,0)=D_a f_j(a^{-1}y)^{-1}f_j(-a^{-1}y)D_a^*.
 \tag{PI.10}\label{p1b:eq:natural-pi-10}
\]
In particular \eqref{p1b:eq:natural-pi-9} cannot be deleted by the first-cut cancellation.

Outside the cap put $C=\mathsf V^{\#*}F_Y$, which is independent
of $a$, and $p_a=\operatorname{diag}(a^2,1)$.
Before the actual parameter divisions its second-cut phase-normalized
matrix has the exact factorization
\[
 M=p_a^{-1}G_a(C^*p_aC)G_a^{-1}
   =\underbrace{p_a^{-1}G_ap_aG_a^{-1}}_{P(p_a,G_a)}
           \;G_aM_0G_a^{-1},\qquad M_0=p_a^{-1}C^*p_aC.
 \tag{PI.11}\label{p1b:eq:natural-pi-11}
\]
The actual sphere-base and circle-base divisions must be applied to
this whole product.  It has both the phase commutator and the complete
conjugation of $M_0$; the latter is not omitted.

\paragraph{A relative kernel for the separated phase commutator.}
For each cut parameter normalize the actual $G_a$ by its actual
sphere-base value:
$\overline G_a(y,n)=G_a(-y,n)G_a(-1,n)^{-1}$.
The neck and sphere-base faces are identity and the polar centre is
continuous by \eqref{p1b:eq:natural-pi-10}.  Extending by identity across the outer normal
stage makes this a literal map
$S^3_y\wedge(D^3/\partial D^3)=S^6\to SU(2)$.
On a primitive cut, the phase commutator
$P(p_{a(s)},\overline G_{a(s)})$ is based at both time endpoints,
since $p_{a(s)}$ is scalar there.  In this whole commutator replace
only $\overline G_{a(s)}$ by $\overline G_{a((1-h)s)}$.
All its relative faces remain identity.  Its terminal map is the
adjoint operation on one actual $S^6$ map, so
\[
 {}[P(p_a,\overline G_a)]
     \in\eta_3E\pi_6SU(2)=0.
 \tag{PI.12}\label{p1b:eq:natural-pi-12}
\]
The zero is the established \eqref{p1b:eq:natural-pg-7} source relation, not an arbitrary
$S^4$ filling choice.  This evaluates only the separated normalized
phase commutator.  The next MAIN comparison \eqref{p1b:eq:natural-sc-8}--\eqref{p1b:eq:natural-sc-14} retains and
evaluates the whole product \eqref{p1b:eq:natural-pi-11}: its second factor is the strict
universal word $[Z,P(t,C^{-1})]$, whose frozen actual input is a
relative map to $S^6$ of even degree.  The complete actual divisions
are carried in \eqref{p1b:eq:natural-sc-13}--\eqref{p1b:eq:natural-sc-14}.  That comparison proves the full
second-cut correction zero and hence completes the two-cut reduction
to the coherent untoggled boundary, with its original projective tail.

No coefficient of $h_t$, $\zeta_{\rm nat}$, or the original $r_J$ is assigned
by this reduction.  Evaluation of the resulting complete coherent
boundary would give $h_t$ directly; the proved original equation
\eqref{p1b:eq:natural-pg-16} would then give $r_J=h_t+1$.  The original stabilized residual
is not identified with a new intrinsic exterior by choice invariance.

\paragraph{Correlated algebra for the entire sphere-base division.}
The nonidentity original circle-base matrix is supported in the real
root cap.  Since $G=I$ there, its division preserves the exact form
$M_{\rm new}=P(p,G)GM_0G^{-1}$, now with $M_0$ denoting the complete
actual circle-divided coherent matrix.  Work in the rebased sphere
coordinate and put
\[
 \begin{gathered}
 A_0=M_0|_{y=1},\quad E_0=M_0A_0^{-1},\quad
 G_b=G|_{y=1},\quad H=GG_b^{-1},\quad p_b=G_bpG_b^{-1},\\
 T=p^{-1}p_b,\qquad
 E_b=G_bE_0G_b^{-1},\qquad A_b=G_bA_0G_b^{-1}.
 \end{gathered}
\]
Direct multiplication, with the actual sphere-base matrix of the new
target rather than a chosen replacement, gives
\[
 \boxed{E_{\rm new}
    =T\,P(p_b,H)\,(HE_bH^{-1})\,[H,A_b]\,T^{-1}.}
 \tag{PI.13}\label{p1b:eq:natural-pi-13}
\]
For example its uncombined form is
$p^{-1}Hp_bE_bA_bH^{-1}A_b^{-1}p_b^{-1}p$;
substituting $T=p^{-1}p_b$ proves the displayed order literally.
Thus the phase term, both actual lower values, and the whole conjugation
remain correlated.  On this second-cut comparison $T$ is identity at
both time seams since $p$ is scalar there, and is identity at the neck.
Extend it by identity on the first cut and on the outer normal stages.
It is then a genuine ordinary boundary gauge on the full original
base.  Conjugation by this particular $T$ therefore preserves the
normalized top coefficient by \eqref{p1b:eq:rt-54-4}.

There is a useful exact support bound for the remaining normal factor
$G_b$.  In the original sphere-base slice $y=-1$,
$x=-e^{\alpha u}$ and $p=-e^{-\alpha u}$.  Both relevant cap centres
are purely imaginary.  Either cap being nonidentity consequently
requires $\pi/2-\alpha<\delta$, so
\[
 \operatorname{supp}(G_b)\subset
            \{\alpha>\pi/2-\delta\}.\tag{PI.14}\label{p1b:eq:natural-pi-14}
\]
This is a neck collar disjoint from the real root cap.  At its neck
the coherent matrix, after the actual \eqref{p1b:eq:natural-nd-7} positive-line comparison,
is $C_0$ and has actual sphere degree two by \eqref{p1b:eq:natural-ob-6}--\eqref{p1b:eq:natural-ob-8}.
These facts supply the actual source and degree retained in the whole
comparison \eqref{p1b:eq:natural-sc-9}--\eqref{p1b:eq:natural-sc-12}.  The proof there uses the universal whole
word and its even relative degree; it does not infer vanishing of
mixed commutators from a cell-dimension claim in \eqref{p1b:eq:natural-pi-13}.
In particular the intermediate time face of $E_0$ retains the
nonconstant root-cap and projective comparison tracks.  The closed
$\Delta_b$ of \eqref{p1b:eq:natural-sc-14}, with its opposite-arc compression, carries
those entire faces.  The coherent boundary, with both primitive cuts and its original
projective tail, is identified in (CI) and evaluated in (AJ) before
applying $r_J=h_t+1$.



\subsection{The relative second-cut comparison has zero top class}
\label{p1b:sec:natural-second-cut-phase}

The following comparison retains both factors of the proposed two-cut
boundary difference and their actual lower values.  It proves that this
difference has zero seven-sphere class.  Evaluation of the resulting
coherent exterior, including its original projective tail, is a separate
calculation.  No value of the original joining coefficient is inferred.

Let $N=D^3/\partial D^3=S^3$, with its collapsed boundary as
basepoint, and let $Y=S^3$, with the retained sphere basepoint $y_b$.
Put
\[
 q_t=\operatorname{diag}(e^{2\pi it},1),\qquad
 P(t,G)=q_t^{-1}Gq_tG^{-1},\qquad 0\le t\le1.
                                                        \tag{SC.1}\label{p1b:eq:natural-sc-1}
\]
Conjugation is in the original intrinsic two-block.  The scalar part
of $q_t$ does not affect it.  Thus $P$ is the winding-one adjoint
map $S^1\wedge SU(2)\to SU(2)$; in the retained axes its class is
$\eta_3$ (a sign convention is immaterial below).

\begin{proposition}[Actual moving input and its actual lower division]
\label{p1b:proposition:second-cut-phase-zero}
Let $Z_t:Y\times N\to SU(2)$ be a continuous family which is
identity on $Y\times\{*\}$.  No constancy is required on either
time end or on $\{y_b\}\times N$.  Then the literal lower-divided
map
\[
 \mathcal P(t,y,n)=P(t,Z_t(y,n))
                         P(t,Z_t(y_b,n))^{-1}
                                                        \tag{SC.2}\label{p1b:eq:natural-sc-2}
\]
is a based map $S^1\wedge Y\wedge N=S^7\to SU(2)$ and has
class zero.  It admits the explicit relative comparisons \eqref{p1b:eq:natural-sc-3}--\eqref{p1b:eq:natural-sc-6}
to a map of class $\eta_3 E\xi$, for a specifically defined
$\xi\in\pi_6SU(2)$.
\end{proposition}
\begin{proof}
Both time faces in \eqref{p1b:eq:natural-sc-2} are literally identity because $q_0=q_1=I$,
regardless of the input $Z_t$ there.  The sphere-base and normal-base
faces are also identity.  In the entire word, including the lower
factor, replace its input time by
\[
                 Z_t\longmapsto Z_{(1-h)t},\quad 0\le h\le1.
                                                        \tag{SC.3}\label{p1b:eq:natural-sc-3}
\]
This is an endpoint-preserving homotopy of the closed map: it keeps
$q_t$ fixed, so both time faces stay identity at every $h$.  It also
keeps the sphere and normal base faces.  Its endpoint uses $Z_0$.
Write its actual normal-only lower value and its based difference as
\[
 B(n)=Z_0(y_b,n),\qquad G(y,n)=Z_0(y,n)B(n)^{-1}.
                                                        \tag{SC.4}\label{p1b:eq:natural-sc-4}
\]
The full product has not been replaced by $G$.  Rather, the exact
multiplication identity is
\[
 P(t,UV)=P(t,U)\,U P(t,V)U^{-1},
 \quad\Longrightarrow\quad
 \mathcal P_0=P(t,G)\,[G,P(t,B)].
                                                        \tag{SC.5}\label{p1b:eq:natural-sc-5}
\]
Here $[U,V]=UVU^{-1}V^{-1}$.  In particular the omitted-lower-value
inference would be false as a pointwise equation.

To remove the last closed commutator, use the actual normal sphere
$N$, rather than the parameter sphere.  Choose its stereographic
coordinate $x\in\mathbb R^3$ from $*$ and the two based degree-one
compressions
\[
 c_h^\pm(x)=
 \begin{cases}
 \mathrm{st}^{-1}\!\left(
 {x\mp2h e_1\over1-h|x\mp2h e_1|^2}\right),
                &1-h|x\mp2h e_1|^2>0,\\
 *,&1-h|x\mp2h e_1|^2\le0.
 \end{cases}                                             \tag{SC.6}\label{p1b:eq:natural-sc-6}
\]
They begin at identity and end with disjoint supports, the two unit
balls centred at $\pm2e_1$.  In the complete right side of \eqref{p1b:eq:natural-sc-5}
replace $G(y,n)$ by $G(y,c_h^-(n))$ and $B(n)$ by
$B(c_h^+(n))$.  This is a homotopy of the whole based seven-sphere
word.  On the time faces $P(t,G)=P(t,B)=I$, even though $G$ itself
need not be identity there; on the sphere-base face $G=I$; on the
normal-base face both inputs are identity.  At $h=1$ the supports
are disjoint, so $[G,P(t,B)]=I$ pointwise.  Reverse only the
compression of $G$ in the remaining factor $P(t,G)$.  This ends
at the literal map $P(t,G(y,n))$ and preserves every collapsed face.

By \eqref{p1b:eq:natural-sc-4}, $G$ is identity on the entire wedge and therefore defines
$G:Y\wedge N=S^6\to SU(2)$.  Put $\xi=[G]$.  The resulting
seven-sphere class is $\eta_3E\xi$.  The retained relation
\[
 {}[\iota_4,\iota_4]=\pm(2\nu_4-E\nu'),\quad
 \eta_3\nu_4=\mu,\quad 2\mu=0
 \quad\Longrightarrow\quad
                  \eta_3E\pi_6SU(2)=0                 \tag{SC.7}\label{p1b:eq:natural-sc-7}
\]
proves the claim.  The implication uses naturality of Whitehead
products into the topological group $SU(2)$; it does not replace
$E\nu'$ by $2\nu_4$ in $\pi_7S^4$.
\end{proof}

\paragraph{The precise second-cut application and its whole word.}
Let $X_+$ be the auxiliary first-cut map obtained from \eqref{p1b:eq:natural-nc-12}, and
define its second-cut value by exchange,
$X_-(a,x,p)=\mathsf J X_+(-a,p,x)\mathsf J$.
It is $X_+D_Y^0X_+^{-1}$, with this exchange prescription, that
can define a descended gauge; $X_+$ itself is not a descended gauge.
In the original odd inner frame write
$A=F_Y^*X_-F_Y$ and $Z=R^{-1}A$.  On the second primitive cut
the restriction of the new gauge is $Z D Z^{-1}$, where
$D=F_Y^*D_Y^0F_Y$.  The two maps $A$ and $R$ agree at the neck,
so $Z=I$ there.  Their supports miss the original real root cap, so
$Z=I$ on that cap.  Their values at the normal origin are independent
of the absent polar direction.  Thus this $Z_t$ satisfies the
normal-relative hypotheses of the proposition.

On the inner stage the actual logarithm factor is constant, and after
the $a=1$ rebase the normalized phase is
$q_t=\operatorname{diag}(a^2,1)$, with
$a=-e^{\pi it}$.  Put $M_0=q_t^{-1}D$ and
$M_{\rm new}=q_t^{-1}ZDZ^{-1}$.  Direct multiplication gives the
exact whole word
\[
 \boxed{M_{\rm new}M_0^{-1}
       =P(t,Z)\,\bigl(ZM_0Z^{-1}M_0^{-1}\bigr).}
                                                        \tag{SC.8}\label{p1b:eq:natural-sc-8}
\]
The first factor, with its own actual sphere-base division, is zero
by the proposition.  The second factor is retained and evaluated next.

\begin{proposition}[The entire conjugation factor, with the exterior frame]
\label{p1b:proposition:second-cut-conjugation-zero}
Let $M=Y\times D^3$ be the actual inner normal domain, and let
$B\subset\operatorname{int}M$ be the original embedded real six-ball.
The actual circle-independent exterior matrix is
$C=V^*F_Y$ on $M\setminus B$.  On the support of $Z$,
\[
 M_0=P(t,C^{-1}),\qquad
 Q=[Z_t,P(t,C^{-1})].                                  \tag{SC.9}\label{p1b:eq:natural-sc-9}
\]
Extend this word by identity where $Z_t=I$, in particular on a
neighbourhood of $B$ and on the neck.  Its actual sphere-base
division has zero seven-sphere class.
\end{proposition}
\begin{proof}
The first equality is the actual untoggled exterior equation
$D_Y^0=V\operatorname{diag}(a^2,1)V^*$, expressed in $F_Y$.
It holds wherever $Z$ is supported: each of the two odd-clutch tubes
requires $x$ or $p$ near a purely imaginary cap, while the root ball
has both coordinates near the real pole.  Thus no root interior
formula has been replaced by an exterior formula.

In the entire second word of \eqref{p1b:eq:natural-sc-9}, freeze its actual input by
$Z_t\mapsto Z_{(1-h)t}$.  Since $P(0,C^{-1})=P(1,C^{-1})=I$,
both time faces are fixed without requiring $Z_t$ constant there.
The neck, root and polar faces are preserved.  Divide at every $h$
by the moving actual value at $y_b$.  This is a homotopy of the
based sphere word, ending at the normalization of
$[Z_0,P(t,C^{-1})]$.

Write $q:S^3\times S^3\to S^3\wedge S^3=S^6$ for the smash
quotient.  The universal whole word
\[
 \Pi(t,U,W)=[U,P(t,W)]
 \quad\text{factors strictly through}\quad
 S^1\wedge S^3\wedge S^3=S^7.                          \tag{SC.10}\label{p1b:eq:natural-sc-10}
\]
Consequently the frozen word is $\Pi$ evaluated on
$T(y,n)=q(Z_0(y,n),C(y,n)^{-1})$.  This $T$ is an actual
continuous map on all of $M$, despite the missing frame on $B$:
on a neighbourhood of $B$ the first input is literally identity,
so the smash quotient is its basepoint.  On the neck it is also
basepoint.  There is no extension of $C$ across the nonzero clutch
implicit in this statement.

Its degree after the sphere-base collapse is even.  Here is the
relative calculation, including the lower normal component.  The
only degree-three class of $M\setminus\operatorname{int}B$ is the
parameter-sphere class: excision of a six-ball does not change
$H^3(M;\mathbb Z)$.  In this group
\[
 C^*u_3=\pm2u_Y,
 \qquad Z_0^*u_3=N u_N\in H^3(M,Y\times\partial D^3;\mathbb Z)
                                                        \tag{SC.11}\label{p1b:eq:natural-sc-11}
\]
for an actual integer $N$.  The second statement follows because
$Z_0$ is identity on the neck; the relative group has only the
normal three-cell generator.  It does not presume $N=0$.
For the first statement, evaluate at the neck.  There
$C=g_w^{-1}C_0$.  The actual line comparison \eqref{p1b:eq:natural-nd-7} homotopes the
first line of $g_w=Q_w^*V(w,-w)$ to the constant first line, with
strictly positive aligned overlap.  Lift that comparison using the
same finite unitary line transport; its terminal matrix is diagonal.
A diagonal map on $S^3$ has a real logarithm and degree zero.
Thus $g_w$ has degree zero.  The actual based $C_0$ slice has
absolute degree two by \eqref{p1b:eq:natural-ob-6}--\eqref{p1b:eq:natural-ob-10}, and
$y\mapsto(-y)u$ has degree one.
This proves $C^*u_3=\pm2u_Y$, with all actual frame factors retained.

The smash quotient obeys $q^*u_6=u_3\smile u_3$, up to orientation.
Equations \eqref{p1b:eq:natural-sc-11} therefore give
\[
       \deg T=\pm2N.                                  \tag{SC.12}\label{p1b:eq:natural-sc-12}
\]
This cup product is well defined relatively: the $Z_0$ factor is
zero near $B$ and on the neck, so the undefined value of $C$ on
$B$ cannot contribute.

To make the sphere-base collapse literal, retain its original
restriction $T|_{\{y_b\}\times D^3}$.  It is a based map
$S^3\to S^6$ and has a based nullhomotopy, since $\pi_3S^6=0$.
Extend that nullhomotopy by homotopy extension, relative to the
neck and the neighbourhood of $B$.  It ends at a map factoring
through $Y\wedge(D^3/\partial D^3)=S^6$ and preserves \eqref{p1b:eq:natural-sc-12}.
At every point of this homotopy apply $\Pi$, and divide by its
actual moving sphere-base value.  Thus its initial value is the
literal normalization of $Q$, not an assumed based replacement;
its endpoint is $\Pi\circ\Sigma T$ on $S^7$.

The latter class is $(\pm2N)[\Pi]$.  The whole target group
$\pi_7SU(2)=C_2$ has exponent two, so it is zero.  No value of
$[\Pi]$ is needed to infer the vanishing, and in particular a
single possible Samelson $\mu$ has not been substituted for this
even coefficient.
\end{proof}

\paragraph{Carry both actual parameter divisions.}
For two matrix families $U,V$ and their actual sphere-base values
$U_b,V_b$, multiplication gives
\[
 (UV)(U_bV_b)^{-1}
 =(UU_b^{-1})\,U_b(VV_b^{-1})U_b^{-1}.                 \tag{SC.13}\label{p1b:eq:natural-sc-13}
\]
Apply this to the two actual factors of \eqref{p1b:eq:natural-sc-8}.  The resulting
conjugation by $P_b$ contributes zero: $P_b$ is based on the time
and normal faces, and the normalized $Q$ is also based on those
faces and on the sphere-base face.  Compress their actual time-circle
supports independently into the two disjoint based arcs, using
\eqref{p1b:eq:natural-sc-6} in dimension one.  The closed commutator is identity at the
disjoint-support endpoint.  Every collapsed face is retained during
this homotopy.  Thus the class of the actual sphere-base division
of $PQ$ is the sum of the two zero classes just proved.

Finally let $G_0$ and $E_0=G_0G_{0,b}^{-1}$ be the actual coherent
circle-divided and sphere-divided supported maps, including their
unchanged projective collar.  The first-cut boundary cancellation is
literal, so its circle-base value is unchanged.  If
$\Delta=M_{\rm new}M_0^{-1}$ denotes the entire second-cut word,
extended by identity on the other cut, its first divided numerator
is exactly $G_{\rm new}=\Delta G_0$.  Its sphere division is
\[
 E_{\rm new}=\Delta E_0\Delta_b^{-1},\qquad
 E_{\rm new}E_0^{-1}
 =(\Delta\Delta_b^{-1})\,
      (\Delta_bE_0\Delta_b^{-1}E_0^{-1}).              \tag{SC.14}\label{p1b:eq:natural-sc-14}
\]
The second factor is another closed commutator with the same based
time circle.  Extend $\Delta_b$ by identity off the inner stage;
the neck value is identity.  The actual $E_0$ already has both
parameter-base faces and the terminal projective face identity.
The same opposite-arc compression contracts this entire commutator,
including the original projective tail.  The first factor in \eqref{p1b:eq:natural-sc-14}
is precisely the normalized whole word evaluated using \eqref{p1b:eq:natural-sc-13}.
This proves that the full second-cut boundary difference has zero
seven-sphere class.  It certifies the two-cut reduction to the
coherent exterior; it does not evaluate that exterior or assign
$r_J$, $h_t$, or $\zeta_{\rm nat}$.



\subsection{The complete paired comparison of the coherent joining tail}
\label{p1b:sec:natural-joining-coherent-tail}

The original joining coefficient is the coefficient of the whole map
\(\mathcal R_J\) in \eqref{p1b:eq:natural-om-6}--\eqref{p1b:eq:natural-om-7}, with its actual \eqref{p1b:eq:natural-nd-8},
ordinary frames, lower divisions and projective comparison.  The proved
annular and characteristic reductions give
\[
 h_t=1+r_J,\qquad \zeta_{\rm nat}=r_J,\qquad
 [\mathcal R_J]=r_J\mu .                                  \tag{JT.1}\label{p1b:eq:natural-jt-1}
\]
The equations below isolate a specific complete comparison required by
an intrinsic reduction of its coherent exterior.  They do not identify
an isolated commutator or one hemisphere with \(\mathcal R_J\), and
its absolute value follows from \eqref{p1b:eq:natural-ci-8}, \eqref{p1b:eq:natural-aj-4}--\eqref{p1b:eq:natural-aj-5}.

\paragraph{The retained maps and the two ordered disks.}
Use precisely \eqref{p1b:eq:natural-ob-1}--\eqref{p1b:eq:natural-ob-16}, with their actual based frame
\(c(y,u)=C_b(u)^{-1}C_0(-y,u)\).  Put
\(d=\operatorname{diag}(1,-1)\), so
\(c(y,-u)=dc(y,u)d\), and write
\(\mathcal C(B,c)=B^{-1}c^{-1}Bc\).  Suppress the unchanged
arguments \((\tau,s,y)\) when this causes no ambiguity.  Define
\[
 \begin{aligned}
 b_+(\tau,s)&=b_\tau(s,0),&
 b_-(\tau,s)&=b_\tau(s,\pi)=d b_+(\tau,s)d,\\
 N_\pm(\tau,s,y,u)&=\mathcal C(b_\pm(\tau,s),c(y,u)),&
 \mathcal A(\tau,s,y,u)&=N_-(\tau,s,y,u)N_+(\tau,s,y,u)^{-1}.
 \end{aligned}                                                \tag{JT.2}\label{p1b:eq:natural-jt-2}
\]
Thus \(N_+\) is exactly the original ordered-disk comparison
\(N_D\), and \(N_-\) is its actual deck transform.  Both have
initial value \(E_Q\), terminal value identity, and identity
\(s=0\) and \(y=1\) faces.  The ratio \(\mathcal A\) is
identity also on both \(\tau\) faces.

Write \(u(r,\theta)\) and \(u_\theta=u(1,\theta)\) exactly
as in \eqref{p1b:eq:natural-ob-2}.  These are the northern ordered hemisphere and its
boundary.  Let
\[
 X(\tau,s,y,\theta)
    =N_1'(\tau,s,y,\theta)
                   N_+(\tau,s,y,u_\theta)^{-1}
    =\Pi(0,\tau,s,y,\theta).                            \tag{JT.3}\label{p1b:eq:natural-jt-3}
\]
Here \(N_1'=N_1R\) retains the actual \eqref{p1b:eq:natural-ob-13} correction, and
\(\Pi\) is the full nine-leg map of \eqref{p1b:eq:natural-ob-14}--\eqref{p1b:eq:natural-ob-15}, with
\(\Pi(1)=I\).  In particular neither inverse copy of the actual
\(K=\mathcal C(D,Q_+)\) in \(R\) is omitted.

\paragraph{The exact paired map.}
Let \(p:S^2\to\mathbb {RP}^2\) be the actual attaching double
cover of the affine normal three-cell, and set
\[
 Z(\tau,s,y,u)
   =\mathcal N_Q(\tau,s,y,p(u))
                   N_+(\tau,s,y,u)^{-1}.                \tag{JT.4}\label{p1b:eq:natural-jt-4}
\]
The notation for \(\mathcal N_Q\) includes its ordering-line
transition on the ordered cover.  Set
\[
 \chi(r)=\begin{cases}
  1,&0\le r\le1/2,\\
  2(1-r),&1/2\le r\le1.
 \end{cases}
\]
Direct substitution into the original \eqref{p1b:eq:natural-ob-16} gives both complete
hemisphere formulas:
\[
 \boxed{\begin{aligned}
 Z(\tau,s,y,u(r,\theta))
       &=\Pi(\chi(r),\tau,s,y,\theta),\\
 Z(\tau,s,y,-u(r,\theta))
       &=d\Pi(\chi(r),\tau,s,y,\theta)d\,
                     \mathcal A(\tau,s,y,-u(r,\theta)).
 \end{aligned}}                                               \tag{JT.5}\label{p1b:eq:natural-jt-5}
\]
The second expression includes the fixed-axis change from \(+j\)
to \(-j\).  It is generally not just the deck conjugate of the
first expression.

For the same equator point \(u_\theta\), the southern
representative has parameter \(\theta+\pi\).  The exact seam is
\[
 \boxed{X(\theta)=dX(\theta+\pi)d\,
                         \mathcal A(u_\theta).}          \tag{JT.6}\label{p1b:eq:natural-jt-6}
\]
Indeed the actual one-cell comparison descends:
\(dN_1'(\theta+\pi)d=N_1'(\theta)\), while
\[
 dN_+(u_{\theta+\pi})d=N_-(u_\theta).
\]
Consequently the right side of \eqref{p1b:eq:natural-jt-6} is exactly
\[
 N_1'(\theta)N_-(u_\theta)^{-1}
                  N_-(u_\theta)N_+(u_\theta)^{-1}
   =X(\theta).
\]
This proves the paired seam by full-matrix cancellation, retaining its
nonconstant intermediate value.

Every value of \eqref{p1b:eq:natural-jt-5} is intrinsic.  Both \(\tau\) faces and the
\(s=0\), \(y=1\) faces are identity.  At the northern centre
\(u_*=(0,1,0)\), \eqref{p1b:eq:natural-ob-16} is literally \(N_D\), so
\(Z(\tau,s,y,u_*)=I\).  Thus \eqref{p1b:eq:natural-jt-4}, with this actual meridian
as basepoint, is a well-defined based map
\[
 Z:S^1_\tau\wedge S^1_s\wedge S^3_y\wedge S^2_u
          =S^7\longrightarrow SU(2).                    \tag{JT.7}\label{p1b:eq:natural-jt-7}
\]
This does not make the \(\theta=0\) meridian of the southern
hemisphere identity: it retains \(d\Pi(\chi,\pi)d\) and the
actual \(\mathcal A\) factor.  That meridian cannot be collapsed
when dividing the two hemisphere contributions.

\paragraph{The actual boundary-axis path in the complete evaluation.}
The two boundary values of \(\Pi\) in \eqref{p1b:eq:natural-jt-6} are connected by
the explicit path
\[
 L(\kappa,\theta)
   =N_1'(\theta)\,
       \mathcal C\bigl(b_\tau(s,\kappa\pi),
                              c(y,u_\theta)\bigr)^{-1},
       \qquad 0\le\kappa\le1.                           \tag{JT.8}\label{p1b:eq:natural-jt-8}
\]
Its initial value is \(X(\theta)\), and its terminal value is
\(dX(\theta+\pi)d\).  All its \(\tau\), \(s\), and \(y\)
base faces remain identity.  The actual \(R\) in \(N_1'\) remains
in this path.  Equations \eqref{p1b:eq:natural-jt-5}--\eqref{p1b:eq:natural-jt-8}, rather than separate copies
of \(\Pi\), therefore specify the closed comparison to be evaluated.

The known degree-two frame endpoint \eqref{p1b:eq:natural-ob-8} does not by itself evaluate
this closed map.  In particular the two inverse-image caps
\(Q_+,Q_-\), the actual \eqref{p1b:eq:natural-ob-11}--\eqref{p1b:eq:natural-ob-12} relative source path, the
fixed-axis change in \eqref{p1b:eq:natural-jt-8}, and both inverse \(K\) factors in
\eqref{p1b:eq:natural-ob-13} have to be transported together.  An unclosed hemisphere
cannot be assigned a separate \(\pi_7SU(2)\) coefficient and then
doubled.

There are also two different unstable operations.  The actual lens
circle-cofibre operation of \eqref{p1b:eq:rt-25-5} sends a coefficient \(\xi\)
to \(A_*\xi=\eta_3E\xi\), so \(A_*\nu'=0\).  The normal-frame
precomposition sends \(\nu'\) to \(\nu'\eta_6=\mu\).  An
identification of \eqref{p1b:eq:natural-jt-7}--\eqref{p1b:eq:natural-jt-8} with one of these operations must
retain the stated paired seam and actual meridian; the primitive
class of \(K\) alone makes no such identification.

The complete identification with the original residual uses the
matched joining replacement (PI), the entire second-cut transport
(SC), and the actual entrance null in \eqref{p1b:eq:natural-om-6}. These are applied in
\eqref{p1b:eq:natural-ci-1}--\eqref{p1b:eq:natural-ci-8}, followed by the absolute evaluation \eqref{p1b:eq:natural-aj-1}--\eqref{p1b:eq:natural-aj-5}.



\subsection{The coherent boundary and its complete fixed-axis/projective difference}
\label{p1b:sec:natural-coherent-boundary-identification}

Retain the complete matched triple \eqref{p1b:eq:natural-pi-3}--\eqref{p1b:eq:natural-pi-7}, the whole
second-cut zero comparison \eqref{p1b:eq:natural-sc-8}--\eqref{p1b:eq:natural-sc-14}, the original ordinary
frames and both divisions \eqref{p1b:eq:natural-oa-8}--\eqref{p1b:eq:natural-oa-9}, and the original projective
null \eqref{p1b:eq:natural-ob-16}. No normal framing is rotated. Write $E^{\mathrm{coh}}$
for the boundary after the whole SC comparison, with its unchanged
real root cap and unchanged projective collar. We identify this whole
boundary with the sum of a localized real cap and the actual paired
map \eqref{p1b:eq:natural-jt-4}. The identification below is followed by the full paired evaluation \eqref{p1b:eq:natural-aj-1}--\eqref{p1b:eq:natural-aj-5}.

\paragraph{An actual circle-independent frame throughout the exterior.}
Use $C_\rho=V^{\#*}F_Y$ in the root-free inner region. On the
first half of the PI middle joining retain the actual positive-line
path \eqref{p1b:eq:natural-nd-7}. Its operands $\widehat v,b,\gamma$ are independent of
$a$. Choose a continuous real logarithm $\vartheta(w)$ of
$\gamma(w)$ and replace its first vector by
$\widehat r_\lambda=e^{-i\lambda\vartheta(w)}r_\lambda$.
This changes no line projection and hence changes no joining matrix.
It starts at $\widehat v$ and ends at $b$. The logarithm can be
chosen with the deck symmetry: the actual covariance in \eqref{p1b:eq:natural-nd-7}
gives $\gamma(-w)=\gamma(w)$; a map $\mathbb {RP}^3\to S^1$
lifts to $\mathbb R$, since a homomorphism $C_2\to\mathbb Z$
is zero. Its additive constant can be fixed at the actual base value.
Put
\[
 W_\lambda=[\widehat r_\lambda,
                  s\mathbin\times\widehat r_\lambda,s],
 \qquad C_\rho=W_\lambda^*F_b .                       \tag{CI.1}\label{p1b:eq:natural-ci-1}
\]
Thus $W_0=V^\#$ on the neck and $W_1=Q_w$; all three columns
are the actual ordered columns. On the phase half of PI joining use
$C_\rho=Q_w^*F_b=C_0$, and use that same $C_0$ on the outer
stage. These definitions agree at both middle seams. Their deck law
is $C_\rho(y,-u)=dC_\rho(y,u)d$. This is an auxiliary frame on
the ordered normal ball, not a claim that $X_a$ descends as a gauge.

In the actual frame $F_N=X_aF_b$ the conjugations in the PI joining
cancel on both primitive cuts: \eqref{p1b:eq:natural-pi-9} gives $X_aF_b=F_N$ at their
neck. Consequently its middle first-leg matrix is the untoggled
positive-line determinant in $F_b$, and its phase leg is
$C_0^{-1}\operatorname{diag}(a^2e^{i\pi\lambda},1)C_0$.
The outer determinant is
$C_0^{-1}\operatorname{diag}(e^{2is},1)C_0$.
All these are actual matrix equations. In particular the phase leg
does not introduce a circle-dependent exterior frame.

\paragraph{Both literal lower divisions.}
Put $p_s=\operatorname{diag}(e^{2is},1)$ and retain the actual
diagonal $B_\rho$ from \eqref{p1b:eq:natural-oa-8}. The raw exterior has the form
$C_\rho^{-1}\operatorname{diag}(e^{2is}e^{i\phi_\rho},1)C_\rho$,
where $\phi_\rho=\pi$ on the inner stage, advances from $\pi$
to $2\pi$ on the phase leg, and is $2\pi$ on the outer stage.
The first line of \eqref{p1b:eq:natural-oa-9} cancels this exact phase with its actual
$s=0$ value, giving
\[
 G^{\mathrm{coh}}=B_\rho
     p_s^{-1}C_\rho^{-1}p_sC_\rho B_\rho^{-1}.
                                                               \tag{CI.2}\label{p1b:eq:natural-ci-2}
\]
At $y=-1$ every section of the actual coherent inner path
has second complex coordinate zero, as used in \eqref{p1b:eq:natural-ob-1}. Its first
ordinary frame column is in $\mathbb Ce_2$. The exterior positive
line is also $\mathbb Ce_2$, and the middle \eqref{p1b:eq:natural-nd-7} path is constant
on that line. Thus $C_b(\rho,u)=C_\rho(-1,u)$ is diagonal on
every root-free stage and commutes with $p_s$. Its actual first
divided value is identity. Equivalently, retain
$c_\rho=C_b^{-1}C_\rho(-y,u)$; then the second line of \eqref{p1b:eq:natural-oa-9}
and \eqref{p1b:eq:natural-ci-2} give the exact exterior equation
\[
 E^{\mathrm{coh}}=B_\rho\mathcal C(a_s,c_\rho)B_\rho^{-1},
 \quad a_s=e^{is},\quad c_\rho(1,u)=I.                 \tag{CI.3}\label{p1b:eq:natural-ci-3}
\]
The scalar part of $p_s$ does not affect conjugation, so its operation
is precisely the retained $a_s$ commutator. At the normal boundary
$B_\rho=I$ and $c_\rho$ is the literal \eqref{p1b:eq:natural-ob-1} frame $c(y,u)$.
The lower values in CI.2--CI.3 have been cancelled by actual diagonal
commutation, not by dropping a sphere slice.

\paragraph{The fixed-axis null and the actual localized cap.}
Replace $a_s$ in \eqref{p1b:eq:natural-ci-3} by the original B32 $b_\tau(s,0)$,
holding $c_\rho$ and $B_\rho$ fixed. Apply exactly the \eqref{p1b:eq:natural-oj-6}
collar extension across the original real ball $B$. It fixes the
root support and agrees at the outer and inner collar radii. Call
the resulting whole family $H_\tau$. Its initial value is the whole
coherent boundary on the uncompressed normal ball, its endpoint is
an actual cap map $\Gamma_{\mathrm{nat}}$ extended by identity,
and its normal boundary is exactly the ordered-disk null $N_+$
of \eqref{p1b:eq:natural-jt-2}. All parameter base faces are fixed.

The cap class equals the already evaluated old cap class. Here is
the retained baseline comparison. On a neighbourhood of this real
ball both odd frame factors are identity for every circle
value, including their actual sphere-base slices. Those lower slices
are exterior diagonal-frame values and their first divisions are
identity. Write the circle as $a=e^{i\theta}$ and put
$m(\theta,z)=\operatorname{diag}(e^{-2i\theta},1)
 F_Y^*D_Y^0(e^{i\theta},z)F_Y$. The native cap numerator and
the rebased old cap numerator are literally
\[
 G_{\mathrm{nat}}(s,z)=m(s+\pi/2,z)m(\pi/2,z)^{-1},\qquad
 G_{\mathrm{old}}(s,z)=m(s,z)m(0,z)^{-1}.               \tag{CI.4}\label{p1b:eq:natural-ci-4}
\]
The actual inner $B_\rho=\operatorname{diag}(-1,1)$ cancels
the constant $e^{i\pi}$ in this equality. Hence
$G_{\mathrm{nat}}=T_{\pi/2}G_{\mathrm{old}}$, where
$T_\sigma F(s)=F(s+\sigma)F(\sigma)^{-1}$. Apply that same
translation to every entire matrix in the old \eqref{p1b:eq:natural-oj-6} localized family.
On its exterior let $S_\tau=b_\tau(\sigma,0)$ and
$B_\tau(s)=b_\tau(s+\sigma,0)S_\tau^{-1}$. Direct multiplication
gives its boundary null
\[
 T_\sigma\mathcal C(b_\tau(s,0),c)
       =S_\tau^{-1}\mathcal C(B_\tau(s),c)S_\tau .     \tag{CI.5}\label{p1b:eq:natural-ci-5}
\]
Remove the moving constant conjugation relative to its initial value
$S_0=e^{i\sigma}$: this is a path homotopy on an interval, with
the initial endpoint fixed; its terminal output is identity and is
therefore unaffected by the terminal conjugator. Both $B_\tau$ and
the B32 disk begin at the same based circle $a_s$ and end at identity.
Their closed difference is a map $S^2\to SU(2)$ and has a based
null since $\pi_2SU(2)=0$. A relative disk comparison therefore
identifies the two complete exterior nulls. Carry it into the \eqref{p1b:eq:natural-oj-6}
collar by the same literal formula; its initial cap and all lower
slices remain fixed. Constant conjugation changes no based cap class.
Finally $T_{h\sigma}\Gamma_B$ is an actual homotopy of cap maps,
since each translated value is divided by its actual moving
circle-base value and $\Gamma_B=I$ on $\partial B$.
For a literal check of the root words, apply the already constructed
common actual lower/frame comparisons OK.6--OK.7 in both localized
caps. Write $q=Q_s^*H^3Q_s$ and $k(t,z)=Q_s^*K(t,z)Q_s$ in
the intrinsic block. At the rebased value $a=1$, the
two retained cuts become exactly
\[
 W(\pi t,z)=A_t^{-1}k(t,z)q(z)^{-2},\qquad
 W(\pi(t+1),z)=A_t^{-1}k^\sigma(t,z).                \tag{CI.5a}\label{p1b:eq:natural-ci-5a}
\]
Their common midpoint is $q^{-2}$ and their two circle endpoints are
identity. Their actual spatial boundary is
$A_t^{-1}gA_tg^{-1}$, closed with the same \eqref{p1b:eq:natural-oj-6} fixed-axis collar.
These are the literal old-cap words retained in \eqref{p1b:eq:natural-ol-5} and subsequent
whole-cap reductions. The first-cut squared-root carry has therefore
not disappeared in the native comparison. The OG cap multipliers are
exactly $B_0E'_1$: its direction multiplier and $L$ are identity on
this real ball, while \eqref{p1b:eq:natural-oj-8} removes $B_0$ by its specified whole
two-support null. Consequently the old annular contribution of OH
is outside this localized cap comparison and introduces no extra
characteristic summand in \eqref{p1b:eq:natural-ci-5a}.
Thus every baseline and null-disk track in CI.4--CI.5 is carried,
and
\[
 {}[\Gamma_{\mathrm{nat}}]=[\Gamma_B]=[W_{\mathrm{old}}]=0.
                                                               \tag{CI.6}\label{p1b:eq:natural-ci-6}
\]
This does not replace the native cap by the characteristic norm alone;
the old cap value includes its actual canonical squared-root carry.

\paragraph{Paste with the actual projective null.}
Use $H_\tau$ on the original normal ball and keep \eqref{p1b:eq:natural-ob-16} on its
compressed collar. Their common initial boundary is $E_Q$. The
usual collar pasting of two nulls is literal multiplication of their
actual values: on the attaching ordered cover it is
\[
 Z(\tau,s,y,u)=\mathcal N_Q(\tau,s,y,p(u))
                           N_+(\tau,s,y,u)^{-1}.       \tag{CI.7}\label{p1b:eq:natural-ci-7}
\]
This is exactly \eqref{p1b:eq:natural-jt-4}, with its complete northern and southern formulas
\eqref{p1b:eq:natural-jt-5} and exact seam \eqref{p1b:eq:natural-jt-6}. The northern meridian is its actual identity
point; no southern meridian is collapsed. Shrink the real cap and this
outer collar into disjoint closed supports using radial collar
coordinates, holding their actual boundary identities fixed. The
endpoint is $\Gamma_{\mathrm{nat}}$ times the literal \eqref{p1b:eq:natural-ci-7} map.
This constructs an intrinsic homotopy on the whole original
$S^1\wedge S^3\wedge(D^3/S^2)$; no isolated hemisphere is
assigned a coefficient. It carries the middle frame \eqref{p1b:eq:natural-ci-1} and the
two divisions CI.2--CI.4 throughout. Therefore
\[
 h_t\mu=[\Gamma_{\mathrm{nat}}]+[Z]=[Z],\qquad
 r_J\mu=\mu+[Z],\qquad\zeta_{\rm nat}=r_J.                    \tag{CI.8}\label{p1b:eq:natural-ci-8}
\]
The last two equalities apply the already proved whole original
\eqref{p1b:eq:natural-pg-16} equation. Only evaluation of the entire paired map \eqref{p1b:eq:natural-ci-7} may
assign either absolute bit; the isolated primitive $K$ is insufficient.



\subsection{The reduced actual neck family retains the exotic clutch}
\label{p1b:sec:natural-reduced-neck-clutch}

Retain CI.1--CI.3 on the actual punctured normal domain
$M\setminus\operatorname{int}B$, where $M=S^3_y\times D^3_n$,
and any fixed actual normal boundary point $u_*$. In particular the
proof applies both to the northern meridian $j$ of \eqref{p1b:eq:natural-jt-7} and to the
equator vertex $k=u(1,0)$ used in FS and TN. These are evaluations
at actual points of the same normal sphere; the original normal
framing and source coordinates remain fixed. Write $c_\rho$ for
that actual circle-independent based frame and $c$ for its normal
boundary value. Define the entire reduced neck map by
\[
 H(y)=c(y,u_*),\qquad h(y,u)=c(y,u)H(y)^{-1}.
                                                               \tag{CN.1}\label{p1b:eq:natural-cn-1}
\]
Both $h(1,u)$ and $h(y,u_*)$ are identity; thus $h$ is literally
a based map $S^3\wedge S^2=S^5\to SU(2)$. We identify its whole
class, rather than treating the degree-two slice $H$ as a clutch.

Extend $H(y)$ over $M$ by independence of $n$. The actual normal-only
diagonal $C_b(\rho,u)$ from \eqref{p1b:eq:natural-ci-3} is defined on all of $D^3$, including
its polar origin, because the actual sphere-base slice misses $B$.
Consequently
\[
 \widehat c(y,n)=c_\rho(y,n)H(y)^{-1}
                                                               \tag{CN.2}\label{p1b:eq:natural-cn-2}
\]
is an actual map on the punctured domain and is identity on
$A=(S^3\times\{u_*\})\cup(\{1\}\times D^3)$.
This $A$ is disjoint from $B$. The first part is a boundary slice;
its inclusion into the punctured product carries the parameter
three-cell, and its normal-only part is contractible relative to
that slice. The inclusion of $A$ into the full $M$ is a homotopy
equivalence: normal contraction onto $u_*$ retracts $M$ to its
parameter slice, and $A$ retracts to that same slice. Thus $M/A$
is contractible. The inclusion of $A$ into the punctured product
induces the degree-three homology isomorphism. Excision of the
interior six-ball consequently gives a simply connected punctured
quotient whose only reduced homology is $H_5=\mathbb Z$. Iterated
Hurewicz shows that it is four-connected, and a generator $S^5$
map is a homology equivalence between simply connected CW spaces,
hence a homotopy equivalence by Whitehead. Therefore
\[
 (M\setminus\operatorname{int}B)/A\simeq S^5.
                                                               \tag{CN.3}\label{p1b:eq:natural-cn-3}
\]
A relative CW decomposition realizes this equivalence; it is a source
comparison retaining $A$, not a trivialization of the missing frame
on $B$. The outer boundary quotient
$(S^3\times S^2)/((S^3\times\{u_*\})\cup(\{1\}\times S^2))$
is the actual $S^5$ smash in \eqref{p1b:eq:natural-cn-1}. Its inclusion has degree of
absolute value one onto \eqref{p1b:eq:natural-cn-3}. The inner boundary $\partial B$
also has unit degree, with the boundary orientation sign. This
follows directly from the relative fundamental chain of the
punctured product, whose two boundary components have opposite
orientations. Four-connectivity implies that these two source maps
agree up to that degree sign by the Hurewicz isomorphism.

On the actual root sphere use $g=Q_s^*V^\#$ and the actual cap
frame comparison $e=Q_s^*F_Y$. Then the literal inner value of
\eqref{p1b:eq:natural-cn-2} is
\[
 \widehat c|_{\partial B}=C_b^{-1}g^{-1}eH^{-1}.
                                                               \tag{CN.4}\label{p1b:eq:natural-cn-4}
\]
The three remaining factors $C_b^{-1}$, $e$ and $H^{-1}$
are defined over all of the actual ball: the first is normal-only,
the second compares two actual frames defined there, and the third
is the parameter-only extension already specified. Contract each
restriction radially in its own actual ball, in the displayed order,
retaining its moving meridian value by division. The resulting
constants can be removed by paths in $SU(2)$. This is a whole-map
comparison of \eqref{p1b:eq:natural-cn-4} to $g^{-1}$, with its actual cap frame retained.
The based meridian of $g$ is the exact identity value in \eqref{p1b:eq:natural-pg-1}.
Combining CN.2--CN.4 and the two source degrees therefore gives
\[
 {}[h]=\pm[g^{-1}]=\eta_3\eta_4\in\pi_5SU(2).
                                                               \tag{CN.5}\label{p1b:eq:natural-cn-5}
\]
Both signs and inversion agree on this order-two class. In particular
the literal equator-vertex reduction
$c(y,u)c(y,k)^{-1}$ has the same proved class $\eta_3\eta_4$.
This does
not provide a null of $h$, an extension of the original $c$ over
the missing cap, or an evaluation of the complete paired \eqref{p1b:eq:natural-jt-7} map.
It supplies the actual reduced neck operand for a further relative
comparison, with its degree-two lower slice and cap frame explicitly
retained.



\subsection{Independence of the paired comparison from its seam null}
\label{p1b:sec:paired-joining-null-independence}

Retain exactly \eqref{p1b:eq:natural-jt-2}--\eqref{p1b:eq:natural-jt-7}, including the fixed-axis ratio
\(\mathcal A\), the actual northern centre as sphere basepoint, and
both lower base faces.  Let \(\Pi_0,\Pi_1\) be any two intrinsic
nullhomotopies of the same corrected seam \(X\), each relative to
both \(\tau\) faces, \(s=0\), \(y=1\), and
\(\theta=0,2\pi\), and each starting at \(X\) and ending at
identity.  Write \(Z_i\) for the complete two-hemisphere map
\eqref{p1b:eq:natural-jt-5} built with \(\Pi_i\) and the same \(\mathcal A\).

\begin{lemma}[Paired seam-null independence]
\label{p1b:lemma:paired-joining-null-independence}
The complete based maps \(Z_0,Z_1:S^7\to SU(2)\) have the
same homotopy class.
\end{lemma}
\begin{proof}
Put \(V=\Pi_1\Pi_0^{-1}\).  This ratio is identity at
both endpoints of the null parameter, on every retained base face,
and on \(\theta=0,2\pi\).  The pointwise ratio of the two
complete maps is, exactly,
\[
 Z_1 Z_0^{-1}=
 \begin{cases}
 V(\chi(r),\tau,s,y,\theta),&\text{northern hemisphere},\\
 dV(\chi(r),\tau,s,y,\theta)d,&\text{southern hemisphere}.
 \end{cases} \tag{PN.1}\label{p1b:eq:natural-pn-1}
\]
The fixed-axis factor \(\mathcal A\) cancels pointwise in the
second expression.  In particular \eqref{p1b:eq:natural-pn-1} is identity on the actual
equator, both inner collar boundaries and the northern sphere
basepoint.  The two annular contributions can therefore be separated
by a based pinch.  Since \(V\) is identity along its angular
meridian, collapsing each annulus boundary and that meridian gives
\[
 S^1_\tau\wedge S^1_s\wedge S^3_y
      \wedge S^1_\theta\wedge S^1_{\rm null}=S^7. \tag{PN.2}\label{p1b:eq:natural-pn-2}
\]
The northern contribution is its class \(\xi\); the southern
contribution is the same class up to its source orientation and
constant target conjugation.  Constant conjugation is based
homotopic to identity.  Thus the full difference is \(\xi\pm\xi\).
Because \(\pi_7 SU(2)=C_2\), it is zero for either orientation.
This argument collapses meridians only in the difference map \eqref{p1b:eq:natural-pn-1},
where they are actually identity.  It does not collapse the
nonconstant meridian of either original map \(Z_i\).
\end{proof}

\paragraph{An exact identity for the retained moving lower value.}
For \(B_\theta=b_\tau(s,\theta)\), fix an axis parameter
\(\kappa\) and put
\[
 D_\kappa(\theta)=B_\theta B_{\kappa\pi}^{-1},\qquad
 e_\kappa=B_0 B_{\kappa\pi}^{-1},\qquad
 D_0(\theta)=B_\theta B_0^{-1}.
                                                        \tag{PN.3}\label{p1b:eq:natural-pn-3}
\]
These are the actual fixed-axis changes occurring in \eqref{p1b:eq:natural-jt-8}.
Their own angular lower value is \(D_\kappa(0)=e_\kappa\).
Full-matrix multiplication gives
\[
 D_\kappa(\theta)D_\kappa(0)^{-1}=D_0(\theta).
                                                        \tag{PN.4}\label{p1b:eq:natural-pn-4}
\]
For any common intrinsic operand \(c\), use the same commutator
\(\mathcal C(B,c)=B^{-1}c^{-1}Bc\).  Since
\(D_\kappa(\theta)=D_0(\theta)e_\kappa\), its complete
commutator division obeys
\[
 \boxed{\mathcal C(D_\kappa(\theta),c)
          \mathcal C(e_\kappa,c)^{-1}
       =e_\kappa^{-1}\mathcal C(D_0(\theta),c)e_\kappa.}
                                                        \tag{PN.5}\label{p1b:eq:natural-pn-5}
\]
Indeed the left side is
\(e_\kappa^{-1}D_0^{-1}c^{-1}D_0e_\kappa c
 c^{-1}e_\kappa^{-1}ce_\kappa\), whose adjacent factors
cancel to the right side.  Thus after the actual circle-diagonal
separation of \eqref{p1b:eq:natural-ob-14}, when the frame operand is exactly
\(c(y,k)\), this part of the fixed-axis transport is target
conjugation with the moving angular lower division retained.  It is
not a precomposition of the \(D_0\) operand.

These positive lemmas do not evaluate \eqref{p1b:eq:natural-jt-7}.  In particular \eqref{p1b:eq:natural-pn-5}
is asserted for a common frame operand: applying it before the
actual diagonal separation, with unequal \(c(y,u_\theta)\)
and \(c(y,k)\), would omit a frame term.  The complete paired homotopy in (TN), (FS), (BG) and (AJ) carries
all nine legs of \eqref{p1b:eq:natural-ob-15}, the corrected \(K\) factors and \eqref{p1b:eq:natural-jt-8}
together. Its absolute coefficients are \eqref{p1b:eq:natural-aj-5}.



\subsection{The paired joining comparison is independent of the frame interior}
\label{p1b:sec:natural-joining-interior-frame-kernel}

Use the complete original paired map \eqref{p1b:eq:natural-jt-4}--\eqref{p1b:eq:natural-jt-7}.  This comparison
retains its southern fixed-axis factor; it does not collapse either
southern meridian.  Let $c^0,c^1$ be two actual based ordered-hemisphere
frame families, with the same complete equator value,
\[
 c^0(y,u_\theta)=c^1(y,u_\theta),\qquad
 c^i(1,u)=I.                                             \tag{IK.1}\label{p1b:eq:natural-ik-1}
\]
Extend each to the southern ordered hemisphere by the same actual deck
conjugation $d$.  Retain the same actual equator frame homotopy $H_3$,
the same cap $Q_+$, the same correction $R$, and the same full seam null
$\Pi$.  They are legitimate for both families because \eqref{p1b:eq:natural-ik-1} preserves
the entire seam \eqref{p1b:eq:natural-ob-5}, its corrected endpoint and its specified null.
The two resulting complete paired maps are denoted $Z^0,Z^1$.

Write $b_\kappa=b_\tau(s,\kappa\pi)$, $b_+=b_0$, $b_-=b_1$,
and $\mathcal C(B,c)=B^{-1}c^{-1}Bc$.  For an arbitrary actual frame
value $c$ the literal universal path
\[
 L_\kappa(c)=\mathcal C(b_\kappa,c)
                         \mathcal C(b_+,c)^{-1},\qquad
 L_0(c)=I,\quad L_1(c)=\mathcal A(c)                     \tag{IK.2}\label{p1b:eq:natural-ik-2}
\]
retains the actual boundary-axis rotation.  Both $\tau$ faces and the
$s=0$ face are identity: at $\tau=0$ both circle factors are the same
$a_s$, and at $\tau=1$ both are identity.  The $y=1$ face is identity
as well.  Hence
\[
 B_\kappa(y,u)=L_\kappa(c^1(y,u))
                         L_\kappa(c^0(y,u))^{-1}          \tag{IK.3}\label{p1b:eq:natural-ik-3}
\]
is an exact relative path from identity to
$\mathcal A(c^1)\mathcal A(c^0)^{-1}$.  On the complete equator it is
identity for every $\kappa$, by \eqref{p1b:eq:natural-ik-1}.  No null of either individual
frame family is used.

On the northern hemisphere $Z^1=Z^0=\Pi(\chi(r),\theta)$.
On the southern hemisphere put $P=d\Pi(\chi(r),\theta)d$.  The actual
product order of \eqref{p1b:eq:natural-jt-5} gives
\[
 Z^1(Z^0)^{-1}
   =P\mathcal A(c^1)\mathcal A(c^0)^{-1}P^{-1}.           \tag{IK.4}\label{p1b:eq:natural-ik-4}
\]
Define the whole closed relative homotopy to be identity in the north
and $PB_\kappa P^{-1}$ in the south.  Equation \eqref{p1b:eq:natural-ik-3} makes their
equator seam literally identity at every parameter.  It is continuous
at the southern pole because $P=I$ on the central half-disc, and the
actual frame values have the prescribed polar continuity.  The
northern centre remains the original based point, and all original
time and parameter-sphere faces remain identity.  Its initial whole
map is identity; its terminal whole map is \eqref{p1b:eq:natural-ik-4}.  Therefore
\[
 \boxed{[Z^1]=[Z^0]\quad\hbox{in }\pi_7SU(2).}          \tag{IK.5}\label{p1b:eq:natural-ik-5}
\]

This kernel evaluates the complete change, including the southern
factor and its nonconstant meridian.  It permits a different interior
extension of the same actual projective one-cell frame path.  It does
not permit changing that path or assigning its relative
$\pi_4SU(2)$ class.  In particular degree two of one frame slice is
not a substitute for the retained equator path.



\subsection{A descended moving-cap reference and its complete ordered comparison}
\label{p1b:sec:natural-joining-moving-cap-reference}

This evaluates a specified reference frame.  Its use for the original
frame still requires the complete relative frame comparison.  Let
$q_\delta:S^3\to SU(2)$ be a cap map of signed degree
$\varepsilon\in\{1,-1\}$, identity off a cap within
distance $\delta<\pi/2$ of $1$, in the same actual quaternion axes.
For a unit imaginary quaternion $u$, define
\[
 f_+(y,u)=q_\delta(u^{-1}y),\qquad
 f_-(y,u)=d q_\delta((-u)^{-1}y)d,\qquad
 c_{\rm mod}(y,u)=f_+(y,u)f_-(y,u).                  \tag{MR.1}\label{p1b:eq:natural-mr-1}
\]
Their supports in $y$ are disjoint caps at $u$ and $-u$; both miss
$y=1$.  Hence $c_{\rm mod}(1,u)=I$, every fixed-$u$ slice has degree
$2\varepsilon$, and its exact ordering equation is
\[
 f_+(y,-u)=d f_-(y,u)d,\quad
 f_-(y,-u)=d f_+(y,u)d,\quad
 c_{\rm mod}(y,-u)=d c_{\rm mod}(y,u)d.              \tag{MR.2}\label{p1b:eq:natural-mr-2}
\]
The actual \eqref{p1b:eq:natural-ob-8} slice has signed degree $-2$ in the usual quaternion
orientation, as computed in \eqref{p1b:eq:natural-pj-7}--\eqref{p1b:eq:natural-pj-9}.  For its application choose
the inverse of the positive cap map, so $\varepsilon=-1$.  Every
formula below holds for either sign; inversion preserves the actual
cap support, sphere-base values and ordering equation.
In the last equality the two factors commute pointwise, since at
least one is identity at every source point.

Retain the actual circle disks $b_+=b_\tau(s,0)$ and
$b_-=b_\tau(s,\pi)=d b_+d$ of \eqref{p1b:eq:natural-jt-2}.  Put
\[
 N_{\rm proj}(\tau,s,y,u)=
    \mathcal C(b_+,f_+(y,u))
    \mathcal C(b_-,f_-(y,u)),\qquad
 N_D^{\rm mod}=\mathcal C(b_+,c_{\rm mod}).           \tag{MR.3}\label{p1b:eq:natural-mr-3}
\]
The first is a literal projective null: \eqref{p1b:eq:natural-mr-2} swaps its two factors
under $u\mapsto-u$, and deck conjugates their circle disks.  Their
disjoint supports permit the exact swap back.  Thus
$N_{\rm proj}(-u)=dN_{\rm proj}(u)d$.  At $\tau=0$ both disks are
$a_s$, so the same disjoint-support identity gives
$N_{\rm proj}(0)=\mathcal C(a_s,c_{\rm mod})$; at $\tau=1$ it is
identity.  The $s=0$ and $y=1$ faces are identity throughout.

The complete ordered comparison has the exact product order
\[
 N_{\rm proj}(N_D^{\rm mod})^{-1}
   =\mathcal C(b_-,f_-)\mathcal C(b_+,f_-)^{-1}
   =L_1(f_-),                                           \tag{MR.4}\label{p1b:eq:natural-mr-4}
\]
where $L_\kappa$ is the entire axis path \eqref{p1b:eq:natural-ik-2}.  No positive-cap
factor survives: the two positive-cap letters cancel after commuting
only disjoint-supported factors.  Equation \eqref{p1b:eq:natural-mr-4} is generally
nonidentity at a chosen normal basepoint $u_*$.  Retain its actual
lower value by defining
\[
 W_\kappa(\tau,s,y,u)=
  L_\kappa(f_-(y,u))
  L_\kappa(f_-(y,u_*))^{-1}.                            \tag{MR.5}\label{p1b:eq:natural-mr-5}
\]
Its initial map is identity; its terminal map is precisely the
complete comparison \eqref{p1b:eq:natural-mr-4}, divided by its actual normal-base value.
It is identity at $u=u_*$, on both $\tau$ faces, and at $s=0,y=1$.
Thus it is a relative intrinsic homotopy of complete based maps
\[
 S^1_\tau\wedge S^1_s\wedge S^3_y\wedge S^2_u=S^7
       \longrightarrow SU(2),\qquad
 \boxed{[W_1]=0.}                                      \tag{MR.6}\label{p1b:eq:natural-mr-6}
\]
In particular neither the original axis change nor its actual normal
lower value has been deleted.  If a projective null literally equal
to $N_D^{\rm mod}$ at $[u_*]$ is needed, use the exact path
$L_\kappa(f_-(y,u_*))$ to align that vertex by homotopy extension,
keeping both time ends and both parameter-base faces.  The vertex
difference is null by this supplied path, rather than an arbitrary
$\pi_5SU(2)$ choice.

\paragraph{The full-cover reduced frame clutch of the same reference.}
Choose any fixed unit imaginary $u_*$, and define the exact based contraction in $S^3$
\[
 F_t(u)=\cos(\pi t/2)u+\sin(\pi t/2),\qquad
 U_t(u)=F_t(u)F_t(u_*)^{-1}u_*.
                                                               \tag{MR.7}\label{p1b:eq:natural-mr-7}
\]
Every value has norm one; $U_0(u)=u$, $U_1(u)=u_*$ and $U_t(u_*)=u_*$.
The moved cap maps keep their actual sphere-base values:
\[
 \begin{aligned}
 f_{+,t}(y,u)&=q_\delta(U_t(u)^{-1}y)
                            q_\delta(U_t(u)^{-1})^{-1},\\
 f_{-,t}(y,u)&=d q_\delta((-U_t(u))^{-1}y)
                            q_\delta((-U_t(u))^{-1})^{-1}d,\\
 c_t&=f_{+,t}f_{-,t},\\
 h_t^{\rm mod}(y,u)&=c_t(y,u)c_t(y,u_*)^{-1}.
 \end{aligned}                                                \tag{MR.8}\label{p1b:eq:natural-mr-8}
\]
At $t=0$ the retained lower factors are identity, because imaginary
centres miss the based point, so $c_0=c_{\rm mod}$.  At every $t$
the $y=1$ and $u=u_*$ faces of $h_t^{\rm mod}$ are literally identity.
At $t=1$ the two operands are independent of $u$, hence
$h_1^{\rm mod}=I$.  These homotopies need not descend projectively:
they are a calculation of the complete ordered-cover frame clutch,
on precisely $S^3\wedge S^2=S^5$.  Consequently
\[
 \boxed{[c_{\rm mod}(y,u)c_{\rm mod}(y,u_*)^{-1}]=0
             \quad\hbox{in }\pi_5SU(2).}             \tag{MR.9}\label{p1b:eq:natural-mr-9}
\]
The original frame has not been substituted by this reference.  Its
relative projective one-cell path and complete change of the paired
joining comparison are identified in (TN), (FS), (BG) and (AJ)
before \eqref{p1b:eq:natural-mr-6} contributes to the
absolute natural coefficient.



\subsection{The relative two-disc norm with its actual deck rectangle}
\label{p1b:sec:natural-joining-twisted-cell-norm}

Let $P=S^p$ with $p\ge3$, and let $k:P\times S^2\to SU(2)$ obey
\[
 k(*,u)=I,\quad k(z,u_0)=I,\quad
 k(z,-u)=d k(z,u)d,\qquad d=\operatorname{diag}(1,-1).
                                                               \tag{TN.1}\label{p1b:eq:natural-tn-1}
\]
Choose $u_0=u(1,0)=k$ in the actual ordered northern hemisphere
\eqref{p1b:eq:natural-ob-2}.  The first equator half is the literal based map
\[
 m(z,\theta)=k(z,u_\theta),\quad0\le\theta\le\pi,
 \quad m(z,0)=m(z,\pi)=I,
 \qquad m:P\wedge S^1=S^{p+1}\to SU(2).             \tag{TN.2}\label{p1b:eq:natural-tn-2}
\]
The second half is exactly $d m(z,\theta-\pi)d$; it is not replaced
by $m$ without its target-conjugation rectangle.  The northern map
is a null of the whole boundary $m*(dmd)$.  Its based null form can
be read directly using the contraction of the unit disk to its actual
boundary basepoint, $F_h(e^{i\theta})=(1-h)e^{i\theta}+h$.
This contraction fixes $\theta=0$ throughout; hence the actual lower
meridian is retained in the null.  The southern disk is the actual
deck conjugate, with the antipodal boundary parametrization.

Put
\[
 q_t=\operatorname{diag}(e^{\pi it},1),\quad t\in\mathbb R,
 \qquad q_0=I,\quad q_1=d.                            \tag{TN.3}\label{p1b:eq:natural-tn-3}
\]
Transport the second boundary half from $dmd$ to $m$ along
$q_{1-h}m q_{1-h}^{-1}$, keeping the first half fixed.  At both
ends of that half $m=I$, so this is a based homotopy of the entire
boundary loop.  Choose one northern homotopy extension $N_h$ of this
entire boundary homotopy, and use \emph{exactly} $dN_hd$ in the southern
disk.  Its terminal northern disk is denoted $N'$ and has boundary
$m*m$.  There is no independently chosen southern homotopy extension.

The southern equator parameter is $\theta+\pi$.  Consequently
on the first half the northern value is $m$, while the
southern value is $\operatorname{Ad}_{q_{2-h}}m
=\operatorname{Ad}_{q_{-h}}m$.  On the second half the
northern value is $\operatorname{Ad}_{q_{1-h}}m$, and the southern
value is $\operatorname{Ad}_{q_1}m$.  Insert the following literal
collars, parameterized from north to south by $v\in[0,1]$:
\[
 \begin{array}{c|c}
  \hbox{first half}&\operatorname{Ad}_{q_{-hv}}m\\
  \hbox{second half}&\operatorname{Ad}_{q_{1-h+hv}}m.
 \end{array}                                                \tag{TN.3a}\label{p1b:eq:natural-tn-3a}
\]
Each endpoint is exactly its adjoining disk value.  At $h=0$ the
collars are constant in $v$, so inserting their source collars is a
relative source comparison of the original map.  At every half-circle
end $m=I$, so the collars paste there literally.  They retain the
nonconstant normal meridians and all parameter-base faces.

At $h=1$ the northern disk is $N'$, the southern disk is $dN'd$,
and the two collar phases are respectively $0\to-1$ and $0\to1$.
Now remove $d$ from the \emph{entire} southern disk by
$\operatorname{Ad}_{q_{1-t}}N'$, $0\le t\le1$.  Continue the first
collar endpoint as $-1-t$, and the second as $1-t$; that is, use
\[
 \begin{array}{c|c}
  \hbox{first half}&\operatorname{Ad}_{q_{-(1+t)v}}m\\
  \hbox{second half}&\operatorname{Ad}_{q_{(1-t)v}}m.
 \end{array}                                                \tag{TN.3b}\label{p1b:eq:natural-tn-3b}
\]
Their southern endpoints are the same target conjugation because the
phases differ by two.  At $t=1$ both entire disks are \emph{literally}
$N'$ with the original antipodal shift.  The first collar is one
negative full winding $0\to-2$; the second is constant.  Thus no
reversed-disk null or independent corner path is hidden in the
transport.

The two residual target-conjugation rectangles are literally
\[
 \begin{aligned}
 R_+(t,w)&=q_t m(w)q_t^{-1}m(w)^{-1},\\
 R_-(t,w)&=d q_t m(w)q_t^{-1}d m(w)^{-1}.
 \end{aligned}                                                \tag{TN.4}\label{p1b:eq:natural-tn-4}
\]
Their full edge data are
\[
 R_+(0)=I,\quad R_+(1)=dmdm^{-1}=R_-(0),\quad R_-(1)=I.
                                                               \tag{TN.5}\label{p1b:eq:natural-tn-5}
\]
Thus neither rectangle is treated as closed separately.  Pasting
them along the retained middle edge gives the complete commutator
\[
 q(t)m(w)q(t)^{-1}m(w)^{-1},\quad
 q(t)=\operatorname{diag}(e^{2\pi it},1),
 \qquad S^1\wedge S^{p+1}=S^{p+2}.                   \tag{TN.6}\label{p1b:eq:natural-tn-6}
\]
All parameter-base and equator-base faces remain identity.  This is
the source-identified adjoint operation $A_*[m]=\eta_3E[m]$ of
\eqref{p1b:eq:rt-23-2}, including its nonconstant halfway seam.

To separate the final collar from the remaining piece, compare the
whole final map with the same two disks and constant seam collars.
Their whole pointwise ratio is identity off the first collar, and on
it is exactly $q_{-2v}m q_{-2v}^{-1}m^{-1}$.  This is the complete
negative-winding form of \eqref{p1b:eq:natural-tn-6}; both collar edges are identity in the
ratio.  The remaining two-disk piece is the
ordinary, untwisted pullback of the based map on $P\wedge\mathbb{RP}^2$
whose one-cell is $m$ and whose two-cell is the just-retained null
of $m*m$.  It therefore factors through
\[
 \Sigma^p p:S^{p+2}\longrightarrow\Sigma^p\mathbb{RP}^2.
                                                               \tag{TN.7}\label{p1b:eq:natural-tn-7}
\]
This map is null.  The G21 technical comparison gives Atiyah's stable
splitting and hence the stable null of $p$.  Here
$\Sigma^p\mathbb{RP}^2$ is $p$-connected; Freudenthal puts its
$(p+2)$nd group in the stable range for $p\ge3$.  Any change of
either disk extension gives two deck-paired copies of a class in
$\pi_{p+2}SU(2)$, so it contributes only its doubled value.

Consequently this complete transport identifies, up to the common
orientation of the first collar,
\[
 {}[k]=\pm A_*[m]=\pm\eta_3E[m].                         \tag{TN.8}\label{p1b:eq:natural-tn-8}
\]
For the two applications relevant here, the paired extension changes
vanish: $\pi_5SU(2)=C_2$ and $\pi_7SU(2)=C_2$.  Thus for $p=3$
the relative frame-edge value $m=\eta_3$ gives
$[k]=\eta_3\eta_4$, while for $p=5$ every relative frame-edge
$m\in\pi_6SU(2)$ gives $[k]=0$ by \eqref{p1b:eq:rt-23-3}.

This proves the relative norm for a literal zero vertex.  An arbitrary
nonzero vertex cannot be removed by this argument: its correlated
lower transport can have the characteristic secondary contribution.
The lemma does not assign the original joining coefficient without
the complete original-to-reference comparison.



\subsection{A relative frame product splits the complete paired comparison}
\label{p1b:sec:natural-joining-frame-product-splitting}

Let $c^0,c^1$ be actual based frame families with the same pointwise
normal-vertex slice $H(y)$.  Retain their ordering equation
$c^i(y,-u)=dc^i(y,u)d$ and sphere-base value $c^i(1,u)=I$.
Put
\[
 k(y,u)=c^1(y,u)c^0(y,u)^{-1},\quad
 k(y,u_0)=I,\quad k(1,u)=I,\quad
 k(y,-u)=d k(y,u)d.                                   \tag{FS.1}\label{p1b:eq:natural-fs-1}
\]
The full product order is $c^1=k c^0$.  The actual comparison for
$c^0$ includes its complete equator frame homotopy, both actual caps,
the correction $R$, and its projective seam null.  For $k$ the vertex
is literally identity; its cap correction is therefore identity.
Its entire seam has the specified null obtained from \eqref{p1b:eq:natural-ob-5} by
removing its conjugator and then using the actual \eqref{p1b:eq:natural-ob-14} circle
diagonal separation: the first separated circle has $k(y,u_0)=I$,
and the second has $D=I$.  Every prescribed based face is fixed.

Choose the standard based degree-one sphere compressions
$a_t^\pm:S^3\to S^3$, $a_0^\pm=\operatorname{id}$,
$a_t^\pm(1)=1$, whose terminal nonbase supports are disjoint
three-balls.  They are precisely the stereographic formula \eqref{p1b:eq:natural-pg-6}
in dimension three, with centres $\pm2e_1$.  Use the whole relative
frame homotopy
\[
 c_t(y,u)=k(a_t^-(y),u)c^0(a_t^+(y),u).               \tag{FS.2}\label{p1b:eq:natural-fs-2}
\]
This preserves the complete ordering equation and sphere-base face.
Its moving normal-vertex value is the actual
$H(a_t^+(y))$; it is not replaced by a constant.  Compose \emph{all}
reference vertex data, including the full $H_3$, both caps and
$Q_+$ in $K$ and $R$, with the same $a_t^+$.  This is a based
source homotopy of the complete reference construction.

The initial corrected seam of the whole product consequently moves
through a homotopy $X_t$ with every original time, sphere and circle
face fixed.  Its actual null is carried by homotopy extension.  More
explicitly, the null of $X_t$ can first follow the full path $X_r$
from $r=t$ back to $r=0$, and then follow the original seam null.
This retains the changing initial endpoint instead of attaching a
null with a different start.  At the terminal support-separated
configuration one may use the product of the reference and $k$
seam nulls: it has that same literal initial seam.  The complete
paired seam-null independence (PN) evaluates the change between
these two full nulls as zero.  Neither hemisphere is assigned a
coefficient separately.

At $t=1$, write $c_+=c^0(a_1^+y,u)$ and $k_-=k(a_1^-y,u)$.
At every $y$ at least one is identity.  Every corresponding reference
letter, including $R_+$ and $\Pi_+$, is supported in the first
ball; every corresponding $k$ letter and $\Pi_-$ is supported in the
second.  Thus the literal entire words satisfy
\[
 \begin{aligned}
 N_1^{\rm tot}&=N_1^+N_1^-,\\
 N_\pm^{\rm tot}&=N_\pm^+N_\pm^-,\\
 R^{\rm tot}&=R_+,\\
 X^{\rm tot}&=X_+X_-,\\
 \Pi^{\rm tot}&=\Pi_+\Pi_-,\\
 \mathcal A^{\rm tot}&=\mathcal A_+\mathcal A_-.
 \end{aligned}                                                \tag{FS.3}\label{p1b:eq:natural-fs-3}
\]
All commutations used here involve only these disjoint supports.
In particular the southern fixed-axis word is present in the last
line.  Substitution into both actual hemisphere formulas \eqref{p1b:eq:natural-jt-5}
and their complete equator equation \eqref{p1b:eq:natural-jt-6} gives
\[
 Z^{\rm tot}=Z_+Z_-\quad\hbox{on the entire original }S^7,
 \qquad
 \boxed{[Z(c^1)]=[Z(c^0)]+[Z(k)].}                   \tag{FS.4}\label{p1b:eq:natural-fs-4}
\]
The source compressions have degree one and are supplied whole-map
homotopies, so the supported factors retain their original classes.
Equation \eqref{p1b:eq:natural-fs-4} includes the complete reference cap calibration and
correction, the actual south meridian, and the moving normal-vertex
slice.  It does not infer linearity by discarding cross words.

For the moving-cap reference \eqref{p1b:eq:natural-mr-1}, the explicit axis-null
\eqref{p1b:eq:natural-mr-4}--\eqref{p1b:eq:natural-mr-6} gives its ordered comparison zero.  The vertex
difference between that descended projective null and $N_D$ has the
supplied entire null $L_\kappa(f_-(y,u_0))$.  Align it by homotopy
extension, retaining its full moving lower value.  Any other
projective null with that same literal vertex differs by a section
with zero vertex on the parameter sphere $S^5$.  The complete
two-disc norm \eqref{p1b:eq:natural-tn-1}--\eqref{p1b:eq:natural-tn-8}, with $p=5$, sends its one-cell class
through $A_*\pi_6SU(2)=0$.  Hence the original paired comparison
constructed for this reference also has class zero.  This use of
null replacement is restricted to the supplied null vertex
difference; no assertion is made about arbitrary nonzero
$\pi_5SU(2)$ vertex changes.

The full product comparison therefore reduces to the complete
zero-vertex $Z(k)$.  Its equator loop is the actual relative
$m:S^4\to SU(2)$; the lower frame norm (TN), with $p=3$, supplies
$[k]=A_*[m]$.  The complete two-bigon calculation \eqref{p1b:eq:natural-bg-1}--\eqref{p1b:eq:natural-bg-11}
then evaluates this whole zero-vertex map after the actual relative
half-edge has been determined below.

\paragraph{Signed vertex alignment and the actual relative edge.}
For the actual application choose the reference cap of degree $-1$
as stated in MR, so its vertex slice has signed degree $-2$, equal
to the actual $C_2$ endpoint of \eqref{p1b:eq:natural-ob-8}.  The actual $H_3$ supplies
the based comparison from the actual slice $H$ to $C_2$.
Since $\pi_3SU(2)=\mathbb Z$ is classified by signed degree,
the reference vertex slice admits a based comparison to this same
$C_2$, and hence to $H$.  Extend that vertex homotopy through the
actual ordering-line bundle by homotopy extension.  Transport the
reference projective null, its actual moving normal-base division,
and all $H_3$, cap, $K$, $R$ and $\Pi$ data along the whole
comparison.  At the aligned endpoint the vertex slice is literally
$H$ and one may choose the same actual $H_3$ as for the original
frame.  This is a based homotopy of the full reference construction;
its complete paired coefficient and its full-cover reduced clutch
remain the zeros computed in MR.  The alignment uses signed degree,
not only its absolute value.

At the aligned endpoint put
$h_i(y,u)=c^i(y,u)H(y)^{-1}$.  Their exact relation is
$h_1=k h_0$; all three are based on $S^3\wedge S^2=S^5$.
The complete actual reduced neck clutch \eqref{p1b:eq:natural-cn-5} is
$[h_1]=\eta_3\eta_4$, while MR gives $[h_0]=0$.
Apply the source-identified zero-vertex norm TN with $p=3$:
\[
 {}[k]=[h_1]-[h_0]=\eta_3\eta_4=A_*[m],\qquad
 \boxed{[m]=\eta_3\in\pi_4SU(2)=C_2.}               \tag{FS.5}\label{p1b:eq:natural-fs-5}
\]
The last step uses the established isomorphism
$A_*:\pi_4SU(2)\to\pi_5SU(2)$ from \eqref{p1b:eq:rt-23-3}.
This calculates the actual relative half-edge.  It does not assume
that an arbitrary degree-two frame can replace the original frame.



\subsection{The actual two bigons for a zero-vertex frame}
\label{p1b:sec:natural-joining-zero-vertex-bigons}

Let $k(y,u)$ be an actual ordered frame family with
\[
 k(1,u)=I,\qquad k(y,-u)=dk(y,u)d,\qquad
 k(y,k)=I,\qquad d=\operatorname{diag}(1,-1).
 \tag{BG.1}\label{p1b:eq:natural-bg-1}
\]
Use exactly the OB northern hemisphere $u(r,\theta)$ and put
$m(y,\theta)=k(y,u_\theta)$.  In particular
$m(y,0)=m(y,\pi)=m(y,2\pi)=I$, and
$m(y,\theta+\pi)=dm(y,\theta)d$.  Write
\[
 B_\alpha=b_\tau(s,\alpha),\quad B_0=b_\tau(s,0),\quad
 D_\alpha=B_\alpha B_0^{-1},\quad
 \mathcal C(B,c)=B^{-1}c^{-1}Bc.
 \tag{BG.2}\label{p1b:eq:natural-bg-2}
\]
The zero vertex means that its cap primitive and correction $R$ are
identity. Its full corrected seam is therefore
$X_k=\mathcal C(B_\theta,m)\mathcal C(B_0,m)^{-1}$.
Choose the following entire null of that exact seam:
\[
 \Pi_k(v,\tau,s,y,\theta)
 =\mathcal C(B_{(1-v)\theta},m(y,\theta))
                  \mathcal C(B_0,m(y,\theta))^{-1}.
 \tag{BG.3}\label{p1b:eq:natural-bg-3}
\]
It starts at $X_k$ and ends at identity.  It is identity on both
$\tau$ faces since $B_\alpha=a_s$ at $\tau=0$ and $B_\alpha=I$
at $\tau=1$, independently of $\alpha$. It is identity at $s=0$
and $y=1$. At $\theta=0,\pi,2\pi$, it is identity because
$m=I$. Thus the nonperiodic intermediate axis $(1-v)\theta$ does
not invalidate the map at its angular seam: its full matrix is
literally identity there. The paired null-independence lemma PN
permits this null of the same zero-vertex corrected seam.

Flatten only the northern frame interior so that
$k(y,u(r,\theta))=m(y,\theta)$ for $r\ge1/2$.
For example precompose its radial argument by $\min(2r,1)$ and
extend to the south by the actual deck law. This preserves the complete
equator path, so IK permits it without changing the paired class.
Put $t=2r-1$ on this outer collar. Direct substitution in both \eqref{p1b:eq:natural-jt-5}
formulas, retaining their southern fixed-axis factor, gives
\[
 Z_k=B_0^{-1} W_k B_0,\qquad
 W_k=\begin{cases}
 I,&\text{northern inner disk},\\
 \mathcal C(D_{t\theta},m(y,\theta)),&\text{northern outer collar},\\
 \mathcal C(D_{\pi+t\phi},d m(y,\phi)d),
                         &\text{southern outer collar},\\
 \mathcal C(D_\pi,k(y,u)),&\text{southern inner disk}.
 \end{cases}
 \tag{BG.4}\label{p1b:eq:natural-bg-4}
\]
Here $\phi$ is the original southern ordered angle, and its
equator point has northern angle $\theta=\phi+\pi$ modulo $2\pi$.
For the southern outer equation the exact adjacent cancellation is
\[
 \mathcal C(B_{\pi+t\phi},dmd)
 \mathcal C(B_\pi,dmd)^{-1}
 \mathcal C(B_\pi,dmd)\mathcal C(B_0,dmd)^{-1}
 =\mathcal C(B_{\pi+t\phi},dmd)\mathcal C(B_0,dmd)^{-1}.
 \tag{BG.5}\label{p1b:eq:natural-bg-5}
\]
For every common operand $c$,
$\mathcal C(B_\alpha,c)\mathcal C(B_0,c)^{-1}
=B_0^{-1}\mathcal C(D_\alpha,c)B_0$.
This proves \eqref{p1b:eq:natural-bg-4} in its full matrix order. The whole word $W_k$ is
identity on both $\tau$ faces because every $D_\alpha$ is identity
there. Consequently replace the conjugator $B_0(\tau)$ in \eqref{p1b:eq:natural-bg-4} by
$B_0(\tau+h(1-\tau))$. This ends at identity and keeps every
collapsed face, proving $[Z_k]=[W_k]$ without an extra conjugation
term.

\paragraph{One actual cap of the intermediate axis, used on both sides.}
The map $D_\pi:S^1_\tau\wedge S^1_s=S^2\to SU(2)$ is based.
Choose one based null $J_\lambda$, with $J_0=D_\pi$ and $J_1=I$,
which exists since $\pi_2SU(2)=0$. On the entire southern inner
disk replace its first operand by $J_\lambda$, holding its actual
frame operand fixed. Insert a collar between that disk and its outer
annulus; its frame operand is the actual common equator value. With
collar coordinate $a$ increasing from its inner to outer boundary,
use $J_{\lambda(1-a)}$. At $\lambda=0$ this is the constant
$D_\pi$ collar; inserting it is a relative radial source comparison.
At every $\lambda$ its endpoints match the changed inner disk and
the retained outer annulus. At $\lambda=1$ the entire inner disk
is identity; traversing the inserted collar inward appends precisely
the path $J_0\to J_1$ to the outer axis path. All $\tau,s,y$
base faces remain identity. The north inner disk remains identity.

The two normal meridians $\theta=0,\pi$ are now literally
identity throughout these collars: their frame operands are identity.
The whole normal sphere therefore pinches into two closed bigons,
each with coordinates an actual half-equator angle and a radial path
from the north identity region to the south identity region. This is
a pinch of the already assembled normal sphere, not separate open
hemisphere coefficients.

\paragraph{The two closed radial axis loops.}
On the first half, $0\le\theta\le\pi$, the southern
representative has $\phi=\theta+\pi$ and its actual frame operand
is $d m(\theta+\pi)d=m(\theta)$. Its radial axis loop, oriented
from north to south, is exactly
\[
 L^{(1)}_\theta:\quad
 D_{0\to\theta}\ *\ D_{\theta+2\pi\to\pi}\ *\ J.
 \tag{BG.6}\label{p1b:eq:natural-bg-6}
\]
On the second half, $\pi\le\theta\le2\pi$, use
$\phi=\theta-\pi$. Its frame operand is again $m(\theta)$,
and the loop is
\[
 L^{(2)}_\theta:\quad
 D_{0\to\theta}\ *\ D_{\theta\to\pi}\ *\ J.
 \tag{BG.7}\label{p1b:eq:natural-bg-7}
\]
The seam between the first two paths in \eqref{p1b:eq:natural-bg-6} is exact by the
literal periodicity $D_{\theta+2\pi}=D_\theta$; no fractional
circle null is used. All loops start and end at identity and are
identity on the $\tau,s$ base faces.

Within each complete bigon commutator freeze only the angle occurring
in $L^{(j)}_\theta$ to $0$ on the first closed half interval and to
$\pi$ on the second closed half interval,
while retaining $m(y,\theta)$ pointwise. Linear interpolation of
that angle is a homotopy of the whole based word: the two angular
faces stay identity because $m=I$, the radial faces stay identity
because the axis loops are based, and the remaining parameter faces
are those already proved. The endpoint therefore is exactly the
quaternion commutator of a based map
$S^1_\tau\wedge S^1_s\wedge S^1_{\rm radial}=S^3$
with the actual half-edge map
$S^3_y\wedge S^1_{\rm half}=S^4$.
Each closed normal bigon gives the literal two-circle smash source
of this $S^3\wedge S^4=S^7$ word.

Let $L_+=D_{0\to\pi}*J$ and $L_-=\bar J*D_{\pi\to2\pi}$.
They are based three-sphere maps. Concatenation cancels the same
$J*\bar J$ and gives the original full axis map $D$, so
\[
 \deg L_+ +\deg L_- =\deg D=\pm1.
 \tag{BG.8}\label{p1b:eq:natural-bg-8}
\]
The second frozen loop \eqref{p1b:eq:natural-bg-7} is $L_+$ up to a retraced path. The
first frozen loop \eqref{p1b:eq:natural-bg-6} at $\theta=0$ is the reversal of $L_-$.
Thus its degree is $-\deg L_-$, and the degree sum of the two
closed bigon operands is
\[
 \deg L_+-\deg L_-=2\deg L_+-\deg D,
 \quad\text{an odd integer}.
 \tag{BG.9}\label{p1b:eq:natural-bg-9}
\]
Source orientation choices of the two bigon pinches change signs,
but retain this parity. Both half-edge operands have the same class
because the second is the actual deck conjugate of the first.

It follows for the literal complete zero-vertex paired map that
\[
 {}[Z_k]=[\mathfrak c]\circ(1_{S^3}\wedge m)
       =[\mathfrak c]\circ E^3[m]
       \quad\text{in }\pi_7SU(2)=C_2,
 \tag{BG.10}\label{p1b:eq:natural-bg-10}
\]
up to signs immaterial in this group, where
$\mathfrak c(a,b)=aba^{-1}b^{-1}$. The convention
$\mathcal C(a,b)=a^{-1}b^{-1}ab$ changes only operand inversions
and their signs. If $[m]=\eta_3$, then the retained commutator
generator theorem and $\nu'\eta_6=\mu$ give
\[
 \boxed{[Z_k]=\mu.}\tag{BG.11}\label{p1b:eq:natural-bg-11}
\]
If $[m]=0$, both complete bigon terms are zero. This calculation
concerns a zero-vertex frame only. Its application to the original
frame requires the full original-to-reference splitting and the
actual relative half-edge evaluation; it makes no direct original
$r_J$ assignment by itself.



\subsection{The absolute natural joining coefficient and completed top tests}
\label{p1b:sec:natural-joining-absolute-evaluation}

Retain the original determinant gauge, its actual odd clutch and
ordinary frames, the complete \eqref{p1b:eq:natural-nd-8} joining, the two actual lower
divisions, and the original projective correction \eqref{p1b:eq:natural-ob-13}--\eqref{p1b:eq:natural-ob-16}.
The calculation below evaluates its original coefficient $r_J$.
It uses the proved complete matched replacement (PI), the entire
zero second-cut comparison (SC), and the original-source coherent
identification (CI), rather than identifying an isolated hemisphere
or primitive with the original residual.

\paragraph{Identify the actual complete operand before evaluating it.}
Let $c(y,u)$ be the actual based neck frame of \eqref{p1b:eq:natural-ob-1}, with
$c(y,-u)=dc(y,u)d$.  Use the normal point $u_0=k$
appearing at $\theta=0$ in the original equator \eqref{p1b:eq:natural-ob-2}.
The reduced-neck proof \eqref{p1b:eq:natural-cn-1}--\eqref{p1b:eq:natural-cn-5} works at this point directly:
the actual outer slice $S^3_y\times\{k\}$ misses the original
root ball, and the entire punctured-frame comparison retains its
actual lower value.  Thus
\[
 {}[c(y,u)c(y,k)^{-1}]=\eta_3\eta_4\in\pi_5SU(2).
                                                               \tag{AJ.1}\label{p1b:eq:natural-aj-1}
\]
No normal-coordinate rotation has been performed.

Choose the moving-cap reference of \eqref{p1b:eq:natural-mr-1} with cap degree $-1$.
Its signed vertex degree is $-2$, equal to the actual endpoint
$C_2$ of \eqref{p1b:eq:natural-ob-8}.  The signed vertex alignment and its full ordering
homotopy extension are those of \eqref{p1b:eq:natural-fs-5}.  After that whole comparison,
the reference slice is literally $c(y,k)$.  Its complete paired
comparison has class zero by MR.4--MR.6, the supplied axis-null of
its actual normal lower value, and the zero-vertex projective-null
replacement TN with parameter sphere $S^5$.  Its full-cover reduced
frame clutch is zero by MR.7--MR.9.  Put
$k_0=c c_{\rm ref}^{-1}$; it is literally identity on the original
$u=k$ slice and on the sphere-base face, with the same deck equation.
The complete relative two-disc norm and the actual reduced clutch
give
\[
 {}[k_0]=\eta_3\eta_4=A_*[m],\qquad
 {}[m]=\eta_3\in\pi_4SU(2),
                                                               \tag{AJ.2}\label{p1b:eq:natural-aj-2}
\]
where $m$ is the actual first projective half-edge of $k_0$.
This uses $A_*\eta_3=\eta_3\eta_4$, and hence the isomorphism
on these two order-two groups.  It assigns no arbitrary frame loop.

\paragraph{Evaluate the whole paired map, including its southern meridian.}
The support-separated full product comparison \eqref{p1b:eq:natural-fs-1}--\eqref{p1b:eq:natural-fs-4}
carries the reference $H_3$, both cap data, the entire $K$ and $R$,
and the whole corrected seam throughout its source homotopy.  At
its terminal endpoint both original hemisphere words and the
nonconstant southern fixed-axis factor split pointwise into their
two disjoint supported factors.  Therefore
\[
 {}[Z(c)]=[Z(c_{\rm ref})]+[Z(k_0)]=[Z(k_0)].           \tag{AJ.3}\label{p1b:eq:natural-aj-3}
\]
The zero-vertex two-bigon calculation BG retains the entire
southern inner word $\mathcal C(D_\pi,k_0)$ and closes it by one
actual based null $J$ of $D_\pi$.  That \emph{same} $J$ closes
both complete bigons.  Their two closed radial axis loops have total
degree $2\deg L_+-\deg D$, an odd integer, because the original
full axis map $D$ has degree $\pm1$.  Their edge operands are the
actual $m$ and its deck conjugate.  Thus BG evaluates the complete
assembled map as
\[
 {}[Z(k_0)]=[\mathfrak c]\circ E^3[m]
         =\nu'\eta_6=\mu,
 \qquad \boxed{[Z(c)]=\mu.}                          \tag{AJ.4}\label{p1b:eq:natural-aj-4}
\]
The primitive commutator theorem is applied only after these two
closed source bigons and their odd degree have been constructed.
The intermediate $D_\pi$ meridian and its null are retained; they
are not assigned a separate parity and doubled.  The adjoint
operation $\eta_3E\nu'=0$ is used only in the separate
zero-vertex null-replacement lemma TN.  It does not replace the
source composition $\nu'\eta_6$ in \eqref{p1b:eq:natural-aj-4}.

\paragraph{Return to the original joining and its natural normalization.}
The native whole-boundary identification \eqref{p1b:eq:natural-ci-8}, including its
complete original cap and carried circle-base disk, states
$h_t\mu=[Z]$ and $r_J\mu=\mu+[Z]$.
The already evaluated original annular and characteristic maps give
the consistent original equation \eqref{p1b:eq:natural-pg-16}:
$h_t=1+r_J$, $\zeta_{\rm nat}=r_J$, and $w_t=1$.  Substitute \eqref{p1b:eq:natural-aj-4}:
\[
 \boxed{h_t=1,\qquad r_J=0,\qquad\zeta_{\rm nat}=0,\qquad w_t=1.}
                                                               \tag{AJ.5}\label{p1b:eq:natural-aj-5}
\]
In particular the original complete residual $\mathcal R_J$
has zero class in $\pi_7SU(2)=C_2$.  This is an evaluation of
that original coefficient through the entire licensed comparison,
not a claim that its matrix-valued representative is everywhere
identity.  The three normalized detector statements are distinct:
the original $\varphi_D$ of \eqref{p1b:eq:natural-ob-18}--\eqref{p1b:eq:natural-ob-20} has class
$h_t\mu=\mu$ and $\operatorname{Sq}^4$ coefficient $h_t=1$;
the original joining $\mathcal R_J$ has coefficient $r_J=0$;
the combined/rebased $\psi'$ has class $\zeta_{\rm nat}\mu=0$, and any
complete filling $K'$ with that retained boundary has
$\operatorname{Sq}^4$ coefficient
$\zeta_{\rm nat}=0$.  The zero rebased coefficient does not change the
original determinant normalization $h_t=1$.

\paragraph{Apply the values to the exact named extension tests.}
Use the already proved formulas \eqref{p1b:eq:rt-41-9}--\eqref{p1b:eq:rt-41-12},
\eqref{p1b:eq:rt-43-6}, and \eqref{p1b:eq:rt-55-7}.  They now become
\begingroup\small
\[
 \begin{array}{c|c|c|c|c}
 \text{candidate}&\text{domain}&O_7&O_8\text{ after }O_7=0&
                         \text{whole extension condition}\\\hline
 D_{e_0,t_0}&\text{odd exotic}&e_0\Gamma_\chi&(t_0+1)\Theta&e_0=0,\ t_0=1\\
 S_{e_0,t_0}&\text{retained odd sphere}&e_0\Gamma_\chi&(t_0+1)\Theta&e_0=0,\ t_0=1\\
 C_\tau&A=1,\ \epsilon=1&0&\tau\Theta&\tau=0\\
 B_{d,p}&\text{matched odd exotic}&d\Gamma_\chi&p\Theta&d=0,\ p=0.
 \end{array}                                                \tag{AJ.6}\label{p1b:eq:natural-aj-6}
\]
\endgroup
The sphere rows retain their original odd domain; the determinant
and Born rows retain the matched odd exotic domain. If \(O_7\ne0\),
extension fails before \(O_8\) is reached. Together with $S_{0,1}$, the
other natural YES entries are $D_{0,1}$, $C_0$ and
$B_{0,0}$; their opposite top alternatives are NO.  With the already
proved sphere test and all degree-seven failures, the natural
fourteen-entry ledger is $4$ YES and $10$ NO, with no unevaluated
entry.  The inherited thirty-two-entry ledger is $4$ YES and
$28$ NO.  In its retained odd exotic determinant row the actual
class is $a+t$, with no $e$ component.
\[
 \boxed{j_h=a+t,\qquad
        r_A=2b+z_Az+t,\qquad j_\eta=t,\qquad
        z_A=5-2A\pmod6.}                                \tag{AJ.7}\label{p1b:eq:natural-aj-7}
\]
The previously proved quotient groups, fixed-row fibre groups and
counts are unchanged.  This completes the finite natural rank-two
top classification frontier.  



\subsection{Complete natural boundary theorem and retained extension ledger}
\begin{theorem}[Complete natural rank-two boundary coefficients and extension tests]
\label{p1b:thm:natural-rank-two-completion}
On the retained source domain and normalization above, including the
complete joining \eqref{p1b:eq:natural-nd-8}, the original projective
correction \eqref{p1b:eq:natural-ob-13}--\eqref{p1b:eq:natural-ob-16},
and both actual lower divisions,
\begin{equation}
 [Z]=\nu'\eta_6=\mu,\qquad
 (h_t,w_t,r_J,\zeta_{\rm nat})=(1,1,0,0).
 \label{p1b:eq:natural-final-coefficients}
\end{equation}
The natural odd exotic normal forms are
\(j_h=a+t\), \(j_\eta=t\),
\(r_A=2b+z_Az+t\), \(z_A=5-2A\pmod6\).
The retained coefficients \(h_e=w_e=\omega=0\) are unchanged.
The scoped tests in \eqref{p1b:eq:natural-aj-6} decide all fourteen
natural entries: precisely \(S_{0,1},D_{0,1},C_0,B_{0,0}\) extend,
giving four YES and ten NO. All thirty-two inherited entries are
decided, giving four YES and twenty-eight NO. When the degree-seven
obstruction is nonzero, degree eight is not reached.
\end{theorem}
\begin{proof}
The bridge-controlled carry is fixed without a separate root-class
choice by Lemma~\ref{p1c:lem:v263-bridge-controlled-carry}; its
actual old-cap cancellation is Corollary~\ref{p1c:cor:v263-old-cap-cancellation}.

The whole Cartesian annulus contributes zero by
\eqref{p1b:eq:natural-pf-3}--\eqref{p1b:eq:natural-pf-9}, and
the restricted characteristic term is \(\mu\) by
\eqref{p1b:eq:natural-pg-14}. Thus
\eqref{p1b:eq:natural-pg-16} gives \(h_t=1+r_J\).
The complete cap and carried circle-base disk in
\eqref{p1b:eq:natural-ci-8} identify
\(h_t\mu=[Z]\), \(r_J\mu=\mu+[Z]\).
The exact neck clutch, common-null comparison, frame-interior
invariance, moving-cap reference, paired deck norm, whole product
split and both closed normal bigons are retained in the included
CN, PN, IK, MR, TN, FS and BG proofs. Their full evaluation is
\eqref{p1b:eq:natural-aj-1}--\eqref{p1b:eq:natural-aj-4}; the
closed axis degree is odd, giving \([Z]=\mu\).
\eqref{p1b:eq:natural-aj-5}--\eqref{p1b:eq:natural-aj-7} return
to the original coefficients, normal forms and scoped tests. These
tests give the two completed ledger totals; they are extension tests,
not the full mixed-class counts.
\end{proof}
This closes the natural rank-two classification frontier on its
declared geometry domain. 




\section{Analytic, finite, and boundary-response calculations}
\label{p1c:sec:spectral}

\subsection{The coupled Dirac family and reciprocal spectral symmetry}
\label{p1b:app:dirac}
The positive native result is stated in Section~\ref{p1b:sec:native-cartan-premises}; this technical section retains the supplementary spectral profiles and comparisons.







\begin{theorem}[Transverse onset and complete reciprocal-odd shell generation]
\label{p1b:thm:dirac-transverse-generation}
Let \(u=\epsilon h\) with real \(h\circ\iota=-h\), and use the
normalization of Theorem~\ref{p1b:thm:dirac-zero-resolvent}. Write
\[
 \psi_\epsilon=(1,0)+\sum_{n\ge1}\epsilon^n(\phi_n,\alpha_n),
 \qquad C=*\dd.
\]
Let \(\Delta_0^{-1}\) denote the inverse scalar Laplacian on mean-zero
functions and \(\Delta_1^{-1}\) the inverse one-form Laplacian on
\(S_R^3\). Then
\[
 \alpha_1=\frac{3}{2R}\dd\Delta_0^{-1}h,
 \qquad
 \boxed{\alpha_2=\frac{3}{4R^2}
 C\Delta_1^{-1}\!\left(h\,\dd\Delta_0^{-1}h\right).}
\]
Consequently, with
\[
 \Omega_h:=\dd h\wedge\dd\Delta_0^{-1}h,
\]
one has the exact criterion
\[
 \boxed{\alpha_2=0\quad\Longleftrightarrow\quad\Omega_h=0.}
\]
Let \(E_k\) be the degree-\(k\) scalar harmonic space,
\(E_k^-:=\{g\in E_k:\iota^*g=-g\}\), and let \(\mathcal C_k^-\)
be the reciprocal-odd part of the full coexact one-form shell with
Laplacian eigenvalue \((k+1)^2/R^2\). For every \(k\ge1\),
\[
 \boxed{\operatorname{span}\{*(\dd f\wedge\dd g):
 f\in E_1^-,\ g\in E_k^-\}=\mathcal C_k^-.}
\]
For \(k\ge2\), if \(f\in E_1^-\), \(g\in E_k^-\), and
\(h=af+bg\), the coefficient of \(ab\epsilon^2\) in the protected
one-form is
\[
 \boxed{\alpha_{2,\times}(f,g)=
 -\frac{R^2(k-1)(k+3)}{4k(k+2)(k+1)^2}
 *(\dd f\wedge\dd g).}
\]
Hence dipole--shell mixing spans every reciprocal-odd coexact vector
mode of the same level \(k\). At \(k=1\) the bilinear harmonic map is
still onto \(\mathcal C_1^-\), but the protected second-order response
vanishes because the two scalar Laplace eigenvalues coincide.
\end{theorem}
\begin{proof}
The recursion from the norm-convergent reduced-inverse formula is
\(\psi_n=-GM_h\psi_{n-1}\).  The first right-hand side is
\((3h/(2R),0)\), whose reduced inverse has one-form component
\(3d\Delta_0^{-1}h/(2R)\) and zero scalar component.
On a one-form input, the one-form component of \(G\) is
\(-C\Delta_1^{-1}\).  Applying the recursion once more gives the
formula for \(\alpha_2\).  Curl commutes with the Hodge inverse;
therefore \(\alpha_2=0\) exactly when
\(C(hd\Delta_0^{-1}h)=*(dh\wedge d\Delta_0^{-1}h)=0\).

We give an explicit surjectivity proof to retain the full shell claim.
It suffices to work at unit radius.  Write ambient coordinates as
\((t,x_1,x_2,x_3)\), so inversion is \((t,x)\mapsto(t,-x)\).
A degree-\(k\) coexact shell is represented by homogeneous harmonic
vector polynomials of degree \(k\), tangent to the sphere and
Euclidean divergence-free; the Gauss computation in the preceding
spectral proof gives eigenvalue \((k+1)^2\).
For a linear \(f=a\cdot x\) and a harmonic polynomial \(g\) of
 degree \(k\), the ambient vector representing
\(T(f,g)=*(df\wedge dg)\) is, up to the fixed orientation sign,
\[
 T_\mu=\epsilon_{\mu\nu\rho\sigma}q_\nu a_\rho\partial_\sigma g.
\]
It is tangential and divergence-free.  Its Euclidean Laplacian is
zero: the only cross term contains
\(\epsilon_{\mu\nu\rho\sigma}\partial_\nu\partial_\sigma g=0\).
Thus it lies in precisely that shell.  If \(f,g\) are inversion
odd, inversion's orientation sign makes \(T(f,g)\) an odd one-form.

Now let \(\beta\in\mathcal C_k^-\) be orthogonal to every such
\(T(f,g)\).  Curl is self-adjoint and \(d^*C\beta=0\), so
\[
 0=\langle\beta,C(f\,dg)\rangle
   =-\langle\langle df,C\beta\rangle,g\rangle
       \quad(g\in E_k^-).
\]
For each of the three imaginary coordinate functions \(f=x_i\),
the scalar \(\langle df,C\beta\rangle\) is the \(i\)-th ambient
component of its tangential harmonic vector polynomial.  It belongs
to \(E_k^-\): curl reverses inversion parity and \(df\) is odd.
Hence that component is zero.  Tangency then gives
\(t(C\beta)_t=0\); continuity gives the remaining component zero
also at \(t=0\).  Curl is invertible on a coexact shell, whence
\(\beta=0\).  Finite dimensionality proves the asserted spanning
identity, including \(k=1\).

Finally, in the mixed term the two scalar inverse eigenvalues are
\(R^2/3\) and \(R^2/[k(k+2)]\).  Curl kills the exact sum
\(d(fg)\), so their difference multiplies \(*(df\wedge dg)\).
The one-form inverse contributes \(R^2/(k+1)^2\).
Substitution in \(3C\Delta_1^{-1}/(4R^2)\) gives exactly the
printed coefficient.  It vanishes at \(k=1\) and is nonzero
for every \(k\ge2\).
\end{proof}


\begin{corollary}[Reciprocal helicity balance]
\label{p1b:cor:dirac-helicity-balance}
In the simple-kernel domain of Theorem~\ref{p1b:thm:dirac-zero-resolvent},
write the coexact part of the protected one-form as
\[
 \alpha^\perp=\sum_{L\ge1}(\alpha_{L,+}+\alpha_{L,-}),\qquad
 C\alpha_{L,\pm}=\pm\frac{L+1}{R}\alpha_{L,\pm}.
\]
Then reciprocal inversion exchanges the two curl branches and
\[
 \boxed{\|\alpha_{L,+}\|=\|\alpha_{L,-}\|\quad\text{for every }L.}
\]
In particular
\[
 \boxed{\int_{S_R^3}\alpha\wedge\dd\alpha=0.}
\]
Thus a nonzero protected one-form is neither a single Beltrami branch
nor a global contact form. This distinguishes the reciprocal protected
response from the one-handed Hopf contact/Beltrami field.
\end{corollary}
\begin{proof}
The normalized protected zero has positive reciprocal grade, so
\(\iota^*\alpha=-\alpha\).  Hodge projection commutes with the
isometry \(\iota\); moreover \(\iota^*C=-C\iota^*\).
It therefore maps the positive curl component in each shell
isometrically to the negative component, with a minus sign.
Their norms agree.  Exact forms contribute zero to
\(\int\alpha\wedge d\alpha=\langle\alpha,C\alpha\rangle\);
the two curl contributions cancel shell by shell.  Smoothness
justifies summation.  A nonzero single curl eigenbranch could not
have this balance.  If a real one-form were contact everywhere,
\(\alpha\wedge d\alpha\) would have a fixed nonzero sign on the
connected oriented sphere and a nonzero integral.  This proves
the last assertion as well.
\end{proof}


The explicit profile below lies in the longitudinal subclass
\(\Omega_h=0\); it is a closed-form representative of the general
protected zero line rather than an exhaustion of its profiles.

The explicit nonconstant zero is stated and directly verified in Section~\ref{p1b:sec:native-cartan-premises}; the supplementary profile and response comparisons remain here.


\begin{proposition}[Isolated-mode compression]
\label{p1b:prop:feshbach}
Let \(D_0=B_R\), \(P\) its kernel projector, \(Q=I-P\), and let
\(V=V^*\) be bounded with \(v=\|V\|<g/2\).
For \(|z|+v<g\), the eigenvalue equation near zero is equivalent to
the scalar Schur equation
\[
 K(z)=PVP-zP-PVQ(QD_0Q+QVQ-z)^{-1}QVP=0.
\]
Replacing the inverse by \((QD_0Q)^{-1}\) has error at most
\[
 \frac{v^2(v+|z|)}{g(g-v-|z|)}.
\]
For \(|z|\le v\), this is at most \(2v^3/[g(g-2v)]\).
For a spin-scalar fibre multiplier \(M\),
\[
 PMP=\left(\frac1{2V_R}\int_{S_R^3}\tr_E M\,\dd\mu_R\right)P.
\]
\end{proposition}
\begin{proof}
The complementary equation uniquely gives
\(Q\psi=-(QD_0Q+QVQ-z)^{-1}QVP\psi\), because its inverse norm is
at most \((g-v-|z|)^{-1}\).  Substitution in the retained equation
is exactly the displayed Schur equation, and reconstructs every
low eigenvector.  The resolvent identity bounds the difference
from \((QD_0Q)^{-1}\) by
\((v+|z|)/[g(g-v-|z|)]\); the two exterior factors contribute
\(v^2\).  For \(|z|\le v\) this gives the second bound.
The pointwise normalized parallel singlet has reduced twisting
 density \(I_E/2\).  Its constant spatial density is \(1/V_R\).
Taking the expectation of a spin-scalar multiplier proves the
last formula, with its ordinary doublet trace retained.
\end{proof}





The energy variable in \(K(\eta,z)\) remains part of the eigenvalue
equation. For example, with
\[
 D_0=\begin{pmatrix}0&0&0\\0&0&2\\0&2&0\end{pmatrix},\quad
 M=\begin{pmatrix}3&1&0\\1&5&0\\0&0&7\end{pmatrix},\quad
 U=\diag(1,1,-1),
\]
one gets \(K(\eta,0)=3\eta+\tfrac74\eta^3+O(\eta^5)\), whereas
the low eigenvalue is \(3\eta+\eta^3+O(\eta^5)\).
An insertion without the stated parity, or a kernel with both grades,
need not have a vanishing second-order compressed matrix.


\begin{proposition}[Exact heat trace and parameter evolution]
\label{p1b:result:exact-heat-trace-and-parameter-evolution}
For constant \(\delta\), put \(h=t/R^2\),
\(b_n=n(n+2)\), \(c_n=(n+1)^2\). Then
\[
\begin{split}
 \Tr e^{-tD_f^2}
 ={}&e^{-9h\delta^2}
 +2\sum_{n\ge1}c_ne^{-h(b_n+5\delta^2)}
       \cosh(2h\delta\sqrt{b_n+4\delta^2})\\
 &+2\sum_{n\ge1}b_ne^{-h(c_n+\delta^2)}
       \cosh(2h\delta\sqrt{c_n}).
\end{split}
\]
For a supplied time-dependent constant-on-the-sphere parameter
\(\delta(\tau)\), each exact/gradient pair evolves by the
two-by-two matrix
\[
 -\delta I+\sqrt{b_n}\sigma_x-2\delta\sigma_z .
\]
The polar angle \(\theta_n\) of the vector
\((\sqrt{b_n},0,-2\delta)\) has speed
\[
 |\theta_n'|=\frac{2\sqrt{b_n}|\delta'|}{b_n+4\delta^2}.
\]
For the rank-one spectral projectors \(P_{n,\pm}\), the corresponding
eigenline speed in the projector operator norm is
\[
 \|P_{n,\pm}'\|=\frac{|\theta_n'|}{2}
 =\frac{\sqrt{b_n}|\delta'|}{b_n+4\delta^2}.
\]
\end{proposition}
\begin{proof}
Square each root of the constant-profile spectral table and sum the
two signs using \(e^a+e^{-a}=2\cosh a\).  The singlet contributes
the first term and the two nonconstant sectors give the two sums.
Their Gaussian decay justifies convergence and differentiation on
compact positive-\(h\) intervals.  The exact/gradient matrix in the
spectral proof is
\(-\delta I+\sqrt{b_n}\sigma_x-2\delta\sigma_z\).
For its nonzero vector \(v=(\sqrt{b_n},0,-2\delta)\), the unit
vector derivative has length
\(2\sqrt{b_n}|\delta'|/(b_n+4\delta^2)\).
Since \(P_\pm=(I\pm\widehat v\cdot\sigma)/2\), their derivative
operator norm is half this length.
\end{proof}




\subsection{Stable refinements and geometric obstructions}
\label{p1b:app:obstructions}
\begin{corollary}[Real K-theory table]
\label{p1b:cor:ko}
The additive stable splitting of Theorem~\ref{p1b:thm:stable}
gives the following groups, periodically in \(q\) modulo eight:
\[
\begin{array}{c|c@{\qquad}c|c}
q&KO^{-q}(M_8)&q&KO^{-q}(M_8)\\\hline
0&\Z^6\oplus(\Z/2)^3&4&\Z^6\oplus(\Z/2)^3\\
1&\Z^4\oplus(\Z/2)^4&5&\Z^4\oplus(\Z/2)^6\\
2&\Z^2\oplus(\Z/2)^4&6&\Z^2\oplus(\Z/2)^6\\
3&\Z^4\oplus(\Z/2)^3&7&\Z^4\oplus(\Z/2)^3
\end{array}
\]
\end{corollary}
\begin{proof}
The stable product splitting has one sphere summand for every
subset of the factor dimensions \(1,1,3,3\), including the empty
subset for unreduced theory.  Their multiplicities by dimension
are the coefficients of
\((1+t)^2(1+t^3)^2\), namely
\((1,2,1,2,4,2,1,2,1)\).
Thus the answer is the direct sum of \(b_d\) copies of
\(KO^{-q-d}(\mathrm{pt})\).  Substitution of the Bott-periodic
coefficient groups
\((\mathbb Z,\mathbb Z/2,\mathbb Z/2,0,\mathbb Z,0,0,0)\)
gives every entry of the table.  A stable wedge splitting gives
actual direct sums, so there is no unresolved extension.
\end{proof}



\begin{proposition}[Symplectic obstruction and four-dimensional LCS forms]
\label{p1b:prop:lcs}
The compact manifold \(M_8\) admits no symplectic form.
On each \(K=S^3\times S^1\), however, a Hopf contact form \(\alpha\)
and the closed scalar form \(\tau\) give a locally conformally
symplectic form
\[
 \omega=\dd\alpha-\tau\wedge\alpha,\qquad
 \dd\omega=\tau\wedge\omega,\qquad
 \omega^2=-2\tau\wedge\alpha\wedge\dd\alpha .
\]
On the scalar cover \(\tau=\dd s\),
\(e^{-s}\omega=\dd(e^{-s}\alpha)\).
The Lie group \(M_8\) admits no left-invariant locally conformally
symplectic form.
\end{proposition}
\begin{proof}
The integral cohomology ring gives \(H^2(M_8;\mathbb R)=\mathbb R\eta\)
with \(\eta^2=0\).  Hence a closed two-form has zero fourth
cohomology power.  A symplectic form on a compact oriented
connected eight-manifold has a nonzero integral fourth power,
a contradiction.  The four-dimensional formulas follow directly
from \(d\tau=0\), \(d^2\alpha=0\) and the contact condition;
the covering identity is the ordinary product rule.

For the last assertion use the Lie algebra
\(\mathfrak g=\mathfrak{su}(2)^2\oplus\mathbb R^2\).
Any invariant closed Lee form \(\theta\) vanishes on the derived
algebra.  If \(\theta\ne0\), choose a central \(Z\) with
\(\theta(Z)=1\).  For \(d_\theta=d-\theta\wedge\), the Cartan
identity is
\(d_\theta\iota_Z+\iota_Zd_\theta=-I\).
Thus a \(d_\theta\)-closed two-form is \(d_\theta\alpha\) for
an invariant one-form \(\alpha\).  On each three-dimensional
simple ideal \(d\alpha\) has rank at most two; hence its rank
on the two ideals is at most four.  Adding \(-\theta\wedge\alpha\)
raises rank by at most two, so the result has rank at most six.
If \(\theta=0\), the invariant second cohomology of the semisimple
part is zero (or compute directly with its Maurer--Cartan equations).
An invariant closed two-form is therefore \(d\alpha\) plus a
central two-form, again of rank at most six.  Neither case is LCS.
\end{proof}



\begin{theorem}[Specified degree-three twists]
\label{p1b:thm:twisted-k}
For \(h=ku_x+lu_p\ne0\) in \(H^3(M_8;\Z)\), let
\(g=\gcd(|k|,|l|)>0\). Ordinary complex K-theory with this
degree-three twist has
\[
 K_h^0(M_8)\cong K_h^1(M_8)\cong(\Z/g)^4.
\]
For \(h=0\), both groups are \(\Z^8\).
For a nonzero real three-class \(h\), the corresponding
two-periodic twisted de Rham cohomology is zero.
\end{theorem}
\begin{proof}
On a three-sphere a degree-\(k\) twist gives the two-term integral
complex \(\mathbb Z\xrightarrow{k}\mathbb Z\) in degrees zero
and three.  Equivalently this follows by trivializing the twist
on the two hemispheres and applying the Mayer--Vietoris boundary
map; it is multiplication by its clutching degree \(k\).
Consequently its twisted groups are \(K^0_k=0,K^1_k=\mathbb Z/|k|\)
for \(k\ne0\), and \(K^0=K^1=\mathbb Z\) for \(k=0\).
The external-product Kunneth sequence for the sum of the two
pullback twists has tensor and Tor terms
\(\mathbb Z/|k|\otimes\mathbb Z/|l|\) and
\(\operatorname{Tor}(\mathbb Z/|k|,\mathbb Z/|l|)\).
Both are \(\mathbb Z/g\), in opposite parities, with just one
nonzero term in each parity, so no extension is possible.
If one of \(k,l\) is zero the free sphere groups give the same
answer directly.  The free even and odd K-groups of \(T^2\)
each have rank two; taking its product therefore gives four
copies in each parity.  The untwisted answer follows in the
same manner.
For real twisted de Rham theory, averaging on this compact group
reduces the spectral sequence to the exterior algebra on the
closed degree-one and degree-three generators.  Its first
nonzero differential is multiplication by \(h\).  Contraction
with a real linear functional on \(\operatorname{span}(u_x,u_p)\)
which takes value one on \(h\) contracts that complex.  Its
cohomology is zero, so the finite filtered twisted complex has
zero limiting cohomology.
\end{proof}



\begin{proposition}[Compatibility test for the two declared groupoids]
\label{p1b:result:compatibility-test-for-the-two-declared-groupoids}
For the proposed pair-groupoid and bi-Hopf structures, preservation
of the pair unit \((q,q)\) by a Hopf target map
\((q,q)\mapsto(qh,hq)\) requires \(qh=hq\).
Thus the commutator obstructs that specific simultaneous unit
compatibility outside the commuting locus.
\end{proposition}
\begin{proof}
A pair-groupoid unit is exactly a diagonal pair.  The proposed
image lies on the diagonal if and only if \(qh=hq\).
Multiplying this equality by inverses is equivalent to vanishing
of the group commutator; no infinitesimal approximation is used.
\end{proof}



\subsection{Finite Cartan links, overlap algebra and spectral transport}
\label{p1b:app:finite-operators}
\begin{theorem}[Exact Cartan edge and face transport]
\label{p1b:thm:cartan-links}
On the Hurwitz graph, with adjacent vertices having round angle
\(\theta=\pi/3\), set \(g(q)=\diag(q,1)\in\SU(3)\).
For the right-frame connection \(A_s=-s\,\dd g\,g^{-1}\), its
inverse transport on the shorter edge is
\[
 U_{qq'}(s)=\exp\{s\log[g(q)g(q')^{-1}]\}.
\]
The spatial triangular holonomy has eigenvalues
\(e^{i\varphi(s)},e^{-i\varphi(s)},1\), with
\[
 \cos\varphi(s)=w(s)=\cos(\theta s)
                  +\frac4{3\sqrt3}\sin^3(\theta s).
\]
For \(0\le s\le1\), \(w(s)=w(1-s)\) and
\[
 \max_s\|I-P_s\|=\sqrt{2-\frac{10}{3\sqrt3}} .
\]
An identity-linked ribbon between profiles \(s_j,s_{j+1}\)
has defect \(2|\sin(\theta(s_{j+1}-s_j)/2)|\).
\end{theorem}
\begin{proof}
On a shorter geodesic the right logarithmic derivative of \(g\)
is constant in its parallel edge frame.  Solving its constant
coefficient transport equation gives the exponential displayed.
For the face calculation it suffices, by left and right isometries,
to use the Hurwitz triangle
\(1,(1+i+j+k)/2,(1+i+j-k)/2\).
Each edge logarithm has length \(\theta\).  Multiplying its
three quaternion exponentials, with
\(\exp(xv)=\cos x+v\sin x\) for a unit imaginary \(v\), gives
scalar part
\(\cos x+4\sin^3x/(3\sqrt3)\), \(x=\theta s\).
The two-by-two eigenvalues are therefore the stated conjugate
pair; the third defining coordinate is fixed.
The elementary triple-angle identity gives \(w(s)=w(1-s)\).
Also
\[
 \frac{dw}{dx}=\sin x\left[-1+\frac2{\sqrt3}\sin(2x)\right].
\]
On \([0,\pi/3]\) its minimum occurs at \(x=\pi/6\), with
value \(5/(3\sqrt3)\).  For a unitary conjugate eigenpair the
operator defect is \(\sqrt{2-2\cos\varphi}\), giving the
maximum.  The identity temporal links leave a single relative
edge exponential of angle \(\theta(s_{j+1}-s_j)\); its unitary
defect is the final sine formula.
\end{proof}


For any positive $C^1$ central functions \(K_s,K_\tau\) on \(\SU(3)\),
define \(\Psi_s(s)=-\log K_s(P_s)\) and
\(\Psi_\tau(d)=-\log K_\tau(\diag(e^{i\theta d},e^{-i\theta d},1))\).
Counting the \(96\) faces per layer and \(96\) ribbons per step gives
the restricted function
\[
 S=96\sum_j\Psi_s(s_j)+96\sum_j\Psi_\tau(s_{j+1}-s_j)
       +\text{constant}.
\]
Its interior derivative equation is
\(\Psi_s'(s_j)+\Psi_\tau'(s_j-s_{j-1})
-\Psi_\tau'(s_{j+1}-s_j)=0\).
For the three-layer endpoints \(0,1\), if \(\Psi_s(1-s)=\Psi_s(s)\),
the midpoint profile \((0,\frac12,1)\) is stationary for this
restricted function.

\begin{theorem}[Finite chiral regulator from a supplied transfer graph]
\label{p1b:thm:overlap}
Supply a cyclic transfer graph of \(N\) layers and a unitary rank-\(r\)
connection on its product with the Hurwitz graph. On
\(\ell^2(C_N\times2T)\otimes\C^4\otimes\C^r\), write spatial
links \(U_{qq'}\), temporal links \(V_{j,q}\), and reverse links
as adjoints. Put
\[
 \nabla_a\psi_{j,q}=
 \frac3{8\ell}\sum_{q'\sim q}n^a_{qq'}U_{qq'}^{(j)}\psi_{j,q'},
 \quad
 \nabla_4\psi_{j,q}=
 \frac{V_{j,q}\psi_{j+1,q}-V_{j-1,q}^*\psi_{j-1,q}}{2a_\tau}.
\]
Here \(n_{qq'}=\operatorname{Im}(q^{-1}q')/
|\operatorname{Im}(q^{-1}q')|\) and \(\ell,a_\tau>0\).
Let \(\Gamma_\mu\) be Hermitian Euclidean Clifford matrices,
\(\Gamma_5=\Gamma_1\Gamma_2\Gamma_3\Gamma_4\), and
\(D^{(0)}=\sum_\mu\Gamma_\mu\nabla_\mu\).
The covariant graph Laplacians
\[
 \Delta_H\psi_q=\frac3{4\ell^2}\sum_{q'\sim q}
       (\psi_q-U_{qq'}\psi_{q'}),\quad
 \Delta_\tau\psi_{j,q}=
 \frac{2\psi_{j,q}-V_{j,q}\psi_{j+1,q}
                   -V_{j-1,q}^*\psi_{j-1,q}}{a_\tau^2}
\]
are positive. For \(r_H,r_\tau\ge0\), \(\rho\in\R\), set
\[
 W=\tfrac12(r_H\ell\Delta_H+r_\tau a_\tau\Delta_\tau),\quad
 H=\Gamma_5(D^{(0)}+W-\rho I).
\]
Then \(H=H^*\). On its open invertible set, with \(a_R>0\),
\[
 D_{\rm ov}=a_R^{-1}(I+\Gamma_5\sgn H)
\]
satisfies
\[
 \Gamma_5D_{\rm ov}+D_{\rm ov}\Gamma_5
 =a_RD_{\rm ov}\Gamma_5D_{\rm ov},\qquad
 D_{\rm ov}^*=\Gamma_5D_{\rm ov}\Gamma_5 .
\]
Its integer index is
\(\nu=\Tr[\Gamma_5(I-a_RD_{\rm ov}/2)]
=-\frac12\Tr\sgn H\), constant on each connected gapped component.
\end{theorem}
\begin{proof}
Reversing a spatial edge negates its imaginary direction and
replaces its link by its adjoint.  Hence every \(\nabla_\mu\)
is skew-adjoint.  Pairing the two orientations of an edge gives
\(\langle\psi,\Delta\psi\rangle\) as a positive weighted sum
of \(\|\psi_x-U_{xy}\psi_y\|^2\); thus both Laplacians and
\(W\) are Hermitian.  The Clifford identities give
\(\Gamma_5^*=\Gamma_5\), \(\Gamma_5^2=I\), anticommutation
with \(D^{(0)}\), and commutation with \(W\).  These prove
\(H^*=H\).  On its invertible set put \(\varepsilon=\operatorname{sgn}H\).
Then \(\varepsilon^*=\varepsilon\), \(\varepsilon^2=I\), and
\(V=\Gamma_5\varepsilon\) is unitary with
\(\Gamma_5V\Gamma_5=V^*\).  Expanding
\(D=a_R^{-1}(I+V)\) proves both displayed relations.
The nonreal eigenvalues of \(V\) occur in conjugate pairs and
give no chiral index; at \(V=-1\), \(D=0\), and at \(V=1\)
the factor \(I-a_RD/2\) is zero.  Thus the trace is the integer
chiral zero-mode difference.  Since \(\operatorname{Tr}\Gamma_5=0\),
it also equals \(-\operatorname{Tr}\varepsilon/2\).
Continuous gapped deformations cannot change its inertia.
\end{proof}


The transfer graph and its link data are supplied in this theorem.
For this four-dimensional product graph a temporal sign twist can
be inserted in the closing ribbon. The second scalar circle of
\(M_8\) is absent from this graph construction.

\begin{proposition}[Clock covers and certified-neighborhood transfer]
\label{p1b:prop:cover-transfer}
Suppose a three-layer rank-three kernel \(H(\phi)\), with a flat
phase on its closing ribbon, has index one and gap exceeding
\(g_0\) for all \(\phi\). Its \(k\)-fold cyclic repetition has
dimension \(864k\), index \(k\), and gap exceeding \(g_0\), for
either periodic or antiperiodic temporal signs.
At unit stencil parameters, if changes of every spatial and
temporal link have norms at most \(\epsilon_s,\epsilon_t\) and
\(|\Delta\rho|\le\epsilon_\rho\), then
\[
 \|\Delta H\|\le b=6\epsilon_s+2\epsilon_t+\epsilon_\rho .
\]
If \(b<g_0\), the perturbed index is \(k\) and the gap exceeds
\(g_0-b\). If the links are holonomies on the given degree-\(k\)
clutched bundle, this agrees with that bundle's charge.
\end{proposition}
\begin{proof}
Discrete Fourier transformation in the repetition index decomposes
the repeated kernel into \(k\) three-layer kernels, at phases
\((\phi+2\pi m)/k\), \(0\le m<k\), where the total closing
phase \(\phi\) is zero or \(\pi\) for the two signs.
Every summand has the assumed index and gap.  Their dimensions,
indices and gap bounds give the first assertion.
For the perturbation estimate use the norm of the Clifford
contraction of a unit edge direction, which is one.
The eight spatial neighbors contribute at most \(3\epsilon_s\)
to the first-order stencil and \(3\epsilon_s\) to its Wilson
term.  The two temporal neighbors contribute at most
\(\epsilon_t\) to each of those terms.  The mass term contributes
\(\epsilon_\rho\).  The triangle inequality gives \(b\).
Hermitian eigenvalues move by at most the perturbation norm,
so none crosses zero if \(b<g_0\), and the gap is at least
\(g_0-b\).  Charge agreement in the last statement uses the
specified clutched interpolation; a bare link array was not
assigned a topological charge by this proof.
\end{proof}


For example, suppose an explicitly specified three-layer kernel has a
certified sixteen-point phase mesh with inertia $(n_-,n_+)=(433,431)$,
gap greater than $0.256$, and closing-block Lipschitz bound
$\|H(\phi)-H(\psi)\|\le|\phi-\psi|$. Then
\[
 g_0=0.256-\frac{11}{56}=\frac{417}{7000},
 \qquad \pi/16<11/56,
\]
is a strictly positive all-phase gap certificate. This is a conditional
certificate-transfer statement: the finite sample enclosures and their
roundoff bounds are part of the supplied certificate. For arbitrary bare
link arrays a separate cellwise interpolation is required to define a
clutching number.

\begin{theorem}[Locality and smooth spectral bundles]
\label{p1b:thm:locality}
For a range-one Hermitian graph kernel with
\(0<\gamma\le|H|\le M\), \(\gamma<M\), put
\(q=(M-\gamma)/(M+\gamma)\). At distinct vertices,
\[
 \|D_{\rm ov}(x,y)\|
 \le\frac{M}{a_R\gamma}\,q^{\lfloor d(x,y)/2\rfloor}.
\]
On a smooth gapped family,
\[
 \|D\sgn_H[E]\|\le\|E\|/\gamma,\qquad
 \|D^n\sgn_H[E_1,\ldots,E_n]\|
 \le\frac{n!\,\Gamma(n/2)}
 {\sqrt\pi\,\Gamma((n+1)/2)}
 \gamma^{-n}\prod_j\|E_j\|.
\]
For \(P=(I+\sgn H)/2\), the image bundle has connection \(P\dd\),
curvature \(P(\dd P\wedge\dd P)P\), and path transport
\[
 U'=[P',P]U,\quad U(0)=I,\qquad P(s)U(s)=U(s)P(0).
\]
\end{theorem}
\begin{proof}
Put \(c=(M+\gamma)^2/4\) and
\(x=(c(1+q^2)-H^2)/(2cq)\); its spectrum is in \([-1,1]\).
The Legendre generating function gives the norm-convergent expansion
\[
 (H^2)^{-1/2}=c^{-1/2}\sum_{n\ge0}q^nP_n(x).
\]
Since \(|P_n|\le1\) on that interval, truncation after \(n\)
has error at most \(q^{n+1}/\gamma\).
Multiplication by \(H\) produces a polynomial of range
\(2n+1\) with error at most \(Mq^{n+1}/\gamma\).
Choose \(n+1=\lfloor d(x,y)/2\rfloor\); for distance one use
the zero polynomial.  This proves the kernel bound.
For derivatives use the resolvent representation
\[
 \operatorname{sgn}H=\frac1\pi\int_0^\infty
          [(H-it)^{-1}+(H+it)^{-1}]\,dt.
\]
An \(n\)-th derivative has \(n!\) ordered resolvent products,
with norm bounded by
\(\prod\|E_j\|(\gamma^2+t^2)^{-(n+1)/2}\).
The beta integral gives exactly the printed constant, including
\(1/\gamma\) at \(n=1\).  In finite dimension differentiation
under the resulting integrable resolvent majorant is legitimate.
Differentiating \(P^2=P\) gives the connection curvature
\(P(dP\wedge dP)P\).  Finally \([P',P]\) is skew-adjoint;
differentiating \(P(s)U(s)-U(s)P(0)\) and using \(P^2=P\)
shows that it remains zero.  Thus the solution is unitary and
intertwines the projector ranges.
\end{proof}


The polynomial-locality argument is the finite-graph form of
the method in \cite{HernandezJansenLuscher1999Locality}.
For the unit Hurwitz stencil \(\|H\|\le|4-\rho|+8\).
All constants remain uniform in the number of layers if the gap
margin and this norm bound remain uniform.


\subsection{Further marked-line and scalar completions}
\label{p1b:app:marked-alternatives}
\begin{theorem}[Absolute Hurwitz invariants and the degree-twelve pencil]
\label{p1b:thm:invariant-pencil}
For defining spin coordinates \(x,y\), put
\begin{align*}
 f_6&=xy(x^4-y^4),& f_8&=x^8+14x^4y^4+y^8,\\
 g_{12}&=x^{12}-33x^8y^4-33x^4y^8+y^{12},&
 u_{12}&=f_6^2.
\end{align*}
They are \(2T\)-invariant. The invariant dimensions in degrees
\(0,\ldots,12\) are
\[
 (1,0,0,0,0,0,1,0,1,0,0,0,2).
\]
Moreover,
\[
 f_8^3=g_{12}^2+108f_6^4,\quad
 \|u_{12}\|^2=8/231,\quad \|g_{12}\|^2=32/5,\quad
 \langle u_{12},g_{12}\rangle=0.
\]
Evaluation gives a nowhere-zero map
\(\epsilon_L:L^{12}\to I_{12}^*\) for every spin line
\(L\subset\C^2\), where \(I_{12}=\operatorname{span}(u_{12},g_{12})\).
Its image \(A_{12}\) is a line bundle over \(\mathbb{CP}^1\), and
the polar map
\[
 U_{12,L}=\epsilon_L(\epsilon_L^*\epsilon_L)^{-1/2}
 \colon L^{12}\longrightarrow A_{12}(L)
\]
is unitary, with \(c_1(A_{12})=-12\).
Degree twelve is the first positive base-point-free system of
absolute \(2T\)-invariants.
\end{theorem}
\begin{proof}
Use the generators \(i,j,(1+i+j+k)/2\) in their defining doublet
matrices.  Substitution in the three displayed polynomials fixes each
of them.  The Molien sum, grouped by the scalar parts \(1,-1,0,1/2,-1/2\), is
\[
 \frac1{24}\left\{\frac1{(1-t)^2}+\frac1{(1+t)^2}
 +\frac6{1+t^2}+\frac8{1-t+t^2}+\frac8{1+t+t^2}\right\}.
\]
Expansion through degree twelve gives the dimension list.
The relation follows by multiplying the polynomials.
For degree-\(d\) binary forms the inherited symmetric-tensor norm is
\(\|\sum a_rx^{d-r}y^r\|^2=\sum |a_r|^2/\binom dr\).
Thus the two norms are
\(2/\binom{12}2+4/\binom{12}6=8/231\) and
\(2+2\cdot33^2/\binom{12}4=32/5\); their monomial supports are
disjoint.  If \(f_6=0\), then either one coordinate vanishes or
\(x^4=y^4\); substitution gives a nonzero \(g_{12}\) in either
case for \((x,y)\ne0\).  Hence \(u_{12},g_{12}\) have no common
projective zero.  Evaluation is therefore injective on each line
\(L^{12}\); its positive polar factor is smooth and gives the unitary
isomorphism with its image.  Since \(L=\mathcal O(-1)\), its Chern
class is \(-12\).  In every smaller positive degree the Molien list
is zero or one-dimensional; a single positive-degree homogeneous
polynomial has a projective zero.  None is base-point-free.
Differentiating in an affine coordinate gives the displayed
Wronskian.  The zeros of \(f_6\) and \(f_8\) are disjoint and simple,
so they give respectively six simple and eight double ramification
points.  Projection curvature is the pullback Fubini--Study curvature
and vanishes at these critical points.  Direct substitution of the
normalizer gives its two stated signs; substitution of an invariant
polynomial gives the evaluation naturality formula.
\end{proof}


In an affine coordinate,
\(u_{12}'g_{12}-u_{12}g_{12}'=-2f_6f_8^2\).
The projective evaluation map has six simple and eight double
ramification points. Its projection curvature vanishes there,
so its projection connection generally differs from the connection
transported from \(L^{12}\) by \(U_{12}\).
The normalizer element \((1+\mathbf i)/\sqrt2\) acts on
\((u_{12},g_{12})\) as \((u_{12},-g_{12})\).
For \(h\in2T\), naturality is
\(\epsilon_{\rho(h)L}\rho(h)^{12}=\epsilon_L\).
Consequently closing the corresponding twelfth-power transport
through an axial line \(\ell=\C h\) leaves the covector
\(\lambda_h^{12}\in\ell^{*12}\); tracing the image-line
endomorphism has already contracted the Hopf factor.

\begin{proposition}[Opposite Hopf lines and lifted axial characters]
\label{p1b:result:opposite-hopf-lines-and-lifted-axial-characters}
On real \(\HH\), let \(Jq=\mathbf i q\), and on its complexification
let \(P_\pm=(1\mp iJ)/2\). For a real unit \(q\), set
\(z_\pm=\sqrt2P_\pm q\). Then
\[
 z_\pm(e^{\theta J}q)=e^{\pm i\theta}z_\pm(q),\quad
 b(z_+,z_-)=1,\quad b(z_\pm,z_\pm)=0 ,
\]
where \(b\) is the complex bilinear extension of the metric.
Thus the opposite Hopf lines are dual.
For the moving axial line \(\C q\), a lifted action
\[
 T_k(\theta)(cq)=e^{ik\theta}c\,e^{\theta J}q
\]
has invariant negative projection exactly for \(k=1\), and
invariant positive projection exactly for \(k=-1\).
With \(c_1(E_{{\rm ax},k}/\U(1))=k\),
\[
 c_1(A_{12}\otimes E_{{\rm ax},k}^{12})=12(k-1).
\]
\end{proposition}
\begin{proof}
Since \(J^2=-I\), the two \(P_\pm\) are its complementary complex
eigenspace projectors and commute with the circle action, giving the
two weights.  Orthogonality of \(J\) implies that its same-sign
eigenspaces are isotropic for the complex bilinear metric.
Writing \(q=P_+q+P_-q\) and using \(b(q,q)=1\) gives
\(2b(P_+q,P_-q)=1\), proving the three pairings after the
\(\sqrt2\) normalization.  The projected lifted weights are
\(k-1\) and \(k+1\), so their vanishing gives exactly the two
characters stated.  Finally first Chern classes add under tensor
product and multiply by twelve under the twelfth power; with
\(c_1(A_{12})=-12\) this gives \(12(k-1)\).
\end{proof}



\begin{proposition}[Angular moments and comparison heat]
\label{p1b:result:angular-moments-and-comparison-heat}
For \(|n|=1\), \(\epsilon\in\{\pm1\}^3\), let
\(f_\epsilon(n)=[3-(\epsilon\cdot n)^2]/4\).
Then
\[
 \sum_\epsilon f_\epsilon=4,\qquad
 \sum_\epsilon f_\epsilon^2=3-\sum_i n_i^4,\qquad
 2\sum_\epsilon e^{-\tau f_\epsilon}\ge16e^{-\tau/2}\quad(\tau>0).
\]
The quadratic sum has minima \(2\) on the six coordinate axes,
maxima \(8/3\) on the eight body diagonals, and saddles on the
twelve edge directions. Equality in the heat inequality occurs
exactly on the axes. Holding \(h=(1+\mathbf i+\mathbf j+\mathbf k)/2\)
fixed while varying the spin line gives a quadratic section with
two simple projective zeros.
\end{proposition}
\begin{proof}
The eight independent sign choices have
\(\sum(\epsilon\cdot n)^2=8\) and
\(\sum(\epsilon\cdot n)^4=8(3-2\sum n_i^4)\).
Expanding \(f_\epsilon\) and its square gives the two moments.
Jensen's strict convexity inequality applied to their mean \(1/2\)
gives the heat inequality.  Equality means all squared sign sums
are one; their sign-Fourier coefficients are the products \(n_in_j\),
so this is equivalent to at most one nonzero coordinate, the axes.
On \(\sum n_i^2=1\), Lagrange multipliers for \(\sum n_i^4\)
force all nonzero squared coordinates equal.  One, two or three
nonzero coordinates give respectively the axes, edge saddles and
body-diagonal extrema, with the printed values.  For fixed \(h\),
the quadratic spin-line section is proportional to
\(z\wedge\rho(h)z\).  The two distinct eigenvalues
\(e^{\pm i\pi/3}\) of \(\rho(h)\) give two distinct eigenlines,
hence exactly two simple projective zeros.
\end{proof}



\begin{theorem}[A specified theta-series completion of the scalar lattice]
\label{p1b:thm:theta}
Supply \(R>0\), \(\tau\) in the upper half-plane, and the
oscillator factor \(\eta_{\rm D}(\tau)^{-1}\), where
\(\eta_{\rm D}(\tau)=q^{1/24}\prod_{m\ge1}(1-q^m)\),
\(q=e^{2\pi i\tau}\). Set
\[
 p_L=n/R+wR,\qquad p_R=n/R-wR,
\quad
 Z(R;\tau)=\frac1{|\eta_{\rm D}(\tau)|^2}
       \sum_{n,w\in\Z}q^{p_L^2/4}\bar q^{p_R^2/4}.
\]
The series is invariant under \(\tau\mapsto\tau+1\),
\(\tau\mapsto-1/\tau\), and \((n,w,R)\mapsto(w,n,R^{-1})\).
At \(R=1\), its two charges are the same-parity eigencharges
\(p_L=n+w,p_R=n-w\).
\end{theorem}
\begin{proof}
The summand is
\(\exp[-\pi\tau_2(n^2/R^2+w^2R^2)+2\pi i\tau_1nw]\).
It is absolutely locally uniformly summable in the upper half-plane.
Translation of \(\tau_1\) by one multiplies it by \(e^{2\pi inw}=1\),
while \(|\eta_{\rm D}|\) is unchanged.  Two-dimensional Poisson
summation for this Gaussian gives
\(\Theta_R(-1/\tau)=|\tau|\Theta_R(\tau)\).
For completeness its Gaussian matrix is
\(\left(\begin{smallmatrix}\tau_2/R^2&-i\tau_1\\-i\tau_1&\tau_2R^2\end{smallmatrix}\right)\),
with determinant \(|\tau|^2\); its inverse is the transformed
matrix after interchanging the two integer coordinates.
The Dedekind product transformation
\(|\eta_{\rm D}(-1/\tau)|^2=|\tau|\,|\eta_{\rm D}(\tau)|^2\)
then cancels this factor.  It follows also by the one-dimensional
Gaussian Poisson identity applied to the standard eta product.
Interchanging \(n,w\) and replacing \(R\) by \(1/R\) fixes
\(p_L\) and negates \(p_R\), so fixes the sum.  At \(R=1\)
the two integral sums \(n\pm w\) have equal parity, and conversely
every such pair gives integral \(n,w\).
\end{proof}



\begin{proposition}[Local coframes, cones and quadratic variations]
\label{p1b:prop:local-coframe}
Supply an isomorphism from the dual relative Lie algebra to the
cotangent bundle of a local four-dimensional chart.
The relative solder then gives a rank-four coframe. A positive
scalar coframe homothety \(a\) has determinant \(a^4\).
If a positive spatial metric \(h\), a scalar coordinate \(t\),
and \(c>0\) are given, the cone
\[
 \{(\tau,\xi):\tau\ge0,\ h(\xi,\xi)\le c^2\tau^2\}
\]
determines the \(O(3)\)-invariant conformal class of
\(-c^2\dd t^2+h\), up to an overall sign.
For a coframe variation
\(\delta e^A=M^A{}_Be^B\) and
\(M_{AB}=\lambda\eta_{AB}+S_{AB}+\Omega_{AB}\), with \(S\)
symmetric trace-free and \(\Omega\) antisymmetric,
\[
 \delta g=2\lambda g+2S,\qquad \delta_\Omega g=0,\qquad
 S_0^2(\R\oplus\R^3)\downarrow\SO(3)
   =S_0^2(\R^3)\oplus\R^3\oplus\R .
\]
\end{proposition}
\begin{proof}
An isomorphism sends a basis of the relative dual to four linearly
independent one-forms.  Scaling each by \(a\) scales their top wedge
by \(a^4\).  An \(O(3)\)-invariant quadratic form on the specified
scalar-plus-spatial decomposition has no mixed term and is
\(a_0\tau^2+b_0h(\xi,\xi)\).  Its stated null cone forces
\(a_0=-b_0c^2\), proving uniqueness up to scale and overall sign.
Differentiation of \(g=\eta_{AB}e^Ae^B\) kills the antisymmetric
part and doubles the symmetric part, giving the variation formula.
Under spatial rotations a symmetric four-by-four matrix splits
into a spatial traceless symmetric tensor, its three scalar-spatial
entries, and two scalars; removing its total trace leaves one scalar.
These are precisely the displayed \(5+3+1\) summands.
\end{proof}



\begin{proposition}[Optional cylinder kink and its chiral mode]
\label{p1b:prop:kink}
On the additional infinite cylinder \(\R_\tau\times S_R^3\),
let \(f=(1+\tanh(\tau-\tau_0))/2\), and use the doublet family
of Theorem~\ref{p1b:thm:dirac-spectrum}. For
\(W=\partial_\tau+RD_f\) on its natural \(L^2\) domain,
\[
 \ker W=0,\qquad\dim\ker W^*=1,\qquad\operatorname{ind}W=-1.
\]
With \(z=\tau-\tau_0\), the unit mode for measure
\(R^4\dd\tau\,\dd\mu_{S^3_1}\) is
\(\Psi_0=\operatorname{sech}^{3/2}z\,\xi_0/(\pi^{3/2}R^2)\),
where \(\xi_0\) is the pointwise unit singlet.
The singlet partner potentials are
\[
 -\partial_z^2+\frac94-\frac{15}4\operatorname{sech}^2z,\qquad
 -\partial_z^2+\frac94-\frac34\operatorname{sech}^2z,
\]
with bound values \(\{0,2\}\) and \(\{2\}\), respectively,
and continuous spectrum beginning at \(9/4\).
\end{proposition}
\begin{proof}
Use the orthogonal harmonic blocks in the constant-profile spectral
proof, which are independent of \(\tau\).  The singlet has
\(RD_f=-3\tanh z/2\).  Thus \(W^*\psi=0\) has the single
square-integrable solution \(\operatorname{sech}^{3/2}z\), while
\(W\psi=0\) has no square-integrable solution.
There are no hidden nonconstant zero modes.  Indeed on an
exact/gradient block put \(b=n(n+2)\ge3\), \(t=\tanh z\), and
\(D=-tI/2+\sqrt b\sigma_x-t\sigma_z\).  Direct multiplication gives
\[
 D^2-D'=(b+\tfrac12+\tfrac34t^2)I-\sqrt b\,t\sigma_x+\sigma_z,
\]
\[
 D^2+D'=(b-\tfrac12+\tfrac74t^2)I-\sqrt b\,t\sigma_x+(2t^2-1)\sigma_z.
\]
Their lowest pointwise eigenvalues are at least
\(b+1/2-\sqrt{b+1}>0\) and
\(b-1/2-\sqrt{b+1}>0\), respectively.
On a coexact block \(D=t/2\pm(n+1)\), both partner potentials
are bounded below by \((n+1-1/2)^2-1/2>0\).
Consequently only the singlet contributes to either kernel.
The closed realizations are the closures of the smooth compactly
supported blocks in the first Sobolev graph norm; the uniform
positive bounds and the invertible limits at both cylinder ends
make them Fredholm.  Hence the index is \(-1\).
The integral \(\int_{\mathbb R}\operatorname{sech}^3z\,dz=\pi/2\)
and sphere volume \(2\pi^2\) give the normalization.
Finally factor with \(Q_a=\partial_z+a\tanh z\):
\(Q_a^*Q_a=-\partial_z^2+a^2-a(a+1)\operatorname{sech}^2z\).
At \(a=3/2\) its partner is \(Q_{1/2}^*Q_{1/2}+2\).
Their normalizable kernels insert values zero and two; the next
partner has no well below the free threshold.  The continuum
starts at \(9/4\), proving the full singlet table.
\end{proof}




\subsection{Finite normalization, stationary reduction and matrix invariants}
\label{p1b:app:finite-calculus}
This technical section collects the finite mathematical tools used by descendants
of the reciprocal construction. Each matrix, quadratic form, measure,
or function mentioned here is part of the stated input.

\begin{proposition}[Positive-metric normalization]
\label{p1b:prop:kinetic}
For a positive metric \(Z\) and Hermitian quadratic form \(K\),
the normalized endomorphism is \(Z^{-1/2}KZ^{-1/2}\).
For the coefficient \(D\) of a sesquilinear coupling \(x_L^*D x_R\),
with positive Hilbert metrics \(Z_L,Z_R\), the normalized coefficient is
\(\widehat D=Z_L^{-1/2}DZ_R^{-1/2}\).
An actual linear map \(T:E_R\to E_L\), expressed in the orthonormal
coordinates \(y_i=Z_i^{1/2}x_i\), instead has matrix
\(Z_L^{1/2}T Z_R^{-1/2}\).
If \(Z_L=cI_L,Z_R=cI_R,D=cD_0\), then
\(\widehat D=D_0\). More generally common powers
\(Z\sim f^z,K\sim f^k\) give quadratic weight \(f^{k-z}\),
and \(Z_L\sim f^l,Z_R\sim f^r,D\sim f^d\) give coupling-coefficient weight
\(f^{d-(l+r)/2}\).
For a smooth moving positive metric, writing
\(x=Z^{-1/2}y\) transforms \(\dot x\) to
\[
 Z^{-1/2}(\dot y+\Omega y),\qquad
 \Omega=Z^{1/2}\frac{\dd}{\dd t}Z^{-1/2}.
\]
\end{proposition}
\begin{proof}
Substitute \(x=Z^{-1/2}y\) into \(x^*Kx\), and independently
into \(x_L^*Dx_R\).  This gives the two congruences.
For a linear map use \(y_L=Z_L^{1/2}Tx_R\), which is the
stated change of matrix rather than a coupling congruence.
The scalar cancellations and weights follow by multiplying the
powers.  Finally the product rule gives
\(\dot x=Z^{-1/2}\dot y+(dZ^{-1/2}/dt)y\); factoring the
leftmost \(Z^{-1/2}\) gives exactly \(\Omega\), without
assuming that \(Z\) commutes with its derivative.
\end{proof}



\begin{theorem}[Stationary reduction through fourth order]
\label{p1b:thm:schur-jets}
Let \(S(w,z)\) be a real \(C^4\) function with \(S_z=0\) and
invertible \(C=S_{zz}\) at a base point. Let \(z_*(w)\) be its
local stationary branch, \(\Phi(w)=S(w,z_*(w))\), and
\[
 A=S_{ww},\quad B=S_{zw},\quad Z=-C^{-1}B,\quad
 La=(a,Za),\quad T_j=D^jS .
\]
Then
\[
 D^2\Phi=A-B^TC^{-1}B,\qquad
 D^3\Phi[a,b,c]=T_3[La,Lb,Lc],
\]
and, with \(T_{3,I}[v,w](\zeta)=T_3[(0,\zeta),v,w]\),
\[
\begin{split}
 D^4\Phi[a,b,c,d]={}&T_4[La,Lb,Lc,Ld]\\
 &-\langle T_{3,I}[La,Lb],C^{-1}T_{3,I}[Lc,Ld]\rangle\\
 &-\langle T_{3,I}[La,Lc],C^{-1}T_{3,I}[Lb,Ld]\rangle\\
 &-\langle T_{3,I}[La,Ld],C^{-1}T_{3,I}[Lb,Lc]\rangle .
\end{split}
\]
For \(C>0\), the repeated-direction correction is
\(-3\|C^{-1/2}T_{3,I}[La,La]\|^2\).
The full Hessian may be singular.
\end{theorem}
\begin{proof}
The implicit function theorem applies to \(S_z=0\).
Its first derivative is \(z_*'a=-C^{-1}Ba\), giving the lift
\(La\) and the Schur Hessian.  Differentiating the stationary
equation twice gives
\(z_*''[a,b]=-C^{-1}T_{3,I}[La,Lb]\).
In the third derivative of \(S(w,z_*(w))\), every term involving
\(z_*''\) pairs a vertical vector against \(D^2S\,La\), whose
vertical part is zero; the term with \(S_z\) also vanishes.
This proves the cubic formula.  Differentiate it once more.
There is one fourth derivative term and three terms differentiating
the lifted arguments.  Substituting the displayed \(z_*''\) gives
the three contractions written in the statement, with no fourth
or third derivative of the stationary branch remaining.
For a repeated direction their sum is the stated negative squared
norm when \(C>0\).  Only the interior Hessian was inverted.
\end{proof}



\begin{theorem}[Partial elimination and Legendre restriction]
\label{p1b:thm:partial-elimination-legendre-duality}
Let \(S(x,z)\) and \(F(x)=\inf_z S(x,z)\) be proper extended-real
functions. Their convex conjugates satisfy
\[
 F^*(p)=S^*(p,0).
\]
This equality itself needs neither convexity nor an attained infimum.
For a separate local \(C^2\) stationary version, suppose \(S_z=0\) and \(S_{zz}\)
and the full Hessian \(H\) are invertible at the point in question.
Writing
\[
 H=\begin{pmatrix}A&B_0^T\\B_0&C_0\end{pmatrix},
 \qquad K=A-B_0^TC_0^{-1}B_0,
\]
the local stationary reduction and local Legendre maps are defined and
\[
 \boxed{K^{-1}=(H^{-1})_{\partial\partial}.}
\]
No positivity assumption is needed for this local identity. If \(A\)
is additionally invertible, eliminating the dual interior block instead
gives \(\operatorname{Schur}_I(H^{-1})=A^{-1}\).
\end{theorem}
\begin{proof}
By the extended-real supremum rules,
\(F^*(p)=\sup_x\{p\cdot x-\inf_zS(x,z)\}
=\sup_{x,z}\{p\cdot x-S(x,z)\}=S^*(p,0)\).
This argument does not exchange an infimum and a supremum.
Locally the implicit function theorem gives the stationary branch
with Hessian \(K\).  Block Gaussian elimination factors \(H\)
into invertible triangular factors and \(\operatorname{diag}(K,C_0)\).
Since \(H,C_0\) are invertible, so is \(K\); the upper-left
block of its inverse is \(K^{-1}\).  Applying the same block
inversion formula to \(H^{-1}\), now with the additional
invertibility of \(A\), gives its interior Schur complement
\(A^{-1}\).  No step requires positive definiteness.
\end{proof}


For \(H=\left(\begin{smallmatrix}2&1\\1&2\end{smallmatrix}\right)\),
the correct inverse reduced Hessian is \(2/3\), while eliminating twice
would give \(1/2\). At
\(H=\left(\begin{smallmatrix}1&1\\1&1\end{smallmatrix}\right)\),
\(C_0\) remains invertible but \(K=0\); there is no ordinary inverse.
For this positive semidefinite quadratic example the conjugate encodes a
domain constraint instead. Critical stationary reduction therefore does
not inherit the regular inverse assertion.

For any differentiable matrix family \(K=A-B^TC^{-1}B\),
all parameter derivatives are computed with
\((C^{-1})'=-C^{-1}C'C^{-1}\), including the derivatives of
the moving base point when present. A simple eigenvalue
\(\lambda\) of \(K\), with normalized eigenvector \(v\), satisfies
\[
 \lambda'=v^*K'v,\qquad
 \lambda''=v^*K''v+
     2\sum_{\mu\ne\lambda}
       \frac{|\langle v_\mu,K'v\rangle|^2}{\lambda-\mu}.
\]
Differentiating \(Kv=\lambda v\) in an orthonormal eigenbasis
proves the formulas. They require an isolated simple eigenvalue.

\begin{theorem}[Matrix reciprocal rigidity]
\label{p1b:thm:matrix-rigidity}
Let \(\mathcal K(y)\) be Hermitian and operator-concave on
\([0,Y]\), \(Y\ge1\), with
\(\mathcal K(0)=0\), \(\mathcal K'_+(0)=I\), and
\(\mathcal K(y)>0\) for \(y>0\).
If a fixed unitary \(B\) obeys
\[
 \mathcal K(1/y)=B\mathcal K(y)^{-1}B^*,
 \qquad 1/Y\le y\le Y,
\]
then \(\mathcal K(y)=yI\) on that reciprocal window.
\end{theorem}
\begin{proof}
Operator concavity and the tangent \(\mathcal K'_+(0)=I\)
give \(\mathcal K(y)\preceq yI\) throughout \([0,Y]\).
At the reciprocal argument this says
\(B\mathcal K(y)^{-1}B^*\preceq I/y\).
Unitary conjugation and reversal of order by positive inversion
give \(\mathcal K(y)\succeq yI\).  Both inequalities hold on
the specified reciprocal window, forcing equality there.
\end{proof}


This is an optional comparison within the matrix-family atlas; its
reciprocal identity is not imposed on every response in this paper.
The conditions concern the same normalized matrix family.
For example, the positive affine block
\(\left(\begin{smallmatrix}y&y\\y&1+y\end{smallmatrix}\right)\)
has Schur response \(y/(1+y)\), so a normalized first derivative
and positivity alone do not imply the conclusion.

\begin{proposition}[Determinant derivatives and relative-frame stationarity]
\label{p1b:result:determinant-derivatives-and-relative-frame-stationarity}
On a smooth invertible matrix branch with a local logarithm,
\[
 \partial_a(-\log\det D)=-\Tr(D^{-1}D_a),\quad
 \partial_{ab}(-\log\det D)
 =\Tr(D^{-1}D_bD^{-1}D_a)-\Tr(D^{-1}D_{ab}).
\]
For positive Hermitian \(A,B\), \(z,w,t>0\), set
\(D_A=zI+A,D_B=wI+B,C=D_B-tD_A^{-1}>0\).
The function
\(\Phi=-\log\det C+\log\det D_B\) is stationary under every
relative unitary rotation of \(B\) exactly when \([A,B]=0\).
\end{proposition}
\begin{proof}
Jacobi's identity \(d\log\det D=\operatorname{Tr}(D^{-1}dD)\)
and \(dD^{-1}=-D^{-1}(dD)D^{-1}\) give both derivative formulas.
For a skew-Hermitian generator \(X\), varying the relative frame
gives \(\delta B=[X,B]\) and hence
\(\delta\Phi=\operatorname{Tr}[(D_B^{-1}-C^{-1})[X,B]]\).
The term with \(D_B^{-1}\) vanishes since it commutes with \(B\).
Vanishing for every \(X\) is exactly \([B,C^{-1}]=0\),
equivalently \([B,C]=0\).  Since \(t>0\), this is
\([B,D_A^{-1}]=0\), and invertibility of \(D_A\) makes it
\([A,B]=0\).  Reversing these implications proves sufficiency.
\end{proof}



\begin{theorem}[Generic rank-three matrix-pair reconstruction]
\label{p1b:thm:rank-three}
Let \(A,B\) be positive Hermitian three-by-three matrices with
distinct ordered eigenvalues \(\alpha_i,\beta_j\), and let \(V\)
be their relative unitary eigenframe. Their spectral moments
\(\tr A^r,\tr B^r\), \(r=1,2,3\), and four mixed moments
\(M_{pq}=\tr A^pB^q\), \(p,q=1,2\), determine
\[
 W_{ij}=|V_{ij}|^2,\qquad
 W=S_A^{-1}MS_B^{-T},\quad
 (S_A)_{pi}=\alpha_i^p,\ (S_B)_{qj}=\beta_j^q,\quad p,q=0,1,2,
\]
where \(M_{00}=3,M_{p0}=\tr A^p,M_{0q}=\tr B^q\).
A nonnegative matrix \(W\) with unit row and column sums is
unistochastic exactly when
\[
 J^2=\frac14\left(2\sum_{i<j}r_ir_j-\sum_i r_i^2\right)\ge0,
 \qquad r_j=W_{1j}W_{2j}.
\]
On the strictly positive, positive-area locus, the pair is
determined up to simultaneous unitary conjugacy and the two
complex-conjugate orientations. For
\(\chi=\tr[A,B]^3/(6i)\),
\[
 \chi^2=\Delta_A^2\Delta_B^2J^2,\qquad
 \Delta_A=(\alpha_1-\alpha_2)(\alpha_2-\alpha_3)(\alpha_3-\alpha_1),
\]
and similarly for \(\Delta_B\).
The mixed resolvent is the positive function
\[
 \tr[(zI+A)^{-1}(wI+B)^{-1}]
 =\sum_{ij}\frac{W_{ij}}{(z+\alpha_i)(w+\beta_j)}
 \quad(z,w>0).
\]
\end{theorem}
\begin{proof}
Newton's identities recover the ordered spectra from their first
three power sums.  In their eigenbases
\(M_{pq}=\sum_{ij}\alpha_i^p\beta_j^q|V_{ij}|^2\).
The two Vandermonde matrices are invertible because the eigenvalues
are distinct; the moment inversion formula follows.
To test unistochasticity, choose phases on the first row to make
it real nonnegative.  Orthogonality of rows one and two requires
three complex vectors of lengths \(\sqrt{r_j}\) to close.
Their triangle condition is exactly nonnegativity of the Heron
expression displayed.  Conversely any such triangle supplies their
relative phases; a unit vector orthogonal to these two rows then
has squared entries \(1-W_{1j}-W_{2j}=W_{3j}\), so completes
a unitary.  On the positive-area locus the two triangles differ
only by reflection, after the irrelevant row and column phases.
This proves generic uniqueness up to conjugation and its reflected
orientation.  The triangle area is \(|J|/2\); Heron's formula is
therefore the printed \(J^2\).  Expanding the six cyclic products
in the determinant of the trace-zero commutator gives
\(\det[A,B]=2iJ\Delta_A\Delta_B\), up to the common orientation
sign.  Since a trace-zero three-by-three matrix has
\(\operatorname{tr}C^3=3\det C\), squaring gives the stated
\(\chi\) relation.  Finally insert both spectral decompositions
in the resolvent trace.  Its summands are nonnegative and their
sum is strictly positive, as \(\sum_{ij}W_{ij}=3\).
\end{proof}



\begin{proposition}[Explicit chiral bundle charts]
\label{p1b:result:explicit-chiral-bundle-charts}
Let \(H_0\) be a finite Hermitian matrix of gap \(\gamma_0>0\),
\(\|H-H_0\|\le\delta<\gamma_0\), \(\gamma=\gamma_0-\delta\),
\(\|H\|\le M\), and \(P_0,P\) the positive spectral projectors.
On \(\im P_0\),
\[
 G=P_0PP_0\ge\frac{2\gamma}{M+\gamma}I,\qquad
 T=PP_0G^{-1/2},\qquad T^*T=P_0,\quad TT^*=P .
\]
The map is smooth and covariant under simultaneous conjugation.
Local orthonormal frames \(V_\alpha\) give determinant transitions
\[
 g_{\alpha\beta}=\det(V_\alpha^*V_\beta)\in\U(1),\qquad
 g_{\alpha\beta}g_{\beta\gamma}=g_{\alpha\gamma}.
\]
\end{proposition}
\begin{proof}
For a unit \(v\in\operatorname{im}P_0\),
\(\langle v,Hv\rangle\ge\gamma_0-\delta=\gamma\).
The negative spectrum of \(H\) is at most \(-\gamma\) and its
positive spectrum is at most \(M\).  Thus
\(\gamma\le M\|Pv\|^2-\gamma\|(I-P)v\|^2\), giving the
lower bound on \(G\).  The inertia is unchanged along the norm
ball, so \(P,P_0\) have equal rank.  Direct multiplication gives
\(T^*T=P_0\); equal ranks give \(TT^*=P\).
The positive square root and gapped spectral projector are smooth
and conjugation covariant.  On overlaps the frame comparison is
unitary; composing comparisons proves the determinant cocycle.
\end{proof}



\begin{corollary}[Chart bounds on the inherited finite-link family]
\label{p1b:cor:certified-chart}
On the finite-link family of
Proposition~\ref{p1b:prop:cover-transfer}, retain its certificate bounds
\(g_*=417/7000\), perturbation size \(b<g_*\), and
\(M=11+\epsilon_\rho\). The positive-projector chart satisfies
\[
 P_0P(H)P_0\ge
 \frac{2(g_*-b)}{11+\epsilon_\rho+g_*-b}\,I
 \quad\hbox{on }\im P_0.
\]
Local triviality on this domain follows from these already supplied
bounds.
\end{corollary}
\begin{proof}
The reference gap is \(g_*\), its norm perturbation is at most
\(b\), and its norm bound is \(11+\epsilon_\rho\).
Substitute these three numbers into the preceding proposition.
The resulting strictly positive Gram bound gives its explicit
polar frame chart throughout the stated domain.
\end{proof}



\begin{proposition}[Compression defect]
\label{p1b:result:compression-defect}
For an orthogonal projector \(P\), \(Q=I-P\), and finite operators
\(A,B\),
\[
 [PAP,PBP]-P[A,B]P=-PAQBP+PBQAP .
\]
Thus a relation among parent operators descends without correction
when the omitted cross terms vanish.
\end{proposition}
\begin{proof}
Insert \(I=P+Q\) between the two factors of \(PABP\) and
\(PBAP\), then subtract.  The terms with \(P\) are exactly
the compressed commutator and the remaining two terms are those
printed.  This also proves the last conditional assertion.
\end{proof}



\begin{proposition}[A complete positive massive determinant branch]
\label{p1b:prop:positive-massive-branch}
Let \(V\) be a finite unitary matrix and \(\Gamma\) a Hermitian
involution satisfying \(\Gamma V\Gamma=V^*\).
For every \(a>0\), \(m>0\), define
\[
 D_m=(a^{-1}+m/2)I+(a^{-1}-m/2)V.
\]
Then
\[
 \sigma_{\min}(D_m)\ge\min(m,2/a)>0,\qquad \det D_m>0.
\]
In particular this applies to \(V=\Gamma_5\sgn H\) from the finite
overlap construction and to arbitrary finite products of the displayed
positive-mass determinants.
\end{proposition}
\begin{proof}
Put \(a_0=a^{-1}+m/2\) and \(b_0=a^{-1}-m/2\).
The unitary spectral theorem gives singular values
\(|a_0+b_0e^{i\theta}|\ge a_0-|b_0|=\min(m,2/a)\).
The relation with \(\Gamma\) pairs every nonreal eigenvalue of
\(V\) with its conjugate, whose two determinant factors have
positive product.  The remaining real eigenvalues are \(1,-1\),
with factors \(2/a,m\), both positive.  This proves the entire
mass range, not only \(0<m<2/a\); products remain positive.
\end{proof}




\subsubsection{Supplementary exceptional responses and polar normalization}
\label{p1b:app:exroot-response-calculus}

All adjoints and operator Hilbert--Schmidt norms in this section use
the declared real carrier metrics and ordinary counting on labels.
They are not replaced by a normalized regular trace, a complex spinor
trace or a stationary CAR source pairing. The protected coefficient-metric carrier is defined in Paper I-B; the root sources here are specified
mathematical descendants of
Theorem~\ref{p1:thm:exroot-explicit-e876}.

\begin{theorem}[Vector, scalar-operator and polar response norms]
\label{p1b:thm:exroot-response-norms}
For \(\Phi_r\), let \(V_r=\operatorname{span}\Phi_r\), \(P_r\) its
orthogonal projector in \(A\), \(N_r=|\Phi_r|\), and define
\[
 \mathsf R_r:A\longrightarrow\mathbb R^{\Phi_r},\qquad
 (\mathsf R_ra)_\alpha=\langle\alpha,a\rangle.
\]
Then
\[
\begin{array}{c|c|c|c}
r&N_r&h_r=N_r/r&\mathsf R_r^*\mathsf R_r\\ \hline
6&72&12&12P_6\\
7&126&18&18P_W\\
8&240&30&30I_A
\end{array}
\]
where \(W=1^\perp\). For the scalar source into operator space,
\(\mathsf S_r(s)=s\mathsf R_r\),
\[
 \mathsf S_r^*\mathsf S_r=N_r,\qquad
 \mathsf S_r^+(X)=
 \frac{\langle\mathsf R_r,X\rangle_{\rm HS}}{N_r},\qquad
 \|\mathsf S_r^+\|=N_r^{-1/2}.
\]
In contrast, the vector inverse and polar map satisfy
\[
 \|\mathsf R_r^+\|=h_r^{-1/2},\qquad
 T_r=\mathsf R_r/\sqrt{h_r},\quad
 T_r^*T_r=P_r,\quad\|T_r\|_{\rm HS}^2=r.
\]
The scalar source \(s\mapsto sT_r\) therefore has inverse norm \(r^{-1/2}\).
\end{theorem}
\begin{proof}
The frame commutes with every root reflection because that reflection
permutes the roots. It preserves each root line. If two roots have
nonzero inner product, self-adjointness forces their frame eigenvalues
to agree. Connectedness and spanning therefore make the frame scalar on
\(V_r\), with scalar \(N_r/r\) from its trace. This proves all three
frame identities. For \(E_8\) the direct plane count gives each diagonal
\(2+7\cdot16/4=30\); for \(E_7\) it gives \(2+4\cdot16/4=18\);
off-diagonal terms cancel by signs.

The scalar source has Gram
\(\|\mathsf R_r\|_{\rm HS}^2=\operatorname{Tr}
(\mathsf R_r^*\mathsf R_r)=N_r\). Its displayed rank-one inverse
follows by orthogonal projection onto the source line. The nonzero
singular values of \(\mathsf R_r\) are \(\sqrt{h_r}\), giving the
vector inverse and polar identities.
\end{proof}

\begin{corollary}[Unit operator norm and individual rows]
\label{p1b:cor:exroot-unit-hs-rows}
The unit Hilbert--Schmidt operator
\(\widehat{\mathsf R}_r=\mathsf R_r/\sqrt{N_r}=T_r/\sqrt r\)
has singular values \(r^{-1/2}\) and individual row norms
\(N_r^{-1/2}\). The image projector \(T_rT_r^*\) has rank \(r\) and
diagonal entries \(r/N_r\).
\end{corollary}
\begin{proof}
Each root row initially has norm one. Rescaling and the preceding
singular values prove the first statements, while the squared row norm
of \(T_r\) is \(1/h_r=r/N_r\).
\end{proof}

The scalar inverse $N_r^{-1/2}$ and the vector inverse are different maps on different domains. The scalar expression is the inverse square root of the displayed finite counting coefficient.

\begin{proposition}[Native 240 labels and an antipodal source test]
\label{p1b:prop:exroot-native-240}
In the native real coefficient metric, \(U\cup U\omega\) consists of
240 unit vectors of type \(H_4\oplus H_4\) and has frame \(30I_A\).
Consequently its scalar operator source also has inverse norm
\(240^{-1/2}\), without the golden \(E_8\) lift.
For either antipodal root presentation, the flat label vector
\(\Omega=(1,\ldots,1)\) satisfies
\[
 \mathsf R^*\Omega=0,\qquad
 \frac{|\Omega\rangle\langle\Omega|}{N}\mathsf R=0.
\]
\end{proposition}
\begin{proof}
Independent coordinate signs cancel every off-diagonal frame entry.
The permutation symmetry makes the four native diagonal entries equal;
their trace is \(120\), so each is \(30\). The two orthogonal sheets
give \(30I_A\). In an antipodal set the vector sum is zero, proving
both flat-label identities.
\end{proof}

The scalar column in operator space is not the flat vector in the label
space. A collective quantum source would need its own generated map;
the natural vector evaluation does not excite this flat mode.

\begin{proposition}[Projection and retention have different Grams]
\label{p1b:prop:exroot-project-retain}
Projecting all 240 golden roots onto \(V_6\) gives the operator
\(\mathsf R_8P_6\) with squared Hilbert--Schmidt norm \(180\).
Retaining the \(E_6\) section gives squared norm \(72\).
The native section \(U\cup U\omega\) orthogonal to \(1\) and \(i+j+k\)
has 126 elements but frame
\[
 3P_2\oplus30I_4
\]
of rank six, rather than the \(E_7\) frame \(18I_7\).
\end{proposition}
\begin{proof}
The projected frame is \(30P_6\), of trace \(30\cdot6=180\).
The retained frame is \(12P_6\), of trace \(72\).
In the native section all 120 odd-sheet units remain and the even
section is the six imaginary-plane half-unit roots. Their
two-dimensional frame has trace six and eigenvalue three.
The odd frame remains \(30I_4\).
\end{proof}

\begin{proposition}[Repeated presentations preserve the polar carrier]
\label{p1b:prop:exroot-repeated-presentation}
Repeat each row \(k\) times and let
\(J_k(y)=(y,\ldots,y)/\sqrt k\). Then
\[
 \mathsf R^{(k)}=\sqrt k J_k\mathsf R,\quad
 T^{(k)}=J_kT,\quad J_k^*J_k=I.
\]
The raw scalar Gram becomes \(kN\), while the polar carrier Gram
remains \(P_r\). The transported operator and ports on the polar image
are isometrically equivalent.
\end{proposition}
\begin{proof}
The formulas follow by summing the \(k\) identical row contributions.
Transport of the carrier operator \(D\) gives \(J_kDJ_k^*\) on the
image, and transport of the source gives \(J_kj\); their resolvent
pairing is unchanged.
\end{proof}

This proposition concerns repeated presentations of one carrier with
transported action and ports. It does not exclude genuinely independent
generated labels or a factor inside a generated action member.

The normalized primitive trace is likewise distinct from the ordinary
operator norm, as already fixed by
Corollary~\ref{p1b:cor:v130-owner-amplification-invariance} and
Remark~\ref{p1b:rem:v130-trace-owner-firewall}.



\subsection{Positive finite blocking and exact source calculus}
\label{p1b:app:blocking}
\begin{theorem}[Finite shell covariance, affine support and Legendre flow]
\label{p1b:thm:finite-flow}
For a supplied finite real realization of \(P=B_R^2\), set
\[
 C_{\Lambda,k}=\int_{\Lambda^{-2}}^{k^{-2}}e^{-sP}\,\dd s,
 \qquad 0<k<\Lambda .
\]
Its spectral value is
\((e^{-\lambda/\Lambda^2}-e^{-\lambda/k^2})/\lambda\) for
\(\lambda>0\), and \(k^{-2}-\Lambda^{-2}\) at zero.
If \(\lambda_{\max}\) is finite,
\[
 (k^{-2}-\Lambda^{-2})e^{-\lambda_{\max}/k^2}I
 \le C_{\Lambda,k}\le(k^{-2}-\Lambda^{-2})I .
\]
It commutes with the reciprocal involution. Gaussian convolution
therefore preserves that involution for invariant interactions.

Let \(\mu\) be a nonzero positive finite measure on \(\R^N\),
with minimal full-measure affine support \(\mathcal A=a+E\).
On an open domain with finite dominated source derivatives, define
\[
 W_k(J)=\log\int e^{-x^TR_kx/2+J^Tx}\,\dd\mu(x).
\]
Then \(\ker W_k^{(2)}=E^\perp\), and \(W_k^{(2)}|_E>0\).
For an orthonormal \(T:\R^d\to E\), put
\(R_{k,E}=T^TR_kT\), \(j=T^T(J-R_ka)\). One has
\[
 W_k(J)=J^Ta-\tfrac12a^TR_ka+W_{k,E}(j).
\]
The modified Legendre transform on \(E\) obeys, at fixed mean
coordinate and for differentiable \(R_k\),
\[
 \partial_\ell\Gamma_{k,E}
 =\frac12\Tr_E[(\Gamma_{k,E}^{(2)}+R_{k,E})^{-1}
                  \partial_\ell R_{k,E}] .
\]
\end{theorem}
\begin{proof}
Diagonalize the nonnegative finite matrix \(P\).  Integrating its
scalar exponential gives both spectral formulas; bounding each
integrand between \(e^{-\lambda_{\max}/k^2}\) and one gives
the two operator bounds.  Since \(P=B_R^2\) commutes with the
reciprocal involution, so do this covariance and its centered
Gaussian.  Change variables by the involution in the convolution
to prove its symmetry.
For the measure assertion, its strictly positive tilt has the same
full-measure affine support as \(\mu\).  Differentiating the
logarithm under the assumed dominated integrals gives its covariance
matrix.  Its null directions are precisely those linear functions
constant almost surely, namely \(E^\perp\); every nonzero direction
in \(E\) has strictly positive variance.  Write \(x=a+T\xi\)
and expand its quadratic and linear terms to obtain exactly the
affine reduction formula.  On \(E\) the positive covariance
locally inverts the gradient map.  For
\(\Gamma=j\cdot\phi-W-\phi^TR_{k,E}\phi/2\), its derivative
at fixed \(\phi\) is
\(\frac12\mathbb E[\xi^T\dot R\xi]-\frac12\phi^T\dot R\phi
=\frac12\operatorname{Tr}(\operatorname{Cov}\xi\,\dot R)\).
Legendre inversion gives
\(\operatorname{Cov}\xi=(\Gamma^{(2)}+R_{k,E})^{-1}\), proving
the flow on precisely its regular source/mean domain.
\end{proof}



\begin{corollary}[Invariant Legendre vertices and linearized flow]
\label{p1b:cor:invariant-legendre-vertex-parity}
In Theorem~\ref{p1b:thm:finite-flow}, suppose a real orthogonal
involution \(J\) preserves the complete positive measure and commutes
with the regulator, on an invariant domain of dominated source
derivatives. Then
\[
 Z_k(Jj)=Z_k(j),\quad W_k(Jj)=W_k(j),\quad
 \Gamma_k(J\phi)=\Gamma_k(\phi).
\]
These statements hold on the minimal affine support with its induced
action; an invariant origin may be chosen there. At a fixed background
\(J\bar\phi=\bar\phi\), for \(Jv_i=\varepsilon_i v_i\),
\[
 D^n\Gamma_k|_{\bar\phi}[v_1,\ldots,v_n]=0
 \quad\hbox{if}\quad \prod_i\varepsilon_i=-1,
\]
whenever those derivatives exist. The mixed \(+\!-\) Hessian vanishes;
a \(+\!-\!-\) vertex is allowed. If a differentiable normalized flow
map \(\mathcal R\) is equivariant and has a \(J\)-fixed fixed point
\(\Gamma_*\), its derivative commutes with \(J^*\), the induced
action on perturbations, and splits into the two parity blocks.
\end{corollary}
\begin{proof}
The orthogonal change of integration variable \(x\mapsto Jx\)
preserves the complete measure and quadratic regulator, giving
\(Z_k(Jj)=Z_k(j)\).  It then gives the logarithm equality and,
by changing the conjugate variable in its supremum, the Legendre
equality.  An invariant affine support has an invariant origin:
average any of its points with its image.  Repeated differentiation
at a fixed background multiplies the multilinear vertex by
\(\prod\varepsilon_i\), forcing zero for negative product.
Finally differentiating \(\mathcal R\circ J^*=J^*\circ\mathcal R\)
at its fixed point gives the commuting derivative and its two
invariant eigenspaces.
\end{proof}




\begin{theorem}[Positive blocking contracts oscillation]
\label{p1b:thm:blocking}
For fixed conditional probability measures, define
\(\mathcal B(S)(x)=-\log\int e^{-S(y)}\nu(\dd y\mid x)\),
assuming nonzero finite fibre integrals.
If \(m\le T-S\le M\) essentially, then
\[
 m\le\mathcal B(T)-\mathcal B(S)\le M,\qquad
 \operatorname{TV}(\mu_S,\mu_T)
 \le\tanh(\osc(T-S)/4),
\]
with \(\operatorname{TV}=\frac12\|\mu_S-\mu_T\|_1\).
The same bound holds after exact blocking.
For a compatible tower with
\(\osc(\mathcal B_{n+1\to n}S_{n+1}-S_n)\le\epsilon_n\),
summability \(\sum_n\epsilon_n<\infty\) gives consistent
total-variation limits at each fixed coarse level.
For bounded observables \(|O_i|\le M_i\),
\[
 |\kappa_k^\mu(O_1,\ldots,O_k)-\kappa_k^\nu(O_1,\ldots,O_k)|
 \le2\operatorname{TV}(\mu,\nu)
    \left(\sum_{r=1}^k r!\,\Stirling{k}{r}\right)\prod_iM_i .
\]
\end{theorem}
\begin{proof}
The inequality \(e^{-M}e^{-S}\le e^{-T}\le e^{-m}e^{-S}\)
under each fixed conditional integral gives the first bounds.
For normalized positive measures their likelihood ratio has
maximum/minimum ratio at most \(e^{M-m}\).
For a ratio \(r\in[a,b]\) with mean one, convexity of \(|r-1|\)
puts its largest possible total variation on the two endpoints;
its value is \((1-a)(b-1)/(b-a)\).  Maximizing over the common
scale of \(a,b\), with \(b/a\le e^{M-m}\), gives
\((e^{(M-m)/2}-1)/(e^{(M-m)/2}+1)\), the printed hyperbolic tangent.
A probability pushforward or conditional blocking is a contraction
of total variation, proving the blocked bound.  In the compatible
tower, telescope the pushed measures from successive fine levels.
Their coarse total-variation distances are bounded by the summable
\(\tanh(\epsilon_n/4)\).  Completeness gives a probability limit;
continuity of pushforward gives consistency.
For the last estimate use the partition formula for cumulants.
A partition with \(r\) blocks has coefficient \((r-1)!\).
Each product of block moments changes by at most
\(2r\operatorname{TV}(\mu,\nu)\prod_iM_i\), by telescoping
its \(r\) factors.  Counting partitions with \(\mathrm{Stirling}(k,r)\)
gives exactly the stated finite sum.
\end{proof}



\begin{proposition}[Heat Hessian and a positive block bound]
\label{p1b:prop:heat-block}
On \(\SU(3)\), normalize the quadratic Casimir by
\(C_{\rm fund}=4/3,C_{\rm adj}=3\), and set
\[
 K_t(g)=\sum_\lambda d_\lambda\chi_\lambda(g)e^{-tC_\lambda/2},
 \qquad
 h(t)=\frac18
 \frac{\sum_\lambda d_\lambda^2 C_\lambda e^{-tC_\lambda/2}}
      {\sum_\lambda d_\lambda^2e^{-tC_\lambda/2}} .
\]
Then \(h'(t)=-\operatorname{Var}_t(C)/16<0\).
In the compatible quadratic norm \(q_C\),
\[
 -\log\frac{K_t(e^X)}{K_t(e)}
 =\frac{h(t)}2q_C(X)
 +\frac{h(t)^2+h'(t)+h(t)/4}{40}q_C(X)^2
 +O(\|X\|^6).
\]
If \(0<\kappa=h(t)<1/12\), then
\[
 a=\frac{1-12\kappa}{1-6\kappa}\le K_t(g)
 \le b=\frac1{1-6\kappa}.
\]
A normalized Haar integral of a product of \(N\) such face
factors is consequently between \(a^N\) and \(b^N\).
\end{proposition}
\begin{proof}
Write \(Z(t)=\sum d_\lambda^2e^{-tC_\lambda/2}\) and let
\(\mathbb E_t\) denote its probability weights.  Differentiation
 gives \(h=\mathbb E_tC/8\) and
\(h'=-\operatorname{Var}_t(C)/16<0\); the variance is strictly
positive since the trivial and fundamental terms both occur.
The invariant symmetric quadratic tensor of a representation is
\(C_\lambda/8\) times the metric.  In \(\mathfrak{su}(3)\)
there is no independent symmetric quartic invariant.  Contracting
the symmetrized trace of four generators and using
\(\sum_aT_aT_bT_a=(C_\lambda-3/2)T_b\) gives
\[
 d_\lambda^{-1}\operatorname{Tr}_\lambda X^4
       =\frac{3C_\lambda(2C_\lambda-1)}{160}\,q_C(X)^2.
\]
Thus the quadratic and quartic terms of \(K_t(e^X)/Z(t)\)
are \(-\mathbb E_tC\,q_C/16\) and
\(\mathbb E_t[3C(2C-1)]q_C^2/(160\cdot24)\).
Expanding the negative logarithm and substituting the mean and
variance above gives exactly the stated coefficient.
Odd terms vanish by inversion symmetry.
Every nontrivial Casimir is at least \(4/3\), so
\(Z-1\le(3/4)Z\mathbb E_tC=6\kappa Z\).
Hence \(Z\le(1-6\kappa)^{-1}=b\).
The character bound \(|\chi_\lambda(g)|\le d_\lambda\) gives
\(2-Z\le K_t(g)\le Z\); thus \(K_t(g)\ge2-b=a>0\).
Pointwise multiplication and integration against normalized Haar
measure give the final bounds.
\end{proof}



\begin{proposition}[Bounded complex sources]
\label{p1b:result:bounded-complex-sources}
For a positive probability measure and bounded real observables,
let \(F(J)=\mathbb E e^{\sum J_iO_i}\) and
\(L=\sum_i|J_i|\|O_i\|_\infty\). Then
\[
 |F(J)-1|\le e^L-1 .
\]
In particular \(L<\log2\) gives a common zero-free domain
and an analytic logarithm continuous from \(J=0\).
Fixed-degree source derivatives converge under uniform
bounded-observable moment convergence.
\end{proposition}
\begin{proof}
For \(Y=\sum J_iO_i\), \(|Y|\le L\) almost surely.
The exponential power series gives \(|e^Y-1|\le e^L-1\);
integration proves the estimate.  If \(L<\log2\), the image
lies in the open disk of radius one centered at one, where the
logarithm continuous from one is analytic.  Its source derivatives
are polynomials in bounded-observable moments divided by powers
of the nonzero partition function.  Uniform moment convergence
and the same positive distance from zero prove their convergence
at every fixed degree on compact smaller source domains.
\end{proof}



\subsection{Hopf trace and reciprocal response families}
\label{p1b:app:response-families}
This technical section consolidates mathematical functions and operators used in
the later framework papers. Their scales, functions and domains are
defined here, so their analytic consequences can be studied independently
of their later interpretations. The trace and response families are
supplementary to the main product classifications.

\begin{theorem}[Primitive Hopf heat trace and the full dyadic jet family]
\label{p1b:thm:dyadic-trace}
On the unit sphere, the primitive Hopf charge \(q=1\) horizontal
scalar Laplacian has levels \(4j^2-2\) of multiplicity \(2j\),
\(j\ge1\). Its heat trace and zeta value are
\[
 Z_H(t)=2e^{2t}\sum_{j\ge1}j e^{-4tj^2}
       =\frac1{4t}+\frac13+O(t),\qquad \zeta_H(0)=\frac13.
\]
Define
\(\mathcal J_H=\tfrac12(\Delta_H+2I)^{1/2}\) and
\(D_H=2^{-\mathcal J_H}\). Then \(D_H\) is positive trace class,
\(\Tr D_H=4\), and, for every \(x>0\),
\[
 S_H(x)=\Tr(x^{D_H}-I)
       =\sum_{j\ge1}2j(x^{2^{-j}}-1).
\]
The function \(y\mapsto S_H(e^y)\) is entire and has the complete Taylor
expansion
\[
 S_H(e^y)=\sum_{r=1}^\infty
       \frac{2^{r+1}}{(2^r-1)^2}\frac{y^r}{r!}.
\]
The Fredholm determinant satisfies \(\det(e^{yD_H})=e^{4y}\).
On the positive real axis,
\[
 S_H(x)=2x^{1/2}+4x^{1/4}
             +O(x^{1/8}\log x)\qquad(x\to\infty).
\]
\end{theorem}
\begin{proof}
The charge-one condition retains odd scalar degrees \(k=2j-1\),
with horizontal eigenvalue \(k(k+2)-1=4j^2-2\) and multiplicity
\(k+1=2j\).  Euler--Maclaurin applied to \(xe^{-4tx^2}\) gives
\(\sum_{j\ge1}je^{-4tj^2}=1/(8t)-1/12+O(t)\).
Multiplying by \(2e^{2t}\) gives the constant \(1/3\).
There is no kernel, so the Mellin heat formula gives that constant
as \(\zeta_H(0)\).  Functional calculus gives \(\mathcal J_H=j\)
and the stated eigenvalues \(2^{-j}\) of \(D_H\).
For any bounded complex \(y\),
\(\sum2j|e^{y2^{-j}}-1|\le |y|e^{|y|}\sum2j2^{-j}<\infty\),
locally uniformly.  Thus the function is entire and its Taylor
coefficients are
\(\operatorname{Tr}D_H^r=\sum2j2^{-rj}=2^{r+1}/(2^r-1)^2\).
At \(r=1\) this is four.  The determinant identity is
\(\det e^{yD_H}=\exp(y\operatorname{Tr}D_H)\).
For \(y=\log x>0\), the tail from \(j=3\) is bounded by
\(2ye^{y/8}\sum_{j\ge3}j2^{-j}=2ye^{y/8}\).
Separating the first two terms, including their finite constants,
proves the final asymptotic with its displayed error.
\end{proof}



\begin{theorem}[Reciprocal Laurent and secant families]
\label{p1b:thm:laurent-family}
For every pair \(0<C<U\), define on \(R>0\)
\[
 h(R)=\frac{R^2+CU}{(C+U)R},\qquad
 v(R)=\frac{\dd\log h}{\dd\log R},\qquad
 n(R)=\frac{\log h(R)}{\log(R/U)}.
\]
At \(R=U\), define \(n(U)=(U-C)/(U+C)\). Then
\[
 v(R)=\frac{R^2-CU}{R^2+CU}
     =\tanh\log\frac{R}{\sqrt{CU}},\qquad
 \frac{\dd v}{\dd\log R}=1-v^2.
\]
The response obeys \(h(CU/R)=h(R)\); the index \(n\) extends
smoothly across \(U\), is strictly increasing in \(R\), and
\[
 n(R)\ge0\ \Longleftrightarrow\ R\ge C,\qquad
 -1<n(R)<1.
\]
For \(A=\log\alpha>0\), \(C_j=U\alpha^{-j}\) and
\(R=U\alpha^{-q}\), \(j>0\), the entire indexed family is
\[
 n_j(q)=
 \frac{\log\cosh(jA/2)-\log\cosh((j/2-q)A)}{qA}
       =\int_0^1\tanh[A(j/2-qt)]\,\dd t.
\]
This formula also defines \(q=0\). On \(q\ge0\), its positive part
\(n_j^+=\max(n_j,0)\) vanishes exactly for \(q\ge j\), and on
\(0\le q\le j\),
\[
 R(R/U)^{n_j(R)}=\frac{R^2+UC_j}{U+C_j}.
\]
\end{theorem}
\begin{proof}
Write \(s=\log(R/\sqrt{CU})\).  Then
\(h=2\sqrt{CU}\cosh s/(C+U)\), so its logarithmic derivative
is \(\tanh s\), with derivative \(1-\tanh^2s\), and inversion
sends \(s\) to \(-s\).  Since \(h(U)=h(C)=1\), \(n\) is a
secant slope of the strictly convex function \(\log\cosh s\)
with one endpoint fixed at \(U\).  Its integral of the strictly
increasing derivative proves smooth removable continuation,
strict monotonicity and the strict bounds \(-1<n<1\).
The two roots of \(h-1\) and the sign of \(\log(R/U)\)
 give exactly the stated nonnegative region.  Substitute
\(s=(j/2-q)A\) into that secant formula to obtain the quotient
and its integral representation, including \(q=0\).
For \(q\ge0\), the numerator is positive exactly when
\(|j/2-q|<j/2\), namely \(0<q<j\); its endpoint limit at zero
is positive.  Finally \((R/U)^{n_j}=h\) on the specified
positive branch; multiply by \(R\) to get the polynomial identity.
\end{proof}



\subsubsection{Weighted reciprocal inversion and Euler reconstruction}
\label{p1b:sec:born-euler-calculus}
The following calculus retains both the quadratic and constant parts of the
reciprocal family. All statements concern scalar functions with the displayed
parameter data; the variable ranges over the stated mathematical interval.

\begin{definition}[Weighted reciprocal pullback]
\label{p1b:def:weighted-born-pullback}
Fix \(B>0\) and an interval \(I\subset(0,\infty)\) preserved by
\(i_B(R)=B/R\). For \(F\in C^1(I)\), set
\[
 (\mathcal T_BF)(R)=\frac{R^2}{B}F(B/R),\qquad
 \mathcal E F=R F'-F.
\]
Writing \(A(R)=F(R)/R\), one has
\((\mathcal T_BF)(R)/R=A(B/R)\). Thus the weighted parity of \(F\)
is the ordinary reciprocal parity of its normalized coefficient \(A\).
\end{definition}

\begin{theorem}[Reciprocal involution and Euler anticommutation]
\label{p1b:thm:born-euler-anticommutation}
On the stated domain,
\begin{equation}
 \label{p1b:eq:born-euler-anticommutation}
 \boxed{\mathcal T_B^2=\mathrm{id},\qquad
        \mathcal E\mathcal T_B=-\mathcal T_B\mathcal E.}
\end{equation}
In particular, \(\mathcal E\) reverses weighted reciprocal parity.
On the quadratic polynomial space \(\mathcal P_2=\operatorname{span}\{1,R,R^2\}\),
\[
 \mathcal T_B(1)=R^2/B,\quad \mathcal T_B(R)=R,\quad
 \mathcal T_B(R^2)=B,
 \qquad
 \mathcal E(1)=-1,\quad \mathcal E(R)=0,\quad
 \mathcal E(R^2)=R^2.
\]
For \(Q_B=R^2+B\), its uniquely determined normalized odd partner is
\begin{equation}
 \label{p1b:eq:born-even-odd-euler-pair}
 \boxed{P_B:=\mathcal E Q_B=R^2-B,\qquad
        \mathcal E P_B=Q_B.}
\end{equation}
With \(u=\log(R/\sqrt B)\),
\[
 Q_B/R=2\sqrt B\cosh u,\qquad
 P_B/R=2\sqrt B\sinh u,\qquad
 \dd(P_B/R)=(Q_B/R)\,\dd u.
\]
For the family of Theorem~\ref{p1b:thm:laurent-family}, \(B=CU\) and
\(P_B/R=(C+U)h(R)v(R)\).
\end{theorem}
\begin{proof}
Substitute \(B/R\) twice to obtain \(\mathcal T_B^2=I\).
Differentiation gives
\(\mathcal E\mathcal T_BF=(R^2/B)F(B/R)-R F'(B/R)\),
which is \(-\mathcal T_B\mathcal EF\).
Applying these operations to \(1,R,R^2\) proves their tables
and the two partners.  Dividing those polynomials by \(R\)
and writing \(R=\sqrt B e^u\) gives their hyperbolic expressions,
whose differentiation proves the final differential identity.
The last relation is the preceding Laurent formula substituted
with \(B=CU\).
\end{proof}



\begin{proposition}[Different kernels of Hessian and Euler reconstruction]
\label{p1b:prop:born-hessian-quotient-obstruction}
On \(\mathcal P_2\), the subspace
\(\ker D^2=\operatorname{span}\{1,R\}\) is not preserved by
\(\mathcal T_B\); hence \(\mathcal T_B\) induces no map on
\(\mathcal P_2/\ker D^2\). In contrast,
\(\ker\mathcal E=\operatorname{span}\{R\}\) is preserved, and
\(\mathcal E\) induces an invertible parity-reversing map on
\(\mathcal P_2/\operatorname{span}\{R\}\), whose square is the identity.
Among operators \(aR D+b\,\mathrm{id}\) with constant real coefficients,
the conditions \(D_0[R]=0\), \(D_0[R^2]=R^2\) uniquely select
\(D_0=\mathcal E\).
\end{proposition}
\begin{proof}
The second derivative kills exactly the constant and linear
polynomials, but \(\mathcal T_B1=R^2/B\) is not killed.
Thus that kernel is not an invariant subspace.
The Euler kernel solves \(RF'=F\), hence is precisely the line
\(\mathbb RR\).  The displayed polynomial table preserves this
line and induces an involution on its complementary constant and
quadratic classes; anticommutation gives parity reversal.
Finally \((aRD+b)[R]=(a+b)R=0\) and its value on \(R^2\)
is \((2a+b)R^2=R^2\), forcing \(a=1,b=-1\).
\end{proof}



\begin{theorem}[Minimal reciprocal refinement of a linear readout]
\label{p1b:thm:reciprocal-readout-refinement}
Let \(J^2=I\) on a real vector space \(V\), and \(L:V\to W\) be linear.
The joint readout
\[
 \widehat L(v)=(Lv,LJv),\qquad
 \widehat L(Jv)=\operatorname{swap}\widehat L(v)
\]
has kernel \(\ker L\cap J(\ker L)\). This is the largest \(J\)-invariant
linear subspace contained in \(\ker L\), so its quotient is the coarsest
linear quotient retaining \(L\) on which \(J\) descends.
\end{theorem}
\begin{proof}
The two components vanish exactly on
\(\ker L\cap J\ker L\).  This subspace is \(J\)-invariant.
Every \(J\)-invariant subspace \(N\subset\ker L\) also satisfies
\(N=JN\subset J\ker L\), so is contained in that intersection.
The maximal kernel gives precisely the coarsest stated quotient;
the swap formula follows immediately from \(J^2=I\).
\end{proof}


For a nonlinear readout the two joint equalities define the corresponding
invariant equivalence relation at the set level. A smooth quotient additionally
requires its regularity and constant-rank hypotheses.

\begin{theorem}[Exact Euler reconstruction from paired Hessians]
\label{p1b:thm:paired-hessian-euler}
Let \(I\subset(0,\infty)\) be a connected reciprocal interval with nonempty
interior, and let \(F\in C^2(I)\), using a \(C^2\) extension at closed
endpoints. The paired Hessian
\(\mathscr H_BF=(F'',(\mathcal T_BF)'')\) satisfies
\[
 \boxed{\ker\mathscr H_B=\operatorname{span}\{R\},\qquad
 \mathcal EF(R)=\tfrac12\bigl[R^2F''(R)
              -B(\mathcal T_BF)''(B/R)\bigr].}
\]
In particular \(F=\lambda R+aR^2+\beta\) has paired readout
\((2a,2\beta/B)\).
\end{theorem}
\begin{proof}
Twice differentiating the weighted transform gives
\[
 (\mathcal T_BF)''(s)=\frac2B F(B/s)
       -\frac2s F'(B/s)+\frac B{s^2}F''(B/s).
\]
Set \(s=B/R\) and rearrange to obtain the reconstruction identity.
If both second derivatives vanish, first \(F=aR+b\), and
then \((\mathcal T_BF)''=2b/B\) forces \(b=0\).
Conversely a multiple of \(R\) is in both kernels.
Substitution of the quadratic family proves its paired readout.
\end{proof}


This refinement retains a constant discarded by one Hessian. It does not
assign the value of the second readout from the first.


\begin{corollary}[Unique quadratic reconstruction and endpoint criterion]
\label{p1b:cor:born-euler-quadratic-reconstruction}
Let \(0<C<U\), \(B=CU\), and \(F(R)=\lambda R+R^2+\beta\).
Set \(W=\mathcal EF\) and \(w=W/R\). The following statements are equivalent:
\[
 \begin{gathered}
 w(B/R)=w(R)\quad(R>0),\qquad w(U)=w(C),\\
 \beta=-B,\qquad \mathcal EF=Q_B.
 \end{gathered}
\]
Moreover the complete \(C^1\) solution family of
\(\mathcal EF=Q_B\) on a connected positive interval is
\[
 \boxed{F(R)=R^2-B+\lambda R.}
\]
Thus fixing the linear coefficient fixes a unique solution. The reciprocal
endpoint criterion above applies to the quadratic family, not to arbitrary
functions or to an unspecified dressed continuation.
\end{corollary}
\begin{proof}
Here \(\mathcal EF=R^2-\beta\) and \(w=R-\beta/R\).
Reciprocal equality is therefore equivalent to \(\beta=-B\).
Equality only at \(C,U\) gives
\((U-C)(1+\beta/(CU))=0\), the same condition since \(U>C\).
For the full differential equation divide by \(R^2\):
\((F/R)'=1+B/R^2\).  Integration on a connected interval gives
\(F/R=R-B/R+\lambda\), including exactly the stated kernel.
\end{proof}



\subsubsection{Evaluation and reciprocal contact lifts}
\label{p1b:sec:reciprocal-contact-lifts}
These are local results for evaluation pairs and coefficient graphs.
The scalar torus supplies local such charts, but its nonzero primitive
two-class precludes a global real primitive.

\begin{proposition}[Evaluation exchange and the Legendre contact lift]
\label{p1b:prop:evaluation-legendre-contact-lift}
On a finite-dimensional vector/dual pair with coordinates \(q,p\), set
\(\vartheta=dz-p\,dq\). The exchange \(q'=p,\ p'=q\) has the lifts
\[
 \mathfrak L_c(q,z,p)=(p,pq-z+c,q),\qquad
 \mathfrak L_c^*\vartheta^\vee=-\vartheta,\qquad
 \mathfrak L_c^2=I .
\]
On a connected chart these are all lifts with these coordinate and
exact-form requirements. For \(F\in C^2\) with locally invertible
gradient, the image of \(z=F(q),p=dF_q\) is the graph of
\[
 F^\vee(p)=p\,q(p)-F(q(p))+c,\qquad dF^\vee=q\,dp.
\]
\end{proposition}
\begin{proof}
Write an arbitrary lift with these exchanged coordinates as
\((p,Z(q,z,p),q)\).  The exact-form requirement is
\(dZ-q\,dp=-dz+p\,dq\), so \(d(Z+z-pq)=0\).
Connectedness gives the single constant \(c\), proving both
existence and exhaustion.  Applying the formula twice returns
\((q,z,p)\).  On the graph differentiate \(p q-F(q)+c\);
the terms \(p\,dq-dF\) cancel, leaving \(q\,dp\).
Local gradient invertibility supplies \(q(p)\) as stated.
\end{proof}



\begin{proposition}[Connection-compatible homogeneous reciprocal lifts]
\label{p1b:prop:reciprocal-connection-lifts}
On the real coefficient line with \(\nabla=d-du\), let a homogeneous
lift covering \(u\mapsto-u\) act as \(TF(u)=a(u)F(-u)\).
It preserves the connection and squares to the identity exactly when
\[
 a(u)=\varepsilon e^{2u},\qquad \varepsilon\in\{+1,-1\}.
\]
Preservation of the positive cone selects \(+1\).
\end{proposition}
\begin{proof}
On coefficient sections the pullback of a one-form changes sign
under \(u\mapsto-u\).  Equating
\(\nabla(aF(-u))\) with the pulled-back \(a\nabla F\)
gives \(a'=2a\), so \(a=c e^{2u}\).
Squaring the lift gives \(c^2=1\).  These solutions are nowhere
zero; precisely its positive sign preserves the positive cone.
\end{proof}



\begin{proposition}[Affine reciprocal contact lifts]
\label{p1b:prop:reciprocal-affine-contact-lifts}
For \(h\in C^2(I)\) and \(\varepsilon=\pm1\), define
\[
 R'=\frac BR,\quad
 z'=\varepsilon\frac B{R^2}z+h(B/R),\quad
 p'=\varepsilon(2z/R-p)+h'(B/R).
\]
Then, for \(\vartheta=dz-p\,dR\),
\[
 \vartheta'=\varepsilon\frac B{R^2}\vartheta,\qquad
 \mathcal A_{\varepsilon,h}F=\varepsilon\mathcal T_BF+h .
\]
The lift is involutive exactly when \(h+\varepsilon\mathcal T_Bh=0\).
Its action on the Euler remainder is
\[
 \mathcal E(\mathcal A_{\varepsilon,h}F)
 =-\varepsilon\mathcal T_B(\mathcal EF)+\mathcal Eh.
\]
\end{proposition}
\begin{proof}
Use \(dR'=-B\,dR/R^2\) in the three coordinate formulas.
The \(z\,dR\) terms and the derivatives of \(h\) cancel,
leaving \(dz'-p' dR'=\varepsilon B(dz-pdR)/R^2\).
On a graph set \(R=B/R'\); its new value is
\(\varepsilon R'^2F(B/R')/B+h(R')\).
Applying this action twice adds \(h+\varepsilon\mathcal T_Bh\)
to \(F\), proving the involution criterion.
Euler anticommutation and linearity give its final action.
\end{proof}


The homogeneous sign classification therefore does not exhaust affine
action-value translations.

\begin{corollary}[Affine reciprocity fixing a linear reference section]
\label{p1b:cor:eh-fixed-affine-reciprocity}
Fix \(U\in\R\) and the entire reference section \(UR\). If
\(\mathcal JF=-\mathcal T_BF+h\) fixes this section, then
\[
 h=2UR,\qquad \mathcal JF=2UR-\mathcal T_BF,\qquad
 \mathcal J^2=I,\qquad
 \mathcal E(\mathcal JF)=\mathcal T_B(\mathcal EF).
\]
For \(F=UR+R^2+\beta\), \(\mathcal JF=F\) is equivalent to
\(\beta=-B\).
\end{corollary}
\begin{proof}
The weighted transform fixes \(UR\), so fixing that entire
section forces \(h=2UR\), with no remaining affine freedom.
This \(h\) satisfies the previous involution criterion at sign
minus.  Since \(\mathcal E(UR)=0\), the Euler formula reduces
to the displayed one.  For the quadratic family compare the
constant and quadratic coefficients in \(2UR-\mathcal T_BF=F\);
both give \(\beta=-B\).
\end{proof}



\subsubsection{Reciprocal convexity and response dressing}
\label{p1b:sec:reciprocal-convexity}
The following classification concerns functions on their stated intervals;
global and finite-interval conclusions are kept distinct.

\begin{theorem}[Convexity on reciprocal partners]
\label{p1b:thm:reciprocal-convexity-test}
Let \(F=UR+F_{\rm nl}\in C^2(I)\) and
\(\mathcal T_BF_{\rm nl}=-F_{\rm nl}\). Put
\[
 u=\log(R/\sqrt B),\qquad A(u)=F_{\rm nl}(R)/R,\qquad
 h(u)=\mathcal EF(R)/R=A'(u).
\]
Then \(A\) is odd and \(h\) is even. Strict positivity of the Hessian
at both reciprocal arguments is equivalent to
\[
 \boxed{h(u)>|h'(u)|,\qquad F''(R)=(h+h')/R.}
\]
If \(F_{\rm nl}\) is a finite integer Laurent polynomial, it has precisely
the form
\[
 F_{\rm nl}(R)=\sum_{m=1}^N a_m
       \bigl(R^{m+1}-B^mR^{1-m}\bigr).
\]
\end{theorem}
\begin{proof}
The weighted oddness is \(A(-u)=-A(u)\), so \(h=A'\) is even.
Differentiate \(F_{\rm nl}=RA(u)\) twice with \(du/dR=1/R\)
to obtain \(F''=(h+h')/R\).  At the reciprocal argument its
numerator is \(h-h'\), so both are strictly positive exactly
when \(h>|h'|\).  A Laurent monomial \(R^r\) is paired by the
weighted transform with \(B^{r-1}R^{2-r}\).  The fixed exponent
\(r=1\) has no odd coefficient.  Pairing the remaining integer
exponents as \(r=1\pm m\) gives exactly the finite list displayed.
\end{proof}



\begin{corollary}[Exact first higher-pair convexity window]
\label{p1b:cor:reciprocal-cubic-window}
For \(a_{\rm nl}>0,\ B=UC,\ 0<C<U\), take
\[
 F=UR+a_{\rm nl}(R^2-B)+b(R^3-B^2/R).
\]
If \(U>\sqrt3\,C\), the condition \(F''(R)>0\) for every
\(R\in[C,U]\) is equivalent to
\[
 \boxed{-\frac{a_{\rm nl}U}{3U^2-C^2}
       <b<\frac{a_{\rm nl}C}{U^2-3C^2}.}
\]
\end{corollary}
\begin{proof}
The Hessian is \(2a_{\rm nl}+2b\psi(R)\), where
\(\psi(R)=3R-B^2/R^3\) has strictly positive derivative.
Its endpoint values are \(-(U^2-3C^2)/C<0\) and
\((3U^2-C^2)/U>0\).  Positivity of this affine function of
\(\psi\) on the entire interval is equivalent to its positivity
at both endpoints.  Solving those two strict inequalities for
\(b\) gives exactly the window stated.
\end{proof}



\begin{theorem}[Global finite integer Laurent rigidity]
\label{p1b:thm:global-integer-laurent-rigidity}
If the finite integer Laurent family in
Theorem~\ref{p1b:thm:reciprocal-convexity-test} has
\(F''(R)>0\) for every \(R>0\), then
\[
 F=UR+a_{\rm nl}(R^2-B),\qquad a_{\rm nl}>0.
\]
\end{theorem}
\begin{proof}
If a highest nonzero pair has index \(N\ge2\), its second
derivative at infinity has leading term
\(a_NN(N+1)R^{N-1}\), so positivity forces \(a_N>0\).
At zero the same pair instead has leading term
\(-a_NB^NN(N-1)R^{-N-1}\), which is negative.
This contradiction removes every \(N\ge2\).
The remaining pair is \(a_1(R^2-B)\), whose Hessian is the
constant \(2a_1\); strict positivity forces \(a_1>0\).
\end{proof}


The hypotheses are essential. The preceding corollary permits nonzero
higher pairs on a finite interval. Outside the integer Laurent class,
\(R^{3/2}-\sqrt B\,R^{1/2}\) is weighted odd and has Hessian
\((3R+\sqrt B)/(4R^{3/2})>0\); \(R\log(R/\sqrt B)\) has
Hessian \(1/R>0\). Thus no global functional uniqueness follows
without the specified class.

\begin{proposition}[Amplitude and dressing in Euler response reconstruction]
\label{p1b:prop:euler-response-dressing-compatibility}
For \(F=UR+a_{\rm nl}R^2+\beta\), weighted evenness of
\(\mathcal EF\) is equivalent to \(\beta=-a_{\rm nl}B\);
then \(\mathcal EF=a_{\rm nl}(R^2+B)\).
For a \(C^2\) response \(Q_d\) on a connected positive interval,
the reconstruction \(\mathcal EF=Q_d\) implies
\[
 F''=\frac{Q_d'}R.
\]
Its agreement with the separate Hessian assignment \(F''=Q_d''\)
holds identically exactly when \(Q_d=aR^2+b\). Finally, for \(D\in C^1\),
\[
 \boxed{\mathcal E[D(R)(R^2-B)]
       =D(R)(R^2+B)+RD'(R)(R^2-B).}
\]
\end{proposition}
\begin{proof}
On the quadratic family \(\mathcal EF=a_{\rm nl}R^2-\beta\);
comparison of its weighted reciprocal coefficients gives
\(\beta=-a_{\rm nl}B\).  Differentiate \(RF'-F=Q_d\)
once to get \(RF''=Q_d'\).  Agreement with \(Q_d''\)
is the differential equation \(RQ_d''=Q_d'\), whose complete
solutions on a connected positive interval are \(aR^2+b\).
The product rule for \(R\,d/dR-1\) gives the last formula,
including the derivative of the dressing.
\end{proof}






\begin{theorem}[Dressed response, convexity and junction measures]
\label{p1b:thm:dressed-convexity}
Let \(\alpha>1\), \(A=\log\alpha\), \(\lambda>0\),
\(c_j=\alpha^{-j}\), \(x=R/U>0\), and \(q=-\log x/A\).
On \(x\ge c_7\) define
\[
 D(x)=\exp[\lambda\max(n_2(q),0)],\qquad
 \widehat f(x)=x+D(x)\frac{x^2+c_7}{1+c_7}-2c_7.
\]
On \(0<x\le c_7\) extend by \(\widehat f(x)=x-c_7\).
This is a continuous function with derivative jumps at \(c_2,c_7\).
For \(c_2<x\le1\), put \(m(q)=n_2(q)\), \(r=c_7/x^2\). Then
\[
 \frac{1+c_7}{D}\widehat f''
 =2-\frac{(3-r)\lambda m'}A+
       \frac{1+r}{A^2}\{\lambda m''+\lambda^2(m')^2\}.
\]
If \(\alpha^3>3\), then
\[
 \widehat f''\ge\frac{2D}{1+c_7}>0
       \quad(c_2<x\le1),\qquad
 \widehat f''=\frac2{1+c_7}>0\quad(c_7<x<c_2).
\]
On \(0<x\le1\), its distributional second derivative is a positive
measure. The two atomic weights are
\[
 w_2=\frac{\lambda\tanh A}{2Ac_2}
               \frac{c_2^2+c_7}{1+c_7},\qquad
 w_7=\frac{2c_7}{1+c_7}.
\]
Convolution locally with any nonnegative smooth approximate identity
produces smooth convex approximants.
\end{theorem}
\begin{proof}
On the dressed interval \(q\in[0,2)\) write
\(m(q)=\int_0^1\tanh[A(1-qt)]dt\).  Then
\(m'=-A\int_0^1t\operatorname{sech}^2[A(1-qt)]dt\le0\) and
\[
 m''=-2A^2\int_0^1t^2\tanh[A(1-qt)]
                \operatorname{sech}^2[A(1-qt)]dt\ge2A m'.
\]
The last bound uses \(t^2\tanh\le t\) and is valid even
when the integrand changes sign.  Differentiating
\(D=e^{\lambda m}\) with \(dq/dx=-1/(Ax)\) gives the printed
Hessian.  Here \(r\le e^{-3A}<1/3\), so the combined two
terms linear in \(m',m''\) are bounded below by
\(\lambda(-1+3r)m'/A\ge0\); the square term is nonnegative.
This proves the lower bound.  Below \(c_2\), \(D=1\) gives
the other Hessian.  At \(q=2\), \(m=0\) and
\(m'=-\tanh A/2\).  The jump in \(D'\) is therefore
\(\lambda\tanh A/(2Ac_2)\), giving \(w_2\).
At \(c_7\), the two displayed values agree and the derivative
jump is \(2c_7/(1+c_7)=w_7\).
Both atoms and both smooth densities are nonnegative; hence the
distributional second derivative is a positive measure.
Convolution with a nonnegative smooth kernel preserves that
positivity on every compact interior interval, as asserted.
\end{proof}



\begin{theorem}[All-order trace differentiation and controlled inverse]
\label{p1b:thm:trace-calculus}
On an interval let \(f\) be smooth, \(F=f'\), \(G=f''>0\),
\(\mathcal T(R)=2f(R)-RF(R)\), and
\(N=\mathcal T'=F-RG\).
Define \(W\) by \(W'(F)=c^2\mathcal T(R(F))/3\), where \(c>0\).
Writing \(H_3=f'''\), \(H_4=f''''\), one obtains
\begin{align*}
 W''&=\frac{c^2}3\left(\frac FG-R\right),&
 W'''&=-\frac{c^2FH_3}{3G^3},\\
 W''''&=-\frac{c^2}{3G^5}
           (G^2H_3+FGH_4-3FH_3^2).
\end{align*}
Every higher derivative is given recursively by
\(W^{(k+1)}=G^{-1}\dd_R W^{(k)}\).
If \(|N-1|\le\eta<1\), then \(\mathcal T\) is strictly
increasing and
\[
 (1-\eta)|R_2-R_1|
 \le|\mathcal T(R_2)-\mathcal T(R_1)|
 \le(1+\eta)|R_2-R_1|.
\]
Its inverse on its image has derivative \(1/N\).
For \(J=[R_c-d,R_c+d]\) inside that interval and
\(|\sigma-\mathcal T(R_c)|\le(1-\eta)d\), the iteration
\[
 R_{k+1}=R_k+\sigma-\mathcal T(R_k),\quad R_0=R_c,
 \qquad
 |R_k-R_*|\le\frac{\eta^k}{1-\eta}|R_1-R_0|
\]
stays in \(J\) and converges to the unique root
\(\mathcal T(R_*)=\sigma\).

For \(f(R)=U\widehat f(R/U)\) from
Theorem~\ref{p1b:thm:dressed-convexity}, its trace response has
the exact downward junction jump
\[
 \mathcal T(C_2^+)-\mathcal T(C_2^-)
 =-U\,\frac{p_2(c_2^2+c_7)}{1+c_7},\qquad
 p_2=\frac{\lambda\tanh A}{2A}>0.
\]
The inverse assertion applies on its separate smooth intervals.
\end{theorem}
\begin{proof}
Since \(dF/dR=G>0\), differentiation in \(F\) is
\(G^{-1}d/dR\).  Apply it successively to
\(W'=c^2(2f-RF)/3\) to obtain all three formulas and the
recursion.  The bound on \(N=\mathcal T'\) integrates to the
two Lipschitz bounds.  For the iteration map \(\Psi(R)=R+\sigma-
\mathcal T(R)\), \(|\Psi'|\le\eta\) and
\(|\Psi(R_c)-R_c|\le(1-\eta)d\); therefore \(\Psi(J)\subset J\).
The contraction theorem proves its unique fixed point and the
geometric a posteriori error bound.  At the dressed junction
\(f\) is continuous and \(f'\) jumps upward by
\(p_2(c_2^2+c_7)/[c_2(1+c_7)]\).
Only \(-Rf'\) jumps in the trace, with \(R=Uc_2\),
giving the downward jump printed.  It is not removed by the
interval inverse theorem.
\end{proof}



\begin{theorem}[Parameter-uniform positivity of the trace derivative]
\label{p1b:thm:uniform-trace-bound}
For the dressed family of Theorem~\ref{p1b:thm:dressed-convexity},
on \(c_2<x\le1\) put \(N=\widehat f'-x\widehat f''\).
Then
\[
 |N-1|\le
 \eta(A,\lambda):=
 e^{\lambda-A}(1+e^{-3A})
       \left(\frac{20}{3}\lambda+4\lambda^2\right).
\]
Whenever \(\alpha^3>3\) and \(\eta(A,\lambda)<1\), the
trace map is strictly increasing on each smooth active interval and
\((f'-Rf'')/(3f'')>0\) there.
For the framework's mathematical parameter member
\[
 A=4\pi^2+\frac1{6\sqrt2},\qquad
 \lambda=\log8+\frac A2,
\]
both inequalities follow from these definitions; no separate positivity
assumption is needed. In particular \(\eta(A,\lambda)<10^{-3}\).
\end{theorem}
\begin{proof}
On \(q\in[0,2]\), \(x=e^{-Aq}\), \(r=c_7/x^2\le e^{-3A}\),
and \(0\le m\le1\), so \(D\le e^\lambda\).
The integral formulas in the preceding proof give the uniform
bounds
\[
 x|m'|/A\le2e^{-A},\qquad
 x|m''|/A^2\le\tfrac83e^{-A},\qquad
 x(m')^2/A^2\le4e^{-A}.
\]
Indeed for \(q\le1\), use
\(\operatorname{sech}^2[A(1-qt)]\le4e^{-2A(1-q)}\)
and integrate \(t,t^2\).  Multiplication by \(x\) gives these
three bounds, including the squared one.  For \(1\le q\le2\)
use \(x\le e^{-A}\), \(\int t=1/2\), \(\int t^2=1/3\),
and \(\operatorname{sech}^2\le1\), which give smaller bounds.
Direct differentiation yields
\[
 N-1=\frac{Dx}{1+c_7}
 \left[\frac{2\lambda(1-r)m'}A
       -\frac{(1+r)\{\lambda m''+\lambda^2(m')^2\}}{A^2}\right].
\]
Substituting the three bounds, using \(|1-r|\le1+r\), gives
exactly \(\eta(A,\lambda)\).  Below \(c_2\) the smooth
undressed interval has \(N=1\).  Positivity of \(N\) and
of \(f''\) gives both stated conclusions on each interval.
For the specific member, \(39<A<40\), \(\lambda<23\),
and \(e^{\lambda-A}=8e^{-A/2}\).
Thus
\(\eta<16e^{-19.5}((20/3)23+4\cdot23^2)<10^{-3}\);
these elementary numerical inequalities can for example be checked
by the positive exponential series.  Also \(e^{3A}>3\).
\end{proof}



\begin{corollary}[Exact quadratic lower secant-comparison response]
\label{p1b:cor:quadratic-response}
On \(C_7<R<C_2\), the exact reciprocal/secant comparison family of this technical section has
\[
 f(R)=R+\frac{R^2+UC_7}{U+C_7}-2C_7,\qquad
 R_0=\frac{2C_7(U+2C_7)}{U+C_7},
\]
and
\[
 \mathcal T(R)=R-R_0,\qquad
 W(F)=\frac{c^2(U+C_7)}{12}(F-F_0)^2+\mathrm{constant},
 \qquad F_0=1+\frac{2R_0}{U+C_7}.
\]
All \(W\)-derivatives of order at least three vanish on that interval.  This
corollary is a finite-scale secant-comparison identity only.  A secant derivative and a local parameter jet are distinct operations and are identified only when their displayed parameter families agree.
\end{corollary}
\begin{proof}
On this interval \(D=1\); substitute its polynomial into
\(2f-Rf'\) to get \(R-R_0\).
Since \(F=1+2R/(U+C_7)\), this trace is
\((U+C_7)(F-F_0)/2\).  Integrating
\(W'=c^2\mathcal T/3\) gives the displayed quadratic.
Consequently its third and all higher derivatives vanish.
No equality with a separately declared fixed-depth Hessian was
used in this algebraic comparison.
\end{proof}



\subsection{Full conformal and boundary-response calculus}
\label{p1b:app:full-response}
The local four-dimensional coframe of
Section~\ref{p1b:app:marked-alternatives} supports a mathematical curvature
functional after a metric and density are specified. The following
identities keep that density intact. They also describe such
functionals on any four-dimensional Riemannian manifold.

\begin{theorem}[Complete conformal Hessian and factorization]
\label{p1b:thm:conformal-hessian}
Let \(\bar g\) be a four-dimensional Riemannian metric of constant
scalar curvature \(R\). Let \(\Phi\) be a \(C^2\) real density
near \(R\), and take a compact manifold without boundary or
compactly supported variations. For
\[
 I[g]=\int\Phi(R[g])\,\dd\vol_g,\qquad
 g=e^{2\sigma}\bar g,\qquad P=-\bar\Box,
\]
the Hessian defined by
\(\tfrac12\int\sigma\mathcal H_\sigma\sigma\,\dd\vol_{\bar g}\)
is
\[
 \mathcal H_\sigma
 =16\Phi-12R\Phi'+12\Phi'P+4\Phi''(R-3P)^2.
\]
On a background satisfying \(R\Phi'-2\Phi=0\) and
\(\Phi''\ne0\), it factors as
\[
 \mathcal H_\sigma=36\Phi''(P-R/3)(P+\mu^2),\qquad
 \mu^2=\frac{\Phi'-R\Phi''}{3\Phi''}.
\]
For \(\Phi_k(R)=k^4\varphi(R/k^2)\), substituting \(P=k^2\zeta\)
gives the corresponding dimensionless polynomial
\[
 k^{-4}\mathcal H_\sigma
 =16\varphi-12r\varphi'+12\zeta\varphi'
          +4\varphi''(r-3\zeta)^2,\qquad r=R/k^2.
\]
\end{theorem}
\begin{proof}
The exact conformal transformation is
\(R[e^{2\sigma}\bar g]=e^{-2\sigma}(R+6P\sigma-6|d\sigma|^2)\)
and its volume factor is \(e^{4\sigma}\).
The linear curvature variation is \(-2R\sigma+6P\sigma\);
its quadratic term is
\(2R\sigma^2-12\sigma P\sigma-6|d\sigma|^2\).
Expand \(e^{4\sigma}\Phi\) through degree two.
The \(\Phi'\) terms combine, after integration by parts, to
\(-6R\Phi'\sigma^2+6\Phi'\sigma P\sigma\);
the other terms are \(8\Phi\sigma^2\) and
\(2\Phi''[(R-3P)\sigma]^2\).
Doubling the quadratic coefficient gives the stated Hessian.
Under \(R\Phi'=2\Phi\), direct multiplication gives the two
factors and their mass parameter.  Finally
\(\Phi_k'=k^2\varphi'\), \(\Phi_k''=\varphi''\), and
\(R=k^2r,P=k^2\zeta\) give the dimensionless polynomial.
\end{proof}



\begin{theorem}[Exact Schur residual at every parameter order]
\label{p1b:thm:all-order-schur}
For a Hermitian block family
\[
 H(y)=\begin{pmatrix}A(y)&B(y)^*\\B(y)&C(y)\end{pmatrix},
 \qquad C(y)\text{ invertible},\qquad
 K(y)=A(y)-B(y)^*C(y)^{-1}B(y),
\]
put \(Z_0=-C(0)^{-1}B(0)\), \(L_0=\binom I{Z_0}\) and
\(R_0(y)=B(y)+C(y)Z_0\). Then
\[
 K(y)=L_0^*H(y)L_0-R_0(y)^*C(y)^{-1}R_0(y).
\]
If the family is differentiable at zero,
\[
 K(y)-K(0)-yK'(0)
 =L_0^*[H(y)-H(0)-yH'(0)]L_0
       -R_0(y)^*C(y)^{-1}R_0(y).
\]
This identity is exact throughout the common invertibility interval.
If \(C(y)>0\) and \(L_0^*H(y)L_0\) is affine in \(y\),
then \(K(y)\) is its matching affine function exactly when
\(R_0(y)=0\) throughout that interval.
\end{theorem}
\begin{proof}
Complete the block square for an arbitrary retained vector:
\[
 \begin{pmatrix}I\\Z_0\end{pmatrix}^{\!*}H
 \begin{pmatrix}I\\Z_0\end{pmatrix}
 =A-B^*C^{-1}B+(B+CZ_0)^*C^{-1}(B+CZ_0).
\]
This is the first identity.  At zero its residual is zero,
so the residual quadratic has value and first derivative zero.
Subtracting those two jets proves the second identity, with
only first differentiability required there.  If \(C>0\),
the residual quadratic vanishes exactly when \(R_0=0\).
For affine compressed data, this is exactly the criterion for
the Schur function to be its matching affine function.
\end{proof}



\subsubsection{Boundary-restricted tangent groupoid and exact Gaussian return gluing}
\label{p1b:subsec:v124-boundary-groupoid-gluing}

The pair- and tangent-groupoid constructions above admit a canonical
restriction to any supplied closed embedded boundary.  Combined with
stationary Schur elimination, this gives an exact mathematical boundary
process calculus.  

\begin{theorem}[Boundary restriction of the pair and tangent groupoids]
\label{p1b:thm:v124-boundary-tangent-groupoid}
Let $K$ be a Riemannian manifold and let $\Sigma\hookrightarrow K$ be a closed
totally geodesic embedded submanifold.  Then
\[
 \Sigma\times\Sigma\rightrightarrows\Sigma
\]
is a closed Lie subgroupoid of the pair groupoid $K\times K\rightrightarrows K$.
Moreover its tangent groupoid
\[
 \boxed{
 \mathcal G_\Sigma
 =T\Sigma\times\{0\}
 \sqcup
 (\Sigma\times\Sigma)\times(0,1]
 }
\]
embeds as a closed Lie subgroupoid of the tangent groupoid $\mathcal G_K$.
For $t>0$, convolution is ordinary boundary-kernel composition,
\[
 (f*g)(x,z)=\int_\Sigma f(x,y)g(y,z)\,\dd y.
\]
At $t=0$, the product is fibrewise addition in $T\Sigma$ and the Fourier
picture is pointwise multiplication on $T^*\Sigma$.
\end{theorem}

\begin{proof}
The pair-groupoid statement follows by restricting source, target, units,
inverse and multiplication to $\Sigma$.  Because $\Sigma$ is totally geodesic,
the exponential map of $K$ restricts to the exponential map of $\Sigma$ on a
neighbourhood of the zero section of $T\Sigma$.  The standard exponential
charts defining the tangent-groupoid smooth structure therefore restrict to
charts for $\mathcal G_\Sigma$, proving the subgroupoid statement.  The
convolution and zero-fibre descriptions are the restrictions of the ambient
ones.
\end{proof}

\begin{definition}[Positive return coordinate and Gaussian boundary kernel]
\label{p1b:def:v124-gaussian-return-kernel}
Let $G>0$ and $0\le\rho<1$.  Define
\[
 \boxed{
 M_G(\rho)
 :=\frac{G}{1-\rho}
 \begin{pmatrix}
 1+\rho&-2\sqrt\rho\\
 -2\sqrt\rho&1+\rho
 \end{pmatrix}.}
\]
For endpoint values $f_0,f_1\in\mathbb R$, set
\[
 S_{G,\rho}(f_0,f_1)
 :=\frac12
 \begin{pmatrix}f_0&f_1\end{pmatrix}
 M_G(\rho)
 \binom{f_0}{f_1}.
\]
\end{definition}

\begin{theorem}[Exact return-coordinate representation of the collar action]
\label{p1b:thm:v124-return-coordinate-collar}
Let $L>0$ and let $u$ solve
\[
 -u''+G^2u=0
\]
on $[0,L]$ with endpoint values $u(0)=f_0$, $u(L)=f_1$.  The on-shell
quadratic action
\[
 \frac12\int_0^L\bigl((u')^2+G^2u^2\bigr)\,\dd s
\]
is $S_{G,\rho}(f_0,f_1)$ with
\[
 \boxed{\rho=e^{-2GL}.}
\]
Equivalently,
\[
 M_G(\rho)
 =G
 \begin{pmatrix}
 \coth(GL)&-\operatorname{csch}(GL)\\
 -\operatorname{csch}(GL)&\coth(GL)
 \end{pmatrix}.
\]
The matrix is positive definite for $0\le\rho<1$.
\end{theorem}

\begin{proof}
The unique solution is a linear combination of $e^{Gs}$ and $e^{-Gs}$.
Integrating by parts on shell reduces the quadratic action to the endpoint
flux pairing.  Solving for the coefficients yields the displayed hyperbolic
matrix.  Using
\[
 \coth(GL)=\frac{1+e^{-2GL}}{1-e^{-2GL}},\qquad
 \operatorname{csch}(GL)=\frac{2e^{-GL}}{1-e^{-2GL}}
\]
gives $M_G(\rho)$.  Its eigenvalues are
\[
 G\frac{1+\sqrt\rho}{1-\sqrt\rho},\qquad
 G\frac{1-\sqrt\rho}{1+\sqrt\rho},
\]
which are positive.
\end{proof}

\begin{theorem}[Exact stationary gluing in the multiplicative return coordinate]
\label{p1b:thm:v124-return-gluing}
Let $\star$ denote stationary Schur reduction of the common endpoint in two
quadratic boundary actions.  Then
\[
 \boxed{
 M_G(\rho_1)\star M_G(\rho_2)
 =M_G(\rho_1\rho_2).}
\]
Thus $\rho$ is the multiplicative composition coordinate of the exact
Gaussian boundary process.
\end{theorem}

\begin{proof}
For $0<\rho_i<1$, write $\rho_i=e^{-2GL_i}$.  The two actions are the on-shell actions on
adjacent intervals of lengths $L_1,L_2$.  Stationary elimination of their
common boundary value is equivalent to solving the same Euler--Lagrange
equation on the glued interval of length $L_1+L_2$.  Therefore the resulting
kernel is $M_G(e^{-2G(L_1+L_2)})=M_G(\rho_1\rho_2)$.  Equivalently, direct
Schur complementation of the common endpoint gives the same matrix identity.
The cases with $\rho_i=0$ follow by continuity from the displayed formula.
\end{proof}

\begin{corollary}[Capped Dirichlet-to-Neumann response]
\label{p1b:cor:v124-capped-dtn-cayley}
Stationarily eliminating the far endpoint gives the one-boundary Hessian
\[
 \boxed{
 K_G(\rho)
 =G\frac{1-\rho}{1+\rho}.}
\]
For $\rho=e^{-2GL}$ this is
\[
 K_G=G\tanh(GL).
\]
Thus the positive return contraction and the boundary impedance are related by
the Cayley transform.
\end{corollary}
\begin{proof}
The far diagonal entry of \(M_G\) is positive, so its Schur
complement exists.  Subtract its off-diagonal square divided by
that diagonal entry:
\[
 \frac{G(1+\rho)}{1-\rho}
 -\frac{4G\rho}{(1-\rho)(1+\rho)}
 =G\frac{1-\rho}{1+\rho}.
\]
Substitution of \(\rho=e^{-2GL}\) is the hyperbolic identity stated.
\end{proof}


\begin{corollary}[Riccati form]
\label{p1b:cor:v124-boundary-riccati}
With $t=-\tfrac12\log\rho$, the normalized response
\[
 \eta(t):=K_G/G=\tanh t
\]
obeys
\[
 \boxed{\frac{\dd\eta}{\dd t}=1-\eta^2,\qquad \eta(0)=0.}
\]
Equivalently, the multiplicative return law $\rho(t_1+t_2)=\rho(t_1)\rho(t_2)$
is Cayley-conjugate to the hyperbolic addition law
\[
 \eta(t_1+t_2)
 =\frac{\eta(t_1)+\eta(t_2)}{1+\eta(t_1)\eta(t_2)}.
\]
\end{corollary}
\begin{proof}
The Cayley expression is
\((1-e^{-2t})/(1+e^{-2t})=\tanh t\).
Differentiate to obtain \(1-\eta^2\); its continuous value at
zero is zero.  Substitute \(e^{-2(t_1+t_2)}=e^{-2t_1}e^{-2t_2}\)
into the same Cayley expression and simplify to obtain the
addition law, on the stated nonnegative parameter domain.
\end{proof}


\begin{definition}[Mathematical boundary-gluing constructor]
\label{p1b:def:v124-boundary-gluing-constructor}
For a supplied positive spectral root $G$ and supplied positive multiplicative
return coordinate $\rho$, define
\[
 \mathfrak B(G,\rho):=M_G(\rho).
\]
The preceding theorems make $\mathfrak B$ a functorial Gaussian boundary
constructor with respect to stationary gluing.
\end{definition}

\begin{remark}[Distinct boundary semigroups]
\label{p1b:rem:v124-boundary-gluing-owner-firewall}
The pair/tangent-groupoid restriction and the Gaussian gluing law use the supplied $G$, $\rho$, domain and endpoint data. A determinant-depth scalar semigroup is identified with a Poisson semigroup only by an explicit equality of generators and domains; agreement of a scalar exponent alone does not identify the operators.
\end{remark}


\subsection{Exact finite calculations and sharpness witnesses}
\label{p1b:app:finite-checks}
In unit quaternion coordinates, the finite sets used in
Theorem~\ref{p1b:thm:design} are
\[
\begin{split}
 2T={}&\{(\pm1,0,0,0)\text{ and coordinate permutations}\}
       \cup\{\tfrac12(\pm1,\pm1,\pm1,\pm1)\},\\
 2O={}&2T\cup
       \{\text{permutations of }(\pm1/\sqrt2,\pm1/\sqrt2,0,0)\},\\
 2I={}&2T\cup
       \{\text{even permutations of }(0,\pm1/2,\pm\varphi/2,
                                      \pm1/(2\varphi))\},
 \qquad \varphi=(1+\sqrt5)/2 .
\end{split}
\]
Signs in each displayed set vary independently. Their cardinalities
are \(24,48,120\). Coordinate sign symmetries make every monomial
with an odd coordinate exponent integrate to zero. The even
monomial formula in the main text therefore reduces exact design
verification to finitely many calculations in quadratic number fields.

The following table gives a real field \(f\) in the first excluded
complete band, the order \(r\), and the distinct means of \(f^r\).
The two \(2T\) rows with \(r=3,4\) use \(f=4x_0^2-1\).
Every failure persists in all larger complete bands.
\[
\begin{array}{c|c|c|c|c}
\Gamma&r&f&|\Gamma|^{-1}\sum f^r&\int_{S^3}f^r\,\dd\sigma\\\hline
2T&2&x_0^3&3/32&5/64\\
2T&3&4x_0^2-1&2&1\\
2T&4&4x_0^2-1&7&3\\
2O&2&x_0^4&15/256&7/128\\
2O&3&x_0^3+x_0^2&65/256&31/128\\
2O&4&4x_0^2-1&4&3\\
2I&2&x_0^6&133/4096&132/4096\\
2I&3&x_0^4&133/4096&132/4096\\
2I&4&x_0^3&133/4096&132/4096
\end{array}
\]
These values follow by expansion into even monomials and direct
finite summation; the supplied exact-arithmetic script replays all
entries and all moments through the claimed design strengths.
It also checks the torus integer matrix, intersection signs,
exchange ranks, KO groups, selected Clifford identities, Hilbert
series, invariant-pencil relation, and Chern-index arithmetic.
Those checks support the extraction and its normalizations; they
do not replace the stated proofs or certify arbitrary extra data.

\begin{proposition}[Sign of the homogeneous heat curvature]
\label{p1b:prop:heat-curvature}
For the heat trace \(Z(h,\delta)\) in
Section~\ref{p1b:app:dirac}, \(Z_\delta(h,0)=0\), and
\[
 Z_{\delta\delta}(h,0)<
 \begin{cases}
 -2\sqrt{\pi h},&0<h\le1,\\
 -(462/47)h,&h\ge1 .
 \end{cases}
\]
\end{proposition}
\begin{proof}
The exact trace is even in \(\delta\).  Differentiating its
absolutely convergent sums gives, with \(b_n=n(n+2),c_n=(n+1)^2\),
\[
 Z_{\delta\delta}/h=-18+4\sum_{n\ge1}
  \{c_n(2hb_n-5)e^{-hb_n}+b_n(2hc_n-1)e^{-hc_n}\}.
\]
For \(h\ge1\), discard the negative terms inside braces and
use \(he^{-h\lambda}\le e^{-\lambda}\) for \(\lambda\ge3\).
Writing \(v=e^{-1}<3/8\) gives the upper bound
\[
 Z_{\delta\delta}/h\le-18+
 8\sum_{m\ge2}m^2(m^2-1)(v^{m^2-1}+v^{m^2}).
\]
The \(m=2\) term is at most \(96[(3/8)^3+(3/8)^4]\).
For \(m\ge3\), consecutive terms have ratio at most
\((4/3)^4(3/8)^7<1/300\); bounding their polynomial coefficient
by \(m^4\), their total is at most
\(8\cdot81[(3/8)^8+(3/8)^9]/(1-1/300)\).
The sum of these two rational bounds is less than \(384/47\).
This proves the large-\(h\) constant.

For \(0<h\le1\) put \(F(h)=\sum_{m\in\mathbb Z}e^{-hm^2}\).
Regrouping the same differentiated series gives exactly
\[
 Z_{\delta\delta}=2hF+4h^2(e^h+1)F''
       +2h[e^h(5+2h)+1+2h]F'.
\]
Poisson summation is
\(F=\sqrt{\pi/h}\sum_{l\in\mathbb Z}e^{-\pi^2l^2/h}\).
For its \(l=0\) term the expression is
\(2\sqrt\pi h^{-1/2}[1-(1+h)e^h]\le-4\sqrt{\pi h}\).
For a nonzero term put \(t=\pi^2/h\).  After factoring
\(2\sqrt\pi h^{-1/2}e^{-tl^2}\), its bracket is
\[
 1-(1+h)e^h+tl^2[e^h(2h-1)+2h-5]
       +2t^2l^4(e^h+1).
\]
The first two terms are negative on \(0<h\le1\): the square
bracket is increasing and its value at one is \(e-3<0\).
The sum of the remaining positive terms is bounded by
\(4(e^h+1)t^2\sum_{l\ge1}l^4e^{-tl^2}\).
Since \(t\ge\pi^2>9\),
\[
 4(e^h+1)t^3\sum_{l\ge1}l^4e^{-tl^2}
 \le\frac{16\cdot9^3(3/8)^9}{1-16(3/8)^{27}}<9<\pi^2.
\]
Here consecutive summands have ratio at most \(16e^{-3t}\),
and \(t^3e^{-t}\) decreases for \(t\ge9\).
Thus that bracket sum is strictly less than \(h\).
The remainder in \(Z_{\delta\delta}\) is less than
\(2\sqrt{\pi h}\), proving the small-\(h\) bound.
\end{proof}


At this stationary point the different function \(-\log Z\)
has positive second derivative. The two functions should be
kept distinct when a spectral functional is specified.

\subsection{All-level bundle sampling and positive quadrature}
\label{p1b:app:all-level-sampling}
This technical section supplies the explicit realization and longer quadrature
argument for Theorem~\ref{p1b:thm:all-level-sampling}.  The homogeneous
bundle model defined here is a uniform mathematical model for all $H_m$.
At $m=1$ the earlier form/spin model remains available with its own
geometric symbol.  Equivariant isometry of the two representation models
does not identify untransported pointwise operations or connections.

\begin{proof}[Proof and explicit model for Theorem~\ref{p1b:thm:all-level-sampling}]
Write \(G=SU(2)\times SU(2)\), \(H=SU(2)_{\rm diag}\),
and \(U_{a,b}=V_{a/2}\otimes V_{b/2}\).
For \(a+b\le2m\), the diagonal Clebsch--Gordan decomposition
contains its highest spin \(s=(a+b)/2\) once.  Let
\(T_{a,b}:U_{a,b}\to V_s\) be that orthogonal coisometry and
\(d_{a,b}=(a+1)(b+1)\).  The equivariant function
\[
 F_v(g)=\sqrt{\frac{d_{a,b}}{a+b+1}}\,
                       T_{a,b}\pi_{a,b}(g^{-1})v
\]
satisfies \(F_v(gh)=\rho_s(h^{-1})F_v(g)\), so defines a
section of the stated homogeneous bundle.  Haar averaging and
Schur's lemma give its Gram as the identity: the averaged
\(\pi(g)T^*T\pi(g)^{-1}\) is
\((a+b+1)I/d_{a,b}\).
Distinct irreducibles are orthogonal, including those placed in
the same fiber-spin summand.  Taking their direct sum therefore
gives an isometric realization of every summand of \(H_m\).
In the section \((q,1)\), its entries are matrix coefficients of
\(\pi_{a/2}(q^{-1})\), hence polynomials of degree \(a\le2m\)
in the quaternion coordinates.  The fiber rank is
\(\sum_{j=0}^{2m}(j+1)=\binom{2m+2}{2}=v_m\).
Products for the pairing have degree at most \(4m\);
products of \(r\) fields with the stated coefficient have degree
at most \(2rm+d\).  These facts prove each design bound,
and taking independent product rules proves their factorwise forms.

For the general quadrature assertion normalize its supplied positive
measure and choose a real basis \(1,f_2,\ldots,f_D\) of the test
space.  The vector of integrals of \((f_2,\ldots,f_D)\) lies
in the convex hull of its compact image.  Caratheodory's theorem
in \(\mathbb R^{D-1}\) writes it as a convex combination of at
most \(D\) evaluation vectors.  Deleting any zero weights gives
a positive finite quadrature; restore the total mass if necessary.
Including all real and imaginary pairing coefficients makes its
weighted evaluation Gram the original Gram, hence an isometry.
The already proved isometric image-transport identities then apply
without identifying any untransported pointwise operator.
\end{proof}






\section{Supplementary calculations for the product atlas}
\label{p1c:sec:atlas}

\hypertarget{p1b:app:g02}{%
\subsection{Coordinate ancestry and retained geometric squares: supplementary calculations}\label{p1b:app:g02}}

\hypertarget{p1b:app:g02:1}{%
\subsubsection{First-order sign convention}\label{p1b:app:g02:1}}

The adjoint Leibniz rule gives

\[
[d,M_f]=\varepsilon(df),\qquad
[d^\dagger,M_f]=-\iota_{\nabla f}.
\]

It follows that

\[
i[+i(d-d^\dagger),M_f]=-(\varepsilon(df)+\iota_{\nabla f}),
\]

while the exterior convention \(-i(d-d^\dagger)\) gives the positive sign. Paper I writes \(+i(d-d^\dagger)\) alongside the positive Clifford-symbol commutator convention. The literal displayed equality has this sign mismatch on the exterior summand.

The unitary exterior grading \((-1)^{\deg}\) anticommutes with \(d-d^\dagger\) and with \(\varepsilon+\iota\), and commutes with scalar multiplication and the retained projection. Its restriction is precisely a permitted sector phase change: \(+1\) on degrees zero and two, \(-1\) on degree one. Extend it by the identity on the spinor summand. It conjugates the two exterior conventions while leaving their squares unchanged. Therefore the coordinate phase classes, both fingerprints, and every LC-square statement proved here are unaffected. During later integration, use one consistent first-order convention rather than reproduce the literal sign mismatch.

\hypertarget{p1b:app:g02:2}{%
\subsubsection{Squaring and compression are different operations}\label{p1b:app:g02:2}}

Let \(P=P_R\), \(Q=I-P\), and \(\mathscr D\) be the self-adjoint LC representative. All retained modes are smooth, so the following finite compression is well defined:

\[
P\mathscr D^2P-(P\mathscr DP)^2
=P\mathscr DQ\mathscr DP
=(Q\mathscr DP)^*(Q\mathscr DP)\ge0.
\]

The defect can be determined exactly here. The derivative of the constant is zero. The retained exact one-forms map under \(d^\dagger\) to first scalar harmonics, which are outside the retained scalar line. The retained two-forms are closed; their \(d^\dagger\) images are Killing/coexact one-forms, orthogonal to the retained exact one-forms. The two retained spinor spaces are invariant under the Dirac operator itself. Hence

\[
R^2(J_R^*\mathscr DJ_R)^2
=(0,\;0,\;0,\;0,\;9/4,\;9/4),
\] \[
\boxed{R^2\left(K_R-(J_R^*\mathscr DJ_R)^2\right)
=(0,\;3,\;4,\;4,\;0,\;0).}
\]

The defect has rank ten. The compressed first-order square has an eleven-dimensional kernel; the retained geometric square has a one-dimensional kernel. This proves why the latter is required for the scalar projector.

On a tensor product there is an additional first-order issue: an ungraded sum \(\mathscr D_x\otimes I+I\otimes\mathscr D_p\) has square containing \(2\mathscr D_x\otimes\mathscr D_p\). A genuine product Dirac construction requires its declared Clifford/graded bundle conventions. The second-order tensor sum in Section~\ref{p1b:der:g02:s5} is defined directly and has the displayed additive spectrum without making that extra identification.

\hypertarget{p1b:app:g02:3}{%
\subsubsection{Quadratic energy and summed Hopf charge}\label{p1b:app:g02:3}}

Choose the \(q_x=q_p=2\) extreme right-bivector state. It has \(A_x=A_p=4\), so

\[
\sqrt{A_x+A_p}=2\sqrt2<4=|N_x+N_p|.
\]

Using opposite extreme weights gives the same witness for \(|N_x-N_p|\). Thus an asserted inequality \(\sqrt{A_x+A_p}\ge|N_x\pm N_p|\) would be false on this very carrier. The valid additive-energy and vector-norm comparisons are the ones stated in Section~\ref{p1b:der:g02:s6}.

\subsubsection{Polynomial form of the retained one-factor square}
\label{p1b:app:g02:casimir-polynomial}
This is a supplementary expression for the relation already proved
in Section~\ref{p1b:der:g02:s5:sub2}. Use that section's retained
one-factor carrier, representation Casimir \(\mathcal C\),
spin projector \(\Pi_{\rm sp}\), and LC square \(A_R=R^2K_R\).
Then
\[
 \Pi_{\rm sp}
 =\frac{64}{45}\mathcal C
   (\mathcal C-\tfrac32I)(\mathcal C-2I),
\]
and hence
\[
 A_R=\frac2{15}\mathcal C
     (8\mathcal C^2-28\mathcal C+39I).
\]
To prove the first formula, evaluate its polynomial on
\(\operatorname{spec}(\mathcal C)=\{0,3/4,3/2,2\}\):
it equals one at \(3/4\) and zero at the other three values.
Substitution into \(A_R=2\mathcal C+3\Pi_{\rm sp}/4\)
proves the second formula. The equality is restricted to this
retained spectrum. On a product it applies separately to
\(\mathcal C_x\) and \(\mathcal C_p\); it does not make
\(A_x+A_p\) a function of the summed Casimir alone.



\hypertarget{p1b:app:g04}{%
\subsection{Hurwitz, trimer, Walsh and resolution structures: supplementary calculations}\label{p1b:app:g04}}

\hypertarget{p1b:app:g04:1}{%
\subsubsection{Product adjacency kernels}\label{p1b:app:g04:1}}

This comparison explains why a single product adjacency does not retain all 25 factorwise shell labels. It adds no new geometric assumptions. If the factor eigenvalue-rank pairs are \((\lambda_i,r_i)\), then \[
A_{\rm sum}=A_+\otimes I+I\otimes A_+,\qquad
A_{\rm ten}=A_+\otimes A_+
\] have eigenvalues \(\lambda_i+\lambda_j\) and \(\lambda_i\lambda_j\), respectively, on \(P_i\otimes P_j\). For the sum, the zero pairs are \((0,0),(4,-4),(-4,4)\). Hence \[
\dim\ker A_{\rm sum}=9^2+2(4)(2)=97.
\] For the tensor adjacency, at least one factor eigenvalue must be zero, so \[
\dim\ker A_{\rm ten}=24^2-15^2=351.
\] The shell product \(P_9\otimes P_9\) has rank 81 and is a specified proper subprojection in each case. The exact full additive spectrum is included in the accompanying derivation and verification record's output.

\hypertarget{p1b:app:g04:2}{%
\subsubsection{Sharp natural-action comparison and the tensor-rank distinction}\label{p1b:app:g04:2}}

This supplementary argument establishes the precise scope of the scalar bridge without demanding that different 15-dimensional constructions carry the same representation.

Under \(SU(2)_L\) restricted to \(2T\), write \(\rho\) for its defining two-dimensional representation and \(\vartheta\) for the three-dimensional tetrahedral rotation representation. The source's sector decomposition gives \[
H_1|_{2T}=6\mathbf1\oplus3\rho\oplus\vartheta.
\] For a unit quaternion \(q=(q_0,\mathbf q)\), their characters are \[
\chi_{\mathbf1}=1,\qquad \chi_\rho=2q_0,\qquad
\chi_\vartheta=4q_0^2-1.
\] Summing products over the explicit 24 nodes gives \[
\frac1{24}\sum_{q\in2T}
\overline{\chi_i(q)}\chi_j(q)=\delta_{ij}.
\] Thus these are three pairwise inequivalent irreducible representations. The regular representation contains them with multiplicities \(1,2,3\).

The neighbor set of the identity is the conjugation-invariant set of eight Hurwitz quaternions with \(q_0=1/2\). The adjacency is convolution by this central set. On \(\vartheta\), its scalar is \[
\frac1{\dim\vartheta}\sum_{\substack{q\in2T\\q_0=1/2}}
\chi_\vartheta(q)=0.
\] The full \(\vartheta\)-isotypic part of the regular representation has dimension \(3\cdot3=9\), equal to \(\dim\ker A_+\). Therefore this isotypic part is exactly \(\mathscr E_9\); \(\mathscr E_{15}\) has no \(\vartheta\) summand. The explicit character projector \[
\frac3{24}\sum_{h\in2T}\chi_\vartheta(h^{-1})L_h=P_9
\] is also verified from quaternion multiplication in the accompanying derivation and verification record.

For semisimple representations \(V=\bigoplus_i V_i\otimes\mathbb C^{m_i}\) and \(U=\bigoplus_i V_i\otimes\mathbb C^{n_i}\), every intertwiner has blocks \(I_{V_i}\otimes f_i\). Its maximum possible rank is \[
\sum_i(\dim V_i)\min(m_i,n_i);
\] choose each multiplicity-space map \(f_i\) of maximal rank to attain the bound. Applying this to the three source types gives \[
\max\operatorname{rank}(H_1\to\mathbb C^{2T})
=1\min(6,1)+2\min(3,2)+3\min(1,3)=8,
\] \[
\max\operatorname{rank}(H_1\to\mathscr E_{15})
=1\min(6,1)+2\min(3,2)=5.
\] The source's non-injection result is thus sharpened to kernel bounds seven and ten for these two respective targets. The scalar five-dimensional evaluation image attains the latter bound.

For the independent product group, the multiplicity of \(V_i\otimes V_j\) is \(m_im_j\) in the source and \(n_in_j\) in a product target. The maximum rank of a \textbf{general} product-group intertwiner is therefore \[
\sum_{i,j}(\dim V_i)(\dim V_j)\min(m_im_j,n_in_j).
\] For the full regular product, with \(m=(6,3,1)\) and \(n=(1,2,3)\), the nine contributions are \[
\begin{pmatrix}
1&4&9\\
4&16&18\\
9&18&9
\end{pmatrix},
\qquad\text{whose sum is }88.
\] For the coarse product targets, the maximum ranks are \[
\begin{array}{c|cc}
&\mathscr E_9&\mathscr E_{15}\\ \hline
\mathscr E_9&9&27\\
\mathscr E_{15}&27&25
\end{array}.
\] Their sum is also 88; their relevant irreducible target types partition the regular product. In particular, every natural product-group map from the 225-dimensional \(H_1\otimes H_1\) to the equally 225-dimensional \(\mathscr E_{15}\otimes\mathscr E_{15}\) has rank at most 25 and kernel dimension at least 200.

A simple tensor \(F_x\otimes F_p:H_1^{\otimes2}\to (\mathbb C^{2T})^{\otimes2}\) has maximum rank \(8^2=64\). Linear combinations of tensor intertwiners can attain the larger rank 88. Multiplying one-factor maximum ranks is therefore insufficient for the general intertwiner problem. This distinction does not alter the scalar bridge, which is a specific tensor isometry of rank 25.

\hypertarget{p1b:app:g04:3}{%
\subsubsection{Graded resolution exchange and the protected Clifford transport}\label{p1b:app:g04:3}}

This refinement uses the native \(\mathbb Z_2\)-grading of \(\mathrm{Cl}_{0,3}\). It is separate from the ordinary product convention used for the principal resolution formulas.

For \(A=A^0\oplus A^1\), the dimensions are \(4,4\). For \(W=A/\mathbb R1\), the dimensions are \(3,4\). Graded exchange on \(A\widehat\otimes A\) has trace zero. Its induced exchange on \(W_\times\) has trace \(-1\), because the removed scalar line was fixed. The conjugation action on \(\operatorname{Hom}_{\mathbb R}(W_\times,A_\times)\) thus has trace \(0\cdot(-1)=0\), giving exchange ranks \(2016,2016\). On the mixed domain \(W\widehat\otimes W\), graded exchange also has trace \(3-4=-1\). The mixed Hom space therefore has exchange trace zero and ranks \(1568,1568\). These are computed directly as conjugation actions on Hom spaces; an identification with a graded tensor product of Hom spaces must include its Koszul signs.

\begin{longtable}[]{@{}
  >{\raggedright\arraybackslash}p{(\columnwidth - 4\tabcolsep) * \real{0.3333}}
  >{\raggedright\arraybackslash}p{(\columnwidth - 4\tabcolsep) * \real{0.3333}}
  >{\raggedright\arraybackslash}p{(\columnwidth - 4\tabcolsep) * \real{0.3333}}@{}}
\toprule\noalign{}
\begin{minipage}[b]{\linewidth}\raggedright
Resolution domain
\end{minipage} & \begin{minipage}[b]{\linewidth}\raggedright
Ordinary exchange \(+\), \(-\)
\end{minipage} & \begin{minipage}[b]{\linewidth}\raggedright
Graded exchange \(+\), \(-\)
\end{minipage} \\
\midrule\noalign{}
\endhead
\bottomrule\noalign{}
\endlastfoot
Mixed, dimension 3136 & 1596, 1540 & 1568, 1568 \\
Full, dimension 4032 & 2044, 1988 & 2016, 2016 \\
\end{longtable}

The protected one-factor Clifford comparison is also explicit. For four positive Clifford generators \(\gamma_\alpha\gamma_\beta+\gamma_\beta\gamma_\alpha =2\delta_{\alpha\beta}\), put \[
E_i=\gamma_i\gamma_4,\qquad i=1,2,3.
\] Then \(E_i^2=-1\) and \(E_iE_j=-E_jE_i\). Their monomials span the eight-dimensional real even algebra: the unit, six bivectors, and the volume element. Thus \[
\mathrm{Cl}_{0,3}\cong\mathrm{Cl}_{4,0}^{\,0}
\] as ungraded real algebras, relative to the chosen frame and real form. The generator \(E_i\) has native \(\mathrm{Cl}_{0,3}\) odd parity but ambient \(\mathrm{Cl}_{4,0}\) even parity. This is why the preceding graded product and exchange are not obtained by silently reusing the ambient protected spinor grading of G03.

For two protected sphere factors, the ordinary tensor algebra maps to the commuting factor-even subalgebra \[
\mathrm{Cl}_{4,0}^{\,0}\otimes\mathrm{Cl}_{4,0}^{\,0}
\subset\mathrm{Cl}_{8,0}^{\,0}.
\] Its real dimension is 64, whereas the full even algebra has dimension 128. After complexification its image on the protected 16-dimensional complex spinor product is the four-block \(M_4(\mathbb C)^{\oplus4}\) algebra in G03. The real and complex dimension statements refer to the respective scalar fields.

\hypertarget{p1b:app:g04:4}{%
\subsubsection{Flags, coordinate changes, and the meaning of the carrier counts}\label{p1b:app:g04:4}}

This is a scope refinement for optional ordered submodules. The bare character-label symmetry \(GL(3,2)\) acts transitively on the seven nonzero characters: extend either nonzero vector to a basis and map one basis to the other. Hence there is no invariant proper nonempty subset of those seven labels. A particular ordered seven-step flag requires more data.

For the factorized product label action \((GL(3,2)\times GL(3,2))\rtimes\mathbb Z_2\), the nonconstant labels have two orbits: the 14 marginal labels and the 49 mixed labels. Without exchange the two marginal seven-element sets are separate orbits. The full \(GL(6,2)\) label action is transitive on all 63 nonzero labels, but it does not preserve the chosen factorization. These are statements about label sets. A label transformation is not automatically an automorphism of the fixed Clifford-twisted multiplication; that requires preserving or correctly transporting the chosen cocycle.

The finite-support module \(D=\operatorname{Hom}_{\mathbb R}(W,A)\) and the real three-form space \(\Lambda^3T^*M_8\) both have dimension 56, since \(\binom83=56\). They have different definitions and natural actions. The dimension equality alone supplies no identification with the alternating-torsion parameter space already discussed in Paper 1B.

Likewise, the product Walsh carriers of dimension 64 concern the two three-dimensional sphere tangent packages. If one explicitly chooses the full real Clifford algebra of an eight-dimensional tangent space, its dimension is \(2^8=256\); its scalar-free quotient has dimension 255 and its analogous full Hom module has dimension \(256\cdot255=65280\). This is a different construction, not a change to the previously defined 56-, 3136-, or 4032-dimensional modules.



\hypertarget{p1b:app:g05}{%
\subsection{Contact pairs and product control geometry: supplementary calculations}\label{p1b:app:g05}}

\hypertarget{p1b:app:g05:1}{%
\subsubsection{Conjugate Jacobian and what singularity does not imply}\label{p1b:app:g05:1}}

This calculation supplies the coefficient check for Section~\ref{p1b:der:g05:s33}. Write \(z_2=\rho e^{i\psi}\), \(z_1=\sqrt{1-\rho^2}e^{i\xi}\). In this Hopf chart the unit-round volume density is \(\rho\,d\rho\,d\xi\,d\psi\). From the normal curve, \[
\rho=\sqrt{1-\beta^2}\sin s,\qquad
\xi=-\beta s+\arg(\cos s+i\beta\sin s),\qquad
\psi=\phi-\beta s.
\] Put \(D=\cos^2s+\beta^2\sin^2s\). On a regular coordinate chart, \[
\partial_s\xi=-\beta+\frac\beta D,\qquad
\partial_\beta\xi=-s+\frac{\sin s\cos s}{D}.
\] Because \(\partial_\phi\psi=1\), the endpoint-volume Jacobian reduces, up to orientation, to \[
\rho\left(
\partial_s\rho\,\partial_\beta\xi
-\partial_\beta\rho\,\partial_s\xi
\right)
=(1-\beta^2)\sin s(\sin s-s\cos s).
\] The final expression is smooth across chart boundaries and is also the quaternion volume determinant there by continuity. For \(v>0\), the change from fixed-time covector parameters to \((s,\beta,\phi)\) is nonsingular: \[
\det\frac{\partial(s,\beta)}{\partial(v,c)}
=\frac{tv}{R(v^2+c^2)}\ne0.
\] Thus the displayed zero set gives the regular conjugate times.

The absence of multiple minimizing paths does not by itself imply smooth squared distance at endpoints with a stationary sphere factor. Already on a sphere, for a signed small fiber angle \(\theta\), \[
d_{CC,R}(1,e^{i\theta})^2
=R^2(2\pi|\theta|-\theta^2).
\] This has a cusp at zero in the smooth vertical coordinate. The product squared-distance formula inherits this cusp at a stationary sphere component, even if the other components have unique minimizing paths. Likewise, a stationary sphere admits different normal covector lifts of exactly the same constant curve. Neither phenomenon creates a second minimizing path. This explains the path-based cut-set definition used in Section~\ref{p1b:der:g05:s32}.

The infimum of cut times over unit-speed sphere geodesics is zero, since \(\pi R/\sqrt{1+c^2}\to0\) as \(|c|\to\infty\). Thus no positive Riemannian-style uniform injectivity radius should be inferred for this control exponential.

On \(M_8\), the two pulled-back contact forms alone also allow a scalar tangent field to be added to either Reeb lift without altering their contraction equations. The source's product splitting selects the sphere-factor lifts used throughout G05. The contact-pair uniqueness statement in Section~\ref{p1b:der:g05:s29} is on \(P_v\).

\hypertarget{p1b:app:g05:2}{%
\subsubsection{Intrinsic exchange-fixed distance can exceed ambient distance}\label{p1b:app:g05:2}}

This example records a global distinction that is invisible in the fixed tangent metric. Take the exchange-invariant torus metric \[
G=\begin{pmatrix}2&1\\1&2\end{pmatrix}.
\] Both \((0,0)\) and \((\pi,\pi)\) lie on its exchange-fixed circle. A shortest lift constrained to that circle has displacement \(\pm(\pi,\pi)\), of squared length \(6\pi^2\). An ambient shortest lift has displacement \((\pi,-\pi)\) or its negative, with squared length \(2\pi^2\). To verify minimality, all allowed components are odd multiples of \(\pi\), and \(G\) has smallest eigenvalue one; their squared Euclidean norm is at least \(2\pi^2\), attained by these opposite-sign vectors. Thus the intrinsic circle distance is strictly larger. This does not affect the locally induced control structure or its intrinsic Hausdorff dimension in Section~\ref{p1b:der:g05:s36}.

\hypertarget{p1b:app:g05:3}{%
\subsubsection{The control distance as a large-vertical-cost limit}\label{p1b:app:g05:3}}

This supplementary limit connects G05 to the Berger metric family to be treated in G07. For \(\varepsilon>0\), set \[
g_\varepsilon=
g_{\mathcal D}
+\varepsilon^{-2}(\alpha_x^2+\alpha_p^2)
\] on the orthogonal splitting \(TM_8=\mathcal D\oplus\operatorname{span}(Z_x,Z_p)\). It is a smooth Riemannian metric. Horizontal paths have unchanged length, so \[
d_{g_\varepsilon}\le d_{\mathcal D},
\] and these Riemannian distances increase as \(\varepsilon\downarrow0\). They converge to \(d_{\mathcal D}\), uniformly on \(M_8\times M_8\).

For the lower bound in the limit, take a sequence of minimizing \(g_\varepsilon\)-paths, parametrized with constant speed on \([0,1]\). Their energies are bounded by \(d_{\mathcal D}^2\). For \(\varepsilon\le1\) this also bounds their energy in a fixed background product metric. Compactness and weak compactness of the controls give a uniformly convergent subsequence with a weakly \(L^2\)-convergent velocity. Moreover, \[
\int_0^1\bigl(|\alpha_x(\dot\gamma_\varepsilon)|^2+
|\alpha_p(\dot\gamma_\varepsilon)|^2\bigr)\,dt
\le C\varepsilon^2\longrightarrow0 .
\] Passing to the limit through the smooth frame shows the limiting curve is horizontal. Lower semicontinuity of horizontal energy therefore gives \(d_{\mathcal D}^2\le\liminf d_{g_\varepsilon}^2\), proving equality. The monotone convergence of continuous distance functions to the continuous control distance is uniform on the compact endpoint space by Dini's theorem, first along a decreasing sequence and then for all \(\varepsilon\) by monotonicity.

\hypertarget{p1b:app:g05:4}{%
\subsubsection{One-form constraint pencil and its horizontal scalar operator}\label{p1b:app:g05:4}}

This optional family uses the same two marked forms but imposes only one specified linear constraint. It is a different control structure from the simultaneous two-constraint geometry in the main derivation. For real constants \((a,b)\ne(0,0)\), put \[
\mathcal D_{a,b}=\ker(a\alpha_x+b\alpha_p).
\] Its orthogonal decomposition is \[
\mathcal D_{a,b}
=\mathcal D\oplus
\mathbb R\frac{bZ_x-aZ_p}{\sqrt{a^2+b^2}}.
\] Hence its growth on \(M_8\) is \((7,8)\) and its homogeneous dimension is nine. The remaining normal is \((aZ_x+bZ_p)/\sqrt{a^2+b^2}\); the two horizontal bracket coefficients in this normal direction are \[
-\frac{2a}{R_x\sqrt{a^2+b^2}},\qquad
-\frac{2b}{R_p\sqrt{a^2+b^2}}.
\] At least one is nonzero, so step-two bracket generation follows.

If \(ab\ne0\), the graded algebra is \(\mathfrak h_2\oplus\mathbb R^3\): four horizontal directions generate one degree-two center, and the remaining three horizontal directions lie in the Levi radical. Its two nonzero Levi frequencies are the absolute values of the displayed coefficients, which need not agree. Thus the abstract Heisenberg algebra type alone does not identify its horizontal metric with the isotropic model. If exactly one coefficient is zero, the graded algebra instead is \(\mathfrak h_1\oplus\mathbb R^5\). On \(P_v\), remove the scalar directions: the growth is \((5,6)\), the homogeneous dimension is seven, and the corresponding symbols are \(\mathfrak h_2\oplus\mathbb R\) or \(\mathfrak h_1\oplus\mathbb R^3\). For \(ab\ne0\) this is the usual even-contact Pfaffian structure on the six-dimensional product; its characteristic line is spanned by \(bZ_x-aZ_p\).

The distance comparison is immediate from inclusion of admissible velocities: \[
d_g\le d_{a,b}\le d_{\mathcal D}.
\] For \(a=0\), the \(x\)-sphere is fully Riemannian and only the \(p\)-sphere is constrained; the squared distance is \(d_{g_x}^2+d_{CC,p}^2+d_T^2\). For \(b=0\), exchange the factors. For two nonzero coefficients no uncoupled product-distance identity is asserted.

The added unit field is divergence-free. Consequently its scalar operator is exactly \[
L_{a,b}
=L-\frac{(bZ_x-aZ_p)^2}{a^2+b^2}
=\Delta_x+\Delta_p+\Delta_T
+\frac{(aZ_x+bZ_p)^2}{a^2+b^2}.
\] On the same scalar Peter--Weyl/Fourier basis this gives \[
\lambda_{a,b}
=\frac{k_x(k_x+2)}{R_x^2}
+\frac{k_p(k_p+2)}{R_p^2}
+n^{\mathsf T}G^{-1}n
-\frac{(a q_x/R_x+b q_p/R_p)^2}{a^2+b^2}.
\] Equivalently it adds \((bq_x/R_x-aq_p/R_p)^2/(a^2+b^2)\) to the eigenvalue of \(L\). This proves positivity and the scalar spectral export directly. The symmetric and antisymmetric pencils \(a=b\) and \(a=-b\) have exchange-invariant kernels when the sphere markings are matched.

Thus, for constraints chosen within the marked two-form span, the full \(M_8\) ranks and homogeneous dimensions are:

\begin{longtable}[]{@{}
  >{\raggedright\arraybackslash}p{(\columnwidth - 6\tabcolsep) * \real{0.2143}}
  >{\raggedleft\arraybackslash}p{(\columnwidth - 6\tabcolsep) * \real{0.2857}}
  >{\raggedright\arraybackslash}p{(\columnwidth - 6\tabcolsep) * \real{0.2143}}
  >{\raggedleft\arraybackslash}p{(\columnwidth - 6\tabcolsep) * \real{0.2857}}@{}}
\toprule\noalign{}
\begin{minipage}[b]{\linewidth}\raggedright
Independent imposed one-form constraints
\end{minipage} & \begin{minipage}[b]{\linewidth}\raggedleft
Horizontal rank
\end{minipage} & \begin{minipage}[b]{\linewidth}\raggedright
Growth
\end{minipage} & \begin{minipage}[b]{\linewidth}\raggedleft
Homogeneous dimension
\end{minipage} \\
\midrule\noalign{}
\endhead
\bottomrule\noalign{}
\endlastfoot
None & 8 & \((8)\) & 8 \\
One nonzero linear combination & 7 & \((7,8)\) & 9 \\
Both independent constraints & 6 & \((6,8)\) & 10 \\
\end{longtable}

This table classifies the bracket and dimension data for that specified constraint family. The complete distance and cut synthesis in the main G05 result concerns the two-constraint orthogonal product.



\hypertarget{p1b:app:g06}{%
\subsection{CR complexes, Hardy spaces and bi-Hopf quantization: supplementary calculations}\label{p1b:app:g06}}

\hypertarget{p1b:app:g06:1}{%
\subsubsection{Closed operators and smooth cores}\label{p1b:app:g06:1}}

Every scalar operator displayed above is diagonal on the complete orthogonal product of the \(\mathcal H^{r,s}\) spaces and, when present, torus Fourier spaces. If its nonnegative eigenvalue on a component is \(\lambda_\gamma\), its self-adjoint spectral domain is \[
\operatorname{Dom}T=
\left\{f=\sum_\gamma f_\gamma:
\sum_\gamma\lambda_\gamma^2\|f_\gamma\|^2<\infty\right\}.
\] Its quadratic-form domain uses \(\lambda_\gamma\) in place of \(\lambda_\gamma^2\). Finite component sums form a core, since truncation converges in the graph norm. On that core the differential formulas agree exactly with the diagonal operators. Hence the nonnegative closures, or equivalently the associated Friedrichs realizations, have precisely these domains. An infinite zero eigenspace is fully included.

For an expression such as \(AB_c\), the operator intended is this closure of the smooth differential product. This statement avoids an unnecessary additional graph-domain restriction from requiring a vector in its zero eigenspace to lie separately in the domain of \(B_c\). The explicit paired basis from Section~\ref{p1b:der:g06:s39} similarly defines the closed Dolbeault differential with weights \(\sqrt{\lambda_\gamma}\). The product nilpotence and Hodge formulas can first be checked on the common finite core and then extended with those domains.

For the full untwisted CR complex, closed range follows from its positive spectral gap and the paired decomposition, even though its cohomology is infinite dimensional. For a fixed charge the quotient is compact and its positive spectral multiplicities in (G06.17) are finite. For the unrestricted CR product some positive multiplicities are infinite. This is why the finite index statements were made only for the basic or fixed-charge quotient complexes.

\hypertarget{p1b:app:g06:2}{%
\subsubsection{Other CR orientations and scalar choices}\label{p1b:app:g06:2}}

Reversing a sphere CR orientation replaces \(J_i\) by \(-J_i\) and interchanges \(Z_{1,i}\) with \(Z_{\bar1,i}\). The Hardy and anti-Hardy spaces are interchanged, as are \(A_i\) and \(B_{c,i}\). For positive Levi normalization use \(-\vartheta_i\) together with the reversed CR structure. The coframe then has geometric charge \(+2\).

Let \(\varepsilon_i=\pm1\) record the chosen sphere CR orientation while keeping the original geometric charge coordinate \(q_i\). All fixed-charge cohomology formulas apply with \(q_i\) replaced by \(\varepsilon_i q_i\), because the holomorphic base orientation and its positive line convention are also reversed. Thus positive holomorphic quantization requires \(\varepsilon_i q_i\ge0\) and has factor dimension \(\varepsilon_i q_i+1\). The one-factor Paneitz operator is unchanged; the joint operator must use the two Kohn choices matching the declared product CR structure.

Choosing the opposite orientation on the scalar torus replaces \(E\) by its conjugate complex structure. The dimensions, Fourier spectrum, and ordinary Dolbeault index above are unchanged. The natural reflection still exchanges the two complex structures rather than preserving either one.

These are explicit variations of the declared geometric data. They do not introduce an unmarked preferred CR orientation or identify the basic complex with the full complex family.

\hypertarget{p1b:app:g06:3}{%
\subsubsection{Constants and factorwise Dirichlet-to-Neumann operators}\label{p1b:app:g06:3}}

For the ordinary Euclidean ball bounded by \(S^3(R)\), a scalar spherical harmonic of degree \(k\) extends harmonically as \((r/R)^kY_k\). Its outward normal derivative at \(r=R\) is \[
\mathsf D_RY_k=\frac{k}{R}Y_k.
\] Thus \(\mathsf D_R\) is nonnegative, has only the constants in its kernel, and commutes with the Hardy projector and the one-factor Paneitz operator. On holomorphic degree \(n\) its eigenvalue is \(n/R\).

On \(P_v\) one may define the factorwise sum \[
\mathsf D_{xp}=\mathsf D_{R_x}\otimes I+
 I\otimes\mathsf D_{R_p}.
\] It has eigenvalues \(k_x/R_x+k_p/R_p\) and a unique scalar constant zero mode. This proves the product version of the constant-mode relation associated with Paper I, p1:thm:v167-hardy-constant-dtn-zero. The definition is a sum of two specified boundary operators. It does not identify \(P_v\) with a smooth hypersurface bounding a single ball.

The constant is also in the kernels of all Paneitz operators in Section~\ref{p1b:der:g06:s44} and in the Hardy range. Consequently any ordinary inverse taken on a Paneitz kernel complement cannot be applied to it. A generalized inverse with zero on that kernel is a different explicitly defined operator.

\hypertarget{p1b:app:g06:4}{%
\subsubsection{Unoriented axes and finite extremal projectors}\label{p1b:app:g06:4}}

For \(n>0\) the extremal-pair projector of Section~\ref{p1b:der:g06:s45} can be written without choosing a phase or an orientation for the axis: \[
P_{n,[\mathbf u]}
=\mathbf1_{\{n^2/4\}}\bigl((\mathbf u\cdot\mathbf J)^2\bigr).
\] It has rank two, is unchanged by \(\mathbf u\mapsto-\mathbf u\), and is a polynomial in \((\mathbf u\cdot\mathbf J)^2\) because the spectrum is finite. The product \[
P_{a,[\mathbf u_x]}\otimes P_{b,[\mathbf u_p]}
\] has rank four for \(a,b>0\). For \(n=0\) the one-dimensional identity is the appropriate degenerate version. These projectors are functions of the stated matrices within the specified quantization carriers.

A noncentral element of \(SU(2)\) can supply such an unoriented axis: write it as \(\cos\chi+\mathbf u\sin\chi\), \(0<\chi<\pi\). An element commutes with it exactly when it commutes with \(\mathbf u\), so its centralizer is the unique circle of quaternions in the real span of \(1,\mathbf u\). Inversion changes \(\mathbf u\) to \(-\mathbf u\) and leaves that circle fixed. The central elements \(+1,-1\) instead have the full group as centralizer and supply no such axis. Applying this argument separately to two specified noncentral elements provides the two-axis input for the product extremal-pair construction.

This establishes the mathematical interface to the source's noncentral-element and Borel--Weil axis results, including the zero-spin and central-element exceptions. The construction is defined relative to the chosen axis and representation.



\hypertarget{p1b:app:g07}{%
\subsection{Global metric geometry and deformations: supplementary calculations}\label{p1b:app:g07}}

\hypertarget{p1b:app:g07:1}{%
\subsubsection{Deriving the Berger diameter and injectivity extrema}\label{p1b:app:g07:1}}

It suffices by axial symmetry to take \(0\le h\le1\). For \(-1<\eta\le0\), (G07.25) is simply \(t_{\rm cut}=\pi R\sqrt{1+\eta h^2}\). Its maximum is \(\pi R\), at \(h=0\), and its minimum is \(\pi R\sqrt{1+\eta}=\pi R/\varepsilon\), at \(h=1\).

For \(\eta>0\), \(h>0\), use the continuous argument along the upper-half ellipse to define \[
\Phi(\tau,h)=\eta h\tau+
          \operatorname{Arg}(\cos\tau+i h\sin\tau),
\quad 0\le\tau\le\pi.
\] The first cut root is exactly \(\Phi(\tau,h)=\pi\). Indeed \(\Phi(0,h)=0\), \(\Phi(\pi,h)=(1+\eta h)\pi>\pi\), and \[
\partial_\tau\Phi
=h\left(\eta+\frac1{\cos^2\tau+h^2\sin^2\tau}\right)>0.
\] Put \(D=\cos^2\tau+h^2\sin^2\tau\) and \(L=1+\eta h^2\). Implicit differentiation gives \[
\tau'(h)
=-\frac{\eta\tau D+\sin\tau\cos\tau}{h(1+\eta D)},
\] and consequently \[
\frac1R\frac{dt_{\rm cut}}{dh}
=-\frac{\cos\tau\,F_{\eta,h}(\tau)}
        {\sqrt L\,h(1+\eta D)}.
\tag{G07.53}\label{p1b:der:g07:eq-G07-53}
\] At a cut point with \(h>0\), the first conjugate point is strictly later for \(\eta>0\); this property is included in the SU(2) cut/conjugate description summarized in \href{https://arxiv.org/abs/1710.02945}{Podobryaev, §2}. Thus \(F_{\eta,h}(\tau)>0\). Also \(\Phi(\pi/2,h)=\pi(1+\eta h)/2\). Therefore the derivative in (G07.53) is positive when \(\eta h<1\), zero when \(\eta h=1\), and negative when \(\eta h>1\). There is at most one interior extremum, and it is a maximum.

The endpoint values are \[
t_{\rm cut}(0)=R\tau_c(0),\qquad
t_{\rm cut}(1)=\frac{\pi R}{\sqrt{1+\eta}}
              =\pi R\varepsilon.
\] The second value follows from \(\tau_3(1)=\pi/(1+\eta)\). For \(0<\eta\le1\), the maximum is the \(h=1\) value. For \(\eta\ge1\), the maximum occurs at \(h=1/\eta\), where \(\tau=\pi/2\), giving \[
\max t_{\rm cut}
=\frac{\pi R}{2}\sqrt{\frac{1+\eta}{\eta}}
=\frac{\pi R}{2\sqrt{1-\varepsilon^2}}.
\] The minimum is always the smaller of the two endpoint values. Homogeneity identifies the maximum cut distance with diameter and the minimum with injectivity radius. This proves both (G07.26) and (G07.27), including all transition values by continuity.

\hypertarget{p1b:app:g07:2}{%
\subsubsection{Products of round Clifford tori}\label{p1b:app:g07:2}}

In one round sphere use Hopf coordinates \[
q=(\cos\rho\,e^{is},\,\sin\rho\,e^{it}),\qquad
0<\rho<\pi/2,\qquad s,t\pmod{2\pi}.
\] The ambient metric is \[
R^2(d\rho^2+\cos^2\rho\,ds^2+\sin^2\rho\,dt^2).
\] Thus a fixed-\(\rho\) torus is intrinsically flat, with induced metric \(R^2\cos^2\rho\,ds^2+ R^2\sin^2\rho\,dt^2\). For unit normal \(n=R^{-1}\partial_\rho\) and shape operator \(A_n=-\nabla n\), differentiation of these two warping factors gives principal curvatures \(\tan\rho/R\) and \(-\cot\rho/R\).

Consequently \[
N=T_{\rho_x}^2\times T_{\rho_p}^2\times T_G^2
       \ \subset\ (M_8,g_0)
\] is a flat six-dimensional submanifold of codimension two. With the two unit normals oriented toward increasing \(\rho_i\), its mean curvature vector, using the dimension-averaged convention, is \[
H_N=\frac16\left[
\frac{\tan\rho_x-\cot\rho_x}{R_x}n_x+
\frac{\tan\rho_p-\cot\rho_p}{R_p}n_p
\right].
\tag{G07.54}\label{p1b:der:g07:eq-G07-54}
\] Independence of the normals proves minimality exactly when \(\rho_x=\rho_p=\pi/4\). At that point its volume is \[
\operatorname{Vol}N
=16\pi^6R_x^2R_p^2\sqrt{\det G}.
\] It is not area-stable: varying just one constant \(\rho_i\) multiplies the corresponding area by \(\sin(2\rho_i)\), which has a strict maximum at \(\pi/4\). These are two independent negative second-variation directions. This assertion concerns the embedded minimal torus and does not assert instability of the ambient round sphere metric.

\hypertarget{p1b:app:g07:3}{%
\subsubsection{Homogeneous Einstein obstruction and relative-scale flow}\label{p1b:app:g07:3}}

There is no homogeneous Einstein metric on the compact manifold \(S^3\times S^3\times T^2\). To prove this, suppose \(\mathrm{Ric}=\lambda g\). If \(\lambda>0\), the Bonnet--Myers theorem forces finite fundamental group, contradicting \(\pi_1=\mathbb Z^2\). For a Killing field \(X\), the integrated Killing-field identity is \[
\int|\nabla X|^2=\int\mathrm{Ric}(X,X).
\] If \(\lambda<0\), it excludes every nonzero Killing field, contradicting homogeneity. If \(\lambda=0\), all Killing fields are parallel. Homogeneity makes their values span the tangent space, so the curvature vanishes. A complete flat simply connected cover is Euclidean, whereas this manifold has universal cover \(S^3\times S^3\times\mathbb R^2\), with nonzero third homotopy group. This excludes the remaining case. The argument uses homogeneity and makes no assertion about nonhomogeneous Einstein metrics.

For round sphere factors in the full normalized flow (G07.43), let \[
\delta=\frac{A_x-A_p}{A_x+A_p}.
\] Since \(\dot A_i=-4+\gamma A_i\), direct differentiation cancels the common normalization term and gives \[
\dot\delta=\frac{8\delta}{A_x+A_p}.
\tag{G07.55}\label{p1b:der:g07:eq-G07-55}
\] Unequal relative sphere scales move away from equality. This is independent of the separate shape-rounding property.

The six-dimensional leaf \(P_v\) has a different global normalization, \(\dot g=-2\mathrm{Ric}+\tfrac13 Sg\). Its equal-radius round product is Einstein and stationary. On the fixed-volume round subfamily write \[
A_x=a_*e^v,\qquad A_p=a_*e^{-v}.
\] The sphere-volume product is constant because \(A_xA_p=a_*^2\). The normalized equations give \[
\dot v=\frac4{a_*}\sinh v
       =\frac4{a_*}v+O(v^3).
\tag{G07.56}\label{p1b:der:g07:eq-G07-56}
\] Thus even this Einstein product is unstable to relative factor rescaling. The first-order variation at \(v=0\) is parallel, trace-free, and divergence-free; it is an actual volume-preserving metric direction. It is not a zero mode of the fixed-Einstein-constant deformation equation. In particular it is compatible with the one-sphere infinitesimal Einstein rigidity statement at p1:prop:infinitesimal-einstein-deformations.

\hypertarget{p1b:app:g07:4}{%
\subsubsection{Algebraic formulas for the full invariant metric cones}\label{p1b:app:g07:4}}

This supplement makes precise what extends when the orthogonal product assumption is dropped. Use the fixed eight-element invariant frame, with \(c_{ij}^{\,k}\) equal to the two SU(2) bracket blocks \(-2\epsilon_{ijk}\), and zero for torus or cross brackets. Let \(Q=(Q_{ab})\in\mathrm{Sym}_8^+\) be its metric matrix. The Koszul formula is the finite algebraic expression \[
\Gamma_{ij}^{\,k}
=\frac12Q^{k\ell}
 \left(c_{ij}^{\,m}Q_{m\ell}
       -c_{j\ell}^{\,m}Q_{mi}
       +c_{\ell i}^{\,m}Q_{mj}\right).
\tag{G07.57}\label{p1b:der:g07:eq-G07-57}
\] Because the coefficients are invariant, their frame derivatives vanish. In the convention \(R(\xi_i,\xi_j)\xi_k=R_{ijk}^{\ \ \ell}\xi_\ell\), \[
R_{ijk}^{\ \ \ell}
=\Gamma_{jk}^{\,m}\Gamma_{im}^{\,\ell}
 -\Gamma_{ik}^{\,m}\Gamma_{jm}^{\,\ell}
 -c_{ij}^{\,m}\Gamma_{mk}^{\,\ell},\quad
\mathrm{Ric}_{jk}=R_{ijk}^{\ \ i},\quad
S=Q^{jk}\mathrm{Ric}_{jk}.
\tag{G07.58}\label{p1b:der:g07:eq-G07-58}
\] These are exact rational formulas on the positive cone. The homogeneous flow is \(\dot Q=-2\mathrm{Ric}(Q)\), or \(\dot Q=-2\mathrm{Ric}(Q)+(S(Q)/4)Q\) in dimension-eight volume normalization. Any stated symmetry-fixed cone, including the Hopf-invariant and exchange-fixed cones, is preserved by equivariance and uniqueness of this ODE. This does not imply that every such metric has an orthogonal product splitting.

The invariant volume is a constant multiple of Haar measure. Unimodularity makes the invariant vector fields divergence-free, so the scalar Laplacian is \[
\Delta_Q=-Q^{ab}\xi_a\xi_b
         =Q^{ab}T_aT_b,\qquad T_a=-i\xi_a.
\tag{G07.59}\label{p1b:der:g07:eq-G07-59}
\] On the Peter--Weyl sector \((k_x,k_p,n)\), the two triples of \(T_a\) are \(2J_a\) in the spin-\(k_i/2\) representations and the torus pair acts by \(n_1,n_2\). Consequently (G07.59) is an explicit Hermitian matrix of size \((k_x+1)(k_p+1)\), repeated with that same multiplicity from the other matrix-coefficient indices. Positivity follows either from the gradient form or by factoring the positive matrix \(Q^{-1}\).

The sphere Casimirs commute with every \(T_a\), so the independent scalar degree bands remain invariant for all these invariant metrics, even with cross terms. Individual Hopf charges require the additional Hopf symmetry; they are not conserved for an arbitrary \(Q\in\mathrm{Sym}_8^+\). Inversion transports the right-invariant convention here to Paper I-C's left-invariant metric convention. These formulas supply the larger-cone curvature, flow, and scalar-operator construction. They do not supply closed-form distances, cut loci, or global flow limits for arbitrary mixed metrics.



\hypertarget{p1b:app:g08}{%
\subsection{Full bundle spectra, constraints and analytic invariants: supplementary calculations}\label{p1b:app:g08}}

\hypertarget{p1b:app:g08:1}{%
\subsubsection{Exact TT exchange trace}\label{p1b:app:g08:1}}

For equal round factors, the exchange trace on the full symmetric-two-tensor bundle has contributions only from two types of diagonal summands. The \(T_x^*\otimes T_p^*\) summand gives the ordinary flip trace \(C_R(2t)\Theta_{\rm diag}(t)\). The symmetric torus summand gives \(A_R(2t)\operatorname{tr}(\operatorname{Sym}^2S_2) \Theta_{\rm diag}(t)\), with \(\operatorname{tr}(\operatorname{Sym}^2S_2)=1\). The pure \(x/p\) symmetric tensors and the sphere/torus mixed tensors are exchanged in pairs and have zero trace.

Removing the full metric trace subtracts \(A_R(2t)\Theta_{\rm diag}\), cancelling the torus term. Both the one-form exchange trace and the Killing exchange trace vanish: their sphere summands are exchanged and \(\operatorname{tr}S_2=0\) on the torus summands. The equivariant version of the rank--nullity argument in Section~\ref{p1b:der:g08:s64} therefore gives \[
\boxed{
Z_{{\rm TT},B}(t)=C_R(2t)\Theta_{\rm diag}(t),\qquad
Z_{{\rm TT},\pm}(t)=\tfrac12[Z_{\rm TT}(t)\pm Z_{{\rm TT},B}(t)].
}
\tag{G08.58}\label{p1b:der:g08:eq-G08-58}
\] Its limit as \(t\to\infty\) is zero, so the four parallel TT zero modes split \(2+2\) by exchange. Its small-time leading coefficient is three times the scalar fixed-set coefficient in (G08.47).

\hypertarget{p1b:app:g08:2}{%
\subsubsection{Product overlap coefficients and selection rules}\label{p1b:app:g08:2}}

Fix an actual equivariant pointwise bundle contraction \(\mathcal C:E_b\otimes E_c\to E_a\), using duals or Clifford multiplication where required by the types. For factorized normalized harmonics, Fubini gives \[
\int_M\langle Y_a,\mathcal C(Y_b,Y_c)\rangle\,dV
=\left(\int_{S_x^3}\cdots\right)
 \left(\int_{S_p^3}\cdots\right)
 \left(\int_{T^2_G}\cdots\right).
\tag{G08.59}\label{p1b:der:g08:eq-G08-59}
\] Each sphere integral is a product of left/right Clebsch--Gordan coefficients and the finite recoupling coefficient for \(\mathcal C\). This follows by decomposing the pointwise tensor product and applying Schur orthogonality. Thus all four SU(2) spin triangles and their integrality conditions must hold, as must their magnetic-weight conservation rules. The torus integral enforces addition of Fourier labels, including any compatible shifts of the input and output bundles. For scalar multiplication it is simply \(n_a=n_b+n_c\), with the usual volume normalization. These are necessary selection rules; an allowed triple can still have a zero recoupling coefficient. No scalar overlap is assigned to a tensor product whose bundle type has not been contracted to the stated target.

\hypertarget{p1b:app:g08:3}{%
\subsubsection{Conformal/Yamabe second variation on the round product}\label{p1b:app:g08:3}}

The sphere's global Yamabe minimization statement does not follow merely by tensoring factors. For the eight-dimensional product, the conformal functional on positive smooth \(u\) is \[
\mathcal Y_g(u)=
\frac{\int_M[(14/3)|\nabla u|^2+Su^2]\,dV}
     {(\int_Mu^{8/3}\,dV)^{3/4}}.
\tag{G08.60}\label{p1b:der:g08:eq-G08-60}
\] The powers are those of \(\widehat g=u^{4/(8-2)}g\). Since \(S\) is constant, \(u=1\) is a critical point and \(\mathcal Y_g(1)=SV^{1/4}\). For a variation \(u=1+s\varphi\) with \(\int\varphi=0\), expanding numerator and denominator to second order gives \[
\frac{d^2}{ds^2}\mathcal Y_g(1+s\varphi)\big|_{s=0}
=\frac{28}{3}V^{-3/4}
\int_M\left[|\nabla\varphi|^2-\frac S7\varphi^2\right]\,dV.
\tag{G08.61}\label{p1b:der:g08:eq-G08-61}
\] Therefore nonnegative second variation is equivalent to \[
\lambda_1(\Delta_0)\ge\frac{6(a+b)}7.
\] Using the exact product gap, this is equivalent to \[
\frac25\le\frac ab\le\frac52,\qquad
\lambda_{T,1}\ge\frac{6(a+b)}7.
\] Strict inequalities give strict second-variation stability modulo constants; equality requires higher-order analysis. Enlarging the torus eventually violates the last condition, producing an explicit negative direction from its first scalar mode. This establishes a conformal stability criterion, not global minimization or a value for the Yamabe invariant of the entire manifold.

\hypertarget{p1b:app:g08:4}{%
\subsubsection{Independent finite certificates for the product constraints}\label{p1b:app:g08:4}}

The matrices in the accompanying derivation and verification record use \[
\alpha_1=\sigma_1\otimes I,\quad
\alpha_2=\sigma_2\otimes I,\quad
\beta_+=\sigma_3\otimes\sigma_1,\quad
\beta_-=\sigma_3\otimes\sigma_2.
\] The full Clifford generators are \[
c_{x,i}=(-i\sigma_i)\otimes I\otimes\alpha_1,\quad
c_{p,i}=I\otimes(-i\sigma_i)\otimes\alpha_2,\quad
c_{T,\pm}=-iI\otimes I\otimes\beta_\pm.
\] These satisfy \(c_\alpha c_\beta+c_\beta c_\alpha =-2\delta_{\alpha\beta}\). Their four reversed-normal product equals the factor flip times \(U\) in (G08.16).

For the periodic torus, zero Fourier label, \(R_x=1\), and \(R_p=1/2\), take constant coefficient spinors and vector-spinors in this frame. Let \(P_c\) be the resulting 128-by-16 twistor matrix, \(Q_c\) the gamma-projected Rarita--Schwinger matrix, and \[
\mathcal M=
\left[\Pi_\gamma-P_c(P_c^*P_c)^{-1}P_c^*\right]Q_cP_c.
\] The exact rational matrix calculation gives \[
\boxed{\operatorname{tr}(\mathcal M^*\mathcal M)=384/5>0.}
\tag{G08.62}\label{p1b:der:g08:eq-G08-62}
\] This is a concrete certificate that \(Q\) sends some twistor-image fields into the transverse space on the product. The full curvature identity (G08.31), the twistor norm identity, and the positive-energy identity (G08.34) are checked in the same independently constructed connection matrices.

For TT tensors the script constructs the complete symmetric-tensor connection, divergence, and trace maps in a 36-dimensional fiber. It verifies \(\operatorname{tr}\Delta_L=\Delta_0\operatorname{tr}\) and \(\delta\Delta_L=\Delta_1\delta\) as operator matrix identities. At the same radius values, the following coefficient-block spectra are obtained by solving the simultaneous eigenvalue, trace, and divergence equations:

\begin{longtable}[]{@{}
  >{\raggedright\arraybackslash}p{(\columnwidth - 4\tabcolsep) * \real{0.3333}}
  >{\raggedright\arraybackslash}p{(\columnwidth - 4\tabcolsep) * \real{0.3333}}
  >{\raggedright\arraybackslash}p{(\columnwidth - 4\tabcolsep) * \real{0.3333}}@{}}
\toprule\noalign{}
\begin{minipage}[b]{\linewidth}\raggedright
Orbital block \((j_x,j_p)\)
\end{minipage} & \begin{minipage}[b]{\linewidth}\raggedright
TT eigenvalue: multiplicity
\end{minipage} & \begin{minipage}[b]{\linewidth}\raggedright
Total coefficient dimension
\end{minipage} \\
\midrule\noalign{}
\endhead
\bottomrule\noalign{}
\endlastfoot
\((0,0)\) & \(0:4,\ 4:6,\ 12:5,\ 16:6,\ 20:9,\ 48:5\) & 35 \\
\((1/2,0)\) & \(3:6,\ 9:8,\ 19:18,\ 25:12,\ 51:10\) & 54 \\
\end{longtable}

The spectator multiplicities are respectively one and two. These are finite checks of the universal formula (G08.24), obtained directly from covariant derivatives rather than from that multiplicity formula itself.

\hypertarget{p1b:app:g08:5}{%
\subsubsection{Heat coefficients of the positive Rarita constraint operator}\label{p1b:app:g08:5}}

The coupled operator \(\mathcal A_{\rm RS}\) has a parallel gamma projection, so it has induced connection and endomorphism \[
E_\gamma=
\frac S4I_\gamma+
\Pi_\gamma\left(\sum_{\alpha<\beta}
c(e_\alpha)c(e_\beta)\otimes
\Omega^T_{\alpha\beta}\right)\Pi_\gamma .
\tag{G08.63}\label{p1b:der:g08:eq-G08-63}
\] Tensor ordering in this formula is immaterial if used consistently; the accompanying derivation and verification record uses vector first. It follows either from the twisted Lichnerowicz formula or by squaring the connection matrices.

For the round product, the isotropy decomposition (G08.28) gives \[
\operatorname{tr}E_\gamma=180(a+b),\qquad
\operatorname{tr}E_\gamma^2
=345(a^2+b^2)+594ab,\qquad
\operatorname{tr}\Omega^\gamma_{\alpha\beta}
                 \Omega^\gamma_{\alpha\beta}
=-360(a^2+b^2).
\] Here is a trace derivation of the non-scalar endomorphism term. On the two rank-32 high-spin summands, \(E_\gamma\) is respectively \((5a+3b)/2\) and \((3a+5b)/2\). On the remaining rank-48 summand it is \(\tfrac32(a+b)I-\mathrm{Ric}_\gamma\). Before gamma projection the latter comes from two rank-16 blocks with Ricci values \(2a,2b\), and two rank-16 torus blocks with value zero. The Clifford contraction weights of these four blocks are \(3,3,1,1\). Since \(\mu\mu^*=8I\), the projected Ricci trace is \(20(a+b)\) and its squared trace is \(16(a^2+b^2)+9(a+b)^2\). Inserting these into the rank-48 expression, and adding the rank-32 contributions, gives the displayed \(E_\gamma\) traces. The curvature trace follows from summing \(-2s(s+1)a^2\), respectively \(-2t(t+1)b^2\), on the same isotropy summands.

Substitution in (G08.44) yields \[
\boxed{
\begin{aligned}
A_0(\mathcal A_{\rm RS})&=112V,\\
A_2(\mathcal A_{\rm RS})&=-68V(a+b),\\
A_4(\mathcal A_{\rm RS})
 &=V\left[\frac{37}{2}(a^2+b^2)+49ab\right].
\end{aligned}}
\tag{G08.64}\label{p1b:der:g08:eq-G08-64}
\] These are coefficients of the positive coupled constraint operator (G08.34). They must not be substituted for coefficients of \(Q^2\), whose principal speeds differ as shown in (G08.50). The exact \(Q^2\) heat trace itself is fully specified by (G08.33): \[
Z_{Q^2}(t)=
\sum_{j_x,j_p,n}(2j_x+1)(2j_p+1)
\operatorname{tr}\exp[-t\,\mathbb Q(j_x,j_p,n)^2].
\] Ellipticity and compactness give absolute convergence for \(t>0\). This formula retains its full matrix coupling at every level.



\hypertarget{p1b:app:g09}{%
\subsection{Reproducing kernels, evaluation measures and exact moments: supplementary calculations}\label{p1b:app:g09}}

\hypertarget{p1b:app:g09:1}{%
\subsubsection{\texorpdfstring{Exact \(1/48\) witness, spectrum, and covariance diagnostic}{Exact 1/48 witness, spectrum, and covariance diagnostic}}\label{p1b:app:g09:1}}

This technical section supplies the extended finite calculations behind the correction already stated in Section~\ref{p1b:der:g09:s78}. It introduces no new carrier or quadrature premise.

Let \(u=q_0^2-q_1^2\) and \(f=q_0^2\). The sphere moments are \[
\int q_0^6=\frac5{64},\qquad
\int q_0^4q_1^2=\int q_0^2q_1^4=\frac1{64},\qquad
\int f u^2=\frac1{16},\qquad \int u^2=\frac16.
\tag{G09.53}\label{p1b:der:g09:eq-G09-53}
\] For the 16 half-coordinate vertices, \(u=0\). Of the eight axis vertices, only \(q=\pm e_0\) contribute to \(fu^2\). Thus the discrete value is \(2/24=1/12\) and \[
\frac1{24}\sum_{q\in2T}fu^2-\int fu^2
=\frac1{12}-\frac1{16}=\frac1{48}.
\tag{G09.54}\label{p1b:der:g09:eq-G09-54}
\] This is a matrix element, not an eigenvalue claim. It becomes \(1/8\) on the unit vector \(\sqrt6\,u\). The distinct test \(q_0^6\) alone has discrepancy \(3/32-5/64=1/64\), so it cannot replace the compiler witness.

For the full correction spectrum, choose the top-degree basis \[
h_i=q_i^2-q_3^2\quad(i=0,1,2),\qquad
b_{ij}=q_iq_j\quad(0\le i<j\le3).
\] The two sets are orthogonal. On the \(h_i\) block the continuum Gram and discrepancy bilinear matrices are \[
G_h=\frac1{12}(I_3+\mathbf 1\mathbf 1^T),\qquad
A_h=\frac1{96}
\begin{pmatrix}2&-1&-1\\-1&-2&-1\\-1&-1&-2\end{pmatrix},
\qquad
D_h=G_h^{-1}A_h.
\tag{G09.55}\label{p1b:der:g09:eq-G09-55}
\] Its eigenvalues are \(1/4,-1/8,-1/8\). For \(b_{ij}\), the Gram entry is \(1/24\); the discrepancy entry is \(-1/192\) if \(0\in\{i,j\}\) and \(1/192\) otherwise. All these entries are diagonal by sign symmetry. Together with the five lower-degree zero modes, \[
\boxed{
\operatorname{spec}D_{q_0^2}
=\{0^{[5]},(1/4)^{[1]},(1/8)^{[3]},(-1/8)^{[5]}\},
\quad
\operatorname{rank}=9,\quad
\|D_{q_0^2}\|=\frac14,\quad\operatorname{tr}D_{q_0^2}=0.
}
\tag{G09.56}\label{p1b:der:g09:eq-G09-56}
\] The zero trace also follows from the constant kernel diagonal and the exact quadrature of the symbol itself. The finite and continuum quadratic moments remain positive operators; their nonzero difference need not be positive.

The covariance error can be seen without a floating-point calculation. The two normalized test directions \(q_1^2-q_2^2\) and \(2q_1q_2\) have equal norm and are rotated into one another by a \(\pi/4\) rotation in the \((q_1,q_2)\) plane. This rotation fixes \(f=q_0^2\). Their discrepancy eigenvalues are \(-1/8\) and \(+1/8\), respectively. Hence \(D_f\) cannot commute with that rotation; the uncorrected sampled quadratic does not have the full continuum stabilizer covariance. Subtracting \(D_f\) restores it.

If \(v\) is a unit vector in every other factor with isometric sampling, (G09.39) reproduces (G09.54) on \(u\otimes v\). One may take \(v=1\) when its Fourier set includes zero, or any retained unit character when it does not. For \(f=x_0^2=R_x^2q_0^2\) and the dimensionless unit test \(\sqrt6\,u\), the discrepancy is \(R_x^2/8\). If the unnormalized test is \(x_0^2-x_1^2=R_x^2u\), the matrix element is \(R_x^6/48\). These are changes of symbol and test-vector units, not new normalization factors in the original \(1/48\).

\hypertarget{p1b:app:g09:2}{%
\subsubsection{The exact state ambiguity and observable reconstruction}\label{p1b:app:g09:2}}

This is a scope proof for Section~\ref{p1b:der:g09:s76}. It describes which component of an operator the already constructed POVM can determine; it adds no extra measurement family.

On a nonconstant sphere band, choose real orthonormal functions \(a=1\) and \(b=2q_0\). The nonzero Hermitian operator \[
A=i(|a\rangle\langle b|-|b\rangle\langle a|)
\] has \(\langle K_q,AK_q\rangle=0\) at every point, since the scalar basis functions are real. Its trace is zero. For sufficiently small \(\epsilon>0\), \(I/D_m\pm\epsilon A\) are different positive density matrices with exactly the same outcome density. Tensoring with any fixed state on the other factors retains this ambiguity.

For a torus band with \(|F|>1\), every Fourier basis state has pointwise squared modulus one. Thus even its distinct diagonal pure states have the same outcome density. This is also seen from \[
|F-F|\le |F|(|F|-1)+1<|F|^2\quad(|F|>1).
\] Combined with the sphere argument, the scalar product POVM is informationally complete only in the trivial case \(N=1\). The exact number of complex operator directions it determines is nevertheless (G09.27); for one \(\mathcal P_2\) it is \(D_4=55\) out of \(14^2=196\), and for two such bands with a singleton \(F\) it is \(55^2=3025\) out of \(196^2=38416\).

Write \(\mathcal B_{\mathcal V}^{-1}\) for the inverse of the strictly positive transform on \(\mathcal V\). For an arbitrary operator \(A\) on \(W\), the Hilbert--Schmidt orthogonal projection onto the moment image is \[
\boxed{
\mathcal R(A)=
T_{\mathcal B_{\mathcal V}^{-1}(T^*A/N)}.
}
\tag{G09.57}\label{p1b:der:g09:eq-G09-57}
\] The function \(T^*A\) belongs to \(\mathcal V\). If \(A=T_f\) with \(f\in\mathcal V\), then \(T^*A/N=\mathcal Bf\), so \(\mathcal R(A)=A\). If \(A\) is orthogonal to the image of \(T\), then \(T^*A=0\). These two statements prove the orthogonal-projection formula. For an observed density \(p_\rho=T^*\rho\), it reconstructs \(\mathcal R(\rho)\), not necessarily \(\rho\). The inverse is a finite algebraic operation; positivity of \(\mathcal R(\rho)\) is not asserted.

The product formula for second-moment noise also follows without any commutativity assumption on compressed operators. For arbitrary bounded scalar symbols \(f,g\), \[
T_{\bar f g}-T_f^*T_g
=PM_{\bar f}(I-P)M_gP.
\tag{G09.58}\label{p1b:der:g09:eq-G09-58}
\] For a list of symbols this is a positive block matrix, the Gram matrix of their leaked multiplication channels. For a real \(f\) in a state \(\rho\), the outcome variance decomposes as \[
\operatorname{tr}(\rho T_{f^2})
 -(\operatorname{tr}\rho T_f)^2
=\left[\operatorname{tr}(\rho T_f^2)
 -(\operatorname{tr}\rho T_f)^2\right]
 +\operatorname{tr}\rho(T_{f^2}-T_f^2).
\] This supplies a direct interpretation of (G09.19) as the positive compression contribution to a scalar outcome variance.

\hypertarget{p1b:app:g09:3}{%
\subsubsection{Irregular frequency sets and positive-polar normalization}\label{p1b:app:g09:3}}

These are comparisons requiring a chosen frequency sequence or an alternative injective evaluation, so they are supplementary.

For \(F_N=\{(2k,0):|k|\le N\}\), \(|F_N|\to\infty\), yet \(d_{F_N}(1,0)=d_{F_N}(0,1)=0\). The coherent transform never recovers either first circle character. It is periodic under \(\theta_t\mapsto\theta_t+\pi\) and is independent of \(\theta_E\). This gives the exact failure of replacing the frequency-overlap condition in Section~\ref{p1b:der:g09:s80} by cardinality growth.

For an arbitrary weighted evaluation \(S\) that is injective, \(\mathsf G=S^*S>0\) and the positive-polar isometry is \[
V=S\mathsf G^{-1/2},\qquad
\widetilde E_a=
\mathsf G^{-1/2}E_a\mathsf G^{-1/2},\qquad
\sum_a\widetilde E_a=I.
\tag{G09.59}\label{p1b:der:g09:eq-G09-59}
\] Its sampled first moment is \(\mathsf G^{-1/2}S^*\operatorname{diag}(f)S\mathsf G^{-1/2}\). The normalization by itself proves no equality to \(T_f\). When \(S\) is already isometric, \(\mathsf G=I\) and this operation does not remove any higher-symbol quadrature error. If evaluation is not injective, the inverse in (G09.59) does not exist on the whole proposed carrier.

For a rectangular tensor-product evaluation, \(\mathsf G=\mathsf G_x\otimes\mathsf G_p\otimes\mathsf G_F\); if all factors are positive definite, its inverse square root and hence its polar isometry tensor in the same order. Exchange covariance holds when the selected evaluation and its weights are exchange-compatible.



\hypertarget{p1b:app:g10}{%
\subsection{Finite quotients and marked reference geometry: supplementary calculations}\label{p1b:app:g10}}

\hypertarget{p1b:app:g10:1}{%
\subsubsection{Exact exceptional support covering radii}\label{p1b:app:g10:1}}

This technical section proves the finite metric constants stated in Section~\ref{p1b:der:g10:s85}. All supports are the explicit unit-quaternion sets already specified in Paper I.

For a finite support spanning \(\mathbb R^4\) with the origin in the interior of its convex hull, put \(h(q)=\max_{v\in\Gamma}q\cdot v\). Its angular covering radius has cosine \(\min_{|q|=1}h(q)\). The polar body \[
\mathcal C_\Gamma=\{x:x\cdot v\le1\text{ for all }v\in\Gamma\}
\] therefore gives \[
\cos\rho_{\rm unit}(\Gamma)=
\left(\max_{x\in\mathcal C_\Gamma}|x|\right)^{-1}.
\tag{G10.58}\label{p1b:der:g10:eq-G10-58}
\] A convex norm attains its maximum over a compact polytope at a vertex, which reduces the computation to a finite linear-constraint problem.

For \(2T\), signed coordinate permutations reduce to \(a\ge b\ge c\ge d\ge0\). Its supporting constraints are \(a\le1\) and \(a+b+c+d\le2\). All vertices of this ordered chamber are intersections of four of its six bounding hyperplanes. There are seven feasible distinct vertices; their squared norms are at most \(2\), attained at \((1,1,0,0)\). Hence \(\cos\rho=1/\sqrt2\).

For \(2O\), add \(a+b\le\sqrt2\), from its two-coordinate vertices. The seven chamber inequalities give thirteen feasible vertices after all four-hyperplane intersections are checked. The maximum squared norm is \[
24-16\sqrt2,
\] attained, for example, at \((2\sqrt2-2,\,2-\sqrt2,\,2-\sqrt2,\,0)\). Its inverse square root is \((2+\sqrt2)/4\). The exact script checks feasibility and each norm inequality in the quadratic field; it does not use approximate optimization to decide these constants.

For \(2I\), let \(\phi=(1+\sqrt5)/2\). Its neighboring vertices have pairwise inner product \(\phi/2\). Four such vertices \(v_1,\ldots,v_4\) form a regular tetrahedral supporting facet. For \(s=v_1+\cdots+v_4\), \[
s\cdot v_i=1+3\phi/2,\qquad
|s|^2=4(1+3\phi/2),\qquad
\frac{s\cdot v_i}{|s|}
=\sqrt{\frac{7+3\sqrt5}{16}}.
\tag{G10.59}\label{p1b:der:g10:eq-G10-59}
\] The script constructs every such facet through the identity, checks its supporting inequality against all 120 vertices, and finds exactly twenty facets there. Their left translates give 600 distinct facets. Their 1200 triangular ridges each occur twice. Thus this set is closed under facet adjacency; the facet adjacency graph of a full-dimensional convex polytope is connected, so there are no omitted facets. One can see that connectivity by radially identifying the boundary with \(S^3\) and joining facet interiors by a path transverse to its ridges.

Every facet has the same distance in (G10.59) from the origin. A ball of that radius lies in the convex hull, and a facet normal realizes the minimum support value. This proves both directions of the covering-radius formula without a numerical extremization assumption.

\hypertarget{p1b:app:g10:2}{%
\subsubsection{The projective quotient pencil and its connections}\label{p1b:app:g10:2}}

Take the affine coordinate \(x=z,y=1\), so \[
f_6=z(z^4-1),\quad f_8=z^8+14z^4+1,\quad
g_{12}=z^{12}-33z^8-33z^4+1,\quad u_{12}=f_6^2.
\] The source invariance and direct expansion give \[
f_8^3=g_{12}^2+108u_{12}^2,\qquad
u_{12}'g_{12}-u_{12}g_{12}'=-2f_6f_8^2.
\tag{G10.60}\label{p1b:der:g10:eq-G10-60}
\] The two degree-twelve homogeneous forms \(u_{12},g_{12}\) have no common projective zero. In the affine chart \(\gcd(u_{12},g_{12})=1\); at infinity the homogeneous value of \(g_{12}\) is nonzero. Their ratio gives a map of degree twelve.

The five finite simple roots of \(f_6\), together with infinity, are six points of ramification index two. The eight distinct roots of \(f_8\) have ramification index three, split between two branch values by the first identity in (G10.60). There are no overlaps among these roots. The ramification degree is \(6+2\cdot8=22=2\cdot12-2\). The map is invariant under \(A_4\); generic orbit size twelve and map degree twelve identify its generic fibres with the group orbits. The resulting finite degree-one map from the projective quotient to \(\mathbb{CP}^1\) is an isomorphism. Its branch signature is \((2,3,3)\).

For the product, the map has covering degree \(12^2\); the Segre embedding has hyperplane pullback \(h_x+h_p\). Therefore the complete map pulls the hyperplane class back to \(12(h_x+h_p)\). Since \(h_x^2=h_p^2=0\) and \(\int h_xh_p=1\), its square is \(288\), as in (G10.55). The ramification locus consists of the six/eight ramification sets in either factor times the other projective line, with their indicated multiplicities. Its divisor class is \(22h_x+22h_p\).

The unitary polar map identifies the source line bundle with the image line bundle. It does not identify arbitrary connections on those bundles. The projection connection has the curvature pulled back from the projective target; that curvature component vanishes at a ramification point. The connection transported from the twelfth tautological power has its nonzero constant-curvature component there. For the product, the difference is the sum of the two factor differences; at a pair of ramification points both projection-curvature components vanish. The equality of Chern classes follows from the bundle isomorphism, while pointwise connection equality fails in general.

\hypertarget{p1b:app:g10:3}{%
\subsubsection{Spectral order, finite transport, and its scope}\label{p1b:app:g10:3}}

In sector order \((0,0),(\frac12,\frac12),(1,0),(0,1), (\frac12,0),(0,\frac12)\), the one-factor dimensionless LC energy has values \[
E=(0,\sqrt3,2,2,3/2,3/2).
\] The maximal right-Hopf weights in absolute value are \((0,1,0,2,0,1)\). Each satisfies the required inequality. Equality \(E=N\) holds only on the scalar and the weight-\(+2\) line in \((0,1)\); equality \(E=-N\) holds on the scalar and the weight-\(-2\) line. All these equality states have acted left spin zero, so compatible left finite averaging preserves them.

For a positive sum of commuting nonnegative operators, the kernel is the intersection of the individual kernels: the vanishing quadratic form forces every summand to annihilate the vector. Applied to \((E_x-N_x)+(E_p+N_p)\) and its opposite, this gives the four-dimensional kernels in (G10.57). No comparison of eigenvalue labels alone replaces this operator positivity argument.

If \(\mathcal S\) is a specified isometric evaluation, the sampled operators are \(\mathcal S E\mathcal S^*\) and \(\mathcal S N\mathcal S^*\) on its image. Conjugation preserves order, spectrum and kernel dimension. An arbitrary positive-polar normalization of a different evaluation provides a comparison only after that evaluation has been specified, as in G09.59.

\hypertarget{p1b:app:g10:4}{%
\subsubsection{The quotient heat comparison and evidence boundaries}\label{p1b:app:g10:4}}

The independent numerical oracle uses the Poisson expression for the covering scalar kernel. For \(a=t/R^2>0\) and \(0<\theta<\pi\), in probability measure it is \[
K^\mu(a;\theta)=
\frac{e^a\sqrt\pi}{4a^{3/2}\sin\theta}
\sum_{j\in\mathbb Z}(\theta-2\pi j)
e^{-(\theta-2\pi j)^2/(4a)}.
\tag{G10.61}\label{p1b:der:g10:eq-G10-61}
\] Differentiate the circle Gaussian Poisson sum to obtain this formula from the sine spectral kernel of G08. At \(0,\pi\) take its derivative limit; the exact script uses those limits rather than dividing by zero. The quotient heat trace is the average of these values over the quaternion angles. It is compared with the integer multiplicity sum obtained from the group characters. The two evaluated routes agree at 70-digit precision within the reported errors. Their tails at the listed times are smaller than the \(10^{-50}\) comparison threshold.

The tests also compute the group-averaged kernels directly on a degree-one product space, where the common projector has rank five and the independent projector rank one. This is a finite counterexample to conflating the two actions, not an objection to either correctly typed quotient.

The supplied verification and appendix-placement rules were applied to the mathematical source and derivation record. No PDF was generated or rendered, and no TeX compilation or appearance review was used. External mathematical references were read as text. The supplied paper sources and archive remain unchanged. The first script run exposed conversion of Python integer literals into an algebraic field; sympifying those literals corrected the implementation, with no change to the mathematical claims. Only the completed run in the accompanying derivation and verification record is active evidence.

The checks of preserved earlier text, equation labels, file hashes and the persistent artifact identity are integrity checks, separate from the mathematical predicate count. The successful persistent replacement and exact-target reopening are recorded as \texttt{WRT-G10} after the write, before reporting the artifact as saved.




\section{Geometric functionals, bundles, and Dolbeault calculations}
\label{p1c:sec:functionals}

\subsection{Cartan transgression, the cylinder functional and its fluctuations}
\label{p1b:app:g11}

Use \(T_a=-i\sigma_a/2\), \(\tr(T_aT_b)=-\delta_{ab}/2\), and
\(\Theta=g^{-1}\dd g=2T_ae^a\), where the unit round coframe satisfies
\(\dd e^1=-2e^2\wedge e^3\) cyclically. Orient \(S^3\) by
\(-\int\tr\Theta^3/(24\pi^2)=1\), and the cylinder by
\(\dd\tau\wedge\vol_{S^3}\). Here \(g_{\rm YM}>0\) is a parameter
of a specified functional.

\begin{proposition}[Chern--Simons polynomial and endpoint charge]
\label{p1b:prop:g11-transgression}
For a real \(C^1\) profile \(f\) and \(A_f=f\Theta\), define
\[
 N_{\rm CS}(A)=\frac1{8\pi^2}\int_{S^3}
       \tr(A\wedge\dd A+\tfrac23A^3).
\]
Then
\[
 N_{\rm CS}(A_f)=N(f):=3f^2-2f^3,\qquad
 \frac1{8\pi^2}\int_{[\tau_-,\tau_+]\times S^3}\tr F_f^2
      =N(f(\tau_+))-N(f(\tau_-)).
\]
Complementing the profile sends \(N(f)\) to \(1-N(f)\) and reverses
this endpoint charge. A smooth \(0\)-to-\(1\) profile constant near its
ends glues to a smooth connection of charge one on the clutched bundle
over \(S^1\times S^3\).
\end{proposition}
\begin{proof}
The spatial Maurer--Cartan equation gives
\(\tr(A\dd A+\frac23A^3)=(-f^2+\frac23f^3)\tr\Theta^3\).
Moreover
\[
 F_f=\dot f\,\dd\tau\wedge\Theta+f(f-1)\Theta^2,\qquad
 \tr F_f^2=2\dot f f(f-1)\dd\tau\wedge\tr\Theta^3.
\]
The spatial fourth-degree term vanishes by dimension, and
\(N'(f)=6f(1-f)\). These identities prove both integrals and the
complement rule. The two endpoint potentials satisfy
\(A_1=g^{-1}A_0g+g^{-1}\dd g\). Constancy near the cuts makes all normal
jets match under this transition, which defines the claimed bundle
and connection.
\end{proof}
The displayed value of \(N_{\rm CS}\) uses this trivialization; large
gauge changes shift it by an integer. A periodic profile in one global
trivialization has zero net charge. At the constant midpoint
\(f=1/2\), \(N_{\rm CS}=1/2\) and the four-dimensional charge is zero.
The glued connection need not be a critical point of a functional.

\begin{theorem}[Reduced variational problem and unit-charge minimum]
\label{p1b:thm:g11-functional}
On \(\R\times S^3\) choose
\(g_{\rm cyl}=R^2(\dd\tau^2+g_{S^3})\), \(R>0\), and
\[
 S_{\rm YM}(A)=-g_{\rm YM}^{-2}\int\tr(F_A\wedge *F_A).
\]
For the homogeneous family,
\[
 S_{\rm YM}(f)=\frac{12\pi^2}{g_{\rm YM}^2}
       \int_\R\{\dot f^{\,2}+4f^2(1-f)^2\}\dd\tau,\qquad
 \ddot f=4f(1-f)(1-2f).
\]
Among finite-action profiles with limits \(f(-\infty)=0\) and
\(f(+\infty)=1\), the minimum is \(8\pi^2/g_{\rm YM}^2\).
Every minimizer is a translate of
\[
 f_*(\tau)=\frac{1+\tanh(\tau-\tau_0)}2.
\]
Its connection is anti-self-dual in the stated orientation and therefore
solves the full Yang--Mills equation on the cylinder.
\end{theorem}
\begin{proof}
The electric curvature coefficients are \(2\dot fT_a\), and the magnetic
coefficients on cyclic spatial two-planes are \(4f(f-1)T_a\).
The trace norm is thus \(6\{\dot f^2+4f^2(1-f)^2\}\).
Integrating the unit sphere gives \(2\pi^2\); the factor \(R\) cancels
because the Hodge star on two-forms is conformally invariant in dimension
four. Compactly supported variation yields the stated equation.
The exact square completion is
\[
 \dot f^2+4f^2(1-f)^2
  =\{\dot f-2f(1-f)\}^2+\frac23\,\frac{\dd}{\dd\tau}N(f).
\]
This proves the bound without imposing \(0\le f\le1\) in advance.
Equality forces the first-order equation
\(\dot f=2f(1-f)\). Separation, together with the two limits, gives the
displayed translates and no others. Their electric coefficient is the
negative of the magnetic one, proving anti-self-duality. The Bianchi
identity then implies \(\dd_A*F_A=0\).
\end{proof}

\begin{proposition}[Complete homogeneous fluctuation spectrum]
\label{p1b:prop:g11-fluctuations}
For \(f=f_*+\varepsilon h\), the coefficient of \(\varepsilon^2\)
in the action is
\[
 \frac{12\pi^2}{g_{\rm YM}^2}\langle h,\mathcal Lh\rangle,\qquad
 \mathcal L=-\partial_z^2+4-6\operatorname{sech}^2z,\quad
 z=\tau-\tau_0.
\]
On \(L^2(\R,\dd z)\), with domain \(H^2(\R)\), the point spectrum is
\(0,3\), both simple, with normalized eigenfunctions
\[
 \psi_0=\frac{\sqrt3}{2}\operatorname{sech}^2z,\qquad
 \psi_1=\sqrt{\frac32}\operatorname{sech}z\tanh z.
\]
The remaining spectrum is absolutely continuous on \([4,\infty)\).
\end{proposition}
\begin{proof}
Half the second derivative of \(4f^2(1-f)^2\) at \(f_*\) is
\(4-6\operatorname{sech}^2z\). Set
\(Q_j=\partial_z+j\tanh z\). Direct multiplication gives
\[
 \mathcal L=Q_2^*Q_2,\qquad
 Q_2Q_2^*=Q_1^*Q_1+3,\qquad Q_1Q_1^*+3=-\partial_z^2+4.
\]
The last operator has only its free absolutely continuous spectrum.
The two intervening normalizable kernels insert precisely the stated
eigenvalues; the intertwining maps preserve the positive spectral parts.
The first kernel is \(\operatorname{sech}^2z\); pulling back
\(\operatorname{sech}z\) from the partner gives
\(\operatorname{sech}z\tanh z\). Substitution \(u=\tanh z\) gives squared
norms \(4/3\) and \(2/3\), proving the constants. The free threshold has
no square-integrable state.
\end{proof}
The zero mode is the translation derivative of \(f_*\). For the constant
midpoint the corresponding differential expression is
\(-\partial_\tau^2-2\); on a periodic parameter interval it has a
negative constant mode. In proper-length coordinates the displayed
fluctuation eigenvalues scale by \(R^{-2}\). With \(r=e^\tau\), the kink
is \(r^2/(r^2+r_0^2)\). These are statements about homogeneous profile
variations; the operator is not the complete gauge-field Hessian.

\subsubsection{Full gauge-field stability of the cylinder kink}
\label{p1b:app:full-cylinder-kink}

Let \(X=\mathbb R\times S^3\),
\(g_c=R^2(\dd\tau^2+g_{S^3})\), and let
\(A_*=f_*\Theta\) be the charge-one anti-self-dual connection of
Theorem~\ref{p1b:thm:g11-functional}.  Use the positive trace inner
product \(\langle u,v\rangle=-\int\operatorname{tr}(u\wedge *v)\).
Write \(d_A^+=\tfrac12(1+*)d_A\).

\begin{theorem}[The full second variation and its Friedrichs domain]
\label{p1b:thm:full-kink-hessian}
For every compactly supported adjoint-valued one-form \(a\),
the coefficient of \(\varepsilon^2\) in
\(S_{\rm YM}(A_*+\varepsilon a)\) is
\[
 \frac{2}{g_{\rm YM}^2}\|d_{A_*}^+a\|_{L^2(g_c)}^2.
\]
Adding the background-gauge term gives the nonnegative form
\[
 q_c(a)=2\|d_{A_*}^+a\|_2^2+\|d_{A_*}^{*,c}a\|_2^2,
 \qquad
 L_c=2(d_{A_*}^+)^{*,c}d_{A_*}^+
                    +d_{A_*}d_{A_*}^{*,c}.
\]
On \(L^2(X,\Lambda^1\otimes\operatorname{ad}E;g_c)\), its
Friedrichs form domain is \(H^1\), and its operator domain is
\(H^2\), defined using a fixed smooth bounded-geometry trivialization
at each end.  In particular
\[
 L_c\ge0,\qquad
 \ker L_c=\{a\in H^1:d_{A_*}^+a=0,
                          \ d_{A_*}^{*,c}a=0\}.
\]
There are no negative gauge-field modes, including inhomogeneous
ones.  This statement concerns the complete one-form perturbation
space and not merely the homogeneous scalar profile.
\end{theorem}
\begin{proof}
For the fixed charge and the stated anti-self-dual orientation,
\[
 S_{\rm YM}(A)=\frac{8\pi^2}{g_{\rm YM}^2}
                    +\frac2{g_{\rm YM}^2}\|F_A^+\|_2^2.
\]
Compact support preserves the charge.  Since \(F_{A_*}^+=0\) and
\(F_{A_*+\varepsilon a}^+=\varepsilon d_{A_*}^+a
+\varepsilon^2(a\wedge a)^+\), expansion proves the asserted
coefficient.  The ordinary Hessian is twice this coefficient.
The gauge term gives the displayed factorization of \(L_c\).
Its principal symbol is \(|\xi|^2\operatorname{id}_{\Lambda^1}\):
indeed \(2|(\xi\wedge a)^+|^2+|\iota_\xi a|^2
=|\xi|^2|a|^2\).  The cylinder has bounded geometry and the kink
connection and all its derivatives are bounded in its two smooth
end gauges.  Standard global elliptic estimates therefore identify
the closure of the form norm with \(H^1\), and the self-adjoint
Friedrichs operator domain with \(H^2\).  Nonnegativity and the
kernel characterization follow directly from the closed sum of
squares.  In a Coulomb slice the only remaining quadratic term is
the first square.  Infinitesimal compactly supported gauges are
annihilated by the ungauged Hessian because
\(d_A^+d_A\epsilon=[F_A^+,\epsilon]=0\).
\end{proof}

\begin{theorem}[A complete five-dimensional kernel in the compactified category]
\label{p1b:thm:compactified-kink-kernel}
Choose the following precise admissibility and gauge convention.
Put \(r=e^\tau\), fill in \(r=0,\infty\), and require the
connection and variations to extend smoothly, in their regular bundle
charts, to the charge-one bundle on \(S^4\).  Quotient by all smooth
\(SU(2)\) gauge transformations on that compact bundle; no frame at
either pole is fixed.  Impose Coulomb gauge for
\[
 g_s=\operatorname{sech}^2\tau\,g_c,
\]
the round metric of radius \(R\) on the compactification.  Its
gauge-fixed operator has a five-dimensional real kernel, no negative
eigenvalues, and a strictly positive first eigenvalue on the orthogonal
complement of its kernel.  Locally the moduli space is smooth of
dimension five.  The homogeneous translation mode is one of these
five modes; the other four move the Euclidean instanton centre.
\end{theorem}
\begin{proof}
The profile becomes
\(f_*=r^2/(r^2+\rho^2)\), \(\rho=e^{\tau_0}>0\).  Thus
\(A_*\) is the centred BPST connection in regular gauge.  Its
apparent singularity at infinity is removed by the original
degree-one clutch transformation.  With
\(\Omega=\operatorname{sech}\tau\), define
\[
 C_Aa=d_A^{*,c}a-2\iota_{\nabla_c\log\Omega}a.
\]
The exact conformal identities in four dimensions are
\[
 d_A^{*,s}a=\Omega^{-2}C_Aa,\qquad
 \|d_A^+a\|_{L^2(g_s)}=\|d_A^+a\|_{L^2(g_c)},\qquad
 \|d_A^{*,s}a\|_{L^2(g_s)}=\|C_Aa\|_{L^2(g_c)}.
\]
Hence the compactified gauge-fixed form, written on the cylinder, is
\[
 q_s(a)=2\|d_A^+a\|_{2,c}^2+\|C_Aa\|_{2,c}^2,
 \qquad
 \|a\|_{2,s}^2=\int_X\Omega^2|a|_c^2\,\mathrm{vol}_c.
\]
Its exact operator in this weighted Hilbert space is
\[
 \widehat L_s=\Omega^{-2}
       \{2(d_A^+)^{*,c}d_A^++C_A^{*,c}C_A\}.
\]
Its form and operator domains are, by definition, the pullbacks of
\(H^1(S^4)\) and \(H^2(S^4)\), with the stated smooth bundle charts
at the poles.  These formulas retain the conformal weight; the
compactified operator is not silently identified with \(L_c\).

The ASD deformation complex is
\[
 0\longrightarrow\Omega^0(\operatorname{ad}E)
 \xrightarrow{d_A}\Omega^1(\operatorname{ad}E)
 \xrightarrow{d_A^+}\Omega^{2,+}(\operatorname{ad}E)
 \longrightarrow0.
\]
Its degree-zero cohomology is zero: the BPST curvature spans
\(\mathfrak{su}(2)\) at a point, so a parallel adjoint section
commutes with all of \(\mathfrak{su}(2)\) and vanishes.  Its degree-two
cohomology is zero by the Weitzenbock identity on self-dual two-forms,
\[
 2d_A^+(d_A^+)^*=
 \nabla_A^*\nabla_A+\frac{\operatorname{Scal}_{g_s}}3-2W^+
 =\nabla_A^*\nabla_A+\frac4{R^2}.
\]
Here \(F_A^+=0\), \(W^+=0\), and
\(\operatorname{Scal}_{g_s}=12/R^2\).  Thus the deformation index
is the kernel dimension.  The Atiyah--Hitchin--Singer index gives
\[
 \dim H_A^1=8k-3(1-b_1+b_2^+)=8-3=5.
\]
Surjectivity of the linearized ASD equation and the slice implicit
function theorem give the asserted smooth local moduli space.
Compact elliptic spectral theory gives compact resolvent and a
positive gap after the kernel is removed.

The explicit translated and rescaled BPST family has four centre
parameters \(b\in\mathbb R^4\) and one scale \(\rho>0\).  Its
derivatives are independent modulo gauge: its gauge-invariant Euclidean
curvature density is a nonzero constant times
\[
 \frac{\rho^4}{(|x-b|^2+\rho^2)^4}.
\]
The four centre derivatives are independent odd functions in the four
coordinates around \(b\), while the nonzero scale derivative is
radial.  These five tangent directions therefore exhaust the kernel
after Coulomb projection.  The change of \(\tau_0=\log\rho\) is
exactly the earlier homogeneous translation zero mode.  The index and
deformation theory used here are the primary results of Atiyah,
Hitchin and Singer, \emph{Deformations of instantons}, PNAS 74 (1977),
2662--2663, doi:10.1073/pnas.74.7.2662, and \emph{Self-duality in
four-dimensional Riemannian geometry}, Proc. Roy. Soc. Lond. A 362
(1978), 425--461, doi:10.1098/rspa.1978.0143.
\end{proof}

\subparagraph{End framing and spectrum scope.}
The five-dimensional assertion uses all smooth compact gauge
transformations.  If a fibre frame at one specified pole is retained,
the local moduli space gains the three-dimensional orientation fibre
\(SU(2)/\{\pm1\}\), and has dimension eight.  These three frame
coordinates are not additional zero eigenforms of the unframed
compact Coulomb operator.  Fixing independent gauges at both cylinder
ends defines another problem.  Neither its dimension nor an entire
continuous spectrum follows from the homogeneous
\(0,3,[4,\infty)\) spectrum.  The theorem proves full-field stability
on the complete cylinder and supplies an exact kernel and local
moduli theorem in a separately specified compactified category.



\subsection{Full midpoint gauge Hessian and gauge subtraction}
\label{p1b:app:native-weak-morse-bridge}
This technical section supplies the representation-theoretic derivation used by
Theorem~\ref{p1b:thm:native-midpoint-gauge-morse}.  It concerns the classical
pure \(SU(2)\) Yang--Mills Hessian on the round sphere with the stated invariant metric and connection.

Use the unit-radius left-invariant frame and restore \(R^{-2}\) at the end.
At \(A_{\rm B}=\Theta/2\), let \(L\) be the orbital generators, \(S\) the
spin-one generators on one-forms, and \(T\) the adjoint color generators.
The midpoint covariant derivative on adjoint one-forms may be written
\[
 \nabla^{A_{\rm B}}=-i(2L+S+T).
\]
The background-gauge-fixed Yang--Mills Hessian is the standard curvature
Weitzenbock operator
\[
 H_{\rm gf}=(\nabla^{A_{\rm B}})^*\nabla^{A_{\rm B}}+2+2S\cdot T.
\]
Putting \(J=L+S+T\), using \(S^2=2\) and \(T^2=t(t+1)\), direct expansion gives
\[
 H_{\rm gf}=2J^2+2L^2-T^2,
\]
and hence
\[
 R^2 H_{\rm gf}=2J(J+1)+2j(j+1)-t(t+1).
\]
For the weak adjoint, \(t=1\), this is
\eqref{p1b:eq:native-midpoint-gauge-hessian}.

On adjoint scalars there is no one-form spin.  Coupling \(L+T=K\) gives
\[
 R^2d_{A_{\rm B}}^*d_{A_{\rm B}}
 =2j(j+1)+2k(k+1)-t(t+1),
\]
which is \eqref{p1b:eq:native-midpoint-ghost} at \(t=1\).  For each positive
ghost eigenvalue, the map \(d_{A_{\rm B}}\) injects the corresponding scalar
representation into the vector representation with the same total spin
\(J=k\); that copy is the infinitesimal gauge orbit and is removed in the quotient by infinitesimal gauge transformations.

For \(j=0\),
\[
 V_0\otimes V_1\otimes V_1=V_0\oplus V_1\oplus V_2,
\]
with vector eigenvalues \(-2,2,10\).  The ghost space is \(V_1\), with
eigenvalue \(2\), so precisely the middle copy is gauge.  This gives the
homogeneous split in \eqref{p1b:eq:native-homogeneous-morse-split}.

For \(j=\frac12\),
\[
 V_{1/2}\otimes V_1\otimes V_1
 =2V_{1/2}\oplus2V_{3/2}\oplus V_{5/2}.
\]
The corresponding vector eigenvalues are \(1,7,17\), while the ghost
representations \(V_{1/2}\oplus V_{3/2}\) have positive eigenvalues
\(1,7\).  One copy of each is removed, leaving a quotient
\(V_{1/2}\) at \(1/R^2\).  For \(j\ge1\),
\[
 2J(J+1)+2j(j+1)-2\ge2j(j+1)-2\ge2,
\]
so no further nonpositive vector eigenvalue occurs even before gauge
subtraction.  This proves the full Morse count and the sharp positive gap.

For completeness, the left-invariant matrix description identifies the
\(J=0,1,2\) summands with the scalar, antisymmetric, and symmetric-tracefree
parts of \(b_a{}^c\), respectively.  The scalar ghost derivative maps onto
the antisymmetric part, so the representation-theoretic subtraction agrees
with the explicit background-gauge projection.



\subsection{Compact-double conventions, exchange and reduction}
\label{p1b:app:g12}
The unfused double uses the \(K\times K\) action
\[
 (g_1,g_2)\cdot(a,b)=(g_1ag_2^{-1},g_2bg_1^{-1}),
 \qquad(\mu_1,\mu_2)=(ab,a^{-1}b^{-1}),
\]
and the first two terms of \(\omega_D\), denoted \(\omega_0\).
For multiplication \(m\) and inversion \(j\), Maurer--Cartan gives
\[
 m^*\theta^L=\Ad_{b^{-1}}a^*\theta^L+b^*\theta^L,\qquad
 j^*\theta^L=-\theta^R,\qquad j^*\chi=-\chi,
\]
\[
 m^*\chi=\operatorname{pr}_1^*\chi+
          \operatorname{pr}_2^*\chi
          -\tfrac12\dd\langle\operatorname{pr}_1^*\theta^L,
                                 \operatorname{pr}_2^*\theta^R\rangle .
\]
Inserting the two moment maps proves
\(\dd\omega_0=-\mu_1^*\chi-\mu_2^*\chi\); the extra fusion term
then gives \(\dd\omega_D=-(\mu_1\mu_2)^*\chi\).
Contraction and minimal degeneracy are the remaining assertions of
the double and fusion theorems \cite{AlekseevMalkinMeinrenken1998}.

\subparagraph{Exchange signs.}
Set \(u=ab\), \(v=a^{-1}b^{-1}\).
Under exchange, \(u\) becomes \(v^{-1}\) and \(v\) becomes \(u^{-1}\).
Thus \(u^*\theta^L\) becomes \(-v^*\theta^R\) and
\(v^*\theta^R\) becomes \(-u^*\theta^L\).
Symmetry of the inner product and antisymmetry of the wedge product
make this pairing negative. The other two pairings exchange with a
minus sign. Also \(ba b^{-1}a^{-1}=(ab a^{-1}b^{-1})^{-1}\).
This proves both exchange formulas on the whole product.

\subparagraph{A matrix realization and the exceptional kernel.}
Identify the vector adjoint representation with \(\R^3\), and put
\(A=\Ad_a\), \(B_0=\Ad_b\), \(L_a=(I\ \ 0)\), \(L_b=(0\ \ I)\).
For matrix-valued one-forms set
\(\mathcal W(U,V)=U^TV-V^TU\). The matrix of the vector part of the
form in left trivialization, for an orthonormal invariant pairing, is
\[
 W=\tfrac12\{\mathcal W(L_a,B_0L_b)+\mathcal W(AL_a,L_b)
       +\mathcal W(B_0^TL_a+L_b,-L_a-A^TL_b)\}.
\]
This gives an explicit coordinate formula for evaluation or rank checks.
At \(a=b=1\) it is the standard off-diagonal alternating pairing of
rank six. At \(a_v=i,\ b_v=(1+j)/\sqrt2\), the commutator is a
pure imaginary unit quaternion and the vector rank is four.
At \(a_v=i,\ b_v=j\), the commutator is \(-1\) and the vector rank is
six. The scalar part always has rank two.

\subparagraph{Normalization relative to the reciprocal two-form.}
On the simple ideal, choose
\(\langle X,Y\rangle=-\Tr(XY)/(4\pi^2)\). Then
\(\chi=-\Tr(\theta^L)^3/(24\pi^2)=\nu\).
At the identity pair the vector part of \(\omega_D\) is
\(-\Tr(X\wedge Y)/(4\pi^2)=2\beta\), with \(\beta\) as in
Theorem~\ref{p1b:thm:transgression}. Even their identity values
therefore require the factor two at this common three-form
normalization. The central pairing can be normalized independently.

\subparagraph{Local quotient structure.}
The Weyl involution on \((\alpha,\beta)\) fixes exactly
\(\{0,\pi\}^2\). Near each fixed point the quotient is
\((\R^2/\{\pm1\})\times\T^2\). The invariant two-form on the covering
torus descends to the four-dimensional smooth stratum and restricts
to a symplectic form on every two-dimensional singular stratum.
The central \(U(1)\) acts trivially throughout. At the other central
value \(\mu_v=-1\), the effective \(SO(3)\) action on the orthogonal
quaternion pair is free and transitive. These descriptions establish
the claimed quotient spaces without imposing an incorrect freeness
assumption at \(\mu=1\).



\subsection{Theta-function eigenstates and magnetic translations}
\label{p1b:app:g13}

Use \(s=\theta_t/(2\pi)\), \(t=\theta_E/(2\pi)\), each of period one.
Every oriented flat metric has the form
\[
 g=\frac{A}{\tau_2}|\dd s+\tau\dd t|^2,\qquad
 \tau=\tau_1+i\tau_2,\quad \tau_2>0.
\]
For the metric \(G\), \(\tau_1=G_{tE}/G_{tt}\) and
\(\tau_2=\sqrt{\det G}/G_{tt}\). Its volume is \(A\,\dd s\,\dd t\).
In this technical section initially \(q>0\).

\begin{lemma}[Connection, boundary conditions and normalized ground states]
\label{p1b:lem:g13-theta}
Represent sections on \(\R^2\) by functions satisfying
\[
 \psi(s+1,t)=e^{2\pi iqt}\psi(s,t),\qquad
 \psi(s,t+1)=\psi(s,t),
\]
with connection
\(\nabla_0=\dd-2\pi iq\,s\,\dd t\). Define
\[
 N_q=\frac{(2q\tau_2)^{1/4}}{A^{1/2}|\tau|^{1/2}},\qquad
 \psi_j(s,t)=N_q\sum_{n\in\Z}
   \exp\left\{2\pi i(j+qn)t-
       \frac{\pi iq}{\tau}(s-n-j/q)^2\right\},
 \quad 0\le j<q.
\]
These functions are an orthonormal basis of \(\ker a\).
The series and all derivatives converge on compact sets. Equivalently,
with \(\vartheta(z\mid\kappa)=\sum_{n\in\Z}
e^{\pi i\kappa n^2+2\pi inz}\),
\[
 \psi_j=N_q e^{\,2\pi ijt-\pi iq(s-j/q)^2/\tau}
       \vartheta\left(qt+\frac{q(s-j/q)}{\tau}
                         \,\middle|\,-\frac q\tau\right).
\]
\end{lemma}
\begin{proof}
The connection respects the stated boundary conditions and has
\(F_{\nabla_0}=-2\pi iq\,\dd s\wedge\dd t\), so its degree is \(q\).
An oriented orthonormal frame gives
\[
 \nabla_1=\sqrt{\frac{\tau_2}{A}}\partial_s,\qquad
 \nabla_2=\frac{\partial_t-\tau_1\partial_s-2\pi iqs}
                    {\sqrt{A\tau_2}} .
\]
Thus the equation \(a\psi=0\) is
\[
 (-i\tau\partial_s+i\partial_t+2\pi qs)\psi=0.
\]
For the Fourier coefficient \(e^{2\pi imt}f_m(s)\), it becomes
\(f_m'=2\pi i(m-qs)f_m/\tau\); its solutions are the displayed
Gaussians. The \(s\)-boundary condition identifies coefficients with
indices differing by \(q\), leaving precisely \(q\) constants.
Since \(\Re(i/\tau)=\tau_2/|\tau|^2>0\), these series are smooth.
Integration in \(t\) removes cross terms, and different residues
\(j\pmod q\) are orthogonal. For each residue, unfolding the \(s\)
integral gives
\[
 \|\psi_j\|^2
 =N_q^2 A\int_\R e^{-2\pi q\tau_2 u^2/|\tau|^2}\dd u
 =N_q^2\,\frac{A|\tau|}{\sqrt{2q\tau_2}}=1.
\]
Expanding the square in the series gives the theta expression;
\(\Im(-q/\tau)>0\) ensures its convergence.
\end{proof}

\begin{proposition}[Complete ladder and connection transport]
\label{p1b:prop:g13-ladder}
The functions
\[
 \psi_{j,r}=\frac{(a^*)^r}{\sqrt{r!}}\psi_j,\qquad
 0\le j<q,\quad r\ge0,
\]
form a complete orthonormal eigenbasis for \(\Delta_{\nabla_0}\).
Every other connection of this curvature is gauge equivalent to
\[
 \nabla_{\alpha,\beta}=\nabla_0+
                  2\pi i(\alpha\,\dd s+\beta\,\dd t),
 \qquad(\alpha,\beta)\in\R^2/\Z^2,
\]
and is unitarily equivalent for the Bochner and Dirac operators.
\end{proposition}
\begin{proof}
The commutator normalizes the ladder and makes different levels
orthogonal. Compact ellipticity supplies a discrete complete
eigensystem for the self-adjoint Bochner operator. If
\(\Delta\phi=\lambda\phi\), then
\(\|a\phi\|^2=(\lambda-b)\|\phi\|^2/(2b)\) and every nonzero lowering
step reduces the eigenvalue by \(2b\). Positivity forces the process
to stop exactly at \(\lambda=b\); otherwise a last nonzero step would
have eigenvalue below \(b\). Raising recovers the whole eigenspace.
The ground-state calculation therefore proves completeness and all
multiplicities.

Two fixed-curvature unitary connections differ by an imaginary closed
one-form. Exact terms are gauge, and integral-period forms are large
gauge changes. This yields the displayed torsor.
For real \(h,k\), the unitary operator
\[
 (T_{h,k}\psi)(s,t)=e^{-2\pi iqks}\psi(s+h,t+k)
\]
respects the same boundary conditions and satisfies
\[
 T_{h,k}^{-1}\nabla_0T_{h,k}
   =\nabla_0+2\pi i(-qk\,\dd s+qh\,\dd t).
\]
Taking \(h=\beta/q,\ k=-\alpha/q\) proves the equivalence.
It also transports the eigenbasis by \(T_{h,k}^{-1}\).
\end{proof}

\begin{proposition}[Finite magnetic-translation representation]
\label{p1b:prop:g13-magnetic}
For \(\nabla_{\alpha,\beta}\), the unitary operators
\[
 (U\psi)(s,t)=e^{-2\pi it}\psi(s+1/q,t),\qquad
 (V\psi)(s,t)=\psi(s,t+1/q)
\]
commute with both covariant derivatives and satisfy
\[
 U^q=V^q=I,\qquad UV=e^{2\pi i/q}VU.
\]
Every Landau eigenspace is the irreducible \(q\)-dimensional
representation with this central character. For the reference
connection,
\[
 U\psi_{j,r}=\psi_{j-1,r},\qquad
 V\psi_{j,r}=e^{2\pi ij/q}\psi_{j,r},
\]
with residues taken modulo \(q\).
\end{proposition}
\begin{proof}
Substitution verifies automorphy, the two derivative commutators and
the powers. Changing the order of the translations shifts the
multiplier \(e^{-2\pi it}\) by \(e^{-2\pi i/q}\), proving their
commutation relation. The series gives the shift and clock matrices
on the ground states; commutation with \(a^*\) gives every level.
The distinct clock eigenvalues and cyclic shift prove irreducibility.
For any flat part, these operators still commute with the derivatives;
the same commutation relation on a \(q\)-dimensional level forces the
same irreducible representation. For \(q=1\) it is one-dimensional.
\end{proof}

\subparagraph{Negative degree, spin and product checks.}
Complex conjugation takes the degree \(q\) data to degree \(-q\);
it preserves the Bochner spectrum and reverses the Dirac chirality.
With \(\Gamma=\diag(1,-1)\) and
\(D=\sigma_1\pi_1+\sigma_2\pi_2\), the index convention is
\(\operatorname{ind}D^+=\dim\ker(D|_{\Gamma=+1})
-\dim\ker(D|_{\Gamma=-1})=q\).
The torus heat trace for its Dirac square is
\[
 \Tr e^{-tD_{\T^2}^2}=|q|\coth(bt).
\]
Consequently the full spin-square trace is
\[
 8|q|\coth(bt)
 \left(\sum_{\ell\ge0}(\ell+1)(\ell+2)
             e^{-t(\ell+3/2)^2/R_x^2}\right)
 \left(\sum_{\ell\ge0}(\ell+1)(\ell+2)
             e^{-t(\ell+3/2)^2/R_p^2}\right).
\]
This checks the odd-factor doubling and the separate zero-level
multiplicity. The scalar-line product index zero is distinct from the
index of the spherical virtual Poincar\'e class in
Theorem~\ref{p1b:thm:spherical-pairing}.



\subsection{Closed-torsion signs and the generalized Kähler bracket}
\label{p1b:app:g14}
This technical section uses \(gX\) for the one-form \(g(X,\,\cdot\,)\).
The bracket is the Dorfman bracket with the sign already fixed in
Theorem~\ref{p1b:thm:dirac-graph}:
\[
 [X+\xi,Y+\zeta]_{H_g}
 =[X,Y]+\mathcal L_X\zeta-\iota_Y\dd\xi+\iota_Y\iota_XH_g.
\]
All generalized-complex assertions can be checked with this bracket;
on an isotropic subbundle it equals its skew-symmetrization.

\subparagraph{Hermitian sign calculation.}
For \(\omega(X,Y)=g(JX,Y)\) and \(d^c=i(\bar\partial-\partial)\),
\[
 d^c\omega(X,Y,Z)=-\dd\omega(JX,JY,JZ).
\]
The invariant-form differential, the Nijenhuis identity and ad-invariance
give, on left-invariant fields,
\(\dd\omega_L(J_LX,J_LY,J_LZ)=-g([X,Y],Z)\).
Thus \(d^c\omega_L=H_g\), while the reversed bracket of right-invariant
fields gives \(d^c\omega_R=-H_g\).
Consequently \(c_L=-d^c\omega_L=-H_g\) and
\(c_R=-d^c\omega_R=H_g\).
For example, on a unit \(SU(2)\times U(1)\) block with
\(JX_1=X_2,\ JX_3=Z\),
\[
 \omega=e^1\wedge e^2+e^3\wedge z,\qquad
 \dd\omega=-2e^1\wedge e^2\wedge z,\qquad
 d^c\omega=2e^1\wedge e^2\wedge e^3.
\]
This example fixes the sign independently of naming conventions for
\(d^c\) or the Bismut connection.

\subparagraph{Closure, periods and exchange.}
The invariant differential of \(H_g\) is a sum of contractions
\(g([X,Y],[Z,W])\), which cancels by Jacobi. The simple ideals are
orthogonal to each other and to the center, so no mixed terms occur.
Its de Rham class is
\[
 [H_g]=4\pi^2(R_x^2u_x+R_p^2u_p).
\]
For an arbitrary bi-invariant metric,
\(B^*H_g=H_{B^*g}\). In particular \(B^*H_g=H_g\) for an
exchange-invariant metric; this is a different parity from that of
\(\Omega\). Inversion has \(j^*H_g=-H_g\) and exchanges the left and
right parallel connections. Flatness of these connections does not
assert flatness of the Levi-Civita connection.

\subparagraph{Direct proof of generalized integrability.}
Let \(X^R\) and \(Y^L\) denote the right- and left-invariant fields
generated by \(X,Y\in\mathfrak g\). Define sections
\[
 r(X)=X^R+gX^R\in C_+,\qquad
 l(Y)=Y^L-gY^L\in C_-.
\]
The group brackets, metric invariance and the displayed Dorfman
formula give
\[
\begin{gathered}
 [r(X),r(Y)]_{H_g}=r(-[X,Y]),\qquad
 [l(X),l(Y)]_{H_g}=l([X,Y]),\\
 [r(X),l(Y)]_{H_g}=[l(Y),r(X)]_{H_g}=0.
\end{gathered}
\]
For the first identity the untwisted one-form contribution is twice
the metric dual of the right-field bracket, while the \(H_g\) term
subtracts one copy; the left identity has the opposite graph sign.
For the cross identity the fields commute, the Lie-derivative term
vanishes, and the remaining exterior derivative cancels the twist.
These identities can also be verified on any third left or right
field, since such fields span the tangent space.

The \(+i\) eigenspace of \(\mathcal J_1\) is the sum of the \(+i\)
eigenspaces of \(J_R\) in \(C_+\) and \(J_L\) in \(C_-\).
For \(\mathcal J_2\) the second summand is the \(-i\) eigenspace.
Integrability of the two invariant ordinary complex structures makes
each constant subspace a complex Lie subalgebra. The bracket identities
therefore close both generalized eigenspaces. They are isotropic; the
Leibniz identities extend closure from these frames to arbitrary smooth
sections. This proves the generalized integrability assertion without
an implicit sign change of \(H_g\). No semisimplicity assumption is
needed: the same argument includes the two central directions.



\subsection{The unstable bundle generator and its stabilization}
\label{p1b:app:g15}
We use based homotopy classes; all the targets below are simply
connected, so passing to unbased bundle classifications introduces no
additional action.

\subparagraph{Dimension and suspension checks.}
Quaternionic projective space has one cell in dimensions
\(0,4,8,\ldots\). The inclusion \(S^4=\mathbb H P^1\to BSU(2)\)
is therefore \(7\)-connected and induces a bijection on maps from the
six-dimensional \(P_v\), including their homotopies.
Freudenthal stabilization is bijective in this dimension, since
\(S^4\to QS^4\) is \(7\)-connected. The stable splitting
\[
 \Sigma(S^3\times S^3)\simeq S^4\vee S^4\vee S^7
\]
leaves only the top-sphere contribution to stable cohomotopy in degree
four: both degree-three sphere contributions vanish. The remaining
group is the stable \(2\)-stem, already equal to
\(\pi_6(S^4)=\Z/2\). The suspension range, splitting and generator
are given in \cite[Corollary 4.24, Proposition 4I.1 and Section 4J]
{hatcher2002algebraic}; their specialization here proves the bundle classification.

\subparagraph{An explicit classifying map.}
Let \(q_\wedge:P_v\to S^3\wedge S^3=S^6\) collapse the two coordinate
spheres, and let \(\eta_n:S^{n+1}\to S^n\) be the suspended Hopf map.
For the tautological quaternionic line \(\mathcal H\) on
\(\mathbb H P^1=S^4\), set
\[
 E=(\eta_4\circ\eta_5\circ q_\wedge)^*\mathcal H.
\]
The composite \(\eta_4\eta_5\) is the nonzero element of
\(\pi_6(S^4)\), so this is the nontrivial rank-two complex bundle.
Its restrictions to both coordinate \(S^3\)'s are trivial; the
obstruction is carried by the top attaching sector. The constant
classifying map gives the other bundle and is a negative control
with the same Chern classes.

\subparagraph{Stable cancellation and the pullback sector.}
The reduced stable class vanishes because its Chern character does
and \(\widetilde K^0(P_v)\) has no torsion.
Rank-three stabilization is bijective in dimension six, so adding
one line actually trivializes \(E\), rather than merely making its
character vanish. If \(\pi^*E\) were trivial, restriction along
\(P_v\times\{t_0\}\) would trivialize \(E\), a contradiction.
Both pulled-back bundles have determinant one. The nontrivial one
admits no flat \(SU(2)\) connection, by
Theorem~\ref{p1b:thm:equivariant-flat}.

Finally \(q_\wedge\circ B\) is the smash-factor transposition followed
by \(q_\wedge\); its degree on \(S^6\) is
\((-1)^{3\cdot3}=-1\). Negation fixes the element of order two,
so \(B^*E\cong E\), and hence \(B^*\pi^*E\cong\pi^*E\).
An isomorphism of these bundle classes is weaker than a chosen
isomorphism satisfying the cocycle equation for an exchange action.

\subsubsection{Equivariant Hopf construction for the unstable G15 bundle}
\label{p1b:app:g21-exotic-construction}

Let \(\sigma\) denote the real sign representation of \(C_2\), and write
\(S^V\) for the one-point compactification of a real representation \(V\).
The smash collapse
\[
 q_\wedge:P_v=S_x^3\times S_p^3\longrightarrow
 S_x^3\wedge S_p^3
\]
is equivariant for factor exchange.  On the top cell the swap
representation decomposes into diagonal and anti-diagonal parts, so
\[
 S_x^3\wedge S_p^3\cong S^{3+3\sigma}.
\]

\subparagraph{Equivariant suspended Hopf map.}
On the usual Hopf coordinates, complex conjugation gives the equivariant map
\[
 \eta_{C_2}:S^{1+2\sigma}\longrightarrow S^{1+\sigma},
 \qquad
 (z_0,z_1)\longmapsto
 \bigl(2z_0\overline z_1,\ |z_0|^2-|z_1|^2\bigr).
\]
Equivariance follows directly by conjugating the displayed formula.  Suspend
by the representation \(2+\sigma\):
\[
 \eta^{C_2}_5=\Sigma^{2+\sigma}\eta_{C_2}:
 S^{3+3\sigma}\longrightarrow S^{3+2\sigma}.
\]
After forgetting the action, this is the ordinary threefold suspension
\(\eta_5:S^6\to S^5\).

Now let \(W=\C^3\) with
\(J=\operatorname{diag}(1,1,-1)\).  As a real representation,
\(W\cong4\cdot\mathbf1\oplus2\sigma\).  Choosing a fixed unit vector in
the first positive complex line and stereographically projecting identifies
the based unit sphere equivariantly with
\[
 S(W)\cong S^{3+2\sigma}.
\]
Thus
\[
 s:=\eta^{C_2}_5\circ q_\wedge:P_v\longrightarrow S(W)
\]
is equivariant and has underlying homotopy class
\(\eta_5q_\wedge\).

\subparagraph{The complement bundle.}
Over \(S(W)=S^5\), define
\[
 \xi_v=v^\perp\subset\C^3.
\]
The tautological line \(\C v\) is globally trivialized by \(v\), and
\[
 \xi\oplus\underline\C\cong\underline{\C}^3.
\]
The fibration
\[
 U(2)\longrightarrow U(3)\longrightarrow S^5
\]
is the unitary frame bundle of this complement.  In its long exact homotopy
sequence, \(\pi_4(U(3))=0\), so the connecting map
\(\pi_5(S^5)=\Z\to\pi_4(U(2))=\Z/2\) is onto.  Consequently \(\xi\)
represents the nonzero class of \([S^5,BU(2)]\).  Its determinant is
ordinarily trivial, so the classifying map lifts to \(BSU(2)\); under the
\(7\)-connected model \(S^4=\mathbb HP^1\to BSU(2)\), it is the nonzero
class \(\eta_4:S^5\to S^4\).  Section~\ref{p1b:app:g15} already records
that \(\eta_4\eta_5q_\wedge\) is the nontrivial G15 class.  Hence
\[
 s^*\xi\cong E_{\rm ex}.
\]

\subparagraph{Coherent involution.}
The unitary involution \(J\) sends \(v^\perp\) to \((Jv)^\perp\), so
\[
 \widehat J(v,w)=(Jv,Jw)
\]
defines a complex bundle involution on \(\xi\) covering the target
involution.  Pullback along the equivariant map \(s\) gives the coherent
lift of Theorem~\ref{p1b:thm:g21-exotic-lift}.  The scalar-torus factor of
\(M_8=P_v\times\T^2\) is irrelevant to the pullback bundle, so the same
construction extends over the full exchange by pullback through
\(\operatorname{pr}_{P_v}\).

At a fixed point \(v\in S(W)^{C_2}=S(\C^2_+)\), the complement splits as
one positive and one negative complex line.  Since \(\C v\) is
equivariantly trivial while \(\det_\C J=-1\), the determinant of the
complement lift is the sign line.  This simultaneously proves the mixed
fixed-locus type and the determinant statement in the main text.



\subsection{Dolbeault model, convergence and explicit Hodge tables}
\label{p1b:app:g16}
\subparagraph{Holomorphic torus fibration and global generators.}
The two marked Hopf projections give
\[
 \T^4\longrightarrow (M_8,J)
     \longrightarrow
     (\mathbb{CP}^1,\epsilon_x)\times(\mathbb{CP}^1,\epsilon_p).
\]
The horizontal root planes are complex lines for \(J\), the vertical
space is \((\mathfrak t,J_0)\), and the differential of the projection
is complex linear. Integrability was proved in
Theorem~\ref{p1b:thm:complex-family}, so this smooth submersion is
holomorphic. Right translation by the four-torus preserves \(J\);
its vertical fundamental fields and their \(J\)-images generate
holomorphic translations of the fiber. Locally the holomorphic
submersion theorem and this action give holomorphic torus-bundle
trivializations. This is the principal-torus case discussed more
generally in \cite{PoddarThakur}; no general cohomology formula from
that reference is assumed here.

Choose a real principal connection \(\rho\) with the two Hopf
connections as its nonclosed components and the two scalar angular
forms as its closed components. Its horizontal space is the sum of
the two root planes. For a \((1,0)\) covector \(\ell_i\) on
\(\mathfrak t\), \(u_i=\ell_i(\rho)\) is a global \((1,0)\) form.
The curvature is basic and of type \((1,1)\); after normalization it is
\(X_xx+X_py\). Thus
\[
 \dd u_i=\bar\partial u_i=C_{ix}x+C_{ip}y,\qquad
 \dd\bar u_i=\partial\bar u_i=\bar C_{ix}x+\bar C_{ip}y.
\]
The pulled-back forms \(x,y\) are closed, square to zero as
differential forms, and represent the two base Dolbeault generators.
They are generally exact on the total space, which is fully retained
in the displayed double complex.

\subparagraph{Why noninvariant forms add no Dolbeault classes.}
Let \(\bar Z_1,\bar Z_2\) be a basis of constant vertical \((0,1)\)
translation fields. In a local holomorphic fiber coordinate their
coefficients are constant; globally they are antiholomorphic fields.
The type-wise Cartan identity is
\[
 \bar\partial\iota_{\bar Z_j}+
       \iota_{\bar Z_j}\bar\partial=\mathcal L_{\bar Z_j}.
\]
Decompose smooth forms into characters \(m\in\Z^4\) of the compact
right torus. On a character, \(\mathcal L_{\bar Z_j}\) is the scalar
\(\lambda_j(m)\). If both scalars were zero, the real covector \(m\)
would annihilate both \((0,1)\) and its conjugate \((1,0)\) space,
so \(m=0\). Since this is an injective real linear map in finite
dimensions, for fixed \(J_0\) there is \(c>0\) such that
\[
 \sum_j|\lambda_j(m)|^2\ge c|m|^2 .
\]
For \(m\ne0\) define
\[
 h_m=\frac{\sum_j\overline{\lambda_j(m)}\,\iota_{\bar Z_j}}
                   {\sum_j|\lambda_j(m)|^2}.
\]
It satisfies \(\bar\partial h_m+h_m\bar\partial=I\).
Its coefficients grow at most polynomially (in fact they decay as
\(|m|^{-1}\)); smooth Fourier coefficients decrease faster than any
polynomial with all derivatives. Summing \(h_m\) therefore gives a
smooth chain homotopy from the identity to torus averaging.
This proves a quasi-isomorphism for the full smooth complex, not
only a computation within left-invariant forms.

\subparagraph{Reduction of the base complex.}
Invariant forms identify, using \(\rho\), with the smooth base double
complex tensored with
\(\Lambda(u_1,u_2,\bar u_1,\bar u_2)\), with the curvature
differentials above. The inclusion
\[
 R=\operatorname{span}_\C\{1,x,y,xy\}
      \hookrightarrow\Omega^{*,*}(\mathbb{CP}^1\times\mathbb{CP}^1)
\]
is a Dolbeault quasi-isomorphism and an algebra map: each projective
line has only \(H^{0,0}\) and \(H^{1,1}\), represented by its constant
and area form. Filter the enlarged complexes by the number of
vertical generators. The associated first differential is the base
\(\bar\partial\); curvature lowers the filtration degree by one.
The finite filtration comparison theorem then makes the inclusion
with \(R\) a Dolbeault quasi-isomorphism. It commutes with both
differentials, hence identifies Frölicher sequences from \(E_1\)
onwards. This establishes the convergence and model assumptions
needed for the full calculation.

This model can also be compared with the principal-torus Dolbeault
model of \cite[Theorem 5.4]{PanovUstinovsky}. The square fan with two
ghost vertices has base \(\mathbb{CP}^1\times\mathbb{CP}^1\), total
smooth space \(S^3\times S^3\times\T^2\), and the same two nonzero
curvature classes. Its four vertical generators give the same
differential algebra. The split cases below agree separately with the
Hopf and Calabi--Eckmann algebras in that reference's Examples 5.11
and 5.12. This supplies a toric-bundle comparison for the direct
averaging and curvature proof above.

\subparagraph{Rank normal forms and geometric examples.}
In rank two, \(C^{-1}\) changes the two complex connection forms to
\(\bar\partial u_x=x,\ \bar\partial u_p=y\).
In rank one a basis gives
\(\bar\partial u=x+\tau y,\ \bar\partial v=0\).
Neither curvature column can be zero. If \(\tau\) were real, a
nonzero real curvature vector would be annihilated by every
\((1,0)\) covector, which is impossible. Thus \(\Im\tau\ne0\).
A split member pairing the two Hopf directions and pairing the two
scalar directions is a Calabi--Eckmann threefold times an elliptic
curve and lies in rank one. A member pairing each Hopf direction
with its own scalar direction is a product of two Hopf surfaces and
lies in rank two. The algebraic calculation does not require a
general member to be biholomorphic to either split example.

\subparagraph{All Hodge numbers.}
Rows below are \(q=0,\ldots,4\), columns are \(p=0,\ldots,4\);
an entry means \(h^{p,q}_{\bar\partial}\).
\[
 \rho=1:\quad
 \begin{pmatrix}
 1&1&0&0&0\\
 2&3&2&1&0\\
 1&3&4&3&1\\
 0&1&2&3&2\\
 0&0&0&1&1
 \end{pmatrix},
 \qquad
 \rho=2:\quad
 \begin{pmatrix}
 1&0&0&0&0\\
 2&0&2&0&0\\
 1&0&4&0&1\\
 0&0&2&0&2\\
 0&0&0&0&1
 \end{pmatrix}.
\]
Both satisfy Serre symmetry \(h^{p,q}=h^{4-p,4-q}\).
In rank one the \(E_1\) algebra has dimension \(32\); the derivation
\(d_1a=(\bar\tau-\tau)\zeta\) has total rank eight, so \(E_2\)
has dimension \(16\). In rank two \(E_1\) already has dimension
\(16\). Both agree at the limiting page with
\[
 (b_0,\ldots,b_8)=(1,2,1,2,4,2,1,2,1).
\]
For example the \(E_2\) class \(a\zeta\) in rank one admits the
total-degree closed representative \((\bar u-u)y\), while
\(u(x-\tau y)\), \(v\) and \(\bar v\) are already closed.
These representatives also check the two degree-three generators
against the underlying sphere-product topology.



\subsubsection{Double-cone scaling and fixed-radius extrinsic geometry}
\label{p1b:app:g18}

\begin{proposition}[Independent cone dilations]
\label{p1b:prop:g18-dilations}
For \(\lambda_x,\lambda_p>0\), let
\(\delta_{\lambda_x,\lambda_p}(q_x,q_p)
=(\lambda_xq_x,\lambda_pq_p)\).  Writing
\(G_{\mathcal C}=G_x+G_p\),
\[
 \delta_{\lambda_x,\lambda_p}^*G_{\mathcal C}
 =\lambda_x^2G_x+\lambda_p^2G_p,
 \qquad
 \delta_{\lambda_x,\lambda_p}^*\vol_8
 =\lambda_x^4\lambda_p^4\vol_8.
\]
Hence the diagonal dilation is a global homothety, while unequal factor
dilations are not.
\end{proposition}
\begin{proof}
Each quaternionic factor has real dimension four, and Euclidean dilation
multiplies its metric by \(\lambda_i^2\) and volume by \(\lambda_i^4\).
\end{proof}

\begin{proposition}[Second fundamental form of the radius leaf]
\label{p1b:prop:g18-second-fundamental}
Give the two radial normal fields the outward orientations \(n_x,n_p\).
For tangent vectors \(X,Y\) to the \(x\)-sphere and \(U,V\) to the
\(p\)-sphere,
\[
 \mathrm{II}(X,Y)=-\frac1{R_x}g(X,Y)n_x,
 \qquad
 \mathrm{II}(U,V)=-\frac1{R_p}g(U,V)n_p,
 \qquad
 \mathrm{II}(X,U)=0,
\]
with the sign reversed if the convention
\(\mathrm{II}(X,Y)=(\nabla_XY)^\perp\) is paired with inward normals.
Thus the trace mean-curvature vector for the displayed convention is
\[
 H_{\rm tr}=-\frac3{R_x}n_x-\frac3{R_p}n_p,
\]
and the averaged mean curvature is \(H_{\rm tr}/6\).
If \(R_x=R_p=R\), the Born fixed diagonal inside the radius leaf is
\(\{(q,q):|q|=R\}\), with induced metric \(2g_{S^3(R)}\), hence isometric
to the round sphere \(S^3(\sqrt2R)\). For unequal radii the exchange does
not preserve the chosen leaf.
\end{proposition}
\begin{proof}
This is the standard sphere shape operator in each Euclidean quaternionic
factor, together with the vanishing of mixed product terms.  The fixed-leaf
statement follows from exchanging the two radius functions.  For the
diagonal inclusion \(i(q)=(q,q)\), one has
\(i^*(g_{S^3(R)}\oplus g_{S^3(R)})=2g_{S^3(R)}\), which gives the
claimed induced radius.
\end{proof}



\subsubsection{Low-rank obstruction theory and unstable rank-two enumeration}
\label{p1b:app:g20}

\subsubsection*{Yang's rank-three congruences on \(M_8\)}
Let
\[
 u_1=n\eta,\quad
 u_2=Aa_tu_x+Ba_tu_p+Ca_Eu_x+Da_Eu_p,\quad
 u_3=2m u_xu_p,\quad u_4=0.
\]
Theorem~1.1 and Theorem~1.2 of \cite{Yang2023Bundles8SpinC} apply with the
spin characteristic class chosen to be zero.  Since \(p_1(TM_8)=0\),
\(\eta^2=0\), and
\[
 u_2^2=2\Delta\,\mathcal O,\qquad
 u_1u_3=2nm\,\mathcal O,
\]
the mod-three condition becomes
\[
 \Delta+2nm\equiv0\pmod3.
\]
The two-primary condition gives \(u_3\equiv0\pmod2\) and, after that
normalization, the top congruence becomes
\[
 \Delta\equiv0\pmod2.
\]
Because \(2nm\) is even, these combine to
\(6\mid(\Delta+2nm)\).

For the fixed-Chern fibre, Yang's Bockstein factor vanishes on the
torsion-free cohomology of \(M_8\).  Using
\[
 v_t=a_tu_xu_p,\qquad v_E=a_Eu_xu_p
\]
as a basis of \(H^7(M_8;\Z)\), the three relevant odd stable-unitary
operations contribute
\[
 6H^7,
 \qquad
 u_2H^3=\langle(A,C),(B,D)\rangle,
 \qquad
 u_3H^1=2m\Z^2.
\]
Thus
\[
 \mathfrak T
 =\frac{\Z^2}{\gcd(6,2m)\Z^2+\langle(A,C),(B,D)\rangle}.
\]
The stated Smith factors follow from the gcd of the entries and of the
\(2\times2\) minors.

\subsubsection*{Even-determinant rank two}
For determinant zero, Corollary~8.4 of \cite{CadekCrabbVanzura2008}
requires
\[
 Sq^2\rho_2(c_2)=0,
 \qquad
 \frac14\left\langle p_1(TM_8)c_2+2c_2^2,[M_8]\right\rangle
 \equiv0\pmod{12}.
\]
The first expression vanishes on the exterior cohomology generators by the
Cartan formula and instability, while the second equals \(\Delta\).  The
criterion is therefore both necessary and sufficient.

\subsubsection*{Odd determinant: twisted-symplectic desuspension}
For odd \(n\), tensoring by a line reduces to \(c_1=\eta\) without changing
\(c_2\).  Let \(L\) be the determinant line.  Following
\cite{AtiyahRees1976}, a complex two-plane with determinant \(L\) is
canonically \(L\)-symplectic and stabilizes by the plane
\[
 E_0=L\oplus\mathbf1.
\]
In the exact exterior K-lattice set
\[
 e=\xi_t\xi_E,\quad
 y_1=\xi_t\xi_x,\ y_2=\xi_t\xi_p,\
 y_3=\xi_E\xi_x,\ y_4=\xi_E\xi_p,
\]
\[
 z=\xi_x\xi_p,\qquad w=\xi_t\xi_E\xi_x\xi_p.
\]
Then \([L]=1+e\), \(z^*=-z\), and \(ez=w\).  The four \(y_i\) have
actual determinant-\(L\) rank-two representatives on the five-complexes
\(T^2\times S_x^3\) and \(T^2\times S_p^3\): the homotopy fibre of
\[
 (c_1,c_2):BU(2)\to K(\Z,2)\times K(\Z,4)
\]
is four-connected with first \(\pi_5\cong\Z/2\), so the first lifting
obstruction is in \(H^6\), which vanishes on a five-complex
\cite{Akbulut1996GaugeRelations}.

For the top generator use the twisted hyperbolic construction
\[
 H_L(Q)=Q\oplus(L\otimes Q^*).
\]
Its underlying K-class satisfies
\[
 F(H_L(z))=z+(1+e)z^*=-w.
\]
Thus the top generator is hit primitively.  If \(k=\Delta/6\), the stable
\(L\)-symplectic class with underlying complex K-class
\[
 x=(2+e)-(Ay_1+By_2+Cy_3+Dy_4)+kw
\]
exists exactly when \(k\in\Z\).

Represent the stabilized class by a complex rank-four \(L\)-symplectic
bundle \(V\), using the stabilization and splitting theorem of
\cite{AtiyahRees1976}.  Its Chern character is
\[
 \operatorname{ch}(V)=4+2\eta-c_2+\frac{\Delta}{6}\mathcal O.
\]
The degree-eight identity
\[
 \operatorname{ch}_4(V)=\frac{c_2^2-2c_4(V)}{12}
\]
and \(c_2^2=2\Delta\mathcal O\) give \(c_4(V)=0\).  For a complex
rank-four bundle on an eight-complex, \(c_4\) is the complete obstruction
to a nowhere-zero section.  Such a section splits an \(L\)-symplectic copy
of \(E_0\), leaving the required actual rank-two bundle.  Conversely an
actual rank-two bundle has
\[
 \operatorname{ch}_4(E)=\frac{\Delta}{6}\mathcal O,
\]
so the integral K-lattice forces \(6\mid\Delta\).

\subsubsection{Determinant trivializations and ordinary complex isomorphism}
\label{p1b:app:g20-determinant-forgetting}
This subsection proves the category boundary in
Remark~\ref{p1b:rem:g20-determinant-forgetting}.  It supplies the detailed
counterexample to forgetting the determinant trivialization; it does not
give a replacement eight-dimensional numerical enumeration.

The determinant fibration
\[
 BSU(2)\longrightarrow BU(2)\xrightarrow{\det}BU(1)
\]
identifies a lift to \(BSU(2)\) with a complex rank-two bundle together
with a chosen trivialization of its determinant.  For a fixed underlying
complex bundle the choices are acted on by
\([M_8,U(1)]=H^1(M_8;\Z)\), by changing the trivialization.  Its orbits
are precisely the ordinary complex isomorphism classes.  Fixing only the
determinant line up to isomorphism does not remove its automorphisms.
Restriction to \(P_v\) preserves the G15 bit because \(H^1(P_v;\Z)=0\).

For an explicit nontrivial orbit take
\[
 Y=S^1_a\times S^1_b\times S^3_u,\qquad c_1=0,\qquad c_2=au.
\]
The generators \(a,b,u\) have degrees \(1,1,3\).  All positive Steenrod
squares on \(H^*(Y;\Z/2)\) vanish by pullback from the sphere factors
and the Cartan formula.  The first relative Postnikov stage of
\((c_1,c_2):BU(2)\to K(\Z,2)\times K(\Z,4)\) has fibre
\(K(\Z/2,5)\) and obstruction
\[
 k_6=Sq^2\bar c_2+\bar c_1\bar c_2
\]
\cite[Proposition~3.1]{Akbulut1996GaugeRelations}.  On the five-complex
\(Y\) there is no degree-six existence obstruction.  Lifts of a fixed
map representing \((c_1,c_2)\) form a torsor for
\(H^5(Y;\Z/2)=\Z/2\).

Loops of that map have coordinates
\[
 c_1(s)=c_1+sh,\qquad c_2(s)=c_2+sv,
 \qquad h\in H^1(Y;\Z),\quad v\in H^3(Y;\Z),
\]
where \(s\) generates the cohomology of the loop circle.  Slanting
\(k_6\) with that circle gives their translation on the lift torsor:
\[
 \delta(h,v)=Sq^2\bar v+\bar c_1\bar v+\bar h\bar c_2.
\]
In the principal-\(SU(2)\) category the determinant trivialization is
fixed, so the \(h\)-loops are absent.  The remaining action
\(Sq^2\bar v\) vanishes here.  Hence there are two principal-\(SU(2)\)
classes with \(c_2=au\).  For ordinary complex isomorphism, the loop
\(h=b,v=0\) acts by
\[
 \delta(b,0)=\bar b\bar a\bar u\ne0,
\]
and interchanges them.  Thus they give one ordinary complex bundle on
\(Y\).  Higher Postnikov stages do not affect this five-dimensional
calculation.

Now write \(M_8=Y\times S_p^3\), using \(a=a_t,b=a_E,u=u_x\).
The projection to \(Y\) has a section.  Pullback therefore preserves
the distinction of the two \(SU(2)\) classes, while their complex
isomorphism also pulls back.  Both have \((A,B,C,D)=(1,0,0,0)\), and
both restrict trivially to \(P_v\), since that restriction factors
through \(S_x^3\).  This proves noninjectivity on the exact manifold and
in the same fixed-Chern, fixed-G15 sector used in the main theorem.

\subsubsection*{The determinant-trivialized mapping 2-type}
The remaining mapping-space and orbit calculations count principal
\(SU(2)\) bundles, with determinant-preserving isomorphisms throughout.
Their unframing operation concerns slice trivializations within this
category; it does not include the additional \(H^1(M_8;\Z)\) quotient
just described.

Put \(P=S_x^3\times S_p^3\), let \(x,p:P\to S^3\) be the two projections,
and set
\[
 \kappa=[x,p].
\]
The Samelson product \(\langle\iota_3,\iota_3\rangle\) generates
\(\pi_6(S^3)\cong\Z/12\) \cite{IwaseMimura2004}.  With the frozen normal
form
\[
 (a,b;z)=p^b x^a\kappa^z,
\]
one has
\[
 (a,b;z)(a',b';z')
 =(a+a',b+b',z+z'+ab'),
\]
\[
 [(a,b;z),(a',b';z')]=(0,0,ab'-a'b).
\]
Thus \([P,S^3]\) is the level-twelve Heisenberg extension of \(\Z^2\).

Evaluation and the suspension splitting
\[
 \Sigma P\simeq S^4\vee S^4\vee S^7
\]
give
\[
 \pi_1\operatorname{Map}(P,S^3)\cong(\Z/2)^3
 =\langle\alpha_x,\alpha_p,\beta\rangle.
\]
Here \(\beta=\nu'\eta_6\) generates \(\pi_7(S^3)\cong\Z/2\).  Naturality
of the Samelson product gives
\[
 \langle\iota_3,\eta_3\rangle
 =\langle\iota_3,\iota_3\rangle\circ\eta_6=\beta,
\]
and therefore
\[
 \rho(a,b;z)(r,s,e)=(r,s,e+br+as).
\]
Same-factor corrections vanish because the reduced diagonal of a suspension
is null, and the central top-cell coordinate acts trivially in this 2-type
by dimension.

The two mapping components corresponding to the trivial and G15-exotic
restrictions have equivalent Postnikov 2-types: applying based maps into
\(BSU(2)\) to the Whitehead-product cofibration of \(P\) gives two homotopy
fibre components, and the canonical loop action translates one to the other.

\subsubsection*{Vanishing of the affine unframing term}
Factor restriction to the individual \(S^3\)'s kills the
\(\alpha_x,\alpha_p\) components of the affine centralizer cocycle.  The
only possible remaining value is the top line \(\langle\beta\rangle\).
The top-cell subgroup is the 1-type of \(\Omega^6S^3\), hence a Picard
1-type with component group \(\Z/12\) and automorphism group \(\Z/2\).
Its nonzero quadratic class is the self-braiding \(\nu'\eta_6\), while the
ordinary cyclic associator is zero by Proposition~\ref{p1b:prop:cyclic-associator-closure}; an order-twelve carry therefore produces
no affine translation.  Mixed top-cell/degree commutators factor through
\(S^6\wedge S^3=S^9\) (and their homotopies through \(S^{10}\)), above the
seven-dimensional range of the mapping 2-type.  Consequently the affine
unframing class vanishes and only the displayed linear action remains.

\subsubsection*{Orbit counts}
For \(c_2=0\), the two scalar monodromies are central; there are
\(12^2\) primary pairs and five secondary orbits, giving
\(144\cdot5=720\) classes in each restriction component.

If \(\Delta=0\) and \(c_2\ne0\), write
\((A,B)=m(a,b)\), \((C,D)=n(a,b)\) with \((a,b)\) primitive, and put
\(d=\gcd(|m|,|n|)\), \(r=\gcd(12,d)\).  The primary quotient has
\(12r\) classes.  The secondary linear orbit count is respectively
\(4,6,5\) for \(d\) odd, \(d\equiv2\pmod4\), and \(4\mid d\), yielding
\(48r,72r,60r\).

If \(\Delta\ne0\), changes of slice trivialization act on
\((z_t,z_E)\in(\Z/12)^2\) by the matrix
\[
 N_c=\begin{pmatrix}B&-A\\D&-C\end{pmatrix}.
\]
Its Smith invariants are
\[
 s_1=\gcd(A,B,C,D),\qquad s_2=|\Delta|/s_1,
\]
so the primary quotient has size
\[
 \gcd(12,s_1)\gcd(12,s_2).
\]
If some Chern coefficient is odd, the top \(\beta\)-line is already
quotiented out and the secondary count is four.  If all coefficients are
even, the stabilizer parity is the kernel of
\[
 \begin{pmatrix}B/2&A/2\\D/2&C/2\end{pmatrix}\pmod2,
\]
and the orbit count is
\[
 4+2^{\rho_2}.
\]
These are exactly the cases in Theorem~\ref{p1b:thm:g20-rank2-fibre-even}.



\subsubsection{Born-equivariant rank-two details}
\label{p1b:app:g21}

\subsubsection*{Born action on the principal-\(SU(2)\) Heisenberg coordinates}
This count uses the determinant-trivialized category in
Theorem~\ref{p1b:thm:g21-born-fixed-count}; Born acts on both the bundle
and its chosen determinant trivialization by pullback.  The determinant
quotient is not part of this orbit calculation.
Use the convention of Section~\ref{p1b:app:g20}:
\[
 \kappa=[x,p],\qquad (a,b;z)=p^b x^a\kappa^z,
\]
\[
 (a,b;z)(a',b';z')=(a+a',b+b',z+z'+ab').
\]
Sphere exchange sends \(x\leftrightarrow p\) and \(\kappa\mapsto\kappa^{-1}\),
so
\[
 \sigma(a,b;z)=(b,a,ab-z).
\]
For Born-fixed Chern data \(C=B,D=A\), this gives
\[
 (z_t,z_E)\longmapsto(AB-z_E,AB-z_t).
\]
On the secondary module,
\[
 \alpha_x\leftrightarrow\alpha_p,\qquad \beta\mapsto\beta.
\]
The abstract Postnikov 2-types of the trivial and exotic G15 restriction
components are equivalent by the homotopy-fibre translation in
Section~\ref{p1b:app:g20}.  Proposition~\ref{p1b:prop:ordinary-born-component-translation} gives an equivalence of mapping two-types intertwining the single Born map up to homotopy. The nonzero square obstruction of Proposition~\ref{p1b:prop:exotic-born-square-obstruction} is retained separately. Hence no
intrinsic additional affine Born shift occurs in the exotic component.

For the primary fixed quotient, let
\[
 N_B=\begin{pmatrix}B&-A\\A&-B\end{pmatrix}.
\]
A primary class is fixed iff
\[
 (AB-z_t-z_E)(1,1)\in\operatorname{im}(N_B\bmod12).
\]
The quotient by the diagonal class \([(1,1)]\) has lattice index equal to
the gcd of the \(2\times2\) minors of
\[
 (B,A),\ (-A,-B),\ (12,0),\ (0,12),\ (1,1).
\]
Using \(12\mid(B^2-A^2)\), that index is
\[
 \gcd(12,B-A)=g_B.
\]
The secondary action then gives two fixed classes if some coefficient is
odd.  In the all-even case the parity stabilizer is the kernel of
\[
 \begin{pmatrix}B/2&A/2\\A/2&B/2\end{pmatrix}\pmod2.
\]
If \(A+B\equiv2\pmod4\), the correcting parity is diagonal and four
secondary classes are fixed.  If \(A+B\equiv0\pmod4\), the two correcting
parity cosets occur equally often and contribute four and two fixed
secondary orbits, averaging to three per primary class.  This proves
Theorem~\ref{p1b:thm:g21-born-fixed-count}.

Changing a determinant trivialization by \(h\in H^1(M_8;\Z)\) is
compatible with exchange through \(h\mapsto B^*h\).  Consequently the
ordinary complex Born-fixed fibre consists of the exchange-stable
determinant orbits, which need not contain a fixed principal-\(SU(2)\)
class.  Enumerating those orbits is an additional problem.  The following
index argument instead applies to any actual coherent complex lift and
does not use the numerical count above.

\subsubsection*{Equivariant Dirac replay}
Regroup \(M_8=K\times K\), \(K=S^1\times S^3\).  The exchange fixed locus is
\[
 F=\Delta K\cong S^1\times S^3,
\]
and the normal bundle is the anti-diagonal copy of \(TK\).  Since \(K\) is
parallelizable, all positive Pontryagin corrections vanish.  The derivative
is \(+I\) tangentially and \(-I\) on the real rank-four normal bundle.
For either Born-fixed spin structure, an involutive spin lift exists; the two
choices change only the overall sign of the nontrivial equivariant character.
The classical Atiyah--Segal--Singer fixed-point denominator therefore has
absolute value
\[
 \det(1-(-I_4))^{1/2}=4
\]
\cite{AtiyahSegal1968,AtiyahSinger1968III}.

On Born-fixed determinant-zero Chern data,
\[
 i^*c_2(E)=2(A+B)au,
\]
so a pure fixed-fibre action gives
\[
 t=\operatorname{ind}_B(D_E)=\pm\frac{A+B}{2}.
\]
Independently, \(TM_8\) is framed and
\[
 s=\operatorname{ind}(D_E)=\frac{B^2-A^2}{6}.
\]
For an actual coherent pure lift, \(t\in\Z\) implies \(A+B\) is even, so
\(4\mid(B^2-A^2)\); also \(s\in\Z\) gives \(6\mid(B^2-A^2)\).  Hence
\(12\mid(B^2-A^2)\) and \(s\) is even.  Since a virtual \(C_2\)-character
has \(s\equiv t\pmod2\), \(t\) is even and
\[
 4\mid(A+B).
\]
Moreover
\[
 m_\pm=\frac{s\pm t}{2}
 =\frac{B^2-A^2}{12}\pm\frac{A+B}{4}.
\]
For mixed eigenranks \((1,1)\), the two eigenbundles are complex lines on
\(F\); since \(H^2(F;\Z)=0\), they are trivial and
\(i^*c_2(E)=0\), giving \(A+B=0\).

\subsubsection*{Symmetric-square no-go for the G15 exotic class}
Let \(P=S^3\times S^3\) with swap quotient
\[
 q_P:P\to Q_P=\operatorname{SP}^2(S^3).
\]
A principal \(SU(2)\) fixed-fibre involution is \(\pm I\).  Multiplication by
central \(-I\) reduces a hypothetical lift to trivial fixed isotropy, which
descends locally and globally to an ordinary principal \(SU(2)\) bundle on
\(Q_P\).

The symmetric-square cofibration is
\[
 S^3\longrightarrow\operatorname{SP}^2(S^3)
 \longrightarrow\Sigma^4\mathbb{RP}^2
\]
\cite{JamesThomasTodaWhitehead1963}.  After collapsing the coordinate wedge,
the induced orbit map is
\[
 r=\Sigma^4p:S^6\to\Sigma^4\mathbb{RP}^2,
\]
where \(p:S^2\to\mathbb{RP}^2\) is the antipodal double cover, equivalently
the top-cell attaching map of \(\mathbb{RP}^3\).  Atiyah's stable splitting
\(\mathbb{RP}^3\simeq_s\mathbb{RP}^2\vee S^3\) makes \(p\) stably null
\cite{Atiyah1961ThomComplexes}; in the present suspension range this gives
\(r\simeq*\).  Since \(BSU(2)\) is three-connected, every class on
\(Q_P\) factors through the reduced symmetric square for the purpose of its
pullback to \(P\), and therefore
\[
 q_P^*:[Q_P,BSU(2)]\to[P,BSU(2)]
\]
is zero.  The G15 class is the nonzero element of
\([P,BSU(2)]\cong\pi_6(BSU(2))\cong\Z/2\), so it cannot descend.  This
proves Theorem~\ref{p1b:thm:g21-exotic-nogo}.

\subsubsection*{Exact principal-\(SU(2)\) descent reduction}
For the full exchange quotient
\[
 q:M_8\to Q=\operatorname{SP}^2(S^1\times S^3),
\]
define
\[
 \operatorname{Desc}(E)
 =\{[\bar E]\in[Q,BSU(2)]:q^*\bar E\cong E\}.
\]
Every coherent principal lift has fixed isotropy \(\pm I\).  In type \(+I\),
the stabilizer acts trivially on the fibre.  On an equivariant slice the
local model is therefore
\[
 (U\times SU(2))/C_2\cong(U/C_2)\times SU(2),
\]
so the quotient is an ordinary principal bundle \(\bar E\to Q\).  Conversely
\(q^*\bar E\) has its canonical exchange lift.  Two pullback identifications
produce gauge-equivalent lifts, while the central signs \(+I\) and \(-I\)
cannot be conjugated.  Thus
\[
 \mathscr L_{SU(2)}(E)\cong\operatorname{Desc}(E)\times\{+,-\}.
\]
This identifies the full nonabelian \(H^1\) set after one splitting is chosen.

A mixed \(U(2)\) lift does not have trivial isotropy, so the coarse quotient
is not its classifying object.  The correct universal object is the genuine
equivariant classifying space retaining the local homomorphism
\(C_2\to U(2)\), as in \cite{LuckUribe2014}.  Its forgetful fibre is the
structural endpoint stated in Remark~\ref{p1b:rem:g21-mixed-fibre}.


\subsubsection{The cyclic associator and the scope of component translation}
\label{p1b:app:rank-two-coherence-closure}

\begin{proposition}[An explicit zero-associator model for the top line]
\label{p1b:prop:cyclic-associator-closure}
The fundamental Picard groupoid of the top-cell mapping space
\(\Omega^6S^3\) has a skeletal model with objects \(C_{12}\),
automorphism group \(C_2\{\beta\}\) at every object, tensor product
given by addition, identity associator, and braiding
\[
 c(a,b)=(a\bmod2)(b\bmod2)\beta.
\]
In particular the ordinary cyclic associator class vanishes, although
the self-braiding of an odd object is \(\beta\ne0\).
\end{proposition}
\begin{proof}
Sixfold loop multiplication is homotopy commutative with a symmetry
coherent through the fundamental groupoid.  Its two groups are
\(\pi_6S^3=C_{12}\{\nu'\}\) and
\(\pi_7S^3=C_2\{\beta\}\).  The trace of its braiding is the
Hopf precomposition operation
\[
 q(a\nu')=a\nu'\eta_6=(a\bmod2)\beta.
\]
For the displayed model, \(c\) is well defined modulo twelve and is
bilinear; hence the pentagon and both hexagons hold with associator
zero.  It is symmetric because \(c(a,b)+c(b,a)=0\), and its trace is
\(c(a,a)=(a\bmod2)\beta=q(a)\).
The Eilenberg--Mac Lane trace classification of braided categorical
groups identifies models with equal trace up to braided monoidal
equivalence.  Thus this is a model of the actual fundamental Picard
groupoid, not a replacement of its nonzero self-braiding by zero.
The trace classification and its explicit zero-associator construction are \cite[Theorem D]{Braunling2019Strictification}; here the quadratic trace is polar with zero polarization.

For example, adding three representatives in \(\{0,\ldots,11\}\)
in either order has precisely the same tensor identification.  The
integer carry at twelve is therefore not an additional automorphism.
This proves the top-line assertion used in affine unframing.  It
does not set to zero the separate braiding correction when actual
top insertions are interchanged.
\end{proof}

\begin{proposition}[Ordinary Born intertwining of the two components]
\label{p1b:prop:ordinary-born-component-translation}
Let \(P=S^3\times S^3\), let \(s\) interchange its factors, and set
\(Y=BSU(2)\).  Write \(\mathcal M_\epsilon\) for the two components
of \(\operatorname{Map}_*(P,Y)\), indexed by the top class
\(\epsilon\in\pi_6Y=C_2\).  There is a homotopy equivalence
\[
 T:\mathcal M_0\longrightarrow\mathcal M_1,
 \qquad s^*T\simeq Ts^*.
\]
Only one homotopy intertwining the two maps is asserted; no coherent
\(C_2\) action on this equivalence is asserted.  Consequently it
intertwines the ordinary Born action on maps from the labelled
two-torus, including exchange of its two circle coordinates.  The two G15 sectors therefore have the same Born-fixed multiplicity among determinant-preserving isomorphism classes; no coherent lift or determinant-forgetting equivalence is asserted.
\end{proposition}
\begin{proof}
Use the product CW presentation
\[
 S^5\xrightarrow{[\iota_x,\iota_p]}
 W=S^3\vee S^3\longrightarrow P\longrightarrow S^6.
\]
Pinching the cone parameter gives the usual coaction
\(c:P\to P\vee S^6\).  It is natural under interchange of the
two factors: on its sphere summand the action is the interchange of
the two three-dimensional top-cell coordinates.  To work with based
cofibres one may choose the boundary basepoint on the fixed diagonal;
the cone pinch fixes that line.  Thus the source comparison is
equivariant before any map into \(Y\) is applied.  Equivalently one
can pinch a small invariant ball about a diagonal interior point and
use a diagonal arc to the original basepoint.  The arc is fixed
pointwise, so the resulting based coaction has the same equivariance.

Choose \(e:S^6\to Y\) representing the nonzero class.  The coaction
defines
\(T_e(f)=(f\vee e)c\).  Its coassociativity gives
\(T_eT_{-e}\simeq T_0\simeq\operatorname{id}\), so it is a homotopy
equivalence between these two mapping components.  If \(S\) is the
interchange map on \(S^3\wedge S^3=S^6\), naturality gives
\[
 s^*T_e\simeq T_{eS}s^*.
\]
The degree of \(S\) is \((-1)^{3\cdot3}=-1\).  Therefore
\([eS]=-[e]=[e]\) in \(\pi_6Y=C_2\).  Choose one based homotopy
\(eS\simeq e\), and apply it in the sphere summand of the same
coaction.  This gives the required homotopy of maps of entire mapping
spaces; it is independent of the varying \(f\).

For clarity about unbased bundles, evaluation
\(\operatorname{Map}_*(P,Y)\to\operatorname{Map}(P,Y)\to Y\)
induces an equivalence on Postnikov two-types because \(Y\) is
three-connected.  A two-torus and homotopies between its maps see only
that two-type.  Thus the argument applies to the determinant-preserving
bundle classification used here.  The coaction is the identity on
\(W\), and restriction to \(T^2\times W\) detects all four
degree-four Chern coefficients of \(T^2\times P\).  The equivalence
therefore preserves the specified \((A,B,C,D)\), as well as its
labelling of the two torus circles.  Composing a torus family with
\(T_e\), and composing its Born transform with the supplied single
homotopy, gives a bijection of ordinary homotopy classes intertwining
Born.  No square coherence is needed to infer a bijection of fixed
isomorphism classes.  The finite orbit arithmetic in the zero sector
therefore transports to the exotic sector without supplying any
coherent lift.
\end{proof}

\begin{proposition}[The exotic translation has a nonzero coherence square]
\label{p1b:prop:exotic-born-square-obstruction}
The preceding equivalence cannot be strengthened to a coherently
\(C_2\)-equivariant equivalence of these two mapping two-types.  At the
exotic top-cell object the obstruction to the square coherence is
\[
 \alpha\eta_6\in\pi_7BSU(2)=C_{12},\qquad
 \alpha=\operatorname{adj}^{-1}(\eta_3\eta_4),\qquad
 \operatorname{adj}(\alpha\eta_6)
       =\eta_3\eta_4\eta_5=6\nu'\ne0.
\]
In particular a single homotopy intertwining Born must not be used as
a proof of the cocycle equation for a Born lift.
\end{proposition}
\begin{proof}
The actual representation on the top cell is
\[
 V=\mathbb R^3_x\oplus\mathbb R^3_p
    \cong\mathbb R^3\oplus3\sigma,
\]
where \(\sigma\) is the sign representation.  It is not a single
sign coordinate.  A homotopy fixed top-cell object, through the square
coherence, would extend the map \(\alpha:S^6\to Y\) over the
two-dimensional Borel skeleton of this representation sphere.  That
skeleton is the Thom space
\[
 \operatorname{Th}(\mathbb R^3\oplus3\gamma\mid\mathbb{RP}^2)
 \cong\Sigma^3(\mathbb{RP}^5/\mathbb{RP}^2).
\]
It has cells in dimensions six, seven and eight.  The first two form
\(M=S^6\cup_2e^7\): the orientation character of \(3\gamma\) is
nontrivial.  Its final attaching map lies in
\(\pi_7M=C_2\{i\eta_6\}\).  Indeed the exact sequence of the
degree-two cofibre gives this group, with generator the bottom Hopf
map.  The Thom identity gives
\[
 Sq^2U=w_2(3\gamma)U=\binom32x^2U=x^2U\ne0.
\]
Thus the final attaching map is the nonzero \(i\eta_6\), since the
zero attaching map would make this bottom-to-top square vanish.
The Moore-cell extension of \(\alpha\) exists because \(2\alpha=0\).
For every choice of its nullhomotopy the final obstruction is exactly
\(\alpha\eta_6\).  Changing that null cannot alter the value because
the last attaching map factors through the bottom sphere.  Adjunction
and the integral relation already used in RT.23 give the displayed
nonzero \(6\nu'\).

This top obstruction cannot be cancelled by changing the comparison
on either lower sphere.  In the frozen Heisenberg coordinates their
possible change is \(h=(a,b;z)\), with
\(s(h)=(b,a;ab-z)\).  Its norm is
\[
 h\,s(h)=(a+b,a+b;a(a+b)).
\]
To preserve the zero lower components one must have \(a+b=0\), and
then its central norm is zero.  Hence it cannot cancel
\((0,0;6)\).  This proves nonvanishing of the square obstruction in
the full mapping two-type, rather than just in its isolated top line.
It is consistent with the no-go for a coherent principal-\(SU(2)\)
lift in the exotic sector, and does not contradict the ordinary fixed
class count proved above.
\end{proof}



\subsubsection{Oriented cohomotopy correction and its exact scope}
\label{p1b:app:v2170-ranktwo-reconciliation}
The oriented adjoint cohomotopy extension uses \(Sq^2+\bar\tau\),
not the untwisted \(Sq^2\) in isolation.  The derivation is
Lemma~\ref{p1b:rt:lemma124}; in particular the correction term is
needed for the actual \(C_4\) circle square.  Its mod-two
Moore and doubling/Bockstein inputs are documented in
\cite{Schwede2001TwoLocal,LiPanWu2022Modular,KirbyMelvinTeichner2012}.
The complete resulting rank-two classification is
Theorem~\ref{p1b:thm:complete-mixed-rank-two}; former partial ledgers
and unassigned comparison bits are not alternative conclusions.

\section{Classification, natural-extension ledgers, and geometric crosschecks}
\label{p1c:sec:ledgers}

\subsection{Expanded natural extension ledgers and verification scope}
\label{p1b:app:rank-two-closure-ledgers}
The main theorem states the exact obstruction criterion and its domain.
This section records all named rows for reproducibility; it adds no
hypothesis. Degree-eight entries are \emph{not reached} when the
degree-seven obstruction already fails. The listed source constructions specify exact maps and finite-word
calculations. Availability of the separately cited scripts and checksums
is recorded in the accompanying evidence ledger; finite matrix replay
cannot replace a homotopy-parity certificate.
\subsubsection{Fourteen natural targets}
\begin{longtable}{@{}lllll@{}}
\toprule Candidate&Domain&$O_7$&$O_8$&Extension\\\midrule\endhead
S\_E0\_T0&odd sphere&$0$&$\Theta$&NO\\
S\_E0\_T1&odd sphere&$0$&$0$&YES\\
S\_E1\_T0&odd sphere&$\Gamma_\chi$&not reached&NO\\
S\_E1\_T1&odd sphere&$\Gamma_\chi$&not reached&NO\\
D\_E0\_T0&odd exotic&$0$&$\Theta$&NO\\
D\_E0\_T1&odd exotic&$0$&$0$&YES\\
D\_E1\_T0&odd exotic&$\Gamma_\chi$&not reached&NO\\
D\_E1\_T1&odd exotic&$\Gamma_\chi$&not reached&NO\\
C\_T0&$A=1$ odd exotic&$0$&$0$&YES\\
C\_T1&$A=1$ odd exotic&$0$&$\Theta$&NO\\
B\_D0\_P0&odd exotic&$0$&$0$&YES\\
B\_D0\_P1&odd exotic&$0$&$\Theta$&NO\\
B\_D1\_P0&odd exotic&$\Gamma_\chi$&not reached&NO\\
B\_D1\_P1&odd exotic&$\Gamma_\chi$&not reached&NO\\
\bottomrule\end{longtable}
\subsubsection{Thirty-two inherited candidates}
\begin{longtable}{@{}lllllll@{}}
\toprule Sector&$d$&$p$&$a$&$e$&$t$&Extension\\\midrule\endhead
0&0&1&0&0&0&YES\\
0&0&1&0&0&1&NO\\
0&0&1&0&1&0&NO\\
0&0&1&0&1&1&NO\\
0&0&1&1&0&0&NO\\
0&0&1&1&0&1&NO\\
0&0&1&1&1&0&NO\\
0&0&1&1&1&1&NO\\
0&1&0&0&0&0&YES\\
0&1&0&0&0&1&NO\\
0&1&0&0&1&0&NO\\
0&1&0&0&1&1&NO\\
0&1&0&1&0&0&NO\\
0&1&0&1&0&1&NO\\
0&1&0&1&1&0&NO\\
0&1&0&1&1&1&NO\\
1&0&1&0&0&0&NO\\
1&0&1&0&0&1&YES\\
1&0&1&0&1&0&NO\\
1&0&1&0&1&1&NO\\
1&0&1&1&0&0&NO\\
1&0&1&1&0&1&NO\\
1&0&1&1&1&0&NO\\
1&0&1&1&1&1&NO\\
1&1&0&0&0&0&NO\\
1&1&0&0&0&1&NO\\
1&1&0&0&1&0&NO\\
1&1&0&0&1&1&NO\\
1&1&0&1&0&0&NO\\
1&1&0&1&0&1&YES\\
1&1&0&1&1&0&NO\\
1&1&0&1&1&1&NO\\
\bottomrule\end{longtable}
All inherited rows have the retained odd parity and their source-defined
candidate types. The four YES identifiers are the two sector-zero pure
finite-zero rows, the sector-one $d=0,p=1,a=e=0,t=1$ row, and the
sector-one $d=1,p=0,a=t=1,e=0$ row. Thus the natural ledger is 4 YES/10
NO and the inherited ledger 4 YES/28 NO. The scalar coefficient
$\zeta_{\rm nat}=0$ is distinct from the local circle class
$\zeta=2\kappa$.

\subsubsection{Supplementary classifying-map decidability}
\label{p1b:app:rank-two-decidability}
This generic decidability comparison is not consumed by the completed exact classification. It uses the specified finite classifying-space target and rational H-space comparison hypotheses below, and supplies no practical time bound.
The corresponding classifying-map extension uses
\(W\times I\) relative to
\((W\times\partial I)\cup(R\times I)\).
A finite rank-two target is \(\operatorname{Gr}_2(\mathbb C^7)\):
its homotopy fibre over \(BU(2)\) is the complex two-frame Stiefel
manifold, with sphere base \(S^{13}\) and fibre \(S^{11}\).
The comparison is therefore 11-connected. The Chern map is rationally
an H-space comparison through dimension ten, so Manin's Theorem A
applies to the nine-dimensional relative problem. This theorem gives
decidability, not evaluated ranks or a practical time bound.
The evaluated restrictions, sphere relations and general-Chern fibres
appear in the rank-two proof sections above. The generic decidability
comparison here is not substituted for those computations.





\subsection{Comparison with the \(\mathfrak u(2)\oplus\mathfrak u(2)\) normal form}
\label{p1b:app:g22}

Fix Cartan generators \(H_x,H_p\) and central generators \(z_x,z_p\).
Magnin's matrix normal form acts by the standard complex structure on one
root plane in each \(\mathfrak{su}(2)\) summand and by a matrix
\(M\in GL(4,\R)\), \(M^2=-I\), on
\[
 \operatorname{span}\{H_x,H_p,z_x,z_p\}.
\]
The normalizing automorphism has independent \(SO(3)\) rotations on the two
simple ideals and positive scalings on the two centers.  If \(D\) denotes
the latter diagonal scaling on the four-dimensional Cartan--central block,
then
\[
 DMD^{-1}
\]
is again an arbitrary real complex structure on that block.  Consequently
no central scaling is needed once the paper's normal form allows all
\(J_0\in GL(4,\R)/GL(2,\C)\): it is absorbed by replacing \(M\) with
\(DMD^{-1}\).  What remains is the pair of \(SO(3)\) rotations, hence inner
automorphisms of the compact \(\SU(2)\) factors.  This is the precise
reason the fixed-marking family is an automorphism-saturating normal form
without being itself an automorphism quotient.



\subsection{Exceptional product ADE bookkeeping}
\label{p1b:app:g23}

For later use, the one-factor invariant rings can be written as weighted
hypersurfaces with generator degrees and relation degrees
\[
\begin{array}{c|c|c}
\Gamma&\text{generator degrees}&\text{relation degree}\\\hline
2T&(6,8,12)&24\\
2O&(8,12,18)&36\\
2I&(12,20,30)&60.
\end{array}
\]
Indeed
\[
 \frac{1-q^{24}}{(1-q^6)(1-q^8)(1-q^{12})}
 =\frac{1+q^{12}}{(1-q^6)(1-q^8)},
\]
and similarly for \(2O\) and \(2I\).  For an ordered pair
\((\Gamma_x,\Gamma_p)\), the product invariant ring is therefore a complete
intersection with six generators and the two independent one-factor
relations.  The bigrading remembers which relation and which three
generators belong to each reciprocal factor.

No mixed algebraic relation is created by the direct-product action.  Mixed
invariant polynomials are products of one-factor invariants, which is exactly
the tensor-product statement of Theorem~\ref{p1b:thm:g23-invariants}.



\subsection{Independent-factor rescaling and automorphism/isometry firewalls}
\label{p1b:app:g24}

For
\[
 R_x\mapsto\lambda_xR_x,
 \qquad R_p\mapsto\lambda_pR_p
\]
with the scalar-torus metric fixed, the vector-sphere volume scales by
\(\lambda_x^3\lambda_p^3\), while
\[
 \operatorname{Scal}
 =\frac6{\lambda_x^2R_x^2}+\frac6{\lambda_p^2R_p^2}.
\]
The two sphere Laplacians scale independently by
\(\lambda_x^{-2}\) and \(\lambda_p^{-2}\) on their respective tensor
factors.  Equal-radius symmetry is restored exactly when
\(\lambda_xR_x=\lambda_pR_p\).  The Born map additionally requires the
scalar-torus swap to be an isometry.

The compact Lie-group automorphism group and the metric isometry group
should not be conflated.  At Lie-algebra level the center admits
\(GL(2,\R)\), whereas automorphisms of the compact central torus must
preserve its integral lattice and therefore use \(GL(2,\Z)\).  For a fixed
metric, only the finite subgroup \(O(G_{\rm sc};\Z)\) survives in the torus
point group.  Conversely the round sphere has the full \(O(4)\) isometry
group, larger than the inner \(SO(3)\) automorphisms acting on its Lie
algebra by conjugation.




\subsection{Retained relative-frame and joining reference identities}
The following four compact crosschecks preserve the precise coordinate
identities cited by the positive proof chain. They are not older
unevaluated extension ledgers and do not supersede the exact positive
classification in Sections~\ref{p1c:sec:rank-two-full-proofs} and
\ref{p1c:sec:natural-rank-two-full-proofs}.

\subsubsection{Relative normal frame and its boundary rotation}
\label{p1b:app:rt43-framing}
The ordered normal monodromy and its absolute replacement have matrices
\(D_3=\operatorname{diag}(I_3,-I_3)\) and
\(D_1=\operatorname{diag}(I_5,-1)\). A concrete absolute change of frame is
\[
 A(\theta)=\operatorname{diag}(I_3,R(\theta/2),1),\qquad
 A(2\pi)D_3=D_1A(0).
\]
On the twice-traversed boundary the residual rotation is the
nontrivial loop of \(\pi_1SO(6)=C_2\). Thus an absolute
bundle isomorphism cannot simply discard the specified product-boundary
framing. Its stable \(J\)-class is the nonzero \(\eta_6\).

\subsubsection{Gauge and natural top-coordinate distinction}
\label{p1b:app:rt51-construction}
For the determinant closing equation the two central solutions are
\(z=0,6\) in sector zero and \(z=3,9\) in sector one.
The sphere equation has the solutions \(z=2+2a,8+2a\pmod{12}\).
The square of either joining coefficient \((0,2,0)\) or
\((0,2,1)\) is \((0,4,0)\). This is a \(\pi_4SU(2)\)
identity, not a value of the separate natural \(\pi_7SU(2)\)
coefficient. Those top coefficients are evaluated in the complete
positive source maps, not selected by this lower-stage arithmetic.

\subsubsection{Joining algebra cross-term}
\label{p1b:app:rt49-joining}
The component joining law admits the faithful binary matrix model
\[
 (k,n,\epsilon)\longmapsto
 \begin{pmatrix}1&n&\epsilon\\0&1&k\\0&0&1\end{pmatrix}.
\]
The component matching action retains its \(\ell d\) term:
\[
 \begin{pmatrix}1&0&0\\d&1&0\\\epsilon_c+\ell d&\ell&1\end{pmatrix}.
\]
Neither the joining \(\pi_4\)-coordinate nor its finite matrix
representation replaces the higher natural boundary comparison.

\subsubsection{Actual pushout versus relative top tests}
\label{p1b:app:rt47-pushout}
A trivial bundle on the four-dimensional joining cell does not
trivialize the two original piece bundles. The degree-three image
\(2\mathbb Z\) only provides a parity restriction on Chern
classes; it is not a certificate that every even clutch change is
an actual piece automorphism. A model with absolute CW dimension
seven may still have a relative eight-cell arising from its
seven-dimensional boundary. Thus the natural top comparison is made
on the actual relative pair with its supplied boundary frame.

\section{Exact finite cubic of the native quadratic family embedding}
\label{p1c:sec:v242-native-family-cubic}
\begin{proposition}[Rooted Hurwitz cubic and actual twenty-five-cell deletion]
\label{p1c:prop:v242-native-family-cubic}
Use precisely the rays, operator trimer and functions \(b_j,f_j\) of
Theorem~\ref{p1b:thm:v242-family-quadratic-embedding}. On the root
Hurwitz cell \(H=2T\),
\[
 \boxed{\langle b_jb_kb_l\rangle_H=
 \begin{cases}
 4/9,&j=k=l\text{ or all three indices distinct},\\
 -2/9,&\text{exactly two indices equal}.
 \end{cases}}
\]
On the full icosian set \(I=2I\) this tensor is zero.
For the geometric direction \(g_e=(-1-i,i,1)^T/2\), put
\(f=\mathcal Jg_e\). Then
\[
 \boxed{\langle f^3\rangle_H=\tfrac34(1-i),\qquad
 \langle\bar f f^2\rangle_H=-\tfrac14(1+i),}
\]
whereas both full icosian averages vanish. On the twenty-four retained
registers, with the original \(1/25\) normalization, these values are
respectively \(-3(1-i)/100\) and \((1+i)/100\).
\end{proposition}
\begin{proof}
The root Hurwitz cell consists of the eight coordinate unit vectors in
\(\mathbb R^4\) and the sixteen points \((\pm1,\pm1,\pm1,\pm1)/2\).
The icosian set adds the ninety-six even coordinate permutations of
\(\frac12(0,\pm1,\pm\varphi,\pm\varphi^{-1})\), with
\(\varphi^2=\varphi+1\). Odd coordinate moments vanish by sign symmetry.
For distinct coordinate indices the complete required sixth moments are
\[
\begin{array}{c|ccc}
 &\langle x_1^6\rangle&\langle x_1^4x_2^2\rangle
                         &\langle x_1^2x_2^2x_3^2\rangle\\\hline
 H&3/32&1/96&1/96\\
 I&5/64&1/64&1/192
\end{array}
\]
For example the ninety-six additional points contribute
\(24(1+\varphi^6+\varphi^{-6})/64=57/8\) to
\(\sum x_1^6\). Their contribution to \(\sum x_1^4x_2^2\)
is \((8\cdot4-19)/8=13/8\), and to
\(\sum x_1^2x_2^2x_3^2\) is \(3/8\). Adding the Hurwitz sums
and dividing by \(120\) gives the second row. The first row follows
immediately from the axis points and half-coordinate points.

Write \(y=\operatorname{Im}q\), \(s=y_0+y_1+y_2\) and
\(t=y_0^2+y_1^2+y_2^2\). Direct substitution of the rooted rays gives
\[
 b_j=\tfrac23(s-3y_j)^2-2(t-s^2/3).
\]
Expansion of the three quadratic factors using the sixth-moment table
gives the two displayed cubic tensors. For \(\sum a_j=0\),
\(\mathcal Ja=\sum a_jb_j\); contraction yields \(6a_0a_1a_2\).
The same trilinear formula, with the first coefficient conjugated, gives
the mixed cubic. Substituting \(g_e\) gives the stated values.
The full icosian sixth moments equal the Haar moments; the Haar cubic
on these augmentation modes is also zero, since a quarter rotation in
\(E\) negates every trace-free symmetric quadratic while preserving
Haar measure.

Finally take the actual cells \(g_iHg_j^{-1}\). For fixed \(j\),
the five left cosets \(g_iH\), right-translated by \(g_j^{-1}\),
partition \(I\). Thus every vertex belongs to exactly five of the
twenty-five registers. Their full normalized average equals the uniform
icosian average. Deleting the root register therefore gives
\((25C_I-C_H)/25=-C_H/25\), or \(-C_H/24\) when renormalized
over the retained registers. The quadratic Gram is \((24/25)I_G\)
by the exact frame theorem. This proves the claimed raw values and
normalization without treating overlapping cells as disjoint vertices.
\end{proof}

\begin{remark}[Reproducible exact scope]
The coefficient computation uses Gaussian integers and exact fractions in
\(\mathbb Q(\sqrt2,\sqrt3,\sqrt5,i)\). Its accompanying checker
reconstructs the signed Fock tensor, joint Hopf ratios, real-vector seed
isometry, both relative/common-frame returns, all 120 vertices and 25
cells, and the tensors above. These finite checks corroborate the displayed
proofs. They neither evaluate an actual physical Yukawa vertex nor replace
the same-action source-incidence certificate.
\end{remark}

\section{One-root comparison: finite character, moment, and boundary proofs}
\label{p1c:sec:v243-one-root-replays}
The results here prove the mathematical provider statements of
Section~\ref{p1b:sec:v243-one-root-geometric-closures}.  No finite
host operation below is by itself a physically licensed action insertion.

\begin{proposition}[Nested subgroup ranks and the quartic-exactness test]
\label{p1c:prop:v243-reynolds-character-replay}
On the spin-$n/2$ representation $V_n$, for $G=2T,2I$,
\[
 \mu_n^G=\frac1{|G|}\sum_{g\in G}U_n(\operatorname{Re}g),
\]
where $U_n$ is the second-kind Chebyshev polynomial.  Their exact
ranks for the decisive degrees are
\[
\begin{array}{c|ccccc}
 n&0&1\text{--}5&6&8&12\\\hline
 \mu_n^{2T}&1&0&1&1&2\\
 \mu_n^{2I}&1&0&0&0&1.
\end{array}
\]
The nested projection $P_{2I}-P_{2T}/25$ therefore has a negative
one-dimensional channel first at $n=6$ and distinct positive and
negative one-dimensional channels at $n=12$.
For $u=q_0^2-q_1^2$ the subgroup and round fourth moments satisfy
\[
 \tau_{2T}(u^4)=\tfrac16,\quad
 \tau_{2O}(u^4)=\tfrac5{48},\quad
 \tau_{2I}(u^4)=\int_{S^3}u^4d\sigma=\tfrac3{40}.
\]
A finite subgroup other than $2I$ fails exactness of all quartic
powers on the complete $\mathcal P_2$ band.
\end{proposition}
\begin{proof}
The $24$ Hurwitz quaternions are the eight coordinate axis units
and sixteen half-units $(\pm1,\pm1,\pm1,\pm1)/2$; the full $120$
icosians add the $96$ even coordinate permutations of
$(0,\pm1,\pm\varphi,\pm\varphi^{-1})/2$ with
$\varphi=(1+\sqrt5)/2$.  Substitute their exact real parts into
the displayed character sum using $U_0=1$, $U_1=2x$ and
$U_{n+1}=2xU_n-U_{n-1}$, obtaining the rank table.  It is
independent of any spectral or physical interpretation of a level.
For $2T$, the eight axis points contribute $4$ nonzero unit values
to the sum of $u^4$ and the sixteen half-units contribute zero,
yielding $4/24=1/6$.  The $2O$ set adds $24$ two-coordinate
half-root units: its direct squared-coordinate enumeration gives
$\tau_{2O}(u^4)=5/48$.  The $2I$ evaluation yields $3/40$,
also obtained independently from the round spherical fourth moment
of $q_0^2-q_1^2$.  Cyclic subgroups lie on a great circle, so a
transverse linear function has vanishing group fourth power but
positive Haar fourth moment.  Every binary dihedral subgroup is
conjugate into the union of two orthogonal great circles; a
coordinate product vanishing on both has the same contradiction.
Finite-subgroup classification then leaves $2I$ alone.  Conversely
the $2I$ 11-design property evaluates every degree-eight product
exactly.
\end{proof}

\begin{proposition}[Four-atom classification and complete low-order moments]
\label{p1c:prop:v243-four-atom-moment-proof}
For real $a\in G_\perp$ write $f_a=\mathcal Ja$ and
$\ell_j=-2v_j+v_{j+1}+v_{j+2}$ for $v=\operatorname{Im}q$.
Then
\[
 f_a(q)=\frac23\sum_{j=0}^2a_j\ell_j(q)^2,
 \qquad \sum_ja_j=0.
\]
The $24$ rooted Hurwitz values comprise six zeros and six
occurrences of each value $2a_0$, $2a_1$, $2a_2$, with repetitions
counted even when some values coincide.  The induced counting law is
therefore the four-atom law in
Theorem~\ref{p1b:thm:v243-four-atom-family-law}.
Its moments satisfy
\[
\begin{array}{c|rrrr}
 &k=2&k=3&k=4&k=5\\\hline
 \tau_{2T}(f_a^k)&r^2&6c&2r^4&20r^2c\\
 \tau_{2I}(f_a^k)&r^2&0&27r^4/10&0,
\end{array}
\]
where $r^2=\sum a_j^2$ and $c=a_0a_1a_2$.
\end{proposition}
\begin{proof}
For the two real-axis units and the four half-units whose three
imaginary signs agree, $\ell_j^2$ does not depend on $j$, so
$f_a=0$.  For each $j$ the two pure imaginary axis units $\pm e_j$
and four half-units with one deviant imaginary sign at $j$ give
$f_a=2a_j$.  Thus there are exactly four groups of six labelled
values.  For $k\ge1$ the moment is
$\tfrac14\sum_j(2a_j)^k=2^{k-2}\sum_ja_j^k$.
Since $\sum_ja_j=0$, Newton's relations give
$\sum a_j^3=3c$, $\sum a_j^4=r^4/2$ and
$\sum a_j^5=(5/2)r^2c$, which prove the first row.
The $2I$ 11-design identifies moments through order five with
round integrals.  For a trace-free two-plane quadratic an orthogonal
rotation on the round sphere yields
$f_a=\sqrt6\,r(X^2-Y^2)$ for suitable orthonormal real coordinates.
The normalized round moments give
$\mathbb E(X^2-Y^2)^4=3/40$, and exchange $X\leftrightarrow Y$
kills every odd power.  Hence the fourth moment is
$36r^4(3/40)=27r^4/10$.  The second moment follows from
Theorem~\ref{p1b:thm:v242-family-quadratic-embedding}.
Subtracting $1/25$ of the first row from the second establishes
all the asserted rooted moments.
\end{proof}

\begin{proposition}[First icosian sextic defect and Riesz harmonic]
\label{p1c:prop:v243-icosian-sextic-riesz-proof}
For $a=(1,-1,0)$, exact enumeration of the $120$ icosian
quadratic values gives
$\tau_{2I}(f_a^6)=1241/16$, while normalized
round integration gives $540/7$.  On
$\mathcal H_{12}\cong V_{12}\otimes V_{12}^*$, the Riesz vector
of $2I$ evaluation is
$h_{12}(q)=(13/120)\sum_{g\in2I}\chi_{12}(g^{-1}q)$,
with $h_{12}|_{2I}=13$ and $\|h_{12}\|^2=13$.
\end{proposition}
\begin{proof}
Use the same $8+16+96$ exact icosian-coordinate enumeration as
in the preceding proposition, now retaining all monomials in
$f_a^6$ of total polynomial degree twelve.  Since
$\varphi^2=\varphi+1$, reduction to rational coefficients and the
$s\mapsto-s$ pairing gives the rational result $1241/16$.
Independently, $f_a=\sqrt6\,\sqrt2(X^2-Y^2)$ and the round moment
$\mathbb E(X^2-Y^2)^6=5/112$ yield
$6^3\cdot2^3\cdot5/112=540/7$.
The degree-$12$ scalar reproducing kernel on $S^3=SU(2)$ is
$13\chi_{12}(p^{-1}q)$; averaging its first variable over the
finite subgroup gives the displayed Riesz vector.  The Reynolds
rank is $\mu_{12}^{2I}=1$, so $h_{12}(g)=13$ for $g\in2I$;
by the Riesz identity its squared norm is
$\tau_{2I}(h_{12})=13$.  The difference $47/112$ is nonzero,
therefore the first unsuppressed degree-$12$ sample defect is also
present in the actual polynomial $f_a^6$.
\end{proof}

\begin{proposition}[Six-port augmentation isometry and unrooted source firewall]
\label{p1c:prop:v243-neutral-boundary-source-proof}
For the diagonal-line representative $u=(1,1,1)/\sqrt3$ the six
rooted boundary half-units can be indexed by
$h_\varepsilon=(1+\varepsilon_1i+\varepsilon_2j+
\varepsilon_3k)/2$, where
$\varepsilon\in\{\pm1\}^3\setminus\{(+,+,+),(-,-,-)\}$.
Let $\operatorname{ev}_\Omega$ denote unrestricted evaluation
on these six points, and let $V(w)(h_\varepsilon)=
(\varepsilon\cdot w)/2$.  Then
\[
 \operatorname{ev}_\Omega(\mathcal Ja)\in W_{2,+},\quad
 \|\operatorname{ev}_\Omega(\mathcal Ja)\|^2=8\|a\|^2,
 \qquad V^\dagger V=2I_3-\tfrac12J_3.
\]
The maps $\operatorname{ev}_\Omega\mathcal J/\sqrt8$ and
$V(V^\dagger V)^{-1/2}$ are orthogonal onto the even two-dimensional
and odd three-dimensional mean-zero boundary spaces.  Their sum,
followed by the already constructed isometry $B^\dagger$, is the
complete $\mathcal I_5$ of
\eqref{p1b:eq:v243-five-augmentation-isometry}.  Conversely the
Riesz $h_{12}$ has no augmentation image:
$P_{\rm aug}h_{12}=0$ and
$\|P_{\rm sing}h_{12}\|^2=13/3$.
\end{proposition}
\begin{proof}
The six sign triples occur in three pairs $\varepsilon,-\varepsilon$.
By direct evaluation the even family values on these three pairs are
$2a_0,2a_1,2a_2$, so their unnormalized squared norm is
$2\sum_j|2a_j|^2=8\|a\|^2$.  The odd coordinates sum to zero;
their Gram is
$\tfrac14\sum_{\varepsilon\in\Omega}
\varepsilon\varepsilon^T=2I_3-J_3/2$, with eigenvalues
$2,2,1/2$.  Even/odd orthogonality and dimension sum prove
surjectivity.  The degree-$12$ restriction coisometry satisfies
$E_\Omega E_\Omega^\dagger=156I_6+13J_6$ and
$BB^\dagger=I_6$ by
Theorems~\ref{p1b:thm:v2208-boundary-gram} and
\ref{p1b:thm:v2208-boundary-coisometry}.  As
$h_{12}(h_\varepsilon)=13$ on all six points,
$Bh_{12}=(13/\sqrt{234})\mathbf1_6$ has only the
singlet component.  Its squared norm is
$6\cdot169/234=13/3$; subtraction from
$\|h_{12}\|^2=13$ leaves $26/3$ in $\ker B$.
\end{proof}

\section{Exact finite Petersen--Hopf certificate and reciprocal tensor check}
\label{p1c:sec:petersen-exact-certificate}

This section provides independent, explicitly finite verification data
for Paper~I-B, Section~\ref{p1:sec:petersen-finite-resolution}, and
Paper~I-C, Section~\ref{p1b:sec:petersen-reciprocal-packets}. It does
not introduce a physical mass/charge source. Set
$\varphi=(1+\sqrt5)/2$, use the ten vectors $u_{ij}$ printed in the
proof of Paper~I-B's ten-axis exact sequence, and perform all entries
in the field $\mathbb Q(\sqrt5)$.

For each unordered pair $v,w$ set $A_{vw}=1$ iff the labels are disjoint,
$S_{vw}=A_{vw}\operatorname{sgn}(u_v\cdot u_w)$, and
$M_v=u_v^\sigma u_v^{\mathsf T}$. These finite identities provide
short certificates for the central new ranks:
\[
 A^2=2I-A+\mathsf J_{10},\quad
 S^3=5S,\quad
 \sum_vM_v=0,\quad
 (\operatorname{tr}M_v^{\mathsf T}M_w)_{vw}
 ={8I+\mathsf J_{10}-6A\over9}.
\]
The matrix on the right has eigenvalues
$0^{(1)},(20/9)^{(4)},(2/9)^{(5)}$ and proves that the ten
rank-one maps span the whole nine-dimensional mirror-operator space.
On the oriented edge space, the four explicit orthogonal projectors
$\Pi_4,\Pi_5,\Pi_3,\Pi_{3'}$ prove
\[
 \operatorname{rank}(\Pi_4,\Pi_5,\Pi_3,\Pi_{3'})=(4,5,3,3),
 \qquad\sum_\alpha\Pi_\alpha=I_{15}.
\]
The exact coordinates of every axis, its Galois conjugate, and all
15 oriented disjoint-pair edges suffice to reproduce every matrix
without a discretization choice. Companion deterministic scripts
\path{Paper1A-PetersenHopfGeometry-v1_3-verify.py} and
\path{Paper1A-PetersenHopfGeometry-v1_4-verify.py} verify the
signed cycles, all 32 subset polar metrics, the cubic harmonic
portal and the source-free operator Gram identities. Exact-arithmetic
verification of these finite statements does not itself verify any
physical occurrence hypothesis.

For the reciprocal product of two marked fifteen-dimensional spaces,
write $r=(4,5,3,3)$. Then
\[
 \sum_{\alpha,\beta}r_\alpha r_\beta=225,\qquad
 \sum_\alpha{r_\alpha(r_\alpha+1)\over2}
 +\sum_{\alpha<\beta}r_\alpha r_\beta=120,
\]
\[
 \sum_\alpha{r_\alpha(r_\alpha-1)\over2}
 +\sum_{\alpha<\beta}r_\alpha r_\beta=105.
\]
The exchange dimensions are \emph{tensor-space} dimensions, not counts
of independent $2I$ group elements or automatically active
$M_8$ excitations. No conclusion in the geometry papers implies
an $SU(3)$ color identification, a physical $3'$ kinetic term, or a
hadron pole without subsequent same-owner action proofs.


\section{The fifteen-edge $A_5$-set and the two $\mathbb H P^2_{15}$ triangulations}
\label{p1c:sec:v340-hp2-petersen-equivariance}
\begin{theorem}[Exact vertex-set bridge, with its orientation obstruction]
\label{p1c:thm:v340-hp2-edge-orbits}
Let $E(P)$ be the fifteen \emph{unoriented} edges of the Petersen
graph $KG(5,2)$. The even-label group $A_5$ acts transitively with
edge stabilizer $V_4$, so $E(P)\cong A_5/V_4$ as an $A_5$-set.
The vertex-transitive Brehm--K\"uhnel triangulation
$K_{\rm BK}$ of $\mathbb H P^2$ has an $A_5$-action on fifteen
vertices with a conjugate $V_4$ stabilizer. Therefore, after a
choice of stabilizer alignment, there exists an $A_5$-equivariant
bijection
\[
 \boxed{E(P)\ \cong_{A_5}\ V(K_{\rm BK})}.
\]
The ordinary unordered pairs of Petersen edges form 105 elements
with four $A_5$-orbits of lengths $15,30,30,30$.
All fifteen-vertex $\mathbb H P^2$ triangulations are
five-neighborly, with the same initial face numbers
\[
(f_0,f_1,f_2,f_3,f_4)
=(15,105,455,1365,3003).
\]
The relation above identifies the \emph{vertex sets} and therefore
their complete one-skeleta, but does not identify their collections
of 8-simplex facets or the full simplicial complexes.
\end{theorem}
\begin{proof}
An edge of $KG(5,2)$ is a partition of four of the five letters
into two unordered pairs, with one unused letter. Its stabilizer
in $A_5$ has order four and is a Klein four group. The cited
Brehm--K\"uhnel action is vertex-transitive; orbit--stabilizer gives
the same order-four stabilizer, necessarily conjugate to $V_4$ in
$A_5$. Matching stabilizers yields the equivariant bijection.
The orbit decomposition of $\binom{E(P)}2$ is obtained by the
finite $A_5$ action and independently checked in the supplied
exact permutation certificate. Five-neighbourliness is a standard
property of minimal fifteen-vertex $\mathbb H P^2$ triangulations.
The Brehm--K\"uhnel $A_5$ example and Gaifullin's different
$C_6\times C_2$ example are both documented in
A.~A.~Gaifullin, \emph{New examples and partial classification of
15-vertex triangulations of the quaternionic projective plane}
(2024), \url{https://arxiv.org/abs/2311.11309}.
\end{proof}

\begin{proposition}[Unsigned versus oriented 105 states]
\label{p1c:prop:v340-hp2-local-coeff-obstruction}
The fifteen-dimensional \emph{unoriented} edge-permutation module is
\[
\mathbb R[E(P)]\cong\mathbf1\oplus\mathbf4\oplus2\mathbf5,
\]
whereas the \emph{oriented} Petersen edge-cochain module is
\[
W_P\cong\mathbf4\oplus\mathbf5\oplus\mathbf3\oplus\mathbf3'.
\]
For the diagonal $A_5$ action,
\[
 \dim(\Lambda^2 W_P)^{A_5}=0,
 \qquad
 \dim\bigl(\mathbb R[\tbinom{E(P)}2]\bigr)^{A_5}=4.
\]
Thus the numerical equality of 105 simplicial edges and
$\dim\Lambda^2W_P=105$ is not an equivariant identification.
The orientation-sign local coefficient is an actual obstruction.
Gaifullin's separately labelled $K_0$ with automorphism group
$C_6\times C_2$ must \emph{not} be renamed $K_{\rm BK}$.
Nor is $C_3\times C_2\cong S_3$: the direct product is $C_6$,
while $S_3\cong C_3\rtimes C_2$.
\end{proposition}
\begin{proof}
The unsigned action is induced from the trivial character of
$V_4$; the oriented action is induced from the nontrivial
one-dimensional sign character recording reversal of an edge.
The $A_5$ character table gives the two stated decompositions.
Distinct real irreducibles have no invariant alternating bilinear
pairing, and each multiplicity in $W_P$ is one; therefore
$\Lambda^2 W_P$ has no singlet. The four orbits of unsigned
edge pairs yield four invariant orbit sums. These statements
refer only to the one-skeleton, not to 8-dimensional topology.
\end{proof}


\appendix
\section{Spectral, sampling, and three-manifold comparisons}
\label{p1:app:mathematical-appendices}



\subsection{Spectral comparisons of the rooted trimer embeddings}
\label{p1:app:trimer-overlap-spectra}
This subsection records principal angles and cyclic multiplicities
of the maps already constructed in
Theorem~\ref{p1:thm:rooted-trimer-vector-overlap} and
Corollary~\ref{p1:cor:rooted-trimer-bivector-overlap}.
It introduces no additional geometric choices.

The overlap eigenvalues on \(Q_0\oplus Q_\perp\) are
\[
 K_V:(1/2,-1),\qquad K_{\Lambda^2}:(1,-1/2),
\]
where the second eigenvalue has multiplicity two.
The singular values are therefore \((1/2,1,1)\) and
\((1,1/2,1/2)\), respectively. Thus the principal angles of the
two vector images are \((\pi/3,0,0)\), and those of the two
exterior-square images are \((0,\pi/3,\pi/3)\).
Both vector images contain \(P\); the exterior-square images
meet in the line \(\Lambda^2P\).

The positive polar decompositions of
\(C_\triangle,K_V,K_{\Lambda^2}\) have the same orthogonal factor
\[
 S=Q_0-Q_\perp,
\]
but respective positive factors
\[
 |C_\triangle|=2Q_0+Q_\perp,\qquad
 |K_V|=\tfrac12Q_0+Q_\perp,\qquad
 |K_{\Lambda^2}|=Q_0+\tfrac12Q_\perp.
\]
These formulas follow directly by squaring on the two orthogonal
projectors. Passing to the polar factor discards the different
singular-value data of the incidence and overlap maps.

For the common rooted cycle, let
\(A=\operatorname{span}_{\mathbb R}\{1,(0,u)\}\), so that
\[
 V=A\oplus P,\qquad
 \Lambda^2V=\Lambda^2A\oplus(A\otimes P)\oplus\Lambda^2P.
\]
Writing the complex cyclic characters in the order
\((1,\omega,\overline\omega)\), their multiplicities are
\[
\begin{array}{c|ccc}
 &1&\omega&\overline\omega\\\hline
 G_{\mathbb C}&1&1&1\\
 P_{\mathbb C}&0&1&1\\
 V_{\mathbb C}&2&1&1\\
 (\Lambda^2V)_{\mathbb C}&2&2&2
\end{array}
\]
Indeed, the cycle fixes \(A\), rotates \(P\) by \(2\pi/3\),
and fixes the two real exterior lines \(\Lambda^2A\) and
\(\Lambda^2P\). The mixed summand supplies two copies of the
rotation plane. These statements concern the indicated cyclic
restriction, not the full rooted operation algebra.

\subsection{Sampling sharpness and a geometric positive-polar reference}
\label{p1:app:v203-sampling-reference}
This appendix supplies a quadrature diagnostic and an optional geometric
spectral comparison of two commuting finite operators.

\begin{proposition}[First quadratic sampling discrepancy]
\label{p1:prop:v203-quadratic-alias-witness}
Uniform Hurwitz quadrature on the comparison band has the exact discrepancy
\[
 \frac1{24}\sum_{q\in2T}q_0^2(q_0^2-q_1^2)^2
 -\int_{S^3}q_0^2(q_0^2-q_1^2)^2\,d\sigma=\frac1{48}.
\]
Thus the bound for linear symbols cannot be raised to include all
quadratic symbols without the correction above.
\end{proposition}
\begin{proof}
The discrete average is $1/12$. The round moments are
$\int q_0^6=5/64$ and $\int q_0^4q_1^2=\int q_0^2q_1^4=1/64$,
so the continuum integral is $1/16$. Subtraction gives $1/48$.
The packaged exact matrix check additionally records rank nine for the
whole discrepancy of $f=q_0^2$, and tests its Hopf covariance failure.
Those finite checks supplement the displayed nonzero witness.
\end{proof}

\begin{corollary}[Dimensionless spectral comparison for the right Hopf generator]
\label{p1:cor:v203-geometric-positive-polar}
Write $R=\Rmom$. Let $N_H$ be the integer-weight right Hopf generator on $\mathcal H_1$ and
put $E_R=R\sqrt{K_{\rm LC}}$, using
Corollary~\ref{p1:cor:canonical-retained-dirac-square}. Then
\[
 [E_R,N_H]=0,\qquad E_R\ge|N_H|,\qquad
 N_\pm=(E_R\pm N_H)/\sqrt2\ge0.
\]
Each $N_\pm$ has a two-dimensional kernel. These identities transport
exactly to $\mathcal V_R$ through $\mathsf S_RJ_R$.
\end{corollary}
\begin{proof}
The square is Spin(4)-invariant. The sector order is
\[
 (0,0),(1/2,1/2),(1,0),(0,1),(1/2,0),(0,1/2).
\]
In this order, the maximal absolute right Hopf weights
$0,1,0,2,0,1$ are bounded by the $E_R$ eigenvalues
$0,\sqrt3,2,2,3/2,3/2$. Equality occurs only on the scalar and the
appropriate one-dimensional extreme weight in $(0,1)$. Isometric transport
preserves the spectrum and order.
\end{proof}

\subsection{Sharpness witnesses for the exceptional complete-band frontier}
\label{p1:app:v181-exceptional-frontier-witnesses}
The negative part of Theorem~\ref{p1:thm:v181-exceptional-complete-band-frontier}
is witnessed by the following moment computations.  On $2T$,
$f=x_0^3$ gives
\[
 \langle f^2\rangle_{2T}=\frac3{32}\ne\frac5{64}
 =\langle f^2\rangle_{S^3},
\]
while $f=4x_0^2-1$ gives respectively
$\langle f^3\rangle_{2T}=2\ne1$ and
$\langle f^4\rangle_{2T}=7\ne3$.
On $2O$, $f=x_0^4$ gives
\[
 \langle f^2\rangle_{2O}=\frac{15}{256}\ne\frac7{128},
\]
$f=x_0^3+x_0^2$ gives
\[
 \langle f^3\rangle_{2O}=\frac{65}{256}\ne\frac{31}{128},
\]
and $f=4x_0^2-1$ gives
$\langle f^4\rangle_{2O}=4\ne3$.
On $2I$, the already computed twelfth moment
\[
 \langle x_0^{12}\rangle_{2I}=\frac{133}{4096}
 \ne\frac{132}{4096}=\langle x_0^{12}\rangle_{S^3}
\]
is obtained as $f^2,f^3,f^4$ for $f=x_0^6,x_0^4,x_0^3$, respectively.
Each profile belongs to the first complete band above the corresponding
entry in the main-text table.  These witnesses establish maximality of those
complete-band frontiers; they do not assert failure of every individual
higher-band polynomial.

\subsection{Comparison with Other 3-Manifolds}
\label{p1:app:comparison}

\begin{table}[htbp]
\centering
\small
\begin{tabular}{|l|l|l|l|}
\hline
\textbf{Manifold} & $\pi_1$ & \textbf{Parallelizable} & \textbf{Einstein} \\
\hline
$S^3$ & $0$ & Yes & Yes (round) \\
$\RR P^3$ & $\ZZ_2$ & Yes & Yes (quotient of round) \\
$S^2 \times \Sone$ & $\ZZ$ & Yes & No \\
Lens space $L(p,q)$ & $\ZZ_p$ & Yes  & Yes (quotient) \\
Nil manifold & Discrete Heisenberg $\Gamma$ & Yes & No \\
\hline
\end{tabular}
\caption[Comparison with compact orientable three-manifolds]{Comparison with compact orientable three-manifolds.}
\label{p1:tab:manifold-comparison}
\TableNotes{The round three-sphere is distinguished here by simple connectivity and its constant-positive-curvature geometry. Parallelizability is included for comparison, but it is not unique among orientable three-manifolds; indeed, all orientable three-manifolds are parallelizable.}
\end{table}


\subsection{Trace factorization and Hilbert--Schmidt pairing}
\label{p1:app:trace-toolkit}

\begin{theorem}[Hilbert--Schmidt pairing and conjugation invariance]
\label{p1:thm:HS-invariance}
Let $V$ be a finite-dimensional complex Hilbert space. Define the Hilbert--Schmidt pairing
\[
\langle A,B\rangle_{\mathrm{HS}}:=\Tr_V(A^\dagger B),\qquad A,B\in\End(V).
\]
Then $\langle \cdot,\cdot\rangle_{\mathrm{HS}}$ is unitarily invariant: for every $U\in U(V)$,
\[
\langle UAU^{-1},\,UBU^{-1}\rangle_{\mathrm{HS}}=\langle A,B\rangle_{\mathrm{HS}}.
\]
\end{theorem}

\begin{proof}
$\Tr((UAU^{-1})^\dagger(UBU^{-1}))=\Tr(U A^\dagger U^{-1} U B U^{-1})=\Tr(U A^\dagger B U^{-1})
=\Tr(A^\dagger B)$.
\end{proof}

\begin{theorem}[Tensor trace factorization]
\label{p1:thm:tensor-trace}
Let $V,W$ be finite-dimensional complex vector spaces, and let $A\in\End(V)$. Then
\[
\Tr_{V\otimes W}(A\otimes \mathbf 1_W)=\dim(W)\,\Tr_V(A).
\]
More generally, for $A\in\End(V)$ and $B\in\End(W)$ one has
\[
\Tr_{V\otimes W}(A\otimes B)=\Tr_V(A)\,\Tr_W(B).
\]
\end{theorem}

\begin{proof}
Choose bases and compute on elementary tensors; both identities are standard.
\end{proof}

\begin{theorem}[Commutant trace action on multiplicity spaces]
\label{p1:thm:commutant-trace-mult}
Let $V\cong \bigoplus_\pi \CC^{m_\pi}\otimes V_\pi$ be the isotypic decomposition of a unitary $G$-module as
in Theorem~\ref{p1:thm:isotypic-commutant}. Then for $A\in\End_G(V)$ corresponding to blocks
$A_\pi\in M_{m_\pi}(\CC)$ one has
\[
\Tr_V(A)=\sum_{\pi} \dim(V_\pi)\,\Tr_{\CC^{m_\pi}}(A_\pi).
\]
\end{theorem}

\begin{proof}
On each isotypic block, $A$ acts as $A_\pi\otimes \mathbf 1_{V_\pi}$. Apply Theorem~\ref{p1:thm:tensor-trace}.
\end{proof}




\section{Coset--Hodge and six-axis parity certificate ledger}
\label{p1c:sec:v341-petersen-bridge-certificate}
\begin{proposition}[Exact rank ledger and non-identification of four fifteens]
\label{p1c:prop:v341-coefficient-rank-and-topology}
For binary counting measure on a marked $2I$ packet,
$\operatorname{rank}B=17$, while the symmetric and skew ten-column
indicator maps have ranks $10$ and $6$ and orthogonal images.
The two conjugation-invariant character lines in
$\operatorname{im}B$ leave the rank-fifteen coefficient carrier
$\operatorname{im}\Pi_{15}$.  The three independent 15-dimensional
objects remain: the real $\operatorname{im}\Pi_{15}\cong W_P$,
the \emph{complex} momentum module $H_1$, and the 24-cell
Bose--Mesner block.  The equivalence of $\operatorname{im}\Pi_{15}$ with $W_P$
is the explicitly typed real finite intertwiner of Paper~I-B,
not a complex-carrier identification.

Separately, the Brehm--K\"uhnel $A_5$-triangulation has a
15-element \emph{unoriented} vertex set that can be matched to
the 15 \emph{unoriented} Petersen graph edges after stabilizer
alignment (Theorem~\ref{p1c:thm:v340-hp2-edge-orbits}).
This does not turn its 105 simplicial edges into the oriented
$\Lambda^2W_P$ exchange-odd 105-space, nor does it identify the
separate Gaifullin $C_6\times C_2$ example with that
$A_5$ triangulation.  In particular no $\mathbb H P^2$
topological calibration is a physical source operator.
\end{proposition}
\begin{proof}
The exact integer identities are recorded in
Theorem~\ref{p1:thm:v341-coset-gram-hodge}, and the coefficient
projector is Corollary~\ref{p1:cor:v341-coset-pi15}.
The two double covers and all 120 label-parity checks are in
Theorem~\ref{p1:thm:v341-desargues-rp2-parity}.
The independent executable certificate and its JSON output are
\path{Proofs/Paper00-PetersenBridges-v1_3.py} and
\path{Proofs/Paper00-PetersenBridges-v1_3-output.json}.
The previously proved unsigned-versus-oriented $105$ obstruction is
Proposition~\ref{p1c:prop:v340-hp2-local-coeff-obstruction}.
\end{proof}


\section{Three-cycle overlap, cycle-space and Lie-word certificates}
\label{p1c:sec:v342-threecycle-proof-ledger}
\begin{proposition}[Finite certificate domain and proof-status separation]
\label{p1c:prop:v342-threecycle-exact-certificate}
The signed inverse-pair construction of Paper~I-B is evaluated
on the $60$ even permutations of five labels, in the rational
$5\times5$ zero-sum model.  It certifies all ten-by-ten Gram
entries, all ten-by-fifteen signed triangle--Petersen incidences,
the $15$-dimensional frame projectors, the exact two five-cycle
class splits over $\mathbb Q(\sqrt5)$, the four orbit types of
three-generator subsets, and the Lie-word depth sequence
$(3,3,5,3,1)$.  The original $120$-vertex binary support is
obtained by the centre-even twofold pullback with row
normalization $1/\sqrt2$; the $60$ and $120$ counting
conventions cannot be silently mixed.  The spectral frame
on $\mathfrak{sl}(4,\mathbb R)$ is not an effect-valued
measurement, and the $60$ path triples do not specify an
unmarked distinguished physical path.
\end{proposition}
\begin{proof}
The exact independent SymPy script
\path{Proofs/Paper00-ThreeCycleReconstruction-v1_2.py}
checks the rational Gram/commutator identities, integer support
matrix and singular-square $10$, and the class spectral projectors
including $D_5^2=(16/45)P_6$.  The associated JSON report records
its checks.  For the earlier $15$-mode edge matrices and
rank-$17$ transporter image consult
Theorem~\ref{p1:thm:v339-petersen-sl4-isometry} and
Corollary~\ref{p1:cor:v341-coset-pi15}.  No statement of the
script checks a physical $2I$ occurrence, an action Hessian,
experimental measurement or GP3 comparison.
\end{proof}


\section{Signed reciprocal tensor, parity twist and quadratic Casimir certificate}
\label{p1c:sec:v343-signed-coupling-proof-ledger}
\begin{proposition}[Exact finite verification and nonidentification of the Lie algebras]
\label{p1c:prop:v343-signed-coupling-certificate}
The finite proof certificate
\path{Proofs/Paper00-ReciprocalSignedCoupling-v1_1.py}
constructs the signed oriented fifteen-edge representation from
all $120$ permutations of $S_5$, checks the two five-cycle
class sums and polynomial $K_W^3=80K_W$, and proves the
$225$-dimensional $\mathcal C_{xp}$ exchange multiplicities
of Theorem~\ref{p1b:thm:v343-doubleodd-contrast}.
It checks the $S_5$ characters, the $15+90$ projector in
Theorem~\ref{p1b:thm:v343-canonical-selected15} and the
$\mathfrak{so}(3)\oplus\mathfrak{so}(3)$ adjoint Casimir
in Theorem~\ref{p1b:thm:v343-s5-twist-so6}.
The independent Paper I-B frame calculation fixes the
class normalization $D_5^2=(16/45)P_6$; it does not make
$\mathfrak{sl}(4,\mathbb R)$ and $\mathfrak{so}(6)$
isomorphic as real Lie algebras.
\end{proposition}
\begin{proof}
The conjugacy-class difference is scalar on each $A_5$ irreducible;
character orthogonality determines its eigenvalues, and the signed
edge matrices are orthogonal.  Exchange traces follow from
$\operatorname{tr}(\tau(A\otimes B))=\operatorname{tr}(AB)$.
For the selected real skew matrices $Z\in\mathfrak{so}(6)$,
conjugating by $\operatorname{diag}(I_3,-I_3)$ is an involutive
Lie automorphism.  Its $+/-$ eigenspaces have dimensions $6/9$,
with direct-adjoint Casimir eigenvalues $2/4$ in the stated
skew-generator normalization.  The script's machine-readable
results are archived as
\path{Proofs/Paper00-ReciprocalSignedCoupling-v1_1.json}.
An orthogonal $A_5$-module isomorphism can preserve neither
Lie bracket nor real form automatically, and none is assumed.
\end{proof}


\section{Finite certificate for the twisted Hurwitz--Petersen bridge}
\label{p1c:sec:v344-hurwitz-petersen-certificate}
\begin{proposition}[Reproducibility and exact/conditional separation]
\label{p1c:prop:v344-hurwitz-petersen-verification}
The executable certificate
\path{Proofs/Paper00-HadronScaleTest-v1_7-certificate.py},
with accompanying JSON output and an explicit $10\times5$ complex
intertwiner, verifies the group orders $|2I|=120$, $|2T|=24$,
the order-three character kernel, the exact twisted Gram
$F^\dagger F=24I_{25}$, the vanishing untwisted/twisted cross
Gram, the rank-$17$ untwisted coefficient image, and the
$P_{\rm Pet}$ isometry of
Proposition~\ref{p1:prop:v344-petersen-quintet-intertwiner}.
It also enumerates $96$ unoriented and $192$ oriented Hurwitz
edges and the eight free oriented-edge orbits.  The polynomial
$SU(3)$ one-link witness checks nonvanishing of the already
existing Paper~VII-A polarized determinant after its $C_3$
Fourier projection; it makes no assertion about an executed
physical Hamiltonian, source residue or hadronic mass curvature.
\end{proposition}
\begin{proof}
The character and finite-incidence identities are verified by exact
cyclic-phase exponent counts over the relevant finite groups.
The intertwiner is checked numerically to a maximum residual below
$1.2\times10^{-15}$ in the supplied convention.  The gauge witness
reduces the determinant to $(1+2\cos\theta)/3$ on the affected
edge fans and checks an exact rational total mode norm $4/9$ at
$\theta=\pi$.  Its extension to actual admissible link domains,
physical states and source correlators is not certified by this script.
\end{proof}


\clearpage
\section{Exact F40 arrow and cycle reconstruction certificate}
\label{p1c:sec:v345-f40-proof}
\begin{proposition}[Full intrinsic readout and source-boundary certificate]
\label{p1c:prop:v345-f40-certificate}
The reproducible source
\path{Proofs/Paper00-PetersenHodgePlucker-v1_11-verify.py},
run with its v1\_10 predecessor, constructs the actual forty
unit-quaternion vertices, their sixty directed edges, all sixty
proper five-label permutations and all five tetrahedral cores.
It checks in exact integer/rational and $\mathbb Q(\sqrt5)$
arithmetic: the edge-to-$(g,m)$ bijection, the 60-element regular
$A_5$ action and five-cycle core transport; the ranks $11,16,21$;
the frame polynomial
\[
 (z-1)^5(z^2-12z+31)^3(z^2-8z+6)^5;
\]
the balanced arrow projection
$AP_ZA^{\mathsf T}=2P_{\rm sym}+P_{\rm skew}$;
and the mixed quaternionic left/right Gram formula of
Proposition~\ref{p1:prop:v345-f40-balanced-mirror}.
The predecessor tests separately check the twelve-shell transport
and the three binary-deck visibility sectors.  The mathematical
certificate does not test a physical QCD source embedding or a
same-action hadronic spectrum.
\end{proposition}
\begin{proof}
The scripts directly reconstruct the finite incidence matrices and
compute exact ranks, products and characteristic polynomials;
the printed output records all passing comparison groups.  All
source and physical distinctions stated above remain separate
from the algebraic equality certificates.
\end{proof}

\clearpage
\section{The exact 44-channel compatibility and rational inverse certificate}
\label{p1c:sec:v349-f40-compatibility}
\begin{theorem}[Twenty-three necessary and sufficient native F40 consistency equations]
\label{p1c:thm:v349-f40-23-equations}
Let $u_g,v_g$ be arbitrary real right/left 12-shell coordinates,
$w_{ij}$ arbitrary directed 20-core-arrow coordinates,
and let $d_g(i,j)$ be the cyclic distance of
Theorem~\ref{p1:thm:v349-f40-mirror-gluing}.
There exists a \emph{unique} F40 cycle $z$ with
$(u,v,w)=(T_Rz,T_Lz,Az)$ if and only if all four classes
of linear conditions hold:
\begin{enumerate}
\item The two shell sums vanish:
$\sum_g u_g=\sum_gv_g=0$.
\item The five arrow row sums and any four arrow column sums
vanish: $\sum_{j\ne i}w_{ij}=0$ and
$\sum_{i\ne j}w_{ij}=0$.  These give nine independent equations.
\item The shell sum is inversion-even:
$u_g+v_g=u_{g^{-1}}+v_{g^{-1}}$
for one member of each of the six inverse pairs.
\item For the six chords $2\leq i<j\leq5$,
\[
 \frac{w_{ij}-w_{ji}}2
 =\sum_g\frac{5-2d_g(i,j)}{40}(u_g-v_g).
\]
\end{enumerate}
These conditions have rank $2+9+6+6=23$.
Writing $\mathcal R=(T_R,T_L,A)|_{\mathcal Z_{40}}$,
the $44\times21$ native frame has rank $21$, and
\[
 \operatorname{im}\mathcal R=\ker C_{23},\qquad
 \dim\ker C_{23}=21.
\]
\end{theorem}
\begin{proof}
Native cycle conservation gives the two shell sums and
doubly-balanced arrows.  The shell parity and skew-arrow
relations follow from the right/left and Hodge-mirror
identities.  An exact integer/rational construction of the
$23\times44$ constraint matrix $C_{23}$ and the
$44\times21$ cycle readout $\mathcal R$ verifies
$C_{23}\mathcal R=0$, $\operatorname{rank}C_{23}=23$ and
$\operatorname{rank}\mathcal R=21$.  The rank-nullity theorem
therefore proves necessity, sufficiency and uniqueness.
The full finite construction is reproducible with
\path{Proofs/Paper00-F40MirrorGluing-v1_1-verify.py};
the exact matrices are provided in the proof package.
\end{proof}

\begin{proposition}[Exact polynomial inverse of the F40 frame]
\label{p1c:prop:v349-f40-rational-frame-inverse}
For $F=\mathcal R^*\mathcal R$ on $\mathcal Z_{40}$,
\[
 \chi_F(x)=(x-1)^5(x^2-12x+31)^3(x^2-8x+6)^5,
\]
and its five distinct spectral values are the roots of
$m_F(x)=x^5-21x^4+153x^3-453x^2+506x-186$.
Hence
\[
 F^{-1}=\frac{F^4-21F^3+153F^2-453F+506I}{186},
 \qquad P_{44}=\mathcal R F^{-1}\mathcal R^*.
\]
$P_{44}$ is the orthogonal rank-21 projector onto the
geometrically compatible channel packets and
$z=F^{-1}\mathcal R^*(u,v,w)$ reconstructs a compatible
cycle without an eigensolver or numerical inverse.
\end{proposition}
\begin{proof}
The exact v2.345 frame spectrum yields the characteristic
and minimal polynomials.  Cayley--Hamilton on each eigenspace
implies $m_F(F)=0$ and therefore the displayed inverse.
The projector properties follow from the full-column-rank
frame factorization.
\end{proof}

\begin{proposition}[Research certificate provenance and strict limits]
\label{p1c:prop:v349-comparison-proof-register}
The distributed F40 module, mirror-gluing, twisted-harmonic,
and multi-gap certificate scripts independently check the relevant exact finite matrices, character
ranks, graph constraints, harmonic moments and invariant tensor
coefficients within their declared domains.  Their finite
checks supplement, rather than replace, the analytic
representation and cubature proofs.  None certifies a physical
$SU(3)$-singlet source, Hamiltonian leakage, or mass pole.
\end{proposition}


\clearpage
\section{Spinor Fischer decomposition and a cumulative comparison filtration}
\label{p1c:sec:v357-clifford-fischer}
\begin{theorem}[Exact homogeneous spinor-polynomial sequence]
\label{p1c:thm:v357-spinor-fischer}
Let $F\cong V_1$ be the complex three-dimensional vector
representation of a specified diagonal $\mathrm{Spin}(3)$ and
$S=V_{1/2}$. This identification with the $F_\sigma$ of Paper~I-B
is an \emph{additional, explicitly typed equivariance hypothesis},
not a consequence of the oriented $C_3$ trimer alone. Put
$P_n=\Sym^nF$, and give polynomial spinors the Fischer inner
product.  With Pauli matrices define
\[
 C=\boldsymbol\sigma\cdot\mathbf q,\qquad
 D=\boldsymbol\sigma\cdot\nabla_{\mathbf q},\qquad
 C=D^\dagger,\qquad
 N=\mathbf q\cdot\nabla_{\mathbf q}.
\]
For every $n\ge1$, $D_n:P_n\otimes S\to P_{n-1}\otimes S$
is onto, and its kernel is the irreducible highest-spin module
$M_n\cong V_{n+1/2}$. There is a rotation-equivariant exact sequence
\[
 0\longrightarrow M_n\longrightarrow P_n\otimes S
 \xrightarrow{D_n}P_{n-1}\otimes S\longrightarrow0,
\]
and an orthogonal decomposition
\[
 P_n\otimes S=C(P_{n-1}\otimes S)\oplus M_n.
\]
\end{theorem}
\begin{proof}
Since $(\boldsymbol\sigma\cdot\mathbf q)^2=
(\mathbf q\cdot\mathbf q)I$, multiplication by $C$ is injective
on polynomial spinors. Its Fischer adjoint $D_n$ is therefore
surjective. The kernel has dimension
$2\binom{n+2}{2}-2\binom{n+1}{2}=2n+2$.
The Clebsch--Gordan decomposition
$P_n\otimes S\cong(P_{n-1}\otimes S)\oplus V_{n+1/2}$
as $\mathrm{Spin}(3)$ modules identifies the kernel with the unique
highest-spin summand.  The image of $C=D^\dagger$ is the orthogonal
complement of the kernel, proving the second decomposition.
\end{proof}

\begin{theorem}[Explicit nested rotation-and-parity comparison supports]
\label{p1c:thm:v357-spinor-filtration}
On the direct-sum spinor-polynomial comparison space
$\mathscr P=\bigoplus_{n\ge0}P_n\otimes S$, define
\[
 \mathscr F_r=\begin{cases}
 \bigoplus_{n=0}^{r}P_n\otimes S,&0\le r\le2,\\
 \bigl(\bigoplus_{n=0}^{r-1}P_n\otimes S\bigr)\oplus M_r,&r\ge3.
 \end{cases}
\]
Then $\mathscr F_r\subset\mathscr F_{r+1}$ and the cumulative
complex dimensions through degree $7$ are
\[
 2,8,20,28,50,82,126,184.
\]
For $r\ge3$ there is a rotation-equivariant linear isomorphism
$\mathcal N_r\otimes S\cong\mathscr F_r$, provided Paper~I-A's
normal--vector decomposition is identified with $V_0\oplus V_1$.
The increment for $r\ge4$ is
\[
 \mathscr F_r\ominus\mathscr F_{r-1}
 \cong C(P_{r-2}\otimes S)\oplus M_r,
\]
with ordinary polynomial parity $(-1)^{r-1}$ on the first and
$(-1)^r$ on the second summand.  The isomorphism from the Hurwitz
quotient is a \emph{rotation} isomorphism; it is not asserted to
preserve an unprovided physical parity assignment or QCD source map.
\end{theorem}
\begin{proof}
By Paper~I-B's boundary definition,
$\mathcal N_r\cong\bigoplus_{k=0}^{r-2}P_k\oplus P_r$ for $r\ge3$.
Apply the preceding Fischer sequence to replace
$P_r\otimes S$ as a rotation module by
$(P_{r-1}\otimes S)\oplus M_r$.  This identifies the resulting
representation with $\mathscr F_r$.  The filtration inclusion
holds since $M_r\subset P_r\otimes S$, while the complement
follows from the orthogonal Fischer decomposition in degree $r-1$.
The dimension is
$2\binom{r+2}{3}+2r+2$ for $r\ge3$, and direct substitution
gives the displayed sequence. Polynomial inversion gives the
parities stated; no parity equivariance of the quotient isomorphism
was used.
\end{proof}
\begin{remark}[Mathematical comparison, not physical nuclear selection]
\label{p1c:rem:v357-spinor-filtration-scope}
The comparison filtration is an explicitly nested representation
realization of the capacity sequence. It is not an identification of
$\mathscr F_r$ with the source image of the actual nuclear QCD
Hamiltonian and does not turn degreewise ranks into experimentally
established physical closures. Its spin-coupled multiplicity is
logically different from an assumption of unsplit orbital levels
with independent twofold spin occupation.
\end{remark}


\section{Exact 600-cell and Hurwitz fan incidence verification}
\label{p1c:sec:v365-hurwitz-fan-certificate}
The checker \path{Proofs/Paper00-HurwitzHadronFiniteCheck-v1_1.py}
constructs the unit 600-cell from eight coordinate-axis roots,
sixteen half-sign roots and ninety-six even permutations of
$\bigl(0,\tfrac12,\tfrac{\varphi}{2},\tfrac1{2\varphi}\bigr)$,
with $\varphi=(1+\sqrt5)/2$ and independent signs on the nonzero
coordinates.  Distinct vertices are joined precisely for inner
product $\varphi/2$, and four-cliques form the 600 tetrahedra.
The 24 initial Hurwitz vertices form a 96-edge, 96-face subcomplex
with inner product $1/2$ adjacency.  The checker constructs the
$600\times120$ incidence $T$ and the $96\times96$ unsigned face
incidence $M$ and checks
\[
T^{\mathsf T}T=20I+5A_{\rm vert},\quad
\operatorname{rank}_{\mathbb F_{101}}T=120,\qquad
M^{\mathsf T}M+L_F=9I,
\quad\operatorname{rank}_{\mathbb F_{101}}M=73.
\]
The 23-dimensional real nullity of $M$ also follows independently
from the $S^3$ cellular chain ranks proved in
Theorem~\ref{p1:thm:v235-hurwitz-chain-complex}.  A modular lower
bound alone would not prove an \emph{upper} bound on real rank;
that upper bound is supplied by the topological argument.

\begin{remark}[Unverified advanced chamber and Fourier packets]
\label{p1c:rem:v365-hurwitz-candidate-quarantine}
The G05--G08, H03--H08 and R01--R07 chamber spectral/spin
packages in the associated Hurwitz--hadron integration proposal
are research candidates rather than added propositions here.
Likewise, no independent source-matched replay of the proposed
25-chamber complement $U$, its $41$-register closure, cubic
$16/32$ blocks, or $274/326$ decomposition is certified by the
above finite checker.  Those packages cannot supply QCD pole
energies or action-owned source maps without additional proofs.
The $C_6$ regular orbit and $\operatorname{Sym}^2(V)$ quadratic
space are not interchangeable as $C_6$ representations merely
because both have dimension six.
\end{remark}


\section{Chamber--transporter proof certificate and conditional source ledger}
\label{p1c:sec:v366-hurwitz-chamber-certificates}
Theorem~\ref{p1:thm:v366-chamber-transporter-complement}
upgrades the specifically displayed \emph{finite} candidate identities of
Remark~\ref{p1:rem:v365-chamber-register-candidate}, but does not
upgrade the action-dependent hadronic interpretation.  The certificates
below concern a fixed marked unit quaternion $2T\subset2I$; a theorem
about an unmarked arbitrary physical realization would require an
additional matching map.

\begin{proposition}[Exact integer-pair coordinate certificate]
\label{p1c:prop:v366-exact-coordinate-chamber}
Represent each coordinate of a $2I$ unit quaternion as
$(a+b\sqrt5)/4$ with $a,b\in\mathbb Z$.  Form the eight signed
coordinate roots, sixteen half-sign roots, and ninety-six even
permutations of
\[
 (0,\tfrac12,\tfrac{1+\sqrt5}{4},
            \tfrac{\sqrt5-1}{4})
\]
with independent signs on the three nonzero entries.  Two roots
are adjacent if their \emph{exact} inner product is
$(1+\sqrt5)/4$, and the tetrahedral facets are the four-cliques.
Conjugate the two $2T$ projections of a tetrahedron stabilizer
into the fixed embedded $H$ and construct the 25 double-coset
chamber sets as in Theorem~\ref{p1:thm:v366-chamber-transporter-complement}.
The executable
\path{Proofs/Paper00-HurwitzHadronChamberExact-v1_1.py}
checks the following entrywise \emph{over integers}:
\[
 \begin{gathered}
 |G|=120,\quad |E|=720,\quad|\mathscr F|=600,\quad
 \mathscr F=\bigsqcup_{25}\mathscr F_{ij},\quad
 |\mathscr F_{ij}|=24,\\
 T^{\mathsf T}U=J-B,\quad
 TB=J-(I+A)U,\quad
 U^{\mathsf T}AU=6R,\quad
 T^{\mathsf T}T=20I+5A_{\rm vert},\\
  (A^2+A-5I)T\Pi_{15}=0,\qquad
  (A^2+A-5I)U(Q\otimes Q)=0.
 \end{gathered}
\]
The auxiliary $3000\Pi_{15}$ is an \emph{integer matrix} and
satisfies $\Pi_{15}^2=\Pi_{15}$ exactly.  In this certificate,
there is no float-based indexing, no tolerance for root matching,
and no imported physical dynamics.
\end{proposition}
\begin{proof}
All coordinate products and dot products reduce to pairs of
integers by the multiplication rule
$(a,b)(c,d)=(ac+5bd,ad+bc)$ before division by the fixed
coordinate denominator.  In the quaternion product the resulting
numerators are divisible by four and return to the same exact bank.
The complete adjacency matrix, all four-cliques, the conjugated
stabilizer and all $3000$ incidence entries are then checked by
explicit finite loops; the projector numerator is multiplied as
integers.  The executable verification is finite and deterministic.
\end{proof}

\begin{proposition}[Further finite harmonic and group certificates]
\label{p1c:prop:v366-cumulative-hurwitz-certificates}
The source-matched supplemental executable
\path{Proofs/Paper00-HadronChamberSplitCertificate-v1_7.py}
verifies the $41$-register closure, full local Hurwitz Fourier
filters, global signed rank $44$, cubic coefficient rank $16$ and
face-walk dimension $32$, and the complete $274+326$ two-sided
representation ledger.  The separate
\path{Proofs/Paper00-HadronChamberSplitCertificate-v1_9.py}
checks the $A_5$ membership classification, the $3/8$ cross-core
multiplication incidence and the rank-three kernel.  The
\texttt{v1\_10} script checks spin-$2\otimes1/2$ projectors,
$2/5:3/5$ factorized source norms and raw meson/baryon
representation products, while the \texttt{v1\_12} script checks
$C_3$-Hessian diagonalization, finite Wilson-source \emph{linear}
independence and conditional angular distributions.
For the cumulative $v1\_7$ replay, the original coordinate labels
are initially matched numerically and the subsequent rank,
annihilator, Gram and character identities use exact integer,
rational or $\mathbb Q(\sqrt5)$ arithmetic; the independent
integer-pair source-matched check above removes this potential
indexing ambiguity from the core transporter/face identities.
The advanced Fourier and spectrum packets have reproducible finite
certificates in the declared marked model; they are not a
whole-framework $E_4$ physical realization proof.
\end{proposition}
\begin{proof}
The certificates enumerate all elements of the finite quaternion
and permutation banks, use exact identity/pivot checks when
available, and report the completed rank and character relations.
The distinct representation multiplicities are also obtained by
restriction-character inner products.  Neither a dimension equality
nor a numerical agreement with a pre-existing QCD coefficient is
accepted as a physical-operator intertwiner.
\end{proof}

\begin{remark}[Verification and naming firewalls]
\label{p1c:rem:v366-hurwitz-source-corrections}
The \emph{rooted} nearest-eight set, the normal subgroup $Q_8$,
and the eight-dimensional twisted Fourier coefficient sector are
distinct objects.  The Petersen rank-$15$ source and the intrinsic
$15$-dimensional momentum module have distinct types.  A nonzero
eight-by-eight evaluation determinant proves \emph{linear}
independence of the eight Wilson functions, \emph{not algebraic}
independence.  The equal numerical factor $1/(6\sqrt2)$ in a local
signed incidence normalization and an older QCD expression is not
an equality of physical operators.  None of the computed source
projector weights, graph eigenvalues or spectrum dimensions is an
observable decay probability, hadron mass or irreversible width.
\end{remark}

\section{Exact certificate for the golden $E_8$--Hurwitz comparison}
\label{p1c:sec:v370-e8-comparison-certificate}
The constructions in Paper~I-B,
Theorems~\ref{p1:thm:v370-e8-d4-triality} and
\ref{p1:thm:v370-e8-shell-chamber}, and
Proposition~\ref{p1:prop:v370-e6-cubic-support} are purely finite
mathematical consequences of its \emph{previously specified}
$\mathrm{Cl}(0,3)$ coefficient embedding and $2I\supset2T$ bank.
No physical state or new parameter is supplied by the certificate.

\begin{proposition}[Independent finite-coordinate replay and checks]
\label{p1c:prop:v370-e8-finite-certificate}
Use the 120 root coordinates
\[
 \{\pm e_a\}_{a=1}^4\;\sqcup\;
 \{(\pm\tfrac12,\pm\tfrac12,\pm\tfrac12,\pm\tfrac12)\}
 \;\sqcup\;
 \left\{\text{even permutations of }
 (0,\pm\tfrac12,\pm\tfrac\varphi2,
 \pm\tfrac{\varphi^{-1}}2)\right\},
\]
with the standard 96 distinct golden vertices.  Represent each coordinate
as $(a+b\varphi)/2$ with integral $(a,b)$; multiplication uses
$\varphi^2=\varphi+1$.  The following independent exact checks all hold:
\begin{enumerate}
\item The coefficient lift has 240 roots and $D_0\perp D_1$;
  reflection orbit sizes are $(64,64,64)$, with $8\times8$ factorization.
\item The unordered zero-root-sum triple census is
  $2240=32+32+6(192)+1024$.
\item The ordinary-sheet classes have sizes $(24,32,32,32)$,
  each of the four outer $2T$-cosets has intersection $(0,8,8,8)$,
  and every nearest-neighbor row realizes the four-by-four quotient
  of Theorem~\ref{p1:thm:v370-e8-shell-chamber}.
\item On all 60 antipodal classes the central $A_5$ character kernels
  have $\mathscr P_5\psi_\zeta=\psi_\zeta$ and
  $\mathscr P_4\psi_\zeta=0$; the $H$-action has cyclic image
  and kernel exactly $Q_8$.
\item The specific $E_6\perp A_2$ cut has
  $72+6+6(27)=240$ roots, and each 27-weight fiber has the
  $3+8+8+8$ orbits and 45 distinct zero-sum projected-weight
  triples, with support counts $1+4+4+4+32$.
\item Forgetting the golden sheet gives 1840 distinct triples on
  600-cell vertices; their 3360-pair shadow is 56-regular,
  has no nearest-neighbor edge, and both its binary adjacency
  and integer pair-multiplicity matrices lie in $\mathbb Q[A_{600}]$.
\end{enumerate}
\end{proposition}
\begin{proof}
The deterministic exact-integer checker
\path{Proofs/Paper00-LisiE8Compatibility-v1_2.py}
reconstructs the entire 120-root bank from its $8+16+96$ families;
constructs both golden sheets by the coefficient map; performs
integer reflection closure, quaternion coset products and exact
root-triple enumeration; forms the complete nearest-neighbor quotient
and antipodal character kernels; and checks the projected $E_6$
weight triples.  It additionally proves the last claim by multiplying
powers $I,A_{600},\ldots,A_{600}^8$ as integer matrices and checking
entrywise rational polynomial certificates for both pair matrices.
Every asserted identity is an integer equality or a rational equality
cleared of denominators.  The checker has no floating-point threshold
and makes no inference about an action-generated physical carrier.
\end{proof}
\begin{remark}[Certificate provenance and remaining physical obligation]
\label{p1c:rem:v370-e8-provenance}
The replay is a mathematical finite-geometry certificate, not a proof
of physical occurrence.  The Paper~I-B signed Spin(8)
triality phases, any physical $SU(3)_c$ source, and the Paper~VII
same-action pole/residue gates are \emph{not} replaced or redefined.
In particular, the 45-element cubic support does not determine the
signed $E_6$ invariant or a decay amplitude.
\end{remark}


\section{Exact Petersen--600-cell total-graph, tangent, and network certificates}
\label{p1c:sec:v371-petersen-600cell-certificate}

\paragraph{Coordinate object and graph edge rule.}
Work with the standard 120 unit icosian quaternions in $2I$ and the
exact field $\mathbb Q(\sqrt5)$. No floating-point closeness tests are
used. The nearest-edge rule for $q,r\in2I$ is
$\operatorname{Re}(q^{-1}r)=\varphi/2$ with
$\varphi=(1+\sqrt5)/2$, and therefore has 720 unordered edges.
Its real-part shell counts are
\[
1_{1}+12_{\varphi/2}+20_{1/2}+12_{(\varphi-1)/2}
 +30_0+12_{-(\varphi-1)/2}+20_{-1/2}
 +12_{-\varphi/2}+1_{-1}=120.
\]
The real-$1/2$ shell has 30 internal edges; the real-zero shell
has 60; precisely 60 additional nearest edges join these shells.
Their induced graph is the $50$-vertex total graph of the
dodecahedron. In particular, the ten-vertex Petersen graph is a
\emph{quotient}, not a nearest-edge Petersen subgraph. Exhaustive
search over all $720$ starting $600$-cell edges and $147{,}600$
5-cycle candidates finds no Petersen subgraph of the ordinary
$600$-cell $1$-skeleton, even without imposing inducedness.

\paragraph{Quaternionic frames.}
For a dodecahedral edge $q=1/2+w$, $r=1/2+z$, use
$m=(w+z)/\varphi$, $d=\varphi(w-z)$, $k=2w\times z$.
All three belong to $2I$, are imaginary orthonormal unit vectors,
and obey $md=-k$. They generate exactly five distinct $Q_8$
subgroups across the 15 projective Petersen channels, each occurring
for three edges. The $150$ axis--channel pairings split into the
four angular orbits
\[
\begin{array}{c|rrrr}
 &\text{incident}&\text{near}&\text{cross B}&\text{cross A}\\\hline
\text{count}&30&60&30&30\\
|u\cdot m|&\varphi/\sqrt3&1/\sqrt3&(\varphi-1)/\sqrt3&0\\
\operatorname{dist}_{61}(u,m)&1&2&2&3
\end{array}
\]
where $\operatorname{dist}_{61}$ is shortest-path distance in the
complete inversion-orbit quotient, not in a truncated subgraph.

\paragraph{Finite exact tensor and tangent verification.}
Fix the ten normalized projective axes $u_{ij}$ from Paper~I-B and
the exact mirror matrices $J_i=(3/4)U_i^\sigma U_i^\dagger$.
Construct the fifteen outgoing trimer columns
$F_{u,e}=k_{u,e}$ at their two endpoint fibres. Exact multiplication
in $\mathbb Q(\sqrt5)$ verifies
\[
 F^\dagger F=3I-B^\dagger B/2,
 \qquad (J_i+J_j)k_e=(J_k+J_l)k_e
 \quad(ij\sim kl),
\]
\[
 MM^\dagger=\operatorname{diag}(20I_3/3,5I_{3'}/3).
\]
These identities give the exact sequence in
Theorem~\ref{p1:thm:v371-petersen-tangent-exact-sequence}.
With $X=B^\dagger B$ the polynomial
$-(X-6I)(13X^2-2X+25I)/450$ satisfies all four Penrose identities
for $G=F^\dagger F$; its reconstruction projector is
$I-3M_3^\dagger M_3/20-3M_{3'}^\dagger M_{3'}/5$.
For every core star the same exact finite arithmetic gives
\[
 \sum_{u\in H_m}P_u={8\over3}I_3,
 \quad\sum_{u\in H_m}T_uP_uT_u^\dagger={2\over3}I_{3'},
 \quad\sum_{u\in H_m}P_uT_u^\dagger={1\over3}J_m^\dagger.
\]
These are the source of the exact rank-six leakage spectrum and the
five-root sum rules in
Theorem~\ref{p1:thm:v371-petersen-rooted-leakage}.
The exact reference certificate is the standalone standard-library
script \texttt{Paper00-Petersen600CellIncidence-v1\_7.py}, evaluated
on the fixed source coordinates; it introduces no experimental inputs.

\paragraph{Optional scalar resistor comparison (not a quantum amplitude).}
Give every $600$-cell edge unit conductance and quotient by
$q\mapsto q^{-1}$, dividing by the uniform lift multiplicity two.
The resulting graph has 61 nodes and 360 edges. The ten $C_6$
axes and fifteen $C_4$ midpoint axes are 25 terminals; eliminating
the other 36 nodes by the graph-Laplacian Dirichlet-to-Neumann
Schur complement gives a nonnegative $10\times15$ cross block $C$.
Its exact entries are
\[
\begin{array}{c|cccc}
\text{pair type}&\text{incident}&\text{near}&\text{cross B}&\text{cross A}\\\hline
 C_{u,e}&352396/289425&133/1135&33688/289425&5086/289425
\end{array}
\]
with multiplicities $30,60,30,30$. With the standard Petersen ten-axis
projectors $P_0,P_4,P_5$,
\[
 CC^\dagger={777600\over51529}P_0+P_4
       +{961776\over180625}P_5.
\]
Eliminating the 15 midpoint terminals too gives
\[
 L_{10}^{\rm full}={2520\over181}P_4+{25200\over2551}P_5,
 \quad R_{\rm adjacent}={2057\over12600},\quad
 R_{\rm nonadjacent}={32\over175}.
\]
The two crossed types have exact resistance split $1/450$.
These quantities belong to the \emph{unoriented scalar circuit}
comparison, while the canonical Clifford fluxes form an
\emph{oriented} $15$-edge representation; no identification of their
norms, physical couplings, poles or decay rates is made.



\section{Independent rank, geometric, and analytic comparison calculations}
\label{p1d:app:independent-comparisons}
\noindent
The following calculations concern separately specified rank, gauge,
finite-character, complex-structure, spectral, quotient and Hessian
models. Each conclusion has precisely the hypotheses and limitations
stated in its paragraph. No implication for an unrestricted global
classification or physical occurrence follows from these comparisons.

\paragraph{Stable desuspension and fixed-rank scope.}
In the established stable Chern lattice, write
$\Delta=BC-AD$, $c_3=2m\,u_xu_p$ and
$c_4=(\Delta+2nm-6s)\mathcal O$.
If $6\mid(\Delta+2nm)$ choose
$s=(\Delta+2nm)/6$ to obtain the stable rank-four representative with
$c_4=0$. The only obstruction to a nowhere-zero section of an oriented
rank-four complex bundle over an eight-complex is its Euler/top Chern class:
the sphere fiber $S^7$ is $6$-connected and the primary obstruction lies
in $H^8(-;\mathbb Z)$. Splitting the resulting trivial line gives a
rank-three bundle with the required $c_1,c_2,c_3$. Necessity follows from
$c_4(E_3\oplus\underline{\mathbb C})=0$ and the same stable lattice. Thus
\[
 E_3\ \text{exists}\quad\Longleftrightarrow\quad
 6\mid(\Delta+2nm),
\]
\emph{provided} the asserted stable representative exists in the original
stable classification. The independent $\mathbb Z/2$ unstable $SU(2)$ class
on $S^3\times S^3$ is separately detected by
$[S^3\times S^3,BSU(2)]\simeq\pi_6(S^4)\simeq\mathbb Z/2$.
This route does not certify the rank-two existence criterion or exact
fixed-Chern rank-three fiber sizes.

\paragraph{Affine mapping components and determinant forgetting.}
For the prescribed two-sphere-factor mapping component model, composition
has central coordinates
\[
 (a,b;z)(a',b';z')=(a+a',b+b';z+z'+ab'),
 \qquad z\in\mathbb Z/12.
\]
The secondary $(\mathbb Z/2)^3$ module transforms by
$(r,s,e)\mapsto(r,s,e+br+as)$ in the stated model. At zero Chern data
this yields five secondary orbits: two with $(r,s)=(0,0)$ and one for
each nonzero $(r,s)$. The $12^2$ primary choices give $720$
determinant-preserving classes per specified restriction sector.
The independent finite-minor quotient script agrees on 2,401 tested
integer matrices. Determinant forgetting need not preserve this count:
for $Y=S^1_a\times S^1_b\times S^3_u$ and $c_2=a u$, winding
$h=b$ changes the secondary class by
$\bar h\,\bar c_2=\bar b\,\bar a\,\bar u\ne0$.
The complete eight-dimensional affine action and all ordinary/Born orbit
counts are not certified by these restricted checks.

\paragraph{Coherent exchange lifts and determinant restriction.}
The involutive smash sphere is $S^3\wedge S^3\simeq S^{3+3\sigma}$.
Suspending the conjugation-equivariant Hopf map yields
$S^{3+3\sigma}\to S^{3+2\sigma}\simeq S(\mathbb C^3)$,
with target involution $\mathrm{diag}(1,1,-1)$.
Pulling back the orthogonal-complement bundle supplies a coherent
\emph{complex} exchange lift of the nontrivial unstable class.
At the fixed locus the two eigenlines have signs $+1$ and $-1$, so the
equivariant determinant is nontrivial; this lift is not principal
$SU(2)$. For a determinant-preserving principal lift the fixed
isotropy must be $\pm I$, with the trivialized case subject to descent
to the symmetric-square quotient. In the stated index convention,
\[
 m_\pm={B^2-A^2\over12}\pm{A+B\over4},\qquad
 24\mid B^2-A^2+3(A+B)
\]
is the stronger principal-$SU(2)$ necessary arithmetic condition.
Exhaustive fixed-class enumeration and equivariant gauge orbit fibers
require the full eight-dimensional affine determinant action.

\paragraph{Relative mixed gluing on the prescribed boundary chart.}
For the declared mixed $(1,1)$ fixed eigenranks, the diagonal restriction
has $c_1=0$ and
$c_2=A(a_t-a_E)(u_x-u_p)$. Removing a tubular neighborhood of the
fixed locus produces $W$ with boundary; by Poincar\'e--Lefschetz
duality, $H^8(W;G)=0$ for the relevant coefficient systems.
In the adopted \emph{conditional} boundary-coordinate chart,
$T=\mathbb Z/6\oplus(\mathbb Z/2)^2$, $|T|=24$,
with integer image $2\mathbb Z$. Quotient arithmetic gives sizes
$24,12,48,12$ for the four respectively even/trivial,
even/exotic, odd/trivial and odd/exotic cases, and totals $72$ and
$120$ when the two compatible free-region classes are included.
These counts remain conditional on the \emph{actual global} gauge
extensions and the lower-degree equivariant Thom/Cartan--Leray data.
Vanishing top absolute obstruction cannot replace those proofs.

\paragraph{Finite icosahedral character and overlap identities.}
Fix the source's two $2I$ embeddings related by
$g_\theta=\cos\theta+\mathbf i\sin\theta$ and let
$w=\cos^2\theta\sin^2\theta$. Direct averaging of the finite
binary-icosahedral quaternion-coordinate character moments gives
\[
 F_{12}(w)=1-28w+196w^2-420w^3,
\]
\[
 \begin{aligned}
 F_{20}(w)={}&1-\frac{220}{3}w+\frac{4532}{3}w^2
 -\frac{36212}{3}w^3\\
 &+\frac{119680}{3}w^4-44880w^5.
 \end{aligned}
\]
The only positive-degree invariant lines through degree twenty occur
in degrees twelve and twenty, with multiplicity one. Since
$F_{12}(0)=1$, $F_{12}(1/4)=-5/16$ and
$\gcd_{\mathbb Q[w]}(F_{12},F_{20})=1$, the selected zero of $F_{12}$
annihilates Fourier modes through degree nineteen but not degree
twenty. In the twin $2T$-core completion model the overlap is
$-5/16$ and the independent projector eigenvalues are
\[
 a={275\over8192},\qquad b={525\over8192},\qquad
 \|\mathrm{err}_{19}(n)\|=\sqrt{13(a^{2n}+b^{2n})}.
\]
The group-generated density condition gives convergence to Haar on
continuous scalar tests and strong $L^2$ convergence of averaging
operators, not any physical measure selection. Full adjudication
must explicitly recheck the source embedding conventions.

\paragraph{Global complex structures and Dolbeault contraction.}
For a nonzero vertical Fourier character $m\in\mathbb Z^4$, the
Dolbeault differential has vertical symbol $\lambda(m)$; the map
$m\mapsto\lambda(m)$ is injective and bounded below by $c|m|$.
The homotopy
\[
 h_m={\sum_j\overline{\lambda_j(m)}\iota_{\bar Z_j}
 \over\sum_j|\lambda_j(m)|^2},\qquad
 \bar\partial h_m+h_m\bar\partial=1
\]
contracts each nonzero vertical Fourier component and preserves smooth
Fourier convergence. Only the torus-invariant complex survives.
The independent global normal-form route specializes the known
compact-Lie-group classification to
$\mathfrak{su}(2)^2\oplus\mathfrak u(1)^2$, with two root signs and a
complex structure on the real four-dimensional Cartan--central plane.
This does not by itself reprove the external classification or finish
the source's full $H$-twisted generalized-K\"ahler bracket check.

\paragraph{Pluriclosed averaging and canonical descent.}
The right four-torus acts holomorphically throughout the stated
normal-form family: it rotates the two roots and fixes the vertical
complex plane. Averaging an arbitrary pluriclosed Hermitian form under
left translations and this compact right torus preserves both
positivity and $\partial\bar\partial\omega=0$. In rank two,
$\mathrm du_1=x$, $\mathrm du_2=y$ gives
$\partial\bar\partial\omega=i\operatorname{Re}(h_{1\bar2})xy$;
positive diagonal $H$ provides a solution. In rank one,
$\mathrm du=x+\tau y$, $\mathrm dv=0$, $\operatorname{Im}\tau\ne0$,
\[
 \partial\bar\partial\omega
 =i h_{1\bar1}\operatorname{Re}\tau\,xy,
\]
so positivity implies the sharp condition $\operatorname{Re}\tau=0$.
The two holomorphic vertical translation fields globally trivialize
$T^{1,0}_{M_8/\mathbb P^1\times\mathbb P^1}$, and the relative tangent
sequence yields
\[
 K_{M_8}\cong\pi^*\mathcal O(-2,-2),\qquad
 H^0(M_8,K_{M_8}^{\otimes m})=0\quad(m\ge1)
\]
by connected compact-fiber descent. This does not promote the
holomorphic model to an independently selected physical background.

\paragraph{Homogeneous coefficient realization and positive sampling.}
Put $G=SU(2)_L\times SU(2)_R$ and
$K=SU(2)_{\mathrm{diag}}$. Each $G$-type
$U_{a,b}=V_{a/2}\otimes V_{b/2}$ has a unique top $K$-type
$V_s$, $s=(a+b)/2$. For normalized equivariant coisometry
$T_{a,b}:U_{a,b}\to V_s$ define
\[
 F_v(g)=\sqrt{\frac{(a+1)(b+1)}{a+b+1}}\,
 T_{a,b}\pi_{a,b}(g^{-1})v.
\]
This is a well-defined section of $G\times_KV_s$, and Schur
orthogonality gives
\[
 \int_G\pi(g)T^*T\pi(g)^*\,dg
 ={a+b+1\over(a+1)(b+1)}\,I,
 \qquad \|F_v\|_{L^2(G/K)}=\|v\|.
\]
For the declared family $H_m=\bigoplus_{a+b\le2m}U_{a,b}$,
sections in the global chart $(q,1)$ have coordinate degree at most
$2m$; their pairings have degree at most $4m$. Positive spherical
cubature of strength $t\ge4m$ gives isometric evaluation, and an
$r$-linear coefficient expression of degree $d$ is exact when
$t\ge2rm+d$. More generally, if the specified real coefficient
test space (including constants) has dimension $D$, the integration
functional has a positive atomic representation with at most $D$
points by the finite-dimensional convex-hull argument. Only
operators \emph{included in that coefficient test space} are thereby
transported; no claim about untransported pointwise calculus follows.

\paragraph{TT and Rarita--Schwinger spectral comparisons.}
The trace-free symmetric-two-tensor fiber rank is $35$; the eight
independent divergence constraints, adjusted by the conformal-Killing
kernel, reproduce the Clebsch--Gordan TT multiplicities. The confirmed
product TT heat coefficients include
\[
 A_0^{TT}=27V,\qquad A_2^{TT}=-27V(a+b),\qquad
 A_4^{TT}=V\bigl[\tfrac{129}{2}(a^2+b^2)-45ab\bigr].
\]
On a separate explicit constant-orbital Dirac calculation,
the ambient tangent-spinor block has dimension $128$ and its
$\gamma$-traceless physical subspace has dimension $112$. On that
subspace the matrix identity
\[
 Q^2+\tfrac12PP^*=\Pi_\gamma D_T^2\Pi_\gamma
\]
holds in the tested convention. The apparent ambient $16$-dimensional
$Q$ kernel is entirely the \emph{excluded gamma-trace complement}:
there are \emph{no genuine $Q$ zeros in this tested constant-orbital
block}. This is not an all-orbital constrained RS spectral proof;
transversality and the full product zeta/heat family remain open.

\paragraph{Finite quotient kernels and degree-twelve branching.}
For a free isometric finite action the quotient heat kernel relative
to Riemannian volume is the image sum; for probability-normalized
measures it is the corresponding finite average. For
$q\mapsto aqb^{-1}$ a fixed point requires conjugate $a,b\in SU(2)$,
equivalently $\operatorname{Re}a=\operatorname{Re}b$; a simultaneous
affine torus action requires solving its own congruences.
The independent degree-twelve rational pencil has
\[
 f_6=z(z^4-1),\quad f_8=z^8+14z^4+1,\quad
 g_{12}=z^{12}-33z^8-33z^4+1,\quad u_{12}=f_6^2,
\]
with homogeneous lifts implicit. Symbolic identities
\[
 f_8^3=g_{12}^2+108u_{12}^2,\qquad
 u_{12}'g_{12}-u_{12}g_{12}'=-2f_6f_8^2
\]
give a basepoint-free degree-twelve pencil with six index-two and
eight index-three ramification points. Riemann--Hurwitz gives
$6+16=22=2(12)-2$; the product degree is $144$, ramification
class $22(h_x+h_p)$ and squared pulled-back anticanonical class
$288$. Exhaustive nonfree/correlated stabilizer strata, and their
orbifold kernels, remain a distinct independent-proof obligation.

\paragraph{Midpoint Yang--Mills Hessian and gauge quotient.}
Use $[e_a,e_b]=(2/R)\epsilon_{abc}e_c$,
$[I_a,I_b]=\epsilon_{abc}I_c$ and
$A_B(e_a)=I_a/R$. The midpoint curvature is
$F_B(e_a,e_b)=-\epsilon_{abc}I_c/R^2$ in the stated convention.
The gauge-fixed Weitzenb\"ock Hessian and ghost are
\[
 \begin{aligned}
 R^2H_{\rm gf}&=(2L+S+T)^2+2+2S\cdot T
 =2J(J+1)+2j(j+1)-2,\\
 R^2\Delta_{\rm gh}&=(2L+T)^2
 =2k(k+1)+2j(j+1)-2.
 \end{aligned}
\]
Independent $j=0,\tfrac12,1$ spin matrices reproduce the printed
vector and ghost eigenvalues. At a Yang--Mills critical point the
unfixed Hessian annihilates infinitesimal gauge directions; thus
the gauge-fixed Hessian restricted to their image agrees with
the positive scalar ghost, so it is legitimate to remove these
summands. The $j=0$ quotient retains the unique negative singlet
$-2/R^2$ and positive quintet $10/R^2$; one $j=\tfrac12$ copy
survives at $R^{-2}$; and for $j\ge1$,
$R^2\lambda\ge2j(j+1)-2\ge2$. Hence the classical spatial
Morse index is one, there is no quotient zero mode, and the least
positive eigenvalue is $R^{-2}$. This does not imply a
Lorentzian transition or a populated decay history.



\section{Source-relative rank-two and natural-boundary propositions}
\label{sec:p1d-ic19-closures}

\noindent\textbf{Hypotheses and scope.}
The following propositions use the coherent G15 mixed $U(2)$ exchange
lifts, the G20--G21 mapping-torus and relative Thom-space constructions,
and the RT/ND/PI/SC/CI/FS/TN/BG maps specified above. Their conclusions
are relative to those geometric source data and their stated frames and
normalizations. They neither modify G0/GP1--GP3 nor identify a gauge,
bundle or transition with an occurring physical process. Finite algebraic
checks do not independently reprove an inherited theorem.

\begin{proposition}[Fixed-stable-class rank-three realization]
\label{p1d:prop:ic19-stable-rank-three}
For the fixed stable complex $K$-class on the actual eight-dimensional
CW model $M_8$, with its prescribed Chern data, the rank-three
representability obstruction equals its fourth Chern class.  In the
source's integer convention, rank-three representability is therefore
\[
 \boxed{6\mid(BC-AD+2nm).}
\]
This statement concerns existence in the fixed stable class and does
not determine rank-two equivariant lifts or unstable fixed-Chern fibre
multiplicities.
\end{proposition}
\begin{proof}
The stabilisation $BU(4)\to BU$ induces a lift of the prescribed
stable class over an eight-complex.  A complex rank-four bundle admits
a nowhere-zero section precisely when its degree-eight Euler/top
Chern obstruction vanishes; its unit-sphere fibre is $S^7$, so the
obstruction lies in $H^8(M_8;\mathbb Z)$ and there are no higher
obstructions in this dimension.  Such a section splits off a trivial
line, producing a rank-three representative.  Conversely, a rank-three
representative has fourth Chern class zero.  Inserting the source's
stable Chern-coordinate formula gives the displayed divisibility.
\end{proof}

\begin{proposition}[The actual compatible free-region seeds]
\label{p1d:prop:ic19-compatible-seeds}
Work with the source's coherent mixed $U(2)$ G15 restriction sectors
$\varepsilon\in\{0,1\}$, and the actual G20/G21 free-region pushout
$W=Y\cup_J Z_s$.  For each $A\in\mathbb Z$ there is a compatible
seed in both sectors, with
\[
 c_1(V_{A,\varepsilon})=0,\qquad
 c_2(V_{A,\varepsilon}|_R)=(A\bmod2)\tau^2,\qquad
 V_{A,\varepsilon}|_R\cong\mathbf1\oplus\lambda_L.
\]
The supported paired-disc fifth-stage modification gives the second
class over each seed, and the actual two-piece joining census gives
\emph{four} boundary-compatible ordinary free-region bundle classes
at each fixed integer $A$.  This count is a free-region bundle count,
not yet the quotient by every global boundary gauge.
\end{proposition}
\begin{proof}
Write $P=S_x^3\times S_p^3$ and let $\Phi_\varepsilon$ be the
source's coherent exchange map on the actual sector bundle.  Set
$\alpha(f)=\Phi_\varepsilon B^*f\Phi_\varepsilon^{-1}$, so
$\alpha^2=1$.  An actual determinant-one gauge $f_A$ of primary
$(A,0)$ is obtained from the wedge $S^3\vee S^3$ and extended over
the six-cell using the source's vanishing $\pi_5(SU(2))$ projection-
degree obstruction.  The strict coboundary
\[
 g_A=f_A\alpha(f_A)^{-1}
 \quad\text{satisfies}\quad \alpha(g_A)=g_A^{-1}
\]
\emph{as an actual gauge map}.  It therefore defines a reflected
circle clutch with an order-two exchange action.  On the fixed
diagonal its clutching circle is constant and retains the mixed
$(1,1)$ fibre involution.  The Borel degree-four filtration has only
the possible positive-filtration class $\tau^2$; fixed-locus
isotropy determines its coefficient to be $A\bmod2$.  The source's
paired-disc supported $H^5$ modification leaves all these boundary
and Chern data invariant.  Finally, the two $Y$-piece and two $Z_s$-
piece classes have one actual $J$ joining orbit for each piece pair.
The disconnected lens gauge component of degree $k$ shifts the
\emph{global} coefficient $A\mapsto A+2k$ and hence is not an
unaccounted infinite family at fixed $A$.
\end{proof}

\begin{proposition}[First unstable determinant action, with its domain]
\label{p1d:prop:ic19-first-postnikov-action}
For an admissible rank-two class on $M_8$ with the source's Chern
coordinates $(n,A,B,C,D)$, let determinant winding be
$h=ra_t+sa_E$ and the primary degree-three loop be
$v=v_xu_x+v_pu_p$.  The induced action on the \emph{first}
unstable degree-five ambiguity is
\[
 \boxed{\delta(h,v)=
 \begin{pmatrix}nv_x+rC-sA\\nv_p+rD-sB\end{pmatrix}\pmod2.}
\]
On the five-dimensional retract $S^1\times S^1\times S^3$ this
first-stage computation is exhaustive because all higher obstructions
vanish there.  It does \emph{not} determine the complete
higher-secondary affine determinant action on $M_8$.
\end{proposition}
\begin{proof}
Apply the source's first $BU(2)$ Postnikov invariant to the loop
slant of $c_1,c_2$; the $Sq^2$ term and multiplication by $c_1$
produce the two mod-two coordinates displayed.  On the retract
there are no cells in the higher obstruction dimensions, so the
resulting affine orbit calculation is complete.  The eight-dimensional
secondary and top-stage action is explicitly not inferred from it.
\end{proof}

\begin{proposition}[Unpointed lens evaluation and the torsion-primary gauge]
\label{p1d:prop:ic19-lens-torsion}
For $E_L=\mathbf1\oplus\lambda_L$ pulled back through the actual
Hopf-line map $\mathbb{RP}^3\to S^2$, the source's pointed gauge
homotopy-fibre sequence has evaluation maps
\[
 \delta_3:\mathbb Z\to C_2,\ k\mapsto k\bmod2,
 \qquad\delta_4:C_2\hookrightarrow C_2\oplus C_6,
 \qquad\delta_5:C_2\xrightarrow{\,0\,} C_2^2,
\]
where the last equality means $\delta_5$ is the zero map.  Consequently
$\pi_2H=0$, $\pi_3H\cong\mathbb Z\oplus C_6$,
$\pi_4H\cong C_2^2$, and
$\ker(j_6:C_{12}\to C_6)=\{0,6\}$.
For the specified G20/G21 joining and product-framed relative surgery,
the constructed torsion-primary global gauge $g_\eta$ has odd-sector
normal boundary value $j_\eta=t$; the result is unchanged by the
admissible fourth-stage comparison shear.
\end{proposition}
\begin{proof}
The evaluation maps follow by naturality from the Hopf-line gauge
evaluation fibration and the source's pointed homotopy-fibre model;
the exact sequence gives the displayed unpointed groups.  On the
$Z_s$ piece $\eta_Z=\beta_{\mathbb Z}(r_Z)$ has mod-two reduction
$\chi r_Z$ and
$(Sq^2+\chi^2)(\chi r_Z)=0$, supplying an actual torsion-primary
gauge.  Its joining difference is fourth-stage and is removed by
the source's supported toggle on $Y$.  In the product-framed surgery
the bottom comparison is null since $\pi_2H=0$; the nonzero raw top
coefficient $6\nu'\in C_{12}$ maps to zero in $C_6$.  The coherent
nulls have fourth-stage correction $\beta_4=u\chi+\lambda x\chi^3$
with $\lambda\in\mathbb F_2$.  On the boundary its commutator with
the torsion-primary reduction $x\chi^2$ is
\[
 (u\chi+\lambda x\chi^3)(x\chi^2)=xu\chi^3\ne0
 \quad\text{in }H^7(R;\mathbb Z_2).
\]
The extra fifth-stage correction first contributes in dimension
eight; under the stated product frame this is precisely the top
boundary generator $t$.
\end{proof}

\begin{proposition}[Terminal relative obstruction is a parity switch]
\label{p1d:prop:ic19-top-transgression}
For the source's connected, oriented free-region pair $(W,R)$,
with the specified product framing,
\[
 \delta_7:H^7(R;\mathbb Z_2)\xrightarrow{\cong}
 H^8(W,R;\mathbb Z_2)=\mathbb Z_2\{\Theta\},
 \qquad O_8(gt)=O_8(g)+\Theta.
\]
Thus, after fixing the lower-stage ambiguity, exactly one member of
each $t$-paired candidate set can extend.  The aggregate natural
\emph{number} of extensions alone does not fix $(h_t,w_t)$.
\end{proposition}
\begin{proof}
Use the long exact sequence of the pair, the actual source
cohomology of $W$, and Poincar\'e--Lefschetz duality for its
boundary orientation.  The boundary top-cell generator $t$ is the
unit transgression.  Hence multiplication by it changes the terminal
obstruction by $\Theta$.  The natural count $4$ out of $14$ is
unchanged under all four possible assignments of $(h_t,w_t)$;
which \emph{named} candidates extend is the separating datum.
\end{proof}

\begin{theorem}[Original natural determinant coefficient: source-relative full route]
\label{p1d:thm:ic19-natural-determinant-route}
Assume the actual established RT.53.4 matched-choice theorem and odd
pure-image property, the G20/G21 source's coherent exchange and
relative Thom models, the completed source-specific root-cap and
projective-null frame paths (NY--NZ and OB--OL), and the original
product framing of the native two-disc comparison.  Then the
\emph{actual} ND/PI/SC/CI determinant gauge, without a replacement
representative, has
\[
 \boxed{h_t\mu=[Z(c)]=\nu'\eta_6=\mu,\qquad
 (h_t,w_t,r_J,\zeta_{\rm nat})=(1,1,0,0).}
\]
This is a theorem on the declared geometric source data and those
named upstream premises.  It does not establish an unconditional
classification of all rank-two gauge components, nor a physical
occurrence or decay selection.
\end{theorem}
\begin{proof}
The actual first-cut PI conjugation is matched and is related to the
original determinant triple by RT.53.4, after preserving its clutch
frame.  At the second cut retain the \emph{whole} difference word
\[
 G=R_+^{-1}dR_-d,\qquad
 M_{\rm new}M_0^{-1}=P(t,G)[G,M_0],\qquad
 P(t,G)=q_t^{-1}Gq_tG^{-1}.
\]
The moving lower-value correction cannot be dropped.  The full phase
word reduces by support separation and the annihilation
$\eta_3E\nu'=0$ in $\pi_7(S^3)=C_2$.  For the conjugation word, the
actual exterior $C$ transports through ND.7 to the source's $C_2$
with degree $-2$ in its oriented chart, while the supported gauge has
integer degree $N$.  The relative cup product therefore gives
$\deg T=\pm2N$, whose universal seven-sphere class vanishes in
$\pi_7(SU(2))=C_2$.  All crossed base-value corrections admit based
support separation.  The complete second cut is null, not just its
phase factor.

The coherent CI joining frame stays uniformly nonsingular:
for its aligned unit vectors the real overlap is at least $1/4$ and
\[
 \big\|(1-\lambda)\widehat v+\lambda\gamma b\big\|^2
 \ge1-\tfrac32\lambda(1-\lambda)\ge\tfrac58.
\]
The deck-invariant phase has a continuous real logarithm on
$\mathbb{RP}^3$.  The old and translated circle caps are related by
the based relative disk homotopy and the original root-cap frame
bridge; the full southern meridian and both inverse $K$ factors of
the JT seam remain present.  This establishes the typed identification
\[
 [\varphi_D(\mathcal T_{ND})]=[\varphi_D(\mathcal T_{PI})]
 =[Z(c)]\quad\text{in the normalized }\pi_7(SU(2)).
\]
The actual NZ cap-clutch factors through
$\eta_3\eta_4\ne0\in\pi_5(S^3)$; its extendible frame factors are
null on the specified product supports.  The moving-cap reference
has zero complete paired class, so whole-map support splitting
reduces $[Z(c)]$ to the source's reduced neck.  Its two closed normal
bigons have odd combined axis degree, yielding the actual
Samelson--Hopf composition $[Z(c)]=\nu'\eta_6=\mu$.
The separately normalized sphere path supplies $w_t=1$; the original
joining and projective nulls give $r_J=\zeta_{\rm nat}=0$.
No step infers these coefficients from the parity-invariant aggregate
$4$-out-of-$14$ extension count.
\end{proof}

\begin{remark}[Scope of the geometric certificates and unresolved extensions]
The marked $2I$ character and cubature calculation, the stated
Hermitian-family pluriclosed and canonical calculation, the declared
finite coefficient-space homogeneous sampling, and the classical
\emph{spatial} midpoint Yang--Mills Hessian/gauge quotient have the
normalizations and selection boundaries stated in their respective
constructions. The fixed-stable-class rank-three existence criterion
does not classify rank-two lifts. The principal $SU(2)$ exchange-lift
obstruction is retained. Compatible free-region seeds, the actual
low-gauge construction, the unpointed lens evaluation, and the
specified natural gauge comparison do not give an all-components
global classification. The complete eight-dimensional secondary
affine determinant action remains undetermined. The unrestricted
complex-structure, constrained all-orbital operator, and nonfree
quotient conclusions are not supplied by their restricted comparison
calculations. The numerical directed-rounding certificate and the
finite 490-simplex, F40/Hurwitz/$E_8$ executable artifacts remain
separate reproducibility requirements; a small-cap inverse-function
homotopy is not a replacement for those numerical certificates.
No physical constructor, occurrence, particle mass or decay rate
is asserted.
\end{remark}

\begingroup
\small
\printbibliography
\endgroup
\end{document}
